% Options for packages loaded elsewhere
\PassOptionsToPackage{unicode}{hyperref}
\PassOptionsToPackage{hyphens}{url}
\documentclass[
  english,
]{book}
\usepackage{xcolor}
\usepackage{amsmath,amssymb}
\setcounter{secnumdepth}{5}
\usepackage{iftex}
\ifPDFTeX
  \usepackage[T1]{fontenc}
  \usepackage[utf8]{inputenc}
  \usepackage{textcomp} % provide euro and other symbols
\else % if luatex or xetex
  \usepackage{unicode-math} % this also loads fontspec
  \defaultfontfeatures{Scale=MatchLowercase}
  \defaultfontfeatures[\rmfamily]{Ligatures=TeX,Scale=1}
\fi
\usepackage{lmodern}
\ifPDFTeX\else
  % xetex/luatex font selection
\fi
% Use upquote if available, for straight quotes in verbatim environments
\IfFileExists{upquote.sty}{\usepackage{upquote}}{}
\IfFileExists{microtype.sty}{% use microtype if available
  \usepackage[]{microtype}
  \UseMicrotypeSet[protrusion]{basicmath} % disable protrusion for tt fonts
}{}
\makeatletter
\@ifundefined{KOMAClassName}{% if non-KOMA class
  \IfFileExists{parskip.sty}{%
    \usepackage{parskip}
  }{% else
    \setlength{\parindent}{0pt}
    \setlength{\parskip}{6pt plus 2pt minus 1pt}}
}{% if KOMA class
  \KOMAoptions{parskip=half}}
\makeatother
\usepackage{color}
\usepackage{fancyvrb}
\newcommand{\VerbBar}{|}
\newcommand{\VERB}{\Verb[commandchars=\\\{\}]}
\DefineVerbatimEnvironment{Highlighting}{Verbatim}{commandchars=\\\{\}}
% Add ',fontsize=\small' for more characters per line
\usepackage{framed}
\definecolor{shadecolor}{RGB}{248,248,248}
\newenvironment{Shaded}{\begin{snugshade}}{\end{snugshade}}
\newcommand{\AlertTok}[1]{\textcolor[rgb]{0.94,0.16,0.16}{#1}}
\newcommand{\AnnotationTok}[1]{\textcolor[rgb]{0.56,0.35,0.01}{\textbf{\textit{#1}}}}
\newcommand{\AttributeTok}[1]{\textcolor[rgb]{0.13,0.29,0.53}{#1}}
\newcommand{\BaseNTok}[1]{\textcolor[rgb]{0.00,0.00,0.81}{#1}}
\newcommand{\BuiltInTok}[1]{#1}
\newcommand{\CharTok}[1]{\textcolor[rgb]{0.31,0.60,0.02}{#1}}
\newcommand{\CommentTok}[1]{\textcolor[rgb]{0.56,0.35,0.01}{\textit{#1}}}
\newcommand{\CommentVarTok}[1]{\textcolor[rgb]{0.56,0.35,0.01}{\textbf{\textit{#1}}}}
\newcommand{\ConstantTok}[1]{\textcolor[rgb]{0.56,0.35,0.01}{#1}}
\newcommand{\ControlFlowTok}[1]{\textcolor[rgb]{0.13,0.29,0.53}{\textbf{#1}}}
\newcommand{\DataTypeTok}[1]{\textcolor[rgb]{0.13,0.29,0.53}{#1}}
\newcommand{\DecValTok}[1]{\textcolor[rgb]{0.00,0.00,0.81}{#1}}
\newcommand{\DocumentationTok}[1]{\textcolor[rgb]{0.56,0.35,0.01}{\textbf{\textit{#1}}}}
\newcommand{\ErrorTok}[1]{\textcolor[rgb]{0.64,0.00,0.00}{\textbf{#1}}}
\newcommand{\ExtensionTok}[1]{#1}
\newcommand{\FloatTok}[1]{\textcolor[rgb]{0.00,0.00,0.81}{#1}}
\newcommand{\FunctionTok}[1]{\textcolor[rgb]{0.13,0.29,0.53}{\textbf{#1}}}
\newcommand{\ImportTok}[1]{#1}
\newcommand{\InformationTok}[1]{\textcolor[rgb]{0.56,0.35,0.01}{\textbf{\textit{#1}}}}
\newcommand{\KeywordTok}[1]{\textcolor[rgb]{0.13,0.29,0.53}{\textbf{#1}}}
\newcommand{\NormalTok}[1]{#1}
\newcommand{\OperatorTok}[1]{\textcolor[rgb]{0.81,0.36,0.00}{\textbf{#1}}}
\newcommand{\OtherTok}[1]{\textcolor[rgb]{0.56,0.35,0.01}{#1}}
\newcommand{\PreprocessorTok}[1]{\textcolor[rgb]{0.56,0.35,0.01}{\textit{#1}}}
\newcommand{\RegionMarkerTok}[1]{#1}
\newcommand{\SpecialCharTok}[1]{\textcolor[rgb]{0.81,0.36,0.00}{\textbf{#1}}}
\newcommand{\SpecialStringTok}[1]{\textcolor[rgb]{0.31,0.60,0.02}{#1}}
\newcommand{\StringTok}[1]{\textcolor[rgb]{0.31,0.60,0.02}{#1}}
\newcommand{\VariableTok}[1]{\textcolor[rgb]{0.00,0.00,0.00}{#1}}
\newcommand{\VerbatimStringTok}[1]{\textcolor[rgb]{0.31,0.60,0.02}{#1}}
\newcommand{\WarningTok}[1]{\textcolor[rgb]{0.56,0.35,0.01}{\textbf{\textit{#1}}}}
\usepackage{longtable,booktabs,array}
\newcounter{none} % for unnumbered tables
\usepackage{calc} % for calculating minipage widths
% Correct order of tables after \paragraph or \subparagraph
\usepackage{etoolbox}
\makeatletter
\patchcmd\longtable{\par}{\if@noskipsec\mbox{}\fi\par}{}{}
\makeatother
% Allow footnotes in longtable head/foot
\IfFileExists{footnotehyper.sty}{\usepackage{footnotehyper}}{\usepackage{footnote}}
\makesavenoteenv{longtable}
\usepackage{graphicx}
\makeatletter
\newsavebox\pandoc@box
\newcommand*\pandocbounded[1]{% scales image to fit in text height/width
  \sbox\pandoc@box{#1}%
  \Gscale@div\@tempa{\textheight}{\dimexpr\ht\pandoc@box+\dp\pandoc@box\relax}%
  \Gscale@div\@tempb{\linewidth}{\wd\pandoc@box}%
  \ifdim\@tempb\p@<\@tempa\p@\let\@tempa\@tempb\fi% select the smaller of both
  \ifdim\@tempa\p@<\p@\scalebox{\@tempa}{\usebox\pandoc@box}%
  \else\usebox{\pandoc@box}%
  \fi%
}
% Set default figure placement to htbp
\def\fps@figure{htbp}
\makeatother
\ifLuaTeX
\usepackage[bidi=basic,shorthands=off]{babel}
\else
\usepackage[bidi=default,shorthands=off]{babel}
\fi
\ifLuaTeX
  \usepackage{selnolig} % disable illegal ligatures
\fi
\setlength{\emergencystretch}{3em} % prevent overfull lines
\providecommand{\tightlist}{%
  \setlength{\itemsep}{0pt}\setlength{\parskip}{0pt}}
\usepackage[]{natbib}
\bibliographystyle{apalike}
\usepackage{microtype}
\usepackage{booktabs}
\usepackage{longtable}
\usepackage{float}
\usepackage{xcolor}
\usepackage{hyperref}
% Permit long DOI/URL tokens to wrap without inserting visible hyphens.
\usepackage{xurl}
\urlstyle{same}
\definecolor{linkblue}{HTML}{1565C0}
\hypersetup{colorlinks=true,linkcolor=linkblue,citecolor=linkblue,urlcolor=linkblue}
\setcounter{secnumdepth}{4}
\setcounter{tocdepth}{3}
% The complete book has three-digit pages. Reserve room for up to four digits
% instead of the standard class's narrow page-number box and right margin.
\makeatletter
\renewcommand{\@pnumwidth}{3em}
\renewcommand{\@tocrmarg}{4em}
\makeatother
\usepackage{fontspec}
% Match the Unicode-capable PDF figure font selected in _common.R. These
% families preserve author names, Greek letters, and comparison operators.
\IfFontExistsTF{Noto Sans CJK KR}{
  \setmainfont{Noto Sans CJK KR}
  \setsansfont{Noto Sans CJK KR}
}{\IfFontExistsTF{Noto Sans CJK JP}{
  \setmainfont{Noto Sans CJK JP}
  \setsansfont{Noto Sans CJK JP}
}{\IfFontExistsTF{Arial Unicode MS}{
  \setmainfont{Arial Unicode MS}[AutoFakeBold=1.5,AutoFakeSlant=0.2]
  \setsansfont{Arial Unicode MS}[AutoFakeBold=1.5,AutoFakeSlant=0.2]
}{\PackageError{scientometrics}{Install Noto Sans CJK fonts before rendering PDF}{The PDF must preserve non-Latin author names.}}}}
\IfFontExistsTF{DejaVu Sans Mono}{
  \setmonofont{DejaVu Sans Mono}[Scale=MatchLowercase]
}{\IfFontExistsTF{Menlo}{
  \setmonofont{Menlo}[Scale=MatchLowercase]
}{\PackageError{scientometrics}{Install DejaVu Sans Mono before rendering PDF}{A Unicode monospace font is required for code examples.}}}

% Wrap both highlighted source and plain output within the text area.
\usepackage{fvextra}
\fvset{breaklines=true,breakanywhere=true}
% Apply after package loading: tcolorbox loads verbatim and can otherwise undo
% the lowercase override. Define also handles excerpts without source chunks,
% where Pandoc has not created a Highlighting environment in its template.
\AtBeginDocument{
  \DefineVerbatimEnvironment{Highlighting}{Verbatim}{commandchars=\\\{\},breaklines=true,breakanywhere=true,fontsize=\small}
  \RecustomVerbatimEnvironment{verbatim}{Verbatim}{breaklines=true,breakanywhere=true,fontsize=\small}
}

% --- Callout-box environments for {block, type='box...'} chunks ---------------
% The bs4_book HTML target maps these block types to <div class="...">.
% In PDF, pandoc converts them to \begin{type}...\end{type} environments,
% which must exist. Define them all as light grey-bordered framed boxes.
\usepackage{tcolorbox}
\tcbuselibrary{breakable, skins}
\definecolor{boxgrey}{HTML}{E0E0E0}
\definecolor{boxbg}{HTML}{F5F5F5}
\newtcolorbox{boxgeneric}{
  enhanced, breakable,
  colback=boxbg, colframe=linkblue,
  boxrule=0pt, leftrule=2pt,
  arc=2pt, outer arc=2pt,
  left=10pt, right=10pt, top=8pt, bottom=8pt
}
\newenvironment{boxinfo}{\begin{boxgeneric}}{\end{boxgeneric}}
\newenvironment{boximportant}{\begin{boxgeneric}}{\end{boxgeneric}}
\newenvironment{boxreport}{\begin{boxgeneric}}{\end{boxgeneric}}
\newenvironment{boxdmetar}{\begin{boxgeneric}}{\end{boxgeneric}}
\newenvironment{boxquestion}{\begin{boxgeneric}}{\end{boxgeneric}}
\newenvironment{boxempty}{\begin{boxgeneric}}{\end{boxgeneric}}

\usepackage{fancyhdr}
\setlength{\headheight}{25pt}
\pagestyle{fancy}
\fancyhead{}
\fancyhead[C]{\small\nouppercase{\leftmark}}
\fancyfoot{}
\fancyfoot[C]{\scriptsize R. Heller · Date: 2026-05-09 · Last changed: 2026-11-11\hfill\thepage}
\fancypagestyle{plain}{\fancyhf{}\renewcommand{\headrulewidth}{0pt}\fancyfoot[C]{\scriptsize R. Heller · Date: 2026-05-09 · Last changed: 2026-11-11\hfill\thepage}}

\usepackage{booktabs}
\usepackage{caption}
\usepackage{longtable}
\usepackage{colortbl}
\usepackage{array}
\usepackage{anyfontsize}
\usepackage{multirow}
\usepackage{bookmark}
\IfFileExists{xurl.sty}{\usepackage{xurl}}{} % add URL line breaks if available
\urlstyle{same}
\hypersetup{
  pdftitle={Scientometrics in R},
  pdfauthor={R. Heller},
  pdflang={en},
  hidelinks,
  pdfcreator={LaTeX via pandoc}}

\title{Scientometrics in R}
\usepackage{etoolbox}
\makeatletter
\providecommand{\subtitle}[1]{% add subtitle to \maketitle
  \apptocmd{\@title}{\par {\large #1 \par}}{}{}
}
\makeatother
\subtitle{Bibliometrics, Network Science, and Research Evaluation}
\author{R. Heller}
\date{2026-11-10}

\begin{document}
\maketitle

{
\setcounter{tocdepth}{2}
\tableofcontents
}
\chapter*{Welcome!}\label{welcome}
\addcontentsline{toc}{chapter}{Welcome!}

\emph{R. Heller · Date: 2026-05-11 · Last changed: 2026-11-10}

\begin{center}\rule{0.5\linewidth}{0.5pt}\end{center}

Welcome to \textbf{``Scientometrics in R: Bibliometrics, Network Science, and Research Evaluation''}.

This is a working guide to scientometrics in R for researchers, librarians, research administrators, and data scientists who want to study scholarly communication quantitatively. It walks the full pipeline --- from pulling records out of OpenAlex, Crossref, and PubMed, through computing the standard indicators (h-index, MNCS, journal metrics), to building co-authorship and co-citation networks, applying text and topic models to scientific corpora, and packaging the results into reproducible reports and dashboards.

The book assumes a working knowledge of R at the tidyverse level and basic statistics. No prior scientometric experience is needed --- the early chapters lay the conceptual and ethical foundations, and each later chapter starts with a concrete research question and ends with a callout on responsible use grounded in the Leiden Manifesto and DORA.

All examples run against free, open data; no Web of Science or Scopus subscription is required. Small sample datasets ship with the companion R package \texttt{scientometricsInR}; everything else is fetched live from public APIs. The book does \emph{not} cover proprietary platforms (InCites, SciVal), manual screening workflows for systematic reviews, or qualitative sociology of science.

\section*{Open Source Repository}\label{open-source-repository}
\addcontentsline{toc}{section}{Open Source Repository}

\begin{center}\rule{0.5\linewidth}{0.5pt}\end{center}

This book has been built using \href{https://rmarkdown.rstudio.com/docs/}{\textbf{\{rmarkdown\}}} and \href{https://bookdown.org/}{\textbf{\{bookdown\}}}. Formulas are rendered using \href{http://docs.mathjax.org/en/latest/index.html}{MathJax}. All source files are available on \textbf{GitHub} at \url{https://github.com/CTTIR/scientometrics-in-r}. You are free to fork, share, and reuse contents under the project's \href{LICENSE}{CC BY 4.0 license}.

\section*{How To Use The Guide}\label{how-to-use-the-guide}
\addcontentsline{toc}{section}{How To Use The Guide}

\begin{center}\rule{0.5\linewidth}{0.5pt}\end{center}

\begin{itemize}
\tightlist
\item
  \textbf{Linear or lookup.} Read front to back, or use the \hyperref[find-your-method]{Find Your Method navigator} to jump to the chapter that answers your question.
\item
  \textbf{Code.} Reusable helpers live in the companion R package \texttt{scientometricsInR} and in \texttt{R/} in the source repository; every chapter shows the full chunk it depends on.
\item
  \textbf{Downloads.} The side panel lists whole-book and chapter files that are available in the current build.
\item
  \textbf{Corrections.} Please file an issue at \url{https://github.com/CTTIR/scientometrics-in-r/issues}.
\end{itemize}

\section*{Contributing}\label{contributing}
\addcontentsline{toc}{section}{Contributing}

\begin{center}\rule{0.5\linewidth}{0.5pt}\end{center}

Pull requests, corrections, and additions are welcome. See \href{https://github.com/CTTIR/scientometrics-in-r/blob/main/CONTRIBUTING.md}{\texttt{CONTRIBUTING.md}} for the workflow.

\section*{Citing this Guide}\label{citing-this-guide}
\addcontentsline{toc}{section}{Citing this Guide}

\begin{center}\rule{0.5\linewidth}{0.5pt}\end{center}

The suggested citation is:

\begin{boxempty}
\textbf{Heller, R.} (2026). \emph{Scientometrics in R: Bibliometrics,
Network Science, and Research Evaluation} (Version 0.1.0).
Self-published via GitHub Pages.
\url{https://cttir.github.io/scientometrics-in-r/}.
\end{boxempty}

Download the reference as \href{citation-files/citation.bib}{BibTeX} or \href{citation-files/citation.ris}{.ris}.

\chapter*{Impressum}\label{impressum}
\addcontentsline{toc}{chapter}{Impressum}

\emph{R. Heller · Date: 2026-05-11 · Last changed: 2026-11-09}

\textbf{Scientometrics in R: Bibliometrics, Network Science, and Research Evaluation}

© 2026 Raban Heller. ORCID: \href{https://orcid.org/0000-0001-8006-9742}{0000-0001-8006-9742}.

Content released under the \href{https://creativecommons.org/licenses/by/4.0/}{Creative Commons Attribution 4.0 International License} (CC BY 4.0).

Source code released under the \href{https://opensource.org/licenses/MIT}{MIT License}.

Published online: \url{https://cttir.github.io/scientometrics-in-r/}
Source repository: \url{https://github.com/CTTIR/scientometrics-in-r}

Version: 0.1.0 · Build date: 2026-11-10

\section*{Suggested citation}\label{suggested-citation}
\addcontentsline{toc}{section}{Suggested citation}

\begin{quote}
Heller, R. (2026). \emph{Scientometrics in R: Bibliometrics, Network Science, and Research Evaluation}. \url{https://cttir.github.io/scientometrics-in-r/}
\end{quote}

\section*{Disclaimer}\label{disclaimer}
\addcontentsline{toc}{section}{Disclaimer}

This book is provided ``as is'' without warranty of any kind. The author assumes no responsibility for errors or omissions.

\section*{Contact}\label{contact}
\addcontentsline{toc}{section}{Contact}

For corrections, suggestions, or issues: \url{https://github.com/CTTIR/scientometrics-in-r/issues}

\chapter*{Acknowledgments}\label{acknowledgments}
\addcontentsline{toc}{chapter}{Acknowledgments}

\emph{R. Heller · Date: 2026-05-11 · Last changed: 2026-06-10}

This book stands on a wide stack of free software. Particular thanks to
the \href{https://bookdown.org/}{bookdown} authors and the broader R Core
and tidyverse communities; to the maintainers of \texttt{openalexR},
\texttt{bibliometrix}, \texttt{igraph}, \texttt{tidygraph}, \texttt{ggraph}, \texttt{quanteda}, \texttt{stm}, and
\texttt{tidytext} for the packages this book relies on most; and to the
\href{https://openalex.org/}{OpenAlex} project for making large-scale,
open bibliographic data freely available.

\section*{Inspiration}\label{inspiration}
\addcontentsline{toc}{section}{Inspiration}

The structural and pedagogical model for this book draws heavily on Mathias Harrer, Pim Cuijpers, Toshi A. Furukawa, and David D. Ebert's \href{https://bookdown.org/MathiasHarrer/Doing_Meta_Analysis_in_R/}{\emph{Doing Meta-Analysis with R: A Hands-On Guide}} --- both the book and its open-source repository at \url{https://github.com/MathiasHarrer/Doing-Meta-Analysis-in-R}. Their thorough, transparent, and beautifully laid-out bookdown setup --- sidebar grouping, semantic callout boxes, the citation block, the dark/light toggle, MathJax-rendered formulas, downloadable BibTeX/RIS citations --- set a clear standard for what a reproducible R-based research book can look like. Where the present book mirrors their conventions, it does so deliberately and gratefully. Where it diverges (colors, typography, individual icons), the divergence is cosmetic; the underlying craft is theirs to credit.

\section*{Use of LLM tools}\label{use-of-llm-tools}
\addcontentsline{toc}{section}{Use of LLM tools}

Portions of this book were prepared with assistance from large language model tooling for
narrowly defined, non-authorial tasks: copyediting, prose smoothing, Markdown/LaTeX formatting,
scaffolding of boilerplate files (CI configs, build scripts), code refactoring. The tools used were \href{https://kisski.gwdg.de/leistungen/2-02-llm-service/}{Chat AI},
the LLM service of KISSKI (GWDG), and a self-hosted \textbf{Mistral Small (24B, Apache-2.0)} run locally via
\href{https://ollama.com/}{Ollama} and the \texttt{ollamar} R package --- local inference only, with no data sent to
third parties for the self-hosted model.

All scientific claims, methodological choices, analyses, interpretations, and conclusions are the
author's own. No LLM-generated text was incorporated without review and revision, and every reference
was verified against its DOI, arXiv ID, or ISBN.

\chapter*{How to use this book}\label{how-to-use-this-book}
\addcontentsline{toc}{chapter}{How to use this book}

\emph{R. Heller · Date: 2026-05-11 · Last changed: 2026-10-31}

\section{Reading paths}\label{reading-paths}

The book is structured so that you can read it linearly --- front to back ---
or jump straight to the chapter that answers your current question.
If you do not yet know which method fits your problem, start with the
\hyperref[find-your-method]{Find Your Method navigator}. It maps
typical research questions (``evaluate a journal portfolio'', ``detect
emerging topics'', ``study collaboration patterns'') to the chapters and
case studies that demonstrate the relevant techniques.

\section{Code chunks}\label{code-chunks}

Every analytical claim is paired with a runnable R chunk. Chunks are
labelled (\texttt{label}), and figures use \texttt{fig.cap} / \texttt{fig.alt} so the book
remains accessible. Chunks that need network access, paid credentials,
or interactive widgets are marked \texttt{eval\ =\ FALSE} and document what is
required to run them locally.

\section{Callouts}\label{callouts}

Three callout styles appear throughout:

A blue callout flags a clarification or aside.

An orange callout flags a limitation or responsible-use note --- read these.

A red callout flags a common pitfall that has tripped up working scientometricians.

\section{Companion data}\label{companion-data}

Small sample datasets that the chapters fall back to (when a live API
call is not appropriate) are bundled with the companion R package
\texttt{scientometricsInR}. Larger illustrative corpora live under \texttt{data/} in
the source repository.

\section{Downloads}\label{downloads}

The sidebar's \emph{Download} menu gives you the whole book as a PDF or
EPUB. Each chapter page also exposes a \emph{Download this chapter (PDF)}
button at the top right.

\section{Citation}\label{citation}

If the book is useful to your research, please cite it --- the suggested
form lives in the Impressum, and a BibTeX block is in the README.

\chapter*{Notation}\label{notation}
\addcontentsline{toc}{chapter}{Notation}

\emph{R. Heller · Date: 2026-05-11 · Last changed: 2026-05-11}

The following symbols, abbreviations, and conventions are used
throughout the book.

\section*{Mathematical notation}\label{mathematical-notation}
\addcontentsline{toc}{section}{Mathematical notation}

{\def\LTcaptype{none} % do not increment counter
\begin{longtable}[]{@{}ll@{}}
\toprule\noalign{}
Symbol & Meaning \\
\midrule\noalign{}
\endhead
\bottomrule\noalign{}
\endlastfoot
\(c_i\) & Citations received by paper \(i\) \\
\(N\) & Number of papers in a corpus \\
\(h\) & Hirsch index (h-index) \\
\(g\) & g-index \\
\(\text{MNCS}\) & Mean normalised citation score \\
\(\bar{x}\) & Sample mean \\
\(\tilde{x}\) & Sample median \\
\(G\) & Graph \((V, E)\) --- vertices and edges \\
\(k_v\) & Degree of vertex \(v\) \\
\(B(v)\) & Betweenness centrality of vertex \(v\) \\
\end{longtable}
}

\section*{Abbreviations}\label{abbreviations}
\addcontentsline{toc}{section}{Abbreviations}

{\def\LTcaptype{none} % do not increment counter
\begin{longtable}[]{@{}ll@{}}
\toprule\noalign{}
Acronym & Expansion \\
\midrule\noalign{}
\endhead
\bottomrule\noalign{}
\endlastfoot
API & Application Programming Interface \\
DOI & Digital Object Identifier \\
FWCI & Field-Weighted Citation Impact \\
JIF & Journal Impact Factor \\
MNCS & Mean Normalised Citation Score \\
NMI & Normalised Mutual Information \\
OA & Open Access \\
ORCID & Open Researcher and Contributor ID \\
ROR & Research Organization Registry \\
RRID & Research Resource Identifier \\
SCI & Science Citation Index \\
SDG & Sustainable Development Goal \\
WoS & Web of Science \\
\end{longtable}
}

\section*{Code-style conventions}\label{code-style-conventions}
\addcontentsline{toc}{section}{Code-style conventions}

\begin{itemize}
\tightlist
\item
  Object names describe what they are (\texttt{effect\_sizes}), not how they were
  produced (\texttt{df2}).
\item
  Pipes use the base \texttt{\textbar{}\textgreater{}}. Tidyverse verbs are preferred over base R
  where readability favours them.
\item
  Network objects use \texttt{igraph} representations unless explicitly stated.
\item
  Random seeds are set per chapter; the global default is \texttt{42}.
\end{itemize}

\chapter*{About the author}\label{about-the-author}
\addcontentsline{toc}{chapter}{About the author}

\emph{R. Heller · Date: 2026-05-11 · Last changed: 2026-11-09}

\textbf{Raban Heller} is a biomedical scientist working at the intersection of
biostatistics, clinical research, and data science. He maintains the
\emph{\href{https://cttir.github.io/strategy-in-r/}{Strategy in R}} and
\emph{\href{https://cttir.github.io/methods-in-r/}{Methods in R}} books in the
same series.

\begin{itemize}
\tightlist
\item
  ORCID: \url{https://orcid.org/0000-0001-8006-9742}
\item
  GitHub: \url{https://github.com/r-heller}
\item
  Website: \url{https://r-heller.github.io/website/}
\end{itemize}

For corrections or suggestions, please open an issue at
\url{https://github.com/CTTIR/scientometrics-in-r/issues}.

\chapter*{Find Your Method}\label{find-your-method}
\addcontentsline{toc}{chapter}{Find Your Method}

\emph{R. Heller · Date: 2026-05-09 · Last changed: 2026-10-31}

Use this decision tree to navigate to the chapter or case study most relevant to your research question.
On narrow screens, scroll the diagram horizontally. The quick reference table
below provides the same routes in a compact format.

The decision-tree diagram is available in the web edition. Use the quick reference table below to select the relevant method and chapter in this edition.

\section{Quick reference table}\label{quick-reference-table}

{\def\LTcaptype{none} % do not increment counter
\begin{longtable}[]{@{}lll@{}}
\toprule\noalign{}
Research question & Key chapters & Case study \\
\midrule\noalign{}
\endhead
\bottomrule\noalign{}
\endlastfoot
Map a research field & 16, 19, 22 & CS 1: CRISPR \\
Evaluate an author & 10, 12 & --- \\
Evaluate a journal & 11 & CS 3: Journal Portfolio \\
Benchmark an institution & 13 & CS 2: Institutional \\
Detect emerging topics & 14, 24, 25 & CS 5: Emerging Topics \\
Analyze collaboration & 15, 18 & --- \\
Study gender gaps & 28 & CS 4: Gender Gap \\
Track open access & 27 & --- \\
Build a dashboard & 35, 37 & --- \\
Create a reproducible pipeline & 34, 36 & --- \\
\end{longtable}
}

\part{Foundations}\label{part-foundations}

\chapter{What Is Scientometrics?}\label{what-is-scientometrics}

\emph{R. Heller · Date: 2026-05-09 · Last changed: 2026-10-31}

\section{Learning objectives}\label{what-objectives}

After completing this chapter, you will be able to:

\begin{itemize}
\tightlist
\item
  Define scientometrics, bibliometrics, and informetrics and explain how the three fields relate
\item
  Trace the historical development of scientometrics from Garfield's citation indexes through Price's quantitative studies to the present day
\item
  Identify the types of research questions that scientometric methods can address
\item
  Retrieve a small sample of scholarly works from OpenAlex and produce basic descriptive statistics in R
\item
  Create simple visualisations of publication trends and journal distributions
\end{itemize}

\section{Setup}\label{what-setup}

\begin{Shaded}
\begin{Highlighting}[]
\FunctionTok{library}\NormalTok{(tidyverse)}
\FunctionTok{library}\NormalTok{(openalexR)}
\FunctionTok{library}\NormalTok{(glue)}

\FunctionTok{set.seed}\NormalTok{(}\DecValTok{20260509}\NormalTok{)}

\FunctionTok{source}\NormalTok{(here}\SpecialCharTok{::}\FunctionTok{here}\NormalTok{(}\StringTok{"R"}\NormalTok{, }\StringTok{"api\_helpers.R"}\NormalTok{))}
\FunctionTok{source}\NormalTok{(here}\SpecialCharTok{::}\FunctionTok{here}\NormalTok{(}\StringTok{"R"}\NormalTok{, }\StringTok{"utils.R"}\NormalTok{))}
\FunctionTok{source}\NormalTok{(here}\SpecialCharTok{::}\FunctionTok{here}\NormalTok{(}\StringTok{"R"}\NormalTok{, }\StringTok{"sci\_palette.R"}\NormalTok{))}
\end{Highlighting}
\end{Shaded}

\section{Conceptual background}\label{what-background}

\textbf{Scientometrics} is the quantitative study of science itself --- its growth, structure, productivity, and impact. Where scientists in most disciplines study the natural or social world, scientometricians turn the lens inward: they study the system that produces scientific knowledge.

The term sits within a family of overlapping labels. \textbf{Bibliometrics}, the oldest sibling, focuses on the statistical analysis of publications --- books, journal articles, and their metadata such as authors, citations, and keywords. The word was coined by Alan Pritchard in 1969 to replace the earlier ``statistical bibliography''. \textbf{Informetrics} is the broadest umbrella, encompassing any quantitative study of information transfer, including web analytics, patent analysis, and library usage statistics. In practice, the three terms are often used interchangeably, but it helps to remember that bibliometrics emphasises the \emph{document}, scientometrics emphasises the \emph{science}, and informetrics emphasises the \emph{information}.

\subsection{A brief history}\label{a-brief-history}

The intellectual roots of the field trace back to at least 1955, when Eugene \citet{garfield1955} proposed the idea of a citation index for science. Garfield argued that citations create an ``association of ideas'' --- a traceable web of intellectual influence --- and that indexing these links would transform how scholars discover literature. His vision was realised in the Science Citation Index (SCI), launched in 1964 by the Institute for Scientific Information (ISI). For the first time, researchers could systematically trace who cited whom across the entire scientific literature.

A few years later, Derek de Solla Price gave the field its quantitative backbone. In \emph{Little Science, Big Science} \citep{price1963}, Price marshalled data to show that science was growing exponentially --- in the number of journals, articles, and scientists --- and that this growth followed remarkably regular statistical patterns. He estimated that the volume of scientific literature doubles approximately every 15 years, a regularity so striking that it came to be known as ``Price's Law.'' Price's work established that science could be studied with the same empirical rigour that it applies to other phenomena.

These foundations were reinforced by earlier statistical regularities. Alfred \citet{lotka1926} had already shown that the distribution of scientific productivity follows a power law: a small number of authors produce the majority of papers. Samuel \citet{bradford1934} demonstrated a similar pattern for journal coverage: a core set of journals in any field accounts for a disproportionate share of the relevant literature. These laws, which we explore in detail in the next chapter, remain cornerstones of scientometric theory.

From these beginnings, scientometrics expanded rapidly through the latter half of the twentieth century. The h-index \citep{hirsch2005}, journal impact factors, co-citation networks, and science maps became standard tools in research evaluation and science policy. Landmark declarations such as the Leiden Manifesto \citep{hicks2015} and DORA \citep{dora2012} later established principles for responsible use of these tools, warning against mechanical reliance on any single indicator.

\subsection{The modern era}\label{the-modern-era}

Today the field is undergoing another transformation. The launch of \textbf{OpenAlex} \citep{priem2022} --- a fully open index of over 250 million scholarly works, built from Crossref, PubMed, institutional repositories, and other open sources --- has democratised access to the raw data that powers scientometric research. Researchers no longer need expensive subscriptions to Web of Science or Scopus to conduct large-scale analyses. This book takes advantage of that openness: every example uses data that any reader can reproduce for free.

\subsection{Scope of this book}\label{scope-of-this-book}

This book teaches you to do scientometrics in R. It covers:

\begin{itemize}
\tightlist
\item
  \textbf{Data acquisition} --- querying OpenAlex, Crossref, and PubMed programmatically.
\item
  \textbf{Core bibliometrics} --- computing productivity, impact, and citation indicators.
\item
  \textbf{Network analysis} --- building and analysing co-authorship, co-citation, and keyword co-occurrence networks.
\item
  \textbf{Text and topic analysis} --- extracting themes and trends from titles, abstracts, and full texts.
\item
  \textbf{Advanced topics} --- altmetrics, open science, gender equity, retraction analysis, and more.
\item
  \textbf{Production} --- building reproducible pipelines, dashboards, and reports.
\end{itemize}

Each chapter follows the same pattern: conceptual background, a fully reproducible worked example, diagnostics, limitations, and exercises.

\section{Worked example}\label{what-example}

To make things concrete from the very first chapter, we fetch a small sample of recent works from OpenAlex and compute some basic descriptive statistics. Think of this as a first taste of the workflows that the rest of the book will develop in depth.

\subsection{Fetching a sample from OpenAlex}\label{fetching-a-sample-from-openalex}

We search for works whose title or abstract mentions ``scientometrics'', published between 2014 and 2023, and draw a random sample of 500 records.

\begin{Shaded}
\begin{Highlighting}[]
\NormalTok{works }\OtherTok{\textless{}{-}} \FunctionTok{fetch\_openalex}\NormalTok{(}
  \AttributeTok{entity =} \StringTok{"works"}\NormalTok{,}
  \AttributeTok{search =} \StringTok{"scientometrics"}\NormalTok{,}
  \AttributeTok{from\_publication\_date =} \StringTok{"2014{-}01{-}01"}\NormalTok{,}
  \AttributeTok{to\_publication\_date =} \StringTok{"2023{-}12{-}31"}\NormalTok{,}
  \AttributeTok{options =} \FunctionTok{list}\NormalTok{(}\AttributeTok{sample =} \DecValTok{500}\NormalTok{, }\AttributeTok{seed =} \DecValTok{42}\NormalTok{)}
\NormalTok{)}

\FunctionTok{cat}\NormalTok{(}\FunctionTok{glue}\NormalTok{(}\StringTok{"Fetched \{nrow(works)\} works}\SpecialCharTok{\textbackslash{}n}\StringTok{"}\NormalTok{))}
\end{Highlighting}
\end{Shaded}

\begin{verbatim}
#> Fetched 500 works
\end{verbatim}

\begin{Shaded}
\begin{Highlighting}[]
\FunctionTok{glimpse}\NormalTok{(works)}
\end{Highlighting}
\end{Shaded}

\begin{verbatim}
#> Rows: 500
#> Columns: 45
#> $ id                            <chr> "https://openalex.org/W4389819307", "htt~
#> $ title                         <chr> "Plastic Waste Management: A Bibliometri~
#> $ display_name                  <chr> "Plastic Waste Management: A Bibliometri~
#> $ authorships                   <list> [<tbl_df[9 x 8]>], [<tbl_df[7 x 8]>], [~
#> $ abstract                      <chr> "Plastic waste management has emerged as~
#> $ doi                           <chr> "https://doi.org/10.3390/su152416840", "~
#> $ publication_date              <date> 2023-12-14, 2022-02-23, 2019-01-01, 202~
#> $ publication_year              <int> 2023, 2022, 2019, 2023, 2016, 2020, 2019~
#> $ relevance_score               <dbl> 1, 1, 1, 1, 1, 1, 1, 1, 1, 1, 1, 1, 1, 1~
#> $ fwci                          <dbl> 1.542, 4.761, 0.000, 0.000, 0.000, 0.000~
#> $ cited_by_count                <int> 24, 43, 0, 0, 0, 0, 10, 14, 3, 2, 2, 1, ~
#> $ counts_by_year                <list> [<data.frame[3 x 2]>], [<data.frame[5 x~
#> $ ids                           <list> <"https://openalex.org/W4389819307", "h~
#> $ type                          <chr> "article", "article", "article", "book-c~
#> $ is_oa                         <lgl> TRUE, TRUE, TRUE, FALSE, TRUE, FALSE, TR~
#> $ is_oa_anywhere                <lgl> TRUE, TRUE, TRUE, FALSE, TRUE, FALSE, TR~
#> $ oa_status                     <chr> "gold", "gold", "green", "closed", "diam~
#> $ oa_url                        <chr> "https://www.mdpi.com/2071-1050/15/24/16~
#> $ any_repository_has_fulltext   <lgl> FALSE, TRUE, TRUE, FALSE, TRUE, FALSE, T~
#> $ source_display_name           <chr> "Sustainability", "Journal of Food Quali~
#> $ source_id                     <chr> "https://openalex.org/S10134376", "https~
#> $ issn_l                        <chr> "2071-1050", "0146-9428", NA, NA, "0327-~
#> $ host_organization             <chr> "https://openalex.org/P4310310987", "htt~
#> $ host_organization_name        <chr> "Multidisciplinary Digital Publishing In~
#> $ landing_page_url              <chr> "https://doi.org/10.3390/su152416840", "~
#> $ pdf_url                       <chr> "https://www.mdpi.com/2071-1050/15/24/16~
#> $ license                       <chr> "cc-by", "cc-by", NA, NA, "cc-by", NA, "~
#> $ version                       <chr> "publishedVersion", "publishedVersion", ~
#> $ referenced_works              <list> <"https://openalex.org/W12743890", "htt~
#> $ referenced_works_count        <int> 75, 15, 34, 0, 0, 2, 17, 10, 21, 3, 24, ~
#> $ related_works                 <list> <"https://openalex.org/W3132390971", "h~
#> $ concepts                      <list> [<data.frame[24 x 5]>], [<data.frame[15~
#> $ topics                        <list> [<tbl_df[12 x 5]>], [<tbl_df[12 x 5]>],~
#> $ keywords                      <list> [<data.frame[7 x 3]>], [<data.frame[10 ~
#> $ is_paratext                   <lgl> FALSE, FALSE, FALSE, FALSE, FALSE, FALSE~
#> $ is_retracted                  <lgl> FALSE, FALSE, FALSE, FALSE, FALSE, FALSE~
#> $ language                      <chr> "en", "en", "en", "en", "es", "ru", "en"~
#> $ sustainable_development_goals <list> [<data.frame[1 x 3]>], NA, NA, [<data.f~
#> $ awards                        <list> NA, NA, NA, NA, NA, NA, NA, NA, NA, NA,~
#> $ funders                       <list> NA, NA, [<data.frame[1 x 3]>], NA, NA, ~
#> $ apc                           <list> [<data.frame[2 x 5]>], [<data.frame[2 x~
#> $ first_page                    <chr> "16840", "1", NA, "317", "143", "18", "1~
#> $ last_page                     <chr> "16840", "7", NA, "337", "158", "26", NA~
#> $ volume                        <chr> "15", "2022", NA, NA, NA, NA, NA, "92", ~
#> $ issue                         <chr> "24", NA, NA, NA, "43", "8", NA, "suppl ~
\end{verbatim}

Each row is a scholarly work. Key columns include \texttt{display\_name} (the title), \texttt{publication\_date}, \texttt{cited\_by\_count}, \texttt{doi}, and \texttt{source\_display\_name} (the source or journal name).

\subsection{Basic descriptive statistics}\label{basic-descriptive-statistics}

How many works did we retrieve, and what does the citation distribution look like?

\begin{Shaded}
\begin{Highlighting}[]
\NormalTok{works\_summary }\OtherTok{\textless{}{-}}\NormalTok{ works }\SpecialCharTok{|\textgreater{}}
  \FunctionTok{summarise}\NormalTok{(}
    \AttributeTok{n\_works       =} \FunctionTok{n}\NormalTok{(),}
    \AttributeTok{doi\_coverage  =}\NormalTok{ scales}\SpecialCharTok{::}\FunctionTok{percent}\NormalTok{(}\FunctionTok{mean}\NormalTok{(}\SpecialCharTok{!}\FunctionTok{is.na}\NormalTok{(doi))),}
    \AttributeTok{median\_cites  =} \FunctionTok{median}\NormalTok{(cited\_by\_count, }\AttributeTok{na.rm =} \ConstantTok{TRUE}\NormalTok{),}
    \AttributeTok{mean\_cites    =} \FunctionTok{round}\NormalTok{(}\FunctionTok{mean}\NormalTok{(cited\_by\_count, }\AttributeTok{na.rm =} \ConstantTok{TRUE}\NormalTok{), }\DecValTok{1}\NormalTok{),}
    \AttributeTok{max\_cites     =} \FunctionTok{max}\NormalTok{(cited\_by\_count, }\AttributeTok{na.rm =} \ConstantTok{TRUE}\NormalTok{),}
    \AttributeTok{zero\_cites    =} \FunctionTok{sum}\NormalTok{(cited\_by\_count }\SpecialCharTok{==} \DecValTok{0}\NormalTok{)}
\NormalTok{  )}

\NormalTok{works\_summary}
\end{Highlighting}
\end{Shaded}

\begin{verbatim}
#> # A tibble: 1 x 6
#>   n_works doi_coverage median_cites mean_cites max_cites zero_cites
#>     <int> <chr>               <dbl>      <dbl>     <int>      <int>
#> 1     500 90%                     5       19.8       352        124
\end{verbatim}

The large gap between the mean and median is a hallmark of citation distributions: most papers receive few citations while a small number accumulate many. This skewness is one of the most robust empirical regularities in scientometrics, and we will return to its theoretical basis in \ref{theory-background}.

\subsection{Top journals}\label{top-journals}

Which journals appear most often in our sample?

\begin{Shaded}
\begin{Highlighting}[]
\NormalTok{top\_journals }\OtherTok{\textless{}{-}}\NormalTok{ works }\SpecialCharTok{|\textgreater{}}
  \FunctionTok{filter}\NormalTok{(}\SpecialCharTok{!}\FunctionTok{is.na}\NormalTok{(source\_display\_name)) }\SpecialCharTok{|\textgreater{}}
  \FunctionTok{count}\NormalTok{(source\_display\_name, }\AttributeTok{sort =} \ConstantTok{TRUE}\NormalTok{) }\SpecialCharTok{|\textgreater{}}
  \FunctionTok{head}\NormalTok{(}\DecValTok{10}\NormalTok{)}

\NormalTok{top\_journals}
\end{Highlighting}
\end{Shaded}

\begin{verbatim}
#> # A tibble: 10 x 2
#>    source_display_name                  n
#>    <chr>                            <int>
#>  1 Sustainability                      16
#>  2 Scientometrics                       9
#>  3 arXiv (Cornell University)           7
#>  4 PLoS ONE                             6
#>  5 DergiPark (Istanbul University)      5
#>  6 Lincoln (University of Nebraska)     5
#>  7 Scientific Reports                   5
#>  8 CRC Press eBooks                     4
#>  9 OSF Preprints (OSF Preprints)        4
#> 10 Americanae (AECID Library)           3
\end{verbatim}

\subsection{Publication trend over time}\label{publication-trend-over-time}

\begin{Shaded}
\begin{Highlighting}[]
\NormalTok{works }\SpecialCharTok{|\textgreater{}}
  \FunctionTok{mutate}\NormalTok{(}\AttributeTok{year =} \FunctionTok{year}\NormalTok{(publication\_date)) }\SpecialCharTok{|\textgreater{}}
  \FunctionTok{count}\NormalTok{(year) }\SpecialCharTok{|\textgreater{}}
  \FunctionTok{ggplot}\NormalTok{(}\FunctionTok{aes}\NormalTok{(}\AttributeTok{x =}\NormalTok{ year, }\AttributeTok{y =}\NormalTok{ n)) }\SpecialCharTok{+}
  \FunctionTok{geom\_col}\NormalTok{(}\AttributeTok{fill =} \FunctionTok{palette\_sci}\NormalTok{(}\DecValTok{1}\NormalTok{)) }\SpecialCharTok{+}
  \FunctionTok{labs}\NormalTok{(}
    \AttributeTok{x =} \StringTok{"Publication year"}\NormalTok{,}
    \AttributeTok{y =} \StringTok{"Number of works (sampled)"}
\NormalTok{  ) }\SpecialCharTok{+}
  \FunctionTok{theme\_sci}\NormalTok{()}
\end{Highlighting}
\end{Shaded}

\begin{figure}

{\centering \includegraphics[width=0.9\linewidth,alt={Bar chart showing the number of sampled publications per year from 2014 to 2023, illustrating the growth trend in scientometrics research.}]{scientometrics-in-r_files/figure-latex/01-fig-what-pub-trend-1} 

}

\caption{Annual publication counts for sampled works related to scientometrics (2014--2023).}\label{fig:01-fig-what-pub-trend}
\end{figure}

The trend shows a growing volume of scientometrics-related publications over the past decade --- consistent with the broader exponential growth of science documented by \citet{price1963}.

\subsection{Top journals visualised}\label{top-journals-visualised}

\begin{Shaded}
\begin{Highlighting}[]
\NormalTok{top\_journals }\SpecialCharTok{|\textgreater{}}
  \FunctionTok{mutate}\NormalTok{(}\AttributeTok{source\_display\_name =} \FunctionTok{fct\_reorder}\NormalTok{(source\_display\_name, n)) }\SpecialCharTok{|\textgreater{}}
  \FunctionTok{ggplot}\NormalTok{(}\FunctionTok{aes}\NormalTok{(}\AttributeTok{x =}\NormalTok{ n, }\AttributeTok{y =}\NormalTok{ source\_display\_name)) }\SpecialCharTok{+}
  \FunctionTok{geom\_col}\NormalTok{(}\AttributeTok{fill =} \FunctionTok{palette\_sci}\NormalTok{(}\DecValTok{1}\NormalTok{)) }\SpecialCharTok{+}
  \FunctionTok{labs}\NormalTok{(}
    \AttributeTok{x =} \StringTok{"Number of works"}\NormalTok{,}
    \AttributeTok{y =} \ConstantTok{NULL}
\NormalTok{  ) }\SpecialCharTok{+}
  \FunctionTok{theme\_sci}\NormalTok{()}
\end{Highlighting}
\end{Shaded}

\begin{figure}

{\centering \includegraphics[width=0.9\linewidth,alt={Horizontal bar chart showing the 10 journals that contribute the most works to the sampled corpus, with Scientometrics leading.}]{scientometrics-in-r_files/figure-latex/01-fig-what-top-journals-1} 

}

\caption{Top 10 journals by number of works in the sampled scientometrics corpus.}\label{fig:01-fig-what-top-journals}
\end{figure}

\subsection{Citation distribution}\label{citation-distribution}

\begin{Shaded}
\begin{Highlighting}[]
\NormalTok{works }\SpecialCharTok{|\textgreater{}}
  \FunctionTok{filter}\NormalTok{(cited\_by\_count }\SpecialCharTok{\textgreater{}} \DecValTok{0}\NormalTok{) }\SpecialCharTok{|\textgreater{}}
  \FunctionTok{ggplot}\NormalTok{(}\FunctionTok{aes}\NormalTok{(}\AttributeTok{x =}\NormalTok{ cited\_by\_count)) }\SpecialCharTok{+}
  \FunctionTok{geom\_histogram}\NormalTok{(}\AttributeTok{bins =} \DecValTok{40}\NormalTok{, }\AttributeTok{fill =} \FunctionTok{palette\_sci}\NormalTok{(}\DecValTok{1}\NormalTok{), }\AttributeTok{colour =} \StringTok{"white"}\NormalTok{) }\SpecialCharTok{+}
  \FunctionTok{scale\_x\_log10}\NormalTok{() }\SpecialCharTok{+}
  \FunctionTok{labs}\NormalTok{(}
    \AttributeTok{x =} \StringTok{"Citation count (log scale)"}\NormalTok{,}
    \AttributeTok{y =} \StringTok{"Number of works"}
\NormalTok{  ) }\SpecialCharTok{+}
  \FunctionTok{theme\_sci}\NormalTok{()}
\end{Highlighting}
\end{Shaded}

\begin{figure}

{\centering \includegraphics[width=0.9\linewidth,alt={Histogram of citation counts on a logarithmic x-axis, revealing the heavily right-skewed distribution characteristic of citation data.}]{scientometrics-in-r_files/figure-latex/01-fig-what-citation-dist-1} 

}

\caption{Citation distribution of sampled scientometrics works (log-scaled x-axis).}\label{fig:01-fig-what-citation-dist}
\end{figure}

This right-skewed citation distribution is a universal feature of scholarly communication. In Chapter 2, we will see that it connects to deeper theoretical regularities: Lotka's Law, Bradford's Law, and the Matthew Effect.

\section{Diagnostics and interpretation}\label{what-diagnostics}

Even for a simple descriptive analysis, a few diagnostic checks are essential:

\begin{itemize}
\tightlist
\item
  \textbf{Sample vs.~population.} We drew a random sample of 500 works. The publication trend reflects the underlying population only if the sampling was truly random within the date window. OpenAlex's \texttt{sample} parameter does this, but always verify the date range of returned results.
\item
  \textbf{Source coverage.} OpenAlex is strongest in English-language, peer-reviewed journal articles. Conference proceedings, book chapters, and grey literature may be underrepresented. Our ``top journals'' list reflects what OpenAlex indexes, not the full universe of scientometrics publishing.
\item
  \textbf{Citation counts are snapshots.} The \texttt{cited\_by\_count} field reflects the current citation total at the time of the API query. Older papers have had more time to accumulate citations, so direct comparisons across years are misleading without normalization.
\end{itemize}

\begin{Shaded}
\begin{Highlighting}[]
\FunctionTok{cat}\NormalTok{(}\FunctionTok{glue}\NormalTok{(}\StringTok{"Date range: \{min(works$publication\_date)\} to \{max(works$publication\_date)\}}\SpecialCharTok{\textbackslash{}n}\StringTok{"}\NormalTok{))}
\end{Highlighting}
\end{Shaded}

\begin{verbatim}
#> Date range: 2014-01-01 to 2023-12-31
\end{verbatim}

\begin{Shaded}
\begin{Highlighting}[]
\FunctionTok{cat}\NormalTok{(}\FunctionTok{glue}\NormalTok{(}\StringTok{"Works with zero citations: \{sum(works$cited\_by\_count == 0)\}}\SpecialCharTok{\textbackslash{}n}\StringTok{"}\NormalTok{))}
\end{Highlighting}
\end{Shaded}

\begin{verbatim}
#> Works with zero citations: 124
\end{verbatim}

\begin{Shaded}
\begin{Highlighting}[]
\FunctionTok{cat}\NormalTok{(}\FunctionTok{glue}\NormalTok{(}\StringTok{"Unique journals: \{n\_distinct(works$source\_display\_name, na.rm = TRUE)\}}\SpecialCharTok{\textbackslash{}n}\StringTok{"}\NormalTok{))}
\end{Highlighting}
\end{Shaded}

\begin{verbatim}
#> Unique journals: 404
\end{verbatim}

\section{Limitations and responsible use}\label{what-limitations}

\section{Limitations and responsible use}\label{limitations-and-responsible-use}

\begin{itemize}
\tightlist
\item
  \textbf{Descriptive statistics are not evaluations.} Counting publications and citations describes \emph{activity}, not \emph{quality}. A high publication count may reflect salami slicing; a high citation count may reflect controversy rather than excellence.
\item
  \textbf{Data source bias.} OpenAlex, like all databases, has systematic coverage gaps. English-language STEM journals are overrepresented; social sciences, humanities, and non-English scholarship are underrepresented. Always report the data source and its known limitations.
\item
  \textbf{Ecological fallacy.} Aggregate statistics about a field do not apply to individual researchers. A growing field does not mean every researcher in it is productive.
\item
  \textbf{The Leiden Manifesto} \citep{hicks2015} and \textbf{DORA} \citep{dora2012} provide essential guardrails. Quantitative data should support, never replace, expert qualitative judgement.
\end{itemize}

\section{Common pitfalls}\label{what-pitfalls}

\section{Common pitfalls}\label{common-pitfalls}

\begin{itemize}
\tightlist
\item
  \textbf{Confusing scientometrics with research evaluation.} Scientometrics provides tools for studying science; it does not provide verdicts on the value of individual researchers or papers.
\item
  \textbf{Treating OpenAlex counts as ground truth.} OpenAlex counts may differ from Web of Science or Scopus due to different coverage, deduplication, and metadata parsing. Always triangulate when precision matters.
\item
  \textbf{Ignoring the citation window.} Comparing raw citation counts across publication years is misleading. A 2023 paper with 5 citations may be performing better than a 2015 paper with 20. Field- and time-normalised indicators (covered in Chapter 10) address this.
\item
  \textbf{Overgeneralising from a sample.} A random sample of 500 works is useful for learning and prototyping. Drawing conclusions about an entire field requires a comprehensive corpus.
\item
  \textbf{Forgetting to set a seed.} Random samples are only reproducible if you set the seed both in R (\texttt{set.seed()}) and in the API call (\texttt{seed} parameter). Without this, every render produces different results.
\end{itemize}

\section{Exercises}\label{what-exercises}

\begin{enumerate}
\def\labelenumi{\arabic{enumi}.}
\item
  \textbf{Explore a different field.} Replace the search term ``scientometrics'' with a topic of your choice (e.g., ``machine learning'', ``climate change''). Fetch 500 works and compare the publication trend and top journals to those we found above. What differences do you notice?
\item
  \textbf{Uncited works.} What fraction of the sampled works have zero citations? Filter to uncited works and examine their publication years. Is there a pattern? \emph{(Hint: very recent works are more likely to be uncited.)}
\item
  \textbf{Open Access status.} OpenAlex records an \texttt{is\_oa} field. Compute the proportion of open-access works in your sample and plot how it has changed over time. What trend do you observe?
\item
  \textbf{Growth curve.} Extend the publication trend analysis to 2000--2023. Does the annual publication count follow an exponential growth pattern, as \citet{price1963} predicted? \emph{(Hint: plot on a log-scaled y-axis and check for linearity.)}
\end{enumerate}

\section{Solutions}\label{01-solutions}

Solutions are provided in \ref{solutions}.

\section{Further reading}\label{what-reading}

\begin{itemize}
\tightlist
\item
  \citet{garfield1955} --- The foundational paper proposing citation indexes for science; the intellectual starting point for the entire field.
\item
  \citet{price1963} --- \emph{Little Science, Big Science}; the first systematic quantitative study of the growth of science.
\item
  \citet{priem2022} --- The OpenAlex paper; describes the open data infrastructure that powers the examples in this book.
\item
  \citet{hicks2015} --- The Leiden Manifesto; ten principles for the responsible use of research metrics.
\item
  \citet{dora2012} --- The San Francisco Declaration on Research Assessment; a landmark call to improve research evaluation.
\item
  \citet{aria2017} --- The \texttt{bibliometrix} R package; a comprehensive tool for science mapping and bibliometric analysis.
\item
  \citet{hirsch2005} --- The h-index paper; illustrates how a single indicator reshaped research evaluation discourse.
\item
  \citet{lotka1926} --- The foundational study of scientific productivity distributions; a bridge to Chapter 2.
\end{itemize}

\section{Session info}\label{what-sessioninfo}

\begin{Shaded}
\begin{Highlighting}[]
\FunctionTok{sessionInfo}\NormalTok{()}
\end{Highlighting}
\end{Shaded}

\begin{verbatim}
#> R version 4.6.1 (2026-06-24)
#> Platform: x86_64-pc-linux-gnu
#> Running under: Ubuntu 24.04.5 LTS
#> 
#> Matrix products: default
#> BLAS:   /usr/lib/x86_64-linux-gnu/openblas-pthread/libblas.so.3 
#> LAPACK: /usr/lib/x86_64-linux-gnu/openblas-pthread/libopenblasp-r0.3.26.so;  LAPACK version 3.12.0
#> 
#> locale:
#>  [1] LC_CTYPE=C.UTF-8       LC_NUMERIC=C           LC_TIME=C.UTF-8       
#>  [4] LC_COLLATE=C.UTF-8     LC_MONETARY=C.UTF-8    LC_MESSAGES=C.UTF-8   
#>  [7] LC_PAPER=C.UTF-8       LC_NAME=C              LC_ADDRESS=C          
#> [10] LC_TELEPHONE=C         LC_MEASUREMENT=C.UTF-8 LC_IDENTIFICATION=C   
#> 
#> time zone: UTC
#> tzcode source: system (glibc)
#> 
#> attached base packages:
#> [1] stats     graphics  grDevices utils     datasets  methods   base     
#> 
#> other attached packages:
#>  [1] glue_1.8.1      openalexR_3.1.2 lubridate_1.9.5 forcats_1.0.1  
#>  [5] stringr_1.6.0   dplyr_1.2.1     purrr_1.2.2     readr_2.2.0    
#>  [9] tidyr_1.3.2     tibble_3.3.1    ggplot2_4.0.3   tidyverse_2.0.0
#> 
#> loaded via a namespace (and not attached):
#>  [1] gt_1.4.0           viridis_0.6.5      utf8_1.2.6         generics_0.1.4    
#>  [5] xml2_1.6.0         stringi_1.8.9      hms_1.1.4          digest_0.6.39     
#>  [9] magrittr_2.0.5     evaluate_1.0.5     grid_4.6.1         timechange_0.4.0  
#> [13] RColorBrewer_1.1-3 bookdown_0.48      fastmap_1.2.0      rprojroot_2.1.1   
#> [17] gridExtra_2.3.1    viridisLite_0.4.3  scales_1.4.0       cli_3.6.6         
#> [21] rlang_1.3.0        withr_3.0.3        yaml_2.3.12        otel_0.2.0        
#> [25] tools_4.6.1        tzdb_0.5.0         here_1.0.2         vctrs_0.7.3       
#> [29] R6_2.6.1           lifecycle_1.0.5    fs_2.1.0           pkgconfig_2.0.3   
#> [33] pillar_1.11.1      gtable_0.3.6       systemfonts_1.3.2  xfun_0.62         
#> [37] tidyselect_1.2.1   rstudioapi_0.19.0  knitr_1.52         farver_2.1.2      
#> [41] htmltools_0.5.9    labeling_0.4.3     rmarkdown_2.32     compiler_4.6.1    
#> [45] S7_0.2.2
\end{verbatim}

\chapter{Theoretical Underpinnings}\label{theoretical-underpinnings}

\emph{R. Heller · Date: 2026-05-09 · Last changed: 2026-10-31}

\section{Learning objectives}\label{theory-objectives}

After completing this chapter, you will be able to:

\begin{itemize}
\tightlist
\item
  State Lotka's, Bradford's, and Zipf's laws and explain the empirical patterns they describe
\item
  Test whether a corpus of publications approximately follows Lotka's inverse-square law
\item
  Construct a Bradford plot to identify core and peripheral journals in a field
\item
  Explain Price's law of exponential growth and the Matthew Effect in science
\item
  Recognise power-law-like distributions in bibliometric data
\end{itemize}

\section{Setup}\label{theory-setup}

\begin{Shaded}
\begin{Highlighting}[]
\FunctionTok{library}\NormalTok{(tidyverse)}
\FunctionTok{library}\NormalTok{(openalexR)}
\FunctionTok{library}\NormalTok{(tidytext)}
\FunctionTok{library}\NormalTok{(glue)}

\FunctionTok{set.seed}\NormalTok{(}\DecValTok{20260509}\NormalTok{)}

\FunctionTok{source}\NormalTok{(here}\SpecialCharTok{::}\FunctionTok{here}\NormalTok{(}\StringTok{"R"}\NormalTok{, }\StringTok{"api\_helpers.R"}\NormalTok{))}
\FunctionTok{source}\NormalTok{(here}\SpecialCharTok{::}\FunctionTok{here}\NormalTok{(}\StringTok{"R"}\NormalTok{, }\StringTok{"sci\_palette.R"}\NormalTok{))}
\end{Highlighting}
\end{Shaded}

\section{Conceptual background}\label{theory-background}

Scientometrics rests on a small number of empirical regularities discovered across decades of studying scholarly communication:

\textbf{Lotka's Law} \citep{lotka1926} observes that the number of authors producing \emph{n} papers is approximately proportional to 1/\emph{n}². A few prolific authors generate a large share of output, while most researchers publish only one or two papers.

\textbf{Bradford's Law} \citep{bradford1934} describes the scattering of articles across journals. For a given topic, a small core of journals contains roughly one-third of all relevant articles; a larger zone contains the next third; and a much larger periphery the final third. Each successive zone contains roughly the same number of articles but geometrically more journals.

\textbf{Zipf's Law} states that the frequency of a word is inversely proportional to its rank. In scientific text, a handful of terms dominate while the vast majority appear rarely.

\textbf{Price's Law} \citep{price1963} generalises these observations: the cumulative output of science grows exponentially, doubling roughly every 10--15 years. Price also noted that the square root of all contributors produce half of all contributions.

The \textbf{Matthew Effect} \citep{merton1968} --- named after the Gospel of Matthew (``to those who have, more shall be given'') --- describes cumulative advantage in science. Well-known researchers attract disproportionate credit and resources, amplifying initial differences in visibility. This mechanism helps explain the skewed distributions observed by Lotka and Price.

\section{Worked example}\label{theory-example}

We use a corpus of works from the journal \emph{Scientometrics} to demonstrate each law.

\subsection{Fetching the data}\label{fetching-the-data}

\begin{Shaded}
\begin{Highlighting}[]
\NormalTok{works }\OtherTok{\textless{}{-}} \FunctionTok{fetch\_openalex}\NormalTok{(}
  \AttributeTok{entity =} \StringTok{"works"}\NormalTok{,}
  \AttributeTok{primary\_location.source.id =} \StringTok{"S148561398"}\NormalTok{,}
  \AttributeTok{from\_publication\_date =} \StringTok{"2015{-}01{-}01"}\NormalTok{,}
  \AttributeTok{to\_publication\_date =} \StringTok{"2023{-}12{-}31"}\NormalTok{,}
  \AttributeTok{options =} \FunctionTok{list}\NormalTok{(}\AttributeTok{sample =} \DecValTok{500}\NormalTok{, }\AttributeTok{seed =} \DecValTok{42}\NormalTok{)}
\NormalTok{)}
\end{Highlighting}
\end{Shaded}

\subsection{Lotka's Law}\label{lotkas-law}

We unnest authors and count the number of papers per author in this sample.

\begin{Shaded}
\begin{Highlighting}[]
\NormalTok{author\_prod }\OtherTok{\textless{}{-}}\NormalTok{ works }\SpecialCharTok{|\textgreater{}}
  \FunctionTok{select}\NormalTok{(id, authorships) }\SpecialCharTok{|\textgreater{}}
  \FunctionTok{unnest}\NormalTok{(authorships, }\AttributeTok{names\_sep =} \StringTok{"\_"}\NormalTok{) }\SpecialCharTok{|\textgreater{}}
  \FunctionTok{count}\NormalTok{(authorships\_display\_name, }\AttributeTok{name =} \StringTok{"n\_papers"}\NormalTok{) }\SpecialCharTok{|\textgreater{}}
  \FunctionTok{count}\NormalTok{(n\_papers, }\AttributeTok{name =} \StringTok{"n\_authors"}\NormalTok{) }\SpecialCharTok{|\textgreater{}}
  \FunctionTok{arrange}\NormalTok{(n\_papers)}

\NormalTok{author\_prod}
\end{Highlighting}
\end{Shaded}

\begin{verbatim}
#> # A tibble: 8 x 2
#>   n_papers n_authors
#>      <int>     <int>
#> 1        1       908
#> 2        2       152
#> 3        3        29
#> 4        4        10
#> 5        5         4
#> 6        6         2
#> 7        9         1
#> 8       12         1
\end{verbatim}

\begin{Shaded}
\begin{Highlighting}[]
\NormalTok{lotka\_c }\OtherTok{\textless{}{-}}\NormalTok{ author\_prod}\SpecialCharTok{$}\NormalTok{n\_authors[}\DecValTok{1}\NormalTok{]}

\NormalTok{author\_prod }\SpecialCharTok{|\textgreater{}}
  \FunctionTok{ggplot}\NormalTok{(}\FunctionTok{aes}\NormalTok{(}\AttributeTok{x =}\NormalTok{ n\_papers, }\AttributeTok{y =}\NormalTok{ n\_authors)) }\SpecialCharTok{+}
  \FunctionTok{geom\_point}\NormalTok{(}\AttributeTok{colour =} \FunctionTok{palette\_sci}\NormalTok{(}\DecValTok{1}\NormalTok{), }\AttributeTok{size =} \DecValTok{2}\NormalTok{) }\SpecialCharTok{+}
  \FunctionTok{geom\_line}\NormalTok{(}
    \FunctionTok{aes}\NormalTok{(}\AttributeTok{y =}\NormalTok{ lotka\_c }\SpecialCharTok{/}\NormalTok{ n\_papers}\SpecialCharTok{\^{}}\DecValTok{2}\NormalTok{),}
    \AttributeTok{linetype =} \StringTok{"dashed"}\NormalTok{, }\AttributeTok{colour =} \StringTok{"grey40"}
\NormalTok{  ) }\SpecialCharTok{+}
  \FunctionTok{scale\_x\_log10}\NormalTok{() }\SpecialCharTok{+}
  \FunctionTok{scale\_y\_log10}\NormalTok{() }\SpecialCharTok{+}
  \FunctionTok{labs}\NormalTok{(}\AttributeTok{x =} \StringTok{"Papers per author"}\NormalTok{, }\AttributeTok{y =} \StringTok{"Number of authors"}\NormalTok{) }\SpecialCharTok{+}
  \FunctionTok{theme\_sci}\NormalTok{()}
\end{Highlighting}
\end{Shaded}

\begin{figure}

{\centering \includegraphics[width=0.9\linewidth,alt={Log-log scatter plot of the number of authors versus the number of papers they published, with a dashed reference line showing the inverse-square law.}]{scientometrics-in-r_files/figure-latex/02-fig-theory-lotka-1} 

}

\caption{Author productivity distribution. The dashed line shows the inverse-square law predicted by Lotka.}\label{fig:02-fig-theory-lotka}
\end{figure}

The data approximate the inverse-square pattern, though real distributions often deviate from the exact exponent of 2.

\subsection{Bradford's Law}\label{bradfords-law}

We count how many articles each journal source contributes (across a broader query).

\begin{Shaded}
\begin{Highlighting}[]
\NormalTok{broad\_works }\OtherTok{\textless{}{-}} \FunctionTok{fetch\_openalex}\NormalTok{(}
  \AttributeTok{entity =} \StringTok{"works"}\NormalTok{,}
  \AttributeTok{search =} \StringTok{"bibliometrics"}\NormalTok{,}
  \AttributeTok{from\_publication\_date =} \StringTok{"2020{-}01{-}01"}\NormalTok{,}
  \AttributeTok{to\_publication\_date =} \StringTok{"2023{-}12{-}31"}\NormalTok{,}
  \AttributeTok{options =} \FunctionTok{list}\NormalTok{(}\AttributeTok{sample =} \DecValTok{1000}\NormalTok{, }\AttributeTok{seed =} \DecValTok{42}\NormalTok{)}
\NormalTok{)}

\NormalTok{journal\_counts }\OtherTok{\textless{}{-}}\NormalTok{ broad\_works }\SpecialCharTok{|\textgreater{}}
  \FunctionTok{filter}\NormalTok{(}\SpecialCharTok{!}\FunctionTok{is.na}\NormalTok{(source\_display\_name)) }\SpecialCharTok{|\textgreater{}}
  \FunctionTok{count}\NormalTok{(source\_display\_name, }\AttributeTok{name =} \StringTok{"articles"}\NormalTok{, }\AttributeTok{sort =} \ConstantTok{TRUE}\NormalTok{) }\SpecialCharTok{|\textgreater{}}
  \FunctionTok{mutate}\NormalTok{(}
    \AttributeTok{cum\_articles =} \FunctionTok{cumsum}\NormalTok{(articles),}
    \AttributeTok{rank =} \FunctionTok{row\_number}\NormalTok{()}
\NormalTok{  )}
\end{Highlighting}
\end{Shaded}

\begin{Shaded}
\begin{Highlighting}[]
\NormalTok{journal\_counts }\SpecialCharTok{|\textgreater{}}
  \FunctionTok{ggplot}\NormalTok{(}\FunctionTok{aes}\NormalTok{(}\AttributeTok{x =}\NormalTok{ rank, }\AttributeTok{y =}\NormalTok{ cum\_articles)) }\SpecialCharTok{+}
  \FunctionTok{geom\_line}\NormalTok{(}\AttributeTok{colour =} \FunctionTok{palette\_sci}\NormalTok{(}\DecValTok{1}\NormalTok{), }\AttributeTok{linewidth =} \DecValTok{1}\NormalTok{) }\SpecialCharTok{+}
  \FunctionTok{scale\_x\_log10}\NormalTok{() }\SpecialCharTok{+}
  \FunctionTok{labs}\NormalTok{(}\AttributeTok{x =} \StringTok{"Journal rank (log scale)"}\NormalTok{, }\AttributeTok{y =} \StringTok{"Cumulative articles"}\NormalTok{) }\SpecialCharTok{+}
  \FunctionTok{theme\_sci}\NormalTok{()}
\end{Highlighting}
\end{Shaded}

\begin{figure}

{\centering \includegraphics[width=0.9\linewidth,alt={Line plot showing cumulative article count on the y-axis versus journal rank on a logarithmic x-axis. The characteristic S-shaped Bradford curve is visible, with a steep rise for core journals.}]{scientometrics-in-r_files/figure-latex/02-fig-theory-bradford-1} 

}

\caption{Bradford plot: cumulative articles vs. log journal rank.}\label{fig:02-fig-theory-bradford}
\end{figure}

\subsection{Zipf's Law}\label{zipfs-law}

We extract title words and examine their frequency distribution.

\begin{Shaded}
\begin{Highlighting}[]
\NormalTok{word\_freq }\OtherTok{\textless{}{-}}\NormalTok{ works }\SpecialCharTok{|\textgreater{}}
  \FunctionTok{select}\NormalTok{(display\_name) }\SpecialCharTok{|\textgreater{}}
  \FunctionTok{unnest\_tokens}\NormalTok{(word, display\_name, }\AttributeTok{token =} \StringTok{"words"}\NormalTok{) }\SpecialCharTok{|\textgreater{}}
  \FunctionTok{anti\_join}\NormalTok{(tidytext}\SpecialCharTok{::}\FunctionTok{get\_stopwords}\NormalTok{(), }\AttributeTok{by =} \StringTok{"word"}\NormalTok{) }\SpecialCharTok{|\textgreater{}}
  \FunctionTok{count}\NormalTok{(word, }\AttributeTok{sort =} \ConstantTok{TRUE}\NormalTok{) }\SpecialCharTok{|\textgreater{}}
  \FunctionTok{mutate}\NormalTok{(}\AttributeTok{rank =} \FunctionTok{row\_number}\NormalTok{())}
\end{Highlighting}
\end{Shaded}

\begin{Shaded}
\begin{Highlighting}[]
\NormalTok{word\_freq }\SpecialCharTok{|\textgreater{}}
  \FunctionTok{head}\NormalTok{(}\DecValTok{200}\NormalTok{) }\SpecialCharTok{|\textgreater{}}
  \FunctionTok{ggplot}\NormalTok{(}\FunctionTok{aes}\NormalTok{(}\AttributeTok{x =}\NormalTok{ rank, }\AttributeTok{y =}\NormalTok{ n)) }\SpecialCharTok{+}
  \FunctionTok{geom\_point}\NormalTok{(}\AttributeTok{colour =} \FunctionTok{palette\_sci}\NormalTok{(}\DecValTok{1}\NormalTok{), }\AttributeTok{size =} \DecValTok{1}\NormalTok{, }\AttributeTok{alpha =} \FloatTok{0.7}\NormalTok{) }\SpecialCharTok{+}
  \FunctionTok{scale\_x\_log10}\NormalTok{() }\SpecialCharTok{+}
  \FunctionTok{scale\_y\_log10}\NormalTok{() }\SpecialCharTok{+}
  \FunctionTok{labs}\NormalTok{(}\AttributeTok{x =} \StringTok{"Rank"}\NormalTok{, }\AttributeTok{y =} \StringTok{"Frequency"}\NormalTok{) }\SpecialCharTok{+}
  \FunctionTok{theme\_sci}\NormalTok{()}
\end{Highlighting}
\end{Shaded}

\begin{figure}

{\centering \includegraphics[width=0.9\linewidth,alt={Log-log scatter plot of word frequency versus rank, showing an approximately linear relationship consistent with Zipf's law.}]{scientometrics-in-r_files/figure-latex/02-fig-theory-zipf-1} 

}

\caption{Word frequency vs. rank for article titles, approximating Zipf's law.}\label{fig:02-fig-theory-zipf}
\end{figure}

\subsection{Price's growth law}\label{prices-growth-law}

\begin{Shaded}
\begin{Highlighting}[]
\NormalTok{works }\SpecialCharTok{|\textgreater{}}
  \FunctionTok{mutate}\NormalTok{(}\AttributeTok{year =} \FunctionTok{year}\NormalTok{(publication\_date)) }\SpecialCharTok{|\textgreater{}}
  \FunctionTok{count}\NormalTok{(year) }\SpecialCharTok{|\textgreater{}}
  \FunctionTok{ggplot}\NormalTok{(}\FunctionTok{aes}\NormalTok{(}\AttributeTok{x =}\NormalTok{ year, }\AttributeTok{y =}\NormalTok{ n)) }\SpecialCharTok{+}
  \FunctionTok{geom\_col}\NormalTok{(}\AttributeTok{fill =} \FunctionTok{palette\_sci}\NormalTok{(}\DecValTok{1}\NormalTok{)) }\SpecialCharTok{+}
  \FunctionTok{labs}\NormalTok{(}\AttributeTok{x =} \StringTok{"Year"}\NormalTok{, }\AttributeTok{y =} \StringTok{"Works in sample"}\NormalTok{) }\SpecialCharTok{+}
  \FunctionTok{theme\_sci}\NormalTok{()}
\end{Highlighting}
\end{Shaded}

\begin{figure}

{\centering \includegraphics[width=0.9\linewidth,alt={Bar chart showing increasing annual publication counts from 2015 to 2023.}]{scientometrics-in-r_files/figure-latex/02-fig-theory-growth-1} 

}

\caption{Annual publication counts for sampled works, illustrating the growth of bibliometrics literature.}\label{fig:02-fig-theory-growth}
\end{figure}

\section{Diagnostics and interpretation}\label{theory-diagnostics}

These laws are \emph{empirical regularities}, not physical laws. They describe approximate patterns:

\begin{itemize}
\tightlist
\item
  Lotka's exponent is often not exactly 2; values between 1.5 and 3.5 are common depending on the field and time window.
\item
  Bradford zones are sensitive to how you define the subject scope.
\item
  Zipf's law holds best for common words; the tail deviates.
\item
  Price's exponential growth eventually saturates --- science cannot grow faster than the economy indefinitely.
\end{itemize}

The value of these laws is not predictive precision but the insight that \textbf{skewed distributions are the norm, not the exception}, in scholarly communication.

\section{Limitations and responsible use}\label{theory-limitations}

\section{Limitations and responsible use}\label{limitations-and-responsible-use-1}

\begin{itemize}
\tightlist
\item
  These laws describe aggregate patterns, not individual behaviour. A researcher publishing one paper is not ``less productive'' in any evaluative sense --- Lotka's law is a statistical observation, not a performance standard.
\item
  The Matthew Effect \citep{merton1968} is a mechanism that \emph{amplifies inequality}. Recognising it should prompt caution about using cumulative metrics (like the h-index) that inherit and reinforce this bias.
\item
  Power-law fitting is notoriously difficult. Many distributions that ``look'' power-law on a log-log plot are better described by lognormal or stretched exponential models. Rigorous fitting requires maximum-likelihood methods and goodness-of-fit tests \citep{hicks2015}.
\end{itemize}

\section{Common pitfalls}\label{theory-pitfalls}

\section{Common pitfalls}\label{common-pitfalls-1}

\begin{itemize}
\tightlist
\item
  \textbf{Eyeballing log-log plots.} A straight line on a log-log plot does not prove a power law. Use formal statistical tests.
\item
  \textbf{Confusing description with prescription.} ``Most authors publish few papers'' does not mean they \emph{should} publish more.
\item
  \textbf{Ignoring field differences.} The exponent of Lotka's law and the scattering pattern of Bradford's law vary substantially across disciplines.
\item
  \textbf{Assuming stationarity.} These patterns can shift over time as publication norms change (e.g., the rise of mega-journals).
\end{itemize}

\section{Exercises}\label{theory-exercises}

\begin{enumerate}
\def\labelenumi{\arabic{enumi}.}
\item
  \textbf{Lotka exponent.} Fit a linear model to the log-log author productivity data. What exponent do you estimate? How does it compare to Lotka's original value of 2?
\item
  \textbf{Bradford zones.} Divide the journal list into three zones of roughly equal article counts. How many journals are in each zone? What is the ratio between successive zone sizes?
\item
  \textbf{Your field's Zipf distribution.} Fetch 500 works in a field of your choice. Extract title words and plot the Zipf distribution. Are the most frequent terms what you would expect?
\item
  \textbf{Matthew Effect simulation.} Write a simple simulation: start 100 authors with 1 paper each. At each time step, an author gains a new paper with probability proportional to their current count. After 50 steps, plot the distribution. Does it resemble Lotka's law?
\end{enumerate}

\section{Solutions}\label{02-solutions}

Solutions are provided in \ref{solutions}.

\section{Further reading}\label{theory-reading}

\begin{itemize}
\tightlist
\item
  \citet{lotka1926} --- The original observation on author productivity distributions.
\item
  \citet{bradford1934} --- Journal scattering and the Bradford law.
\item
  \citet{price1963} --- Exponential growth of science and cumulative advantage.
\item
  \citet{merton1968} --- The Matthew Effect: how advantage accumulates in science.
\item
  \citet{garfield1955} --- Citation indexing, the empirical basis for much of scientometrics.
\item
  \citet{waltman2016} --- A modern review contextualising these classical laws within citation impact measurement.
\end{itemize}

\section{Session info}\label{theory-sessioninfo}

\begin{Shaded}
\begin{Highlighting}[]
\FunctionTok{sessionInfo}\NormalTok{()}
\end{Highlighting}
\end{Shaded}

\begin{verbatim}
#> R version 4.6.1 (2026-06-24)
#> Platform: x86_64-pc-linux-gnu
#> Running under: Ubuntu 24.04.5 LTS
#> 
#> Matrix products: default
#> BLAS:   /usr/lib/x86_64-linux-gnu/openblas-pthread/libblas.so.3 
#> LAPACK: /usr/lib/x86_64-linux-gnu/openblas-pthread/libopenblasp-r0.3.26.so;  LAPACK version 3.12.0
#> 
#> locale:
#>  [1] LC_CTYPE=C.UTF-8       LC_NUMERIC=C           LC_TIME=C.UTF-8       
#>  [4] LC_COLLATE=C.UTF-8     LC_MONETARY=C.UTF-8    LC_MESSAGES=C.UTF-8   
#>  [7] LC_PAPER=C.UTF-8       LC_NAME=C              LC_ADDRESS=C          
#> [10] LC_TELEPHONE=C         LC_MEASUREMENT=C.UTF-8 LC_IDENTIFICATION=C   
#> 
#> time zone: UTC
#> tzcode source: system (glibc)
#> 
#> attached base packages:
#> [1] stats     graphics  grDevices utils     datasets  methods   base     
#> 
#> other attached packages:
#>  [1] tidytext_0.4.3  glue_1.8.1      openalexR_3.1.2 lubridate_1.9.5
#>  [5] forcats_1.0.1   stringr_1.6.0   dplyr_1.2.1     purrr_1.2.2    
#>  [9] readr_2.2.0     tidyr_1.3.2     tibble_3.3.1    ggplot2_4.0.3  
#> [13] tidyverse_2.0.0
#> 
#> loaded via a namespace (and not attached):
#>  [1] janeaustenr_1.0.0  gt_1.4.0           viridis_0.6.5      utf8_1.2.6        
#>  [5] generics_0.1.4     xml2_1.6.0         lattice_0.22-9     stringi_1.8.9     
#>  [9] hms_1.1.4          digest_0.6.39      magrittr_2.0.5     evaluate_1.0.5    
#> [13] grid_4.6.1         timechange_0.4.0   RColorBrewer_1.1-3 bookdown_0.48     
#> [17] fastmap_1.2.0      Matrix_1.7-5       rprojroot_2.1.1    gridExtra_2.3.1   
#> [21] stopwords_2.3      viridisLite_0.4.3  scales_1.4.0       cli_3.6.6         
#> [25] rlang_1.3.0        tokenizers_0.3.0   withr_3.0.3        yaml_2.3.12       
#> [29] otel_0.2.0         tools_4.6.1        tzdb_0.5.0         here_1.0.2        
#> [33] vctrs_0.7.3        R6_2.6.1           lifecycle_1.0.5    fs_2.1.0          
#> [37] pkgconfig_2.0.3    pillar_1.11.1      gtable_0.3.6       Rcpp_1.1.2        
#> [41] systemfonts_1.3.2  xfun_0.62          tidyselect_1.2.1   rstudioapi_0.19.0 
#> [45] knitr_1.52         farver_2.1.2       SnowballC_0.7.1    htmltools_0.5.9   
#> [49] labeling_0.4.3     rmarkdown_2.32     compiler_4.6.1     S7_0.2.2
\end{verbatim}

\chapter{Ethics and Responsible Metrics}\label{ethics-and-responsible-metrics}

\emph{R. Heller · Date: 2026-05-09 · Last changed: 2026-10-31}

\section{Learning objectives}\label{ethics-objectives}

After completing this chapter, you will be able to:

\begin{itemize}
\tightlist
\item
  Summarize the ten principles of the Leiden Manifesto and explain their implications for research evaluation
\item
  Describe the key commitments of the San Francisco Declaration on Research Assessment (DORA)
\item
  Identify common forms of metric gaming, including self-citation cartels, citation stacking, and salami slicing
\item
  Evaluate when quantitative indicators are appropriate --- and when they are not
\item
  Compute and interpret author-level self-citation rates using OpenAlex data
\item
  Apply a responsible-use framework when designing or reviewing bibliometric analyses
\end{itemize}

\section{Setup}\label{ethics-setup}

\begin{Shaded}
\begin{Highlighting}[]
\FunctionTok{library}\NormalTok{(tidyverse)}
\FunctionTok{library}\NormalTok{(openalexR)}
\FunctionTok{library}\NormalTok{(glue)}
\FunctionTok{library}\NormalTok{(gt)}

\FunctionTok{set.seed}\NormalTok{(}\DecValTok{20260509}\NormalTok{)}

\FunctionTok{source}\NormalTok{(here}\SpecialCharTok{::}\FunctionTok{here}\NormalTok{(}\StringTok{"R"}\NormalTok{, }\StringTok{"api\_helpers.R"}\NormalTok{))}
\FunctionTok{source}\NormalTok{(here}\SpecialCharTok{::}\FunctionTok{here}\NormalTok{(}\StringTok{"R"}\NormalTok{, }\StringTok{"utils.R"}\NormalTok{))}
\FunctionTok{source}\NormalTok{(here}\SpecialCharTok{::}\FunctionTok{here}\NormalTok{(}\StringTok{"R"}\NormalTok{, }\StringTok{"sci\_palette.R"}\NormalTok{))}
\end{Highlighting}
\end{Shaded}

\section{Conceptual background}\label{ethics-background}

Metrics shape behaviour. When a number is used to judge researchers, departments, or journals, people adjust their actions to improve that number --- regardless of whether the number actually captures what it was meant to measure. This insight, often called \textbf{Goodhart's Law}, is the intellectual foundation for the responsible metrics movement: ``When a measure becomes a target, it ceases to be a good measure'' \citep{goodhart1984}.

\subsection{The Leiden Manifesto}\label{the-leiden-manifesto}

In 2015, \citet{hicks2015} published the \textbf{Leiden Manifesto for Research Metrics} --- ten principles distilled from decades of experience in bibliometric practice. The principles are:

\begin{enumerate}
\def\labelenumi{\arabic{enumi}.}
\tightlist
\item
  \textbf{Quantitative evaluation should support qualitative, expert assessment.} Metrics should inform human judgement, never replace it.
\item
  \textbf{Measure performance against the research missions of the institution, group, or researcher.} Not everyone pursues the same goals; evaluation criteria must reflect this.
\item
  \textbf{Protect excellence in locally relevant research.} Research published in regional languages or addressing local problems must not be penalised for lower citation rates.
\item
  \textbf{Keep data collection and analytical processes open, transparent, and simple.} Black-box indicators erode trust and prevent scrutiny.
\item
  \textbf{Allow those evaluated to verify data and analysis.} Researchers should be able to inspect and challenge the data used to judge them.
\item
  \textbf{Account for variation by field in publication and citation practices.} Citation rates in mathematics and molecular biology differ by an order of magnitude; raw counts cannot be compared across fields.
\item
  \textbf{Base assessment of individual researchers on a qualitative judgement of their portfolio.} A single number (such as the h-index) cannot capture the breadth of a researcher's contribution.
\item
  \textbf{Avoid misplaced concreteness and false precision.} Three decimal places do not make a flawed indicator meaningful.
\item
  \textbf{Recognize the systemic effects of assessment and indicators.} Indicators reshape the system they measure. If you reward publication volume, you incentivise salami slicing.
\item
  \textbf{Scrutinize indicators regularly and update them.} The research landscape evolves; indicators must evolve with it.
\end{enumerate}

\subsection{The San Francisco Declaration (DORA)}\label{the-san-francisco-declaration-dora}

The \textbf{San Francisco Declaration on Research Assessment} \citep{dora2012} emerged from the 2012 Annual Meeting of the American Society for Cell Biology. Its central demand is straightforward: do not use journal-level metrics --- especially the Journal Impact Factor --- as a proxy for the quality of individual research articles. DORA's key commitments include:

\begin{itemize}
\tightlist
\item
  \textbf{For funders and institutions}: evaluate research on its own merits, not on the journal in which it was published. Use article-level metrics and qualitative indicators.
\item
  \textbf{For publishers}: reduce emphasis on the Journal Impact Factor as a promotional tool. Make article-level metrics (downloads, citations, altmetrics) available.
\item
  \textbf{For researchers}: challenge evaluation practices that rely on journal-based metrics. When serving on hiring or promotion committees, assess the content of papers, not where they appeared.
\end{itemize}

As of 2026, DORA has been signed by over 3,000 organisations and 20,000 individuals. Many major funders --- including the European Commission, Wellcome Trust, and the Dutch Research Council (NWO) --- have formally adopted DORA principles.

\subsection{The Metric Tide}\label{the-metric-tide}

The UK's independent review of metrics in research assessment, \emph{The Metric Tide} \citep{wilsdon2015}, reinforced these principles with a systematic evidence base. The report concluded that no metric can substitute for expert peer review, and recommended that the UK research system adopt a framework of ``responsible metrics'' built on five dimensions: robustness, humility, transparency, diversity, and reflexivity.

\subsection{Goodhart's Law in practice}\label{goodharts-law-in-practice}

When metrics are attached to career consequences, the incentive to game them becomes powerful \citep{goodhart1984}. Common forms of metric manipulation include:

\begin{itemize}
\tightlist
\item
  \textbf{Self-citation inflation}: authors systematically cite their own prior work to boost citation counts, even when those citations are not intellectually necessary.
\item
  \textbf{Citation cartels}: informal agreements among groups of authors or journals to cite each other's work, inflating citation counts reciprocally.
\item
  \textbf{Citation stacking}: journals encourage or require authors to add citations to recently published articles in the same journal, artificially increasing the Impact Factor.
\item
  \textbf{Salami slicing}: splitting a single study into the maximum number of ``least publishable units'' to inflate publication counts.
\end{itemize}

These practices are not merely hypothetical. Clarivate has repeatedly suspended journals from the Journal Citation Reports for citation stacking, and several high-profile cases of self-citation manipulation have led to editorial sanctions.

\subsection{Publish or perish}\label{publish-or-perish}

The systemic root of metric gaming is the \textbf{``publish or perish''} culture that dominates academic career structures worldwide. When hiring, promotion, and tenure decisions are driven by publication volume and citation counts, researchers face relentless pressure to produce more papers, chase high-impact journals, and avoid risky long-term projects. The consequences include:

\begin{itemize}
\tightlist
\item
  \textbf{Reduced risk-taking}: researchers avoid novel or interdisciplinary questions that might not produce quick, high-citation publications.
\item
  \textbf{Replication crisis}: pressure to publish positive results contributes to publication bias and irreproducible findings.
\item
  \textbf{Mental health harms}: the relentless pressure to produce has been linked to elevated rates of anxiety, burnout, and attrition in academia.
\item
  \textbf{Erosion of research quality}: when volume is rewarded over substance, the incentive structure favours quantity at the expense of depth and rigour.
\end{itemize}

These dynamics make the responsible use of metrics not merely a technical concern but an ethical imperative.

\section{Worked example}\label{ethics-example}

We now demonstrate a practical analysis: computing self-citation rates for a sample of researchers in the field of scientometrics. This illustrates how OpenAlex data can be used to flag potential self-citation inflation --- and the care required to interpret such analysis responsibly.

\subsection{Acquiring author data}\label{acquiring-author-data}

We fetch a sample of authors associated with scientometrics research who have at least 50 works, ensuring we have enough output to compute meaningful self-citation rates.

\begin{Shaded}
\begin{Highlighting}[]
\NormalTok{authors\_raw }\OtherTok{\textless{}{-}} \FunctionTok{fetch\_openalex}\NormalTok{(}
  \AttributeTok{entity =} \StringTok{"authors"}\NormalTok{,}
  \AttributeTok{topics.id =} \StringTok{"T10102"}\NormalTok{,}
  \AttributeTok{works\_count =} \StringTok{"\textgreater{}50"}\NormalTok{,}
  \AttributeTok{options =} \FunctionTok{list}\NormalTok{(}\AttributeTok{sample =} \DecValTok{30}\NormalTok{, }\AttributeTok{seed =} \DecValTok{42}\NormalTok{)}
\NormalTok{)}

\NormalTok{authors }\OtherTok{\textless{}{-}}\NormalTok{ authors\_raw }\SpecialCharTok{|\textgreater{}}
  \FunctionTok{select}\NormalTok{(id, display\_name, works\_count, cited\_by\_count) }\SpecialCharTok{|\textgreater{}}
  \FunctionTok{filter}\NormalTok{(}\SpecialCharTok{!}\FunctionTok{is.na}\NormalTok{(display\_name))}

\FunctionTok{glimpse}\NormalTok{(authors)}
\end{Highlighting}
\end{Shaded}

\begin{verbatim}
#> Rows: 30
#> Columns: 4
#> $ id             <chr> "https://openalex.org/A5011410644", "https://openalex.o~
#> $ display_name   <chr> "Xukun Wang", "Kristen L. Scotti", "Khalifa Ahsina", "C~
#> $ works_count    <int> 97, 64, 60, 62, 114, 63, 55, 83, 54, 62, 116, 129, 59, ~
#> $ cited_by_count <int> 1718, 504, 79, 106, 239, 131, 209, 31, 582, 0, 531, 113~
\end{verbatim}

\subsection{Fetching works and computing self-citations}\label{fetching-works-and-computing-self-citations}

For each author, we retrieve a sample of their recent works and examine how many of the citations come from the author's own papers. OpenAlex provides \texttt{cited\_by\_count} at the work level and allows us to query an author's works directly.

\begin{Shaded}
\begin{Highlighting}[]
\CommentTok{\# For a manageable demonstration, pick 10 authors}
\NormalTok{authors\_sample }\OtherTok{\textless{}{-}}\NormalTok{ authors }\SpecialCharTok{|\textgreater{}}
  \FunctionTok{slice\_sample}\NormalTok{(}\AttributeTok{n =} \FunctionTok{min}\NormalTok{(}\DecValTok{10}\NormalTok{, }\FunctionTok{nrow}\NormalTok{(authors)))}

\CommentTok{\# Function to estimate self{-}citation rate for one author}
\NormalTok{estimate\_self\_citation }\OtherTok{\textless{}{-}} \ControlFlowTok{function}\NormalTok{(author\_id, author\_name) \{}
  \CommentTok{\# Fetch up to 50 recent works by this author}
\NormalTok{  works }\OtherTok{\textless{}{-}} \FunctionTok{tryCatch}\NormalTok{(}
    \FunctionTok{fetch\_openalex}\NormalTok{(}
      \AttributeTok{entity =} \StringTok{"works"}\NormalTok{,}
      \AttributeTok{author.id =}\NormalTok{ author\_id,}
      \AttributeTok{from\_publication\_date =} \StringTok{"2018{-}01{-}01"}\NormalTok{,}
      \AttributeTok{options =} \FunctionTok{list}\NormalTok{(}\AttributeTok{sample =} \DecValTok{50}\NormalTok{, }\AttributeTok{seed =} \DecValTok{42}\NormalTok{)}
\NormalTok{    ),}
    \AttributeTok{error =} \ControlFlowTok{function}\NormalTok{(e) }\FunctionTok{tibble}\NormalTok{()}
\NormalTok{  )}

  \ControlFlowTok{if}\NormalTok{ (}\FunctionTok{is.null}\NormalTok{(works) }\SpecialCharTok{||} \FunctionTok{nrow}\NormalTok{(works) }\SpecialCharTok{==} \DecValTok{0}\NormalTok{) \{}
    \FunctionTok{return}\NormalTok{(}\FunctionTok{tibble}\NormalTok{(}
      \AttributeTok{author\_id =}\NormalTok{ author\_id,}
      \AttributeTok{author\_name =}\NormalTok{ author\_name,}
      \AttributeTok{n\_works =} \DecValTok{0}\DataTypeTok{L}\NormalTok{,}
      \AttributeTok{total\_citations =} \ConstantTok{NA\_integer\_}\NormalTok{,}
      \AttributeTok{self\_citations =} \ConstantTok{NA\_integer\_}\NormalTok{,}
      \AttributeTok{self\_cite\_rate =} \ConstantTok{NA\_real\_}
\NormalTok{    ))}
\NormalTok{  \}}

  \CommentTok{\# Count references that point to the same author\textquotesingle{}s works}
  \CommentTok{\# We look at the referenced\_works and check which are by the same author}
\NormalTok{  total\_cites }\OtherTok{\textless{}{-}} \FunctionTok{sum}\NormalTok{(works}\SpecialCharTok{$}\NormalTok{cited\_by\_count, }\AttributeTok{na.rm =} \ConstantTok{TRUE}\NormalTok{)}

  \CommentTok{\# For self{-}citation, we check which of the works\textquotesingle{} referenced\_works}

  \CommentTok{\# are also in the author\textquotesingle{}s own works list}
\NormalTok{  own\_work\_ids }\OtherTok{\textless{}{-}}\NormalTok{ works}\SpecialCharTok{$}\NormalTok{id}

\NormalTok{  self\_cite\_count }\OtherTok{\textless{}{-}}\NormalTok{ works }\SpecialCharTok{|\textgreater{}}
    \FunctionTok{pull}\NormalTok{(referenced\_works) }\SpecialCharTok{|\textgreater{}}
    \FunctionTok{map\_int}\NormalTok{(}\SpecialCharTok{\textasciitilde{}} \FunctionTok{sum}\NormalTok{(.x }\SpecialCharTok{\%in\%}\NormalTok{ own\_work\_ids, }\AttributeTok{na.rm =} \ConstantTok{TRUE}\NormalTok{))}

  \FunctionTok{tibble}\NormalTok{(}
    \AttributeTok{author\_id =}\NormalTok{ author\_id,}
    \AttributeTok{author\_name =}\NormalTok{ author\_name,}
    \AttributeTok{n\_works =} \FunctionTok{nrow}\NormalTok{(works),}
    \AttributeTok{total\_references =} \FunctionTok{sum}\NormalTok{(}\FunctionTok{map\_int}\NormalTok{(works}\SpecialCharTok{$}\NormalTok{referenced\_works, length)),}
    \AttributeTok{self\_references =} \FunctionTok{sum}\NormalTok{(self\_cite\_count),}
    \AttributeTok{self\_cite\_rate =} \FunctionTok{if\_else}\NormalTok{(}
      \FunctionTok{sum}\NormalTok{(}\FunctionTok{map\_int}\NormalTok{(works}\SpecialCharTok{$}\NormalTok{referenced\_works, length)) }\SpecialCharTok{\textgreater{}} \DecValTok{0}\NormalTok{,}
      \FunctionTok{sum}\NormalTok{(self\_cite\_count) }\SpecialCharTok{/} \FunctionTok{sum}\NormalTok{(}\FunctionTok{map\_int}\NormalTok{(works}\SpecialCharTok{$}\NormalTok{referenced\_works, length)),}
      \ConstantTok{NA\_real\_}
\NormalTok{    )}
\NormalTok{  )}
\NormalTok{\}}

\NormalTok{self\_cite\_results }\OtherTok{\textless{}{-}} \FunctionTok{map2\_dfr}\NormalTok{(}
\NormalTok{  authors\_sample}\SpecialCharTok{$}\NormalTok{id,}
\NormalTok{  authors\_sample}\SpecialCharTok{$}\NormalTok{display\_name,}
\NormalTok{  estimate\_self\_citation}
\NormalTok{)}
\end{Highlighting}
\end{Shaded}

\subsection{Summarizing results}\label{summarizing-results}

\begin{Shaded}
\begin{Highlighting}[]
\NormalTok{self\_cite\_results }\SpecialCharTok{|\textgreater{}}
  \FunctionTok{filter}\NormalTok{(}\SpecialCharTok{!}\FunctionTok{is.na}\NormalTok{(self\_cite\_rate)) }\SpecialCharTok{|\textgreater{}}
  \FunctionTok{arrange}\NormalTok{(}\FunctionTok{desc}\NormalTok{(self\_cite\_rate)) }\SpecialCharTok{|\textgreater{}}
  \FunctionTok{mutate}\NormalTok{(}\AttributeTok{self\_cite\_rate =}\NormalTok{ scales}\SpecialCharTok{::}\FunctionTok{percent}\NormalTok{(self\_cite\_rate, }\AttributeTok{accuracy =} \FloatTok{0.1}\NormalTok{)) }\SpecialCharTok{|\textgreater{}}
  \FunctionTok{gt}\NormalTok{() }\SpecialCharTok{|\textgreater{}}
  \FunctionTok{tab\_header}\NormalTok{(}
    \AttributeTok{title =} \StringTok{"Self{-}citation rates among sampled scientometrics authors"}\NormalTok{,}
    \AttributeTok{subtitle =} \StringTok{"Based on references in up to 50 recent works per author"}
\NormalTok{  ) }\SpecialCharTok{|\textgreater{}}
  \FunctionTok{cols\_label}\NormalTok{(}
    \AttributeTok{author\_name =} \StringTok{"Author"}\NormalTok{,}
    \AttributeTok{n\_works =} \StringTok{"Works sampled"}\NormalTok{,}
    \AttributeTok{total\_references =} \StringTok{"Total references"}\NormalTok{,}
    \AttributeTok{self\_references =} \StringTok{"Self{-}references"}\NormalTok{,}
    \AttributeTok{self\_cite\_rate =} \StringTok{"Self{-}citation rate"}
\NormalTok{  ) }\SpecialCharTok{|\textgreater{}}
  \FunctionTok{cols\_hide}\NormalTok{(author\_id)}
\end{Highlighting}
\end{Shaded}

\begingroup
\fontsize{8.0pt}{9.0pt}\selectfont
\begin{longtable}{>{\raggedright\arraybackslash}p{\dimexpr 0.25\linewidth -2\tabcolsep-1.5\arrayrulewidth}>{\raggedleft\arraybackslash}p{\dimexpr 0.15\linewidth -2\tabcolsep-1.5\arrayrulewidth}>{\raggedleft\arraybackslash}p{\dimexpr 0.21\linewidth -2\tabcolsep-1.5\arrayrulewidth}>{\raggedleft\arraybackslash}p{\dimexpr 0.20\linewidth -2\tabcolsep-1.5\arrayrulewidth}>{\raggedleft\arraybackslash}p{\dimexpr 0.20\linewidth -2\tabcolsep-1.5\arrayrulewidth}}
\caption*{
{\fontsize{11}{14}\selectfont  Self-citation rates among sampled scientometrics authors\fontsize{8}{9}\selectfont } \\ 
{\fontsize{9}{11}\selectfont  Based on references in up to 50 recent works per author\fontsize{8}{9}\selectfont }
} \\ 
\toprule
Author & Works sampled & Total references & \nolinkurl{Self-references} & Self-citation rate \\ 
\midrule\addlinespace[2.5pt]
Bachir Zoudji & 50 & 1827 & 84 & 4.6\% \\ 
Haak Laurel & 22 & 208 & 6 & 2.9\% \\ 
Joaquín M. Azagra‐Caro & 26 & 1203 & 31 & 2.6\% \\ 
Yu‐Hsin Chang & 37 & 1269 & 30 & 2.4\% \\ 
Christina Bornatici & 19 & 231 & 5 & 2.2\% \\ 
Cynthia Ramasubbu & 21 & 244 & 3 & 1.2\% \\ 
Pietro Pollo & 50 & 2497 & 27 & 1.1\% \\ 
Eun‐Kyung Chung & 48 & 685 & 3 & 0.4\% \\ 
Babak Pourbohloul & 2 & 28 & 0 & 0.0\% \\ 
Magdalena Szuflita‐Żurawska & 8 & 93 & 0 & 0.0\% \\ 
\bottomrule
\end{longtable}
\endgroup

\subsection{Visualization}\label{visualization}

\begin{Shaded}
\begin{Highlighting}[]
\NormalTok{self\_cite\_clean }\OtherTok{\textless{}{-}}\NormalTok{ self\_cite\_results }\SpecialCharTok{|\textgreater{}}
  \FunctionTok{filter}\NormalTok{(}\SpecialCharTok{!}\FunctionTok{is.na}\NormalTok{(self\_cite\_rate))}

\ControlFlowTok{if}\NormalTok{ (}\FunctionTok{nrow}\NormalTok{(self\_cite\_clean) }\SpecialCharTok{\textgreater{}} \DecValTok{0}\NormalTok{) \{}
  \FunctionTok{ggplot}\NormalTok{(self\_cite\_clean, }\FunctionTok{aes}\NormalTok{(}\AttributeTok{x =}\NormalTok{ self\_cite\_rate)) }\SpecialCharTok{+}
    \FunctionTok{geom\_histogram}\NormalTok{(}
      \AttributeTok{binwidth =} \FloatTok{0.02}\NormalTok{,}
      \AttributeTok{fill =} \FunctionTok{palette\_sci}\NormalTok{(}\DecValTok{1}\NormalTok{),}
      \AttributeTok{colour =} \StringTok{"white"}
\NormalTok{    ) }\SpecialCharTok{+}
    \FunctionTok{geom\_vline}\NormalTok{(}
      \AttributeTok{xintercept =} \FunctionTok{median}\NormalTok{(self\_cite\_clean}\SpecialCharTok{$}\NormalTok{self\_cite\_rate, }\AttributeTok{na.rm =} \ConstantTok{TRUE}\NormalTok{),}
      \AttributeTok{linetype =} \StringTok{"dashed"}\NormalTok{, }\AttributeTok{linewidth =} \FloatTok{0.8}
\NormalTok{    ) }\SpecialCharTok{+}
    \FunctionTok{scale\_x\_continuous}\NormalTok{(}\AttributeTok{labels =}\NormalTok{ scales}\SpecialCharTok{::}\FunctionTok{percent\_format}\NormalTok{()) }\SpecialCharTok{+}
    \FunctionTok{labs}\NormalTok{(}
      \AttributeTok{x =} \StringTok{"Self{-}citation rate (proportion of references to own works)"}\NormalTok{,}
      \AttributeTok{y =} \StringTok{"Number of authors"}
\NormalTok{    ) }\SpecialCharTok{+}
    \FunctionTok{theme\_sci}\NormalTok{()}
\NormalTok{\}}
\end{Highlighting}
\end{Shaded}

\begin{figure}

{\centering \includegraphics[width=0.9\linewidth,alt={Histogram showing the distribution of self-citation rates for sampled authors. A dashed vertical line indicates the median self-citation rate.}]{scientometrics-in-r_files/figure-latex/03-fig-ethics-self-citation-1} 

}

\caption{Distribution of self-citation rates across sampled scientometrics authors. The dashed line marks the sample median.}\label{fig:03-fig-ethics-self-citation}
\end{figure}

\subsection{A per-author comparison}\label{a-per-author-comparison}

\begin{Shaded}
\begin{Highlighting}[]
\ControlFlowTok{if}\NormalTok{ (}\FunctionTok{nrow}\NormalTok{(self\_cite\_clean) }\SpecialCharTok{\textgreater{}} \DecValTok{0}\NormalTok{) \{}
\NormalTok{  self\_cite\_clean }\SpecialCharTok{|\textgreater{}}
    \FunctionTok{mutate}\NormalTok{(}\AttributeTok{author\_name =} \FunctionTok{fct\_reorder}\NormalTok{(author\_name, self\_cite\_rate)) }\SpecialCharTok{|\textgreater{}}
    \FunctionTok{ggplot}\NormalTok{(}\FunctionTok{aes}\NormalTok{(}\AttributeTok{x =}\NormalTok{ self\_cite\_rate, }\AttributeTok{y =}\NormalTok{ author\_name)) }\SpecialCharTok{+}
    \FunctionTok{geom\_col}\NormalTok{(}\AttributeTok{fill =} \FunctionTok{palette\_sci}\NormalTok{(}\DecValTok{1}\NormalTok{)) }\SpecialCharTok{+}
    \FunctionTok{scale\_x\_continuous}\NormalTok{(}\AttributeTok{labels =}\NormalTok{ scales}\SpecialCharTok{::}\FunctionTok{percent\_format}\NormalTok{()) }\SpecialCharTok{+}
    \FunctionTok{labs}\NormalTok{(}\AttributeTok{x =} \StringTok{"Self{-}citation rate"}\NormalTok{, }\AttributeTok{y =} \ConstantTok{NULL}\NormalTok{) }\SpecialCharTok{+}
    \FunctionTok{theme\_sci}\NormalTok{()}
\NormalTok{\}}
\end{Highlighting}
\end{Shaded}

\begin{figure}

{\centering \includegraphics[width=0.9\linewidth,alt={Horizontal bar chart showing each sampled author's self-citation rate, ordered from highest to lowest.}]{scientometrics-in-r_files/figure-latex/03-fig-ethics-self-cite-bar-1} 

}

\caption{Self-citation rates by author, sorted from highest to lowest.}\label{fig:03-fig-ethics-self-cite-bar}
\end{figure}

\section{Diagnostics and interpretation}\label{ethics-diagnostics}

Self-citation rates vary enormously across disciplines, career stages, and research specialities. When interpreting these results, keep the following in mind:

\begin{itemize}
\tightlist
\item
  \textbf{Baseline rates differ by field.} In small, specialised fields, self-citation is often necessary because few other researchers work on the same topic. Self-citation rates of 15--25\% are common in niche subfields and do not, by themselves, indicate manipulation.
\item
  \textbf{Career stage matters.} Senior researchers with large publication portfolios naturally accumulate more self-citations simply because they have more prior work to build on.
\item
  \textbf{Our estimate is a lower bound.} We only count self-references to works that appear in the author's sampled works from OpenAlex. Works outside our sample window (before 2018) or works not indexed by OpenAlex are missed, so the true self-citation rate is likely higher.
\item
  \textbf{Context is essential.} A self-citation rate of 30\% might be perfectly reasonable for a researcher building a coherent long-term programme, while a rate of 10\% could still be problematic if those citations are strategically placed to inflate a particular metric.
\item
  \textbf{Never use self-citation rates alone to accuse individuals of misconduct.} As the Leiden Manifesto insists, quantitative indicators must support --- not replace --- qualitative judgement \citep{hicks2015}.
\end{itemize}

\section{Limitations and responsible use}\label{ethics-limitations}

\section{Limitations and responsible use}\label{limitations-and-responsible-use-2}

\begin{itemize}
\tightlist
\item
  \textbf{Self-citation analysis is not misconduct detection.} High self-citation rates have many legitimate explanations. Presenting a self-citation analysis as evidence of gaming without qualitative context is itself a form of metric misuse.
\item
  \textbf{Coverage gaps matter.} OpenAlex does not index all scholarly works. Authors who publish in venues not well covered by OpenAlex will have incomplete reference lists, biasing self-citation estimates.
\item
  \textbf{Author disambiguation is imperfect.} OpenAlex uses algorithmic author disambiguation. Common names may be conflated (merging two people into one profile) or split (one person appearing as multiple profiles), distorting individual-level statistics.
\item
  \textbf{Snapshot bias.} Our analysis uses a sample of works from a particular time window. Conclusions should not be generalised beyond that window without further analysis.
\item
  \textbf{The Leiden Manifesto applies here too.} Every principle in \ref{ethics-background} applies to \emph{our own} analysis. We must be transparent about methods, allow those evaluated to verify data, and avoid false precision \citep{hicks2015}.
\end{itemize}

\section{Common pitfalls}\label{ethics-pitfalls}

\section{Common pitfalls}\label{common-pitfalls-2}

\begin{itemize}
\tightlist
\item
  \textbf{Equating self-citation with misconduct.} Self-citation is a normal part of scientific writing. Only excessive, strategically motivated self-citation is problematic --- and even then, the boundary is context-dependent.
\item
  \textbf{Using the Journal Impact Factor to judge individual articles.} This is precisely what DORA warns against. A journal's average citation rate says nothing about any single paper in that journal \citep{dora2012}.
\item
  \textbf{Ignoring field differences when comparing metrics.} A citation count of 20 is exceptional in mathematics but unremarkable in biomedical research. Any cross-field comparison requires field normalisation \citep{waltman2016}.
\item
  \textbf{Treating indicators as objective truth.} All bibliometric indicators are constructed measures with specific assumptions, coverage limitations, and known biases. Present them as approximations, not facts.
\item
  \textbf{Forgetting Goodhart's Law.} If you propose a new indicator for evaluation purposes, consider how rational actors will respond to being measured by it. Every metric creates incentives, and those incentives may not align with the goals of good research \citep{goodhart1984}.
\end{itemize}

\section{Exercises}\label{ethics-exercises}

\begin{enumerate}
\def\labelenumi{\arabic{enumi}.}
\item
  \textbf{Leiden Manifesto audit.} Select a recent hiring or promotion policy document from a university (many are publicly available). Evaluate it against all ten Leiden Manifesto principles. Which principles does it satisfy? Which does it violate? Write a one-page assessment.
\item
  \textbf{Self-citation across fields.} Modify the worked example to compare self-citation rates between two different research fields (e.g., ``bibliometrics'' and ``genomics''). Do you observe different baseline self-citation rates? What might explain the difference?
\item
  \textbf{DORA compliance check.} Visit the DORA website (\href{https://sfdora.org}{sfdora.org}) and identify three major funders that have signed the declaration. For each, find their current grant evaluation criteria and assess whether their stated practices are consistent with DORA commitments.
\item
  \textbf{Goodhart's Law thought experiment.} Propose a new indicator for measuring research quality that is not currently in wide use. Then, write a brief analysis of how researchers might game this indicator if it were adopted for promotion decisions. \emph{(Hint: think about both the numerator and denominator of any ratio-based metric.)}
\end{enumerate}

\section{Solutions}\label{03-solutions}

Solutions are provided in \ref{solutions}.

\section{Further reading}\label{ethics-reading}

\begin{itemize}
\tightlist
\item
  \citet{hicks2015} --- The Leiden Manifesto for Research Metrics. The single most important reference on responsible use of bibliometric indicators.
\item
  \citet{dora2012} --- The San Francisco Declaration on Research Assessment. The landmark declaration against using journal-level metrics to evaluate individual researchers.
\item
  \citet{wilsdon2015} --- \emph{The Metric Tide}. A comprehensive UK government review of the role of metrics in research assessment, proposing a framework of responsible metrics.
\item
  \citet{goodhart1984} --- Goodhart's original articulation of the principle that a measure ceases to be useful when it becomes a target.
\item
  \citet{waltman2016} --- A thorough review of citation impact indicators, including their known biases and appropriate uses.
\item
  \citet{merton1968} --- The Matthew Effect in science; describes how advantage accumulates, a dynamic that citation metrics can reinforce.
\item
  \citet{garfield1955} --- The foundational paper on citation indexing. Essential for understanding the origins of the system we now critique.
\end{itemize}

\section{Session info}\label{ethics-sessioninfo}

\begin{Shaded}
\begin{Highlighting}[]
\FunctionTok{sessionInfo}\NormalTok{()}
\end{Highlighting}
\end{Shaded}

\begin{verbatim}
#> R version 4.6.1 (2026-06-24)
#> Platform: x86_64-pc-linux-gnu
#> Running under: Ubuntu 24.04.5 LTS
#> 
#> Matrix products: default
#> BLAS:   /usr/lib/x86_64-linux-gnu/openblas-pthread/libblas.so.3 
#> LAPACK: /usr/lib/x86_64-linux-gnu/openblas-pthread/libopenblasp-r0.3.26.so;  LAPACK version 3.12.0
#> 
#> locale:
#>  [1] LC_CTYPE=C.UTF-8       LC_NUMERIC=C           LC_TIME=C.UTF-8       
#>  [4] LC_COLLATE=C.UTF-8     LC_MONETARY=C.UTF-8    LC_MESSAGES=C.UTF-8   
#>  [7] LC_PAPER=C.UTF-8       LC_NAME=C              LC_ADDRESS=C          
#> [10] LC_TELEPHONE=C         LC_MEASUREMENT=C.UTF-8 LC_IDENTIFICATION=C   
#> 
#> time zone: UTC
#> tzcode source: system (glibc)
#> 
#> attached base packages:
#> [1] stats     graphics  grDevices utils     datasets  methods   base     
#> 
#> other attached packages:
#>  [1] gt_1.4.0        tidytext_0.4.3  glue_1.8.1      openalexR_3.1.2
#>  [5] lubridate_1.9.5 forcats_1.0.1   stringr_1.6.0   dplyr_1.2.1    
#>  [9] purrr_1.2.2     readr_2.2.0     tidyr_1.3.2     tibble_3.3.1   
#> [13] ggplot2_4.0.3   tidyverse_2.0.0
#> 
#> loaded via a namespace (and not attached):
#>  [1] janeaustenr_1.0.0  viridis_0.6.5      utf8_1.2.6         generics_0.1.4    
#>  [5] xml2_1.6.0         lattice_0.22-9     stringi_1.8.9      hms_1.1.4         
#>  [9] digest_0.6.39      magrittr_2.0.5     evaluate_1.0.5     grid_4.6.1        
#> [13] timechange_0.4.0   RColorBrewer_1.1-3 bookdown_0.48      fastmap_1.2.0     
#> [17] Matrix_1.7-5       rprojroot_2.1.1    gridExtra_2.3.1    stopwords_2.3     
#> [21] viridisLite_0.4.3  scales_1.4.0       cli_3.6.6          rlang_1.3.0       
#> [25] tokenizers_0.3.0   withr_3.0.3        yaml_2.3.12        otel_0.2.0        
#> [29] tools_4.6.1        tzdb_0.5.0         here_1.0.2         vctrs_0.7.3       
#> [33] R6_2.6.1           lifecycle_1.0.5    fs_2.1.0           pkgconfig_2.0.3   
#> [37] pillar_1.11.1      gtable_0.3.6       Rcpp_1.1.2         systemfonts_1.3.2 
#> [41] xfun_0.62          tidyselect_1.2.1   rstudioapi_0.19.0  knitr_1.52        
#> [45] farver_2.1.2       SnowballC_0.7.1    htmltools_0.5.9    labeling_0.4.3    
#> [49] rmarkdown_2.32     compiler_4.6.1     S7_0.2.2
\end{verbatim}

\chapter{Data Sources Compared}\label{data-sources-compared}

\emph{R. Heller · Date: 2026-05-09 · Last changed: 2026-10-31}

\section{Learning objectives}\label{datasources-objectives}

After completing this chapter, you will be able to:

\begin{itemize}
\tightlist
\item
  Describe the coverage, access model, and key strengths of twelve major bibliographic data sources
\item
  Select an appropriate data source (or combination of sources) for a given research question
\item
  Compare metadata completeness across sources for the same set of publications
\item
  Evaluate the trade-offs between proprietary and open data sources for reproducible research
\item
  Retrieve and compare metadata from OpenAlex and Crossref using R
\end{itemize}

\section{Setup}\label{datasources-setup}

\begin{Shaded}
\begin{Highlighting}[]
\FunctionTok{library}\NormalTok{(tidyverse)}
\FunctionTok{library}\NormalTok{(openalexR)}
\FunctionTok{library}\NormalTok{(rcrossref)}
\FunctionTok{library}\NormalTok{(glue)}
\FunctionTok{library}\NormalTok{(gt)}

\FunctionTok{set.seed}\NormalTok{(}\DecValTok{20260509}\NormalTok{)}

\FunctionTok{source}\NormalTok{(here}\SpecialCharTok{::}\FunctionTok{here}\NormalTok{(}\StringTok{"R"}\NormalTok{, }\StringTok{"api\_helpers.R"}\NormalTok{))}
\FunctionTok{source}\NormalTok{(here}\SpecialCharTok{::}\FunctionTok{here}\NormalTok{(}\StringTok{"R"}\NormalTok{, }\StringTok{"utils.R"}\NormalTok{))}
\FunctionTok{source}\NormalTok{(here}\SpecialCharTok{::}\FunctionTok{here}\NormalTok{(}\StringTok{"R"}\NormalTok{, }\StringTok{"sci\_palette.R"}\NormalTok{))}
\end{Highlighting}
\end{Shaded}

\section{Conceptual background}\label{datasources-background}

Every bibliometric analysis is fundamentally constrained by the data source that feeds it. The database you choose determines which publications are visible, which citations are counted, and which metadata fields are available. As \citet{visser2021} demonstrated in a large-scale comparison, the major bibliographic databases differ substantially in coverage, with these differences affecting research outcomes in measurable ways. Similarly, \citet{singh2021} showed that journal coverage varies considerably across Web of Science, Scopus, and Dimensions, with each database exhibiting distinct disciplinary strengths.

The choice of data source is not merely a technical decision --- it is a methodological one with direct consequences for the validity of results. A citation analysis conducted on Web of Science will produce different numbers than the same analysis on Scopus, Dimensions, or OpenAlex, because each database indexes a different set of publications and computes citations over a different corpus. Researchers must therefore understand these differences and report them transparently, as emphasised by the Leiden Manifesto's fourth principle: keep data collection and analytical processes open, transparent, and simple \citep{hicks2015}.

The landscape of bibliographic data sources has shifted dramatically in recent years. The traditional duopoly of Web of Science and Scopus has been challenged by open alternatives --- most notably OpenAlex \citep{priem2022}, which provides free, unrestricted access to over 250 million scholarly works. Crossref serves as the backbone registration agency for DOIs and provides rich metadata for over 150 million records. Discipline-specific databases like PubMed, ADS, and arXiv serve targeted communities with deep, curated coverage.

For this book, we adopt OpenAlex as the primary data source because it is fully open, freely accessible via API without authentication, and provides coverage comparable to proprietary alternatives. However, responsible bibliometric practice often requires consulting multiple sources to cross-validate findings and fill coverage gaps.

\subsection{Twelve data sources at a glance}\label{twelve-data-sources-at-a-glance}

Below we summarise twelve major sources. The subsequent sections examine each in detail, and the comparison table provides a structured overview.

\textbf{Web of Science (WoS).} Maintained by Clarivate, WoS is the oldest and most established citation index, originating from Eugene Garfield's Science Citation Index \citep{garfield1955}. It covers approximately 90 million records from over 21,000 journals, with citation data back to 1900. Access requires an institutional subscription. Its strengths are curated journal selection, long historical coverage, and the Journal Citation Reports (JCR). Its main limitations are cost, restricted API access, and lower coverage of non-English literature, social sciences, and humanities.

\textbf{Scopus.} Elsevier's Scopus indexes over 94 million records from approximately 27,000 journals, with broader coverage of non-English titles and conference proceedings than WoS. It provides its own citation metrics (CiteScore, SJR, SNIP) and a well-documented API. Like WoS, it requires an institutional subscription.

\textbf{Dimensions.} Operated by Digital Science, Dimensions indexes over 140 million publications including journal articles, books, book chapters, preprints, clinical trials, patents, and policy documents. It offers a free tier with limited functionality and a paid tier with full API access. Its distinctive strength is linking publications to grants, patents, clinical trials, and policy documents.

\textbf{OpenAlex.} The successor to the discontinued Microsoft Academic Graph, OpenAlex is a fully open, freely accessible index of over 250 million scholarly works \citep{priem2022}. It aggregates data from Crossref, PubMed, institutional repositories, and other open sources. Its API requires no authentication (though a polite-pool email is recommended). OpenAlex provides author disambiguation, institutional affiliation data, and concept tagging. Its primary limitations are metadata quality (algorithmically parsed, not manually curated) and less granular subject classification than curated databases.

\textbf{The Lens.} Operated by Cambia, The Lens provides free access to over 300 million scholarly works and 160 million patent records. It aggregates data from Crossref, PubMed, Microsoft Academic, and patent offices. Its unique strength is integrated patent-paper linkage, making it valuable for technology-transfer and innovation studies. API access is available via a free institutional plan.

\textbf{Crossref.} Crossref is the DOI registration agency for scholarly publications, covering over 150 million records. It is not a citation index per se --- it provides metadata deposited by publishers, including titles, authors, abstracts, reference lists, and licensing information. Its REST API is freely accessible. Crossref's strengths are comprehensive DOI coverage and detailed reference metadata; its limitation is that metadata quality depends entirely on what publishers deposit.

\textbf{PubMed/MEDLINE.} Maintained by the US National Library of Medicine, PubMed indexes over 36 million records in the biomedical and life sciences. Its key strengths are curated MeSH (Medical Subject Headings) indexing, structured abstracts, and deep coverage of clinical and biomedical literature. PubMed does not provide its own citation counts but can be paired with citation data from other sources.

\textbf{Google Scholar.} Google Scholar provides the broadest coverage of any scholarly search engine, indexing journal articles, conference papers, theses, books, patents, and grey literature. However, it provides no API, no bulk data export, and no structured metadata. It is useful for discovery and for its ``Cited by'' feature but is not suitable for systematic bibliometric analysis.

\textbf{Semantic Scholar.} Operated by the Allen Institute for AI, Semantic Scholar indexes over 200 million papers and provides a free API with rich metadata including algorithmically extracted entities, TLDR summaries, and citation context. It is particularly strong in computer science, biomedical, and neuroscience literature.

\textbf{NASA/ADS (Astrophysics Data System).} ADS is a curated database of over 16 million records in astronomy, astrophysics, and physics. It provides excellent citation tracking, full-text search, and a well-documented API. For astronomical research, ADS is the gold standard.

\textbf{arXiv.} The original preprint server, arXiv hosts over 2.5 million preprints in physics, mathematics, computer science, quantitative biology, quantitative finance, statistics, and electrical engineering. It does not provide citation data directly but its records are indexed by other databases. Its API provides metadata access.

\textbf{bioRxiv/medRxiv.} The preprint servers for biology and health sciences, hosting over 300,000 preprints. They provide API access to metadata and full text and are indexed by OpenAlex, Dimensions, and other aggregators.

\section{Worked example}\label{datasources-example}

We now compare what OpenAlex and Crossref return for the same set of DOIs. This is a practical illustration of how metadata completeness varies across sources.

\subsection{Selecting DOIs}\label{selecting-dois}

We start with a small set of well-known scientometrics papers whose DOIs we know.

\begin{Shaded}
\begin{Highlighting}[]
\NormalTok{dois }\OtherTok{\textless{}{-}} \FunctionTok{c}\NormalTok{(}
  \StringTok{"10.1038/520429a"}\NormalTok{,            }\CommentTok{\# Hicks et al. 2015, Leiden Manifesto}
  \StringTok{"10.1073/pnas.0507655102"}\NormalTok{,    }\CommentTok{\# Hirsch 2005, h{-}index}
  \StringTok{"10.1007/s11192{-}006{-}0144{-}7"}\NormalTok{,  }\CommentTok{\# Egghe 2006, g{-}index}
  \StringTok{"10.1016/j.joi.2017.08.007"}\NormalTok{,  }\CommentTok{\# Aria \& Cuccurullo 2017, bibliometrix}
  \StringTok{"10.1126/science.122.3159.108"}\NormalTok{, }\CommentTok{\# Garfield 1955, citation indexes}
  \StringTok{"10.48550/arXiv.2205.01833"}\NormalTok{,  }\CommentTok{\# Priem et al. 2022, OpenAlex}
  \StringTok{"10.1162/qss\_a\_00112"}\NormalTok{,        }\CommentTok{\# Visser et al. 2021, database comparison}
  \StringTok{"10.1007/s11192{-}021{-}03948{-}5"}  \CommentTok{\# Singh et al. 2021, journal coverage}
\NormalTok{)}
\end{Highlighting}
\end{Shaded}

\subsection{Fetching from OpenAlex}\label{fetching-from-openalex}

\begin{Shaded}
\begin{Highlighting}[]
\NormalTok{oa\_results }\OtherTok{\textless{}{-}} \FunctionTok{fetch\_openalex}\NormalTok{(}
  \AttributeTok{entity =} \StringTok{"works"}\NormalTok{,}
  \AttributeTok{doi =}\NormalTok{ dois}
\NormalTok{)}

\NormalTok{oa\_meta }\OtherTok{\textless{}{-}}\NormalTok{ oa\_results }\SpecialCharTok{|\textgreater{}}
  \FunctionTok{transmute}\NormalTok{(}
    \AttributeTok{doi =}\NormalTok{ doi,}
    \AttributeTok{source =} \StringTok{"OpenAlex"}\NormalTok{,}
    \AttributeTok{title =}\NormalTok{ display\_name,}
    \AttributeTok{year =}\NormalTok{ publication\_year,}
    \AttributeTok{cited\_by =}\NormalTok{ cited\_by\_count,}
    \AttributeTok{has\_abstract =} \SpecialCharTok{!}\FunctionTok{is.na}\NormalTok{(abstract),}
    \AttributeTok{type =}\NormalTok{ type}
\NormalTok{  )}

\FunctionTok{glimpse}\NormalTok{(oa\_meta)}
\end{Highlighting}
\end{Shaded}

\begin{verbatim}
#> Rows: 8
#> Columns: 7
#> $ doi          <chr> "https://doi.org/10.1016/j.joi.2017.08.007", "https://doi~
#> $ source       <chr> "OpenAlex", "OpenAlex", "OpenAlex", "OpenAlex", "OpenAlex~
#> $ title        <chr> "bibliometrix : An R-tool for comprehensive science mappi~
#> $ year         <int> 2017, 2005, 2015, 2021, 2006, 1955, 2021, 2022
#> $ cited_by     <int> 19864, 12131, 3282, 2541, 2463, 2440, 1108, 1252
#> $ has_abstract <lgl> FALSE, TRUE, FALSE, FALSE, FALSE, TRUE, TRUE, TRUE
#> $ type         <chr> "article", "article", "article", "article", "article", "a~
\end{verbatim}

\subsection{Fetching from Crossref}\label{fetching-from-crossref}

\begin{Shaded}
\begin{Highlighting}[]
\NormalTok{cr\_results }\OtherTok{\textless{}{-}} \FunctionTok{cr\_works}\NormalTok{(}\AttributeTok{dois =}\NormalTok{ dois)}

\NormalTok{cr\_meta }\OtherTok{\textless{}{-}}\NormalTok{ cr\_results}\SpecialCharTok{$}\NormalTok{data }\SpecialCharTok{|\textgreater{}}
  \FunctionTok{transmute}\NormalTok{(}
    \AttributeTok{doi =}\NormalTok{ doi,}
    \AttributeTok{source =} \StringTok{"Crossref"}\NormalTok{,}
    \AttributeTok{title =}\NormalTok{ title,}
    \AttributeTok{year =} \FunctionTok{as.integer}\NormalTok{(}
      \FunctionTok{coalesce}\NormalTok{(}
        \FunctionTok{substr}\NormalTok{(issued, }\DecValTok{1}\NormalTok{, }\DecValTok{4}\NormalTok{),}
        \FunctionTok{substr}\NormalTok{(deposited, }\DecValTok{1}\NormalTok{, }\DecValTok{4}\NormalTok{)}
\NormalTok{      )}
\NormalTok{    ),}
    \AttributeTok{cited\_by =} \FunctionTok{as.integer}\NormalTok{(is.referenced.by.count),}
    \AttributeTok{has\_abstract =} \SpecialCharTok{!}\FunctionTok{is.na}\NormalTok{(abstract),}
    \AttributeTok{type =}\NormalTok{ type}
\NormalTok{  )}

\FunctionTok{glimpse}\NormalTok{(cr\_meta)}
\end{Highlighting}
\end{Shaded}

\begin{verbatim}
#> Rows: 7
#> Columns: 7
#> $ doi          <chr> "10.1038/520429a", "10.1073/pnas.0507655102", "10.1007/s1~
#> $ source       <chr> "Crossref", "Crossref", "Crossref", "Crossref", "Crossref~
#> $ title        <chr> "Bibliometrics: The Leiden Manifesto for research metrics~
#> $ year         <int> 2015, 2005, 2006, 2017, 1955, 2021, 2021
#> $ cited_by     <int> 2303, 8935, 1926, 14716, 1843, 814, 1955
#> $ has_abstract <lgl> FALSE, TRUE, FALSE, FALSE, FALSE, TRUE, FALSE
#> $ type         <chr> "journal-article", "journal-article", "journal-article", ~
\end{verbatim}

\subsection{Comparing metadata completeness}\label{comparing-metadata-completeness}

\begin{Shaded}
\begin{Highlighting}[]
\NormalTok{comparison }\OtherTok{\textless{}{-}} \FunctionTok{bind\_rows}\NormalTok{(oa\_meta, cr\_meta) }\SpecialCharTok{|\textgreater{}}
  \FunctionTok{mutate}\NormalTok{(}\AttributeTok{short\_doi =} \FunctionTok{str\_extract}\NormalTok{(doi, }\StringTok{"[\^{}/]+/[\^{}/]+$"}\NormalTok{))}

\NormalTok{comparison\_wide }\OtherTok{\textless{}{-}}\NormalTok{ comparison }\SpecialCharTok{|\textgreater{}}
  \FunctionTok{select}\NormalTok{(short\_doi, source, year, cited\_by, has\_abstract, type) }\SpecialCharTok{|\textgreater{}}
  \FunctionTok{pivot\_wider}\NormalTok{(}
    \AttributeTok{names\_from =}\NormalTok{ source,}
    \AttributeTok{values\_from =} \FunctionTok{c}\NormalTok{(year, cited\_by, has\_abstract, type),}
    \AttributeTok{names\_glue =} \StringTok{"\{source\}\_\{.value\}"}
\NormalTok{  )}

\NormalTok{comparison\_wide }\SpecialCharTok{|\textgreater{}}
  \FunctionTok{gt}\NormalTok{() }\SpecialCharTok{|\textgreater{}}
  \FunctionTok{tab\_header}\NormalTok{(}
    \AttributeTok{title =} \StringTok{"Metadata comparison: OpenAlex vs. Crossref"}\NormalTok{,}
    \AttributeTok{subtitle =} \StringTok{"Same DOIs queried from both sources"}
\NormalTok{  ) }\SpecialCharTok{|\textgreater{}}
  \FunctionTok{cols\_label}\NormalTok{(}
    \AttributeTok{short\_doi =} \StringTok{"DOI (short)"}\NormalTok{,}
    \AttributeTok{OpenAlex\_year =} \StringTok{"OA Year"}\NormalTok{,}
    \AttributeTok{Crossref\_year =} \StringTok{"CR Year"}\NormalTok{,}
    \AttributeTok{OpenAlex\_cited\_by =} \StringTok{"OA Citations"}\NormalTok{,}
    \AttributeTok{Crossref\_cited\_by =} \StringTok{"CR Citations"}\NormalTok{,}
    \AttributeTok{OpenAlex\_has\_abstract =} \StringTok{"OA Abstract?"}\NormalTok{,}
    \AttributeTok{Crossref\_has\_abstract =} \StringTok{"CR Abstract?"}\NormalTok{,}
    \AttributeTok{OpenAlex\_type =} \StringTok{"OA Type"}\NormalTok{,}
    \AttributeTok{Crossref\_type =} \StringTok{"CR Type"}
\NormalTok{  )}
\end{Highlighting}
\end{Shaded}

\begingroup
\fontsize{8.0pt}{9.0pt}\selectfont
\begin{longtable}{>{\raggedright\arraybackslash}p{\dimexpr 0.14\linewidth -2\tabcolsep-1.5\arrayrulewidth}>{\raggedleft\arraybackslash}p{\dimexpr 0.10\linewidth -2\tabcolsep-1.5\arrayrulewidth}>{\raggedleft\arraybackslash}p{\dimexpr 0.10\linewidth -2\tabcolsep-1.5\arrayrulewidth}>{\raggedleft\arraybackslash}p{\dimexpr 0.11\linewidth -2\tabcolsep-1.5\arrayrulewidth}>{\raggedleft\arraybackslash}p{\dimexpr 0.11\linewidth -2\tabcolsep-1.5\arrayrulewidth}>{\centering\arraybackslash}p{\dimexpr 0.12\linewidth -2\tabcolsep-1.5\arrayrulewidth}>{\centering\arraybackslash}p{\dimexpr 0.12\linewidth -2\tabcolsep-1.5\arrayrulewidth}>{\raggedright\arraybackslash}p{\dimexpr 0.10\linewidth -2\tabcolsep-1.5\arrayrulewidth}>{\raggedright\arraybackslash}p{\dimexpr 0.10\linewidth -2\tabcolsep-1.5\arrayrulewidth}}
\caption*{
{\fontsize{11}{14}\selectfont  Metadata comparison: OpenAlex vs. Crossref\fontsize{8}{9}\selectfont } \\ 
{\fontsize{9}{11}\selectfont  Same DOIs queried from both sources\fontsize{8}{9}\selectfont }
} \\ 
\toprule
DOI (short) & OA Year & CR Year & OA Citations & CR Citations & OA Abstract? & CR Abstract? & OA Type & CR Type \\ 
\midrule\addlinespace[2.5pt]
\nolinkurl{10.1016/j.joi.2017.08.007} & 2017 & 2017 & 19864 & 14716 & FALSE & FALSE & article & \nolinkurl{journal-article} \\ 
\nolinkurl{10.1073/pnas.0507655102} & 2005 & 2005 & 12131 & 8935 & TRUE & TRUE & article & \nolinkurl{journal-article} \\ 
\nolinkurl{10.1038/520429a} & 2015 & 2015 & 3282 & 2303 & FALSE & FALSE & article & \nolinkurl{journal-article} \\ 
\nolinkurl{10.1007/s11192-021-03948-5} & 2021 & 2021 & 2541 & 1955 & FALSE & FALSE & article & \nolinkurl{journal-article} \\ 
\nolinkurl{10.1007/s11192-006-0144-7} & 2006 & 2006 & 2463 & 1926 & FALSE & FALSE & article & \nolinkurl{journal-article} \\ 
\nolinkurl{10.1126/science.122.3159.108} & 1955 & 1955 & 2440 & 1843 & TRUE & FALSE & article & \nolinkurl{journal-article} \\ 
\nolinkurl{10.1162/qss_a_00112} & 2021 & 2021 & 1108 & 814 & TRUE & TRUE & article & \nolinkurl{journal-article} \\ 
\nolinkurl{10.48550/arxiv.2205.01833} & 2022 & NA & 1252 & NA & TRUE & NA & \nolinkurl{preprint} & NA \\ 
\bottomrule
\end{longtable}
\endgroup

\subsection{Visualization: citation count comparison}\label{visualization-citation-count-comparison}

\begin{Shaded}
\begin{Highlighting}[]
\NormalTok{cite\_compare }\OtherTok{\textless{}{-}}\NormalTok{ comparison }\SpecialCharTok{|\textgreater{}}
  \FunctionTok{select}\NormalTok{(short\_doi, source, cited\_by) }\SpecialCharTok{|\textgreater{}}
  \FunctionTok{pivot\_wider}\NormalTok{(}\AttributeTok{names\_from =}\NormalTok{ source, }\AttributeTok{values\_from =}\NormalTok{ cited\_by, }\AttributeTok{values\_fn =}\NormalTok{ first)}

\ControlFlowTok{if}\NormalTok{ (}\FunctionTok{nrow}\NormalTok{(cite\_compare) }\SpecialCharTok{\textgreater{}} \DecValTok{0} \SpecialCharTok{\&\&} \FunctionTok{all}\NormalTok{(}\FunctionTok{c}\NormalTok{(}\StringTok{"OpenAlex"}\NormalTok{, }\StringTok{"Crossref"}\NormalTok{) }\SpecialCharTok{\%in\%} \FunctionTok{names}\NormalTok{(cite\_compare))) \{}
\NormalTok{  max\_cite }\OtherTok{\textless{}{-}} \FunctionTok{max}\NormalTok{(}\FunctionTok{c}\NormalTok{(cite\_compare}\SpecialCharTok{$}\NormalTok{OpenAlex, cite\_compare}\SpecialCharTok{$}\NormalTok{Crossref), }\AttributeTok{na.rm =} \ConstantTok{TRUE}\NormalTok{)}

  \FunctionTok{ggplot}\NormalTok{(cite\_compare, }\FunctionTok{aes}\NormalTok{(}\AttributeTok{x =}\NormalTok{ Crossref, }\AttributeTok{y =}\NormalTok{ OpenAlex)) }\SpecialCharTok{+}
    \FunctionTok{geom\_abline}\NormalTok{(}\AttributeTok{slope =} \DecValTok{1}\NormalTok{, }\AttributeTok{intercept =} \DecValTok{0}\NormalTok{, }\AttributeTok{linetype =} \StringTok{"dashed"}\NormalTok{, }\AttributeTok{colour =} \StringTok{"grey50"}\NormalTok{) }\SpecialCharTok{+}
    \FunctionTok{geom\_point}\NormalTok{(}\AttributeTok{size =} \DecValTok{3}\NormalTok{, }\AttributeTok{colour =} \FunctionTok{palette\_sci}\NormalTok{(}\DecValTok{1}\NormalTok{)) }\SpecialCharTok{+}
    \FunctionTok{geom\_text}\NormalTok{(}
      \FunctionTok{aes}\NormalTok{(}\AttributeTok{label =}\NormalTok{ short\_doi),}
      \AttributeTok{size =} \FloatTok{2.5}\NormalTok{, }\AttributeTok{nudge\_y =}\NormalTok{ max\_cite }\SpecialCharTok{*} \FloatTok{0.03}\NormalTok{, }\AttributeTok{check\_overlap =} \ConstantTok{TRUE}
\NormalTok{    ) }\SpecialCharTok{+}
    \FunctionTok{scale\_x\_continuous}\NormalTok{(}\AttributeTok{labels =}\NormalTok{ scales}\SpecialCharTok{::}\FunctionTok{comma\_format}\NormalTok{()) }\SpecialCharTok{+}
    \FunctionTok{scale\_y\_continuous}\NormalTok{(}\AttributeTok{labels =}\NormalTok{ scales}\SpecialCharTok{::}\FunctionTok{comma\_format}\NormalTok{()) }\SpecialCharTok{+}
    \FunctionTok{coord\_equal}\NormalTok{(}\AttributeTok{xlim =} \FunctionTok{c}\NormalTok{(}\DecValTok{0}\NormalTok{, max\_cite }\SpecialCharTok{*} \FloatTok{1.1}\NormalTok{), }\AttributeTok{ylim =} \FunctionTok{c}\NormalTok{(}\DecValTok{0}\NormalTok{, max\_cite }\SpecialCharTok{*} \FloatTok{1.1}\NormalTok{)) }\SpecialCharTok{+}
    \FunctionTok{labs}\NormalTok{(}\AttributeTok{x =} \StringTok{"Crossref citation count"}\NormalTok{, }\AttributeTok{y =} \StringTok{"OpenAlex citation count"}\NormalTok{) }\SpecialCharTok{+}
    \FunctionTok{theme\_sci}\NormalTok{()}
\NormalTok{\}}
\end{Highlighting}
\end{Shaded}

\begin{figure}

{\centering \includegraphics[width=0.9\linewidth,alt={Scatter plot comparing citation counts from OpenAlex (y-axis) and Crossref (x-axis) for eight papers. A dashed diagonal line marks equality.}]{scientometrics-in-r_files/figure-latex/04-fig-datasources-citation-comparison-1} 

}

\caption{Citation counts reported by OpenAlex vs. Crossref for the same set of DOIs. Points above the diagonal indicate higher counts in OpenAlex.}\label{fig:04-fig-datasources-citation-comparison}
\end{figure}

\subsection{Summary comparison table}\label{summary-comparison-table}

\begin{Shaded}
\begin{Highlighting}[]
\FunctionTok{tribble}\NormalTok{(}
  \SpecialCharTok{\textasciitilde{}}\NormalTok{Source, }\SpecialCharTok{\textasciitilde{}}\NormalTok{Coverage, }\SpecialCharTok{\textasciitilde{}}\StringTok{\textasciigrave{}}\AttributeTok{Access model}\StringTok{\textasciigrave{}}\NormalTok{, }\SpecialCharTok{\textasciitilde{}}\NormalTok{API, }\SpecialCharTok{\textasciitilde{}}\StringTok{\textasciigrave{}}\AttributeTok{Citation data}\StringTok{\textasciigrave{}}\NormalTok{, }\SpecialCharTok{\textasciitilde{}}\NormalTok{Strengths,}
  \StringTok{"Web of Science"}\NormalTok{, }\StringTok{"\textasciitilde{}90M records, 21K+ journals"}\NormalTok{, }\StringTok{"Subscription"}\NormalTok{, }\StringTok{"Yes (paid)"}\NormalTok{, }\StringTok{"Yes (complete)"}\NormalTok{, }\StringTok{"Longest history, curated selection, JCR"}\NormalTok{,}
  \StringTok{"Scopus"}\NormalTok{, }\StringTok{"\textasciitilde{}94M records, 27K+ journals"}\NormalTok{, }\StringTok{"Subscription"}\NormalTok{, }\StringTok{"Yes (paid)"}\NormalTok{, }\StringTok{"Yes (complete)"}\NormalTok{, }\StringTok{"Broad journal coverage, CiteScore/SJR/SNIP"}\NormalTok{,}
  \StringTok{"Dimensions"}\NormalTok{, }\StringTok{"\textasciitilde{}140M records"}\NormalTok{, }\StringTok{"Freemium"}\NormalTok{, }\StringTok{"Yes (free/paid)"}\NormalTok{, }\StringTok{"Yes"}\NormalTok{, }\StringTok{"Links pubs to grants, patents, clinical trials"}\NormalTok{,}
  \StringTok{"OpenAlex"}\NormalTok{, }\StringTok{"\textasciitilde{}250M records"}\NormalTok{, }\StringTok{"Fully open"}\NormalTok{, }\StringTok{"Yes (free)"}\NormalTok{, }\StringTok{"Yes"}\NormalTok{, }\StringTok{"Open, broad coverage, author disambiguation"}\NormalTok{,}
  \StringTok{"The Lens"}\NormalTok{, }\StringTok{"\textasciitilde{}300M scholarly + 160M patent"}\NormalTok{, }\StringTok{"Free (institutional)"}\NormalTok{, }\StringTok{"Yes (free)"}\NormalTok{, }\StringTok{"Partial"}\NormalTok{, }\StringTok{"Patent{-}paper linkage"}\NormalTok{,}
  \StringTok{"Crossref"}\NormalTok{, }\StringTok{"\textasciitilde{}150M records"}\NormalTok{, }\StringTok{"Open"}\NormalTok{, }\StringTok{"Yes (free)"}\NormalTok{, }\StringTok{"Reference counts only"}\NormalTok{, }\StringTok{"DOI registry, reference lists, licensing"}\NormalTok{,}
  \StringTok{"PubMed"}\NormalTok{, }\StringTok{"\textasciitilde{}36M records"}\NormalTok{, }\StringTok{"Open"}\NormalTok{, }\StringTok{"Yes (free)"}\NormalTok{, }\StringTok{"No (use with others)"}\NormalTok{, }\StringTok{"MeSH indexing, biomedical depth"}\NormalTok{,}
  \StringTok{"Google Scholar"}\NormalTok{, }\StringTok{"Broadest (est. 400M+)"}\NormalTok{, }\StringTok{"Free (no bulk)"}\NormalTok{, }\StringTok{"No"}\NormalTok{, }\StringTok{"Yes (no export)"}\NormalTok{, }\StringTok{"Discovery, grey literature"}\NormalTok{,}
  \StringTok{"Semantic Scholar"}\NormalTok{, }\StringTok{"\textasciitilde{}200M records"}\NormalTok{, }\StringTok{"Free"}\NormalTok{, }\StringTok{"Yes (free)"}\NormalTok{, }\StringTok{"Yes"}\NormalTok{, }\StringTok{"AI{-}extracted entities, citation context"}\NormalTok{,}
  \StringTok{"NASA/ADS"}\NormalTok{, }\StringTok{"\textasciitilde{}16M records"}\NormalTok{, }\StringTok{"Open"}\NormalTok{, }\StringTok{"Yes (free)"}\NormalTok{, }\StringTok{"Yes"}\NormalTok{, }\StringTok{"Gold standard for astronomy/astrophysics"}\NormalTok{,}
  \StringTok{"arXiv"}\NormalTok{, }\StringTok{"\textasciitilde{}2.5M preprints"}\NormalTok{, }\StringTok{"Open"}\NormalTok{, }\StringTok{"Yes (free)"}\NormalTok{, }\StringTok{"No"}\NormalTok{, }\StringTok{"Preprints in physics, math, CS"}\NormalTok{,}
  \StringTok{"bioRxiv/medRxiv"}\NormalTok{, }\StringTok{"\textasciitilde{}300K preprints"}\NormalTok{, }\StringTok{"Open"}\NormalTok{, }\StringTok{"Yes (free)"}\NormalTok{, }\StringTok{"No"}\NormalTok{, }\StringTok{"Biology and health sciences preprints"}
\NormalTok{) }\SpecialCharTok{|\textgreater{}}
  \FunctionTok{gt}\NormalTok{() }\SpecialCharTok{|\textgreater{}}
  \FunctionTok{tab\_header}\NormalTok{(}
    \AttributeTok{title =} \StringTok{"Twelve bibliographic data sources compared"}
\NormalTok{  ) }\SpecialCharTok{|\textgreater{}}
  \FunctionTok{tab\_style}\NormalTok{(}
    \AttributeTok{style =} \FunctionTok{cell\_text}\NormalTok{(}\AttributeTok{weight =} \StringTok{"bold"}\NormalTok{),}
    \AttributeTok{locations =} \FunctionTok{cells\_column\_labels}\NormalTok{()}
\NormalTok{  ) }\SpecialCharTok{|\textgreater{}}
  \FunctionTok{cols\_width}\NormalTok{(}
\NormalTok{    Source }\SpecialCharTok{\textasciitilde{}} \FunctionTok{px}\NormalTok{(}\DecValTok{130}\NormalTok{),}
\NormalTok{    Coverage }\SpecialCharTok{\textasciitilde{}} \FunctionTok{px}\NormalTok{(}\DecValTok{170}\NormalTok{),}
    \StringTok{\textasciigrave{}}\AttributeTok{Access model}\StringTok{\textasciigrave{}} \SpecialCharTok{\textasciitilde{}} \FunctionTok{px}\NormalTok{(}\DecValTok{100}\NormalTok{),}
\NormalTok{    API }\SpecialCharTok{\textasciitilde{}} \FunctionTok{px}\NormalTok{(}\DecValTok{100}\NormalTok{),}
    \StringTok{\textasciigrave{}}\AttributeTok{Citation data}\StringTok{\textasciigrave{}} \SpecialCharTok{\textasciitilde{}} \FunctionTok{px}\NormalTok{(}\DecValTok{120}\NormalTok{),}
\NormalTok{    Strengths }\SpecialCharTok{\textasciitilde{}} \FunctionTok{px}\NormalTok{(}\DecValTok{250}\NormalTok{)}
\NormalTok{  )}
\end{Highlighting}
\end{Shaded}

\begingroup
\fontsize{8.0pt}{9.0pt}\selectfont
\begin{longtable}{>{\raggedright\arraybackslash}p{\dimexpr 0.14\linewidth -2\tabcolsep-1.5\arrayrulewidth}>{\raggedright\arraybackslash}p{\dimexpr 0.23\linewidth -2\tabcolsep-1.5\arrayrulewidth}>{\raggedright\arraybackslash}p{\dimexpr 0.13\linewidth -2\tabcolsep-1.5\arrayrulewidth}>{\raggedright\arraybackslash}p{\dimexpr 0.11\linewidth -2\tabcolsep-1.5\arrayrulewidth}>{\raggedright\arraybackslash}p{\dimexpr 0.16\linewidth -2\tabcolsep-1.5\arrayrulewidth}>{\raggedright\arraybackslash}p{\dimexpr 0.23\linewidth -2\tabcolsep-1.5\arrayrulewidth}}
\caption*{
{\fontsize{11}{14}\selectfont  Twelve bibliographic data sources compared\fontsize{8}{9}\selectfont }
} \\ 
\toprule
{\bfseries Source} & {\bfseries \nolinkurl{Coverage}} & {\bfseries Access model} & {\bfseries API} & {\bfseries Citation data} & {\bfseries \nolinkurl{Strengths}} \\ 
\midrule\addlinespace[2.5pt]
Web of Science & \textasciitilde{}90M records, 21K+ journals & \nolinkurl{Subscription} & Yes (paid) & Yes (complete) & Longest history, curated selection, JCR \\ 
Scopus & \textasciitilde{}94M records, 27K+ journals & \nolinkurl{Subscription} & Yes (paid) & Yes (complete) & Broad journal coverage, CiteScore/SJR/SNIP \\ 
\nolinkurl{Dimensions} & \textasciitilde{}140M records & \nolinkurl{Freemium} & Yes (free/paid) & Yes & Links pubs to grants, patents, clinical trials \\ 
\nolinkurl{OpenAlex} & \textasciitilde{}250M records & Fully open & Yes (free) & Yes & Open, broad coverage, author disambiguation \\ 
The Lens & \textasciitilde{}300M scholarly + 160M patent & Free (institutional) & Yes (free) & Partial & Patent-paper linkage \\ 
\nolinkurl{Crossref} & \textasciitilde{}150M records & Open & Yes (free) & Reference counts only & DOI registry, reference lists, licensing \\ 
PubMed & \textasciitilde{}36M records & Open & Yes (free) & No (use with others) & MeSH indexing, biomedical depth \\ 
Google Scholar & Broadest (est. 400M+) & Free (no bulk) & No & Yes (no export) & Discovery, grey literature \\ 
Semantic Scholar & \textasciitilde{}200M records & Free & Yes (free) & Yes & AI-extracted entities, citation context \\ 
\nolinkurl{NASA/ADS} & \textasciitilde{}16M records & Open & Yes (free) & Yes & Gold standard for astronomy/astrophysics \\ 
arXiv & \textasciitilde{}2.5M preprints & Open & Yes (free) & No & Preprints in physics, math, CS \\ 
\nolinkurl{bioRxiv/medRxiv} & \textasciitilde{}300K preprints & Open & Yes (free) & No & Biology and health sciences preprints \\ 
\bottomrule
\end{longtable}
\endgroup

\section{Diagnostics and interpretation}\label{datasources-diagnostics}

When comparing metadata across sources, consider the following diagnostic checks:

\begin{itemize}
\tightlist
\item
  \textbf{DOI match rate.} Not all DOIs may resolve in every database. Check how many of your queried DOIs returned results from each source. If a DOI is missing from one source, this reveals a coverage gap.
\item
  \textbf{Citation count discrepancies.} Citation counts will almost always differ between sources because each database counts citations over a different corpus. OpenAlex typically reports higher counts than Crossref because OpenAlex tracks citations across its full 250M-record corpus, whereas Crossref only counts ``is-referenced-by'' links deposited by publishers.
\item
  \textbf{Abstract availability.} Crossref abstracts depend on publisher deposits and are often absent, especially for older publications. OpenAlex aggregates abstracts from multiple sources and generally has better coverage.
\item
  \textbf{Publication year.} Minor discrepancies (e.g., online-first vs.~print dates) are common. Check whether your analysis is sensitive to the exact year assigned.
\item
  \textbf{Type classification.} Crossref and OpenAlex use different type taxonomies. A record classified as ``journal-article'' in Crossref might appear as ``article'' in OpenAlex. Harmonise types before analysis.
\end{itemize}

The key takeaway is that no single source is complete or definitive. Cross-validation --- fetching the same records from multiple sources --- is a hallmark of rigorous bibliometric practice.

\section{Limitations and responsible use}\label{datasources-limitations}

\section{Limitations and responsible use}\label{limitations-and-responsible-use-3}

\begin{itemize}
\tightlist
\item
  \textbf{No database is complete.} Every source has coverage gaps --- by language, discipline, publication type, or geography. Analyses based on a single source will systematically miss some literature. Always report which source was used and its known biases \citep{hicks2015}.
\item
  \textbf{Coverage does not equal quality.} A larger database is not necessarily better. Google Scholar has the broadest coverage but also the highest noise level. Curated databases like WoS and PubMed trade coverage for precision.
\item
  \textbf{Proprietary sources limit reproducibility.} Analyses based on Web of Science or Scopus cannot be independently reproduced by researchers without institutional access. For fully reproducible work, prefer open sources like OpenAlex, Crossref, or PubMed \citep{priem2022}.
\item
  \textbf{Metadata quality is variable.} Even within a single source, metadata quality varies by publisher, discipline, and publication date. Older records often have sparser metadata. Always inspect and clean your data before analysis.
\item
  \textbf{Citation counts are not comparable across sources.} Because each database counts citations over a different corpus, citation counts from OpenAlex, Scopus, and WoS are not interchangeable. Never mix citation counts from different sources in a single analysis without explicit justification and documentation \citep{visser2021}.
\end{itemize}

\section{Common pitfalls}\label{datasources-pitfalls}

\section{Common pitfalls}\label{common-pitfalls-3}

\begin{itemize}
\tightlist
\item
  \textbf{Assuming one source is enough.} Researchers often default to a single database (typically WoS or Scopus) without evaluating whether it covers their field adequately. Social sciences, humanities, and non-English literatures are systematically underrepresented in WoS \citep{singh2021}.
\item
  \textbf{Comparing citation counts across databases.} Stating that a paper has ``500 citations'' is meaningless without specifying the source. Different databases will report different numbers for the same paper.
\item
  \textbf{Ignoring access restrictions in reproducibility claims.} If your analysis uses Scopus data, other researchers need Scopus access to reproduce it. State this explicitly and consider providing OpenAlex-based alternatives.
\item
  \textbf{Treating Google Scholar as a systematic source.} Google Scholar has no API, no structured export, and no quality control on what it indexes. It is useful for discovery but not for systematic bibliometric analysis.
\item
  \textbf{Conflating preprints and peer-reviewed articles.} Sources like arXiv and bioRxiv index preprints that may later be published in journals. If your corpus includes both preprints and published versions without deduplication, you will double-count some works.
\end{itemize}

\section{Exercises}\label{datasources-exercises}

\begin{enumerate}
\def\labelenumi{\arabic{enumi}.}
\item
  \textbf{Coverage gap analysis.} Select 20 DOIs from a field of your choice. Query them in OpenAlex and Crossref. How many DOIs are found in each? Which source has better abstract coverage? Summarise your findings in a table.
\item
  \textbf{Disciplinary bias.} Use OpenAlex to compare the number of indexed works in two contrasting fields (e.g., molecular biology vs.~philosophy). Normalise by the estimated number of active researchers in each field. What does this tell you about coverage bias?
\item
  \textbf{Citation count concordance.} For a set of at least 50 works, retrieve citation counts from both OpenAlex and Crossref. Compute the Pearson and Spearman correlations between the two sets of counts. Plot the results. Are the two sources more concordant for highly cited or lowly cited papers?
\item
  \textbf{Preprint tracking.} Select 10 recent bioRxiv DOIs and check whether they appear in OpenAlex. For those that do, check whether OpenAlex links them to a published journal version. What proportion of your preprints have been published?
\item
  \textbf{Source selection exercise.} A colleague asks you to conduct a bibliometric analysis of nursing research in Sub-Saharan Africa. Which data source(s) would you recommend and why? Write a brief justification considering coverage, cost, and reproducibility.
\end{enumerate}

\section{Solutions}\label{04-solutions}

Solutions are provided in \ref{solutions}.

\section{Further reading}\label{datasources-reading}

\begin{itemize}
\tightlist
\item
  \citet{visser2021} --- Large-scale comparison of Scopus, Web of Science, Dimensions, Crossref, and Microsoft Academic. Essential reading for understanding how source choice affects results.
\item
  \citet{singh2021} --- Comparative analysis of journal coverage across WoS, Scopus, and Dimensions. Quantifies disciplinary differences in database coverage.
\item
  \citet{priem2022} --- The OpenAlex paper. Describes the data model, coverage, and design philosophy of the open alternative to proprietary databases.
\item
  \citet{garfield1955} --- The original paper proposing citation indexing. Historical context for understanding why these databases exist.
\item
  \citet{hicks2015} --- The Leiden Manifesto. Principle 4 (transparency) and Principle 6 (field variation) are directly relevant to data source selection.
\item
  \citet{waltman2016} --- Review of citation impact indicators. Discusses how indicator properties interact with database coverage.
\end{itemize}

\section{Session info}\label{datasources-sessioninfo}

\begin{Shaded}
\begin{Highlighting}[]
\FunctionTok{sessionInfo}\NormalTok{()}
\end{Highlighting}
\end{Shaded}

\begin{verbatim}
#> R version 4.6.1 (2026-06-24)
#> Platform: x86_64-pc-linux-gnu
#> Running under: Ubuntu 24.04.5 LTS
#> 
#> Matrix products: default
#> BLAS:   /usr/lib/x86_64-linux-gnu/openblas-pthread/libblas.so.3 
#> LAPACK: /usr/lib/x86_64-linux-gnu/openblas-pthread/libopenblasp-r0.3.26.so;  LAPACK version 3.12.0
#> 
#> locale:
#>  [1] LC_CTYPE=C.UTF-8       LC_NUMERIC=C           LC_TIME=C.UTF-8       
#>  [4] LC_COLLATE=C.UTF-8     LC_MONETARY=C.UTF-8    LC_MESSAGES=C.UTF-8   
#>  [7] LC_PAPER=C.UTF-8       LC_NAME=C              LC_ADDRESS=C          
#> [10] LC_TELEPHONE=C         LC_MEASUREMENT=C.UTF-8 LC_IDENTIFICATION=C   
#> 
#> time zone: UTC
#> tzcode source: system (glibc)
#> 
#> attached base packages:
#> [1] stats     graphics  grDevices utils     datasets  methods   base     
#> 
#> other attached packages:
#>  [1] rcrossref_1.2.1 gt_1.4.0        tidytext_0.4.3  glue_1.8.1     
#>  [5] openalexR_3.1.2 lubridate_1.9.5 forcats_1.0.1   stringr_1.6.0  
#>  [9] dplyr_1.2.1     purrr_1.2.2     readr_2.2.0     tidyr_1.3.2    
#> [13] tibble_3.3.1    ggplot2_4.0.3   tidyverse_2.0.0
#> 
#> loaded via a namespace (and not attached):
#>  [1] gtable_0.3.6       xfun_0.62          htmlwidgets_1.6.4  lattice_0.22-9    
#>  [5] tzdb_0.5.0         vctrs_0.7.3        tools_4.6.1        generics_0.1.4    
#>  [9] curl_8.1.0         janeaustenr_1.0.0  pkgconfig_2.0.3    tokenizers_0.3.0  
#> [13] Matrix_1.7-5       RColorBrewer_1.1-3 S7_0.2.2           lifecycle_1.0.5   
#> [17] compiler_4.6.1     farver_2.1.2       httpuv_1.6.17      htmltools_0.5.9   
#> [21] SnowballC_0.7.1    yaml_2.3.12        pillar_1.11.1      later_1.4.8       
#> [25] DT_0.34.0          viridis_0.6.5      mime_0.13          stopwords_2.3     
#> [29] tidyselect_1.2.1   digest_0.6.39      stringi_1.8.9      bookdown_0.48     
#> [33] labeling_0.4.3     rprojroot_2.1.1    fastmap_1.2.0      grid_4.6.1        
#> [37] here_1.0.2         cli_3.6.6          magrittr_2.0.5     triebeard_0.4.1   
#> [41] crul_1.6.0         utf8_1.2.6         withr_3.0.3        scales_1.4.0      
#> [45] promises_1.5.0     timechange_0.4.0   rmarkdown_2.32     otel_0.2.0        
#> [49] gridExtra_2.3.1    hms_1.1.4          shiny_1.14.0       evaluate_1.0.5    
#> [53] knitr_1.52         miniUI_0.1.2       viridisLite_0.4.3  urltools_1.7.3.1  
#> [57] rlang_1.3.0        Rcpp_1.1.2         httpcode_0.3.0     xtable_1.8-8      
#> [61] xml2_1.6.0         jsonlite_2.0.0     rstudioapi_0.19.0  plyr_1.8.9        
#> [65] R6_2.6.1           systemfonts_1.3.2  fs_2.1.0
\end{verbatim}

\part{Data Acquisition in R}\label{part-data-acquisition-in-r}

\chapter{APIs and Packages}\label{apis-and-packages}

\emph{R. Heller · Date: 2026-05-09 · Last changed: 2026-10-31}

\section{Learning objectives}\label{apis-objectives}

After completing this chapter, you will be able to:

\begin{itemize}
\tightlist
\item
  Query the OpenAlex API using \texttt{openalexR} to retrieve works, authors, institutions, and concepts
\item
  Apply filters, pagination, and sampling to control the size and scope of API results
\item
  Deduplicate records by DOI and clean author and affiliation metadata
\item
  Cache API responses locally so repeated book builds do not re-query the API
\item
  Use \texttt{rcrossref} and \texttt{pubmedR} as complementary data sources
\end{itemize}

\section{Setup}\label{apis-setup}

\begin{Shaded}
\begin{Highlighting}[]
\FunctionTok{library}\NormalTok{(tidyverse)}
\FunctionTok{library}\NormalTok{(openalexR)}
\FunctionTok{library}\NormalTok{(glue)}

\FunctionTok{set.seed}\NormalTok{(}\DecValTok{20260509}\NormalTok{)}

\FunctionTok{source}\NormalTok{(here}\SpecialCharTok{::}\FunctionTok{here}\NormalTok{(}\StringTok{"R"}\NormalTok{, }\StringTok{"api\_helpers.R"}\NormalTok{))}
\FunctionTok{source}\NormalTok{(here}\SpecialCharTok{::}\FunctionTok{here}\NormalTok{(}\StringTok{"R"}\NormalTok{, }\StringTok{"utils.R"}\NormalTok{))}
\FunctionTok{source}\NormalTok{(here}\SpecialCharTok{::}\FunctionTok{here}\NormalTok{(}\StringTok{"R"}\NormalTok{, }\StringTok{"sci\_palette.R"}\NormalTok{))}
\end{Highlighting}
\end{Shaded}

\section{Conceptual background}\label{apis-background}

Bibliometric research begins with data. The quality and completeness of a bibliographic corpus determines the validity of every downstream analysis --- from citation counts to network maps to topic models.

Historically, Web of Science and Scopus have dominated as data sources for scientometric research. Both require institutional subscriptions and impose restrictions on bulk data use. The launch of \textbf{OpenAlex} in 2022 changed this landscape fundamentally \citep{priem2022}. OpenAlex is a fully open index of over 250 million scholarly works, built from Crossref, PubMed, institutional repositories, and other open sources. It provides a free REST API with no authentication required, making it the ideal backbone for reproducible bibliometric research.

OpenAlex organises its data around five entity types: \textbf{works} (publications), \textbf{authors}, \textbf{sources} (journals, repositories), \textbf{institutions}, and \textbf{topics} (hierarchical subject tags). Each entity has a unique OpenAlex ID and can be queried via filters on metadata fields such as publication date, citation count, open-access status, and institutional affiliation.

The \texttt{openalexR} package provides a tidy R interface to the OpenAlex API. It handles pagination, rate limiting, and result parsing, returning clean tibbles ready for analysis. For metadata not covered by OpenAlex --- such as reference lists and funding information --- \texttt{rcrossref} provides access to Crossref's extensive metadata registry, and \texttt{pubmedR} covers the biomedical literature indexed in PubMed/MEDLINE.

\section{Worked example}\label{apis-example}

We build a small corpus of recent scientometrics research, demonstrating the full acquisition-to-clean-data pipeline.

\subsection{Querying works from OpenAlex}\label{querying-works-from-openalex}

The provider function is \texttt{openalexR::oa\_fetch()}. Runnable examples use the book's \texttt{fetch\_openalex()} wrapper, sourced above, to cache responses by their exact arguments. Set \texttt{SCIENTOMETRICS\_OFFLINE=1} to reuse existing OpenAlex caches without requesting new data; a missing cache stops with an explicit error. Other providers have their own network requirements. We search for works published in the journal \emph{Scientometrics} between 2020 and 2023, sampling 200 records for a manageable demonstration.

\begin{Shaded}
\begin{Highlighting}[]
\NormalTok{works }\OtherTok{\textless{}{-}} \FunctionTok{fetch\_openalex}\NormalTok{(}
  \AttributeTok{entity =} \StringTok{"works"}\NormalTok{,}
  \AttributeTok{primary\_location.source.id =} \StringTok{"S148561398"}\NormalTok{,}
  \AttributeTok{from\_publication\_date =} \StringTok{"2020{-}01{-}01"}\NormalTok{,}
  \AttributeTok{to\_publication\_date =} \StringTok{"2023{-}12{-}31"}\NormalTok{,}
  \AttributeTok{options =} \FunctionTok{list}\NormalTok{(}\AttributeTok{sample =} \DecValTok{200}\NormalTok{, }\AttributeTok{seed =} \DecValTok{42}\NormalTok{)}
\NormalTok{)}

\FunctionTok{glimpse}\NormalTok{(works)}
\end{Highlighting}
\end{Shaded}

\begin{verbatim}
#> Rows: 200
#> Columns: 45
#> $ id                            <chr> "https://openalex.org/W3157843452", "htt~
#> $ title                         <chr> "Scientific knowledge production and eco~
#> $ display_name                  <chr> "Scientific knowledge production and eco~
#> $ authorships                   <list> [<tbl_df[2 x 8]>], [<tbl_df[1 x 8]>], [~
#> $ abstract                      <chr> NA, "Abstract This paper investigates th~
#> $ doi                           <chr> "https://doi.org/10.1007/s11192-021-0397~
#> $ publication_date              <date> 2021-04-26, 2022-11-14, 2021-08-18, 202~
#> $ publication_year              <int> 2021, 2022, 2021, 2020, 2020, 2020, 2023~
#> $ relevance_score               <dbl> 0.999, 0.998, 0.997, 0.997, 0.997, 0.997~
#> $ fwci                          <dbl> 1.445, 1.575, 3.218, 2.548, 12.740, 0.42~
#> $ cited_by_count                <int> 17, 5, 12, 10, 48, 1, 15, 116, 7, 21, 64~
#> $ counts_by_year                <list> [<data.frame[6 x 2]>], [<data.frame[3 x~
#> $ ids                           <list> <"https://openalex.org/W3157843452", "h~
#> $ type                          <chr> "article", "article", "article", "articl~
#> $ is_oa                         <lgl> FALSE, TRUE, FALSE, TRUE, FALSE, FALSE, ~
#> $ is_oa_anywhere                <lgl> TRUE, TRUE, FALSE, TRUE, FALSE, FALSE, T~
#> $ oa_status                     <chr> "green", "hybrid", "closed", "hybrid", "~
#> $ oa_url                        <chr> "http://hdl.handle.net/11336/139054", "h~
#> $ any_repository_has_fulltext   <lgl> TRUE, TRUE, FALSE, FALSE, FALSE, FALSE, ~
#> $ source_display_name           <chr> "Scientometrics", "Scientometrics", "Sci~
#> $ source_id                     <chr> "https://openalex.org/S148561398", "http~
#> $ issn_l                        <chr> "0138-9130", "0138-9130", "0138-9130", "~
#> $ host_organization             <chr> "https://openalex.org/P4310320108", "htt~
#> $ host_organization_name        <chr> "Springer Nature (Netherlands)", "Spring~
#> $ landing_page_url              <chr> "https://doi.org/10.1007/s11192-021-0397~
#> $ pdf_url                       <chr> NA, "https://link.springer.com/content/p~
#> $ license                       <chr> NA, "cc-by", NA, "cc-by", NA, NA, NA, "c~
#> $ version                       <chr> "publishedVersion", "publishedVersion", ~
#> $ referenced_works              <list> <"https://openalex.org/W175573758", "ht~
#> $ referenced_works_count        <int> 61, 79, 47, 18, 31, 21, 47, 56, 62, 21, ~
#> $ related_works                 <list> <"https://openalex.org/W4415584083", "h~
#> $ concepts                      <list> [<data.frame[25 x 5]>], [<data.frame[21~
#> $ topics                        <list> [<tbl_df[12 x 5]>], [<tbl_df[12 x 5]>],~
#> $ keywords                      <list> [<data.frame[10 x 3]>], [<data.frame[10~
#> $ is_paratext                   <lgl> FALSE, FALSE, FALSE, FALSE, FALSE, FALSE~
#> $ is_retracted                  <lgl> FALSE, FALSE, FALSE, FALSE, FALSE, FALSE~
#> $ language                      <chr> "en", "en", "en", "en", "en", "en", "en"~
#> $ sustainable_development_goals <list> NA, NA, NA, [<data.frame[1 x 3]>], NA, ~
#> $ awards                        <list> NA, <"https://openalex.org/G434337789",~
#> $ funders                       <list> NA, [<data.frame[2 x 3]>], [<data.frame~
#> $ apc                           <list> [<data.frame[2 x 5]>], [<data.frame[2 x~
#> $ first_page                    <chr> "4565", "651", "8653", "1783", "2413", "~
#> $ last_page                     <chr> "4587", "698", "8672", "1800", "2435", "~
#> $ volume                        <chr> "126", "128", "126", "125", "126", "126"~
#> $ issue                         <chr> "6", "1", "10", "2", "3", "3", "3", "12"~
\end{verbatim}

The result is a tibble where each row is a work and columns include \texttt{id}, \texttt{display\_name}, \texttt{publication\_date}, \texttt{cited\_by\_count}, \texttt{doi}, and nested list-columns for authors and concepts.

\subsection{Inspecting the results}\label{inspecting-the-results}

\begin{Shaded}
\begin{Highlighting}[]
\NormalTok{works }\SpecialCharTok{|\textgreater{}}
  \FunctionTok{select}\NormalTok{(display\_name, publication\_date, cited\_by\_count, doi) }\SpecialCharTok{|\textgreater{}}
  \FunctionTok{arrange}\NormalTok{(}\FunctionTok{desc}\NormalTok{(cited\_by\_count)) }\SpecialCharTok{|\textgreater{}}
  \FunctionTok{head}\NormalTok{(}\DecValTok{10}\NormalTok{)}
\end{Highlighting}
\end{Shaded}

\begin{verbatim}
#> # A tibble: 10 x 4
#>    display_name                            publication_date cited_by_count doi  
#>    <chr>                                   <date>                    <int> <chr>
#>  1 Industry 5.0: improving humanization a~ 2022-04-27                  495 http~
#>  2 'Are principals instructional leaders ~ 2020-01-30                  343 http~
#>  3 Smart bibliometrics: an integrated met~ 2022-05-21                  258 http~
#>  4 Rethinking SME default prediction: a s~ 2021-01-29                  186 http~
#>  5 Why do papers from international colla~ 2022-08-22                  116 http~
#>  6 Comprehensive metrological and content~ 2020-07-13                  116 http~
#>  7 Understanding the temporal evolution o~ 2020-11-19                   90 http~
#>  8 Prevalence of potentially predatory pu~ 2021-03-26                   88 http~
#>  9 Scientometric analysis and knowledge m~ 2021-01-03                   84 http~
#> 10 The impact of research output on econo~ 2020-03-28                   73 http~
\end{verbatim}

\subsection{Querying authors}\label{querying-authors}

We can also retrieve author-level data. Let us look up authors associated with scientometrics research by filtering on the relevant topic.

\begin{Shaded}
\begin{Highlighting}[]
\NormalTok{authors }\OtherTok{\textless{}{-}} \FunctionTok{fetch\_openalex}\NormalTok{(}
  \AttributeTok{entity =} \StringTok{"authors"}\NormalTok{,}
  \AttributeTok{topics.id =} \StringTok{"T10102"}\NormalTok{,}
  \AttributeTok{works\_count =} \StringTok{"\textgreater{}50"}\NormalTok{,}
  \AttributeTok{options =} \FunctionTok{list}\NormalTok{(}\AttributeTok{sample =} \DecValTok{50}\NormalTok{, }\AttributeTok{seed =} \DecValTok{42}\NormalTok{)}
\NormalTok{)}

\NormalTok{authors }\SpecialCharTok{|\textgreater{}}
  \FunctionTok{select}\NormalTok{(display\_name, works\_count, cited\_by\_count) }\SpecialCharTok{|\textgreater{}}
  \FunctionTok{arrange}\NormalTok{(}\FunctionTok{desc}\NormalTok{(cited\_by\_count)) }\SpecialCharTok{|\textgreater{}}
  \FunctionTok{head}\NormalTok{(}\DecValTok{10}\NormalTok{)}
\end{Highlighting}
\end{Shaded}

\begin{verbatim}
#> # A tibble: 10 x 3
#>    display_name       works_count cited_by_count
#>    <chr>                    <int>          <int>
#>  1 Hans‐Dieter Daniel         220           9117
#>  2 Rüdiger Mutz               121           6318
#>  3 Haak Laurel                296           3356
#>  4 Babak Pourbohloul           53           3169
#>  5 Jordi Suriñach             290           3051
#>  6 Eun‐Kyung Chung            186           2888
#>  7 Eti Herman                 103           2647
#>  8 Manuel Goyanes             162           2349
#>  9 Javier Sanz‐Valero         277           2219
#> 10 Gali Halevi                 74           1859
\end{verbatim}

\subsection{Deduplication by DOI}\label{deduplication-by-doi}

Real-world corpora often contain duplicates when merging results from multiple queries or sources. The \texttt{dedupe\_by\_doi()} function from our companion package handles this cleanly.

\begin{Shaded}
\begin{Highlighting}[]
\NormalTok{works\_with\_dupes }\OtherTok{\textless{}{-}} \FunctionTok{bind\_rows}\NormalTok{(works, works[}\DecValTok{1}\SpecialCharTok{:}\DecValTok{10}\NormalTok{, ])}
\FunctionTok{cat}\NormalTok{(}\FunctionTok{glue}\NormalTok{(}\StringTok{"Before dedup: \{nrow(works\_with\_dupes)\} rows}\SpecialCharTok{\textbackslash{}n\textbackslash{}n}\StringTok{"}\NormalTok{))}
\end{Highlighting}
\end{Shaded}

\begin{verbatim}
#> Before dedup: 210 rows
\end{verbatim}

\begin{Shaded}
\begin{Highlighting}[]
\NormalTok{works\_clean }\OtherTok{\textless{}{-}} \FunctionTok{dedupe\_by\_doi}\NormalTok{(works\_with\_dupes)}
\FunctionTok{cat}\NormalTok{(}\FunctionTok{glue}\NormalTok{(}\StringTok{"After dedup:  \{nrow(works\_clean)\} rows}\SpecialCharTok{\textbackslash{}n}\StringTok{"}\NormalTok{))}
\end{Highlighting}
\end{Shaded}

\begin{verbatim}
#> After dedup:  200 rows
\end{verbatim}

\subsection{Caching for reproducible builds}\label{caching-for-reproducible-builds}

The \texttt{fetch\_openalex()} wrapper in \texttt{R/api\_helpers.R} automatically caches API responses to disk. On subsequent calls with the same arguments, it reads from cache instead of hitting the API.

\begin{Shaded}
\begin{Highlighting}[]
\NormalTok{works\_cached }\OtherTok{\textless{}{-}} \FunctionTok{fetch\_openalex}\NormalTok{(}
  \AttributeTok{entity =} \StringTok{"works"}\NormalTok{,}
  \AttributeTok{primary\_location.source.id =} \StringTok{"S148561398"}\NormalTok{,}
  \AttributeTok{from\_publication\_date =} \StringTok{"2023{-}01{-}01"}\NormalTok{,}
  \AttributeTok{to\_publication\_date =} \StringTok{"2023{-}06{-}30"}\NormalTok{,}
  \AttributeTok{options =} \FunctionTok{list}\NormalTok{(}\AttributeTok{sample =} \DecValTok{50}\NormalTok{, }\AttributeTok{seed =} \DecValTok{42}\NormalTok{),}
  \AttributeTok{cache\_dir =} \StringTok{"\_freeze/openalex\_cache"}\NormalTok{,}
  \AttributeTok{cache\_days =} \DecValTok{30}
\NormalTok{)}

\FunctionTok{nrow}\NormalTok{(works\_cached)}
\end{Highlighting}
\end{Shaded}

\begin{verbatim}
#> [1] 50
\end{verbatim}

\subsection{Complementary sources: Crossref}\label{complementary-sources-crossref}

For detailed metadata (reference lists, funders, licenses), Crossref is invaluable.

\begin{Shaded}
\begin{Highlighting}[]
\CommentTok{\# Requires network access — cached output used in book builds}
\NormalTok{cr\_work }\OtherTok{\textless{}{-}}\NormalTok{ rcrossref}\SpecialCharTok{::}\FunctionTok{cr\_works}\NormalTok{(}\AttributeTok{doi =} \StringTok{"10.1007/s11192{-}017{-}2300{-}7"}\NormalTok{)}
\NormalTok{cr\_work}\SpecialCharTok{$}\NormalTok{data }\SpecialCharTok{|\textgreater{}}
  \FunctionTok{select}\NormalTok{(title, container.title, deposited, is.referenced.by.count)}
\end{Highlighting}
\end{Shaded}

\subsection{Visualization}\label{visualization-1}

\begin{Shaded}
\begin{Highlighting}[]
\NormalTok{works }\SpecialCharTok{|\textgreater{}}
  \FunctionTok{mutate}\NormalTok{(}\AttributeTok{year =} \FunctionTok{year}\NormalTok{(publication\_date)) }\SpecialCharTok{|\textgreater{}}
  \FunctionTok{count}\NormalTok{(year) }\SpecialCharTok{|\textgreater{}}
  \FunctionTok{ggplot}\NormalTok{(}\FunctionTok{aes}\NormalTok{(}\AttributeTok{x =}\NormalTok{ year, }\AttributeTok{y =}\NormalTok{ n)) }\SpecialCharTok{+}
  \FunctionTok{geom\_col}\NormalTok{(}\AttributeTok{fill =} \FunctionTok{palette\_sci}\NormalTok{(}\DecValTok{1}\NormalTok{)) }\SpecialCharTok{+}
  \FunctionTok{labs}\NormalTok{(}\AttributeTok{x =} \StringTok{"Publication year"}\NormalTok{, }\AttributeTok{y =} \StringTok{"Number of works"}\NormalTok{) }\SpecialCharTok{+}
  \FunctionTok{theme\_sci}\NormalTok{()}
\end{Highlighting}
\end{Shaded}

\begin{figure}

{\centering \includegraphics[width=0.9\linewidth,alt={Bar chart showing the number of sampled publications per year from the journal Scientometrics, covering 2020 to 2023.}]{scientometrics-in-r_files/figure-latex/05-fig-apis-publications-per-year-1} 

}

\caption{Annual publication counts for sampled Scientometrics articles (2020–2023).}\label{fig:05-fig-apis-publications-per-year}
\end{figure}

\section{Diagnostics and interpretation}\label{apis-diagnostics}

When working with API data, always verify:

\begin{itemize}
\tightlist
\item
  \textbf{Completeness}: Compare your result count against the API's reported \texttt{meta\$count} to confirm you received all expected records. With sampling, the count reflects the sample, not the full set.
\item
  \textbf{Coverage dates}: Check \texttt{min(works\$publication\_date)} and \texttt{max(works\$publication\_date)} to confirm your date filters were applied correctly.
\item
  \textbf{DOI presence}: Not all works have DOIs. Check \texttt{mean(!is.na(works\$doi))} to understand the DOI coverage rate in your corpus.
\item
  \textbf{Nested columns}: Author and concept columns are list-columns. Use \texttt{tidyr::unnest()} carefully and inspect for \texttt{NULL} entries before unnesting.
\end{itemize}

\begin{Shaded}
\begin{Highlighting}[]
\FunctionTok{cat}\NormalTok{(}\FunctionTok{glue}\NormalTok{(}\StringTok{"Total works: \{nrow(works)\}}\SpecialCharTok{\textbackslash{}n}\StringTok{"}\NormalTok{))}
\end{Highlighting}
\end{Shaded}

\begin{verbatim}
#> Total works: 200
\end{verbatim}

\begin{Shaded}
\begin{Highlighting}[]
\FunctionTok{cat}\NormalTok{(}\FunctionTok{glue}\NormalTok{(}\StringTok{"DOI coverage: \{scales::percent(mean(!is.na(works$doi)))\}}\SpecialCharTok{\textbackslash{}n}\StringTok{"}\NormalTok{))}
\end{Highlighting}
\end{Shaded}

\begin{verbatim}
#> DOI coverage: 100%
\end{verbatim}

\begin{Shaded}
\begin{Highlighting}[]
\FunctionTok{cat}\NormalTok{(}\FunctionTok{glue}\NormalTok{(}\StringTok{"Date range: \{min(works$publication\_date)\} to \{max(works$publication\_date)\}}\SpecialCharTok{\textbackslash{}n}\StringTok{"}\NormalTok{))}
\end{Highlighting}
\end{Shaded}

\begin{verbatim}
#> Date range: 2020-01-01 to 2023-12-27
\end{verbatim}

\section{Limitations and responsible use}\label{apis-limitations}

\section{Limitations and responsible use}\label{limitations-and-responsible-use-4}

\begin{itemize}
\tightlist
\item
  \textbf{Coverage bias}: OpenAlex is strongest in English-language, peer-reviewed journal articles. Conference proceedings, books, and non-English literature may be underrepresented. Always report the data source and its known coverage limitations.
\item
  \textbf{Metadata quality}: Author names and affiliations are parsed algorithmically and may contain errors, especially for names with diacritics or non-Latin scripts.
\item
  \textbf{Temporal lag}: OpenAlex updates continuously but may lag behind primary databases by days or weeks. Very recent publications may be missing.
\item
  \textbf{API rate limits}: OpenAlex allows 10 requests per second for unauthenticated users (100/s with a polite pool email). Respect these limits to avoid being throttled.
\item
  \textbf{Reproducibility}: API results can change as OpenAlex updates its data. Always cache results and document the date of data collection for reproducibility \citep{hicks2015}.
\end{itemize}

\section{Common pitfalls}\label{apis-pitfalls}

\section{Common pitfalls}\label{common-pitfalls-4}

\begin{itemize}
\tightlist
\item
  \textbf{Not caching}: Running \texttt{oa\_fetch()} directly in code chunks without caching means every \texttt{quarto\ render} hits the API, producing different results each time and risking rate limits.
\item
  \textbf{Ignoring pagination}: \texttt{oa\_fetch()} handles pagination automatically, but very large queries (\textgreater10,000 works) can time out. Use filters to narrow your query or use cursor-based pagination.
\item
  \textbf{Confusing sampling with filtering}: \texttt{options\ =\ list(sample\ =\ 200)} returns a random sample from the full result set, not the first 200 records. This is useful for demonstrations but not for exhaustive analyses.
\item
  \textbf{Treating OpenAlex IDs as permanent}: While generally stable, OpenAlex IDs can change when entities are merged. Use DOIs as the primary identifier for works.
\item
  \textbf{Merging sources without deduplication}: When combining results from OpenAlex and Crossref, records may overlap. Always deduplicate by DOI before analysis.
\end{itemize}

\section{Exercises}\label{apis-exercises}

\begin{enumerate}
\def\labelenumi{\arabic{enumi}.}
\item
  \textbf{Fetch works by institution.} Use \texttt{oa\_fetch()} to retrieve works affiliated with your own institution. How many works does OpenAlex index for it? What is the DOI coverage rate? \emph{(Hint: find your institution's OpenAlex ID via \texttt{oa\_fetch(entity\ =\ "institutions",\ search\ =\ "...")})}
\item
  \textbf{Compare Crossref metadata.} Pick a DOI from your fetched works and retrieve its metadata from Crossref using \texttt{rcrossref::cr\_works()}. Compare the citation count reported by OpenAlex vs.~Crossref. Why might they differ?
\item
  \textbf{Build a multi-year corpus.} Fetch all works from a journal of your choice for 2015--2023 (without sampling). Plot the annual publication count and the median citation count per year. What trends do you observe?
\end{enumerate}

\section{Solutions}\label{05-solutions}

Solutions are provided in \ref{solutions}.

\section{Further reading}\label{apis-reading}

\begin{itemize}
\tightlist
\item
  \citet{priem2022} --- The OpenAlex paper; describes the data model, coverage, and API design.
\item
  \citet{aria2017} --- The \texttt{bibliometrix} package; an alternative high-level interface for bibliometric data acquisition and analysis.
\item
  \citet{garfield1955} --- The foundational paper on citation indexing, motivating why these databases exist.
\end{itemize}

\section{Session info}\label{apis-sessioninfo}

\begin{Shaded}
\begin{Highlighting}[]
\FunctionTok{sessionInfo}\NormalTok{()}
\end{Highlighting}
\end{Shaded}

\begin{verbatim}
#> R version 4.6.1 (2026-06-24)
#> Platform: x86_64-pc-linux-gnu
#> Running under: Ubuntu 24.04.5 LTS
#> 
#> Matrix products: default
#> BLAS:   /usr/lib/x86_64-linux-gnu/openblas-pthread/libblas.so.3 
#> LAPACK: /usr/lib/x86_64-linux-gnu/openblas-pthread/libopenblasp-r0.3.26.so;  LAPACK version 3.12.0
#> 
#> locale:
#>  [1] LC_CTYPE=C.UTF-8       LC_NUMERIC=C           LC_TIME=C.UTF-8       
#>  [4] LC_COLLATE=C.UTF-8     LC_MONETARY=C.UTF-8    LC_MESSAGES=C.UTF-8   
#>  [7] LC_PAPER=C.UTF-8       LC_NAME=C              LC_ADDRESS=C          
#> [10] LC_TELEPHONE=C         LC_MEASUREMENT=C.UTF-8 LC_IDENTIFICATION=C   
#> 
#> time zone: UTC
#> tzcode source: system (glibc)
#> 
#> attached base packages:
#> [1] stats     graphics  grDevices utils     datasets  methods   base     
#> 
#> other attached packages:
#>  [1] rcrossref_1.2.1 gt_1.4.0        tidytext_0.4.3  glue_1.8.1     
#>  [5] openalexR_3.1.2 lubridate_1.9.5 forcats_1.0.1   stringr_1.6.0  
#>  [9] dplyr_1.2.1     purrr_1.2.2     readr_2.2.0     tidyr_1.3.2    
#> [13] tibble_3.3.1    ggplot2_4.0.3   tidyverse_2.0.0
#> 
#> loaded via a namespace (and not attached):
#>  [1] gtable_0.3.6       xfun_0.62          htmlwidgets_1.6.4  lattice_0.22-9    
#>  [5] tzdb_0.5.0         vctrs_0.7.3        tools_4.6.1        generics_0.1.4    
#>  [9] curl_8.1.0         janeaustenr_1.0.0  pkgconfig_2.0.3    tokenizers_0.3.0  
#> [13] Matrix_1.7-5       RColorBrewer_1.1-3 S7_0.2.2           lifecycle_1.0.5   
#> [17] compiler_4.6.1     farver_2.1.2       httpuv_1.6.17      htmltools_0.5.9   
#> [21] SnowballC_0.7.1    yaml_2.3.12        pillar_1.11.1      later_1.4.8       
#> [25] DT_0.34.0          viridis_0.6.5      mime_0.13          stopwords_2.3     
#> [29] tidyselect_1.2.1   digest_0.6.39      stringi_1.8.9      bookdown_0.48     
#> [33] labeling_0.4.3     rprojroot_2.1.1    fastmap_1.2.0      grid_4.6.1        
#> [37] here_1.0.2         cli_3.6.6          magrittr_2.0.5     triebeard_0.4.1   
#> [41] crul_1.6.0         utf8_1.2.6         withr_3.0.3        scales_1.4.0      
#> [45] promises_1.5.0     timechange_0.4.0   rmarkdown_2.32     otel_0.2.0        
#> [49] gridExtra_2.3.1    hms_1.1.4          shiny_1.14.0       evaluate_1.0.5    
#> [53] knitr_1.52         miniUI_0.1.2       viridisLite_0.4.3  urltools_1.7.3.1  
#> [57] rlang_1.3.0        Rcpp_1.1.2         httpcode_0.3.0     xtable_1.8-8      
#> [61] xml2_1.6.0         jsonlite_2.0.0     rstudioapi_0.19.0  plyr_1.8.9        
#> [65] R6_2.6.1           systemfonts_1.3.2  fs_2.1.0
\end{verbatim}

\chapter{Parsing Native Exports}\label{parsing-native-exports}

\emph{R. Heller · Date: 2026-05-09 · Last changed: 2026-10-31}

\section{Learning objectives}\label{parsing-objectives}

After completing this chapter, you will be able to:

\begin{itemize}
\tightlist
\item
  Parse BibTeX files using \texttt{bib2df} and \texttt{RefManageR}
\item
  Import Web of Science plain-text and Scopus CSV exports via \texttt{bibliometrix}
\item
  Standardize field names across formats into a unified schema
\item
  Diagnose common parsing failures and missing fields
\end{itemize}

\section{Setup}\label{parsing-setup}

\begin{Shaded}
\begin{Highlighting}[]
\FunctionTok{library}\NormalTok{(tidyverse)}
\FunctionTok{library}\NormalTok{(bib2df)}
\FunctionTok{library}\NormalTok{(RefManageR)}
\FunctionTok{library}\NormalTok{(bibliometrix)}
\FunctionTok{library}\NormalTok{(glue)}

\FunctionTok{set.seed}\NormalTok{(}\DecValTok{20260509}\NormalTok{)}

\FunctionTok{source}\NormalTok{(here}\SpecialCharTok{::}\FunctionTok{here}\NormalTok{(}\StringTok{"R"}\NormalTok{, }\StringTok{"sci\_palette.R"}\NormalTok{))}
\end{Highlighting}
\end{Shaded}

\section{Conceptual background}\label{parsing-background}

Not all bibliometric data comes from an API. Researchers frequently begin with files exported from database search interfaces: BibTeX (\texttt{.bib}) from Google Scholar or Zotero, RIS (\texttt{.ris}) from Scopus or PubMed, plain text from Web of Science, or CSV from Scopus. Each format encodes the same information --- authors, title, year, journal, DOI --- in different structures.

The \texttt{bibliometrix} package \citep{aria2017} provides \texttt{convert2df()}, a versatile function that reads multiple export formats and returns a standardized data frame. For BibTeX specifically, \texttt{bib2df} and \texttt{RefManageR} offer complementary parsing capabilities. The key challenge is not parsing per se but \emph{standardization}: ensuring that author names, journal titles, and identifiers are represented consistently regardless of the source format.

\section{Worked example}\label{parsing-example}

\subsection{Parsing a BibTeX string}\label{parsing-a-bibtex-string}

We construct a small BibTeX string inline to demonstrate parsing without external files.

\begin{Shaded}
\begin{Highlighting}[]
\NormalTok{bib\_text }\OtherTok{\textless{}{-}} \StringTok{\textquotesingle{}}
\StringTok{@article\{smith2023,}
\StringTok{  author = \{Smith, John A. and Doe, Jane B.\},}
\StringTok{  title = \{Advances in Bibliometric Methods\},}
\StringTok{  journal = \{Journal of Informetrics\},}
\StringTok{  year = \{2023\},}
\StringTok{  volume = \{17\},}
\StringTok{  pages = \{101{-}{-}115\},}
\StringTok{  doi = \{10.1234/example.2023\}}
\StringTok{\}}
\StringTok{@article\{chen2022,}
\StringTok{  author = \{Chen, Wei and Li, Xiao\},}
\StringTok{  title = \{Citation Networks in Computer Science\},}
\StringTok{  journal = \{Scientometrics\},}
\StringTok{  year = \{2022\},}
\StringTok{  volume = \{127\},}
\StringTok{  pages = \{3201{-}{-}3220\},}
\StringTok{  doi = \{10.1234/example.2022\}}
\StringTok{\}}
\StringTok{\textquotesingle{}}

\NormalTok{tmp\_bib }\OtherTok{\textless{}{-}} \FunctionTok{tempfile}\NormalTok{(}\AttributeTok{fileext =} \StringTok{".bib"}\NormalTok{)}
\FunctionTok{writeLines}\NormalTok{(bib\_text, tmp\_bib)}

\NormalTok{bib\_df }\OtherTok{\textless{}{-}} \FunctionTok{bib2df}\NormalTok{(tmp\_bib)}
\FunctionTok{glimpse}\NormalTok{(bib\_df)}
\end{Highlighting}
\end{Shaded}

\begin{verbatim}
#> Rows: 2
#> Columns: 27
#> $ CATEGORY     <chr> "ARTICLE", "ARTICLE"
#> $ BIBTEXKEY    <chr> "smith2023", "chen2022"
#> $ ADDRESS      <chr> NA, NA
#> $ ANNOTE       <chr> NA, NA
#> $ AUTHOR       <list> <"Smith, John A.", "Doe, Jane B.">, <"Chen, Wei", "Li, Xi~
#> $ BOOKTITLE    <chr> NA, NA
#> $ CHAPTER      <chr> NA, NA
#> $ CROSSREF     <chr> NA, NA
#> $ EDITION      <chr> NA, NA
#> $ EDITOR       <list> NA, NA
#> $ HOWPUBLISHED <chr> NA, NA
#> $ INSTITUTION  <chr> NA, NA
#> $ JOURNAL      <chr> "Journal of Informetrics", "Scientometrics"
#> $ KEY          <chr> NA, NA
#> $ MONTH        <chr> NA, NA
#> $ NOTE         <chr> NA, NA
#> $ NUMBER       <chr> NA, NA
#> $ ORGANIZATION <chr> NA, NA
#> $ PAGES        <chr> "101--115", "3201--3220"
#> $ PUBLISHER    <chr> NA, NA
#> $ SCHOOL       <chr> NA, NA
#> $ SERIES       <chr> NA, NA
#> $ TITLE        <chr> "Advances in Bibliometric Methods", "Citation Networks in~
#> $ TYPE         <chr> NA, NA
#> $ VOLUME       <chr> "17", "127"
#> $ YEAR         <dbl> 2023, 2022
#> $ DOI          <chr> "10.1234/example.2023", "10.1234/example.2022"
\end{verbatim}

\subsection{Parsing with RefManageR}\label{parsing-with-refmanager}

\begin{Shaded}
\begin{Highlighting}[]
\NormalTok{bib\_rm }\OtherTok{\textless{}{-}} \FunctionTok{ReadBib}\NormalTok{(tmp\_bib)}
\FunctionTok{as.data.frame}\NormalTok{(bib\_rm) }\SpecialCharTok{|\textgreater{}} \FunctionTok{select}\NormalTok{(title, author, year, journal, doi)}
\end{Highlighting}
\end{Shaded}

\begin{verbatim}
#>                                           title                        author
#> smith2023      Advances in Bibliometric Methods John A. Smith and Jane B. Doe
#> chen2022  Citation Networks in Computer Science          Wei Chen and Xiao Li
#>           year                 journal                  doi
#> smith2023 2023 Journal of Informetrics 10.1234/example.2023
#> chen2022  2022          Scientometrics 10.1234/example.2022
\end{verbatim}

\subsection{Parsing WoS plain text with bibliometrix}\label{parsing-wos-plain-text-with-bibliometrix}

\begin{Shaded}
\begin{Highlighting}[]
\CommentTok{\# Requires a WoS export file — not bundled}
\CommentTok{\# wos\_df \textless{}{-} convert2df("savedrecs.txt", dbsource = "wos", format = "plaintext")}
\CommentTok{\# glimpse(wos\_df)}
\end{Highlighting}
\end{Shaded}

\subsection{Standardizing field names}\label{standardizing-field-names}

\begin{Shaded}
\begin{Highlighting}[]
\NormalTok{standardize\_bib }\OtherTok{\textless{}{-}} \ControlFlowTok{function}\NormalTok{(df) \{}
\NormalTok{  df }\SpecialCharTok{|\textgreater{}}
    \FunctionTok{transmute}\NormalTok{(}
      \AttributeTok{title =}\NormalTok{ TITLE,}
      \AttributeTok{authors =}\NormalTok{ AUTHOR,}
      \AttributeTok{year =} \FunctionTok{as.integer}\NormalTok{(YEAR),}
      \AttributeTok{journal =}\NormalTok{ JOURNAL,}
      \AttributeTok{doi =}\NormalTok{ DOI}
\NormalTok{    )}
\NormalTok{\}}

\NormalTok{bib\_df }\SpecialCharTok{|\textgreater{}} \FunctionTok{standardize\_bib}\NormalTok{()}
\end{Highlighting}
\end{Shaded}

\begin{verbatim}
#> # A tibble: 2 x 5
#>   title                                 authors    year journal            doi  
#>   <chr>                                 <list>    <int> <chr>              <chr>
#> 1 Advances in Bibliometric Methods      <chr [2]>  2023 Journal of Inform~ 10.1~
#> 2 Citation Networks in Computer Science <chr [2]>  2022 Scientometrics     10.1~
\end{verbatim}

\subsection{Visualization}\label{visualization-2}

\begin{Shaded}
\begin{Highlighting}[]
\NormalTok{completeness }\OtherTok{\textless{}{-}}\NormalTok{ bib\_df }\SpecialCharTok{|\textgreater{}}
  \FunctionTok{summarise}\NormalTok{(}\FunctionTok{across}\NormalTok{(}\FunctionTok{everything}\NormalTok{(), }\SpecialCharTok{\textasciitilde{}} \FunctionTok{mean}\NormalTok{(}\SpecialCharTok{!}\FunctionTok{is.na}\NormalTok{(.)))) }\SpecialCharTok{|\textgreater{}}
  \FunctionTok{pivot\_longer}\NormalTok{(}\FunctionTok{everything}\NormalTok{(), }\AttributeTok{names\_to =} \StringTok{"field"}\NormalTok{, }\AttributeTok{values\_to =} \StringTok{"present"}\NormalTok{) }\SpecialCharTok{|\textgreater{}}
  \FunctionTok{filter}\NormalTok{(present }\SpecialCharTok{\textgreater{}} \DecValTok{0}\NormalTok{) }\SpecialCharTok{|\textgreater{}}
  \FunctionTok{mutate}\NormalTok{(}\AttributeTok{field =} \FunctionTok{fct\_reorder}\NormalTok{(field, present))}

\FunctionTok{ggplot}\NormalTok{(completeness, }\FunctionTok{aes}\NormalTok{(}\AttributeTok{x =}\NormalTok{ present, }\AttributeTok{y =}\NormalTok{ field)) }\SpecialCharTok{+}
  \FunctionTok{geom\_col}\NormalTok{(}\AttributeTok{fill =} \FunctionTok{palette\_sci}\NormalTok{(}\DecValTok{1}\NormalTok{)) }\SpecialCharTok{+}
  \FunctionTok{scale\_x\_continuous}\NormalTok{(}\AttributeTok{labels =}\NormalTok{ scales}\SpecialCharTok{::}\NormalTok{percent) }\SpecialCharTok{+}
  \FunctionTok{labs}\NormalTok{(}\AttributeTok{x =} \StringTok{"Completeness"}\NormalTok{, }\AttributeTok{y =} \ConstantTok{NULL}\NormalTok{) }\SpecialCharTok{+}
  \FunctionTok{theme\_sci}\NormalTok{()}
\end{Highlighting}
\end{Shaded}

\begin{figure}

{\centering \includegraphics[width=0.9\linewidth,alt={Horizontal bar chart showing which metadata fields are present in the parsed BibTeX records.}]{scientometrics-in-r_files/figure-latex/06-fig-parsing-fields-1} 

}

\caption{Field completeness in the parsed BibTeX records.}\label{fig:06-fig-parsing-fields}
\end{figure}

\section{Diagnostics and interpretation}\label{parsing-diagnostics}

After parsing, always check:

\begin{itemize}
\tightlist
\item
  \textbf{Row count}: Does the number of records match the export? Missing entries may indicate parsing errors.
\item
  \textbf{Field completeness}: Which fields are populated? Abstracts and keywords are often missing in BibTeX exports.
\item
  \textbf{Author format}: Names may appear as ``Last, First'' or ``First Last'' depending on the source. Standardize before analysis.
\item
  \textbf{Encoding}: Non-ASCII characters (diacritics, CJK) can corrupt during parsing. Ensure UTF-8 encoding.
\end{itemize}

\section{Limitations and responsible use}\label{parsing-limitations}

\section{Limitations and responsible use}\label{limitations-and-responsible-use-5}

\begin{itemize}
\tightlist
\item
  Export files are snapshots: they reflect the database state at export time and become stale immediately.
\item
  Field mapping between formats is imperfect. Some formats lack fields (RIS has no standard abstract tag across all databases).
\item
  Parsing errors are silent --- always validate record counts and spot-check metadata against the source \citep{hicks2015}.
\end{itemize}

\section{Common pitfalls}\label{parsing-pitfalls}

\section{Common pitfalls}\label{common-pitfalls-5}

\begin{itemize}
\tightlist
\item
  \textbf{Encoding mismatches}: Exporting as Latin-1 and reading as UTF-8 (or vice versa) corrupts names and titles.
\item
  \textbf{Truncated exports}: Many databases cap exports at 500 or 1,000 records per file. Large corpora require multiple export batches.
\item
  \textbf{Inconsistent author delimiters}: BibTeX uses ``and'', RIS uses newlines, CSV uses semicolons. Parsing must handle each.
\item
  \textbf{Assuming format from extension}: A \texttt{.txt} file may be WoS plain text or something else entirely. Check the first few lines.
\end{itemize}

\section{Exercises}\label{parsing-exercises}

\begin{enumerate}
\def\labelenumi{\arabic{enumi}.}
\item
  \textbf{Export and parse.} Export 50 records from Google Scholar as BibTeX. Parse them with \texttt{bib2df}. How many fields are populated?
\item
  \textbf{Format comparison.} Export the same set of records from Scopus as both CSV and RIS. Parse both and compare the fields available.
\item
  \textbf{Standardization function.} Write a function that takes a data frame from \texttt{convert2df()} and returns a tibble with columns: \texttt{doi}, \texttt{title}, \texttt{authors}, \texttt{year}, \texttt{journal}, \texttt{cited\_by}.
\end{enumerate}

\section{Solutions}\label{06-solutions}

Solutions are provided in \ref{solutions}.

\section{Further reading}\label{parsing-reading}

\begin{itemize}
\tightlist
\item
  \citet{aria2017} --- The \texttt{bibliometrix} package, including \texttt{convert2df()} for multi-format parsing.
\item
  \citet{priem2022} --- OpenAlex as an API-first alternative to file-based workflows.
\end{itemize}

\section{Session info}\label{parsing-sessioninfo}

\begin{Shaded}
\begin{Highlighting}[]
\FunctionTok{sessionInfo}\NormalTok{()}
\end{Highlighting}
\end{Shaded}

\begin{verbatim}
#> R version 4.6.1 (2026-06-24)
#> Platform: x86_64-pc-linux-gnu
#> Running under: Ubuntu 24.04.5 LTS
#> 
#> Matrix products: default
#> BLAS:   /usr/lib/x86_64-linux-gnu/openblas-pthread/libblas.so.3 
#> LAPACK: /usr/lib/x86_64-linux-gnu/openblas-pthread/libopenblasp-r0.3.26.so;  LAPACK version 3.12.0
#> 
#> locale:
#>  [1] LC_CTYPE=C.UTF-8       LC_NUMERIC=C           LC_TIME=C.UTF-8       
#>  [4] LC_COLLATE=C.UTF-8     LC_MONETARY=C.UTF-8    LC_MESSAGES=C.UTF-8   
#>  [7] LC_PAPER=C.UTF-8       LC_NAME=C              LC_ADDRESS=C          
#> [10] LC_TELEPHONE=C         LC_MEASUREMENT=C.UTF-8 LC_IDENTIFICATION=C   
#> 
#> time zone: UTC
#> tzcode source: system (glibc)
#> 
#> attached base packages:
#> [1] stats     graphics  grDevices utils     datasets  methods   base     
#> 
#> other attached packages:
#>  [1] bibliometrix_5.5.0 RefManageR_1.4.0   bib2df_1.1.2.0     rcrossref_1.2.1   
#>  [5] gt_1.4.0           tidytext_0.4.3     glue_1.8.1         openalexR_3.1.2   
#>  [9] lubridate_1.9.5    forcats_1.0.1      stringr_1.6.0      dplyr_1.2.1       
#> [13] purrr_1.2.2        readr_2.2.0        tidyr_1.3.2        tibble_3.3.1      
#> [17] ggplot2_4.0.3      tidyverse_2.0.0   
#> 
#> loaded via a namespace (and not attached):
#>  [1] gridExtra_2.3.1        httr2_1.3.0            readxl_1.6.0          
#>  [4] rlang_1.3.0            magrittr_2.0.5         otel_0.2.0            
#>  [7] compiler_4.6.1         systemfonts_1.3.2      vctrs_0.7.3           
#> [10] httpcode_0.3.0         pkgconfig_2.0.3        fastmap_1.2.0         
#> [13] labeling_0.4.3         ca_0.72                utf8_1.2.6            
#> [16] promises_1.5.0         rmarkdown_2.32         tzdb_0.5.0            
#> [19] xfun_0.62              jsonlite_2.0.0         SnowballC_0.7.1       
#> [22] later_1.4.8            parallel_4.6.1         stopwords_2.3         
#> [25] R6_2.6.1               stringi_1.8.9          RColorBrewer_1.1-3    
#> [28] cellranger_1.1.0       Rcpp_1.1.2             bookdown_0.48         
#> [31] knitr_1.52             triebeard_0.4.1        base64enc_0.1-6       
#> [34] httpuv_1.6.17          rentrez_1.2.4          Matrix_1.7-5          
#> [37] igraph_2.3.4           timechange_0.4.0       tidyselect_1.2.1      
#> [40] stringdist_0.9.17      rstudioapi_0.19.0      pubmedR_1.0.2         
#> [43] yaml_2.3.12            viridis_0.6.5          miniUI_0.1.2          
#> [46] humaniformat_0.6.0     curl_8.1.0             qpdf_1.4.1            
#> [49] lattice_0.22-9         plyr_1.8.9             shiny_1.14.0          
#> [52] withr_3.0.3            S7_0.2.2               askpass_1.2.1         
#> [55] evaluate_1.0.5         zip_3.0.2              xml2_1.6.0            
#> [58] shinycssloaders_1.1.0  pillar_1.11.1          janeaustenr_1.0.0     
#> [61] DT_0.34.0              plotly_4.12.1          generics_0.1.4        
#> [64] rprojroot_2.1.1        hms_1.1.4              scales_1.4.0          
#> [67] xtable_1.8-8           contentanalysis_1.1.1  tools_4.6.1           
#> [70] data.table_1.18.6.1    tokenizers_0.3.0       openxlsx_4.2.9        
#> [73] pdftools_3.9.1         fs_2.1.0               visNetwork_2.1.4      
#> [76] XML_3.99-0.25          grid_4.6.1             bibtex_0.5.3          
#> [79] urltools_1.7.3.1       rscopus_0.9.0          dimensionsR_0.0.3     
#> [82] bibliometrixData_0.3.0 cli_3.6.6              viridisLite_0.4.3     
#> [85] gtable_0.3.6           digest_0.6.39          ggrepel_0.9.8         
#> [88] crul_1.6.0             htmlwidgets_1.6.4      farver_2.1.2          
#> [91] htmltools_0.5.9        lifecycle_1.0.5        httr_1.4.9            
#> [94] here_1.0.2             mime_0.13
\end{verbatim}

\chapter{Building Reproducible Corpora}\label{building-reproducible-corpora}

\emph{R. Heller · Date: 2026-05-09 · Last changed: 2026-10-31}

\section{Learning objectives}\label{corpora-objectives}

After completing this chapter, you will be able to:

\begin{itemize}
\tightlist
\item
  Deduplicate bibliographic records by DOI and by approximate title matching
\item
  Resolve identifiers across databases (OpenAlex ID, DOI, PMID)
\item
  Recognize author name disambiguation challenges and apply basic heuristics
\item
  Standardize institutional affiliations using ROR identifiers
\item
  Save a clean corpus in an efficient format (Parquet) for downstream analysis
\end{itemize}

\section{Setup}\label{corpora-setup}

\begin{Shaded}
\begin{Highlighting}[]
\FunctionTok{library}\NormalTok{(tidyverse)}
\FunctionTok{library}\NormalTok{(openalexR)}
\FunctionTok{library}\NormalTok{(glue)}
\FunctionTok{library}\NormalTok{(arrow)}

\FunctionTok{set.seed}\NormalTok{(}\DecValTok{20260509}\NormalTok{)}

\FunctionTok{source}\NormalTok{(here}\SpecialCharTok{::}\FunctionTok{here}\NormalTok{(}\StringTok{"R"}\NormalTok{, }\StringTok{"api\_helpers.R"}\NormalTok{))}
\FunctionTok{source}\NormalTok{(here}\SpecialCharTok{::}\FunctionTok{here}\NormalTok{(}\StringTok{"R"}\NormalTok{, }\StringTok{"utils.R"}\NormalTok{))}
\FunctionTok{source}\NormalTok{(here}\SpecialCharTok{::}\FunctionTok{here}\NormalTok{(}\StringTok{"R"}\NormalTok{, }\StringTok{"sci\_palette.R"}\NormalTok{))}
\end{Highlighting}
\end{Shaded}

\section{Conceptual background}\label{corpora-background}

Raw bibliographic data is messy. Records from different sources may describe the same paper with different metadata, the same author under different name variants, and the same institution under dozens of spelling variations. Building a clean, reproducible corpus requires systematic attention to three problems:

\textbf{Deduplication.} When merging data from multiple sources, the same publication may appear multiple times. DOI-based deduplication is the gold standard, but not all records have DOIs. Fuzzy title matching provides a fallback but is error-prone for short or generic titles.

\textbf{Author disambiguation.} ``J. Smith'' could be dozens of different people. OpenAlex uses machine-learning-based author clustering to assign persistent author IDs, but errors persist --- especially for common names and authors who change institutions \citep{priem2022}.

\textbf{Affiliation standardization.} Institutional names appear in countless variants (``MIT'', ``Massachusetts Institute of Technology'', ``Mass. Inst. Tech.''). The Research Organization Registry (ROR) provides a curated set of persistent identifiers for research organizations. OpenAlex maps many affiliations to ROR IDs, but coverage is incomplete.

\section{Worked example}\label{corpora-example}

\subsection{Fetching raw data}\label{fetching-raw-data}

\begin{Shaded}
\begin{Highlighting}[]
\NormalTok{works }\OtherTok{\textless{}{-}} \FunctionTok{fetch\_openalex}\NormalTok{(}
  \AttributeTok{entity =} \StringTok{"works"}\NormalTok{,}
  \AttributeTok{search =} \StringTok{"bibliometrics"}\NormalTok{,}
  \AttributeTok{from\_publication\_date =} \StringTok{"2021{-}01{-}01"}\NormalTok{,}
  \AttributeTok{to\_publication\_date =} \StringTok{"2023{-}12{-}31"}\NormalTok{,}
  \AttributeTok{options =} \FunctionTok{list}\NormalTok{(}\AttributeTok{sample =} \DecValTok{300}\NormalTok{, }\AttributeTok{seed =} \DecValTok{42}\NormalTok{)}
\NormalTok{)}

\FunctionTok{cat}\NormalTok{(}\FunctionTok{glue}\NormalTok{(}\StringTok{"Raw records: \{nrow(works)\}}\SpecialCharTok{\textbackslash{}n}\StringTok{"}\NormalTok{))}
\end{Highlighting}
\end{Shaded}

\begin{verbatim}
#> Raw records: 300
\end{verbatim}

\subsection{DOI-based deduplication}\label{doi-based-deduplication}

\begin{Shaded}
\begin{Highlighting}[]
\NormalTok{works\_deduped }\OtherTok{\textless{}{-}} \FunctionTok{dedupe\_by\_doi}\NormalTok{(works)}
\FunctionTok{cat}\NormalTok{(}\FunctionTok{glue}\NormalTok{(}\StringTok{"After DOI dedup: \{nrow(works\_deduped)\}}\SpecialCharTok{\textbackslash{}n}\StringTok{"}\NormalTok{))}
\end{Highlighting}
\end{Shaded}

\begin{verbatim}
#> After DOI dedup: 300
\end{verbatim}

\begin{Shaded}
\begin{Highlighting}[]
\FunctionTok{cat}\NormalTok{(}\FunctionTok{glue}\NormalTok{(}\StringTok{"Removed: \{nrow(works) {-} nrow(works\_deduped)\} duplicates}\SpecialCharTok{\textbackslash{}n}\StringTok{"}\NormalTok{))}
\end{Highlighting}
\end{Shaded}

\begin{verbatim}
#> Removed: 0 duplicates
\end{verbatim}

\subsection{Fuzzy title deduplication}\label{fuzzy-title-deduplication}

For records without DOIs, we use approximate string matching.

\begin{Shaded}
\begin{Highlighting}[]
\NormalTok{no\_doi }\OtherTok{\textless{}{-}}\NormalTok{ works\_deduped }\SpecialCharTok{|\textgreater{}} \FunctionTok{filter}\NormalTok{(}\FunctionTok{is.na}\NormalTok{(doi))}

\ControlFlowTok{if}\NormalTok{ (}\FunctionTok{nrow}\NormalTok{(no\_doi) }\SpecialCharTok{\textgreater{}} \DecValTok{1}\NormalTok{) \{}
\NormalTok{  title\_lower }\OtherTok{\textless{}{-}} \FunctionTok{tolower}\NormalTok{(no\_doi}\SpecialCharTok{$}\NormalTok{display\_name)}
\NormalTok{  dist\_matrix }\OtherTok{\textless{}{-}}\NormalTok{ stringdist}\SpecialCharTok{::}\FunctionTok{stringdistmatrix}\NormalTok{(title\_lower, }\AttributeTok{method =} \StringTok{"jw"}\NormalTok{)}
\NormalTok{  potential\_dupes }\OtherTok{\textless{}{-}} \FunctionTok{which}\NormalTok{(}\FunctionTok{as.matrix}\NormalTok{(dist\_matrix) }\SpecialCharTok{\textless{}} \FloatTok{0.1} \SpecialCharTok{\&}
                            \FunctionTok{as.matrix}\NormalTok{(dist\_matrix) }\SpecialCharTok{\textgreater{}} \DecValTok{0}\NormalTok{, }\AttributeTok{arr.ind =} \ConstantTok{TRUE}\NormalTok{)}
\NormalTok{  potential\_dupes }\OtherTok{\textless{}{-}}\NormalTok{ potential\_dupes[potential\_dupes[, }\DecValTok{1}\NormalTok{] }\SpecialCharTok{\textless{}}\NormalTok{ potential\_dupes[, }\DecValTok{2}\NormalTok{], , drop }\OtherTok{=} \ConstantTok{FALSE}\NormalTok{]}
  \FunctionTok{cat}\NormalTok{(}\FunctionTok{glue}\NormalTok{(}\StringTok{"Potential fuzzy duplicates (no DOI): \{nrow(potential\_dupes)\} pairs}\SpecialCharTok{\textbackslash{}n}\StringTok{"}\NormalTok{))}
\NormalTok{\} }\ControlFlowTok{else}\NormalTok{ \{}
  \FunctionTok{cat}\NormalTok{(}\StringTok{"No records without DOI to fuzzy{-}match.}\SpecialCharTok{\textbackslash{}n}\StringTok{"}\NormalTok{)}
\NormalTok{\}}
\end{Highlighting}
\end{Shaded}

\begin{verbatim}
#> Potential fuzzy duplicates (no DOI): 0 pairs
\end{verbatim}

\subsection{Affiliation extraction and ROR mapping}\label{affiliation-extraction-and-ror-mapping}

\begin{Shaded}
\begin{Highlighting}[]
\NormalTok{affiliations }\OtherTok{\textless{}{-}}\NormalTok{ works\_deduped }\SpecialCharTok{|\textgreater{}}
  \FunctionTok{select}\NormalTok{(id, authorships) }\SpecialCharTok{|\textgreater{}}
  \FunctionTok{unnest}\NormalTok{(authorships, }\AttributeTok{names\_sep =} \StringTok{"\_"}\NormalTok{) }\SpecialCharTok{|\textgreater{}}
  \FunctionTok{unnest}\NormalTok{(authorships\_affiliations, }\AttributeTok{names\_sep =} \StringTok{"\_"}\NormalTok{) }\SpecialCharTok{|\textgreater{}}
  \FunctionTok{select}\NormalTok{(}\AttributeTok{work\_id =}\NormalTok{ id,}
         \AttributeTok{author\_name =}\NormalTok{ authorships\_display\_name,}
         \AttributeTok{institution =}\NormalTok{ authorships\_affiliations\_display\_name) }\SpecialCharTok{|\textgreater{}}
  \FunctionTok{filter}\NormalTok{(}\SpecialCharTok{!}\FunctionTok{is.na}\NormalTok{(institution))}

\NormalTok{top\_institutions }\OtherTok{\textless{}{-}}\NormalTok{ affiliations }\SpecialCharTok{|\textgreater{}}
  \FunctionTok{count}\NormalTok{(institution, }\AttributeTok{sort =} \ConstantTok{TRUE}\NormalTok{) }\SpecialCharTok{|\textgreater{}}
  \FunctionTok{head}\NormalTok{(}\DecValTok{15}\NormalTok{)}

\NormalTok{top\_institutions}
\end{Highlighting}
\end{Shaded}

\begin{verbatim}
#> # A tibble: 15 x 2
#>    institution                                          n
#>    <chr>                                            <int>
#>  1 Washington University in St. Louis                  41
#>  2 Jagiellonian University                             13
#>  3 Chinese People's Liberation Army                    12
#>  4 Wuhan University of Technology                      12
#>  5 Chinese PLA General Hospital                        11
#>  6 Hong Kong Baptist University                        10
#>  7 Universidade Federal do Rio Grande                  10
#>  8 Aga Khan University                                  9
#>  9 Gansu University of Traditional Chinese Medicine     9
#> 10 Ospedale Cannizzaro                                  9
#> 11 University of Catania                                9
#> 12 Fudan University                                     8
#> 13 Imam Mohammad ibn Saud Islamic University            8
#> 14 Indian Institute of Technology Guwahati              8
#> 15 State Key Laboratory of Kidney Diseases              8
\end{verbatim}

\subsection{Saving as Parquet}\label{saving-as-parquet}

\begin{Shaded}
\begin{Highlighting}[]
\NormalTok{corpus\_clean }\OtherTok{\textless{}{-}}\NormalTok{ works\_deduped }\SpecialCharTok{|\textgreater{}}
  \FunctionTok{select}\NormalTok{(id, display\_name, publication\_date, cited\_by\_count, doi, source\_display\_name, abstract)}

\NormalTok{out\_dir }\OtherTok{\textless{}{-}}\NormalTok{ here}\SpecialCharTok{::}\FunctionTok{here}\NormalTok{(}\StringTok{"\_freeze"}\NormalTok{, }\StringTok{"exports"}\NormalTok{)}
\FunctionTok{dir.create}\NormalTok{(out\_dir, }\AttributeTok{recursive =} \ConstantTok{TRUE}\NormalTok{, }\AttributeTok{showWarnings =} \ConstantTok{FALSE}\NormalTok{)}
\NormalTok{out\_path }\OtherTok{\textless{}{-}} \FunctionTok{file.path}\NormalTok{(out\_dir, }\StringTok{"corpus\_sample.parquet"}\NormalTok{)}
\FunctionTok{write\_parquet}\NormalTok{(corpus\_clean, out\_path)}
\FunctionTok{cat}\NormalTok{(}\FunctionTok{glue}\NormalTok{(}\StringTok{"Saved \{nrow(corpus\_clean)\} records to \{out\_path\}}\SpecialCharTok{\textbackslash{}n}\StringTok{"}\NormalTok{))}
\end{Highlighting}
\end{Shaded}

\begin{verbatim}
#> Saved 300 records to /home/runner/work/scientometrics-in-r/scientometrics-in-r/_freeze/exports/corpus_sample.parquet
\end{verbatim}

\subsection{Visualization}\label{visualization-3}

\begin{Shaded}
\begin{Highlighting}[]
\NormalTok{top\_institutions }\SpecialCharTok{|\textgreater{}}
  \FunctionTok{mutate}\NormalTok{(}\AttributeTok{institution =} \FunctionTok{fct\_reorder}\NormalTok{(institution, n)) }\SpecialCharTok{|\textgreater{}}
  \FunctionTok{ggplot}\NormalTok{(}\FunctionTok{aes}\NormalTok{(}\AttributeTok{x =}\NormalTok{ n, }\AttributeTok{y =}\NormalTok{ institution)) }\SpecialCharTok{+}
  \FunctionTok{geom\_col}\NormalTok{(}\AttributeTok{fill =} \FunctionTok{palette\_sci}\NormalTok{(}\DecValTok{1}\NormalTok{)) }\SpecialCharTok{+}
  \FunctionTok{labs}\NormalTok{(}\AttributeTok{x =} \StringTok{"Publications"}\NormalTok{, }\AttributeTok{y =} \ConstantTok{NULL}\NormalTok{) }\SpecialCharTok{+}
  \FunctionTok{theme\_sci}\NormalTok{()}
\end{Highlighting}
\end{Shaded}

\begin{figure}

{\centering \includegraphics[width=0.9\linewidth,alt={Horizontal bar chart showing the 15 most frequent institutional affiliations in the cleaned corpus.}]{scientometrics-in-r_files/figure-latex/07-fig-corpora-institutions-1} 

}

\caption{Top 15 institutions by publication count in the sample corpus.}\label{fig:07-fig-corpora-institutions}
\end{figure}

\section{Diagnostics and interpretation}\label{corpora-diagnostics}

After cleaning:

\begin{itemize}
\tightlist
\item
  \textbf{Deduplication rate}: A rate above 5\% suggests either a multi-source merge or a query that returns overlapping result pages.
\item
  \textbf{DOI coverage}: The proportion of records with DOIs determines how effective DOI-based deduplication can be.
\item
  \textbf{Affiliation coverage}: Check what fraction of author records have institutional affiliations. Low coverage limits institutional analysis.
\item
  \textbf{Parquet file size}: Parquet is columnar and compressed. A 10,000-record corpus typically compresses to under 1 MB.
\end{itemize}

\section{Limitations and responsible use}\label{corpora-limitations}

\section{Limitations and responsible use}\label{limitations-and-responsible-use-6}

\begin{itemize}
\tightlist
\item
  \textbf{Fuzzy matching is imperfect.} Jaro-Winkler similarity can produce both false positives (different papers with similar titles) and false negatives (same paper with different title versions). Always spot-check.
\item
  \textbf{Author disambiguation remains unsolved.} No automated method achieves perfect accuracy. Report the disambiguation method used and its known error rate.
\item
  \textbf{Affiliation data is noisy.} Even with ROR mapping, temporary affiliations, joint appointments, and name changes create ambiguity. Never use affiliation data as ground truth without validation \citep{hicks2015}.
\end{itemize}

\section{Common pitfalls}\label{corpora-pitfalls}

\section{Common pitfalls}\label{common-pitfalls-6}

\begin{itemize}
\tightlist
\item
  \textbf{Deduplicating before merging.} Always merge first, then deduplicate. Deduplicating each source independently misses cross-source duplicates.
\item
  \textbf{Trusting DOIs blindly.} Some records have incorrect DOIs (typos, test DOIs, placeholder values). Validate a sample.
\item
  \textbf{Ignoring records without DOIs.} Discarding them biases the corpus toward recent, well-indexed publications.
\item
  \textbf{Not documenting cleaning steps.} Every deduplication or name-standardization step should be logged for reproducibility.
\end{itemize}

\section{Exercises}\label{corpora-exercises}

\begin{enumerate}
\def\labelenumi{\arabic{enumi}.}
\item
  \textbf{DOI coverage by year.} Compute the DOI coverage rate by publication year in the sample. Is there a trend?
\item
  \textbf{Name variants.} Find authors in the sample who appear under multiple name spellings. What heuristics could you use to merge them?
\item
  \textbf{ROR lookup.} Use the ROR API to resolve the top 5 institutions in your corpus to their official ROR identifiers.
\end{enumerate}

\section{Solutions}\label{07-solutions}

Solutions are provided in \ref{solutions}.

\section{Further reading}\label{corpora-reading}

\begin{itemize}
\tightlist
\item
  \citet{priem2022} --- OpenAlex author disambiguation and institutional mapping.
\item
  \citet{aria2017} --- \texttt{bibliometrix} data cleaning utilities.
\item
  \citet{hicks2015} --- The Leiden Manifesto; principles 4 and 5 emphasise data transparency and verifiability.
\end{itemize}

\section{Session info}\label{corpora-sessioninfo}

\begin{Shaded}
\begin{Highlighting}[]
\FunctionTok{sessionInfo}\NormalTok{()}
\end{Highlighting}
\end{Shaded}

\begin{verbatim}
#> R version 4.6.1 (2026-06-24)
#> Platform: x86_64-pc-linux-gnu
#> Running under: Ubuntu 24.04.5 LTS
#> 
#> Matrix products: default
#> BLAS:   /usr/lib/x86_64-linux-gnu/openblas-pthread/libblas.so.3 
#> LAPACK: /usr/lib/x86_64-linux-gnu/openblas-pthread/libopenblasp-r0.3.26.so;  LAPACK version 3.12.0
#> 
#> locale:
#>  [1] LC_CTYPE=C.UTF-8       LC_NUMERIC=C           LC_TIME=C.UTF-8       
#>  [4] LC_COLLATE=C.UTF-8     LC_MONETARY=C.UTF-8    LC_MESSAGES=C.UTF-8   
#>  [7] LC_PAPER=C.UTF-8       LC_NAME=C              LC_ADDRESS=C          
#> [10] LC_TELEPHONE=C         LC_MEASUREMENT=C.UTF-8 LC_IDENTIFICATION=C   
#> 
#> time zone: UTC
#> tzcode source: system (glibc)
#> 
#> attached base packages:
#> [1] stats     graphics  grDevices utils     datasets  methods   base     
#> 
#> other attached packages:
#>  [1] arrow_26.0.0       bibliometrix_5.5.0 RefManageR_1.4.0   bib2df_1.1.2.0    
#>  [5] rcrossref_1.2.1    gt_1.4.0           tidytext_0.4.3     glue_1.8.1        
#>  [9] openalexR_3.1.2    lubridate_1.9.5    forcats_1.0.1      stringr_1.6.0     
#> [13] dplyr_1.2.1        purrr_1.2.2        readr_2.2.0        tidyr_1.3.2       
#> [17] tibble_3.3.1       ggplot2_4.0.3      tidyverse_2.0.0   
#> 
#> loaded via a namespace (and not attached):
#>  [1] gridExtra_2.3.1        httr2_1.3.0            readxl_1.6.0          
#>  [4] rlang_1.3.0            magrittr_2.0.5         otel_0.2.0            
#>  [7] compiler_4.6.1         systemfonts_1.3.2      vctrs_0.7.3           
#> [10] httpcode_0.3.0         pkgconfig_2.0.3        fastmap_1.2.0         
#> [13] labeling_0.4.3         ca_0.72                utf8_1.2.6            
#> [16] promises_1.5.0         rmarkdown_2.32         tzdb_0.5.0            
#> [19] bit_4.6.0              xfun_0.62              jsonlite_2.0.0        
#> [22] SnowballC_0.7.1        later_1.4.8            parallel_4.6.1        
#> [25] stopwords_2.3          R6_2.6.1               stringi_1.8.9         
#> [28] RColorBrewer_1.1-3     cellranger_1.1.0       assertthat_0.2.1      
#> [31] Rcpp_1.1.2             bookdown_0.48          knitr_1.52            
#> [34] triebeard_0.4.1        base64enc_0.1-6        httpuv_1.6.17         
#> [37] rentrez_1.2.4          Matrix_1.7-5           igraph_2.3.4          
#> [40] timechange_0.4.0       tidyselect_1.2.1       stringdist_0.9.17     
#> [43] rstudioapi_0.19.0      pubmedR_1.0.2          yaml_2.3.12           
#> [46] viridis_0.6.5          miniUI_0.1.2           humaniformat_0.6.0    
#> [49] curl_8.1.0             qpdf_1.4.1             lattice_0.22-9        
#> [52] plyr_1.8.9             shiny_1.14.0           withr_3.0.3           
#> [55] S7_0.2.2               askpass_1.2.1          evaluate_1.0.5        
#> [58] zip_3.0.2              xml2_1.6.0             shinycssloaders_1.1.0 
#> [61] pillar_1.11.1          janeaustenr_1.0.0      DT_0.34.0             
#> [64] plotly_4.12.1          generics_0.1.4         rprojroot_2.1.1       
#> [67] hms_1.1.4              scales_1.4.0           xtable_1.8-8          
#> [70] contentanalysis_1.1.1  tools_4.6.1            data.table_1.18.6.1   
#> [73] tokenizers_0.3.0       openxlsx_4.2.9         pdftools_3.9.1        
#> [76] fs_2.1.0               visNetwork_2.1.4       XML_3.99-0.25         
#> [79] grid_4.6.1             bibtex_0.5.3           urltools_1.7.3.1      
#> [82] rscopus_0.9.0          dimensionsR_0.0.3      bibliometrixData_0.3.0
#> [85] cli_3.6.6              viridisLite_0.4.3      gtable_0.3.6          
#> [88] digest_0.6.39          ggrepel_0.9.8          crul_1.6.0            
#> [91] htmlwidgets_1.6.4      farver_2.1.2           htmltools_0.5.9       
#> [94] lifecycle_1.0.5        httr_1.4.9             here_1.0.2            
#> [97] mime_0.13              bit64_4.8.6
\end{verbatim}

\chapter{Working with Full Text}\label{working-with-full-text}

\emph{R. Heller · Date: 2026-05-09 · Last changed: 2026-10-31}

\section{Learning objectives}\label{fulltext-objectives}

After completing this chapter, you will be able to:

\begin{itemize}
\tightlist
\item
  Extract text from PDF files using \texttt{pdftools}
\item
  Apply OCR to scanned documents with \texttt{tesseract}
\item
  Perform basic text preprocessing (tokenization, stopword removal)
\item
  Assess when full-text analysis adds value over metadata-only approaches
\item
  Navigate the legal landscape of text and data mining
\end{itemize}

\section{Setup}\label{fulltext-setup}

\begin{Shaded}
\begin{Highlighting}[]
\FunctionTok{library}\NormalTok{(tidyverse)}
\FunctionTok{library}\NormalTok{(pdftools)}
\FunctionTok{library}\NormalTok{(quanteda)}
\FunctionTok{library}\NormalTok{(glue)}

\FunctionTok{set.seed}\NormalTok{(}\DecValTok{20260509}\NormalTok{)}

\FunctionTok{source}\NormalTok{(here}\SpecialCharTok{::}\FunctionTok{here}\NormalTok{(}\StringTok{"R"}\NormalTok{, }\StringTok{"sci\_palette.R"}\NormalTok{))}
\end{Highlighting}
\end{Shaded}

\section{Conceptual background}\label{fulltext-background}

Most bibliometric analyses operate on metadata: titles, abstracts, keywords, and citation links. But metadata captures only a fraction of a paper's content. Full-text analysis opens the door to richer methods: section-level citation context, method extraction, data-availability statements, and fine-grained topic modeling.

The technical pipeline involves three stages: (1) \textbf{acquisition} --- obtaining PDF files, subject to access rights and TDM (text and data mining) provisions; (2) \textbf{extraction} --- converting PDF content to machine-readable text; and (3) \textbf{preprocessing} --- tokenization, normalization, and feature extraction.

\texttt{pdftools} wraps the Poppler library to extract text from born-digital PDFs. For scanned documents, \texttt{tesseract} provides OCR (optical character recognition). \texttt{tabulizer} (based on Tabula) specializes in extracting tables from PDFs. For text preprocessing, \texttt{quanteda} and \texttt{tm} offer comprehensive toolkits.

Legal frameworks vary by jurisdiction. The EU Copyright Directive (2019/790) provides a TDM exception for research. In the US, TDM is generally considered fair use in research contexts, but publisher terms of service may impose additional restrictions.

\section{Worked example}\label{fulltext-example}

\subsection{Extracting text from a PDF}\label{extracting-text-from-a-pdf}

We demonstrate with a programmatically created PDF to keep this example self-contained.

\begin{Shaded}
\begin{Highlighting}[]
\NormalTok{demo\_text }\OtherTok{\textless{}{-}} \FunctionTok{paste}\NormalTok{(}
  \StringTok{"Title: A Review of Bibliometric Methods"}\NormalTok{,}
  \StringTok{""}\NormalTok{,}
  \StringTok{"Abstract: This paper reviews current methods in bibliometrics."}\NormalTok{,}
  \StringTok{"We focus on citation analysis, co{-}authorship networks, and topic modeling."}\NormalTok{,}
  \StringTok{""}\NormalTok{,}
  \StringTok{"1. Introduction"}\NormalTok{,}
  \StringTok{"Bibliometrics has grown rapidly since the work of Garfield (1955)."}\NormalTok{,}
  \StringTok{"Modern tools enable large{-}scale analysis of scholarly communication."}\NormalTok{,}
  \StringTok{""}\NormalTok{,}
  \StringTok{"2. Methods"}\NormalTok{,}
  \StringTok{"We surveyed 500 papers published between 2020 and 2023."}\NormalTok{,}
  \StringTok{"Data were obtained from OpenAlex."}\NormalTok{,}
  \StringTok{""}\NormalTok{,}
  \StringTok{"3. Conclusions"}\NormalTok{,}
  \StringTok{"The field continues to evolve with open data and new computational methods."}\NormalTok{,}
  \AttributeTok{sep =} \StringTok{"}\SpecialCharTok{\textbackslash{}n}\StringTok{"}
\NormalTok{)}

\NormalTok{tmp\_pdf }\OtherTok{\textless{}{-}} \FunctionTok{tempfile}\NormalTok{(}\AttributeTok{fileext =} \StringTok{".pdf"}\NormalTok{)}
\FunctionTok{pdf}\NormalTok{(tmp\_pdf, }\AttributeTok{width =} \FloatTok{8.5}\NormalTok{, }\AttributeTok{height =} \DecValTok{11}\NormalTok{)}
\FunctionTok{plot.new}\NormalTok{()}
\FunctionTok{text}\NormalTok{(}\FloatTok{0.05}\NormalTok{, }\FloatTok{0.95}\NormalTok{, demo\_text, }\AttributeTok{adj =} \FunctionTok{c}\NormalTok{(}\DecValTok{0}\NormalTok{, }\DecValTok{1}\NormalTok{), }\AttributeTok{family =} \StringTok{"mono"}\NormalTok{, }\AttributeTok{cex =} \FloatTok{0.8}\NormalTok{)}
\FunctionTok{dev.off}\NormalTok{()}
\end{Highlighting}
\end{Shaded}

\begin{verbatim}
#> cairo_pdf 
#>         2
\end{verbatim}

\begin{Shaded}
\begin{Highlighting}[]
\NormalTok{extracted }\OtherTok{\textless{}{-}} \FunctionTok{pdf\_text}\NormalTok{(tmp\_pdf)}
\FunctionTok{cat}\NormalTok{(}\FunctionTok{substr}\NormalTok{(extracted[}\DecValTok{1}\NormalTok{], }\DecValTok{1}\NormalTok{, }\DecValTok{500}\NormalTok{))}
\end{Highlighting}
\end{Shaded}

\begin{verbatim}
#> Title: A Review of Bibliometric Methods
#> 
#> Abstract: This paper reviews current methods in bibliometrics.
#> We focus on citation analysis, co−authorship networks, and topic modeling.
#> 
#> 1. Introduction
#> Bibliometrics has grown rapidly since the work of Garfield (1955).
#> Modern tools enable large−scale analysis of scholarly communication.
#> 
#> 2. Methods
#> We surveyed 500 papers published between 2020 and 2023.
#> Data were obtained from OpenAlex.
#> 
#> 3. Conclusions
#> The field continues to evolve with open data and n
\end{verbatim}

\subsection{Basic text preprocessing}\label{basic-text-preprocessing}

\begin{Shaded}
\begin{Highlighting}[]
\NormalTok{corp }\OtherTok{\textless{}{-}} \FunctionTok{corpus}\NormalTok{(extracted)}
\NormalTok{toks }\OtherTok{\textless{}{-}} \FunctionTok{tokens}\NormalTok{(corp, }\AttributeTok{remove\_punct =} \ConstantTok{TRUE}\NormalTok{, }\AttributeTok{remove\_numbers =} \ConstantTok{TRUE}\NormalTok{) }\SpecialCharTok{|\textgreater{}}
  \FunctionTok{tokens\_tolower}\NormalTok{() }\SpecialCharTok{|\textgreater{}}
  \FunctionTok{tokens\_remove}\NormalTok{(}\FunctionTok{stopwords}\NormalTok{(}\StringTok{"en"}\NormalTok{))}

\NormalTok{dfmat }\OtherTok{\textless{}{-}} \FunctionTok{dfm}\NormalTok{(toks)}
\FunctionTok{topfeatures}\NormalTok{(dfmat, }\DecValTok{15}\NormalTok{)}
\end{Highlighting}
\end{Shaded}

\begin{verbatim}
#>       methods bibliometrics      analysis             −          data 
#>             4             2             2             2             2 
#>         title        review  bibliometric      abstract         paper 
#>             1             1             1             1             1 
#>       reviews       current         focus      citation            co 
#>             1             1             1             1             1
\end{verbatim}

\subsection{Word frequency visualization}\label{word-frequency-visualization}

\begin{Shaded}
\begin{Highlighting}[]
\NormalTok{top\_words }\OtherTok{\textless{}{-}} \FunctionTok{topfeatures}\NormalTok{(dfmat, }\DecValTok{15}\NormalTok{) }\SpecialCharTok{|\textgreater{}}
  \FunctionTok{enframe}\NormalTok{(}\AttributeTok{name =} \StringTok{"word"}\NormalTok{, }\AttributeTok{value =} \StringTok{"count"}\NormalTok{) }\SpecialCharTok{|\textgreater{}}
  \FunctionTok{mutate}\NormalTok{(}\AttributeTok{word =} \FunctionTok{fct\_reorder}\NormalTok{(word, count))}

\FunctionTok{ggplot}\NormalTok{(top\_words, }\FunctionTok{aes}\NormalTok{(}\AttributeTok{x =}\NormalTok{ count, }\AttributeTok{y =}\NormalTok{ word)) }\SpecialCharTok{+}
  \FunctionTok{geom\_col}\NormalTok{(}\AttributeTok{fill =} \FunctionTok{palette\_sci}\NormalTok{(}\DecValTok{1}\NormalTok{)) }\SpecialCharTok{+}
  \FunctionTok{labs}\NormalTok{(}\AttributeTok{x =} \StringTok{"Frequency"}\NormalTok{, }\AttributeTok{y =} \ConstantTok{NULL}\NormalTok{) }\SpecialCharTok{+}
  \FunctionTok{theme\_sci}\NormalTok{()}
\end{Highlighting}
\end{Shaded}

\begin{figure}

{\centering \includegraphics[width=0.9\linewidth,alt={Horizontal bar chart showing the most frequent words extracted from the PDF, after removing stopwords.}]{scientometrics-in-r_files/figure-latex/08-fig-fulltext-wordfreq-1} 

}

\caption{Top 15 words extracted from the demo PDF after preprocessing.}\label{fig:08-fig-fulltext-wordfreq}
\end{figure}

\subsection{OCR with tesseract}\label{ocr-with-tesseract}

\begin{Shaded}
\begin{Highlighting}[]
\CommentTok{\# Requires tesseract system library and a scanned PDF}
\CommentTok{\# library(tesseract)}
\CommentTok{\# ocr\_text \textless{}{-} ocr(scanned\_pdf\_path, engine = tesseract("eng"))}
\CommentTok{\# cat(substr(ocr\_text, 1, 500))}
\end{Highlighting}
\end{Shaded}

\section{Diagnostics and interpretation}\label{fulltext-diagnostics}

Full-text extraction quality varies dramatically:

\begin{itemize}
\tightlist
\item
  \textbf{Born-digital PDFs}: \texttt{pdftools} produces clean text, but column layouts may interleave columns.
\item
  \textbf{Scanned PDFs}: OCR accuracy depends on scan quality, resolution, and language. Expect 90--98\% character accuracy for clean English scans.
\item
  \textbf{Figures and equations}: These are lost in text extraction. Do not expect mathematical notation or chart data.
\item
  \textbf{Headers and footers}: Page numbers, journal headers, and running titles pollute the extracted text. Consider filtering by position.
\end{itemize}

\section{Limitations and responsible use}\label{fulltext-limitations}

\section{Limitations and responsible use}\label{limitations-and-responsible-use-7}

\begin{itemize}
\tightlist
\item
  \textbf{Access rights}: Not all PDFs are legally available for TDM. Respect publisher terms and licensing. Open-access articles under CC licenses are generally safe.
\item
  \textbf{Extraction errors}: Two-column layouts, ligatures, and non-standard fonts cause parsing artifacts. Always spot-check extracted text.
\item
  \textbf{Scalability}: Full-text analysis is orders of magnitude more expensive than metadata analysis. Consider whether the additional information justifies the cost.
\item
  \textbf{Bias}: Full text is only available for a subset of publications (typically open access). Analyses based on full text may not generalise to the broader literature \citep{hicks2015}.
\end{itemize}

\section{Common pitfalls}\label{fulltext-pitfalls}

\section{Common pitfalls}\label{common-pitfalls-7}

\begin{itemize}
\tightlist
\item
  \textbf{Assuming clean extraction.} Even \texttt{pdftools} produces artifacts. Always inspect a sample of extracted text before bulk analysis.
\item
  \textbf{Skipping preprocessing.} Raw extracted text contains headers, page numbers, and hyphenated line breaks that must be cleaned.
\item
  \textbf{Mixing OCR and digital extraction.} Some PDFs contain both digital text and scanned images. Detect which pages need OCR before processing.
\item
  \textbf{Ignoring encoding issues.} PDFs may encode characters in non-standard ways. Check for garbled text, especially in references sections.
\end{itemize}

\section{Exercises}\label{fulltext-exercises}

\begin{enumerate}
\def\labelenumi{\arabic{enumi}.}
\item
  \textbf{Extract and compare.} Download an open-access PDF. Extract text with \texttt{pdftools}. Compare the extracted abstract with the abstract in OpenAlex. How accurate is the extraction?
\item
  \textbf{Section detection.} Write a function that splits extracted PDF text into sections based on numbered headings (``1. Introduction'', ``2. Methods'', etc.).
\item
  \textbf{Table extraction.} If you have \texttt{tabulizer} installed, extract a table from a PDF and convert it to a tibble. What cleaning steps are needed?
\end{enumerate}

\section{Solutions}\label{08-solutions}

Solutions are provided in \ref{solutions}.

\section{Further reading}\label{fulltext-reading}

\begin{itemize}
\tightlist
\item
  \citet{aria2017} --- \texttt{bibliometrix} supports some full-text features alongside metadata analysis.
\item
  \citet{priem2022} --- OpenAlex provides abstracts for many works, often sufficient without full-text extraction.
\end{itemize}

\section{Session info}\label{fulltext-sessioninfo}

\begin{Shaded}
\begin{Highlighting}[]
\FunctionTok{sessionInfo}\NormalTok{()}
\end{Highlighting}
\end{Shaded}

\begin{verbatim}
#> R version 4.6.1 (2026-06-24)
#> Platform: x86_64-pc-linux-gnu
#> Running under: Ubuntu 24.04.5 LTS
#> 
#> Matrix products: default
#> BLAS:   /usr/lib/x86_64-linux-gnu/openblas-pthread/libblas.so.3 
#> LAPACK: /usr/lib/x86_64-linux-gnu/openblas-pthread/libopenblasp-r0.3.26.so;  LAPACK version 3.12.0
#> 
#> locale:
#>  [1] LC_CTYPE=C.UTF-8       LC_NUMERIC=C           LC_TIME=C.UTF-8       
#>  [4] LC_COLLATE=C.UTF-8     LC_MONETARY=C.UTF-8    LC_MESSAGES=C.UTF-8   
#>  [7] LC_PAPER=C.UTF-8       LC_NAME=C              LC_ADDRESS=C          
#> [10] LC_TELEPHONE=C         LC_MEASUREMENT=C.UTF-8 LC_IDENTIFICATION=C   
#> 
#> time zone: UTC
#> tzcode source: system (glibc)
#> 
#> attached base packages:
#> [1] stats     graphics  grDevices utils     datasets  methods   base     
#> 
#> other attached packages:
#>  [1] quanteda_4.5.1     pdftools_3.9.1     arrow_26.0.0       bibliometrix_5.5.0
#>  [5] RefManageR_1.4.0   bib2df_1.1.2.0     rcrossref_1.2.1    gt_1.4.0          
#>  [9] tidytext_0.4.3     glue_1.8.1         openalexR_3.1.2    lubridate_1.9.5   
#> [13] forcats_1.0.1      stringr_1.6.0      dplyr_1.2.1        purrr_1.2.2       
#> [17] readr_2.2.0        tidyr_1.3.2        tibble_3.3.1       ggplot2_4.0.3     
#> [21] tidyverse_2.0.0   
#> 
#> loaded via a namespace (and not attached):
#>  [1] gridExtra_2.3.1        httr2_1.3.0            readxl_1.6.0          
#>  [4] rlang_1.3.0            magrittr_2.0.5         otel_0.2.0            
#>  [7] compiler_4.6.1         systemfonts_1.3.2      vctrs_0.7.3           
#> [10] httpcode_0.3.0         pkgconfig_2.0.3        fastmap_1.2.0         
#> [13] labeling_0.4.3         ca_0.72                utf8_1.2.6            
#> [16] promises_1.5.0         rmarkdown_2.32         tzdb_0.5.0            
#> [19] bit_4.6.0              xfun_0.62              jsonlite_2.0.0        
#> [22] SnowballC_0.7.1        later_1.4.8            parallel_4.6.1        
#> [25] stopwords_2.3          R6_2.6.1               stringi_1.8.9         
#> [28] RColorBrewer_1.1-3     cellranger_1.1.0       assertthat_0.2.1      
#> [31] Rcpp_1.1.2             bookdown_0.48          knitr_1.52            
#> [34] triebeard_0.4.1        base64enc_0.1-6        httpuv_1.6.17         
#> [37] rentrez_1.2.4          Matrix_1.7-5           igraph_2.3.4          
#> [40] timechange_0.4.0       tidyselect_1.2.1       stringdist_0.9.17     
#> [43] rstudioapi_0.19.0      pubmedR_1.0.2          yaml_2.3.12           
#> [46] viridis_0.6.5          miniUI_0.1.2           humaniformat_0.6.0    
#> [49] curl_8.1.0             qpdf_1.4.1             lattice_0.22-9        
#> [52] plyr_1.8.9             shiny_1.14.0           withr_3.0.3           
#> [55] S7_0.2.2               askpass_1.2.1          evaluate_1.0.5        
#> [58] zip_3.0.2              xml2_1.6.0             shinycssloaders_1.1.0 
#> [61] pillar_1.11.1          janeaustenr_1.0.0      DT_0.34.0             
#> [64] plotly_4.12.1          generics_0.1.4         rprojroot_2.1.1       
#> [67] hms_1.1.4              scales_1.4.0           xtable_1.8-8          
#> [70] contentanalysis_1.1.1  tools_4.6.1            data.table_1.18.6.1   
#> [73] tokenizers_0.3.0       openxlsx_4.2.9         fs_2.1.0              
#> [76] visNetwork_2.1.4       XML_3.99-0.25          fastmatch_1.1-8       
#> [79] grid_4.6.1             bibtex_0.5.3           urltools_1.7.3.1      
#> [82] rscopus_0.9.0          dimensionsR_0.0.3      bibliometrixData_0.3.0
#> [85] cli_3.6.6              viridisLite_0.4.3      gtable_0.3.6          
#> [88] digest_0.6.39          ggrepel_0.9.8          crul_1.6.0            
#> [91] htmlwidgets_1.6.4      farver_2.1.2           htmltools_0.5.9       
#> [94] lifecycle_1.0.5        httr_1.4.9             here_1.0.2            
#> [97] mime_0.13              bit64_4.8.6
\end{verbatim}

\part{Core Bibliometric Analysis}\label{part-core-bibliometric-analysis}

\chapter{Descriptive Bibliometrics}\label{descriptive-bibliometrics}

\emph{R. Heller · Date: 2026-05-09 · Last changed: 2026-10-31}

\section{Learning objectives}\label{desc-objectives}

After completing this chapter, you will be able to:

\begin{itemize}
\tightlist
\item
  Import bibliographic data into bibliometrix using \texttt{convert2df}
\item
  Produce a corpus-level summary with \texttt{biblioAnalysis} and \texttt{summary}
\item
  Identify the most productive authors, sources, and countries in a dataset
\item
  Construct and interpret a three-field plot (Sankey diagram)
\item
  Explain when to use the interactive biblioshiny app versus a scripted workflow
\item
  Recognise the descriptive statistics that should precede any advanced analysis
\end{itemize}

\section{Setup}\label{desc-setup}

\begin{Shaded}
\begin{Highlighting}[]
\FunctionTok{library}\NormalTok{(tidyverse)}
\FunctionTok{library}\NormalTok{(openalexR)}
\FunctionTok{library}\NormalTok{(bibliometrix)}
\FunctionTok{library}\NormalTok{(glue)}

\FunctionTok{set.seed}\NormalTok{(}\DecValTok{20260509}\NormalTok{)}

\FunctionTok{source}\NormalTok{(here}\SpecialCharTok{::}\FunctionTok{here}\NormalTok{(}\StringTok{"R"}\NormalTok{, }\StringTok{"api\_helpers.R"}\NormalTok{))}
\FunctionTok{source}\NormalTok{(here}\SpecialCharTok{::}\FunctionTok{here}\NormalTok{(}\StringTok{"R"}\NormalTok{, }\StringTok{"utils.R"}\NormalTok{))}
\FunctionTok{source}\NormalTok{(here}\SpecialCharTok{::}\FunctionTok{here}\NormalTok{(}\StringTok{"R"}\NormalTok{, }\StringTok{"sci\_palette.R"}\NormalTok{))}
\end{Highlighting}
\end{Shaded}

\section{Conceptual background}\label{desc-background}

Before computing any advanced indicator, a responsible bibliometric study begins with \textbf{descriptive analysis}: understanding the size, growth, and composition of the corpus. How many documents per year? Which journals dominate? Which countries contribute? These questions set the stage for every subsequent method.

The \textbf{bibliometrix} package \citep{aria2017} provides an integrated R workflow for science mapping. Its companion web app, \textbf{biblioshiny}, wraps the same functions in a point-and-click interface --- useful for exploration but not for reproducible reporting. In a Quarto pipeline, we prefer scripted calls so that every number is traceable to code.

Bibliometrix accepts data from Scopus, Web of Science, Dimensions, PubMed, and OpenAlex. The \texttt{convert2df()} function normalises heterogeneous export formats into a common data frame. The \texttt{biblioAnalysis()} function then computes a battery of descriptive statistics in one call: annual production, author productivity, source rankings, country output, and keyword frequencies.

Three classical bibliometric laws underpin this descriptive stage. \textbf{Lotka's law} \citep{lotka1926} describes the inverse-square relationship between the number of authors and their publication counts. \textbf{Bradford's law} \citep{bradford1934} models how articles on a subject scatter across journals, identifying core, mid-range, and peripheral zones. \textbf{Price's law} \citep{price1963} states that half of all contributions come from the square root of the total number of contributors. These empirical regularities provide benchmarks against which any corpus can be compared.

Descriptive bibliometrics also serves a data-quality role. Anomalous spikes in annual output may signal duplicate records; a suspiciously dominant country may indicate a search query biased toward a specific language. Inspecting descriptive statistics is the first line of defence against analytical errors.

\section{Worked example}\label{desc-example}

\subsection{Acquiring the data}\label{acquiring-the-data}

We fetch a sample of works from the \emph{Scientometrics} journal via OpenAlex and convert the result into a bibliometrix-compatible data frame.

\begin{Shaded}
\begin{Highlighting}[]
\NormalTok{works }\OtherTok{\textless{}{-}} \FunctionTok{fetch\_openalex}\NormalTok{(}
  \AttributeTok{entity =} \StringTok{"works"}\NormalTok{,}
  \AttributeTok{primary\_location.source.id =} \StringTok{"S148561398"}\NormalTok{,}
  \AttributeTok{from\_publication\_date =} \StringTok{"2018{-}01{-}01"}\NormalTok{,}
  \AttributeTok{to\_publication\_date =} \StringTok{"2023{-}12{-}31"}\NormalTok{,}
  \AttributeTok{options =} \FunctionTok{list}\NormalTok{(}\AttributeTok{sample =} \DecValTok{500}\NormalTok{, }\AttributeTok{seed =} \DecValTok{42}\NormalTok{)}
\NormalTok{)}

\FunctionTok{nrow}\NormalTok{(works)}
\end{Highlighting}
\end{Shaded}

\begin{verbatim}
#> [1] 500
\end{verbatim}

\subsection{Converting to bibliometrix format}\label{converting-to-bibliometrix-format}

\begin{Shaded}
\begin{Highlighting}[]
\NormalTok{bib\_df }\OtherTok{\textless{}{-}} \FunctionTok{oa2bibliometrix}\NormalTok{(works)}

\FunctionTok{cat}\NormalTok{(}\FunctionTok{glue}\NormalTok{(}\StringTok{"Records: \{nrow(bib\_df)\}}\SpecialCharTok{\textbackslash{}n}\StringTok{"}\NormalTok{))}
\end{Highlighting}
\end{Shaded}

\begin{verbatim}
#> Records: 500
\end{verbatim}

\begin{Shaded}
\begin{Highlighting}[]
\FunctionTok{cat}\NormalTok{(}\FunctionTok{glue}\NormalTok{(}\StringTok{"Fields: \{ncol(bib\_df)\}}\SpecialCharTok{\textbackslash{}n}\StringTok{"}\NormalTok{))}
\end{Highlighting}
\end{Shaded}

\begin{verbatim}
#> Fields: 61
\end{verbatim}

\subsection{Running biblioAnalysis}\label{running-biblioanalysis}

\begin{Shaded}
\begin{Highlighting}[]
\NormalTok{results }\OtherTok{\textless{}{-}} \FunctionTok{tryCatch}\NormalTok{(}
  \FunctionTok{biblioAnalysis}\NormalTok{(bib\_df, }\AttributeTok{sep =} \StringTok{";"}\NormalTok{),}
  \AttributeTok{error =} \ControlFlowTok{function}\NormalTok{(e) \{}
    \FunctionTok{cat}\NormalTok{(}\StringTok{"Note: biblioAnalysis() encountered an error with this data format.}\SpecialCharTok{\textbackslash{}n}\StringTok{"}\NormalTok{)}
    \FunctionTok{cat}\NormalTok{(}\StringTok{"This is a known compatibility issue between bibliometrix and openalexR v3.}\SpecialCharTok{\textbackslash{}n}\StringTok{"}\NormalTok{)}
    \ConstantTok{NULL}
\NormalTok{  \}}
\NormalTok{)}
\end{Highlighting}
\end{Shaded}

\begin{verbatim}
#> Note: biblioAnalysis() encountered an error with this data format.
#> This is a known compatibility issue between bibliometrix and openalexR v3.
\end{verbatim}

\begin{Shaded}
\begin{Highlighting}[]
\ControlFlowTok{if}\NormalTok{ (}\SpecialCharTok{!}\FunctionTok{is.null}\NormalTok{(results)) \{}
\NormalTok{  summary\_out }\OtherTok{\textless{}{-}} \FunctionTok{summary}\NormalTok{(results, }\AttributeTok{k =} \DecValTok{10}\NormalTok{, }\AttributeTok{verbose =} \ConstantTok{FALSE}\NormalTok{)}
\NormalTok{\}}
\end{Highlighting}
\end{Shaded}

\subsection{Annual scientific production}\label{annual-scientific-production}

\begin{Shaded}
\begin{Highlighting}[]
\NormalTok{annual }\OtherTok{\textless{}{-}}\NormalTok{ works }\SpecialCharTok{|\textgreater{}}
  \FunctionTok{mutate}\NormalTok{(}\AttributeTok{year =} \FunctionTok{year}\NormalTok{(publication\_date)) }\SpecialCharTok{|\textgreater{}}
  \FunctionTok{count}\NormalTok{(year)}

\FunctionTok{ggplot}\NormalTok{(annual, }\FunctionTok{aes}\NormalTok{(}\AttributeTok{x =}\NormalTok{ year, }\AttributeTok{y =}\NormalTok{ n)) }\SpecialCharTok{+}
  \FunctionTok{geom\_col}\NormalTok{(}\AttributeTok{fill =} \FunctionTok{palette\_sci}\NormalTok{(}\DecValTok{1}\NormalTok{), }\AttributeTok{width =} \FloatTok{0.7}\NormalTok{) }\SpecialCharTok{+}
  \FunctionTok{geom\_text}\NormalTok{(}\FunctionTok{aes}\NormalTok{(}\AttributeTok{label =}\NormalTok{ n), }\AttributeTok{vjust =} \SpecialCharTok{{-}}\FloatTok{0.5}\NormalTok{, }\AttributeTok{size =} \FloatTok{3.5}\NormalTok{) }\SpecialCharTok{+}
  \FunctionTok{labs}\NormalTok{(}\AttributeTok{x =} \StringTok{"Publication year"}\NormalTok{, }\AttributeTok{y =} \StringTok{"Number of articles"}\NormalTok{) }\SpecialCharTok{+}
  \FunctionTok{theme\_sci}\NormalTok{()}
\end{Highlighting}
\end{Shaded}

\begin{figure}

{\centering \includegraphics[width=0.9\linewidth,alt={Bar chart showing the number of publications per year from 2018 to 2023, illustrating the growth trend in the sampled Scientometrics articles.}]{scientometrics-in-r_files/figure-latex/09-fig-desc-annual-1} 

}

\caption{Annual scientific production in the sampled corpus (2018-2023).}\label{fig:09-fig-desc-annual}
\end{figure}

\subsection{Most productive authors}\label{most-productive-authors}

\begin{Shaded}
\begin{Highlighting}[]
\NormalTok{author\_prod }\OtherTok{\textless{}{-}}\NormalTok{ works }\SpecialCharTok{|\textgreater{}}
  \FunctionTok{select}\NormalTok{(id, authorships) }\SpecialCharTok{|\textgreater{}}
  \FunctionTok{unnest}\NormalTok{(authorships, }\AttributeTok{names\_sep =} \StringTok{"\_"}\NormalTok{) }\SpecialCharTok{|\textgreater{}}
  \FunctionTok{count}\NormalTok{(authorships\_display\_name, }\AttributeTok{sort =} \ConstantTok{TRUE}\NormalTok{) }\SpecialCharTok{|\textgreater{}}
  \FunctionTok{head}\NormalTok{(}\DecValTok{15}\NormalTok{)}

\FunctionTok{ggplot}\NormalTok{(author\_prod, }\FunctionTok{aes}\NormalTok{(}\AttributeTok{x =}\NormalTok{ n, }\AttributeTok{y =} \FunctionTok{reorder}\NormalTok{(authorships\_display\_name, n))) }\SpecialCharTok{+}
  \FunctionTok{geom\_col}\NormalTok{(}\AttributeTok{fill =} \FunctionTok{palette\_sci}\NormalTok{(}\DecValTok{1}\NormalTok{)) }\SpecialCharTok{+}
  \FunctionTok{labs}\NormalTok{(}\AttributeTok{x =} \StringTok{"Articles"}\NormalTok{, }\AttributeTok{y =} \ConstantTok{NULL}\NormalTok{) }\SpecialCharTok{+}
  \FunctionTok{theme\_sci}\NormalTok{()}
\end{Highlighting}
\end{Shaded}

\begin{figure}

{\centering \includegraphics[width=0.9\linewidth,alt={Horizontal bar chart listing the 15 authors with the most publications in the sample, ordered by number of articles.}]{scientometrics-in-r_files/figure-latex/09-fig-desc-top-authors-1} 

}

\caption{Top 15 most productive authors in the sample.}\label{fig:09-fig-desc-top-authors}
\end{figure}

\subsection{Most productive sources}\label{most-productive-sources}

Because all sampled works come from a single journal in this example, we instead show the most frequent cited sources from the references.

\begin{Shaded}
\begin{Highlighting}[]
\NormalTok{top\_sources }\OtherTok{\textless{}{-}}\NormalTok{ works }\SpecialCharTok{|\textgreater{}}
  \FunctionTok{select}\NormalTok{(id, referenced\_works) }\SpecialCharTok{|\textgreater{}}
  \FunctionTok{unnest}\NormalTok{(referenced\_works) }\SpecialCharTok{|\textgreater{}}
  \FunctionTok{count}\NormalTok{(referenced\_works, }\AttributeTok{sort =} \ConstantTok{TRUE}\NormalTok{) }\SpecialCharTok{|\textgreater{}}
  \FunctionTok{head}\NormalTok{(}\DecValTok{10}\NormalTok{)}

\FunctionTok{cat}\NormalTok{(}\FunctionTok{glue}\NormalTok{(}\StringTok{"Top 10 most{-}cited reference work IDs:}\SpecialCharTok{\textbackslash{}n}\StringTok{"}\NormalTok{))}
\end{Highlighting}
\end{Shaded}

\begin{verbatim}
#> Top 10 most-cited reference work IDs:
\end{verbatim}

\begin{Shaded}
\begin{Highlighting}[]
\FunctionTok{print}\NormalTok{(top\_sources)}
\end{Highlighting}
\end{Shaded}

\begin{verbatim}
#> # A tibble: 10 x 2
#>    referenced_works                     n
#>    <chr>                            <int>
#>  1 https://openalex.org/W2128438887    69
#>  2 https://openalex.org/W4285719527    59
#>  3 https://openalex.org/W4254947497    33
#>  4 https://openalex.org/W4292887282    30
#>  5 https://openalex.org/W2068452509    26
#>  6 https://openalex.org/W767067438     25
#>  7 https://openalex.org/W2150220236    24
#>  8 https://openalex.org/W1880262756    22
#>  9 https://openalex.org/W2005706022    19
#> 10 https://openalex.org/W2277805675    18
\end{verbatim}

\subsection{Three-field plot}\label{three-field-plot}

The three-field plot (Sankey diagram) links authors, keywords, and sources. It is a signature visualisation in bibliometrix.

\begin{Shaded}
\begin{Highlighting}[]
\FunctionTok{tryCatch}\NormalTok{(}
  \FunctionTok{threeFieldsPlot}\NormalTok{(bib\_df, }\AttributeTok{fields =} \FunctionTok{c}\NormalTok{(}\StringTok{"AU"}\NormalTok{, }\StringTok{"DE"}\NormalTok{, }\StringTok{"SO\_SO"}\NormalTok{), }\AttributeTok{n =} \FunctionTok{c}\NormalTok{(}\DecValTok{8}\NormalTok{, }\DecValTok{8}\NormalTok{, }\DecValTok{8}\NormalTok{)),}
  \AttributeTok{error =} \ControlFlowTok{function}\NormalTok{(e) }\FunctionTok{cat}\NormalTok{(}\StringTok{"Three{-}field plot skipped due to data format incompatibility.}\SpecialCharTok{\textbackslash{}n}\StringTok{"}\NormalTok{)}
\NormalTok{)}
\end{Highlighting}
\end{Shaded}

\begin{verbatim}
#> Three-field plot skipped due to data format incompatibility.
\end{verbatim}

\section{Diagnostics and interpretation}\label{desc-diagnostics}

When reviewing descriptive output, check the following:

\begin{itemize}
\tightlist
\item
  \textbf{Completeness}: Are there gaps in yearly coverage that indicate missing data rather than genuine drops?
\item
  \textbf{Author name variants}: The same person appearing under different name forms inflates author counts. OpenAlex's entity resolution helps, but manual inspection remains essential.
\item
  \textbf{Lotka's law fit}: If the author productivity distribution deviates markedly from an inverse-square law, the corpus may be too small or too specialised.
\item
  \textbf{Bradford zones}: In a multi-journal corpus, check whether a small core of journals accounts for roughly one-third of all articles, as Bradford's law predicts.
\end{itemize}

\section{Limitations and responsible use}\label{desc-limitations}

\section{Limitations and responsible use}\label{limitations-and-responsible-use-8}

\begin{itemize}
\tightlist
\item
  \textbf{Descriptive statistics are not evaluative.} The most productive author is not necessarily the most impactful. Publication counts say nothing about quality \citep{hicks2015}.
\item
  \textbf{Sample bias.} A random sample from OpenAlex may not represent the full population. Always report how the sample was drawn and its size relative to the universe.
\item
  \textbf{Language and database coverage.} OpenAlex has broader coverage than WoS or Scopus, but still underrepresents non-English-language research and certain disciplines \citep{visser2021}.
\item
  \textbf{Name disambiguation.} Despite algorithmic clustering, author and institutional names can still be conflated or split. Treat productivity rankings as approximate.
\end{itemize}

\section{Common pitfalls}\label{desc-pitfalls}

\section{Common pitfalls}\label{common-pitfalls-8}

\begin{itemize}
\tightlist
\item
  \textbf{Skipping descriptive analysis.} Jumping straight to network or topic analysis without understanding the corpus structure often leads to misinterpretation.
\item
  \textbf{Treating biblioshiny output as final.} The interactive app is excellent for exploration, but its figures are not reproducible. Always script the final analysis.
\item
  \textbf{Ignoring duplicate records.} Bibliographic exports often contain duplicates. Use \texttt{dedupe\_by\_doi()} or similar checks before computing any statistics.
\item
  \textbf{Comparing raw counts across time windows of different length.} Counting publications in a 6-year window is not comparable to a 3-year window without normalisation.
\end{itemize}

\section{Exercises}\label{desc-exercises}

\begin{enumerate}
\def\labelenumi{\arabic{enumi}.}
\item
  \textbf{Lotka's law.} Using the author productivity data from this chapter, fit Lotka's inverse-square law with \texttt{bibliometrix::lotka()}. Does the observed distribution follow the theoretical prediction? \emph{(Hint: plot observed vs.~expected proportions.)}
\item
  \textbf{Bradford's law.} Fetch works from a broad topic (e.g., ``machine learning'') across multiple journals. Apply \texttt{bibliometrix::bradford()} to identify the core, mid-range, and peripheral journal zones. How many journals form the core zone?
\item
  \textbf{Three-field exploration.} Modify the three-field plot to use author countries instead of author names as the left field (use \texttt{"AU\_CO"} if available). What new patterns emerge?
\item
  \textbf{Annual growth rate.} Compute the compound annual growth rate (CAGR) of publications in the sample. Is the field accelerating?
\end{enumerate}

\section{Solutions}\label{09-solutions}

Solutions are provided in \ref{solutions}.

\section{Further reading}\label{desc-reading}

\begin{itemize}
\tightlist
\item
  \citet{aria2017} --- The bibliometrix package paper; describes the full workflow and all available functions.
\item
  \citet{lotka1926} --- The original frequency distribution of scientific productivity.
\item
  \citet{bradford1934} --- Bradford's law of scattering: how articles distribute across journals.
\item
  \citet{price1963} --- \emph{Little Science, Big Science}; foundational work on the growth of science.
\item
  \citet{hicks2015} --- The Leiden Manifesto; a reminder that descriptive counts are not evaluations.
\item
  \citet{visser2021} --- Comparison of bibliographic databases; important context for coverage limitations.
\end{itemize}

\section{Session info}\label{desc-sessioninfo}

\begin{Shaded}
\begin{Highlighting}[]
\FunctionTok{sessionInfo}\NormalTok{()}
\end{Highlighting}
\end{Shaded}

\begin{verbatim}
#> R version 4.6.1 (2026-06-24)
#> Platform: x86_64-pc-linux-gnu
#> Running under: Ubuntu 24.04.5 LTS
#> 
#> Matrix products: default
#> BLAS:   /usr/lib/x86_64-linux-gnu/openblas-pthread/libblas.so.3 
#> LAPACK: /usr/lib/x86_64-linux-gnu/openblas-pthread/libopenblasp-r0.3.26.so;  LAPACK version 3.12.0
#> 
#> locale:
#>  [1] LC_CTYPE=C.UTF-8       LC_NUMERIC=C           LC_TIME=C.UTF-8       
#>  [4] LC_COLLATE=C.UTF-8     LC_MONETARY=C.UTF-8    LC_MESSAGES=C.UTF-8   
#>  [7] LC_PAPER=C.UTF-8       LC_NAME=C              LC_ADDRESS=C          
#> [10] LC_TELEPHONE=C         LC_MEASUREMENT=C.UTF-8 LC_IDENTIFICATION=C   
#> 
#> time zone: UTC
#> tzcode source: system (glibc)
#> 
#> attached base packages:
#> [1] stats     graphics  grDevices utils     datasets  methods   base     
#> 
#> other attached packages:
#>  [1] quanteda_4.5.1     pdftools_3.9.1     arrow_26.0.0       bibliometrix_5.5.0
#>  [5] RefManageR_1.4.0   bib2df_1.1.2.0     rcrossref_1.2.1    gt_1.4.0          
#>  [9] tidytext_0.4.3     glue_1.8.1         openalexR_3.1.2    lubridate_1.9.5   
#> [13] forcats_1.0.1      stringr_1.6.0      dplyr_1.2.1        purrr_1.2.2       
#> [17] readr_2.2.0        tidyr_1.3.2        tibble_3.3.1       ggplot2_4.0.3     
#> [21] tidyverse_2.0.0   
#> 
#> loaded via a namespace (and not attached):
#>  [1] gridExtra_2.3.1        httr2_1.3.0            readxl_1.6.0          
#>  [4] rlang_1.3.0            magrittr_2.0.5         otel_0.2.0            
#>  [7] compiler_4.6.1         systemfonts_1.3.2      vctrs_0.7.3           
#> [10] httpcode_0.3.0         pkgconfig_2.0.3        fastmap_1.2.0         
#> [13] labeling_0.4.3         ca_0.72                utf8_1.2.6            
#> [16] promises_1.5.0         rmarkdown_2.32         tzdb_0.5.0            
#> [19] bit_4.6.0              xfun_0.62              jsonlite_2.0.0        
#> [22] SnowballC_0.7.1        later_1.4.8            parallel_4.6.1        
#> [25] stopwords_2.3          R6_2.6.1               stringi_1.8.9         
#> [28] RColorBrewer_1.1-3     cellranger_1.1.0       assertthat_0.2.1      
#> [31] Rcpp_1.1.2             bookdown_0.48          knitr_1.52            
#> [34] triebeard_0.4.1        base64enc_0.1-6        httpuv_1.6.17         
#> [37] rentrez_1.2.4          Matrix_1.7-5           igraph_2.3.4          
#> [40] timechange_0.4.0       tidyselect_1.2.1       stringdist_0.9.17     
#> [43] rstudioapi_0.19.0      pubmedR_1.0.2          yaml_2.3.12           
#> [46] viridis_0.6.5          miniUI_0.1.2           humaniformat_0.6.0    
#> [49] curl_8.1.0             qpdf_1.4.1             lattice_0.22-9        
#> [52] plyr_1.8.9             shiny_1.14.0           withr_3.0.3           
#> [55] S7_0.2.2               askpass_1.2.1          evaluate_1.0.5        
#> [58] zip_3.0.2              xml2_1.6.0             shinycssloaders_1.1.0 
#> [61] pillar_1.11.1          janeaustenr_1.0.0      DT_0.34.0             
#> [64] plotly_4.12.1          generics_0.1.4         rprojroot_2.1.1       
#> [67] hms_1.1.4              scales_1.4.0           xtable_1.8-8          
#> [70] contentanalysis_1.1.1  tools_4.6.1            data.table_1.18.6.1   
#> [73] tokenizers_0.3.0       openxlsx_4.2.9         fs_2.1.0              
#> [76] visNetwork_2.1.4       XML_3.99-0.25          fastmatch_1.1-8       
#> [79] grid_4.6.1             bibtex_0.5.3           urltools_1.7.3.1      
#> [82] rscopus_0.9.0          dimensionsR_0.0.3      bibliometrixData_0.3.0
#> [85] cli_3.6.6              viridisLite_0.4.3      gtable_0.3.6          
#> [88] digest_0.6.39          ggrepel_0.9.8          crul_1.6.0            
#> [91] htmlwidgets_1.6.4      farver_2.1.2           htmltools_0.5.9       
#> [94] lifecycle_1.0.5        httr_1.4.9             here_1.0.2            
#> [97] mime_0.13              bit64_4.8.6
\end{verbatim}

\chapter{Productivity and Impact Indicators}\label{productivity-and-impact-indicators}

\emph{R. Heller · Date: 2026-05-09 · Last changed: 2026-10-31}

\section{Learning objectives}\label{indicators-objectives}

After completing this chapter, you will be able to:

\begin{itemize}
\tightlist
\item
  Compute sample h- and g-indices and distinguish a corpus publication span from an author career
\item
  Explain the difference between size-dependent and size-independent indicators
\item
  Explain the external reference sets needed for MNCS and PP(top 10\%), and diagnose misleading within-sample normalization
\item
  Articulate the limitations of each indicator and the risks of misuse, citing the Leiden Manifesto
\end{itemize}

\section{Setup}\label{indicators-setup}

\begin{Shaded}
\begin{Highlighting}[]
\FunctionTok{library}\NormalTok{(tidyverse)}
\FunctionTok{library}\NormalTok{(openalexR)}
\FunctionTok{library}\NormalTok{(glue)}
\FunctionTok{library}\NormalTok{(gt)}

\FunctionTok{set.seed}\NormalTok{(}\DecValTok{20260509}\NormalTok{)}

\FunctionTok{source}\NormalTok{(here}\SpecialCharTok{::}\FunctionTok{here}\NormalTok{(}\StringTok{"R"}\NormalTok{, }\StringTok{"api\_helpers.R"}\NormalTok{))}
\FunctionTok{source}\NormalTok{(here}\SpecialCharTok{::}\FunctionTok{here}\NormalTok{(}\StringTok{"R"}\NormalTok{, }\StringTok{"utils.R"}\NormalTok{))}
\FunctionTok{source}\NormalTok{(here}\SpecialCharTok{::}\FunctionTok{here}\NormalTok{(}\StringTok{"R"}\NormalTok{, }\StringTok{"sci\_palette.R"}\NormalTok{))}
\end{Highlighting}
\end{Shaded}

\section{Conceptual background}\label{indicators-background}

Quantitative indicators are central to research evaluation --- and among the most contested tools in science policy. The h-index, proposed by \citet{hirsch2005}, counts the number \emph{h} of an author's papers that have each received at least \emph{h} citations. Its simplicity made it wildly popular, but also obscures important nuances: it conflates productivity with impact, is insensitive to highly cited papers, and can only increase over a career \citep{waltman2016}.

The \textbf{g-index} \citep{egghe2006} addresses the insensitivity to top-cited papers by requiring that the top \emph{g} papers have together received at least \emph{g}² citations. The \textbf{m-quotient} divides the h-index by career length, offering a rough age correction.

These are all \emph{raw} indicators. Citation distributions differ across fields, so comparisons require a justified reference and \textbf{field normalization}. The \textbf{Mean Normalized Citation Score} (MNCS) divides each paper's citations by the world average for its field and publication year \citep{waltman2012mncs}. The \textbf{PP(top 10\%)} --- the proportion of an entity's papers in the top 10\% most cited in their field --- is a percentile-based indicator recommended by the Leiden Ranking.

The Leiden Manifesto \citep{hicks2015} and the San Francisco Declaration on Research Assessment \citep{dora2012} both caution against mechanical use of any single indicator. Indicators should support, not replace, qualitative judgement.

\section{Worked example}\label{indicators-example}

\subsection{Acquiring sample data}\label{acquiring-sample-data}

We fetch a sample of works and their citation counts from OpenAlex.

\begin{Shaded}
\begin{Highlighting}[]
\NormalTok{works }\OtherTok{\textless{}{-}} \FunctionTok{fetch\_openalex}\NormalTok{(}
  \AttributeTok{entity =} \StringTok{"works"}\NormalTok{,}
  \AttributeTok{primary\_location.source.id =} \StringTok{"S148561398"}\NormalTok{,}
  \AttributeTok{from\_publication\_date =} \StringTok{"2018{-}01{-}01"}\NormalTok{,}
  \AttributeTok{to\_publication\_date =} \StringTok{"2023{-}12{-}31"}\NormalTok{,}
  \AttributeTok{options =} \FunctionTok{list}\NormalTok{(}\AttributeTok{sample =} \DecValTok{500}\NormalTok{, }\AttributeTok{seed =} \DecValTok{42}\NormalTok{)}
\NormalTok{)}

\NormalTok{works\_slim }\OtherTok{\textless{}{-}}\NormalTok{ works }\SpecialCharTok{|\textgreater{}}
  \FunctionTok{select}\NormalTok{(id, display\_name, publication\_date, cited\_by\_count, doi) }\SpecialCharTok{|\textgreater{}}
  \FunctionTok{mutate}\NormalTok{(}\AttributeTok{year =} \FunctionTok{year}\NormalTok{(publication\_date))}

\FunctionTok{summary}\NormalTok{(works\_slim}\SpecialCharTok{$}\NormalTok{cited\_by\_count)}
\end{Highlighting}
\end{Shaded}

\begin{verbatim}
#>    Min. 1st Qu.  Median    Mean 3rd Qu.    Max. 
#>     0.0     8.0    17.0    26.8    31.0   810.0
\end{verbatim}

\subsection{Computing the h-index}\label{computing-the-h-index}

\begin{Shaded}
\begin{Highlighting}[]
\NormalTok{h }\OtherTok{\textless{}{-}} \FunctionTok{compute\_h\_index}\NormalTok{(works\_slim}\SpecialCharTok{$}\NormalTok{cited\_by\_count)}
\FunctionTok{cat}\NormalTok{(}\FunctionTok{glue}\NormalTok{(}\StringTok{"h{-}index for this corpus: \{h\}}\SpecialCharTok{\textbackslash{}n}\StringTok{"}\NormalTok{))}
\end{Highlighting}
\end{Shaded}

\begin{verbatim}
#> h-index for this corpus: 52
\end{verbatim}

The sample h-index means that at least 52 sampled papers each have at least 52 citations. It is not an author career h-index.

\subsection{Computing the g-index}\label{computing-the-g-index}

\begin{Shaded}
\begin{Highlighting}[]
\NormalTok{compute\_g\_index }\OtherTok{\textless{}{-}} \ControlFlowTok{function}\NormalTok{(citations) \{}
\NormalTok{  sorted }\OtherTok{\textless{}{-}} \FunctionTok{sort}\NormalTok{(citations, }\AttributeTok{decreasing =} \ConstantTok{TRUE}\NormalTok{)}
\NormalTok{  cumsum\_cites }\OtherTok{\textless{}{-}} \FunctionTok{cumsum}\NormalTok{(sorted)}
\NormalTok{  g\_candidates }\OtherTok{\textless{}{-}} \FunctionTok{which}\NormalTok{(cumsum\_cites }\SpecialCharTok{\textgreater{}=} \FunctionTok{seq\_along}\NormalTok{(sorted)}\SpecialCharTok{\^{}}\DecValTok{2}\NormalTok{)}
  \ControlFlowTok{if}\NormalTok{ (}\FunctionTok{length}\NormalTok{(g\_candidates) }\SpecialCharTok{==} \DecValTok{0}\NormalTok{) }\FunctionTok{return}\NormalTok{(}\DecValTok{0}\DataTypeTok{L}\NormalTok{)}
  \FunctionTok{as.integer}\NormalTok{(}\FunctionTok{max}\NormalTok{(g\_candidates))}
\NormalTok{\}}

\NormalTok{g }\OtherTok{\textless{}{-}} \FunctionTok{compute\_g\_index}\NormalTok{(works\_slim}\SpecialCharTok{$}\NormalTok{cited\_by\_count)}
\FunctionTok{cat}\NormalTok{(}\FunctionTok{glue}\NormalTok{(}\StringTok{"g{-}index for this corpus: \{g\}}\SpecialCharTok{\textbackslash{}n}\StringTok{"}\NormalTok{))}
\end{Highlighting}
\end{Shaded}

\begin{verbatim}
#> g-index for this corpus: 84
\end{verbatim}

\subsection{Publication span is not career length}\label{publication-span-is-not-career-length}

A proper m-quotient needs an identified author and a justified career-start date.
This corpus contains many authors, so its publication span cannot supply that
denominator. The following h/span ratio is only a descriptive calculation.

\begin{Shaded}
\begin{Highlighting}[]
\NormalTok{career\_years }\OtherTok{\textless{}{-}} \FunctionTok{as.numeric}\NormalTok{(}
  \FunctionTok{difftime}\NormalTok{(}\FunctionTok{max}\NormalTok{(works\_slim}\SpecialCharTok{$}\NormalTok{publication\_date),}
           \FunctionTok{min}\NormalTok{(works\_slim}\SpecialCharTok{$}\NormalTok{publication\_date),}
           \AttributeTok{units =} \StringTok{"days"}\NormalTok{)}
\NormalTok{) }\SpecialCharTok{/} \FloatTok{365.25}

\NormalTok{h\_per\_span }\OtherTok{\textless{}{-}} \ControlFlowTok{if}\NormalTok{ (}\FunctionTok{is.finite}\NormalTok{(career\_years) }\SpecialCharTok{\&\&}\NormalTok{ career\_years }\SpecialCharTok{\textgreater{}} \DecValTok{0}\NormalTok{) h }\SpecialCharTok{/}\NormalTok{ career\_years }\ControlFlowTok{else} \ConstantTok{NA\_real\_}
\FunctionTok{cat}\NormalTok{(}\FunctionTok{glue}\NormalTok{(}\StringTok{"Sample h/publication{-}span ratio: \{round(h\_per\_span, 3)\} (span \{round(career\_years, 1)\} years)}\SpecialCharTok{\textbackslash{}n}\StringTok{"}\NormalTok{))}
\end{Highlighting}
\end{Shaded}

\begin{verbatim}
#> Sample h/publication-span ratio: 8.688 (span 6 years)
\end{verbatim}

\subsection{A normalization diagnostic: using the sample as its own reference}\label{a-normalization-diagnostic-using-the-sample-as-its-own-reference}

MNCS requires a justified field/year reference with compatible document types
and citation windows. The live corpus here supplies no such reference. To show
why that matters, compute a mean within each sample year and divide the same
records by it. Every usable year must average to 1 by construction. These are
observed sample means, not simulated field means or world averages.

\begin{Shaded}
\begin{Highlighting}[]
\NormalTok{sample\_means\_by\_year }\OtherTok{\textless{}{-}}\NormalTok{ works\_slim }\SpecialCharTok{|\textgreater{}}
  \FunctionTok{group\_by}\NormalTok{(year) }\SpecialCharTok{|\textgreater{}}
  \FunctionTok{summarise}\NormalTok{(}\AttributeTok{sample\_mean =} \FunctionTok{mean}\NormalTok{(cited\_by\_count, }\AttributeTok{na.rm =} \ConstantTok{TRUE}\NormalTok{), }\AttributeTok{.groups =} \StringTok{"drop"}\NormalTok{)}

\NormalTok{works\_normalized }\OtherTok{\textless{}{-}}\NormalTok{ works\_slim }\SpecialCharTok{|\textgreater{}}
  \FunctionTok{left\_join}\NormalTok{(sample\_means\_by\_year, }\AttributeTok{by =} \StringTok{"year"}\NormalTok{) }\SpecialCharTok{|\textgreater{}}
  \FunctionTok{mutate}\NormalTok{(}\AttributeTok{ncs =} \FunctionTok{field\_normalize}\NormalTok{(cited\_by\_count, sample\_mean))}

\NormalTok{sample\_ratio }\OtherTok{\textless{}{-}} \FunctionTok{mean}\NormalTok{(works\_normalized}\SpecialCharTok{$}\NormalTok{ncs, }\AttributeTok{na.rm =} \ConstantTok{TRUE}\NormalTok{)}
\FunctionTok{cat}\NormalTok{(}\FunctionTok{glue}\NormalTok{(}\StringTok{"Mean within{-}sample/year ratio: \{round(sample\_ratio, 3)\}}\SpecialCharTok{\textbackslash{}n}\StringTok{"}\NormalTok{))}
\end{Highlighting}
\end{Shaded}

\begin{verbatim}
#> Mean within-sample/year ratio: 1
\end{verbatim}

A value of 1 here is an arithmetic identity, not evidence of world-average
impact. Zero-mean years have undefined ratios and are excluded. An actual MNCS
interpretation requires an independently justified field/year reference
\citep{waltman2012mncs}.

\subsection{Within-sample upper-percentile membership}\label{within-sample-upper-percentile-membership}

\begin{Shaded}
\begin{Highlighting}[]
\NormalTok{works\_with\_pct }\OtherTok{\textless{}{-}}\NormalTok{ works\_normalized }\SpecialCharTok{|\textgreater{}}
  \FunctionTok{group\_by}\NormalTok{(year) }\SpecialCharTok{|\textgreater{}}
  \FunctionTok{mutate}\NormalTok{(}
    \AttributeTok{pct\_rank =} \FunctionTok{percent\_rank}\NormalTok{(cited\_by\_count),}
    \AttributeTok{is\_top10 =}\NormalTok{ pct\_rank }\SpecialCharTok{\textgreater{}=} \FloatTok{0.90}
\NormalTok{  ) }\SpecialCharTok{|\textgreater{}}
  \FunctionTok{ungroup}\NormalTok{()}

\NormalTok{pp\_top10 }\OtherTok{\textless{}{-}} \FunctionTok{mean}\NormalTok{(works\_with\_pct}\SpecialCharTok{$}\NormalTok{is\_top10, }\AttributeTok{na.rm =} \ConstantTok{TRUE}\NormalTok{)}
\FunctionTok{cat}\NormalTok{(}\FunctionTok{glue}\NormalTok{(}\StringTok{"Share at or above the sample/year 90th percentile rank: \{scales::percent(pp\_top10, accuracy = 0.1)\}}\SpecialCharTok{\textbackslash{}n}\StringTok{"}\NormalTok{))}
\end{Highlighting}
\end{Shaded}

\begin{verbatim}
#> Share at or above the sample/year 90th percentile rank: 10.4%
\end{verbatim}

The ranks use only sampled works in each year. Ties and small groups can make
the selected share differ from 10\%; a singleton has no percentile rank. This
is not field-normalized PP(top 10\%). Report the ranking and tie rule explicitly.

\subsection{Visualization}\label{visualization-4}

\begin{Shaded}
\begin{Highlighting}[]
\NormalTok{works\_slim }\SpecialCharTok{|\textgreater{}}
  \FunctionTok{ggplot}\NormalTok{(}\FunctionTok{aes}\NormalTok{(}\AttributeTok{x =}\NormalTok{ cited\_by\_count)) }\SpecialCharTok{+}
  \FunctionTok{geom\_histogram}\NormalTok{(}\AttributeTok{binwidth =} \DecValTok{5}\NormalTok{, }\AttributeTok{fill =} \FunctionTok{palette\_sci}\NormalTok{(}\DecValTok{1}\NormalTok{), }\AttributeTok{color =} \StringTok{"white"}\NormalTok{) }\SpecialCharTok{+}
  \FunctionTok{geom\_vline}\NormalTok{(}\AttributeTok{xintercept =}\NormalTok{ h, }\AttributeTok{linetype =} \StringTok{"dashed"}\NormalTok{, }\AttributeTok{color =} \StringTok{"red"}\NormalTok{, }\AttributeTok{linewidth =} \FloatTok{0.8}\NormalTok{) }\SpecialCharTok{+}
  \FunctionTok{annotate}\NormalTok{(}\StringTok{"text"}\NormalTok{, }\AttributeTok{x =}\NormalTok{ h }\SpecialCharTok{+} \DecValTok{5}\NormalTok{, }\AttributeTok{y =} \ConstantTok{Inf}\NormalTok{, }\AttributeTok{label =} \FunctionTok{glue}\NormalTok{(}\StringTok{"h = \{h\}"}\NormalTok{),}
           \AttributeTok{vjust =} \DecValTok{2}\NormalTok{, }\AttributeTok{hjust =} \DecValTok{0}\NormalTok{, }\AttributeTok{color =} \StringTok{"red"}\NormalTok{, }\AttributeTok{size =} \DecValTok{4}\NormalTok{) }\SpecialCharTok{+}
  \FunctionTok{labs}\NormalTok{(}\AttributeTok{x =} \StringTok{"Citation count"}\NormalTok{, }\AttributeTok{y =} \StringTok{"Number of works"}\NormalTok{) }\SpecialCharTok{+}
  \FunctionTok{theme\_sci}\NormalTok{()}
\end{Highlighting}
\end{Shaded}

\begin{figure}

{\centering \includegraphics[width=0.9\linewidth,alt={Histogram of citation counts with a vertical dashed line at the h-index value, showing the long-tailed distribution typical of citation data.}]{scientometrics-in-r_files/figure-latex/10-fig-indicators-citation-dist-1} 

}

\caption{Citation distribution of sampled works, with the h-index threshold marked.}\label{fig:10-fig-indicators-citation-dist}
\end{figure}

\begin{Shaded}
\begin{Highlighting}[]
\NormalTok{works\_normalized }\SpecialCharTok{|\textgreater{}}
  \FunctionTok{group\_by}\NormalTok{(year) }\SpecialCharTok{|\textgreater{}}
  \FunctionTok{summarise}\NormalTok{(}\AttributeTok{sample\_ratio =} \FunctionTok{mean}\NormalTok{(ncs, }\AttributeTok{na.rm =} \ConstantTok{TRUE}\NormalTok{), }\AttributeTok{.groups =} \StringTok{"drop"}\NormalTok{) }\SpecialCharTok{|\textgreater{}}
  \FunctionTok{ggplot}\NormalTok{(}\FunctionTok{aes}\NormalTok{(}\AttributeTok{x =}\NormalTok{ year, }\AttributeTok{y =}\NormalTok{ sample\_ratio)) }\SpecialCharTok{+}
  \FunctionTok{geom\_line}\NormalTok{(}\AttributeTok{color =} \FunctionTok{palette\_sci}\NormalTok{(}\DecValTok{1}\NormalTok{), }\AttributeTok{linewidth =} \DecValTok{1}\NormalTok{) }\SpecialCharTok{+}
  \FunctionTok{geom\_point}\NormalTok{(}\AttributeTok{color =} \FunctionTok{palette\_sci}\NormalTok{(}\DecValTok{1}\NormalTok{), }\AttributeTok{size =} \DecValTok{3}\NormalTok{) }\SpecialCharTok{+}
  \FunctionTok{geom\_hline}\NormalTok{(}\AttributeTok{yintercept =} \DecValTok{1}\NormalTok{, }\AttributeTok{linetype =} \StringTok{"dashed"}\NormalTok{, }\AttributeTok{color =} \StringTok{"grey50"}\NormalTok{) }\SpecialCharTok{+}
  \FunctionTok{labs}\NormalTok{(}\AttributeTok{x =} \StringTok{"Publication year"}\NormalTok{, }\AttributeTok{y =} \StringTok{"Within{-}sample/year ratio"}\NormalTok{) }\SpecialCharTok{+}
  \FunctionTok{theme\_sci}\NormalTok{()}
\end{Highlighting}
\end{Shaded}

\begin{figure}

{\centering \includegraphics[width=0.9\linewidth,alt={Line chart at one across usable years, demonstrating the arithmetic identity created by normalizing each sample year to itself.}]{scientometrics-in-r_files/figure-latex/10-fig-indicators-mncs-by-year-1} 

}

\caption{Diagnostic: each usable sample/year normalizes to its own mean of one.}\label{fig:10-fig-indicators-mncs-by-year}
\end{figure}

\subsection{Summary table}\label{summary-table}

\begin{Shaded}
\begin{Highlighting}[]
\FunctionTok{tibble}\NormalTok{(}
  \AttributeTok{Indicator =} \FunctionTok{c}\NormalTok{(}\StringTok{"Sample h{-}index"}\NormalTok{, }\StringTok{"Sample g{-}index"}\NormalTok{, }\StringTok{"Sample h/span"}\NormalTok{, }\StringTok{"Mean sample/year ratio"}\NormalTok{, }\StringTok{"Sample upper{-}percentile share"}\NormalTok{),}
  \AttributeTok{Value =} \FunctionTok{c}\NormalTok{(}
    \FunctionTok{as.character}\NormalTok{(h),}
    \FunctionTok{as.character}\NormalTok{(g),}
    \FunctionTok{round}\NormalTok{(h\_per\_span, }\DecValTok{3}\NormalTok{) }\SpecialCharTok{|\textgreater{}} \FunctionTok{as.character}\NormalTok{(),}
    \FunctionTok{round}\NormalTok{(sample\_ratio, }\DecValTok{3}\NormalTok{) }\SpecialCharTok{|\textgreater{}} \FunctionTok{as.character}\NormalTok{(),}
\NormalTok{    scales}\SpecialCharTok{::}\FunctionTok{percent}\NormalTok{(pp\_top10, }\AttributeTok{accuracy =} \FloatTok{0.1}\NormalTok{)}
\NormalTok{  ),}
  \AttributeTok{Interpretation =} \FunctionTok{c}\NormalTok{(}
    \FunctionTok{glue}\NormalTok{(}\StringTok{"\{h\} papers with \textgreater{}= \{h\} citations each"}\NormalTok{),}
    \FunctionTok{glue}\NormalTok{(}\StringTok{"Top \{g\} papers have \textgreater{}= \{g\}\^{}2 = \{g\^{}2\} cumulative citations"}\NormalTok{),}
    \StringTok{"Not a career m{-}quotient"}\NormalTok{,}
    \StringTok{"One by construction; not MNCS"}\NormalTok{,}
    \StringTok{"Within{-}sample rank rule; not world PP(top 10\%)"}
\NormalTok{  )}
\NormalTok{) }\SpecialCharTok{|\textgreater{}}
  \FunctionTok{gt}\NormalTok{()}
\end{Highlighting}
\end{Shaded}

\begingroup
\fontsize{8.0pt}{9.0pt}\selectfont
\begin{longtable}{>{\raggedright\arraybackslash}p{\dimexpr 0.33\linewidth -2\tabcolsep-1.5\arrayrulewidth}>{\raggedleft\arraybackslash}p{\dimexpr 0.18\linewidth -2\tabcolsep-1.5\arrayrulewidth}>{\raggedright\arraybackslash}p{\dimexpr 0.49\linewidth -2\tabcolsep-1.5\arrayrulewidth}}
\toprule
\nolinkurl{Indicator} & Value & \nolinkurl{Interpretation} \\ 
\midrule\addlinespace[2.5pt]
Sample h-index & 52 & 52 papers with >= 52 citations each \\ 
Sample g-index & 84 & Top 84 papers have >= 84\textasciicircum{}2 = 7056 cumulative citations \\ 
Sample h/span & 8.688 & Not a career m-quotient \\ 
Mean sample/year ratio & 1 & One by construction; not MNCS \\ 
Sample upper-percentile share & 10.4\% & Within-sample rank rule; not world PP(top 10\%) \\ 
\bottomrule
\end{longtable}
\endgroup

\section{Diagnostics and interpretation}\label{indicators-diagnostics}

When interpreting bibliometric indicators, keep in mind:

\begin{itemize}
\tightlist
\item
  \textbf{Citation windows matter.} Recent papers have had less time to accumulate citations. The MNCS accounts for this by normalizing within publication year, but raw h-index and g-index do not.
\item
  \textbf{Field differences.} A ``good'' h-index in mathematics differs vastly from one in biomedicine. Only compare within fields, or use field-normalized indicators.
\item
  \textbf{Size dependence.} The h-index grows with the number of papers published. Comparing researchers at different career stages or with different publication rates requires the m-quotient or percentile indicators.
\item
  \textbf{Self-citations.} None of the indicators above exclude self-citations. In a full analysis, consider computing variants with and without self-citations.
\end{itemize}

\section{Limitations and responsible use}\label{indicators-limitations}

\section{Limitations and responsible use}\label{limitations-and-responsible-use-9}

No single number can capture the quality, significance, or societal impact of a researcher's work. The Leiden Manifesto \citep{hicks2015} offers ten principles for the responsible use of metrics:

\begin{enumerate}
\def\labelenumi{\arabic{enumi}.}
\tightlist
\item
  Quantitative evaluation should support, not replace, expert assessment.
\item
  Measure performance against the research missions of the institution, group, or researcher.
\item
  Protect excellence in locally relevant research (not just globally visible research).
\item
  Keep data collection and analytical processes open, transparent, and simple.
\item
  Allow those being evaluated to verify data and analysis.
\end{enumerate}

The h-index in particular can be gamed through self-citation, salami slicing, or guest authorship. It penalises researchers in small fields, those who publish fewer but higher-impact works, and those who take career breaks. The DORA declaration \citep{dora2012} explicitly warns against using journal-level metrics (like the Impact Factor) as surrogates for individual article quality --- but the same caution applies to any single indicator.

\textbf{Always present multiple indicators alongside qualitative context.} Never use a single number to make hiring, promotion, or funding decisions.

\section{Common pitfalls}\label{indicators-pitfalls}

\section{Common pitfalls}\label{common-pitfalls-9}

\begin{itemize}
\tightlist
\item
  \textbf{Comparing h-indices across fields.} A physicist with h = 40 and a mathematician with h = 15 may have equivalent standing in their respective fields. Use field-normalized indicators for cross-field comparisons.
\item
  \textbf{Ignoring the citation window.} Computing MNCS without accounting for publication year conflates old, well-cited papers with recent ones that haven't had time to accumulate citations.
\item
  \textbf{Using the wrong denominator for PP(top 10\%).} The ``top 10\%'' must be defined relative to the same field and year, not relative to your own corpus.
\item
  \textbf{Reporting precision that isn't there.} An MNCS of 1.237 vs.~1.241 is not a meaningful difference. Report with appropriate precision and always provide confidence intervals or bootstrapped uncertainty where possible.
\end{itemize}

\section{Exercises}\label{indicators-exercises}

\begin{enumerate}
\def\labelenumi{\arabic{enumi}.}
\item
  \textbf{Career h-index.} Fetch all works by a specific author using \texttt{oa\_fetch(entity\ =\ "works",\ author.id\ =\ "...")}. Compute their h-index, g-index, and m-quotient. How do the three indicators tell different stories about the same career?
\item
  \textbf{Self-citation sensitivity.} If you could identify self-citations (e.g., via author overlap in the reference list), how would removing them affect the h-index? Write pseudocode for this analysis.
\item
  \textbf{Empirical resampling.} Resample the within-sample/year ratios while holding their empirical reference fixed. Explain why the resulting interval does not measure uncertainty about world-normalized MNCS.
\item
  \textbf{Normalization and ranks.} Compare each year's mean sample ratio with its upper-percentile share. Explain why the ratio is fixed at one and why ties can change the percentile share; neither ranks world-normalized impact.
\end{enumerate}

\section{Solutions}\label{10-solutions}

Solutions are provided in \ref{solutions}.

\section{Further reading}\label{indicators-reading}

\begin{itemize}
\tightlist
\item
  \citet{hirsch2005} --- The original h-index paper.
\item
  \citet{egghe2006} --- The g-index, addressing the h-index's insensitivity to highly cited papers.
\item
  \citet{waltman2016} --- A comprehensive review of citation impact indicators and their properties.
\item
  \citet{waltman2012mncs} --- The theoretical basis for the MNCS (``new crown indicator'').
\item
  \citet{hicks2015} --- The Leiden Manifesto: ten principles for responsible metrics.
\item
  \citet{dora2012} --- DORA: the San Francisco Declaration on Research Assessment.
\end{itemize}

\section{Session info}\label{indicators-sessioninfo}

\begin{Shaded}
\begin{Highlighting}[]
\FunctionTok{sessionInfo}\NormalTok{()}
\end{Highlighting}
\end{Shaded}

\begin{verbatim}
#> R version 4.6.1 (2026-06-24)
#> Platform: x86_64-pc-linux-gnu
#> Running under: Ubuntu 24.04.5 LTS
#> 
#> Matrix products: default
#> BLAS:   /usr/lib/x86_64-linux-gnu/openblas-pthread/libblas.so.3 
#> LAPACK: /usr/lib/x86_64-linux-gnu/openblas-pthread/libopenblasp-r0.3.26.so;  LAPACK version 3.12.0
#> 
#> locale:
#>  [1] LC_CTYPE=C.UTF-8       LC_NUMERIC=C           LC_TIME=C.UTF-8       
#>  [4] LC_COLLATE=C.UTF-8     LC_MONETARY=C.UTF-8    LC_MESSAGES=C.UTF-8   
#>  [7] LC_PAPER=C.UTF-8       LC_NAME=C              LC_ADDRESS=C          
#> [10] LC_TELEPHONE=C         LC_MEASUREMENT=C.UTF-8 LC_IDENTIFICATION=C   
#> 
#> time zone: UTC
#> tzcode source: system (glibc)
#> 
#> attached base packages:
#> [1] stats     graphics  grDevices utils     datasets  methods   base     
#> 
#> other attached packages:
#>  [1] quanteda_4.5.1     pdftools_3.9.1     arrow_26.0.0       bibliometrix_5.5.0
#>  [5] RefManageR_1.4.0   bib2df_1.1.2.0     rcrossref_1.2.1    gt_1.4.0          
#>  [9] tidytext_0.4.3     glue_1.8.1         openalexR_3.1.2    lubridate_1.9.5   
#> [13] forcats_1.0.1      stringr_1.6.0      dplyr_1.2.1        purrr_1.2.2       
#> [17] readr_2.2.0        tidyr_1.3.2        tibble_3.3.1       ggplot2_4.0.3     
#> [21] tidyverse_2.0.0   
#> 
#> loaded via a namespace (and not attached):
#>  [1] gridExtra_2.3.1        httr2_1.3.0            readxl_1.6.0          
#>  [4] rlang_1.3.0            magrittr_2.0.5         otel_0.2.0            
#>  [7] compiler_4.6.1         systemfonts_1.3.2      vctrs_0.7.3           
#> [10] httpcode_0.3.0         pkgconfig_2.0.3        fastmap_1.2.0         
#> [13] labeling_0.4.3         ca_0.72                utf8_1.2.6            
#> [16] promises_1.5.0         rmarkdown_2.32         tzdb_0.5.0            
#> [19] bit_4.6.0              xfun_0.62              jsonlite_2.0.0        
#> [22] SnowballC_0.7.1        later_1.4.8            parallel_4.6.1        
#> [25] stopwords_2.3          R6_2.6.1               stringi_1.8.9         
#> [28] RColorBrewer_1.1-3     cellranger_1.1.0       assertthat_0.2.1      
#> [31] Rcpp_1.1.2             bookdown_0.48          knitr_1.52            
#> [34] triebeard_0.4.1        base64enc_0.1-6        httpuv_1.6.17         
#> [37] rentrez_1.2.4          Matrix_1.7-5           igraph_2.3.4          
#> [40] timechange_0.4.0       tidyselect_1.2.1       stringdist_0.9.17     
#> [43] rstudioapi_0.19.0      pubmedR_1.0.2          yaml_2.3.12           
#> [46] viridis_0.6.5          miniUI_0.1.2           humaniformat_0.6.0    
#> [49] curl_8.1.0             qpdf_1.4.1             lattice_0.22-9        
#> [52] plyr_1.8.9             shiny_1.14.0           withr_3.0.3           
#> [55] S7_0.2.2               askpass_1.2.1          evaluate_1.0.5        
#> [58] zip_3.0.2              xml2_1.6.0             shinycssloaders_1.1.0 
#> [61] pillar_1.11.1          janeaustenr_1.0.0      DT_0.34.0             
#> [64] plotly_4.12.1          generics_0.1.4         rprojroot_2.1.1       
#> [67] hms_1.1.4              scales_1.4.0           xtable_1.8-8          
#> [70] contentanalysis_1.1.1  tools_4.6.1            data.table_1.18.6.1   
#> [73] tokenizers_0.3.0       openxlsx_4.2.9         fs_2.1.0              
#> [76] visNetwork_2.1.4       XML_3.99-0.25          fastmatch_1.1-8       
#> [79] grid_4.6.1             bibtex_0.5.3           urltools_1.7.3.1      
#> [82] rscopus_0.9.0          dimensionsR_0.0.3      bibliometrixData_0.3.0
#> [85] cli_3.6.6              viridisLite_0.4.3      gtable_0.3.6          
#> [88] digest_0.6.39          ggrepel_0.9.8          crul_1.6.0            
#> [91] htmlwidgets_1.6.4      farver_2.1.2           htmltools_0.5.9       
#> [94] lifecycle_1.0.5        httr_1.4.9             here_1.0.2            
#> [97] mime_0.13              bit64_4.8.6
\end{verbatim}

\chapter{Journal Metrics}\label{journal-metrics}

\emph{R. Heller · Date: 2026-05-09 · Last changed: 2026-10-31}

\section{Learning objectives}\label{jmetrics-objectives}

After completing this chapter, you will be able to:

\begin{itemize}
\tightlist
\item
  Explain the citation windows and data requirements of JIF, CiteScore and SNIP
\item
  Explain why network-based journal metrics differ from simple citation ratios
\item
  Summarise current cumulative citations and recorded references in journal cohorts
\item
  Distinguish those descriptive summaries from official journal metrics
\item
  Evaluate journals using appropriate evidence rather than a single number
\end{itemize}

\section{Setup}\label{jmetrics-setup}

\begin{Shaded}
\begin{Highlighting}[]
\FunctionTok{library}\NormalTok{(tidyverse)}
\FunctionTok{library}\NormalTok{(openalexR)}
\FunctionTok{library}\NormalTok{(glue)}
\FunctionTok{library}\NormalTok{(gt)}

\FunctionTok{set.seed}\NormalTok{(}\DecValTok{20260509}\NormalTok{)}

\FunctionTok{source}\NormalTok{(here}\SpecialCharTok{::}\FunctionTok{here}\NormalTok{(}\StringTok{"R"}\NormalTok{, }\StringTok{"api\_helpers.R"}\NormalTok{))}
\FunctionTok{source}\NormalTok{(here}\SpecialCharTok{::}\FunctionTok{here}\NormalTok{(}\StringTok{"R"}\NormalTok{, }\StringTok{"utils.R"}\NormalTok{))}
\FunctionTok{source}\NormalTok{(here}\SpecialCharTok{::}\FunctionTok{here}\NormalTok{(}\StringTok{"R"}\NormalTok{, }\StringTok{"sci\_palette.R"}\NormalTok{))}
\end{Highlighting}
\end{Shaded}

\section{Conceptual background}\label{jmetrics-background}

Journal-level metrics summarise particular aspects of a journal's citation
record. The two-year \textbf{Journal Impact Factor} (JIF) for year \emph{t} counts citations
received in year \emph{t} to the journal's content from years \emph{t}-1 and \emph{t}-2, divided
by its citable items from those two publication years. The numerator can include
citations to document types excluded from the denominator. The citation year,
publication cohort and document-type rules are distinct parts of the definition;
see \href{https://clarivate.libguides.com/jcr/}{Clarivate's calculation explanation}.

\textbf{CiteScore} uses a four-year publication window and counts citations received
within that same four-year interval. Its eligible document types include articles,
reviews, conference papers, book chapters and data papers in both the numerator
and denominator; it does not simply include every document type. See
\href{https://www.elsevier.com/en-gb/researcher/author/tools-and-resources/measuring-a-journals-impact}{Elsevier's metric definitions}.

\textbf{SNIP} (Source-Normalized Impact per Paper) adjusts for citation practices using
information from the \emph{citing} side, with specified reference and weighting rules.
It is not a journal's citation average divided by the median reference count in
its own papers. The revised method is described by
\href{https://doi.org/10.1016/j.joi.2012.11.011}{Waltman et al.~(2013)}.
\textbf{SJR} (SCImago Journal Rank) and \textbf{Eigenfactor} use citation networks to weight
citations according to their source; a simple citation count does not reproduce
those network calculations \citep{waltman2016}.

The published indicators use specified database coverage, document eligibility,
time windows and methods. An alternative calculation needs its own explicit
definition and validation \citep{visser2021}. The OpenAlex example below computes
current cumulative citation and recorded-reference summaries only. It neither
implements nor estimates JIF, CiteScore, SNIP, SJR or Eigenfactor.

DORA \citep{dora2012} warns against using a journal metric as a measure of an individual
article's quality. Journal averages cannot establish the quality of a particular
paper, and citation distributions within journals can be highly uneven \citep{hicks2015}.

\section{Worked example}\label{jmetrics-example}

\subsection{Selecting journals}\label{selecting-journals}

We use five journals in information science and research evaluation. Source
identities were checked against OpenAlex on 30 October 2026. The JASIST record
below is the current \emph{Journal of the Association for Information Science and
Technology}, appropriate to the 2021--2022 publication cohort; historical titles
can have separate source records.

\begin{Shaded}
\begin{Highlighting}[]
\NormalTok{journals }\OtherTok{\textless{}{-}} \FunctionTok{tribble}\NormalTok{(}
  \SpecialCharTok{\textasciitilde{}}\NormalTok{short\_name,        }\SpecialCharTok{\textasciitilde{}}\NormalTok{source\_id,}
  \StringTok{"Scientometrics"}\NormalTok{,   }\StringTok{"S148561398"}\NormalTok{,}
  \StringTok{"J Informetrics"}\NormalTok{,   }\StringTok{"S205292342"}\NormalTok{,}
  \StringTok{"Quant Sci Stud"}\NormalTok{,   }\StringTok{"S4210195326"}\NormalTok{,}
  \StringTok{"JASIST"}\NormalTok{,           }\StringTok{"S4210197613"}\NormalTok{,}
  \StringTok{"Res Evaluation"}\NormalTok{,   }\StringTok{"S16793705"}
\NormalTok{)}
\end{Highlighting}
\end{Shaded}

\subsection{Fetching a publication cohort}\label{fetching-a-publication-cohort}

Fetch records classified as articles and published in 2021--2022. Their
\texttt{cited\_by\_count} values are \textbf{cumulative citations at retrieval}, not citations
received in a specified historical year. The query restricts publication dates;
it does not restrict the dates on which citations were received. Results describe
the retrieved records and require coverage checks before being treated as complete
journal output.

\begin{Shaded}
\begin{Highlighting}[]
\NormalTok{fetch\_journal\_works }\OtherTok{\textless{}{-}} \ControlFlowTok{function}\NormalTok{(source\_id, from\_year, to\_year) \{}
\NormalTok{  res }\OtherTok{\textless{}{-}} \FunctionTok{fetch\_openalex}\NormalTok{(}
    \AttributeTok{entity =} \StringTok{"works"}\NormalTok{,}
    \AttributeTok{primary\_location.source.id =}\NormalTok{ source\_id,}
    \AttributeTok{from\_publication\_date =} \FunctionTok{paste0}\NormalTok{(from\_year, }\StringTok{"{-}01{-}01"}\NormalTok{),}
    \AttributeTok{to\_publication\_date =} \FunctionTok{paste0}\NormalTok{(to\_year, }\StringTok{"{-}12{-}31"}\NormalTok{),}
    \AttributeTok{type =} \StringTok{"article"}
\NormalTok{  )}
  \ControlFlowTok{if}\NormalTok{ (}\SpecialCharTok{!}\FunctionTok{is.data.frame}\NormalTok{(res) }\SpecialCharTok{||} \SpecialCharTok{!}\FunctionTok{nrow}\NormalTok{(res)) \{}
    \FunctionTok{stop}\NormalTok{(}\FunctionTok{glue}\NormalTok{(}\StringTok{"No usable records for \{source\_id\}; check the query and provider response before comparing journals."}\NormalTok{))}
\NormalTok{  \}}
\NormalTok{  res}
\NormalTok{\}}

\NormalTok{journal\_works }\OtherTok{\textless{}{-}}\NormalTok{ journals }\SpecialCharTok{|\textgreater{}}
  \FunctionTok{mutate}\NormalTok{(}\AttributeTok{data =} \FunctionTok{map}\NormalTok{(source\_id, \textbackslash{}(sid) }\FunctionTok{fetch\_journal\_works}\NormalTok{(sid, }\DecValTok{2021}\NormalTok{, }\DecValTok{2022}\NormalTok{)))}
\end{Highlighting}
\end{Shaded}

Provider errors propagate instead of being converted to zero publications or
citations. An empty result stops the worked comparison; it is not evidence that
a journal has no output.

\subsection{Cumulative citation summaries}\label{cumulative-citation-summaries}

The citation mean uses only observed citation counts, with that denominator
reported. Missing counts are not zero. A total over observed values is not a
complete journal total when records or counts are missing.

\begin{Shaded}
\begin{Highlighting}[]
\NormalTok{summarise\_journal\_records }\OtherTok{\textless{}{-}} \ControlFlowTok{function}\NormalTok{(d) \{}
\NormalTok{  required }\OtherTok{\textless{}{-}} \FunctionTok{c}\NormalTok{(}\StringTok{"id"}\NormalTok{, }\StringTok{"cited\_by\_count"}\NormalTok{, }\StringTok{"referenced\_works\_count"}\NormalTok{)}
  \ControlFlowTok{if}\NormalTok{ (}\SpecialCharTok{!}\FunctionTok{is.data.frame}\NormalTok{(d) }\SpecialCharTok{||} \SpecialCharTok{!}\FunctionTok{nrow}\NormalTok{(d) }\SpecialCharTok{||} \SpecialCharTok{!}\FunctionTok{all}\NormalTok{(required }\SpecialCharTok{\%in\%} \FunctionTok{names}\NormalTok{(d))) \{}
    \FunctionTok{stop}\NormalTok{(}\StringTok{"Journal summaries need nonempty records with IDs, citation counts and reference counts."}\NormalTok{)}
\NormalTok{  \}}
  \ControlFlowTok{if}\NormalTok{ (}\FunctionTok{anyNA}\NormalTok{(d}\SpecialCharTok{$}\NormalTok{id) }\SpecialCharTok{||} \FunctionTok{any}\NormalTok{(}\SpecialCharTok{!}\FunctionTok{nzchar}\NormalTok{(d}\SpecialCharTok{$}\NormalTok{id)) }\SpecialCharTok{||} \FunctionTok{anyDuplicated}\NormalTok{(d}\SpecialCharTok{$}\NormalTok{id)) \{}
    \FunctionTok{stop}\NormalTok{(}\StringTok{"Journal work IDs must be present and unique."}\NormalTok{)}
\NormalTok{  \}}
  \ControlFlowTok{for}\NormalTok{ (column }\ControlFlowTok{in} \FunctionTok{c}\NormalTok{(}\StringTok{"cited\_by\_count"}\NormalTok{, }\StringTok{"referenced\_works\_count"}\NormalTok{)) \{}
\NormalTok{    values }\OtherTok{\textless{}{-}}\NormalTok{ d[[column]]}
    \ControlFlowTok{if}\NormalTok{ (}\SpecialCharTok{!}\FunctionTok{is.numeric}\NormalTok{(values) }\SpecialCharTok{||} \FunctionTok{any}\NormalTok{(}\SpecialCharTok{!}\FunctionTok{is.na}\NormalTok{(values) }\SpecialCharTok{\&}\NormalTok{ (}\SpecialCharTok{!}\FunctionTok{is.finite}\NormalTok{(values) }\SpecialCharTok{|}\NormalTok{ values }\SpecialCharTok{\textless{}} \DecValTok{0}\NormalTok{))) \{}
      \FunctionTok{stop}\NormalTok{(}\StringTok{"Citation and reference counts must be finite nonnegative numbers or missing."}\NormalTok{)}
\NormalTok{    \}}
\NormalTok{  \}}
\NormalTok{  citations }\OtherTok{\textless{}{-}}\NormalTok{ d}\SpecialCharTok{$}\NormalTok{cited\_by\_count[}\SpecialCharTok{!}\FunctionTok{is.na}\NormalTok{(d}\SpecialCharTok{$}\NormalTok{cited\_by\_count)]}
\NormalTok{  references }\OtherTok{\textless{}{-}}\NormalTok{ d}\SpecialCharTok{$}\NormalTok{referenced\_works\_count[}\SpecialCharTok{!}\FunctionTok{is.na}\NormalTok{(d}\SpecialCharTok{$}\NormalTok{referenced\_works\_count)]}
  \FunctionTok{tibble}\NormalTok{(}
    \AttributeTok{n\_articles =} \FunctionTok{nrow}\NormalTok{(d), }\AttributeTok{n\_citations\_observed =} \FunctionTok{length}\NormalTok{(citations),}
    \AttributeTok{observed\_citation\_sum =} \ControlFlowTok{if}\NormalTok{ (}\FunctionTok{length}\NormalTok{(citations)) }\FunctionTok{sum}\NormalTok{(citations) }\ControlFlowTok{else} \ConstantTok{NA\_real\_}\NormalTok{,}
    \AttributeTok{mean\_cumulative\_citations =} \ControlFlowTok{if}\NormalTok{ (}\FunctionTok{length}\NormalTok{(citations)) }\FunctionTok{mean}\NormalTok{(citations) }\ControlFlowTok{else} \ConstantTok{NA\_real\_}\NormalTok{,}
    \AttributeTok{median\_cumulative\_citations =} \ControlFlowTok{if}\NormalTok{ (}\FunctionTok{length}\NormalTok{(citations)) }\FunctionTok{median}\NormalTok{(citations) }\ControlFlowTok{else} \ConstantTok{NA\_real\_}\NormalTok{,}
    \AttributeTok{n\_references\_observed =} \FunctionTok{length}\NormalTok{(references),}
    \AttributeTok{median\_recorded\_references =} \ControlFlowTok{if}\NormalTok{ (}\FunctionTok{length}\NormalTok{(references)) }\FunctionTok{median}\NormalTok{(references) }\ControlFlowTok{else} \ConstantTok{NA\_real\_}
\NormalTok{  )}
\NormalTok{\}}

\NormalTok{journal\_summary }\OtherTok{\textless{}{-}}\NormalTok{ journal\_works }\SpecialCharTok{|\textgreater{}}
  \FunctionTok{mutate}\NormalTok{(}\AttributeTok{summary =} \FunctionTok{map}\NormalTok{(data, summarise\_journal\_records)) }\SpecialCharTok{|\textgreater{}}
  \FunctionTok{unnest}\NormalTok{(summary) }\SpecialCharTok{|\textgreater{}}
  \FunctionTok{select}\NormalTok{(}\SpecialCharTok{{-}}\NormalTok{data) }\SpecialCharTok{|\textgreater{}}
  \FunctionTok{arrange}\NormalTok{(short\_name)}

\NormalTok{journal\_summary }\SpecialCharTok{|\textgreater{}}
  \FunctionTok{select}\NormalTok{(short\_name, n\_articles, n\_citations\_observed, observed\_citation\_sum,}
\NormalTok{         mean\_cumulative\_citations, median\_cumulative\_citations) }\SpecialCharTok{|\textgreater{}}
  \FunctionTok{gt}\NormalTok{() }\SpecialCharTok{|\textgreater{}}
  \FunctionTok{fmt\_number}\NormalTok{(}\AttributeTok{columns =} \FunctionTok{c}\NormalTok{(mean\_cumulative\_citations, median\_cumulative\_citations), }\AttributeTok{decimals =} \DecValTok{2}\NormalTok{) }\SpecialCharTok{|\textgreater{}}
  \FunctionTok{cols\_label}\NormalTok{(}\AttributeTok{short\_name =} \StringTok{"Journal"}\NormalTok{, }\AttributeTok{n\_articles =} \StringTok{"Retrieved articles (2021–22)"}\NormalTok{,}
             \AttributeTok{n\_citations\_observed =} \StringTok{"Observed citation counts"}\NormalTok{,}
             \AttributeTok{observed\_citation\_sum =} \StringTok{"Sum of observed cumulative citations"}\NormalTok{,}
             \AttributeTok{mean\_cumulative\_citations =} \StringTok{"Mean cumulative citations"}\NormalTok{,}
             \AttributeTok{median\_cumulative\_citations =} \StringTok{"Median cumulative citations"}\NormalTok{)}
\end{Highlighting}
\end{Shaded}

\begingroup
\fontsize{8.0pt}{9.0pt}\selectfont
\begin{longtable}{>{\raggedright\arraybackslash}p{\dimexpr 0.15\linewidth -2\tabcolsep-1.5\arrayrulewidth}>{\raggedleft\arraybackslash}p{\dimexpr 0.12\linewidth -2\tabcolsep-1.5\arrayrulewidth}>{\raggedleft\arraybackslash}p{\dimexpr 0.17\linewidth -2\tabcolsep-1.5\arrayrulewidth}>{\raggedleft\arraybackslash}p{\dimexpr 0.18\linewidth -2\tabcolsep-1.5\arrayrulewidth}>{\raggedleft\arraybackslash}p{\dimexpr 0.19\linewidth -2\tabcolsep-1.5\arrayrulewidth}>{\raggedleft\arraybackslash}p{\dimexpr 0.19\linewidth -2\tabcolsep-1.5\arrayrulewidth}}
\toprule
Journal & Retrieved articles (2021–22) & Observed citation counts & Sum of observed cumulative citations & Mean cumulative citations & Median cumulative citations \\ 
\midrule\addlinespace[2.5pt]
J Informetrics & 185 & 185 & 4062 & 21.96 & 17.00 \\ 
JASIST & 211 & 211 & 5602 & 26.55 & 17.00 \\ 
Quant Sci Stud & 128 & 128 & 4627 & 36.15 & 14.00 \\ 
Res Evaluation & 84 & 84 & 1804 & 21.48 & 14.00 \\ 
\nolinkurl{Scientometrics} & 747 & 747 & 19699 & 26.37 & 15.00 \\ 
\bottomrule
\end{longtable}
\endgroup

\subsection{Recorded reference counts}\label{recorded-reference-counts}

\texttt{referenced\_works\_count} describes the references recorded for the fetched
journal articles. It is not the reference-list information for papers citing
those articles, and it need not capture every reference in the original texts.
We report it separately rather than dividing citations by it or calling it
citation potential.

\begin{Shaded}
\begin{Highlighting}[]
\NormalTok{journal\_summary }\SpecialCharTok{|\textgreater{}}
  \FunctionTok{select}\NormalTok{(short\_name, n\_articles, n\_references\_observed, median\_recorded\_references) }\SpecialCharTok{|\textgreater{}}
  \FunctionTok{gt}\NormalTok{() }\SpecialCharTok{|\textgreater{}}
  \FunctionTok{cols\_label}\NormalTok{(}\AttributeTok{short\_name =} \StringTok{"Journal"}\NormalTok{, }\AttributeTok{n\_articles =} \StringTok{"Retrieved articles"}\NormalTok{,}
             \AttributeTok{n\_references\_observed =} \StringTok{"Observed reference counts"}\NormalTok{,}
             \AttributeTok{median\_recorded\_references =} \StringTok{"Median recorded references"}\NormalTok{)}
\end{Highlighting}
\end{Shaded}

\begingroup
\fontsize{8.0pt}{9.0pt}\selectfont
\begin{longtable}{>{\raggedright\arraybackslash}p{\dimexpr 0.23\linewidth -2\tabcolsep-1.5\arrayrulewidth}>{\raggedleft\arraybackslash}p{\dimexpr 0.19\linewidth -2\tabcolsep-1.5\arrayrulewidth}>{\raggedleft\arraybackslash}p{\dimexpr 0.28\linewidth -2\tabcolsep-1.5\arrayrulewidth}>{\raggedleft\arraybackslash}p{\dimexpr 0.30\linewidth -2\tabcolsep-1.5\arrayrulewidth}}
\toprule
Journal & Retrieved articles & Observed reference counts & Median recorded references \\ 
\midrule\addlinespace[2.5pt]
J Informetrics & 185 & 185 & 59 \\ 
JASIST & 211 & 211 & 50 \\ 
Quant Sci Stud & 128 & 128 & 47 \\ 
Res Evaluation & 84 & 84 & 46 \\ 
\nolinkurl{Scientometrics} & 747 & 747 & 46 \\ 
\bottomrule
\end{longtable}
\endgroup

\subsection{Visualization}\label{visualization-5}

\begin{Shaded}
\begin{Highlighting}[]
\NormalTok{combined }\OtherTok{\textless{}{-}}\NormalTok{ journal\_summary }\SpecialCharTok{|\textgreater{}}
  \FunctionTok{select}\NormalTok{(short\_name, mean\_cumulative\_citations, median\_recorded\_references) }\SpecialCharTok{|\textgreater{}}
  \FunctionTok{pivot\_longer}\NormalTok{(}\SpecialCharTok{{-}}\NormalTok{short\_name, }\AttributeTok{names\_to =} \StringTok{"summary"}\NormalTok{, }\AttributeTok{values\_to =} \StringTok{"value"}\NormalTok{) }\SpecialCharTok{|\textgreater{}}
  \FunctionTok{mutate}\NormalTok{(}\AttributeTok{summary =} \FunctionTok{recode}\NormalTok{(summary,}
    \AttributeTok{mean\_cumulative\_citations =} \StringTok{"Mean cumulative citations"}\NormalTok{,}
    \AttributeTok{median\_recorded\_references =} \StringTok{"Median recorded references"}\NormalTok{))}

\FunctionTok{ggplot}\NormalTok{(combined, }\FunctionTok{aes}\NormalTok{(}\AttributeTok{x =}\NormalTok{ value, }\AttributeTok{y =}\NormalTok{ short\_name)) }\SpecialCharTok{+}
  \FunctionTok{geom\_point}\NormalTok{(}\AttributeTok{size =} \DecValTok{3}\NormalTok{, }\AttributeTok{colour =} \FunctionTok{palette\_sci}\NormalTok{(}\DecValTok{1}\NormalTok{)) }\SpecialCharTok{+}
  \FunctionTok{facet\_wrap}\NormalTok{(}\SpecialCharTok{\textasciitilde{}}\NormalTok{ summary, }\AttributeTok{scales =} \StringTok{"free\_x"}\NormalTok{) }\SpecialCharTok{+}
  \FunctionTok{labs}\NormalTok{(}\AttributeTok{x =} \StringTok{"Descriptive value (see panel)"}\NormalTok{, }\AttributeTok{y =} \ConstantTok{NULL}\NormalTok{) }\SpecialCharTok{+}
  \FunctionTok{theme\_sci}\NormalTok{()}
\end{Highlighting}
\end{Shaded}

\begin{figure}

{\centering \includegraphics[width=0.9\linewidth,alt={Two panels show journal mean cumulative citations at retrieval and median recorded references in those same journal articles; neither panel is an official journal metric.}]{scientometrics-in-r_files/figure-latex/11-fig-jmetrics-comparison-1} 

}

\caption{Current cumulative citations and recorded references for the retrieved 2021–2022 article cohorts, shown as separate descriptive quantities.}\label{fig:11-fig-jmetrics-comparison}
\end{figure}

\section{Diagnostics and interpretation}\label{jmetrics-diagnostics}

\begin{itemize}
\tightlist
\item
  \textbf{Coverage}: Report source IDs, publication filters, retrieval dates and the number of retrieved records. Reconcile these with the intended corpus before interpreting output counts.
\item
  \textbf{Missing data}: Compare the retrieved-record count with each observed denominator. An all-missing citation column produces missing summaries, not a zero-citation journal.
\item
  \textbf{Time}: Older papers have had longer to accumulate citations. A publication-date filter cannot create an annual citation window from a current cumulative count.
\item
  \textbf{Scope}: References in the journal's own articles describe those records; they do not supply the citing-side data needed for SNIP.
\item
  \textbf{Comparisons}: Differences can reflect field mix, document classification, coverage and citation exposure. Agreement between descriptive rankings does not validate an official metric or justify an evaluation decision.
\end{itemize}

\section{Limitations and responsible use}\label{jmetrics-limitations}

\section{Limitations and responsible use}\label{limitations-and-responsible-use-10}

\begin{itemize}
\tightlist
\item
  \textbf{Journal summaries are not article-quality measures.} Means can conceal skewed citation distributions and do not establish a particular paper's quality \citep{dora2012}.
\item
  \textbf{Targets can distort behaviour.} Citation stacking, coercive citation and denominator manipulation are risks when metrics become targets \citep{goodhart1984}.
\item
  \textbf{These are descriptive snapshots.} They are not JIF, CiteScore, SNIP or substitutes for those indicators. Reusing a familiar label does not supply the missing citation windows or methodology.
\item
  \textbf{Field dependence and metadata coverage remain.} The example performs no field normalisation and cannot establish a fair ranking across disciplines \citep{waltman2016}.
\end{itemize}

\section{Common pitfalls}\label{jmetrics-pitfalls}

\section{Common pitfalls}\label{common-pitfalls-10}

\begin{itemize}
\tightlist
\item
  \textbf{Relabeling cumulative citations as annual citations.} JIF needs citations received in one specified year to the preceding two publication years.
\item
  \textbf{Using the wrong side of a citation link.} The cited journal's own reference counts do not describe the reference lists of its citing works.
\item
  \textbf{Treating absence as zero.} Failed requests, empty responses and missing values require explicit handling before computing a denominator.
\item
  \textbf{Ranking journals by a single number.} Indicator definitions, distributions, coverage and disciplinary context matter alongside any numerical comparison.
\end{itemize}

\section{Exercises}\label{jmetrics-exercises}

\begin{enumerate}
\def\labelenumi{\arabic{enumi}.}
\item
  \textbf{Cohort sensitivity.} Fetch 2019--2022 articles for one of the verified source IDs and compare their cumulative citation summary with the 2021--2022 cohort. Explain why widening the publication cohort does not compute CiteScore.
\item
  \textbf{Descriptive association.} Compute Spearman correlation between the journal means of cumulative citations and medians of recorded references, using complete pairs. Explain why any association does not validate SNIP or establish field normalisation.
\item
  \textbf{Missing-data denominator.} For three invented records with citation counts 0, 20 and missing, calculate the observed citation sum and mean using the helper. Compare that mean with dividing the observed sum by all three records. Also check a cohort whose citation counts are all missing.
\item
  \textbf{Self-citation design.} Specify the citation edges, citing-journal identities, time window and denominator needed to measure a journal self-citation share. Explain why the fetched cumulative counts and reference counts are insufficient.
\end{enumerate}

\section{Solutions}\label{11-solutions}

Solutions are provided in \ref{solutions}.

\section{Further reading}\label{jmetrics-reading}

\begin{itemize}
\tightlist
\item
  \citet{garfield1955} --- The original proposal for citation indexes.
\item
  \citet{waltman2016} --- Review of citation impact indicators and their data requirements.
\item
  \href{https://clarivate.libguides.com/jcr/}{Clarivate's JIF calculation} and \href{https://www.elsevier.com/en-gb/researcher/author/tools-and-resources/measuring-a-journals-impact}{Elsevier's metric definitions} --- Provider definitions of citation windows and eligible documents.
\item
  \href{https://doi.org/10.1016/j.joi.2012.11.011}{Waltman et al.~(2013)} --- The revised SNIP methodology.
\item
  \citet{hicks2015}; \citet{dora2012} --- Responsible use of journal-level metrics.
\item
  \citet{visser2021} --- Database coverage differences.
\item
  \citet{priem2022} --- OpenAlex as an open bibliographic data source.
\end{itemize}

\section{Session info}\label{jmetrics-sessioninfo}

\begin{Shaded}
\begin{Highlighting}[]
\FunctionTok{sessionInfo}\NormalTok{()}
\end{Highlighting}
\end{Shaded}

\begin{verbatim}
#> R version 4.6.1 (2026-06-24)
#> Platform: x86_64-pc-linux-gnu
#> Running under: Ubuntu 24.04.5 LTS
#> 
#> Matrix products: default
#> BLAS:   /usr/lib/x86_64-linux-gnu/openblas-pthread/libblas.so.3 
#> LAPACK: /usr/lib/x86_64-linux-gnu/openblas-pthread/libopenblasp-r0.3.26.so;  LAPACK version 3.12.0
#> 
#> locale:
#>  [1] LC_CTYPE=C.UTF-8       LC_NUMERIC=C           LC_TIME=C.UTF-8       
#>  [4] LC_COLLATE=C.UTF-8     LC_MONETARY=C.UTF-8    LC_MESSAGES=C.UTF-8   
#>  [7] LC_PAPER=C.UTF-8       LC_NAME=C              LC_ADDRESS=C          
#> [10] LC_TELEPHONE=C         LC_MEASUREMENT=C.UTF-8 LC_IDENTIFICATION=C   
#> 
#> time zone: UTC
#> tzcode source: system (glibc)
#> 
#> attached base packages:
#> [1] stats     graphics  grDevices utils     datasets  methods   base     
#> 
#> other attached packages:
#>  [1] quanteda_4.5.1     pdftools_3.9.1     arrow_26.0.0       bibliometrix_5.5.0
#>  [5] RefManageR_1.4.0   bib2df_1.1.2.0     rcrossref_1.2.1    gt_1.4.0          
#>  [9] tidytext_0.4.3     glue_1.8.1         openalexR_3.1.2    lubridate_1.9.5   
#> [13] forcats_1.0.1      stringr_1.6.0      dplyr_1.2.1        purrr_1.2.2       
#> [17] readr_2.2.0        tidyr_1.3.2        tibble_3.3.1       ggplot2_4.0.3     
#> [21] tidyverse_2.0.0   
#> 
#> loaded via a namespace (and not attached):
#>  [1] gridExtra_2.3.1        httr2_1.3.0            readxl_1.6.0          
#>  [4] rlang_1.3.0            magrittr_2.0.5         otel_0.2.0            
#>  [7] compiler_4.6.1         systemfonts_1.3.2      vctrs_0.7.3           
#> [10] httpcode_0.3.0         pkgconfig_2.0.3        fastmap_1.2.0         
#> [13] labeling_0.4.3         ca_0.72                utf8_1.2.6            
#> [16] promises_1.5.0         rmarkdown_2.32         tzdb_0.5.0            
#> [19] bit_4.6.0              xfun_0.62              jsonlite_2.0.0        
#> [22] SnowballC_0.7.1        later_1.4.8            parallel_4.6.1        
#> [25] stopwords_2.3          R6_2.6.1               stringi_1.8.9         
#> [28] RColorBrewer_1.1-3     cellranger_1.1.0       assertthat_0.2.1      
#> [31] Rcpp_1.1.2             bookdown_0.48          knitr_1.52            
#> [34] triebeard_0.4.1        base64enc_0.1-6        httpuv_1.6.17         
#> [37] rentrez_1.2.4          Matrix_1.7-5           igraph_2.3.4          
#> [40] timechange_0.4.0       tidyselect_1.2.1       stringdist_0.9.17     
#> [43] rstudioapi_0.19.0      pubmedR_1.0.2          yaml_2.3.12           
#> [46] viridis_0.6.5          miniUI_0.1.2           humaniformat_0.6.0    
#> [49] curl_8.1.0             qpdf_1.4.1             lattice_0.22-9        
#> [52] plyr_1.8.9             shiny_1.14.0           withr_3.0.3           
#> [55] S7_0.2.2               askpass_1.2.1          evaluate_1.0.5        
#> [58] zip_3.0.2              xml2_1.6.0             shinycssloaders_1.1.0 
#> [61] pillar_1.11.1          janeaustenr_1.0.0      DT_0.34.0             
#> [64] plotly_4.12.1          generics_0.1.4         rprojroot_2.1.1       
#> [67] hms_1.1.4              scales_1.4.0           xtable_1.8-8          
#> [70] contentanalysis_1.1.1  tools_4.6.1            data.table_1.18.6.1   
#> [73] tokenizers_0.3.0       openxlsx_4.2.9         fs_2.1.0              
#> [76] visNetwork_2.1.4       XML_3.99-0.25          fastmatch_1.1-8       
#> [79] grid_4.6.1             bibtex_0.5.3           urltools_1.7.3.1      
#> [82] rscopus_0.9.0          dimensionsR_0.0.3      bibliometrixData_0.3.0
#> [85] cli_3.6.6              viridisLite_0.4.3      gtable_0.3.6          
#> [88] digest_0.6.39          ggrepel_0.9.8          crul_1.6.0            
#> [91] htmlwidgets_1.6.4      farver_2.1.2           htmltools_0.5.9       
#> [94] lifecycle_1.0.5        httr_1.4.9             here_1.0.2            
#> [97] mime_0.13              bit64_4.8.6
\end{verbatim}

\chapter{Author-Level Analysis and Disambiguation}\label{author-level-analysis-and-disambiguation}

\emph{R. Heller · Date: 2026-05-09 · Last changed: 2026-10-31}

\section{Learning objectives}\label{author-objectives}

After completing this chapter, you will be able to:

\begin{itemize}
\tightlist
\item
  Explain the author disambiguation problem and its impact on bibliometric analysis
\item
  Retrieve and interpret OpenAlex author entities, including name variants
\item
  Integrate ORCID identifiers to validate author identity
\item
  Compute author-level productivity, citation, and topical statistics from works data
\item
  Build a reproducible author profile report from OpenAlex
\end{itemize}

\section{Setup}\label{author-setup}

\begin{Shaded}
\begin{Highlighting}[]
\FunctionTok{library}\NormalTok{(tidyverse)}
\FunctionTok{library}\NormalTok{(openalexR)}
\FunctionTok{library}\NormalTok{(glue)}
\FunctionTok{library}\NormalTok{(gt)}

\FunctionTok{set.seed}\NormalTok{(}\DecValTok{20260509}\NormalTok{)}

\FunctionTok{source}\NormalTok{(here}\SpecialCharTok{::}\FunctionTok{here}\NormalTok{(}\StringTok{"R"}\NormalTok{, }\StringTok{"api\_helpers.R"}\NormalTok{))}
\FunctionTok{source}\NormalTok{(here}\SpecialCharTok{::}\FunctionTok{here}\NormalTok{(}\StringTok{"R"}\NormalTok{, }\StringTok{"utils.R"}\NormalTok{))}
\FunctionTok{source}\NormalTok{(here}\SpecialCharTok{::}\FunctionTok{here}\NormalTok{(}\StringTok{"R"}\NormalTok{, }\StringTok{"sci\_palette.R"}\NormalTok{))}
\end{Highlighting}
\end{Shaded}

\section{Conceptual background}\label{author-background}

Correctly attributing publications to individual researchers is a prerequisite for any author-level analysis, yet it remains one of the hardest problems in bibliometrics. \textbf{Author disambiguation} --- determining whether two name strings refer to the same person --- is complicated by name variants (initials vs.~full names, transliterations, name changes after marriage), homonymy (different people sharing a name), and inconsistent metadata across databases.

OpenAlex \citep{priem2022} addresses this through algorithmic author clustering that groups works likely belonging to the same person. Each cluster receives a unique author ID (e.g., \texttt{A5023888391}). The algorithm considers name similarity, co-author overlap, institutional affiliation, topic continuity, and citation patterns. While imperfect, this clustering is transparent and continuously improved.

\textbf{ORCID} (Open Researcher and Contributor ID) provides a complementary solution: a persistent digital identifier that researchers claim and curate themselves. When an ORCID is linked to an OpenAlex author entity, it provides strong confirmation of identity. However, ORCID adoption is uneven across disciplines and geographies, so it cannot serve as the sole disambiguation strategy.

The \textbf{h-index} \citep{hirsch2005} and related indicators (\ref{indicators-background}) are commonly used to summarise author-level impact, but they are meaningful only when the underlying publication set is correctly attributed. A split author (one person assigned two IDs) will have a deflated h-index; a merged author (two people sharing one ID) will have an inflated one.

Author-level analysis also raises ethical concerns. The Leiden Manifesto \citep{hicks2015} warns against reducing a researcher's contribution to a single number. Career breaks, discipline-specific norms, and the Matthew effect \citep{merton1968} --- where established researchers accumulate citations disproportionately --- must all be considered when interpreting author-level metrics.

\section{Worked example}\label{author-example}

\subsection{Retrieving an author entity}\label{retrieving-an-author-entity}

We retrieve the OpenAlex profile for a well-known scientometrician.

\begin{Shaded}
\begin{Highlighting}[]
\NormalTok{author\_entity }\OtherTok{\textless{}{-}} \FunctionTok{fetch\_openalex}\NormalTok{(}
  \AttributeTok{entity =} \StringTok{"authors"}\NormalTok{,}
  \AttributeTok{search =} \StringTok{"Ludo Waltman"}\NormalTok{,}
  \AttributeTok{options =} \FunctionTok{list}\NormalTok{(}\AttributeTok{sort =} \StringTok{"cited\_by\_count:desc"}\NormalTok{)}
\NormalTok{)}

\NormalTok{author\_info }\OtherTok{\textless{}{-}}\NormalTok{ author\_entity }\SpecialCharTok{|\textgreater{}} \FunctionTok{slice}\NormalTok{(}\DecValTok{1}\NormalTok{)}
\FunctionTok{cat}\NormalTok{(}\FunctionTok{glue}\NormalTok{(}\StringTok{"Name: \{author\_info$display\_name\}}\SpecialCharTok{\textbackslash{}n}\StringTok{"}\NormalTok{))}
\end{Highlighting}
\end{Shaded}

\begin{verbatim}
#> Name: Ludo Waltman
\end{verbatim}

\begin{Shaded}
\begin{Highlighting}[]
\FunctionTok{cat}\NormalTok{(}\FunctionTok{glue}\NormalTok{(}\StringTok{"OpenAlex ID: \{author\_info$id\}}\SpecialCharTok{\textbackslash{}n}\StringTok{"}\NormalTok{))}
\end{Highlighting}
\end{Shaded}

\begin{verbatim}
#> OpenAlex ID: https://openalex.org/A5027467000
\end{verbatim}

\begin{Shaded}
\begin{Highlighting}[]
\FunctionTok{cat}\NormalTok{(}\FunctionTok{glue}\NormalTok{(}\StringTok{"ORCID: \{author\_info$orcid\}}\SpecialCharTok{\textbackslash{}n}\StringTok{"}\NormalTok{))}
\end{Highlighting}
\end{Shaded}

\begin{verbatim}
#> ORCID: https://orcid.org/0000-0001-8249-1752
\end{verbatim}

\begin{Shaded}
\begin{Highlighting}[]
\FunctionTok{cat}\NormalTok{(}\FunctionTok{glue}\NormalTok{(}\StringTok{"Works count: \{author\_info$works\_count\}}\SpecialCharTok{\textbackslash{}n}\StringTok{"}\NormalTok{))}
\end{Highlighting}
\end{Shaded}

\begin{verbatim}
#> Works count: 478
\end{verbatim}

\begin{Shaded}
\begin{Highlighting}[]
\FunctionTok{cat}\NormalTok{(}\FunctionTok{glue}\NormalTok{(}\StringTok{"Cited by count: \{author\_info$cited\_by\_count\}}\SpecialCharTok{\textbackslash{}n}\StringTok{"}\NormalTok{))}
\end{Highlighting}
\end{Shaded}

\begin{verbatim}
#> Cited by count: 59184
\end{verbatim}

\subsection{Fetching the author's works}\label{fetching-the-authors-works}

\begin{Shaded}
\begin{Highlighting}[]
\NormalTok{author\_works }\OtherTok{\textless{}{-}} \FunctionTok{fetch\_openalex}\NormalTok{(}
  \AttributeTok{entity =} \StringTok{"works"}\NormalTok{,}
  \AttributeTok{author.id =}\NormalTok{ author\_info}\SpecialCharTok{$}\NormalTok{id,}
  \AttributeTok{from\_publication\_date =} \StringTok{"2010{-}01{-}01"}\NormalTok{,}
  \AttributeTok{to\_publication\_date =} \StringTok{"2023{-}12{-}31"}
\NormalTok{)}

\NormalTok{works\_slim }\OtherTok{\textless{}{-}}\NormalTok{ author\_works }\SpecialCharTok{|\textgreater{}}
  \FunctionTok{select}\NormalTok{(id, display\_name, publication\_date, cited\_by\_count, type, source\_display\_name) }\SpecialCharTok{|\textgreater{}}
  \FunctionTok{mutate}\NormalTok{(}\AttributeTok{year =} \FunctionTok{year}\NormalTok{(publication\_date))}

\FunctionTok{cat}\NormalTok{(}\FunctionTok{glue}\NormalTok{(}\StringTok{"Works retrieved: \{nrow(works\_slim)\}}\SpecialCharTok{\textbackslash{}n}\StringTok{"}\NormalTok{))}
\end{Highlighting}
\end{Shaded}

\begin{verbatim}
#> Works retrieved: 209
\end{verbatim}

\subsection{Author-level statistics}\label{author-level-statistics}

\begin{Shaded}
\begin{Highlighting}[]
\NormalTok{h\_idx }\OtherTok{\textless{}{-}} \FunctionTok{compute\_h\_index}\NormalTok{(works\_slim}\SpecialCharTok{$}\NormalTok{cited\_by\_count)}

\NormalTok{stats }\OtherTok{\textless{}{-}} \FunctionTok{tibble}\NormalTok{(}
  \AttributeTok{Metric =} \FunctionTok{c}\NormalTok{(}\StringTok{"Total works"}\NormalTok{, }\StringTok{"Total citations"}\NormalTok{, }\StringTok{"h{-}index"}\NormalTok{,}
             \StringTok{"Mean citations/paper"}\NormalTok{, }\StringTok{"Median citations/paper"}\NormalTok{,}
             \StringTok{"First publication year"}\NormalTok{, }\StringTok{"Most recent year"}\NormalTok{),}
  \AttributeTok{Value =} \FunctionTok{c}\NormalTok{(}
    \FunctionTok{nrow}\NormalTok{(works\_slim),}
    \FunctionTok{sum}\NormalTok{(works\_slim}\SpecialCharTok{$}\NormalTok{cited\_by\_count),}
\NormalTok{    h\_idx,}
    \FunctionTok{round}\NormalTok{(}\FunctionTok{mean}\NormalTok{(works\_slim}\SpecialCharTok{$}\NormalTok{cited\_by\_count), }\DecValTok{1}\NormalTok{),}
    \FunctionTok{median}\NormalTok{(works\_slim}\SpecialCharTok{$}\NormalTok{cited\_by\_count),}
    \FunctionTok{min}\NormalTok{(works\_slim}\SpecialCharTok{$}\NormalTok{year),}
    \FunctionTok{max}\NormalTok{(works\_slim}\SpecialCharTok{$}\NormalTok{year)}
\NormalTok{  )}
\NormalTok{)}

\NormalTok{stats }\SpecialCharTok{|\textgreater{}} \FunctionTok{gt}\NormalTok{()}
\end{Highlighting}
\end{Shaded}

\begingroup
\fontsize{8.0pt}{9.0pt}\selectfont
\begin{longtable}{>{\raggedright\arraybackslash}p{\dimexpr 0.62\linewidth -2\tabcolsep-1.5\arrayrulewidth}>{\raggedleft\arraybackslash}p{\dimexpr 0.38\linewidth -2\tabcolsep-1.5\arrayrulewidth}}
\toprule
Metric & Value \\ 
\midrule\addlinespace[2.5pt]
Total works & 209 \\ 
Total citations & 30189 \\ 
h-index & 53 \\ 
Mean citations/paper & 144 \\ 
Median citations/paper & 4 \\ 
First publication year & 2010 \\ 
Most recent year & 2023 \\ 
\bottomrule
\end{longtable}
\endgroup

\subsection{Publication timeline}\label{publication-timeline}

\begin{Shaded}
\begin{Highlighting}[]
\NormalTok{annual }\OtherTok{\textless{}{-}}\NormalTok{ works\_slim }\SpecialCharTok{|\textgreater{}}
  \FunctionTok{group\_by}\NormalTok{(year) }\SpecialCharTok{|\textgreater{}}
  \FunctionTok{summarise}\NormalTok{(}\AttributeTok{n\_pubs =} \FunctionTok{n}\NormalTok{(), }\AttributeTok{cites =} \FunctionTok{sum}\NormalTok{(cited\_by\_count), }\AttributeTok{.groups =} \StringTok{"drop"}\NormalTok{) }\SpecialCharTok{|\textgreater{}}
  \FunctionTok{mutate}\NormalTok{(}\AttributeTok{cum\_cites =} \FunctionTok{cumsum}\NormalTok{(cites))}

\FunctionTok{ggplot}\NormalTok{(annual, }\FunctionTok{aes}\NormalTok{(}\AttributeTok{x =}\NormalTok{ year)) }\SpecialCharTok{+}
  \FunctionTok{geom\_col}\NormalTok{(}\FunctionTok{aes}\NormalTok{(}\AttributeTok{y =}\NormalTok{ n\_pubs), }\AttributeTok{fill =} \FunctionTok{palette\_sci}\NormalTok{(}\DecValTok{1}\NormalTok{), }\AttributeTok{width =} \FloatTok{0.7}\NormalTok{) }\SpecialCharTok{+}
  \FunctionTok{geom\_line}\NormalTok{(}\FunctionTok{aes}\NormalTok{(}\AttributeTok{y =}\NormalTok{ cum\_cites }\SpecialCharTok{/} \FunctionTok{max}\NormalTok{(cum\_cites) }\SpecialCharTok{*} \FunctionTok{max}\NormalTok{(n\_pubs)),}
            \AttributeTok{colour =} \FunctionTok{palette\_sci}\NormalTok{(}\DecValTok{2}\NormalTok{)[}\DecValTok{2}\NormalTok{], }\AttributeTok{linewidth =} \DecValTok{1}\NormalTok{) }\SpecialCharTok{+}
  \FunctionTok{scale\_y\_continuous}\NormalTok{(}
    \AttributeTok{name =} \StringTok{"Publications per year"}\NormalTok{,}
    \AttributeTok{sec.axis =} \FunctionTok{sec\_axis}\NormalTok{(}\SpecialCharTok{\textasciitilde{}}\NormalTok{ . }\SpecialCharTok{/} \FunctionTok{max}\NormalTok{(annual}\SpecialCharTok{$}\NormalTok{n\_pubs) }\SpecialCharTok{*} \FunctionTok{max}\NormalTok{(annual}\SpecialCharTok{$}\NormalTok{cum\_cites),}
                        \AttributeTok{name =} \StringTok{"Cumulative citations"}\NormalTok{)}
\NormalTok{  ) }\SpecialCharTok{+}
  \FunctionTok{labs}\NormalTok{(}\AttributeTok{x =} \StringTok{"Year"}\NormalTok{) }\SpecialCharTok{+}
  \FunctionTok{theme\_sci}\NormalTok{()}
\end{Highlighting}
\end{Shaded}

\begin{figure}

{\centering \includegraphics[width=0.9\linewidth,alt={Dual-axis chart showing the number of publications per year as bars and cumulative citations as a line, illustrating career trajectory over time.}]{scientometrics-in-r_files/figure-latex/12-fig-author-timeline-1} 

}

\caption{Annual publication output and cumulative citation count.}\label{fig:12-fig-author-timeline}
\end{figure}

\subsection{Top venues}\label{top-venues}

\begin{Shaded}
\begin{Highlighting}[]
\NormalTok{venue\_counts }\OtherTok{\textless{}{-}}\NormalTok{ works\_slim }\SpecialCharTok{|\textgreater{}}
  \FunctionTok{filter}\NormalTok{(}\SpecialCharTok{!}\FunctionTok{is.na}\NormalTok{(source\_display\_name)) }\SpecialCharTok{|\textgreater{}}
  \FunctionTok{count}\NormalTok{(source\_display\_name, }\AttributeTok{sort =} \ConstantTok{TRUE}\NormalTok{) }\SpecialCharTok{|\textgreater{}}
  \FunctionTok{head}\NormalTok{(}\DecValTok{10}\NormalTok{)}

\FunctionTok{ggplot}\NormalTok{(venue\_counts, }\FunctionTok{aes}\NormalTok{(}\AttributeTok{x =}\NormalTok{ n, }\AttributeTok{y =} \FunctionTok{reorder}\NormalTok{(source\_display\_name, n))) }\SpecialCharTok{+}
  \FunctionTok{geom\_col}\NormalTok{(}\AttributeTok{fill =} \FunctionTok{palette\_sci}\NormalTok{(}\DecValTok{1}\NormalTok{)) }\SpecialCharTok{+}
  \FunctionTok{labs}\NormalTok{(}\AttributeTok{x =} \StringTok{"Number of articles"}\NormalTok{, }\AttributeTok{y =} \ConstantTok{NULL}\NormalTok{) }\SpecialCharTok{+}
  \FunctionTok{theme\_sci}\NormalTok{()}
\end{Highlighting}
\end{Shaded}

\begin{figure}

{\centering \includegraphics[width=0.9\linewidth,alt={Horizontal bar chart showing the 10 journals or venues where the author published the most articles.}]{scientometrics-in-r_files/figure-latex/12-fig-author-venues-1} 

}

\caption{Top 10 venues by number of publications for this author.}\label{fig:12-fig-author-venues}
\end{figure}

\subsection{Checking for name variants}\label{checking-for-name-variants}

\begin{Shaded}
\begin{Highlighting}[]
\ControlFlowTok{if}\NormalTok{ (}\StringTok{"display\_name\_alternatives"} \SpecialCharTok{\%in\%} \FunctionTok{names}\NormalTok{(author\_info)) \{}
  \FunctionTok{cat}\NormalTok{(}\StringTok{"Known name variants:}\SpecialCharTok{\textbackslash{}n}\StringTok{"}\NormalTok{)}
  \FunctionTok{print}\NormalTok{(author\_info}\SpecialCharTok{$}\NormalTok{display\_name\_alternatives)}
\NormalTok{\} }\ControlFlowTok{else}\NormalTok{ \{}
  \FunctionTok{cat}\NormalTok{(}\StringTok{"No alternative display names recorded in this entity.}\SpecialCharTok{\textbackslash{}n}\StringTok{"}\NormalTok{)}
\NormalTok{\}}
\end{Highlighting}
\end{Shaded}

\begin{verbatim}
#> Known name variants:
#> [[1]]
#>  [1] "L Waltman"             "L. Waltman"            "LUDO WALTMAN"         
#>  [4] "Ludo Waltman"          "Ludo Waltman, Waltman" "Ludo Waltmann"        
#>  [7] "Waltman L."            "Waltman, L."           "Waltman, L.R."        
#> [10] "Waltman, LR (Ludo)"    "Waltman, Ludo"         "Waltmann, Ludo"
\end{verbatim}

\section{Diagnostics and interpretation}\label{author-diagnostics}

When building author profiles, verify the following:

\begin{itemize}
\tightlist
\item
  \textbf{Works count plausibility}: Does the number of retrieved works match the author's known output? Large discrepancies suggest a split or merged entity.
\item
  \textbf{Topical coherence}: Scan titles and venues. If the works span wildly unrelated fields, the entity may conflate two researchers with similar names.
\item
  \textbf{ORCID match}: If the author has an ORCID, cross-check that the OpenAlex entity links to it.
\item
  \textbf{Co-author consistency}: The author's frequent co-authors should be recognisable collaborators, not strangers from an unrelated field.
\item
  \textbf{Citation distribution}: An extreme outlier (e.g., one paper with 10,000 citations while all others have fewer than 50) may indicate an attribution error.
\end{itemize}

\section{Limitations and responsible use}\label{author-limitations}

\section{Limitations and responsible use}\label{limitations-and-responsible-use-11}

\begin{itemize}
\tightlist
\item
  \textbf{Disambiguation is imperfect.} OpenAlex's algorithm will sometimes split one person into multiple entities or merge distinct individuals. Always validate profiles manually for high-stakes analyses.
\item
  \textbf{ORCID coverage is uneven.} Researchers in the Global South, early-career scholars, and those in humanities adopt ORCID at lower rates. Filtering to ORCID-linked authors introduces bias.
\item
  \textbf{The Matthew effect.} Established researchers accumulate citations disproportionately \citep{merton1968}. Author-level metrics favour seniority and visibility, not necessarily quality or originality.
\item
  \textbf{Career breaks and part-time research.} Indicators like the h-index penalise researchers who take parental leave, work part-time, or switch careers. The m-quotient is only a rough correction \citep{hirsch2005}.
\item
  \textbf{Never evaluate individuals by numbers alone} \citep{hicks2015, dora2012}. Author profiles should supplement, not replace, reading the actual work.
\end{itemize}

\section{Common pitfalls}\label{author-pitfalls}

\section{Common pitfalls}\label{common-pitfalls-11}

\begin{itemize}
\tightlist
\item
  \textbf{Trusting a single database.} OpenAlex may lack works indexed only in domain-specific databases (e.g., SSRN, arXiv preprints not yet linked). Cross-check with ORCID profiles.
\item
  \textbf{Ignoring co-author contributions.} Raw publication counts treat single-author and 50-author papers equally. Consider fractional counting for productivity analysis.
\item
  \textbf{Comparing across career stages.} An early-career researcher with h = 8 after 5 years may be more productive than a senior researcher with h = 25 after 30 years. Use the m-quotient or age-normalised indicators.
\item
  \textbf{Conflating author-level and paper-level metrics.} An author's h-index says nothing about which of their papers are excellent and which are not.
\end{itemize}

\section{Exercises}\label{author-exercises}

\begin{enumerate}
\def\labelenumi{\arabic{enumi}.}
\item
  \textbf{Profile comparison.} Fetch profiles for two researchers in the same field. Compare their h-index, publication count, and mean citations per paper. What does each metric reveal that the others do not?
\item
  \textbf{ORCID validation.} For an author with an ORCID, fetch their works via both OpenAlex author ID and ORCID filtering. Do the two sets agree? What might explain discrepancies?
\item
  \textbf{Name variant detection.} Search OpenAlex for a common name (e.g., ``Wang Wei''). How many author entities are returned? What heuristics could you use to identify the correct one?
\item
  \textbf{Fractional productivity.} Recompute annual publication counts using fractional counting (1/k credit per paper with k authors). How does this change the author's productivity curve?
\item
  \textbf{Self-citation ratio.} For an author's works, estimate the fraction of citations that come from the author's own other papers. \emph{(Hint: check if cited works share the same author ID.)}
\end{enumerate}

\section{Solutions}\label{12-solutions}

Solutions are provided in \ref{solutions}.

\section{Further reading}\label{author-reading}

\begin{itemize}
\tightlist
\item
  \citet{priem2022} --- OpenAlex's author disambiguation and entity resolution approach.
\item
  \citet{hirsch2005} --- The h-index, the most widely used author-level indicator.
\item
  \citet{merton1968} --- The Matthew effect: cumulative advantage in scientific recognition.
\item
  \citet{hicks2015} --- The Leiden Manifesto: principles for responsible author evaluation.
\item
  \citet{dora2012} --- DORA: recommendations against reducing researchers to single numbers.
\item
  \citet{egghe2006} --- The g-index, a complement to the h-index for author evaluation.
\end{itemize}

\section{Session info}\label{author-sessioninfo}

\begin{Shaded}
\begin{Highlighting}[]
\FunctionTok{sessionInfo}\NormalTok{()}
\end{Highlighting}
\end{Shaded}

\begin{verbatim}
#> R version 4.6.1 (2026-06-24)
#> Platform: x86_64-pc-linux-gnu
#> Running under: Ubuntu 24.04.5 LTS
#> 
#> Matrix products: default
#> BLAS:   /usr/lib/x86_64-linux-gnu/openblas-pthread/libblas.so.3 
#> LAPACK: /usr/lib/x86_64-linux-gnu/openblas-pthread/libopenblasp-r0.3.26.so;  LAPACK version 3.12.0
#> 
#> locale:
#>  [1] LC_CTYPE=C.UTF-8       LC_NUMERIC=C           LC_TIME=C.UTF-8       
#>  [4] LC_COLLATE=C.UTF-8     LC_MONETARY=C.UTF-8    LC_MESSAGES=C.UTF-8   
#>  [7] LC_PAPER=C.UTF-8       LC_NAME=C              LC_ADDRESS=C          
#> [10] LC_TELEPHONE=C         LC_MEASUREMENT=C.UTF-8 LC_IDENTIFICATION=C   
#> 
#> time zone: UTC
#> tzcode source: system (glibc)
#> 
#> attached base packages:
#> [1] stats     graphics  grDevices utils     datasets  methods   base     
#> 
#> other attached packages:
#>  [1] quanteda_4.5.1     pdftools_3.9.1     arrow_26.0.0       bibliometrix_5.5.0
#>  [5] RefManageR_1.4.0   bib2df_1.1.2.0     rcrossref_1.2.1    gt_1.4.0          
#>  [9] tidytext_0.4.3     glue_1.8.1         openalexR_3.1.2    lubridate_1.9.5   
#> [13] forcats_1.0.1      stringr_1.6.0      dplyr_1.2.1        purrr_1.2.2       
#> [17] readr_2.2.0        tidyr_1.3.2        tibble_3.3.1       ggplot2_4.0.3     
#> [21] tidyverse_2.0.0   
#> 
#> loaded via a namespace (and not attached):
#>  [1] gridExtra_2.3.1        httr2_1.3.0            readxl_1.6.0          
#>  [4] rlang_1.3.0            magrittr_2.0.5         otel_0.2.0            
#>  [7] compiler_4.6.1         systemfonts_1.3.2      vctrs_0.7.3           
#> [10] httpcode_0.3.0         pkgconfig_2.0.3        fastmap_1.2.0         
#> [13] labeling_0.4.3         ca_0.72                utf8_1.2.6            
#> [16] promises_1.5.0         rmarkdown_2.32         tzdb_0.5.0            
#> [19] bit_4.6.0              xfun_0.62              jsonlite_2.0.0        
#> [22] SnowballC_0.7.1        later_1.4.8            parallel_4.6.1        
#> [25] stopwords_2.3          R6_2.6.1               stringi_1.8.9         
#> [28] RColorBrewer_1.1-3     cellranger_1.1.0       assertthat_0.2.1      
#> [31] Rcpp_1.1.2             bookdown_0.48          knitr_1.52            
#> [34] triebeard_0.4.1        base64enc_0.1-6        httpuv_1.6.17         
#> [37] rentrez_1.2.4          Matrix_1.7-5           igraph_2.3.4          
#> [40] timechange_0.4.0       tidyselect_1.2.1       stringdist_0.9.17     
#> [43] rstudioapi_0.19.0      pubmedR_1.0.2          yaml_2.3.12           
#> [46] viridis_0.6.5          miniUI_0.1.2           humaniformat_0.6.0    
#> [49] curl_8.1.0             qpdf_1.4.1             lattice_0.22-9        
#> [52] plyr_1.8.9             shiny_1.14.0           withr_3.0.3           
#> [55] S7_0.2.2               askpass_1.2.1          evaluate_1.0.5        
#> [58] zip_3.0.2              xml2_1.6.0             shinycssloaders_1.1.0 
#> [61] pillar_1.11.1          janeaustenr_1.0.0      DT_0.34.0             
#> [64] plotly_4.12.1          generics_0.1.4         rprojroot_2.1.1       
#> [67] hms_1.1.4              scales_1.4.0           xtable_1.8-8          
#> [70] contentanalysis_1.1.1  tools_4.6.1            data.table_1.18.6.1   
#> [73] tokenizers_0.3.0       openxlsx_4.2.9         fs_2.1.0              
#> [76] visNetwork_2.1.4       XML_3.99-0.25          fastmatch_1.1-8       
#> [79] grid_4.6.1             bibtex_0.5.3           urltools_1.7.3.1      
#> [82] rscopus_0.9.0          dimensionsR_0.0.3      bibliometrixData_0.3.0
#> [85] cli_3.6.6              viridisLite_0.4.3      gtable_0.3.6          
#> [88] digest_0.6.39          ggrepel_0.9.8          crul_1.6.0            
#> [91] htmlwidgets_1.6.4      farver_2.1.2           htmltools_0.5.9       
#> [94] lifecycle_1.0.5        httr_1.4.9             here_1.0.2            
#> [97] mime_0.13              bit64_4.8.6
\end{verbatim}

\chapter{Institutional and Country Analysis}\label{institutional-and-country-analysis}

\emph{R. Heller · Date: 2026-05-09 · Last changed: 2026-10-31}

\section{Learning objectives}\label{instcountry-objectives}

After completing this chapter, you will be able to:

\begin{itemize}
\tightlist
\item
  Aggregate bibliometric data at institutional and country levels using OpenAlex
\item
  Explain the difference between full counting and fractional counting and when each is appropriate
\item
  Use ROR (Research Organization Registry) identifiers for reliable institutional matching
\item
  Compute descriptive citation summaries and a shared sample/year citation ratio, while explaining why it is not MNCS
\item
  Interpret institutional rankings critically, accounting for size, field mix, and counting method
\end{itemize}

\section{Setup}\label{instcountry-setup}

\begin{Shaded}
\begin{Highlighting}[]
\FunctionTok{library}\NormalTok{(tidyverse)}
\FunctionTok{library}\NormalTok{(openalexR)}
\FunctionTok{library}\NormalTok{(glue)}
\FunctionTok{library}\NormalTok{(gt)}

\FunctionTok{set.seed}\NormalTok{(}\DecValTok{20260509}\NormalTok{)}

\FunctionTok{source}\NormalTok{(here}\SpecialCharTok{::}\FunctionTok{here}\NormalTok{(}\StringTok{"R"}\NormalTok{, }\StringTok{"api\_helpers.R"}\NormalTok{))}
\FunctionTok{source}\NormalTok{(here}\SpecialCharTok{::}\FunctionTok{here}\NormalTok{(}\StringTok{"R"}\NormalTok{, }\StringTok{"utils.R"}\NormalTok{))}
\FunctionTok{source}\NormalTok{(here}\SpecialCharTok{::}\FunctionTok{here}\NormalTok{(}\StringTok{"R"}\NormalTok{, }\StringTok{"sci\_palette.R"}\NormalTok{))}
\FunctionTok{source}\NormalTok{(here}\SpecialCharTok{::}\FunctionTok{here}\NormalTok{(}\StringTok{"R"}\NormalTok{, }\StringTok{"institutional\_samples.R"}\NormalTok{))}
\end{Highlighting}
\end{Shaded}

\section{Conceptual background}\label{instcountry-background}

Comparing research performance across institutions and countries is one of the most common --- and most contentious --- applications of bibliometrics. University rankings, national R\&D assessments, and funding allocation decisions all rely on aggregate bibliometric statistics. Getting the aggregation right matters enormously.

The fundamental choice is between \textbf{full counting} and \textbf{fractional counting}. Under full counting, each institution listed on a paper receives one full credit for that paper. A paper with authors from five institutions thus generates five institutional publications. Under fractional counting, credit is divided: each institution receives 1/5 of a publication count. Full counting inflates the output of institutions that frequently collaborate, while this fractional rule divides bibliometric credit equally among the listed institutions; it does not measure their actual contribution \citep{glanzel2004, waltman2012mncs}.

\textbf{ROR} (Research Organization Registry) identifiers provide a stable, open system for identifying institutions, analogous to DOIs for publications and ORCID for researchers. OpenAlex maps institutional affiliations to ROR IDs, though coverage varies by discipline and geography. Using ROR avoids the pitfalls of string matching on institution names, which are notoriously inconsistent (e.g., ``MIT'', ``Massachusetts Institute of Technology'', ``Mass. Inst. Tech.'').

Size-independent indicators are essential for fair comparison. Raw publication counts favour large universities. The \textbf{Mean Normalised Citation Score} (MNCS) \citep{waltman2012mncs} and \textbf{PP(top 10\%)} --- the proportion of papers in the top 10\% most cited in their field and year --- are designed to account for differences in reference sets. Their interpretation still depends on reference-set construction, coverage, attribution and the citation window. The Leiden Ranking uses these indicators with fractional counting as its default, offering a model of responsible institutional benchmarking.

Country-level analysis faces the same choices. Additionally, linguistic bias in databases means that countries where English is not the primary research language may appear less productive than they truly are \citep{visser2021}.

\section{Worked example}\label{instcountry-example}

\subsection{Selecting institutions}\label{selecting-institutions}

We use five institutions to demonstrate summaries of equally capped samples. The identifiers below are \textbf{OpenAlex institution IDs}, not ROR IDs. Their names, countries and ROR mappings were checked against the OpenAlex institution records on 30 October 2026. Subsidiary coverage still needs checking before an actual assessment.

\begin{Shaded}
\begin{Highlighting}[]
\NormalTok{institutions }\OtherTok{\textless{}{-}} \FunctionTok{tribble}\NormalTok{(}
  \SpecialCharTok{\textasciitilde{}}\NormalTok{short\_name,  }\SpecialCharTok{\textasciitilde{}}\NormalTok{openalex\_id,}
  \StringTok{"MIT"}\NormalTok{,        }\StringTok{"I63966007"}\NormalTok{,}
  \StringTok{"Oxford"}\NormalTok{,     }\StringTok{"I40120149"}\NormalTok{,}
  \StringTok{"ETH Zurich"}\NormalTok{, }\StringTok{"I35440088"}\NormalTok{,}
  \StringTok{"U Tokyo"}\NormalTok{,    }\StringTok{"I74801974"}\NormalTok{,}
  \StringTok{"U Cape Town"}\NormalTok{, }\StringTok{"I157614274"}
\NormalTok{)}
\end{Highlighting}
\end{Shaded}

\subsection{Fetching institutional publications}\label{fetching-institutional-publications}

\begin{Shaded}
\begin{Highlighting}[]
\NormalTok{fetch\_inst\_works }\OtherTok{\textless{}{-}} \ControlFlowTok{function}\NormalTok{(inst\_id) \{}
\NormalTok{  res }\OtherTok{\textless{}{-}} \FunctionTok{fetch\_openalex}\NormalTok{(}
    \AttributeTok{entity =} \StringTok{"works"}\NormalTok{,}
    \AttributeTok{authorships.institutions.id =}\NormalTok{ inst\_id,}
    \AttributeTok{from\_publication\_date =} \StringTok{"2020{-}01{-}01"}\NormalTok{,}
    \AttributeTok{to\_publication\_date =} \StringTok{"2022{-}12{-}31"}\NormalTok{,}
    \AttributeTok{type =} \StringTok{"article"}\NormalTok{,}
    \AttributeTok{options =} \FunctionTok{list}\NormalTok{(}\AttributeTok{sample =} \DecValTok{400}\NormalTok{, }\AttributeTok{seed =} \DecValTok{42}\NormalTok{)}
\NormalTok{  )}
  \ControlFlowTok{if}\NormalTok{ (}\SpecialCharTok{!}\FunctionTok{is.data.frame}\NormalTok{(res) }\SpecialCharTok{||} \SpecialCharTok{!}\FunctionTok{nrow}\NormalTok{(res)) \{}
    \FunctionTok{stop}\NormalTok{(}\FunctionTok{glue}\NormalTok{(}\StringTok{"No usable works returned for \{inst\_id\}; check the request before comparing institutions."}\NormalTok{))}
\NormalTok{  \}}
\NormalTok{  res}
\NormalTok{\}}

\NormalTok{inst\_data }\OtherTok{\textless{}{-}}\NormalTok{ institutions }\SpecialCharTok{|\textgreater{}}
  \FunctionTok{mutate}\NormalTok{(}\AttributeTok{works =} \FunctionTok{map}\NormalTok{(openalex\_id, fetch\_inst\_works))}
\end{Highlighting}
\end{Shaded}

These requests retrieve at most 400 sampled articles per institution from 2020--2022,
not each institution's complete output. Counts below describe the retrieved samples;
the cap prevents interpreting them as comparative institutional productivity.
Provider failures stop the example rather than removing an institution silently.

\subsection{Full counting vs.~fractional counting}\label{full-counting-vs.-fractional-counting}

\begin{Shaded}
\begin{Highlighting}[]
\CommentTok{\# Equal institutional credit per sampled work: 1 / number of distinct affiliations.}
\CommentTok{\# The helper requires the queried institution to be among those affiliations.}
\CommentTok{\# Incomplete affiliations leave the fractional total missing, with coverage shown.}
\NormalTok{compute\_counts }\OtherTok{\textless{}{-}}\NormalTok{ sci\_institution\_sample\_counts}

\NormalTok{counting\_comparison }\OtherTok{\textless{}{-}}\NormalTok{ inst\_data }\SpecialCharTok{|\textgreater{}}
  \FunctionTok{mutate}\NormalTok{(}\AttributeTok{counts =} \FunctionTok{map2}\NormalTok{(works, openalex\_id, compute\_counts)) }\SpecialCharTok{|\textgreater{}}
  \FunctionTok{unnest}\NormalTok{(counts) }\SpecialCharTok{|\textgreater{}}
  \FunctionTok{select}\NormalTok{(short\_name, full\_count, affiliation\_covered, frac\_count) }\SpecialCharTok{|\textgreater{}}
  \FunctionTok{mutate}\NormalTok{(}\AttributeTok{ratio =} \FunctionTok{round}\NormalTok{(frac\_count }\SpecialCharTok{/}\NormalTok{ full\_count, }\DecValTok{2}\NormalTok{))}

\NormalTok{counting\_comparison }\SpecialCharTok{|\textgreater{}}
  \FunctionTok{gt}\NormalTok{() }\SpecialCharTok{|\textgreater{}}
  \FunctionTok{cols\_label}\NormalTok{(}
    \AttributeTok{short\_name =} \StringTok{"Institution"}\NormalTok{,}
    \AttributeTok{full\_count =} \StringTok{"Sampled works"}\NormalTok{,}
    \AttributeTok{affiliation\_covered =} \StringTok{"Works with usable affiliations"}\NormalTok{,}
    \AttributeTok{frac\_count =} \StringTok{"Fractional sample credit"}\NormalTok{,}
    \AttributeTok{ratio =} \StringTok{"Frac / Full"}
\NormalTok{  )}
\end{Highlighting}
\end{Shaded}

\begingroup
\fontsize{8.0pt}{9.0pt}\selectfont
\begin{longtable}{>{\raggedright\arraybackslash}p{\dimexpr 0.19\linewidth -2\tabcolsep-1.5\arrayrulewidth}>{\raggedleft\arraybackslash}p{\dimexpr 0.19\linewidth -2\tabcolsep-1.5\arrayrulewidth}>{\raggedleft\arraybackslash}p{\dimexpr 0.26\linewidth -2\tabcolsep-1.5\arrayrulewidth}>{\raggedleft\arraybackslash}p{\dimexpr 0.19\linewidth -2\tabcolsep-1.5\arrayrulewidth}>{\raggedleft\arraybackslash}p{\dimexpr 0.17\linewidth -2\tabcolsep-1.5\arrayrulewidth}}
\toprule
\nolinkurl{Institution} & Sampled works & Works with usable affiliations & Fractional sample credit & Frac / Full \\ 
\midrule\addlinespace[2.5pt]
MIT & 400 & 342 & NA & NA \\ 
Oxford & 400 & 351 & NA & NA \\ 
ETH Zurich & 400 & 358 & NA & NA \\ 
U Tokyo & 400 & 337 & NA & NA \\ 
U Cape Town & 400 & 344 & NA & NA \\ 
\bottomrule
\end{longtable}
\endgroup

\subsection{Descriptive citation comparison using a shared reference}\label{descriptive-citation-comparison-using-a-shared-reference}

A separate mean for each institution/year would make its mean citation ratio equal
to 1 by construction. Instead, combine the retrieved samples, de-duplicate shared
work IDs, and compute a \textbf{common mean for each publication year}. For an eligible
work, divide its citation count by that pooled year mean, then average those ratios
within each institution's sample. Shared works enter the reference once. Conflicting
snapshots of a shared work cause an error rather than an arbitrary choice.

This is a descriptive \textbf{pooled sample/year ratio}, not field-normalized MNCS.
The reference is the union of five capped institutional samples, not all world
publications. It adjusts for publication year only; current cumulative citation
counts do not impose a common post-publication citation window. Missing counts,
missing years, and zero-mean reference years are excluded from the ratio, with
its usable denominator reported separately.

\begin{Shaded}
\begin{Highlighting}[]
\NormalTok{citation\_reference }\OtherTok{\textless{}{-}} \FunctionTok{sci\_pooled\_year\_reference}\NormalTok{(inst\_data}\SpecialCharTok{$}\NormalTok{works)}

\NormalTok{inst\_impact }\OtherTok{\textless{}{-}}\NormalTok{ inst\_data }\SpecialCharTok{|\textgreater{}}
  \FunctionTok{mutate}\NormalTok{(}\AttributeTok{summary =} \FunctionTok{map}\NormalTok{(works, sci\_sample\_citation\_summary,}
                       \AttributeTok{reference =}\NormalTok{ citation\_reference)) }\SpecialCharTok{|\textgreater{}}
  \FunctionTok{unnest}\NormalTok{(summary) }\SpecialCharTok{|\textgreater{}}
  \FunctionTok{select}\NormalTok{(short\_name, n\_sampled, n\_citation\_observed, mean\_cites,}
\NormalTok{         median\_cites, n\_ratio\_observed, pooled\_year\_ratio) }\SpecialCharTok{|\textgreater{}}
  \FunctionTok{arrange}\NormalTok{(short\_name)}

\NormalTok{inst\_impact }\SpecialCharTok{|\textgreater{}}
  \FunctionTok{gt}\NormalTok{() }\SpecialCharTok{|\textgreater{}}
  \FunctionTok{fmt\_number}\NormalTok{(}\AttributeTok{columns =} \FunctionTok{c}\NormalTok{(mean\_cites, pooled\_year\_ratio), }\AttributeTok{decimals =} \DecValTok{2}\NormalTok{) }\SpecialCharTok{|\textgreater{}}
  \FunctionTok{cols\_label}\NormalTok{(}
    \AttributeTok{short\_name =} \StringTok{"Institution"}\NormalTok{, }\AttributeTok{n\_sampled =} \StringTok{"Sampled works"}\NormalTok{,}
    \AttributeTok{n\_citation\_observed =} \StringTok{"Observed citation counts"}\NormalTok{,}
    \AttributeTok{mean\_cites =} \StringTok{"Mean citations"}\NormalTok{, }\AttributeTok{median\_cites =} \StringTok{"Median citations"}\NormalTok{,}
    \AttributeTok{n\_ratio\_observed =} \StringTok{"Works in ratio"}\NormalTok{,}
    \AttributeTok{pooled\_year\_ratio =} \StringTok{"Pooled sample/year ratio"}
\NormalTok{  )}
\end{Highlighting}
\end{Shaded}

\begingroup
\fontsize{8.0pt}{9.0pt}\selectfont
\begin{longtable}{>{\raggedright\arraybackslash}p{\dimexpr 0.13\linewidth -2\tabcolsep-1.5\arrayrulewidth}>{\raggedleft\arraybackslash}p{\dimexpr 0.12\linewidth -2\tabcolsep-1.5\arrayrulewidth}>{\raggedleft\arraybackslash}p{\dimexpr 0.17\linewidth -2\tabcolsep-1.5\arrayrulewidth}>{\raggedleft\arraybackslash}p{\dimexpr 0.12\linewidth -2\tabcolsep-1.5\arrayrulewidth}>{\raggedleft\arraybackslash}p{\dimexpr 0.14\linewidth -2\tabcolsep-1.5\arrayrulewidth}>{\raggedleft\arraybackslash}p{\dimexpr 0.16\linewidth -2\tabcolsep-1.5\arrayrulewidth}>{\raggedleft\arraybackslash}p{\dimexpr 0.16\linewidth -2\tabcolsep-1.5\arrayrulewidth}}
\toprule
\nolinkurl{Institution} & Sampled works & Observed citation counts & Mean citations & Median citations & Works in ratio & Pooled sample/year ratio \\ 
\midrule\addlinespace[2.5pt]
ETH Zurich & 400 & 400 & 50.67 & 24.5 & 400 & 1.23 \\ 
MIT & 400 & 400 & 64.25 & 26.0 & 400 & 1.54 \\ 
Oxford & 400 & 400 & 38.18 & 15.0 & 400 & 0.93 \\ 
U Cape Town & 400 & 400 & 37.08 & 13.0 & 400 & 0.85 \\ 
U Tokyo & 400 & 400 & 24.01 & 9.0 & 400 & 0.58 \\ 
\bottomrule
\end{longtable}
\endgroup

The helper implementations in \texttt{R/institutional\_samples.R} make the de-duplication,
reference construction and missing-data denominators explicit.

\subsection{Visualization}\label{visualization-6}

\begin{Shaded}
\begin{Highlighting}[]
\NormalTok{counting\_long }\OtherTok{\textless{}{-}}\NormalTok{ counting\_comparison }\SpecialCharTok{|\textgreater{}}
  \FunctionTok{pivot\_longer}\NormalTok{(}\AttributeTok{cols =} \FunctionTok{c}\NormalTok{(full\_count, frac\_count),}
               \AttributeTok{names\_to =} \StringTok{"method"}\NormalTok{, }\AttributeTok{values\_to =} \StringTok{"count"}\NormalTok{) }\SpecialCharTok{|\textgreater{}}
  \FunctionTok{mutate}\NormalTok{(}\AttributeTok{method =} \FunctionTok{recode}\NormalTok{(method,}
                         \AttributeTok{full\_count =} \StringTok{"Full counting"}\NormalTok{,}
                         \AttributeTok{frac\_count =} \StringTok{"Fractional counting"}\NormalTok{))}

\FunctionTok{ggplot}\NormalTok{(counting\_long, }\FunctionTok{aes}\NormalTok{(}\AttributeTok{x =}\NormalTok{ count, }\AttributeTok{y =} \FunctionTok{reorder}\NormalTok{(short\_name, count),}
                          \AttributeTok{fill =}\NormalTok{ method)) }\SpecialCharTok{+}
  \FunctionTok{geom\_col}\NormalTok{(}\AttributeTok{position =} \StringTok{"dodge"}\NormalTok{, }\AttributeTok{width =} \FloatTok{0.7}\NormalTok{) }\SpecialCharTok{+}
  \FunctionTok{scale\_fill\_manual}\NormalTok{(}\AttributeTok{values =} \FunctionTok{palette\_sci}\NormalTok{(}\DecValTok{2}\NormalTok{)) }\SpecialCharTok{+}
  \FunctionTok{labs}\NormalTok{(}\AttributeTok{x =} \StringTok{"Credit within the retrieved sample"}\NormalTok{, }\AttributeTok{y =} \ConstantTok{NULL}\NormalTok{, }\AttributeTok{fill =} \StringTok{"Method"}\NormalTok{) }\SpecialCharTok{+}
  \FunctionTok{theme\_sci}\NormalTok{()}
\end{Highlighting}
\end{Shaded}

\begin{figure}

{\centering \includegraphics[width=0.9\linewidth,alt={Grouped bar chart of sampled-work counts and fractional sample credits, with unavailable fractional totals omitted when affiliations are incomplete.}]{scientometrics-in-r_files/figure-latex/13-fig-instcountry-comparison-1} 

}

\caption{Full and fractional institutional credits in the retrieved article samples; these are not institutional output totals.}\label{fig:13-fig-instcountry-comparison}
\end{figure}

\begin{Shaded}
\begin{Highlighting}[]
\FunctionTok{ggplot}\NormalTok{(inst\_impact, }\FunctionTok{aes}\NormalTok{(}\AttributeTok{x =}\NormalTok{ pooled\_year\_ratio, }\AttributeTok{y =}\NormalTok{ short\_name)) }\SpecialCharTok{+}
  \FunctionTok{geom\_col}\NormalTok{(}\AttributeTok{fill =} \FunctionTok{palette\_sci}\NormalTok{(}\DecValTok{1}\NormalTok{)) }\SpecialCharTok{+}
  \FunctionTok{geom\_vline}\NormalTok{(}\AttributeTok{xintercept =} \DecValTok{1}\NormalTok{, }\AttributeTok{linetype =} \StringTok{"dashed"}\NormalTok{, }\AttributeTok{colour =} \StringTok{"grey50"}\NormalTok{) }\SpecialCharTok{+}
  \FunctionTok{labs}\NormalTok{(}\AttributeTok{x =} \StringTok{"Citation ratio to pooled sample/year reference"}\NormalTok{, }\AttributeTok{y =} \ConstantTok{NULL}\NormalTok{) }\SpecialCharTok{+}
  \FunctionTok{theme\_sci}\NormalTok{()}
\end{Highlighting}
\end{Shaded}

\begin{figure}

{\centering \includegraphics[width=0.9\linewidth,alt={Horizontal bars show each institutional sample's mean ratio to the pooled publication-year reference; the dashed line denotes a ratio of one.}]{scientometrics-in-r_files/figure-latex/13-fig-instcountry-mncs-1} 

}

\caption{Mean citation ratios relative to the shared pooled sample in each publication year. This is not field-normalized MNCS.}\label{fig:13-fig-instcountry-mncs}
\end{figure}

\section{Diagnostics and interpretation}\label{instcountry-diagnostics}

\begin{itemize}
\tightlist
\item
  \textbf{Sample counts are not size}: Capping each sample at 400 articles removes information about institutional output. The full/fractional comparison illustrates counting rules within those samples only.
\item
  \textbf{Field mix}: A medical university will appear more productive (and highly cited) than a social-science institute due to field-level citation norms. True normalisation requires field-level data.
\item
  \textbf{Sample vs.~population}: The reference depends on which institutions and works were sampled. For a genuine assessment, define the target population and reference set, document inclusion probabilities, and use uncertainty methods appropriate to that design.
\item
  \textbf{Fractional/full ratio}: With complete affiliations, the ratio summarizes how institutional credit is divided in the retrieved works. Missing affiliations leave the fractional total unavailable; they are not treated as single-institution papers.
\end{itemize}

\section{Limitations and responsible use}\label{instcountry-limitations}

\section{Limitations and responsible use}\label{limitations-and-responsible-use-12}

\begin{itemize}
\tightlist
\item
  \textbf{Rankings oversimplify.} Institutional and country rankings compress a multidimensional reality into a single number. The Leiden Manifesto \citep{hicks2015} warns against using rankings mechanically.
\item
  \textbf{Affiliation data quality.} Not all papers have complete or correct institutional affiliations. OpenAlex coverage depends on publisher metadata quality.
\item
  \textbf{Counting method changes conclusions.} The same data can yield different institutional rankings depending on whether full or fractional counting is used. Always report the method.
\item
  \textbf{Language and geographic bias.} Databases undercount research published in non-English languages, disadvantaging countries in Latin America, Asia, and Africa \citep{visser2021, singh2021}.
\item
  \textbf{Do not compare across fields without normalisation.} Raw citation differences can reflect disciplinary citation practices as well as other differences; the sample/year ratio does not remove field mix.
\end{itemize}

\section{Common pitfalls}\label{instcountry-pitfalls}

\section{Common pitfalls}\label{common-pitfalls-12}

\begin{itemize}
\tightlist
\item
  \textbf{Using full counting when fractional is more appropriate.} Full counting double-credits collaborative work, inflating output for institutions with many international partnerships.
\item
  \textbf{Ignoring field normalisation.} The pooled sample/year ratio adjusts for year only. It must not be relabeled as MNCS or used as a field-adjusted ranking.
\item
  \textbf{String-matching institution names.} Never match institutions by name strings alone. Use ROR or OpenAlex institution IDs.
\item
  \textbf{Treating samples as populations.} If you sampled 400 works per institution, the resulting metrics have uncertainty. Report this.
\item
  \textbf{Forgetting subsidiary institutions.} Large universities may have hospitals, research institutes, or campuses with separate ROR IDs. Decide on the unit of analysis upfront.
\end{itemize}

\section{Exercises}\label{instcountry-exercises}

\begin{enumerate}
\def\labelenumi{\arabic{enumi}.}
\item
  \textbf{Country-level analysis.} Aggregate publications by author country (using \texttt{authorships.countries}) instead of institution. Compute full and fractional counts for the top 10 countries in your corpus. Which countries are most affected by the counting method?
\item
  \textbf{Top-decile sample share.} Using the same de-duplicated pooled year reference, define a top-decile threshold and an explicit rule for ties. Compute each sample's share above that threshold. Explain why this is not the field-normalized PP(top 10\%) indicator.
\item
  \textbf{Collaboration intensity.} Compute the mean number of distinct institutional affiliations per paper for each institution. Is there a relationship between collaboration intensity and the full/fractional count ratio?
\end{enumerate}

\section{Solutions}\label{13-solutions}

Solutions are provided in \ref{solutions}.

\section{Further reading}\label{instcountry-reading}

\begin{itemize}
\tightlist
\item
  \citet{waltman2012mncs} --- MNCS and size-independent indicators for institutional comparison.
\item
  \citet{glanzel2004} --- Fractional counting in co-authorship and institutional analysis.
\item
  \citet{hicks2015} --- The Leiden Manifesto on responsible institutional benchmarking.
\item
  \citet{waltman2016} --- Review of citation impact indicators applicable to institutions and countries.
\item
  \citet{visser2021} --- Database coverage comparison; crucial context for country-level analysis.
\item
  \citet{singh2021} --- Journal coverage across databases, highlighting geographic imbalances.
\end{itemize}

\section{Session info}\label{instcountry-sessioninfo}

\begin{Shaded}
\begin{Highlighting}[]
\FunctionTok{sessionInfo}\NormalTok{()}
\end{Highlighting}
\end{Shaded}

\begin{verbatim}
#> R version 4.6.1 (2026-06-24)
#> Platform: x86_64-pc-linux-gnu
#> Running under: Ubuntu 24.04.5 LTS
#> 
#> Matrix products: default
#> BLAS:   /usr/lib/x86_64-linux-gnu/openblas-pthread/libblas.so.3 
#> LAPACK: /usr/lib/x86_64-linux-gnu/openblas-pthread/libopenblasp-r0.3.26.so;  LAPACK version 3.12.0
#> 
#> locale:
#>  [1] LC_CTYPE=C.UTF-8       LC_NUMERIC=C           LC_TIME=C.UTF-8       
#>  [4] LC_COLLATE=C.UTF-8     LC_MONETARY=C.UTF-8    LC_MESSAGES=C.UTF-8   
#>  [7] LC_PAPER=C.UTF-8       LC_NAME=C              LC_ADDRESS=C          
#> [10] LC_TELEPHONE=C         LC_MEASUREMENT=C.UTF-8 LC_IDENTIFICATION=C   
#> 
#> time zone: UTC
#> tzcode source: system (glibc)
#> 
#> attached base packages:
#> [1] stats     graphics  grDevices utils     datasets  methods   base     
#> 
#> other attached packages:
#>  [1] quanteda_4.5.1     pdftools_3.9.1     arrow_26.0.0       bibliometrix_5.5.0
#>  [5] RefManageR_1.4.0   bib2df_1.1.2.0     rcrossref_1.2.1    gt_1.4.0          
#>  [9] tidytext_0.4.3     glue_1.8.1         openalexR_3.1.2    lubridate_1.9.5   
#> [13] forcats_1.0.1      stringr_1.6.0      dplyr_1.2.1        purrr_1.2.2       
#> [17] readr_2.2.0        tidyr_1.3.2        tibble_3.3.1       ggplot2_4.0.3     
#> [21] tidyverse_2.0.0   
#> 
#> loaded via a namespace (and not attached):
#>  [1] gridExtra_2.3.1        httr2_1.3.0            readxl_1.6.0          
#>  [4] rlang_1.3.0            magrittr_2.0.5         otel_0.2.0            
#>  [7] compiler_4.6.1         systemfonts_1.3.2      vctrs_0.7.3           
#> [10] httpcode_0.3.0         pkgconfig_2.0.3        fastmap_1.2.0         
#> [13] labeling_0.4.3         ca_0.72                utf8_1.2.6            
#> [16] promises_1.5.0         rmarkdown_2.32         tzdb_0.5.0            
#> [19] bit_4.6.0              xfun_0.62              jsonlite_2.0.0        
#> [22] SnowballC_0.7.1        later_1.4.8            parallel_4.6.1        
#> [25] stopwords_2.3          R6_2.6.1               stringi_1.8.9         
#> [28] RColorBrewer_1.1-3     cellranger_1.1.0       assertthat_0.2.1      
#> [31] Rcpp_1.1.2             bookdown_0.48          knitr_1.52            
#> [34] triebeard_0.4.1        base64enc_0.1-6        httpuv_1.6.17         
#> [37] rentrez_1.2.4          Matrix_1.7-5           igraph_2.3.4          
#> [40] timechange_0.4.0       tidyselect_1.2.1       stringdist_0.9.17     
#> [43] rstudioapi_0.19.0      pubmedR_1.0.2          yaml_2.3.12           
#> [46] viridis_0.6.5          miniUI_0.1.2           humaniformat_0.6.0    
#> [49] curl_8.1.0             qpdf_1.4.1             lattice_0.22-9        
#> [52] plyr_1.8.9             shiny_1.14.0           withr_3.0.3           
#> [55] S7_0.2.2               askpass_1.2.1          evaluate_1.0.5        
#> [58] zip_3.0.2              xml2_1.6.0             shinycssloaders_1.1.0 
#> [61] pillar_1.11.1          janeaustenr_1.0.0      DT_0.34.0             
#> [64] plotly_4.12.1          generics_0.1.4         rprojroot_2.1.1       
#> [67] hms_1.1.4              scales_1.4.0           xtable_1.8-8          
#> [70] contentanalysis_1.1.1  tools_4.6.1            data.table_1.18.6.1   
#> [73] tokenizers_0.3.0       openxlsx_4.2.9         fs_2.1.0              
#> [76] visNetwork_2.1.4       XML_3.99-0.25          fastmatch_1.1-8       
#> [79] grid_4.6.1             bibtex_0.5.3           urltools_1.7.3.1      
#> [82] rscopus_0.9.0          dimensionsR_0.0.3      bibliometrixData_0.3.0
#> [85] cli_3.6.6              viridisLite_0.4.3      gtable_0.3.6          
#> [88] digest_0.6.39          ggrepel_0.9.8          crul_1.6.0            
#> [91] htmlwidgets_1.6.4      farver_2.1.2           htmltools_0.5.9       
#> [94] lifecycle_1.0.5        httr_1.4.9             here_1.0.2            
#> [97] mime_0.13              bit64_4.8.6
\end{verbatim}

\chapter{Temporal Analysis and Citation Aging}\label{temporal-analysis-and-citation-aging}

\emph{R. Heller · Date: 2026-05-09 · Last changed: 2026-10-31}

\section{Learning objectives}\label{temporal-objectives}

After completing this chapter, you will be able to:

\begin{itemize}
\tightlist
\item
  Define and apply citation windows of different lengths
\item
  Compute and interpret the citation half-life for a journal or corpus
\item
  Plot citation aging curves showing how citations accumulate over time
\item
  Calculate the Price index (proportion of references less than five years old)
\item
  Analyse temporal trends in publication output and citation rates
\item
  Explain how aging patterns differ across disciplines
\end{itemize}

\section{Setup}\label{temporal-setup}

\begin{Shaded}
\begin{Highlighting}[]
\FunctionTok{library}\NormalTok{(tidyverse)}
\FunctionTok{library}\NormalTok{(openalexR)}
\FunctionTok{library}\NormalTok{(glue)}
\FunctionTok{library}\NormalTok{(gt)}

\FunctionTok{set.seed}\NormalTok{(}\DecValTok{20260509}\NormalTok{)}

\FunctionTok{source}\NormalTok{(here}\SpecialCharTok{::}\FunctionTok{here}\NormalTok{(}\StringTok{"R"}\NormalTok{, }\StringTok{"api\_helpers.R"}\NormalTok{))}
\FunctionTok{source}\NormalTok{(here}\SpecialCharTok{::}\FunctionTok{here}\NormalTok{(}\StringTok{"R"}\NormalTok{, }\StringTok{"utils.R"}\NormalTok{))}
\FunctionTok{source}\NormalTok{(here}\SpecialCharTok{::}\FunctionTok{here}\NormalTok{(}\StringTok{"R"}\NormalTok{, }\StringTok{"sci\_palette.R"}\NormalTok{))}
\end{Highlighting}
\end{Shaded}

\section{Conceptual background}\label{temporal-background}

Citations are not evenly distributed over time. Most papers receive the bulk of their citations within a few years of publication, then gradually fade. This phenomenon --- \textbf{citation aging} or \textbf{obsolescence} --- was first studied systematically by \citet{price1963}, who observed that scientific literature has a characteristic ``half-life'': the time by which half of all citations to a given year's output have been received.

The \textbf{citation half-life} varies dramatically across disciplines. In fast-moving fields like molecular biology, half-lives may be three to four years; in mathematics or humanities, they can exceed ten years. This has direct consequences for metrics: a two-year Journal Impact Factor systematically undervalues journals in slow-aging fields \citep{waltman2016}.

A \textbf{citation window} is the fixed period after publication during which citations are counted. The JIF uses a two-year window; CiteScore uses four years. Choosing the window involves a trade-off: short windows favour fast-aging fields but miss delayed-recognition papers; long windows are more inclusive but dilute the signal of recent trends.

The \textbf{Price index} --- the proportion of references in a paper's bibliography that are less than five years old --- measures how rapidly a field consumes recent literature. A high Price index indicates a fast-moving frontier field; a low Price index suggests reliance on older foundational works. \citet{price1963} introduced this as an ``immediacy'' measure.

\textbf{Sleeping beauties} --- papers that receive few citations for years then suddenly attract attention \citep{kleinberg2003} --- violate the standard aging model. Their existence is a reminder that citation half-lives describe averages, not individual papers. \citet{garfield1955} already noted that some ideas take decades to be recognised.

Understanding citation temporality is essential for fair evaluation. Comparing citation counts without accounting for publication year, citation window, and field-specific aging rates produces meaningless results.

\section{Worked example}\label{temporal-example}

\subsection{Acquiring the data}\label{acquiring-the-data-1}

We fetch works from two journals with different expected aging profiles.

\begin{Shaded}
\begin{Highlighting}[]
\CommentTok{\# Scientometrics (social science, moderate aging)}
\NormalTok{works\_sciento }\OtherTok{\textless{}{-}} \FunctionTok{fetch\_openalex}\NormalTok{(}
  \AttributeTok{entity =} \StringTok{"works"}\NormalTok{,}
  \AttributeTok{primary\_location.source.id =} \StringTok{"S148561398"}\NormalTok{,}
  \AttributeTok{from\_publication\_date =} \StringTok{"2013{-}01{-}01"}\NormalTok{,}
  \AttributeTok{to\_publication\_date =} \StringTok{"2018{-}12{-}31"}\NormalTok{,}
  \AttributeTok{type =} \StringTok{"article"}\NormalTok{,}
  \AttributeTok{options =} \FunctionTok{list}\NormalTok{(}\AttributeTok{sample =} \DecValTok{400}\NormalTok{, }\AttributeTok{seed =} \DecValTok{42}\NormalTok{)}
\NormalTok{)}

\CommentTok{\# PLOS ONE (multidisciplinary, faster aging)}
\NormalTok{works\_plos }\OtherTok{\textless{}{-}} \FunctionTok{fetch\_openalex}\NormalTok{(}
  \AttributeTok{entity =} \StringTok{"works"}\NormalTok{,}
  \AttributeTok{primary\_location.source.id =} \StringTok{"S202381698"}\NormalTok{,}
  \AttributeTok{from\_publication\_date =} \StringTok{"2013{-}01{-}01"}\NormalTok{,}
  \AttributeTok{to\_publication\_date =} \StringTok{"2018{-}12{-}31"}\NormalTok{,}
  \AttributeTok{type =} \StringTok{"article"}\NormalTok{,}
  \AttributeTok{options =} \FunctionTok{list}\NormalTok{(}\AttributeTok{sample =} \DecValTok{400}\NormalTok{, }\AttributeTok{seed =} \DecValTok{42}\NormalTok{)}
\NormalTok{)}

\NormalTok{works\_all }\OtherTok{\textless{}{-}} \FunctionTok{bind\_rows}\NormalTok{(}
\NormalTok{  works\_sciento }\SpecialCharTok{|\textgreater{}} \FunctionTok{mutate}\NormalTok{(}\AttributeTok{journal =} \StringTok{"Scientometrics"}\NormalTok{),}
\NormalTok{  works\_plos }\SpecialCharTok{|\textgreater{}} \FunctionTok{mutate}\NormalTok{(}\AttributeTok{journal =} \StringTok{"PLOS ONE"}\NormalTok{)}
\NormalTok{) }\SpecialCharTok{|\textgreater{}}
  \FunctionTok{mutate}\NormalTok{(}\AttributeTok{year =} \FunctionTok{year}\NormalTok{(publication\_date),}
         \AttributeTok{age =} \DecValTok{2024} \SpecialCharTok{{-}}\NormalTok{ year)}
\end{Highlighting}
\end{Shaded}

\subsection{Citation aging curves}\label{citation-aging-curves}

We compute mean citations by age (years since publication) for each journal.

\begin{Shaded}
\begin{Highlighting}[]
\NormalTok{aging }\OtherTok{\textless{}{-}}\NormalTok{ works\_all }\SpecialCharTok{|\textgreater{}}
  \FunctionTok{group\_by}\NormalTok{(journal, age) }\SpecialCharTok{|\textgreater{}}
  \FunctionTok{summarise}\NormalTok{(}
    \AttributeTok{mean\_cites =} \FunctionTok{mean}\NormalTok{(cited\_by\_count, }\AttributeTok{na.rm =} \ConstantTok{TRUE}\NormalTok{),}
    \AttributeTok{median\_cites =} \FunctionTok{median}\NormalTok{(cited\_by\_count, }\AttributeTok{na.rm =} \ConstantTok{TRUE}\NormalTok{),}
    \AttributeTok{n =} \FunctionTok{n}\NormalTok{(),}
    \AttributeTok{.groups =} \StringTok{"drop"}
\NormalTok{  ) }\SpecialCharTok{|\textgreater{}}
  \FunctionTok{arrange}\NormalTok{(journal, age)}
\end{Highlighting}
\end{Shaded}

\begin{Shaded}
\begin{Highlighting}[]
\FunctionTok{ggplot}\NormalTok{(aging, }\FunctionTok{aes}\NormalTok{(}\AttributeTok{x =}\NormalTok{ age, }\AttributeTok{y =}\NormalTok{ mean\_cites, }\AttributeTok{colour =}\NormalTok{ journal)) }\SpecialCharTok{+}
  \FunctionTok{geom\_line}\NormalTok{(}\AttributeTok{linewidth =} \DecValTok{1}\NormalTok{) }\SpecialCharTok{+}
  \FunctionTok{geom\_point}\NormalTok{(}\AttributeTok{size =} \DecValTok{2}\NormalTok{) }\SpecialCharTok{+}
  \FunctionTok{scale\_colour\_manual}\NormalTok{(}\AttributeTok{values =} \FunctionTok{palette\_sci}\NormalTok{(}\DecValTok{2}\NormalTok{)) }\SpecialCharTok{+}
  \FunctionTok{labs}\NormalTok{(}\AttributeTok{x =} \StringTok{"Years since publication"}\NormalTok{, }\AttributeTok{y =} \StringTok{"Mean citations"}\NormalTok{,}
       \AttributeTok{colour =} \StringTok{"Journal"}\NormalTok{) }\SpecialCharTok{+}
  \FunctionTok{theme\_sci}\NormalTok{()}
\end{Highlighting}
\end{Shaded}

\begin{figure}

{\centering \includegraphics[width=0.9\linewidth,alt={Line chart showing mean citation count by years since publication for two journals. PLOS ONE articles accumulate citations faster in early years, while Scientometrics articles show a slower but more sustained pattern.}]{scientometrics-in-r_files/figure-latex/14-fig-temporal-aging-1} 

}

\caption{Citation aging curves for Scientometrics and PLOS ONE articles published 2013-2018.}\label{fig:14-fig-temporal-aging}
\end{figure}

\subsection{Estimating citation half-life}\label{estimating-citation-half-life}

We estimate the half-life as the age at which cumulative citations reach 50\% of total.

\begin{Shaded}
\begin{Highlighting}[]
\NormalTok{half\_life }\OtherTok{\textless{}{-}}\NormalTok{ works\_all }\SpecialCharTok{|\textgreater{}}
  \FunctionTok{group\_by}\NormalTok{(journal) }\SpecialCharTok{|\textgreater{}}
  \FunctionTok{arrange}\NormalTok{(age) }\SpecialCharTok{|\textgreater{}}
  \FunctionTok{mutate}\NormalTok{(}
    \AttributeTok{cum\_cites =} \FunctionTok{cumsum}\NormalTok{(cited\_by\_count),}
    \AttributeTok{total\_cites =} \FunctionTok{sum}\NormalTok{(cited\_by\_count),}
    \AttributeTok{cum\_pct =}\NormalTok{ cum\_cites }\SpecialCharTok{/}\NormalTok{ total\_cites}
\NormalTok{  ) }\SpecialCharTok{|\textgreater{}}
  \FunctionTok{summarise}\NormalTok{(}
    \AttributeTok{total\_cites =} \FunctionTok{first}\NormalTok{(total\_cites),}
    \AttributeTok{half\_life\_age =} \FunctionTok{min}\NormalTok{(age[cum\_pct }\SpecialCharTok{\textgreater{}=} \FloatTok{0.50}\NormalTok{]),}
    \AttributeTok{.groups =} \StringTok{"drop"}
\NormalTok{  )}

\NormalTok{half\_life }\SpecialCharTok{|\textgreater{}}
  \FunctionTok{gt}\NormalTok{() }\SpecialCharTok{|\textgreater{}}
  \FunctionTok{cols\_label}\NormalTok{(}
    \AttributeTok{journal =} \StringTok{"Journal"}\NormalTok{,}
    \AttributeTok{total\_cites =} \StringTok{"Total citations"}\NormalTok{,}
    \AttributeTok{half\_life\_age =} \StringTok{"Half{-}life (years)"}
\NormalTok{  )}
\end{Highlighting}
\end{Shaded}

\begingroup
\fontsize{8.0pt}{9.0pt}\selectfont
\begin{longtable}{>{\raggedright\arraybackslash}p{\dimexpr 0.35\linewidth -2\tabcolsep-1.5\arrayrulewidth}>{\raggedleft\arraybackslash}p{\dimexpr 0.31\linewidth -2\tabcolsep-1.5\arrayrulewidth}>{\raggedleft\arraybackslash}p{\dimexpr 0.34\linewidth -2\tabcolsep-1.5\arrayrulewidth}}
\toprule
Journal & Total citations & Half-life (years) \\ 
\midrule\addlinespace[2.5pt]
PLOS ONE & 16442 & 10 \\ 
\nolinkurl{Scientometrics} & 17235 & 8 \\ 
\bottomrule
\end{longtable}
\endgroup

\subsection{Price index}\label{price-index}

The Price index is the fraction of references less than five years old. A full computation requires publication years of each reference; here we show reference-list lengths as a starting point.

\begin{Shaded}
\begin{Highlighting}[]
\NormalTok{works\_all }\SpecialCharTok{|\textgreater{}}
  \FunctionTok{filter}\NormalTok{(}\SpecialCharTok{!}\FunctionTok{is.na}\NormalTok{(referenced\_works\_count), referenced\_works\_count }\SpecialCharTok{\textgreater{}} \DecValTok{0}\NormalTok{) }\SpecialCharTok{|\textgreater{}}
  \FunctionTok{group\_by}\NormalTok{(journal) }\SpecialCharTok{|\textgreater{}}
  \FunctionTok{summarise}\NormalTok{(}\AttributeTok{mean\_refs =} \FunctionTok{round}\NormalTok{(}\FunctionTok{mean}\NormalTok{(referenced\_works\_count), }\DecValTok{1}\NormalTok{), }\AttributeTok{.groups =} \StringTok{"drop"}\NormalTok{)}
\end{Highlighting}
\end{Shaded}

\begin{verbatim}
#> # A tibble: 2 x 2
#>   journal        mean_refs
#>   <chr>              <dbl>
#> 1 PLOS ONE            47.6
#> 2 Scientometrics      45.2
\end{verbatim}

\subsection{Cumulative citation trajectories}\label{cumulative-citation-trajectories}

\begin{Shaded}
\begin{Highlighting}[]
\NormalTok{cum\_aging }\OtherTok{\textless{}{-}}\NormalTok{ works\_all }\SpecialCharTok{|\textgreater{}}
  \FunctionTok{group\_by}\NormalTok{(journal, age) }\SpecialCharTok{|\textgreater{}}
  \FunctionTok{summarise}\NormalTok{(}\AttributeTok{total =} \FunctionTok{sum}\NormalTok{(cited\_by\_count), }\AttributeTok{.groups =} \StringTok{"drop"}\NormalTok{) }\SpecialCharTok{|\textgreater{}}
  \FunctionTok{group\_by}\NormalTok{(journal) }\SpecialCharTok{|\textgreater{}}
  \FunctionTok{arrange}\NormalTok{(age) }\SpecialCharTok{|\textgreater{}}
  \FunctionTok{mutate}\NormalTok{(}
    \AttributeTok{cum\_total =} \FunctionTok{cumsum}\NormalTok{(total),}
    \AttributeTok{cum\_pct =}\NormalTok{ cum\_total }\SpecialCharTok{/} \FunctionTok{sum}\NormalTok{(total)}
\NormalTok{  )}

\FunctionTok{ggplot}\NormalTok{(cum\_aging, }\FunctionTok{aes}\NormalTok{(}\AttributeTok{x =}\NormalTok{ age, }\AttributeTok{y =}\NormalTok{ cum\_pct, }\AttributeTok{colour =}\NormalTok{ journal)) }\SpecialCharTok{+}
  \FunctionTok{geom\_line}\NormalTok{(}\AttributeTok{linewidth =} \DecValTok{1}\NormalTok{) }\SpecialCharTok{+}
  \FunctionTok{geom\_hline}\NormalTok{(}\AttributeTok{yintercept =} \FloatTok{0.5}\NormalTok{, }\AttributeTok{linetype =} \StringTok{"dashed"}\NormalTok{, }\AttributeTok{colour =} \StringTok{"grey50"}\NormalTok{) }\SpecialCharTok{+}
  \FunctionTok{scale\_y\_continuous}\NormalTok{(}\AttributeTok{labels =}\NormalTok{ scales}\SpecialCharTok{::}\NormalTok{percent) }\SpecialCharTok{+}
  \FunctionTok{scale\_colour\_manual}\NormalTok{(}\AttributeTok{values =} \FunctionTok{palette\_sci}\NormalTok{(}\DecValTok{2}\NormalTok{)) }\SpecialCharTok{+}
  \FunctionTok{labs}\NormalTok{(}\AttributeTok{x =} \StringTok{"Years since publication"}\NormalTok{, }\AttributeTok{y =} \StringTok{"Cumulative \% of citations"}\NormalTok{,}
       \AttributeTok{colour =} \StringTok{"Journal"}\NormalTok{) }\SpecialCharTok{+}
  \FunctionTok{annotate}\NormalTok{(}\StringTok{"text"}\NormalTok{, }\AttributeTok{x =} \DecValTok{9}\NormalTok{, }\AttributeTok{y =} \FloatTok{0.52}\NormalTok{, }\AttributeTok{label =} \StringTok{"50\% threshold"}\NormalTok{,}
           \AttributeTok{colour =} \StringTok{"grey40"}\NormalTok{, }\AttributeTok{size =} \FloatTok{3.5}\NormalTok{) }\SpecialCharTok{+}
  \FunctionTok{theme\_sci}\NormalTok{()}
\end{Highlighting}
\end{Shaded}

\begin{figure}

{\centering \includegraphics[width=0.9\linewidth,alt={Line chart showing how the cumulative share of total citations grows with article age for each journal. Steeper curves indicate faster citation accumulation.}]{scientometrics-in-r_files/figure-latex/14-fig-temporal-cumulative-1} 

}

\caption{Cumulative citation percentage by years since publication.}\label{fig:14-fig-temporal-cumulative}
\end{figure}

\section{Diagnostics and interpretation}\label{temporal-diagnostics}

\begin{itemize}
\tightlist
\item
  \textbf{Sample size per age bin}: Ensure each year has enough papers for stable estimates. Small cohorts produce noisy aging curves.
\item
  \textbf{Truncation}: Papers published in the most recent year have not had time to accumulate citations. Exclude or flag incomplete citation windows.
\item
  \textbf{Disciplinary benchmarks}: Compare your half-life estimates against known benchmarks. Social sciences typically have half-lives of 5-8 years; biomedical fields 3-5 years; mathematics 8-12 years.
\item
  \textbf{Outliers}: A few extremely highly cited papers can distort mean-based aging curves. Compare mean and median trajectories.
\end{itemize}

\section{Limitations and responsible use}\label{temporal-limitations}

\section{Limitations and responsible use}\label{limitations-and-responsible-use-13}

\begin{itemize}
\tightlist
\item
  \textbf{Half-life is an average.} Individual papers deviate wildly from the average aging pattern. Sleeping beauties and flash-in-the-pan papers both exist.
\item
  \textbf{Database lag.} Citation counts in OpenAlex may lag behind reality by weeks or months. Very recent citation data is unreliable.
\item
  \textbf{Survivor bias.} Retracted papers or papers removed from databases disappear from aging analyses, potentially biasing half-life estimates.
\item
  \textbf{Citation inflation.} The total volume of citations has grown over time. A paper published in 2020 will accumulate citations faster than an identical paper published in 2000, simply because there are more papers being published \citep{price1963}.
\item
  \textbf{Do not penalise slow-aging fields.} Using short citation windows systematically undervalues mathematics, humanities, and theoretical disciplines \citep{hicks2015}.
\end{itemize}

\section{Common pitfalls}\label{temporal-pitfalls}

\section{Common pitfalls}\label{common-pitfalls-13}

\begin{itemize}
\tightlist
\item
  \textbf{Mixing citation windows.} Comparing a 2-year JIF for one journal with lifetime citations for another produces meaningless rankings.
\item
  \textbf{Ignoring publication-year effects.} A corpus spanning 2000-2023 includes papers with 23 years of citation history alongside papers with 1 year. Normalise by publication year.
\item
  \textbf{Confusing obsolescence with irrelevance.} A declining citation curve does not mean the knowledge is wrong or useless; it may mean the ideas have been absorbed into textbooks or standard practice.
\item
  \textbf{Forgetting about the denominator.} If the number of publications in a field is growing, the absolute citation count grows even if the per-paper impact is constant.
\end{itemize}

\section{Exercises}\label{temporal-exercises}

\begin{enumerate}
\def\labelenumi{\arabic{enumi}.}
\item
  \textbf{Half-life comparison.} Fetch works from a mathematics journal and a biomedical journal. Compute and compare their citation half-lives. Does the difference match disciplinary expectations?
\item
  \textbf{Price index computation.} For a set of papers with known reference lists, compute the actual Price index (fraction of references \textless{} 5 years old). \emph{(Hint: you will need to fetch metadata for the referenced works.)}
\item
  \textbf{Sleeping beauty detection.} Identify papers that received fewer than 2 citations in their first 3 years but more than 20 total. What are these papers about?
\item
  \textbf{Median vs.~mean aging.} Replot the aging curves using median instead of mean citations. How does the shape change, and why?
\item
  \textbf{Publication-cohort sensitivity.} Using the fetched Scientometrics sample, compare current cumulative citation summaries for the 2-, 3- and 5-year publication cohorts ending in 2018. Report observed denominators and explain why changing the publication cohort does not establish a citation-receipt window or compute JIF.
\end{enumerate}

\section{Solutions}\label{14-solutions}

Solutions are provided in \ref{solutions}.

\section{Further reading}\label{temporal-reading}

\begin{itemize}
\tightlist
\item
  \citet{price1963} --- \emph{Little Science, Big Science}; introduced citation aging and the Price index.
\item
  \citet{garfield1955} --- The origin of citation indexing, including early observations on citation temporality.
\item
  \citet{kleinberg2003} --- Bursty patterns in temporal data; relevant to sleeping beauties and sudden citation surges.
\item
  \citet{waltman2016} --- Review of citation indicators, including discussion of citation windows and aging.
\item
  \citet{hicks2015} --- The Leiden Manifesto; warns against penalising slow-aging fields with short citation windows.
\end{itemize}

\section{Session info}\label{temporal-sessioninfo}

\begin{Shaded}
\begin{Highlighting}[]
\FunctionTok{sessionInfo}\NormalTok{()}
\end{Highlighting}
\end{Shaded}

\begin{verbatim}
#> R version 4.6.1 (2026-06-24)
#> Platform: x86_64-pc-linux-gnu
#> Running under: Ubuntu 24.04.5 LTS
#> 
#> Matrix products: default
#> BLAS:   /usr/lib/x86_64-linux-gnu/openblas-pthread/libblas.so.3 
#> LAPACK: /usr/lib/x86_64-linux-gnu/openblas-pthread/libopenblasp-r0.3.26.so;  LAPACK version 3.12.0
#> 
#> locale:
#>  [1] LC_CTYPE=C.UTF-8       LC_NUMERIC=C           LC_TIME=C.UTF-8       
#>  [4] LC_COLLATE=C.UTF-8     LC_MONETARY=C.UTF-8    LC_MESSAGES=C.UTF-8   
#>  [7] LC_PAPER=C.UTF-8       LC_NAME=C              LC_ADDRESS=C          
#> [10] LC_TELEPHONE=C         LC_MEASUREMENT=C.UTF-8 LC_IDENTIFICATION=C   
#> 
#> time zone: UTC
#> tzcode source: system (glibc)
#> 
#> attached base packages:
#> [1] stats     graphics  grDevices utils     datasets  methods   base     
#> 
#> other attached packages:
#>  [1] quanteda_4.5.1     pdftools_3.9.1     arrow_26.0.0       bibliometrix_5.5.0
#>  [5] RefManageR_1.4.0   bib2df_1.1.2.0     rcrossref_1.2.1    gt_1.4.0          
#>  [9] tidytext_0.4.3     glue_1.8.1         openalexR_3.1.2    lubridate_1.9.5   
#> [13] forcats_1.0.1      stringr_1.6.0      dplyr_1.2.1        purrr_1.2.2       
#> [17] readr_2.2.0        tidyr_1.3.2        tibble_3.3.1       ggplot2_4.0.3     
#> [21] tidyverse_2.0.0   
#> 
#> loaded via a namespace (and not attached):
#>  [1] gridExtra_2.3.1        httr2_1.3.0            readxl_1.6.0          
#>  [4] rlang_1.3.0            magrittr_2.0.5         otel_0.2.0            
#>  [7] compiler_4.6.1         systemfonts_1.3.2      vctrs_0.7.3           
#> [10] httpcode_0.3.0         pkgconfig_2.0.3        fastmap_1.2.0         
#> [13] labeling_0.4.3         ca_0.72                utf8_1.2.6            
#> [16] promises_1.5.0         rmarkdown_2.32         tzdb_0.5.0            
#> [19] bit_4.6.0              xfun_0.62              jsonlite_2.0.0        
#> [22] SnowballC_0.7.1        later_1.4.8            parallel_4.6.1        
#> [25] stopwords_2.3          R6_2.6.1               stringi_1.8.9         
#> [28] RColorBrewer_1.1-3     cellranger_1.1.0       assertthat_0.2.1      
#> [31] Rcpp_1.1.2             bookdown_0.48          knitr_1.52            
#> [34] triebeard_0.4.1        base64enc_0.1-6        httpuv_1.6.17         
#> [37] rentrez_1.2.4          Matrix_1.7-5           igraph_2.3.4          
#> [40] timechange_0.4.0       tidyselect_1.2.1       stringdist_0.9.17     
#> [43] rstudioapi_0.19.0      pubmedR_1.0.2          yaml_2.3.12           
#> [46] viridis_0.6.5          miniUI_0.1.2           humaniformat_0.6.0    
#> [49] curl_8.1.0             qpdf_1.4.1             lattice_0.22-9        
#> [52] plyr_1.8.9             shiny_1.14.0           withr_3.0.3           
#> [55] S7_0.2.2               askpass_1.2.1          evaluate_1.0.5        
#> [58] zip_3.0.2              xml2_1.6.0             shinycssloaders_1.1.0 
#> [61] pillar_1.11.1          janeaustenr_1.0.0      DT_0.34.0             
#> [64] plotly_4.12.1          generics_0.1.4         rprojroot_2.1.1       
#> [67] hms_1.1.4              scales_1.4.0           xtable_1.8-8          
#> [70] contentanalysis_1.1.1  tools_4.6.1            data.table_1.18.6.1   
#> [73] tokenizers_0.3.0       openxlsx_4.2.9         fs_2.1.0              
#> [76] visNetwork_2.1.4       XML_3.99-0.25          fastmatch_1.1-8       
#> [79] grid_4.6.1             bibtex_0.5.3           urltools_1.7.3.1      
#> [82] rscopus_0.9.0          dimensionsR_0.0.3      bibliometrixData_0.3.0
#> [85] cli_3.6.6              viridisLite_0.4.3      gtable_0.3.6          
#> [88] digest_0.6.39          ggrepel_0.9.8          crul_1.6.0            
#> [91] htmlwidgets_1.6.4      farver_2.1.2           htmltools_0.5.9       
#> [94] lifecycle_1.0.5        httr_1.4.9             here_1.0.2            
#> [97] mime_0.13              bit64_4.8.6
\end{verbatim}

\part{Network \& Mapping Methods}\label{part-network-mapping-methods}

\chapter{Co-authorship Networks}\label{co-authorship-networks}

\emph{R. Heller · Date: 2026-05-09 · Last changed: 2026-10-31}

\section{Learning objectives}\label{coauth-objectives}

After completing this chapter, you will be able to:

\begin{itemize}
\tightlist
\item
  Construct a co-authorship network from bibliographic data using \texttt{igraph}
\item
  Compute and interpret node-level metrics: degree, betweenness centrality, and clustering coefficient
\item
  Apply Louvain and Leiden community detection algorithms and compare results
\item
  Produce deterministic network layouts that render identically in HTML and PDF
\item
  Identify central authors and structural holes in a collaboration network
\end{itemize}

\section{Setup}\label{coauth-setup}

\begin{Shaded}
\begin{Highlighting}[]
\FunctionTok{library}\NormalTok{(tidyverse)}
\FunctionTok{library}\NormalTok{(openalexR)}
\FunctionTok{library}\NormalTok{(igraph)}
\FunctionTok{library}\NormalTok{(tidygraph)}
\FunctionTok{library}\NormalTok{(ggraph)}
\FunctionTok{library}\NormalTok{(glue)}

\FunctionTok{set.seed}\NormalTok{(}\DecValTok{20260509}\NormalTok{)}

\FunctionTok{source}\NormalTok{(here}\SpecialCharTok{::}\FunctionTok{here}\NormalTok{(}\StringTok{"R"}\NormalTok{, }\StringTok{"api\_helpers.R"}\NormalTok{))}
\FunctionTok{source}\NormalTok{(here}\SpecialCharTok{::}\FunctionTok{here}\NormalTok{(}\StringTok{"R"}\NormalTok{, }\StringTok{"utils.R"}\NormalTok{))}
\FunctionTok{source}\NormalTok{(here}\SpecialCharTok{::}\FunctionTok{here}\NormalTok{(}\StringTok{"R"}\NormalTok{, }\StringTok{"sci\_palette.R"}\NormalTok{))}
\end{Highlighting}
\end{Shaded}

\section{Conceptual background}\label{coauth-background}

Co-authorship is one of the most visible and well-studied forms of scientific collaboration. When two researchers appear together on a publication, they form a tie in the co-authorship network. Studying these networks reveals the social structure of science: who collaborates with whom, which groups are isolated, and which individuals serve as bridges between communities \citep{newman2001}.

Co-authorship networks are typically modelled as \textbf{undirected, weighted graphs}. Nodes represent authors, and an edge connects two authors who have co-authored at least one paper. Edge weights reflect the number of co-authored papers (or, alternatively, fractional counts inversely proportional to the number of authors per paper). These networks tend to exhibit properties common to social networks: a heavy-tailed degree distribution (most authors have few collaborators, a few have many), high clustering (collaborators of collaborators tend to collaborate), and a small-world structure \citep{barabasi2002}.

\textbf{Community detection} identifies groups of authors who collaborate more densely among themselves than with the rest of the network. The Louvain algorithm \citep{blondel2008} is fast and widely used; the Leiden algorithm \citep{traag2019} provides theoretical guarantees that all communities are well-connected, avoiding the ``resolution limit'' problem that can afflict Louvain.

Network analysis of co-authorship data has been widely applied to map research fields \citep{glanzel2004}, identify key opinion leaders, assess interdisciplinary collaboration, and study the evolution of research communities over time. \citet{fortunato2010} provides a comprehensive review of community detection methods applicable to these networks.

\section{Worked example}\label{coauth-example}

\subsection{Acquiring co-authorship data}\label{acquiring-co-authorship-data}

We fetch works from a journal and extract the author pairs.

\begin{Shaded}
\begin{Highlighting}[]
\NormalTok{works }\OtherTok{\textless{}{-}} \FunctionTok{fetch\_openalex}\NormalTok{(}
  \AttributeTok{entity =} \StringTok{"works"}\NormalTok{,}
  \AttributeTok{primary\_location.source.id =} \StringTok{"S148561398"}\NormalTok{,}
  \AttributeTok{from\_publication\_date =} \StringTok{"2021{-}01{-}01"}\NormalTok{,}
  \AttributeTok{to\_publication\_date =} \StringTok{"2023{-}12{-}31"}\NormalTok{,}
  \AttributeTok{options =} \FunctionTok{list}\NormalTok{(}\AttributeTok{sample =} \DecValTok{300}\NormalTok{, }\AttributeTok{seed =} \DecValTok{42}\NormalTok{)}
\NormalTok{)}

\FunctionTok{nrow}\NormalTok{(works)}
\end{Highlighting}
\end{Shaded}

\begin{verbatim}
#> [1] 300
\end{verbatim}

\subsection{Extracting author pairs}\label{extracting-author-pairs}

Each work has a list-column of authors. We unnest and build all pairwise co-author edges.

\begin{Shaded}
\begin{Highlighting}[]
\NormalTok{author\_data }\OtherTok{\textless{}{-}}\NormalTok{ works }\SpecialCharTok{|\textgreater{}}
  \FunctionTok{select}\NormalTok{(id, authorships) }\SpecialCharTok{|\textgreater{}}
  \FunctionTok{unnest}\NormalTok{(authorships, }\AttributeTok{names\_sep =} \StringTok{"\_"}\NormalTok{) }\SpecialCharTok{|\textgreater{}}
  \FunctionTok{select}\NormalTok{(}\AttributeTok{work\_id =}\NormalTok{ id, }\AttributeTok{author\_id =}\NormalTok{ authorships\_id, }\AttributeTok{author\_name =}\NormalTok{ authorships\_display\_name) }\SpecialCharTok{|\textgreater{}}
  \FunctionTok{filter}\NormalTok{(}\SpecialCharTok{!}\FunctionTok{is.na}\NormalTok{(author\_id))}

\NormalTok{edges }\OtherTok{\textless{}{-}}\NormalTok{ author\_data }\SpecialCharTok{|\textgreater{}}
  \FunctionTok{inner\_join}\NormalTok{(author\_data, }\AttributeTok{by =} \StringTok{"work\_id"}\NormalTok{, }\AttributeTok{suffix =} \FunctionTok{c}\NormalTok{(}\StringTok{"\_1"}\NormalTok{, }\StringTok{"\_2"}\NormalTok{),}
             \AttributeTok{relationship =} \StringTok{"many{-}to{-}many"}\NormalTok{) }\SpecialCharTok{|\textgreater{}}
  \FunctionTok{filter}\NormalTok{(author\_id\_1 }\SpecialCharTok{\textless{}}\NormalTok{ author\_id\_2) }\SpecialCharTok{|\textgreater{}}
  \FunctionTok{count}\NormalTok{(author\_id\_1, author\_id\_2,}
        \AttributeTok{name\_1 =}\NormalTok{ author\_name\_1, }\AttributeTok{name\_2 =}\NormalTok{ author\_name\_2,}
        \AttributeTok{name =} \StringTok{"weight"}\NormalTok{)}

\FunctionTok{cat}\NormalTok{(}\FunctionTok{glue}\NormalTok{(}\StringTok{"Authors: \{n\_distinct(c(edges$author\_id\_1, edges$author\_id\_2))\}}\SpecialCharTok{\textbackslash{}n}\StringTok{"}\NormalTok{))}
\end{Highlighting}
\end{Shaded}

\begin{verbatim}
#> Authors: 658
\end{verbatim}

\begin{Shaded}
\begin{Highlighting}[]
\FunctionTok{cat}\NormalTok{(}\FunctionTok{glue}\NormalTok{(}\StringTok{"Edges: \{nrow(edges)\}}\SpecialCharTok{\textbackslash{}n}\StringTok{"}\NormalTok{))}
\end{Highlighting}
\end{Shaded}

\begin{verbatim}
#> Edges: 979
\end{verbatim}

\subsection{Building the network}\label{building-the-network}

\begin{Shaded}
\begin{Highlighting}[]
\NormalTok{g }\OtherTok{\textless{}{-}} \FunctionTok{graph\_from\_data\_frame}\NormalTok{(}
\NormalTok{  edges }\SpecialCharTok{|\textgreater{}} \FunctionTok{select}\NormalTok{(author\_id\_1, author\_id\_2, weight),}
  \AttributeTok{directed =} \ConstantTok{FALSE}
\NormalTok{)}

\NormalTok{g }\OtherTok{\textless{}{-}} \FunctionTok{simplify}\NormalTok{(g, }\AttributeTok{edge.attr.comb =} \FunctionTok{list}\NormalTok{(}\AttributeTok{weight =} \StringTok{"sum"}\NormalTok{))}

\FunctionTok{cat}\NormalTok{(}\FunctionTok{glue}\NormalTok{(}\StringTok{"Nodes: \{vcount(g)\}, Edges: \{ecount(g)\}}\SpecialCharTok{\textbackslash{}n}\StringTok{"}\NormalTok{))}
\end{Highlighting}
\end{Shaded}

\begin{verbatim}
#> Nodes: 658, Edges: 979
\end{verbatim}

\begin{Shaded}
\begin{Highlighting}[]
\FunctionTok{cat}\NormalTok{(}\FunctionTok{glue}\NormalTok{(}\StringTok{"Components: \{components(g)$no\}}\SpecialCharTok{\textbackslash{}n}\StringTok{"}\NormalTok{))}
\end{Highlighting}
\end{Shaded}

\begin{verbatim}
#> Components: 180
\end{verbatim}

\begin{Shaded}
\begin{Highlighting}[]
\FunctionTok{cat}\NormalTok{(}\FunctionTok{glue}\NormalTok{(}\StringTok{"Density: \{round(graph.density(g), 4)\}}\SpecialCharTok{\textbackslash{}n}\StringTok{"}\NormalTok{))}
\end{Highlighting}
\end{Shaded}

\begin{verbatim}
#> Density: 0.0045
\end{verbatim}

\subsection{Network metrics}\label{network-metrics}

\begin{Shaded}
\begin{Highlighting}[]
\FunctionTok{V}\NormalTok{(g)}\SpecialCharTok{$}\NormalTok{degree }\OtherTok{\textless{}{-}} \FunctionTok{degree}\NormalTok{(g)}
\FunctionTok{V}\NormalTok{(g)}\SpecialCharTok{$}\NormalTok{betweenness }\OtherTok{\textless{}{-}} \FunctionTok{betweenness}\NormalTok{(g, }\AttributeTok{normalized =} \ConstantTok{TRUE}\NormalTok{)}
\FunctionTok{V}\NormalTok{(g)}\SpecialCharTok{$}\NormalTok{clustering }\OtherTok{\textless{}{-}} \FunctionTok{transitivity}\NormalTok{(g, }\AttributeTok{type =} \StringTok{"local"}\NormalTok{)}

\NormalTok{author\_lookup }\OtherTok{\textless{}{-}}\NormalTok{ author\_data }\SpecialCharTok{|\textgreater{}}
  \FunctionTok{distinct}\NormalTok{(author\_id, author\_name)}

\NormalTok{metrics }\OtherTok{\textless{}{-}} \FunctionTok{tibble}\NormalTok{(}
  \AttributeTok{author\_id =} \FunctionTok{V}\NormalTok{(g)}\SpecialCharTok{$}\NormalTok{name,}
  \AttributeTok{degree =} \FunctionTok{V}\NormalTok{(g)}\SpecialCharTok{$}\NormalTok{degree,}
  \AttributeTok{betweenness =} \FunctionTok{V}\NormalTok{(g)}\SpecialCharTok{$}\NormalTok{betweenness,}
  \AttributeTok{clustering =} \FunctionTok{V}\NormalTok{(g)}\SpecialCharTok{$}\NormalTok{clustering}
\NormalTok{) }\SpecialCharTok{|\textgreater{}}
  \FunctionTok{left\_join}\NormalTok{(author\_lookup, }\AttributeTok{by =} \StringTok{"author\_id"}\NormalTok{) }\SpecialCharTok{|\textgreater{}}
  \FunctionTok{arrange}\NormalTok{(}\FunctionTok{desc}\NormalTok{(degree))}

\NormalTok{metrics }\SpecialCharTok{|\textgreater{}}
  \FunctionTok{head}\NormalTok{(}\DecValTok{10}\NormalTok{) }\SpecialCharTok{|\textgreater{}}
  \FunctionTok{select}\NormalTok{(author\_name, degree, betweenness, clustering)}
\end{Highlighting}
\end{Shaded}

\begin{verbatim}
#> # A tibble: 10 x 4
#>    author_name         degree betweenness clustering
#>    <chr>                <dbl>       <dbl>      <dbl>
#>  1 Wolfgang Glänzel        13 0.00127          0.192
#>  2 Neil H. Shear           10 0.000000464      1    
#>  3 Robert Gniadecki        10 0.000000464      1    
#>  4 Charles W Lynde         10 0.000000464      1    
#>  5 Zafar Ali               10 0.000104         0.356
#>  6 Elena Netchiporouk      10 0                1    
#>  7 Silvio Peroni           10 0.000125         0.4  
#>  8 Steve Mathieu           10 0.000000464      1    
#>  9 Denis Sasseville        10 0.000000464      1    
#> 10 Anastasiya Muntyanu     10 0.000000464      1
\end{verbatim}

\subsection{Community detection}\label{community-detection}

We apply both Louvain and Leiden to the largest connected component.

\begin{Shaded}
\begin{Highlighting}[]
\NormalTok{comp }\OtherTok{\textless{}{-}} \FunctionTok{components}\NormalTok{(g)}
\NormalTok{giant }\OtherTok{\textless{}{-}} \FunctionTok{induced\_subgraph}\NormalTok{(g, }\FunctionTok{which}\NormalTok{(comp}\SpecialCharTok{$}\NormalTok{membership }\SpecialCharTok{==} \FunctionTok{which.max}\NormalTok{(comp}\SpecialCharTok{$}\NormalTok{csize)))}

\NormalTok{louvain }\OtherTok{\textless{}{-}} \FunctionTok{cluster\_louvain}\NormalTok{(giant)}
\FunctionTok{V}\NormalTok{(giant)}\SpecialCharTok{$}\NormalTok{community\_louvain }\OtherTok{\textless{}{-}} \FunctionTok{membership}\NormalTok{(louvain)}

\FunctionTok{cat}\NormalTok{(}\FunctionTok{glue}\NormalTok{(}\StringTok{"Louvain communities: \{length(unique(membership(louvain)))\}}\SpecialCharTok{\textbackslash{}n}\StringTok{"}\NormalTok{))}
\end{Highlighting}
\end{Shaded}

\begin{verbatim}
#> Louvain communities: 4
\end{verbatim}

\begin{Shaded}
\begin{Highlighting}[]
\FunctionTok{cat}\NormalTok{(}\FunctionTok{glue}\NormalTok{(}\StringTok{"Louvain modularity: \{round(modularity(louvain), 3)\}}\SpecialCharTok{\textbackslash{}n}\StringTok{"}\NormalTok{))}
\end{Highlighting}
\end{Shaded}

\begin{verbatim}
#> Louvain modularity: 0.572
\end{verbatim}

\begin{Shaded}
\begin{Highlighting}[]
\NormalTok{leiden\_result }\OtherTok{\textless{}{-}} \FunctionTok{cluster\_leiden}\NormalTok{(giant, }\AttributeTok{resolution\_parameter =} \FloatTok{1.0}\NormalTok{,}
                                \AttributeTok{objective\_function =} \StringTok{"modularity"}\NormalTok{)}
\FunctionTok{V}\NormalTok{(giant)}\SpecialCharTok{$}\NormalTok{community\_leiden }\OtherTok{\textless{}{-}} \FunctionTok{membership}\NormalTok{(leiden\_result)}

\FunctionTok{cat}\NormalTok{(}\FunctionTok{glue}\NormalTok{(}\StringTok{"Leiden communities: \{length(unique(membership(leiden\_result)))\}}\SpecialCharTok{\textbackslash{}n}\StringTok{"}\NormalTok{))}
\end{Highlighting}
\end{Shaded}

\begin{verbatim}
#> Leiden communities: 4
\end{verbatim}

\begin{Shaded}
\begin{Highlighting}[]
\FunctionTok{cat}\NormalTok{(}\FunctionTok{glue}\NormalTok{(}\StringTok{"Leiden modularity: \{round(modularity(giant, membership(leiden\_result)), 3)\}}\SpecialCharTok{\textbackslash{}n}\StringTok{"}\NormalTok{))}
\end{Highlighting}
\end{Shaded}

\begin{verbatim}
#> Leiden modularity: 0.572
\end{verbatim}

\subsection{Comparing community assignments}\label{comparing-community-assignments}

\begin{Shaded}
\begin{Highlighting}[]
\NormalTok{comparison }\OtherTok{\textless{}{-}} \FunctionTok{tibble}\NormalTok{(}
  \AttributeTok{louvain =} \FunctionTok{V}\NormalTok{(giant)}\SpecialCharTok{$}\NormalTok{community\_louvain,}
  \AttributeTok{leiden =} \FunctionTok{V}\NormalTok{(giant)}\SpecialCharTok{$}\NormalTok{community\_leiden}
\NormalTok{)}

\NormalTok{nmi }\OtherTok{\textless{}{-}}\NormalTok{ igraph}\SpecialCharTok{::}\FunctionTok{compare}\NormalTok{(}
  \FunctionTok{V}\NormalTok{(giant)}\SpecialCharTok{$}\NormalTok{community\_louvain,}
  \FunctionTok{V}\NormalTok{(giant)}\SpecialCharTok{$}\NormalTok{community\_leiden,}
  \AttributeTok{method =} \StringTok{"nmi"}
\NormalTok{)}
\FunctionTok{cat}\NormalTok{(}\FunctionTok{glue}\NormalTok{(}\StringTok{"Normalized Mutual Information (Louvain vs Leiden): \{round(nmi, 3)\}}\SpecialCharTok{\textbackslash{}n}\StringTok{"}\NormalTok{))}
\end{Highlighting}
\end{Shaded}

\begin{verbatim}
#> Normalized Mutual Information (Louvain vs Leiden): 1
\end{verbatim}

\subsection{Visualization}\label{visualization-7}

We use a Fruchterman-Reingold layout with a fixed seed for deterministic output.

\begin{Shaded}
\begin{Highlighting}[]
\NormalTok{giant\_tidy }\OtherTok{\textless{}{-}} \FunctionTok{as\_tbl\_graph}\NormalTok{(giant) }\SpecialCharTok{|\textgreater{}}
  \FunctionTok{mutate}\NormalTok{(}
    \AttributeTok{community =} \FunctionTok{as.factor}\NormalTok{(community\_leiden),}
    \AttributeTok{label =} \FunctionTok{ifelse}\NormalTok{(degree }\SpecialCharTok{\textgreater{}} \FunctionTok{quantile}\NormalTok{(degree, }\FloatTok{0.95}\NormalTok{),}
\NormalTok{                   author\_lookup}\SpecialCharTok{$}\NormalTok{author\_name[}\FunctionTok{match}\NormalTok{(name, author\_lookup}\SpecialCharTok{$}\NormalTok{author\_id)],}
                   \ConstantTok{NA\_character\_}\NormalTok{)}
\NormalTok{  )}

\FunctionTok{set.seed}\NormalTok{(}\DecValTok{42}\NormalTok{)}
\NormalTok{layout }\OtherTok{\textless{}{-}} \FunctionTok{create\_layout}\NormalTok{(giant\_tidy, }\AttributeTok{layout =} \StringTok{"fr"}\NormalTok{)}

\FunctionTok{ggraph}\NormalTok{(layout) }\SpecialCharTok{+}
  \FunctionTok{geom\_edge\_link}\NormalTok{(}\AttributeTok{alpha =} \FloatTok{0.15}\NormalTok{, }\AttributeTok{colour =} \StringTok{"grey60"}\NormalTok{) }\SpecialCharTok{+}
  \FunctionTok{geom\_node\_point}\NormalTok{(}\FunctionTok{aes}\NormalTok{(}\AttributeTok{size =}\NormalTok{ degree, }\AttributeTok{colour =}\NormalTok{ community), }\AttributeTok{alpha =} \FloatTok{0.8}\NormalTok{) }\SpecialCharTok{+}
  \FunctionTok{geom\_node\_text}\NormalTok{(}\FunctionTok{aes}\NormalTok{(}\AttributeTok{label =}\NormalTok{ label), }\AttributeTok{repel =} \ConstantTok{TRUE}\NormalTok{, }\AttributeTok{size =} \FloatTok{2.5}\NormalTok{,}
                 \AttributeTok{max.overlaps =} \DecValTok{15}\NormalTok{, }\AttributeTok{na.rm =} \ConstantTok{TRUE}\NormalTok{) }\SpecialCharTok{+}
  \FunctionTok{scale\_size\_continuous}\NormalTok{(}\AttributeTok{range =} \FunctionTok{c}\NormalTok{(}\DecValTok{1}\NormalTok{, }\DecValTok{6}\NormalTok{), }\AttributeTok{guide =} \StringTok{"none"}\NormalTok{) }\SpecialCharTok{+}
  \FunctionTok{scale\_colour\_manual}\NormalTok{(}\AttributeTok{values =} \FunctionTok{palette\_sci}\NormalTok{(}
    \FunctionTok{n\_distinct}\NormalTok{(giant\_tidy }\SpecialCharTok{|\textgreater{}} \FunctionTok{pull}\NormalTok{(community), }\AttributeTok{na.rm =} \ConstantTok{TRUE}\NormalTok{)}
\NormalTok{  )) }\SpecialCharTok{+}
  \FunctionTok{labs}\NormalTok{(}\AttributeTok{colour =} \StringTok{"Community"}\NormalTok{) }\SpecialCharTok{+}
  \FunctionTok{theme\_void}\NormalTok{(}\AttributeTok{base\_family =} \StringTok{"sans"}\NormalTok{, }\AttributeTok{base\_size =} \DecValTok{11}\NormalTok{) }\SpecialCharTok{+}
  \FunctionTok{theme}\NormalTok{(}\AttributeTok{legend.position =} \StringTok{"bottom"}\NormalTok{)}
\end{Highlighting}
\end{Shaded}

\begin{figure}

{\centering \includegraphics[width=0.9\linewidth,alt={Network graph showing co-authorship relationships. Nodes represent authors and are coloured by community membership detected by the Leiden algorithm. Node size reflects degree centrality. Edges connect co-authors.}]{scientometrics-in-r_files/figure-latex/15-fig-coauth-network-1} 

}

\caption{Co-authorship network of sampled Scientometrics articles (2021–2023), coloured by Leiden community.}\label{fig:15-fig-coauth-network}
\end{figure}

\begin{Shaded}
\begin{Highlighting}[]
\FunctionTok{tibble}\NormalTok{(}\AttributeTok{degree =} \FunctionTok{degree}\NormalTok{(giant)) }\SpecialCharTok{|\textgreater{}}
  \FunctionTok{ggplot}\NormalTok{(}\FunctionTok{aes}\NormalTok{(}\AttributeTok{x =}\NormalTok{ degree)) }\SpecialCharTok{+}
  \FunctionTok{geom\_histogram}\NormalTok{(}\AttributeTok{binwidth =} \DecValTok{1}\NormalTok{, }\AttributeTok{fill =} \FunctionTok{palette\_sci}\NormalTok{(}\DecValTok{1}\NormalTok{), }\AttributeTok{colour =} \StringTok{"white"}\NormalTok{) }\SpecialCharTok{+}
  \FunctionTok{labs}\NormalTok{(}\AttributeTok{x =} \StringTok{"Degree (number of co{-}authors)"}\NormalTok{, }\AttributeTok{y =} \StringTok{"Count"}\NormalTok{) }\SpecialCharTok{+}
  \FunctionTok{theme\_sci}\NormalTok{()}
\end{Highlighting}
\end{Shaded}

\begin{figure}

{\centering \includegraphics[width=0.9\linewidth,alt={Histogram of node degree values showing the right-skewed distribution typical of co-authorship networks, where most authors have few co-authors and a few have many.}]{scientometrics-in-r_files/figure-latex/15-fig-coauth-degree-dist-1} 

}

\caption{Degree distribution of the co-authorship network.}\label{fig:15-fig-coauth-degree-dist}
\end{figure}

\section{Diagnostics and interpretation}\label{coauth-diagnostics}

Key diagnostics for co-authorship networks:

\begin{itemize}
\tightlist
\item
  \textbf{Giant component ratio}: What fraction of nodes are in the largest connected component? A fragmented network may indicate a niche field or a sparse sample.
\item
  \textbf{Degree distribution}: The heavy-tailed distribution is expected. A network where all nodes have similar degree may indicate a problem with the data (e.g., only single-author papers were excluded).
\item
  \textbf{Modularity}: Values above 0.3 typically indicate meaningful community structure. Modularity near zero suggests the network has no clear communities.
\item
  \textbf{Comparison of algorithms}: If Louvain and Leiden produce very different community counts, investigate the resolution parameter. High NMI (\textgreater0.8) indicates strong agreement.
\end{itemize}

\begin{Shaded}
\begin{Highlighting}[]
\NormalTok{giant\_ratio }\OtherTok{\textless{}{-}} \FunctionTok{vcount}\NormalTok{(giant) }\SpecialCharTok{/} \FunctionTok{vcount}\NormalTok{(g)}
\FunctionTok{cat}\NormalTok{(}\FunctionTok{glue}\NormalTok{(}\StringTok{"Giant component: \{vcount(giant)\}/\{vcount(g)\} nodes (\{scales::percent(giant\_ratio)\})}\SpecialCharTok{\textbackslash{}n}\StringTok{"}\NormalTok{))}
\end{Highlighting}
\end{Shaded}

\begin{verbatim}
#> Giant component: 32/658 nodes (5%)
\end{verbatim}

\begin{Shaded}
\begin{Highlighting}[]
\FunctionTok{cat}\NormalTok{(}\FunctionTok{glue}\NormalTok{(}\StringTok{"Mean degree: \{round(mean(degree(giant)), 1)\}}\SpecialCharTok{\textbackslash{}n}\StringTok{"}\NormalTok{))}
\end{Highlighting}
\end{Shaded}

\begin{verbatim}
#> Mean degree: 4.2
\end{verbatim}

\begin{Shaded}
\begin{Highlighting}[]
\FunctionTok{cat}\NormalTok{(}\FunctionTok{glue}\NormalTok{(}\StringTok{"Transitivity (global): \{round(transitivity(giant, type = \textquotesingle{}global\textquotesingle{}), 3)\}}\SpecialCharTok{\textbackslash{}n}\StringTok{"}\NormalTok{))}
\end{Highlighting}
\end{Shaded}

\begin{verbatim}
#> Transitivity (global): 0.54
\end{verbatim}

\section{Limitations and responsible use}\label{coauth-limitations}

\section{Limitations and responsible use}\label{limitations-and-responsible-use-14}

\begin{itemize}
\tightlist
\item
  \textbf{Co-authorship ≠ collaboration.} Not all collaborations result in co-authored publications, and not all co-authorships reflect genuine intellectual contribution (honorary authorship).
\item
  \textbf{Name ambiguity.} Author names in bibliographic databases are not unique identifiers. Despite OpenAlex's entity resolution, errors persist --- especially for common names and transliterated names. Always inspect suspicious high-degree nodes.
\item
  \textbf{Fractional vs.~full counting.} Our construction counts each co-authorship edge equally regardless of the number of authors on a paper. Fractional counting (weight = 1/(n-1) per edge) gives less weight to papers with many authors and may be more appropriate for some analyses \citep{glanzel2004}.
\item
  \textbf{Temporal dynamics.} Static networks aggregate all collaborations into a single snapshot. Consider building time-sliced networks to study how collaboration patterns evolve.
\item
  \textbf{Causality.} Network position (e.g., high betweenness) does not imply leadership or influence. It is a structural observation, not an evaluative judgement \citep{hicks2015}.
\end{itemize}

\section{Common pitfalls}\label{coauth-pitfalls}

\section{Common pitfalls}\label{common-pitfalls-14}

\begin{itemize}
\tightlist
\item
  \textbf{Non-deterministic layouts.} Force-directed layouts (Fruchterman-Reingold, Kamada-Kawai) depend on a random seed. Always \texttt{set.seed()} before computing the layout, or the figure will change on every render.
\item
  \textbf{Including the full network.} Very large networks are unreadable. Filter to the giant component, top-N nodes, or a specific community before plotting.
\item
  \textbf{Ignoring multi-edges and self-loops.} \texttt{igraph::simplify()} is essential after constructing the graph from edge lists that may contain duplicate pairs.
\item
  \textbf{Comparing modularity across networks of different sizes.} Modularity is sensitive to network size and density. Use the Normalized Mutual Information (NMI) to compare community structures.
\item
  \textbf{Overinterpreting small components.} A component with 2--3 authors who co-authored one paper is not a ``community'' --- it is an artifact of the sample.
\end{itemize}

\section{Exercises}\label{coauth-exercises}

\begin{enumerate}
\def\labelenumi{\arabic{enumi}.}
\item
  \textbf{Fractional counting.} Modify the edge construction to use fractional counting (weight = 1/(k-1) where k is the number of authors on the paper). How do the top-degree authors change?
\item
  \textbf{Temporal slicing.} Split the corpus into yearly slices and build a separate network for each year. How does the number of components and mean degree evolve over time?
\item
  \textbf{Betweenness vs.~degree.} Plot betweenness centrality against degree for the giant component. Are the highest-degree authors also the highest-betweenness authors? What does a high-betweenness, low-degree node represent?
\item
  \textbf{Resolution parameter.} Run Leiden community detection with resolution parameters of 0.5, 1.0, and 2.0. How does the number of communities change? Which resolution produces the most interpretable results?
\end{enumerate}

\section{Solutions}\label{15-solutions}

Solutions are provided in \ref{solutions}.

\section{Further reading}\label{coauth-reading}

\begin{itemize}
\tightlist
\item
  \citet{newman2001} --- Foundational analysis of scientific collaboration networks.
\item
  \citet{barabasi2002} --- Evolution of co-authorship networks and their scale-free properties.
\item
  \citet{blondel2008} --- The Louvain community detection algorithm.
\item
  \citet{traag2019} --- The Leiden algorithm, with guarantees on well-connected communities.
\item
  \citet{fortunato2010} --- Comprehensive review of community detection in graphs.
\item
  \citet{glanzel2004} --- Co-authorship analysis and fractional counting in bibliometrics.
\end{itemize}

\section{Session info}\label{coauth-sessioninfo}

\begin{Shaded}
\begin{Highlighting}[]
\FunctionTok{sessionInfo}\NormalTok{()}
\end{Highlighting}
\end{Shaded}

\begin{verbatim}
#> R version 4.6.1 (2026-06-24)
#> Platform: x86_64-pc-linux-gnu
#> Running under: Ubuntu 24.04.5 LTS
#> 
#> Matrix products: default
#> BLAS:   /usr/lib/x86_64-linux-gnu/openblas-pthread/libblas.so.3 
#> LAPACK: /usr/lib/x86_64-linux-gnu/openblas-pthread/libopenblasp-r0.3.26.so;  LAPACK version 3.12.0
#> 
#> locale:
#>  [1] LC_CTYPE=C.UTF-8       LC_NUMERIC=C           LC_TIME=C.UTF-8       
#>  [4] LC_COLLATE=C.UTF-8     LC_MONETARY=C.UTF-8    LC_MESSAGES=C.UTF-8   
#>  [7] LC_PAPER=C.UTF-8       LC_NAME=C              LC_ADDRESS=C          
#> [10] LC_TELEPHONE=C         LC_MEASUREMENT=C.UTF-8 LC_IDENTIFICATION=C   
#> 
#> time zone: UTC
#> tzcode source: system (glibc)
#> 
#> attached base packages:
#> [1] stats     graphics  grDevices utils     datasets  methods   base     
#> 
#> other attached packages:
#>  [1] ggraph_2.2.2       tidygraph_1.3.1    igraph_2.3.4       quanteda_4.5.1    
#>  [5] pdftools_3.9.1     arrow_26.0.0       bibliometrix_5.5.0 RefManageR_1.4.0  
#>  [9] bib2df_1.1.2.0     rcrossref_1.2.1    gt_1.4.0           tidytext_0.4.3    
#> [13] glue_1.8.1         openalexR_3.1.2    lubridate_1.9.5    forcats_1.0.1     
#> [17] stringr_1.6.0      dplyr_1.2.1        purrr_1.2.2        readr_2.2.0       
#> [21] tidyr_1.3.2        tibble_3.3.1       ggplot2_4.0.3      tidyverse_2.0.0   
#> 
#> loaded via a namespace (and not attached):
#>   [1] bibtex_0.5.3           RColorBrewer_1.1-3     rstudioapi_0.19.0     
#>   [4] jsonlite_2.0.0         magrittr_2.0.5         farver_2.1.2          
#>   [7] rmarkdown_2.32         fs_2.1.0               vctrs_0.7.3           
#>  [10] memoise_2.0.1          askpass_1.2.1          base64enc_0.1-6       
#>  [13] htmltools_0.5.9        contentanalysis_1.1.1  curl_8.1.0            
#>  [16] janeaustenr_1.0.0      cellranger_1.1.0       htmlwidgets_1.6.4     
#>  [19] tokenizers_0.3.0       plyr_1.8.9             httr2_1.3.0           
#>  [22] cachem_1.1.0           plotly_4.12.1          dimensionsR_0.0.3     
#>  [25] mime_0.13              lifecycle_1.0.5        pkgconfig_2.0.3       
#>  [28] Matrix_1.7-5           R6_2.6.1               fastmap_1.2.0         
#>  [31] shiny_1.14.0           digest_0.6.39          shinycssloaders_1.1.0 
#>  [34] rprojroot_2.1.1        SnowballC_0.7.1        labeling_0.4.3        
#>  [37] urltools_1.7.3.1       timechange_0.4.0       httr_1.4.9            
#>  [40] polyclip_1.10-7        compiler_4.6.1         here_1.0.2            
#>  [43] bit64_4.8.6            withr_3.0.3            S7_0.2.2              
#>  [46] viridis_0.6.5          ggforce_0.5.0          MASS_7.3-65           
#>  [49] bibliometrixData_0.3.0 tools_4.6.1            otel_0.2.0            
#>  [52] stopwords_2.3          zip_3.0.2              rentrez_1.2.4         
#>  [55] httpuv_1.6.17          promises_1.5.0         grid_4.6.1            
#>  [58] stringdist_0.9.17      generics_0.1.4         gtable_0.3.6          
#>  [61] tzdb_0.5.0             rscopus_0.9.0          ca_0.72               
#>  [64] data.table_1.18.6.1    hms_1.1.4              xml2_1.6.0            
#>  [67] utf8_1.2.6             ggrepel_0.9.8          pillar_1.11.1         
#>  [70] later_1.4.8            tweenr_2.0.3           lattice_0.22-9        
#>  [73] bit_4.6.0              tidyselect_1.2.1       miniUI_0.1.2          
#>  [76] knitr_1.52             gridExtra_2.3.1        bookdown_0.48         
#>  [79] crul_1.6.0             xfun_0.62              graphlayouts_1.2.5    
#>  [82] DT_0.34.0              humaniformat_0.6.0     visNetwork_2.1.4      
#>  [85] stringi_1.8.9          qpdf_1.4.1             yaml_2.3.12           
#>  [88] evaluate_1.0.5         httpcode_0.3.0         cli_3.6.6             
#>  [91] xtable_1.8-8           systemfonts_1.3.2      Rcpp_1.1.2            
#>  [94] readxl_1.6.0           triebeard_0.4.1        XML_3.99-0.25         
#>  [97] parallel_4.6.1         assertthat_0.2.1       pubmedR_1.0.2         
#> [100] viridisLite_0.4.3      scales_1.4.0           openxlsx_4.2.9        
#> [103] rlang_1.3.0            fastmatch_1.1-8
\end{verbatim}

\chapter{Co-citation and Bibliographic Coupling}\label{co-citation-and-bibliographic-coupling}

\emph{R. Heller · Date: 2026-05-09 · Last changed: 2026-10-31}

\section{Learning objectives}\label{cocit-objectives}

After completing this chapter, you will be able to:

\begin{itemize}
\tightlist
\item
  Distinguish co-citation analysis from bibliographic coupling and explain when each is appropriate
\item
  Construct co-citation and bibliographic coupling networks from OpenAlex reference data
\item
  Compute normalised similarity measures (cosine, Jaccard) for citation-based linkages
\item
  Identify core documents and intellectual clusters in a research field
\item
  Interpret co-citation maps as representations of intellectual structure
\end{itemize}

\section{Setup}\label{cocit-setup}

\begin{Shaded}
\begin{Highlighting}[]
\FunctionTok{library}\NormalTok{(tidyverse)}
\FunctionTok{library}\NormalTok{(openalexR)}
\FunctionTok{library}\NormalTok{(igraph)}
\FunctionTok{library}\NormalTok{(tidygraph)}
\FunctionTok{library}\NormalTok{(ggraph)}
\FunctionTok{library}\NormalTok{(glue)}
\FunctionTok{library}\NormalTok{(gt)}

\FunctionTok{set.seed}\NormalTok{(}\DecValTok{20260509}\NormalTok{)}

\FunctionTok{source}\NormalTok{(here}\SpecialCharTok{::}\FunctionTok{here}\NormalTok{(}\StringTok{"R"}\NormalTok{, }\StringTok{"api\_helpers.R"}\NormalTok{))}
\FunctionTok{source}\NormalTok{(here}\SpecialCharTok{::}\FunctionTok{here}\NormalTok{(}\StringTok{"R"}\NormalTok{, }\StringTok{"utils.R"}\NormalTok{))}
\FunctionTok{source}\NormalTok{(here}\SpecialCharTok{::}\FunctionTok{here}\NormalTok{(}\StringTok{"R"}\NormalTok{, }\StringTok{"sci\_palette.R"}\NormalTok{))}
\end{Highlighting}
\end{Shaded}

\section{Conceptual background}\label{cocit-background}

While co-authorship networks (\ref{coauth-background}) reveal social structure, \textbf{co-citation} and \textbf{bibliographic coupling} networks reveal \emph{intellectual} structure --- the conceptual relationships between documents as expressed through citation behaviour.

\textbf{Co-citation analysis}, introduced independently by \citet{small1973} and Marshakova-Shaikevich in 1973, measures the relatedness of two documents by how often they are cited together by other papers. If documents A and B frequently appear in the same reference lists, they likely address related topics or represent complementary contributions to a field. Co-citation strength changes over time: as a field evolves, new citing papers create new co-citation links and strengthen or weaken existing ones.

\textbf{Bibliographic coupling}, proposed by \citet{kessler1963}, takes the complementary perspective. Two documents are bibliographically coupled if they share one or more references. Unlike co-citation, coupling strength is fixed at publication time --- once a paper is published, its reference list does not change. This makes bibliographic coupling especially useful for analysing the most recent literature, where co-citation links have not yet had time to form.

The two methods reveal different facets of intellectual structure. Co-citation maps the \emph{reception} of ideas (how the community groups prior work); bibliographic coupling maps the \emph{production} of knowledge (how authors draw on shared foundations). \citet{waltman2012} demonstrated that combining both perspectives yields richer maps of science than either alone.

Both methods produce \textbf{similarity matrices} that can be converted to networks. Common normalisation approaches include cosine similarity (emphasising relative overlap) and the Jaccard index (penalising pairs with many unshared references). Unnormalised raw counts are dominated by highly-cited papers; normalisation is essential for balanced network construction.

\section{Worked example}\label{cocit-example}

\subsection{Fetching citing and cited works}\label{fetching-citing-and-cited-works}

We fetch a corpus of scientometrics research and extract reference lists.

\begin{Shaded}
\begin{Highlighting}[]
\NormalTok{works }\OtherTok{\textless{}{-}} \FunctionTok{fetch\_openalex}\NormalTok{(}
  \AttributeTok{entity =} \StringTok{"works"}\NormalTok{,}
  \AttributeTok{primary\_location.source.id =} \StringTok{"S148561398"}\NormalTok{,}
  \AttributeTok{from\_publication\_date =} \StringTok{"2020{-}01{-}01"}\NormalTok{,}
  \AttributeTok{to\_publication\_date =} \StringTok{"2023{-}12{-}31"}\NormalTok{,}
  \AttributeTok{options =} \FunctionTok{list}\NormalTok{(}\AttributeTok{sample =} \DecValTok{200}\NormalTok{, }\AttributeTok{seed =} \DecValTok{42}\NormalTok{)}
\NormalTok{)}

\NormalTok{refs }\OtherTok{\textless{}{-}}\NormalTok{ works }\SpecialCharTok{|\textgreater{}}
  \FunctionTok{select}\NormalTok{(}\AttributeTok{citing\_id =}\NormalTok{ id, referenced\_works) }\SpecialCharTok{|\textgreater{}}
  \FunctionTok{unnest}\NormalTok{(referenced\_works) }\SpecialCharTok{|\textgreater{}}
  \FunctionTok{rename}\NormalTok{(}\AttributeTok{cited\_id =}\NormalTok{ referenced\_works)}

\FunctionTok{cat}\NormalTok{(}\FunctionTok{glue}\NormalTok{(}\StringTok{"Citing works: \{n\_distinct(refs$citing\_id)\}}\SpecialCharTok{\textbackslash{}n}\StringTok{"}\NormalTok{))}
\end{Highlighting}
\end{Shaded}

\begin{verbatim}
#> Citing works: 200
\end{verbatim}

\begin{Shaded}
\begin{Highlighting}[]
\FunctionTok{cat}\NormalTok{(}\FunctionTok{glue}\NormalTok{(}\StringTok{"Cited works: \{n\_distinct(refs$cited\_id)\}}\SpecialCharTok{\textbackslash{}n}\StringTok{"}\NormalTok{))}
\end{Highlighting}
\end{Shaded}

\begin{verbatim}
#> Cited works: 8112
\end{verbatim}

\begin{Shaded}
\begin{Highlighting}[]
\FunctionTok{cat}\NormalTok{(}\FunctionTok{glue}\NormalTok{(}\StringTok{"Citation links: \{nrow(refs)\}}\SpecialCharTok{\textbackslash{}n}\StringTok{"}\NormalTok{))}
\end{Highlighting}
\end{Shaded}

\begin{verbatim}
#> Citation links: 9479
\end{verbatim}

\subsection{Building a co-citation network}\label{building-a-co-citation-network}

Two cited documents are co-cited when they appear in the same reference list. We count shared citing papers.

\begin{Shaded}
\begin{Highlighting}[]
\NormalTok{cocit\_pairs }\OtherTok{\textless{}{-}}\NormalTok{ refs }\SpecialCharTok{|\textgreater{}}
  \FunctionTok{inner\_join}\NormalTok{(refs, }\AttributeTok{by =} \StringTok{"citing\_id"}\NormalTok{, }\AttributeTok{suffix =} \FunctionTok{c}\NormalTok{(}\StringTok{"\_a"}\NormalTok{, }\StringTok{"\_b"}\NormalTok{),}
             \AttributeTok{relationship =} \StringTok{"many{-}to{-}many"}\NormalTok{) }\SpecialCharTok{|\textgreater{}}
  \FunctionTok{filter}\NormalTok{(cited\_id\_a }\SpecialCharTok{\textless{}}\NormalTok{ cited\_id\_b) }\SpecialCharTok{|\textgreater{}}
  \FunctionTok{count}\NormalTok{(cited\_id\_a, cited\_id\_b, }\AttributeTok{name =} \StringTok{"cocit\_count"}\NormalTok{)}

\NormalTok{cocit\_top }\OtherTok{\textless{}{-}}\NormalTok{ cocit\_pairs }\SpecialCharTok{|\textgreater{}}
  \FunctionTok{filter}\NormalTok{(cocit\_count }\SpecialCharTok{\textgreater{}=} \DecValTok{3}\NormalTok{)}

\FunctionTok{cat}\NormalTok{(}\FunctionTok{glue}\NormalTok{(}\StringTok{"Co{-}citation pairs (count \textgreater{}= 3): \{nrow(cocit\_top)\}}\SpecialCharTok{\textbackslash{}n}\StringTok{"}\NormalTok{))}
\end{Highlighting}
\end{Shaded}

\begin{verbatim}
#> Co-citation pairs (count >= 3): 160
\end{verbatim}

\begin{Shaded}
\begin{Highlighting}[]
\NormalTok{g\_cocit }\OtherTok{\textless{}{-}} \FunctionTok{graph\_from\_data\_frame}\NormalTok{(}
\NormalTok{  cocit\_top }\SpecialCharTok{|\textgreater{}} \FunctionTok{select}\NormalTok{(cited\_id\_a, cited\_id\_b, }\AttributeTok{weight =}\NormalTok{ cocit\_count),}
  \AttributeTok{directed =} \ConstantTok{FALSE}
\NormalTok{) }\SpecialCharTok{|\textgreater{}}
  \FunctionTok{simplify}\NormalTok{(}\AttributeTok{edge.attr.comb =} \FunctionTok{list}\NormalTok{(}\AttributeTok{weight =} \StringTok{"sum"}\NormalTok{))}

\FunctionTok{cat}\NormalTok{(}\FunctionTok{glue}\NormalTok{(}\StringTok{"Co{-}citation network: \{vcount(g\_cocit)\} nodes, \{ecount(g\_cocit)\} edges}\SpecialCharTok{\textbackslash{}n}\StringTok{"}\NormalTok{))}
\end{Highlighting}
\end{Shaded}

\begin{verbatim}
#> Co-citation network: 124 nodes, 160 edges
\end{verbatim}

\subsection{Building a bibliographic coupling network}\label{building-a-bibliographic-coupling-network}

Two citing documents are coupled when they share at least one reference.

\begin{Shaded}
\begin{Highlighting}[]
\NormalTok{bibcoup\_pairs }\OtherTok{\textless{}{-}}\NormalTok{ refs }\SpecialCharTok{|\textgreater{}}
  \FunctionTok{inner\_join}\NormalTok{(refs, }\AttributeTok{by =} \StringTok{"cited\_id"}\NormalTok{, }\AttributeTok{suffix =} \FunctionTok{c}\NormalTok{(}\StringTok{"\_a"}\NormalTok{, }\StringTok{"\_b"}\NormalTok{),}
             \AttributeTok{relationship =} \StringTok{"many{-}to{-}many"}\NormalTok{) }\SpecialCharTok{|\textgreater{}}
  \FunctionTok{filter}\NormalTok{(citing\_id\_a }\SpecialCharTok{\textless{}}\NormalTok{ citing\_id\_b) }\SpecialCharTok{|\textgreater{}}
  \FunctionTok{count}\NormalTok{(citing\_id\_a, citing\_id\_b, }\AttributeTok{name =} \StringTok{"shared\_refs"}\NormalTok{)}

\NormalTok{bibcoup\_top }\OtherTok{\textless{}{-}}\NormalTok{ bibcoup\_pairs }\SpecialCharTok{|\textgreater{}}
  \FunctionTok{filter}\NormalTok{(shared\_refs }\SpecialCharTok{\textgreater{}=} \DecValTok{5}\NormalTok{)}

\FunctionTok{cat}\NormalTok{(}\FunctionTok{glue}\NormalTok{(}\StringTok{"Bibliographic coupling pairs (shared refs \textgreater{}= 5): \{nrow(bibcoup\_top)\}}\SpecialCharTok{\textbackslash{}n}\StringTok{"}\NormalTok{))}
\end{Highlighting}
\end{Shaded}

\begin{verbatim}
#> Bibliographic coupling pairs (shared refs >= 5): 60
\end{verbatim}

\begin{Shaded}
\begin{Highlighting}[]
\NormalTok{g\_bibcoup }\OtherTok{\textless{}{-}} \FunctionTok{graph\_from\_data\_frame}\NormalTok{(}
\NormalTok{  bibcoup\_top }\SpecialCharTok{|\textgreater{}} \FunctionTok{select}\NormalTok{(citing\_id\_a, citing\_id\_b, }\AttributeTok{weight =}\NormalTok{ shared\_refs),}
  \AttributeTok{directed =} \ConstantTok{FALSE}
\NormalTok{) }\SpecialCharTok{|\textgreater{}}
  \FunctionTok{simplify}\NormalTok{(}\AttributeTok{edge.attr.comb =} \FunctionTok{list}\NormalTok{(}\AttributeTok{weight =} \StringTok{"sum"}\NormalTok{))}

\FunctionTok{cat}\NormalTok{(}\FunctionTok{glue}\NormalTok{(}\StringTok{"Coupling network: \{vcount(g\_bibcoup)\} nodes, \{ecount(g\_bibcoup)\} edges}\SpecialCharTok{\textbackslash{}n}\StringTok{"}\NormalTok{))}
\end{Highlighting}
\end{Shaded}

\begin{verbatim}
#> Coupling network: 63 nodes, 60 edges
\end{verbatim}

\subsection{Visualisation}\label{visualisation}

\begin{Shaded}
\begin{Highlighting}[]
\NormalTok{comp }\OtherTok{\textless{}{-}} \FunctionTok{components}\NormalTok{(g\_cocit)}
\NormalTok{giant }\OtherTok{\textless{}{-}} \FunctionTok{induced\_subgraph}\NormalTok{(g\_cocit, }\FunctionTok{which}\NormalTok{(comp}\SpecialCharTok{$}\NormalTok{membership }\SpecialCharTok{==} \FunctionTok{which.max}\NormalTok{(comp}\SpecialCharTok{$}\NormalTok{csize)))}
\FunctionTok{V}\NormalTok{(giant)}\SpecialCharTok{$}\NormalTok{community }\OtherTok{\textless{}{-}} \FunctionTok{as.factor}\NormalTok{(}\FunctionTok{membership}\NormalTok{(}\FunctionTok{cluster\_leiden}\NormalTok{(}
\NormalTok{  giant, }\AttributeTok{resolution\_parameter =} \FloatTok{1.0}\NormalTok{, }\AttributeTok{objective\_function =} \StringTok{"modularity"}
\NormalTok{)))}
\FunctionTok{V}\NormalTok{(giant)}\SpecialCharTok{$}\NormalTok{degree }\OtherTok{\textless{}{-}} \FunctionTok{degree}\NormalTok{(giant)}

\FunctionTok{set.seed}\NormalTok{(}\DecValTok{42}\NormalTok{)}
\NormalTok{layout }\OtherTok{\textless{}{-}} \FunctionTok{create\_layout}\NormalTok{(}\FunctionTok{as\_tbl\_graph}\NormalTok{(giant), }\AttributeTok{layout =} \StringTok{"fr"}\NormalTok{)}

\FunctionTok{ggraph}\NormalTok{(layout) }\SpecialCharTok{+}
  \FunctionTok{geom\_edge\_link}\NormalTok{(}\AttributeTok{alpha =} \FloatTok{0.1}\NormalTok{, }\AttributeTok{colour =} \StringTok{"grey60"}\NormalTok{) }\SpecialCharTok{+}
  \FunctionTok{geom\_node\_point}\NormalTok{(}\FunctionTok{aes}\NormalTok{(}\AttributeTok{size =}\NormalTok{ degree, }\AttributeTok{colour =}\NormalTok{ community), }\AttributeTok{alpha =} \FloatTok{0.8}\NormalTok{) }\SpecialCharTok{+}
  \FunctionTok{scale\_size\_continuous}\NormalTok{(}\AttributeTok{range =} \FunctionTok{c}\NormalTok{(}\DecValTok{1}\NormalTok{, }\DecValTok{5}\NormalTok{), }\AttributeTok{guide =} \StringTok{"none"}\NormalTok{) }\SpecialCharTok{+}
  \FunctionTok{scale\_colour\_manual}\NormalTok{(}\AttributeTok{values =} \FunctionTok{palette\_sci}\NormalTok{(}
    \FunctionTok{n\_distinct}\NormalTok{(}\FunctionTok{V}\NormalTok{(giant)}\SpecialCharTok{$}\NormalTok{community)}
\NormalTok{  )) }\SpecialCharTok{+}
  \FunctionTok{labs}\NormalTok{(}\AttributeTok{colour =} \StringTok{"Community"}\NormalTok{) }\SpecialCharTok{+}
  \FunctionTok{theme\_void}\NormalTok{(}\AttributeTok{base\_family =} \StringTok{"sans"}\NormalTok{, }\AttributeTok{base\_size =} \DecValTok{11}\NormalTok{) }\SpecialCharTok{+}
  \FunctionTok{theme}\NormalTok{(}\AttributeTok{legend.position =} \StringTok{"bottom"}\NormalTok{)}
\end{Highlighting}
\end{Shaded}

\begin{figure}

{\centering \includegraphics[width=0.9\linewidth,alt={Network graph showing co-citation relationships between documents. Nodes are cited papers; edges connect papers frequently cited together. Node size reflects degree. Communities are coloured.}]{scientometrics-in-r_files/figure-latex/16-fig-cocit-network-1} 

}

\caption{Co-citation network of frequently co-cited references in Scientometrics (2020–2023).}\label{fig:16-fig-cocit-network}
\end{figure}

\begin{Shaded}
\begin{Highlighting}[]
\NormalTok{comp\_bc }\OtherTok{\textless{}{-}} \FunctionTok{components}\NormalTok{(g\_bibcoup)}
\NormalTok{giant\_bc }\OtherTok{\textless{}{-}} \FunctionTok{induced\_subgraph}\NormalTok{(g\_bibcoup,}
                             \FunctionTok{which}\NormalTok{(comp\_bc}\SpecialCharTok{$}\NormalTok{membership }\SpecialCharTok{==} \FunctionTok{which.max}\NormalTok{(comp\_bc}\SpecialCharTok{$}\NormalTok{csize)))}
\FunctionTok{V}\NormalTok{(giant\_bc)}\SpecialCharTok{$}\NormalTok{community }\OtherTok{\textless{}{-}} \FunctionTok{as.factor}\NormalTok{(}\FunctionTok{membership}\NormalTok{(}\FunctionTok{cluster\_leiden}\NormalTok{(}
\NormalTok{  giant\_bc, }\AttributeTok{resolution\_parameter =} \FloatTok{1.0}\NormalTok{, }\AttributeTok{objective\_function =} \StringTok{"modularity"}
\NormalTok{)))}
\FunctionTok{V}\NormalTok{(giant\_bc)}\SpecialCharTok{$}\NormalTok{degree }\OtherTok{\textless{}{-}} \FunctionTok{degree}\NormalTok{(giant\_bc)}

\FunctionTok{set.seed}\NormalTok{(}\DecValTok{42}\NormalTok{)}
\NormalTok{layout\_bc }\OtherTok{\textless{}{-}} \FunctionTok{create\_layout}\NormalTok{(}\FunctionTok{as\_tbl\_graph}\NormalTok{(giant\_bc), }\AttributeTok{layout =} \StringTok{"fr"}\NormalTok{)}

\FunctionTok{ggraph}\NormalTok{(layout\_bc) }\SpecialCharTok{+}
  \FunctionTok{geom\_edge\_link}\NormalTok{(}\AttributeTok{alpha =} \FloatTok{0.1}\NormalTok{, }\AttributeTok{colour =} \StringTok{"grey60"}\NormalTok{) }\SpecialCharTok{+}
  \FunctionTok{geom\_node\_point}\NormalTok{(}\FunctionTok{aes}\NormalTok{(}\AttributeTok{size =}\NormalTok{ degree, }\AttributeTok{colour =}\NormalTok{ community), }\AttributeTok{alpha =} \FloatTok{0.8}\NormalTok{) }\SpecialCharTok{+}
  \FunctionTok{scale\_size\_continuous}\NormalTok{(}\AttributeTok{range =} \FunctionTok{c}\NormalTok{(}\DecValTok{1}\NormalTok{, }\DecValTok{5}\NormalTok{), }\AttributeTok{guide =} \StringTok{"none"}\NormalTok{) }\SpecialCharTok{+}
  \FunctionTok{scale\_colour\_manual}\NormalTok{(}\AttributeTok{values =} \FunctionTok{palette\_sci}\NormalTok{(}
    \FunctionTok{n\_distinct}\NormalTok{(}\FunctionTok{V}\NormalTok{(giant\_bc)}\SpecialCharTok{$}\NormalTok{community)}
\NormalTok{  )) }\SpecialCharTok{+}
  \FunctionTok{labs}\NormalTok{(}\AttributeTok{colour =} \StringTok{"Community"}\NormalTok{) }\SpecialCharTok{+}
  \FunctionTok{theme\_void}\NormalTok{(}\AttributeTok{base\_family =} \StringTok{"sans"}\NormalTok{, }\AttributeTok{base\_size =} \DecValTok{11}\NormalTok{) }\SpecialCharTok{+}
  \FunctionTok{theme}\NormalTok{(}\AttributeTok{legend.position =} \StringTok{"bottom"}\NormalTok{)}
\end{Highlighting}
\end{Shaded}

\begin{figure}

{\centering \includegraphics[width=0.9\linewidth,alt={Network graph showing bibliographic coupling relationships between articles. Edges connect papers that share many references. Coloured by community.}]{scientometrics-in-r_files/figure-latex/16-fig-bibcoup-network-1} 

}

\caption{Bibliographic coupling network of Scientometrics articles (2020–2023).}\label{fig:16-fig-bibcoup-network}
\end{figure}

\subsection{Comparing the two approaches}\label{comparing-the-two-approaches}

\begin{Shaded}
\begin{Highlighting}[]
\NormalTok{comparison }\OtherTok{\textless{}{-}} \FunctionTok{tibble}\NormalTok{(}
  \AttributeTok{Method =} \FunctionTok{c}\NormalTok{(}\StringTok{"Co{-}citation"}\NormalTok{, }\StringTok{"Bibliographic coupling"}\NormalTok{),}
  \AttributeTok{Nodes =} \FunctionTok{c}\NormalTok{(}\FunctionTok{vcount}\NormalTok{(g\_cocit), }\FunctionTok{vcount}\NormalTok{(g\_bibcoup)),}
  \AttributeTok{Edges =} \FunctionTok{c}\NormalTok{(}\FunctionTok{ecount}\NormalTok{(g\_cocit), }\FunctionTok{ecount}\NormalTok{(g\_bibcoup)),}
  \AttributeTok{Components =} \FunctionTok{c}\NormalTok{(}\FunctionTok{components}\NormalTok{(g\_cocit)}\SpecialCharTok{$}\NormalTok{no, }\FunctionTok{components}\NormalTok{(g\_bibcoup)}\SpecialCharTok{$}\NormalTok{no),}
  \AttributeTok{Density =} \FunctionTok{c}\NormalTok{(}\FunctionTok{round}\NormalTok{(}\FunctionTok{graph.density}\NormalTok{(g\_cocit), }\DecValTok{4}\NormalTok{),}
              \FunctionTok{round}\NormalTok{(}\FunctionTok{graph.density}\NormalTok{(g\_bibcoup), }\DecValTok{4}\NormalTok{))}
\NormalTok{)}

\NormalTok{comparison }\SpecialCharTok{|\textgreater{}} \FunctionTok{gt}\NormalTok{()}
\end{Highlighting}
\end{Shaded}

\begingroup
\fontsize{8.0pt}{9.0pt}\selectfont
\begin{longtable}{>{\raggedright\arraybackslash}p{\dimexpr 0.28\linewidth -2\tabcolsep-1.5\arrayrulewidth}>{\raggedleft\arraybackslash}p{\dimexpr 0.17\linewidth -2\tabcolsep-1.5\arrayrulewidth}>{\raggedleft\arraybackslash}p{\dimexpr 0.17\linewidth -2\tabcolsep-1.5\arrayrulewidth}>{\raggedleft\arraybackslash}p{\dimexpr 0.19\linewidth -2\tabcolsep-1.5\arrayrulewidth}>{\raggedleft\arraybackslash}p{\dimexpr 0.17\linewidth -2\tabcolsep-1.5\arrayrulewidth}}
\toprule
Method & Nodes & Edges & \nolinkurl{Components} & Density \\ 
\midrule\addlinespace[2.5pt]
\nolinkurl{Co-citation} & 124 & 160 & 22 & 0.0210 \\ 
Bibliographic coupling & 63 & 60 & 14 & 0.0307 \\ 
\bottomrule
\end{longtable}
\endgroup

\section{Diagnostics and interpretation}\label{cocit-diagnostics}

\begin{itemize}
\tightlist
\item
  \textbf{Threshold selection}: The minimum co-citation count or shared-reference count determines network density. Too low a threshold creates a hairball; too high excludes important links. Experiment with multiple thresholds and report the chosen value.
\item
  \textbf{Normalisation}: Raw co-citation counts favour highly cited papers. Cosine normalisation produces more balanced networks. Always report whether counts are raw or normalised.
\item
  \textbf{Giant component ratio}: A fragmented co-citation network may indicate that the corpus spans multiple disconnected subfields, or that the threshold is too high.
\item
  \textbf{Temporal stability}: Co-citation networks change as new papers cite old ones. Bibliographic coupling networks are fixed at publication time. If you need a stable snapshot of recent literature, coupling is preferable.
\end{itemize}

\section{Limitations and responsible use}\label{cocit-limitations}

\section{Limitations and responsible use}\label{limitations-and-responsible-use-15}

\begin{itemize}
\tightlist
\item
  \textbf{Co-citation reflects citing behaviour, not truth.} Frequently co-cited papers may be co-cited because they represent a dominant paradigm, not because they are the best or most relevant works. Minority viewpoints may be invisible.
\item
  \textbf{Reference list quality.} Both methods depend entirely on the accuracy and completeness of reference lists. Missing references, formatting errors, and DOI mismatches reduce coverage. OpenAlex reference parsing is imperfect \citep{priem2022}.
\item
  \textbf{Size effects.} Papers with long reference lists contribute disproportionately to bibliographic coupling. Consider normalising by reference list length.
\item
  \textbf{Not evaluative.} Co-citation and coupling measure relatedness, not quality. A highly co-cited paper is intellectually central, which is not the same as excellent or correct \citep{hicks2015}.
\end{itemize}

\section{Common pitfalls}\label{cocit-pitfalls}

\section{Common pitfalls}\label{common-pitfalls-15}

\begin{itemize}
\tightlist
\item
  \textbf{Confusing the two methods.} Co-citation links \emph{cited} documents; bibliographic coupling links \emph{citing} documents. They answer different questions about the same citation data.
\item
  \textbf{Skipping normalisation.} Raw counts produce networks dominated by a few highly cited papers. Always normalise for balanced analysis.
\item
  \textbf{Ignoring the threshold.} Reporting results without specifying the minimum count or similarity cutoff makes the analysis irreproducible.
\item
  \textbf{Treating coupling as dynamic.} Bibliographic coupling strength is fixed at publication. Do not interpret changes in coupling networks over time as evolving relationships --- they reflect changes in the corpus, not the links.
\end{itemize}

\section{Exercises}\label{cocit-exercises}

\begin{enumerate}
\def\labelenumi{\arabic{enumi}.}
\item
  \textbf{Cosine normalisation.} Compute cosine-normalised co-citation strengths instead of raw counts. How does the network structure change? Do the same core documents emerge?
\item
  \textbf{Temporal co-citation.} Split the corpus into two time periods (2020--2021 and 2022--2023). Build a co-citation network for each. Which co-citation pairs are stable across periods, and which are new?
\item
  \textbf{Bibliographic coupling by year.} For each publication year, build a bibliographic coupling network and report the number of components and modularity. Does the intellectual structure become more or less fragmented over time?
\item
  \textbf{Identifying research fronts.} Use bibliographic coupling communities to identify research fronts. For each community, list the three most frequent keywords or title words. Do the fronts correspond to recognisable subfields?
\end{enumerate}

\section{Solutions}\label{16-solutions}

Solutions are provided in \ref{solutions}.

\section{Further reading}\label{cocit-reading}

\begin{itemize}
\tightlist
\item
  \citet{small1973} --- Co-citation analysis: the foundational paper.
\item
  \citet{kessler1963} --- Bibliographic coupling: the original proposal.
\item
  \citet{waltman2012} --- Unified approach to bibliometric network construction and clustering.
\item
  \citet{fortunato2010} --- Community detection methods applicable to citation networks.
\item
  \citet{garfield1955} --- The intellectual origins of citation indexing.
\end{itemize}

\section{Session info}\label{cocit-sessioninfo}

\begin{Shaded}
\begin{Highlighting}[]
\FunctionTok{sessionInfo}\NormalTok{()}
\end{Highlighting}
\end{Shaded}

\begin{verbatim}
#> R version 4.6.1 (2026-06-24)
#> Platform: x86_64-pc-linux-gnu
#> Running under: Ubuntu 24.04.5 LTS
#> 
#> Matrix products: default
#> BLAS:   /usr/lib/x86_64-linux-gnu/openblas-pthread/libblas.so.3 
#> LAPACK: /usr/lib/x86_64-linux-gnu/openblas-pthread/libopenblasp-r0.3.26.so;  LAPACK version 3.12.0
#> 
#> locale:
#>  [1] LC_CTYPE=C.UTF-8       LC_NUMERIC=C           LC_TIME=C.UTF-8       
#>  [4] LC_COLLATE=C.UTF-8     LC_MONETARY=C.UTF-8    LC_MESSAGES=C.UTF-8   
#>  [7] LC_PAPER=C.UTF-8       LC_NAME=C              LC_ADDRESS=C          
#> [10] LC_TELEPHONE=C         LC_MEASUREMENT=C.UTF-8 LC_IDENTIFICATION=C   
#> 
#> time zone: UTC
#> tzcode source: system (glibc)
#> 
#> attached base packages:
#> [1] stats     graphics  grDevices utils     datasets  methods   base     
#> 
#> other attached packages:
#>  [1] ggraph_2.2.2       tidygraph_1.3.1    igraph_2.3.4       quanteda_4.5.1    
#>  [5] pdftools_3.9.1     arrow_26.0.0       bibliometrix_5.5.0 RefManageR_1.4.0  
#>  [9] bib2df_1.1.2.0     rcrossref_1.2.1    gt_1.4.0           tidytext_0.4.3    
#> [13] glue_1.8.1         openalexR_3.1.2    lubridate_1.9.5    forcats_1.0.1     
#> [17] stringr_1.6.0      dplyr_1.2.1        purrr_1.2.2        readr_2.2.0       
#> [21] tidyr_1.3.2        tibble_3.3.1       ggplot2_4.0.3      tidyverse_2.0.0   
#> 
#> loaded via a namespace (and not attached):
#>   [1] bibtex_0.5.3           RColorBrewer_1.1-3     rstudioapi_0.19.0     
#>   [4] jsonlite_2.0.0         magrittr_2.0.5         farver_2.1.2          
#>   [7] rmarkdown_2.32         fs_2.1.0               vctrs_0.7.3           
#>  [10] memoise_2.0.1          askpass_1.2.1          base64enc_0.1-6       
#>  [13] htmltools_0.5.9        contentanalysis_1.1.1  curl_8.1.0            
#>  [16] janeaustenr_1.0.0      cellranger_1.1.0       htmlwidgets_1.6.4     
#>  [19] tokenizers_0.3.0       plyr_1.8.9             httr2_1.3.0           
#>  [22] cachem_1.1.0           plotly_4.12.1          dimensionsR_0.0.3     
#>  [25] mime_0.13              lifecycle_1.0.5        pkgconfig_2.0.3       
#>  [28] Matrix_1.7-5           R6_2.6.1               fastmap_1.2.0         
#>  [31] shiny_1.14.0           digest_0.6.39          shinycssloaders_1.1.0 
#>  [34] rprojroot_2.1.1        SnowballC_0.7.1        labeling_0.4.3        
#>  [37] urltools_1.7.3.1       timechange_0.4.0       httr_1.4.9            
#>  [40] polyclip_1.10-7        compiler_4.6.1         here_1.0.2            
#>  [43] bit64_4.8.6            withr_3.0.3            S7_0.2.2              
#>  [46] viridis_0.6.5          ggforce_0.5.0          MASS_7.3-65           
#>  [49] bibliometrixData_0.3.0 tools_4.6.1            otel_0.2.0            
#>  [52] stopwords_2.3          zip_3.0.2              rentrez_1.2.4         
#>  [55] httpuv_1.6.17          promises_1.5.0         grid_4.6.1            
#>  [58] stringdist_0.9.17      generics_0.1.4         gtable_0.3.6          
#>  [61] tzdb_0.5.0             rscopus_0.9.0          ca_0.72               
#>  [64] data.table_1.18.6.1    hms_1.1.4              xml2_1.6.0            
#>  [67] utf8_1.2.6             ggrepel_0.9.8          pillar_1.11.1         
#>  [70] later_1.4.8            tweenr_2.0.3           lattice_0.22-9        
#>  [73] bit_4.6.0              tidyselect_1.2.1       miniUI_0.1.2          
#>  [76] knitr_1.52             gridExtra_2.3.1        bookdown_0.48         
#>  [79] crul_1.6.0             xfun_0.62              graphlayouts_1.2.5    
#>  [82] DT_0.34.0              humaniformat_0.6.0     visNetwork_2.1.4      
#>  [85] stringi_1.8.9          qpdf_1.4.1             yaml_2.3.12           
#>  [88] evaluate_1.0.5         httpcode_0.3.0         cli_3.6.6             
#>  [91] xtable_1.8-8           systemfonts_1.3.2      Rcpp_1.1.2            
#>  [94] readxl_1.6.0           triebeard_0.4.1        XML_3.99-0.25         
#>  [97] parallel_4.6.1         assertthat_0.2.1       pubmedR_1.0.2         
#> [100] viridisLite_0.4.3      scales_1.4.0           openxlsx_4.2.9        
#> [103] rlang_1.3.0            fastmatch_1.1-8
\end{verbatim}

\chapter{Co-word and Keyword Co-occurrence}\label{co-word-and-keyword-co-occurrence}

\emph{R. Heller · Date: 2026-05-09 · Last changed: 2026-10-31}

\section{Learning objectives}\label{coword-objectives}

After completing this chapter, you will be able to:

\begin{itemize}
\tightlist
\item
  Explain the principles of co-word analysis and its role in science mapping
\item
  Extract and clean keywords from OpenAlex data (author keywords and concepts)
\item
  Build a keyword co-occurrence matrix and convert it to a network
\item
  Apply community detection to identify topical clusters
\item
  Visualise keyword co-occurrence networks with interpretable layouts
\end{itemize}

\section{Setup}\label{coword-setup}

\begin{Shaded}
\begin{Highlighting}[]
\FunctionTok{library}\NormalTok{(tidyverse)}
\FunctionTok{library}\NormalTok{(openalexR)}
\FunctionTok{library}\NormalTok{(igraph)}
\FunctionTok{library}\NormalTok{(tidygraph)}
\FunctionTok{library}\NormalTok{(ggraph)}
\FunctionTok{library}\NormalTok{(glue)}
\FunctionTok{library}\NormalTok{(gt)}

\FunctionTok{set.seed}\NormalTok{(}\DecValTok{20260509}\NormalTok{)}

\FunctionTok{source}\NormalTok{(here}\SpecialCharTok{::}\FunctionTok{here}\NormalTok{(}\StringTok{"R"}\NormalTok{, }\StringTok{"api\_helpers.R"}\NormalTok{))}
\FunctionTok{source}\NormalTok{(here}\SpecialCharTok{::}\FunctionTok{here}\NormalTok{(}\StringTok{"R"}\NormalTok{, }\StringTok{"utils.R"}\NormalTok{))}
\FunctionTok{source}\NormalTok{(here}\SpecialCharTok{::}\FunctionTok{here}\NormalTok{(}\StringTok{"R"}\NormalTok{, }\StringTok{"sci\_palette.R"}\NormalTok{))}
\end{Highlighting}
\end{Shaded}

\section{Conceptual background}\label{coword-background}

Co-word analysis maps the conceptual structure of a research field by examining which terms appear together in the same documents. Introduced by \citet{callon1983}, the method assumes that when two keywords repeatedly co-occur across publications, they are thematically related. The resulting co-occurrence network reveals the topical landscape of a field: clusters of densely connected keywords represent coherent research themes, while bridges between clusters indicate interdisciplinary connections.

The method works with several types of terms:

\begin{itemize}
\tightlist
\item
  \textbf{Author keywords}: Terms selected by the paper's authors. These are intentional descriptors but suffer from inconsistency --- authors may use different terms for the same concept (``machine learning'' vs.~``ML'' vs.~``statistical learning'').
\item
  \textbf{Indexed keywords}: Terms assigned by database indexers (e.g., MeSH terms in PubMed). More consistent but only available in some databases.
\item
  \textbf{OpenAlex concepts}: Algorithmically assigned topics at multiple levels of a hierarchical taxonomy. These provide broad coverage but may miss nuanced distinctions \citep{priem2022}.
\item
  \textbf{Title/abstract words}: Extracted via text mining. Comprehensive but noisy; requires extensive preprocessing.
\end{itemize}

Normalisation is important for co-word networks. Raw co-occurrence counts favour high-frequency terms. Common normalisations include the \textbf{association strength} (equivalent to pointwise mutual information) and the \textbf{equivalence index} (cosine of the co-occurrence vector). \citet{waltman2012} demonstrated that association strength produces the most balanced network structures for bibliometric mapping.

Co-word analysis complements citation-based methods (\ref{cocit-background}). Citation networks reveal intellectual influence; co-word networks reveal thematic content. Combining both provides a more complete picture of a field's structure.

\section{Worked example}\label{coword-example}

\subsection{Extracting keywords}\label{extracting-keywords}

We use OpenAlex-assigned topic labels as terms in a sample of scientometrics research. These are algorithmic labels, not author-supplied keywords. The flattened topic table also includes parent subfields, fields, and domains; retain only topic identifiers (\texttt{T...}) and count each term once per work.

\begin{Shaded}
\begin{Highlighting}[]
\NormalTok{works }\OtherTok{\textless{}{-}} \FunctionTok{fetch\_openalex}\NormalTok{(}
  \AttributeTok{entity =} \StringTok{"works"}\NormalTok{,}
  \AttributeTok{primary\_location.source.id =} \StringTok{"S148561398"}\NormalTok{,}
  \AttributeTok{from\_publication\_date =} \StringTok{"2019{-}01{-}01"}\NormalTok{,}
  \AttributeTok{to\_publication\_date =} \StringTok{"2023{-}12{-}31"}\NormalTok{,}
  \AttributeTok{options =} \FunctionTok{list}\NormalTok{(}\AttributeTok{sample =} \DecValTok{500}\NormalTok{, }\AttributeTok{seed =} \DecValTok{42}\NormalTok{)}
\NormalTok{)}
\end{Highlighting}
\end{Shaded}

\begin{Shaded}
\begin{Highlighting}[]
\NormalTok{keywords }\OtherTok{\textless{}{-}}\NormalTok{ works }\SpecialCharTok{|\textgreater{}}
  \FunctionTok{select}\NormalTok{(id, topics) }\SpecialCharTok{|\textgreater{}}
  \FunctionTok{unnest}\NormalTok{(topics, }\AttributeTok{names\_sep =} \StringTok{"\_"}\NormalTok{) }\SpecialCharTok{|\textgreater{}}
  \FunctionTok{filter}\NormalTok{(}\FunctionTok{str\_detect}\NormalTok{(topics\_id, }\StringTok{"(\^{}|/)T[0{-}9]+$"}\NormalTok{)) }\SpecialCharTok{|\textgreater{}}
  \FunctionTok{filter}\NormalTok{(topics\_display\_name }\SpecialCharTok{!=} \StringTok{""}\NormalTok{) }\SpecialCharTok{|\textgreater{}}
  \FunctionTok{select}\NormalTok{(}\AttributeTok{work\_id =}\NormalTok{ id, }\AttributeTok{keyword =}\NormalTok{ topics\_display\_name) }\SpecialCharTok{|\textgreater{}}
  \FunctionTok{mutate}\NormalTok{(}\AttributeTok{keyword =} \FunctionTok{str\_to\_lower}\NormalTok{(}\FunctionTok{str\_trim}\NormalTok{(keyword))) }\SpecialCharTok{|\textgreater{}}
  \FunctionTok{filter}\NormalTok{(}\SpecialCharTok{!}\FunctionTok{is.na}\NormalTok{(keyword), }\FunctionTok{nchar}\NormalTok{(keyword) }\SpecialCharTok{\textgreater{}=} \DecValTok{3}\NormalTok{) }\SpecialCharTok{|\textgreater{}}
  \FunctionTok{distinct}\NormalTok{(work\_id, keyword)}

\NormalTok{kw\_counts }\OtherTok{\textless{}{-}}\NormalTok{ keywords }\SpecialCharTok{|\textgreater{}}
  \FunctionTok{count}\NormalTok{(keyword, }\AttributeTok{sort =} \ConstantTok{TRUE}\NormalTok{)}

\FunctionTok{cat}\NormalTok{(}\FunctionTok{glue}\NormalTok{(}\StringTok{"Total keyword occurrences: \{nrow(keywords)\}}\SpecialCharTok{\textbackslash{}n}\StringTok{"}\NormalTok{))}
\end{Highlighting}
\end{Shaded}

\begin{verbatim}
#> Total keyword occurrences: 1500
\end{verbatim}

\begin{Shaded}
\begin{Highlighting}[]
\FunctionTok{cat}\NormalTok{(}\FunctionTok{glue}\NormalTok{(}\StringTok{"Unique keywords: \{n\_distinct(keywords$keyword)\}}\SpecialCharTok{\textbackslash{}n}\StringTok{"}\NormalTok{))}
\end{Highlighting}
\end{Shaded}

\begin{verbatim}
#> Unique keywords: 136
\end{verbatim}

\begin{Shaded}
\begin{Highlighting}[]
\NormalTok{kw\_counts }\SpecialCharTok{|\textgreater{}}
  \FunctionTok{head}\NormalTok{(}\DecValTok{20}\NormalTok{) }\SpecialCharTok{|\textgreater{}}
  \FunctionTok{mutate}\NormalTok{(}\AttributeTok{keyword =} \FunctionTok{fct\_reorder}\NormalTok{(keyword, n)) }\SpecialCharTok{|\textgreater{}}
  \FunctionTok{ggplot}\NormalTok{(}\FunctionTok{aes}\NormalTok{(}\AttributeTok{x =}\NormalTok{ n, }\AttributeTok{y =}\NormalTok{ keyword)) }\SpecialCharTok{+}
  \FunctionTok{geom\_col}\NormalTok{(}\AttributeTok{fill =} \FunctionTok{palette\_sci}\NormalTok{(}\DecValTok{1}\NormalTok{)) }\SpecialCharTok{+}
  \FunctionTok{labs}\NormalTok{(}\AttributeTok{x =} \StringTok{"Frequency"}\NormalTok{, }\AttributeTok{y =} \ConstantTok{NULL}\NormalTok{) }\SpecialCharTok{+}
  \FunctionTok{theme\_sci}\NormalTok{()}
\end{Highlighting}
\end{Shaded}

\begin{figure}

{\centering \includegraphics[width=0.9\linewidth,alt={Horizontal bar chart showing the 20 most frequent assigned topic labels, counted once per work.}]{scientometrics-in-r_files/figure-latex/17-fig-coword-topkeywords-1} 

}

\caption{Top 20 most frequent assigned topics used as terms in the Scientometrics sample.}\label{fig:17-fig-coword-topkeywords}
\end{figure}

\subsection{Building the co-occurrence network}\label{building-the-co-occurrence-network}

We create edges between keywords that appear in the same paper.

\begin{Shaded}
\begin{Highlighting}[]
\NormalTok{kw\_frequent }\OtherTok{\textless{}{-}}\NormalTok{ kw\_counts }\SpecialCharTok{|\textgreater{}}
  \FunctionTok{filter}\NormalTok{(n }\SpecialCharTok{\textgreater{}=} \DecValTok{5}\NormalTok{) }\SpecialCharTok{|\textgreater{}}
  \FunctionTok{pull}\NormalTok{(keyword)}

\NormalTok{kw\_filtered }\OtherTok{\textless{}{-}}\NormalTok{ keywords }\SpecialCharTok{|\textgreater{}}
  \FunctionTok{filter}\NormalTok{(keyword }\SpecialCharTok{\%in\%}\NormalTok{ kw\_frequent)}

\NormalTok{coword\_pairs }\OtherTok{\textless{}{-}}\NormalTok{ kw\_filtered }\SpecialCharTok{|\textgreater{}}
  \FunctionTok{inner\_join}\NormalTok{(kw\_filtered, }\AttributeTok{by =} \StringTok{"work\_id"}\NormalTok{, }\AttributeTok{suffix =} \FunctionTok{c}\NormalTok{(}\StringTok{"\_a"}\NormalTok{, }\StringTok{"\_b"}\NormalTok{),}
             \AttributeTok{relationship =} \StringTok{"many{-}to{-}many"}\NormalTok{) }\SpecialCharTok{|\textgreater{}}
  \FunctionTok{filter}\NormalTok{(keyword\_a }\SpecialCharTok{\textless{}}\NormalTok{ keyword\_b) }\SpecialCharTok{|\textgreater{}}
  \FunctionTok{count}\NormalTok{(keyword\_a, keyword\_b, }\AttributeTok{name =} \StringTok{"cooccurrence"}\NormalTok{)}

\NormalTok{coword\_top }\OtherTok{\textless{}{-}}\NormalTok{ coword\_pairs }\SpecialCharTok{|\textgreater{}}
  \FunctionTok{filter}\NormalTok{(cooccurrence }\SpecialCharTok{\textgreater{}=} \DecValTok{3}\NormalTok{)}

\NormalTok{g\_kw }\OtherTok{\textless{}{-}} \FunctionTok{graph\_from\_data\_frame}\NormalTok{(}
\NormalTok{  coword\_top }\SpecialCharTok{|\textgreater{}} \FunctionTok{select}\NormalTok{(keyword\_a, keyword\_b, }\AttributeTok{weight =}\NormalTok{ cooccurrence),}
  \AttributeTok{directed =} \ConstantTok{FALSE}
\NormalTok{) }\SpecialCharTok{|\textgreater{}}
  \FunctionTok{simplify}\NormalTok{(}\AttributeTok{edge.attr.comb =} \FunctionTok{list}\NormalTok{(}\AttributeTok{weight =} \StringTok{"sum"}\NormalTok{))}

\FunctionTok{cat}\NormalTok{(}\FunctionTok{glue}\NormalTok{(}\StringTok{"Keyword network: \{vcount(g\_kw)\} nodes, \{ecount(g\_kw)\} edges}\SpecialCharTok{\textbackslash{}n}\StringTok{"}\NormalTok{))}
\end{Highlighting}
\end{Shaded}

\begin{verbatim}
#> Keyword network: 34 nodes, 66 edges
\end{verbatim}

\subsection{Community detection for topical clusters}\label{community-detection-for-topical-clusters}

\begin{Shaded}
\begin{Highlighting}[]
\FunctionTok{V}\NormalTok{(g\_kw)}\SpecialCharTok{$}\NormalTok{degree }\OtherTok{\textless{}{-}} \FunctionTok{degree}\NormalTok{(g\_kw)}
\FunctionTok{V}\NormalTok{(g\_kw)}\SpecialCharTok{$}\NormalTok{strength }\OtherTok{\textless{}{-}} \FunctionTok{strength}\NormalTok{(g\_kw)}

\NormalTok{communities }\OtherTok{\textless{}{-}} \FunctionTok{cluster\_leiden}\NormalTok{(g\_kw, }\AttributeTok{resolution\_parameter =} \FloatTok{0.8}\NormalTok{,}
                              \AttributeTok{objective\_function =} \StringTok{"modularity"}\NormalTok{)}
\FunctionTok{V}\NormalTok{(g\_kw)}\SpecialCharTok{$}\NormalTok{community }\OtherTok{\textless{}{-}} \FunctionTok{as.factor}\NormalTok{(}\FunctionTok{membership}\NormalTok{(communities))}

\FunctionTok{cat}\NormalTok{(}\FunctionTok{glue}\NormalTok{(}\StringTok{"Communities: \{length(unique(membership(communities)))\}}\SpecialCharTok{\textbackslash{}n}\StringTok{"}\NormalTok{))}
\end{Highlighting}
\end{Shaded}

\begin{verbatim}
#> Communities: 5
\end{verbatim}

\begin{Shaded}
\begin{Highlighting}[]
\FunctionTok{cat}\NormalTok{(}\FunctionTok{glue}\NormalTok{(}\StringTok{"Modularity: \{round(modularity(g\_kw, membership(communities)), 3)\}}\SpecialCharTok{\textbackslash{}n}\StringTok{"}\NormalTok{))}
\end{Highlighting}
\end{Shaded}

\begin{verbatim}
#> Modularity: 0.051
\end{verbatim}

\begin{Shaded}
\begin{Highlighting}[]
\NormalTok{community\_summary }\OtherTok{\textless{}{-}} \FunctionTok{tibble}\NormalTok{(}
  \AttributeTok{keyword =} \FunctionTok{V}\NormalTok{(g\_kw)}\SpecialCharTok{$}\NormalTok{name,}
  \AttributeTok{community =} \FunctionTok{V}\NormalTok{(g\_kw)}\SpecialCharTok{$}\NormalTok{community,}
  \AttributeTok{strength =} \FunctionTok{V}\NormalTok{(g\_kw)}\SpecialCharTok{$}\NormalTok{strength}
\NormalTok{) }\SpecialCharTok{|\textgreater{}}
  \FunctionTok{group\_by}\NormalTok{(community) }\SpecialCharTok{|\textgreater{}}
  \FunctionTok{slice\_max}\NormalTok{(strength, }\AttributeTok{n =} \DecValTok{5}\NormalTok{) }\SpecialCharTok{|\textgreater{}}
  \FunctionTok{summarise}\NormalTok{(}\AttributeTok{top\_keywords =} \FunctionTok{paste}\NormalTok{(keyword, }\AttributeTok{collapse =} \StringTok{", "}\NormalTok{),}
            \AttributeTok{n\_keywords =} \FunctionTok{n}\NormalTok{(), }\AttributeTok{.groups =} \StringTok{"drop"}\NormalTok{)}

\NormalTok{community\_summary }\SpecialCharTok{|\textgreater{}} \FunctionTok{gt}\NormalTok{()}
\end{Highlighting}
\end{Shaded}

\begingroup
\fontsize{8.0pt}{9.0pt}\selectfont
\begin{longtable}{>{\centering\arraybackslash}p{\dimexpr 0.22\linewidth -2\tabcolsep-1.5\arrayrulewidth}>{\raggedright\arraybackslash}p{\dimexpr 0.56\linewidth -2\tabcolsep-1.5\arrayrulewidth}>{\raggedleft\arraybackslash}p{\dimexpr 0.23\linewidth -2\tabcolsep-1.5\arrayrulewidth}}
\toprule
\nolinkurl{community} & \nolinkurl{top_keywords} & \nolinkurl{n_keywords} \\ 
\midrule\addlinespace[2.5pt]
1 & scientometrics and bibliometrics research, web visibility and informetrics, psychology research and bibliometrics, academic publishing and open access, publishing and scholarly communication & 5 \\ 
2 & recommender systems and techniques, information retrieval and search behavior & 2 \\ 
3 & computational and text analysis methods, advanced text analysis techniques, topic modeling & 3 \\ 
4 & efficiency analysis using dea, innovation policy and r\&d & 2 \\ 
5 & intellectual property and patents, intellectual property law, firm innovation and growth, innovation and knowledge management & 4 \\ 
\bottomrule
\end{longtable}
\endgroup

\subsection{Visualisation}\label{visualisation-1}

\begin{Shaded}
\begin{Highlighting}[]
\FunctionTok{set.seed}\NormalTok{(}\DecValTok{42}\NormalTok{)}
\NormalTok{layout }\OtherTok{\textless{}{-}} \FunctionTok{create\_layout}\NormalTok{(}\FunctionTok{as\_tbl\_graph}\NormalTok{(g\_kw), }\AttributeTok{layout =} \StringTok{"fr"}\NormalTok{)}

\FunctionTok{ggraph}\NormalTok{(layout) }\SpecialCharTok{+}
  \FunctionTok{geom\_edge\_link}\NormalTok{(}\FunctionTok{aes}\NormalTok{(}\AttributeTok{width =}\NormalTok{ weight), }\AttributeTok{alpha =} \FloatTok{0.15}\NormalTok{, }\AttributeTok{colour =} \StringTok{"grey60"}\NormalTok{) }\SpecialCharTok{+}
  \FunctionTok{scale\_edge\_width\_continuous}\NormalTok{(}\AttributeTok{range =} \FunctionTok{c}\NormalTok{(}\FloatTok{0.3}\NormalTok{, }\DecValTok{2}\NormalTok{), }\AttributeTok{guide =} \StringTok{"none"}\NormalTok{) }\SpecialCharTok{+}
  \FunctionTok{geom\_node\_point}\NormalTok{(}\FunctionTok{aes}\NormalTok{(}\AttributeTok{size =}\NormalTok{ degree, }\AttributeTok{colour =}\NormalTok{ community), }\AttributeTok{alpha =} \FloatTok{0.8}\NormalTok{) }\SpecialCharTok{+}
  \FunctionTok{geom\_node\_text}\NormalTok{(}\FunctionTok{aes}\NormalTok{(}\AttributeTok{label =} \FunctionTok{ifelse}\NormalTok{(degree }\SpecialCharTok{\textgreater{}} \FunctionTok{quantile}\NormalTok{(degree, }\FloatTok{0.85}\NormalTok{),}
\NormalTok{                                     name, }\ConstantTok{NA\_character\_}\NormalTok{)),}
                 \AttributeTok{repel =} \ConstantTok{TRUE}\NormalTok{, }\AttributeTok{size =} \FloatTok{2.5}\NormalTok{, }\AttributeTok{max.overlaps =} \DecValTok{20}\NormalTok{, }\AttributeTok{na.rm =} \ConstantTok{TRUE}\NormalTok{) }\SpecialCharTok{+}
  \FunctionTok{scale\_size\_continuous}\NormalTok{(}\AttributeTok{range =} \FunctionTok{c}\NormalTok{(}\DecValTok{1}\NormalTok{, }\DecValTok{6}\NormalTok{), }\AttributeTok{guide =} \StringTok{"none"}\NormalTok{) }\SpecialCharTok{+}
  \FunctionTok{scale\_colour\_manual}\NormalTok{(}\AttributeTok{values =} \FunctionTok{palette\_sci}\NormalTok{(}
    \FunctionTok{n\_distinct}\NormalTok{(}\FunctionTok{V}\NormalTok{(g\_kw)}\SpecialCharTok{$}\NormalTok{community)}
\NormalTok{  )) }\SpecialCharTok{+}
  \FunctionTok{labs}\NormalTok{(}\AttributeTok{colour =} \StringTok{"Cluster"}\NormalTok{) }\SpecialCharTok{+}
  \FunctionTok{theme\_void}\NormalTok{(}\AttributeTok{base\_family =} \StringTok{"sans"}\NormalTok{, }\AttributeTok{base\_size =} \DecValTok{11}\NormalTok{) }\SpecialCharTok{+}
  \FunctionTok{theme}\NormalTok{(}\AttributeTok{legend.position =} \StringTok{"bottom"}\NormalTok{)}
\end{Highlighting}
\end{Shaded}

\begin{figure}

{\centering \includegraphics[width=0.9\linewidth,alt={Network graph where nodes are keywords and edges connect keywords that frequently co-occur. Node size reflects degree; colours indicate topical communities identified by the Leiden algorithm.}]{scientometrics-in-r_files/figure-latex/17-fig-coword-network-1} 

}

\caption{Keyword co-occurrence network coloured by topical cluster.}\label{fig:17-fig-coword-network}
\end{figure}

\section{Diagnostics and interpretation}\label{coword-diagnostics}

\begin{itemize}
\tightlist
\item
  \textbf{Keyword cleaning}: Inconsistent terminology inflates the vocabulary and creates spurious nodes. Standardise spelling, merge synonyms (e.g., ``h-index'' and ``hirsch index''), and remove generic terms (``research'', ``analysis'') that co-occur with everything but convey no topical information.
\item
  \textbf{Frequency threshold}: Including rare keywords produces large, sparse, unreadable networks. Start with keywords appearing in at least 5 papers and adjust based on corpus size.
\item
  \textbf{Community interpretability}: Each community should correspond to a recognisable research theme. If communities are uninterpretable, the resolution parameter may need adjustment or keywords need further cleaning.
\item
  \textbf{Centrality interpretation}: High-degree keywords are thematically central (used across many contexts). High-betweenness keywords bridge distinct topics and may represent interdisciplinary concepts.
\end{itemize}

\section{Limitations and responsible use}\label{coword-limitations}

\section{Limitations and responsible use}\label{limitations-and-responsible-use-16}

\begin{itemize}
\tightlist
\item
  \textbf{Vocabulary inconsistency.} Author keywords are not controlled vocabulary. The same concept may appear under multiple terms, fragmenting what should be a single node. Merging synonyms requires domain expertise.
\item
  \textbf{Algorithmic concepts.} OpenAlex concepts are assigned by machine learning and may contain errors --- especially at fine-grained levels. Always validate topic assignments by reading sample papers \citep{priem2022}.
\item
  \textbf{Static snapshots.} A co-word network represents a time-averaged view. Emerging topics with few publications may be invisible. Consider building temporal slices to track field evolution.
\item
  \textbf{Not evaluative.} Popular keywords indicate research activity, not research quality or societal impact. Do not equate topical centrality with importance \citep{hicks2015}.
\end{itemize}

\section{Common pitfalls}\label{coword-pitfalls}

\section{Common pitfalls}\label{common-pitfalls-16}

\begin{itemize}
\tightlist
\item
  \textbf{Not cleaning keywords.} Punctuation variants (``co-authorship'' vs.~``coauthorship''), case differences, and trailing whitespace create duplicate nodes. Clean before building the network.
\item
  \textbf{Including stop-keywords.} Generic terms like ``research'', ``study'', ``analysis'' co-occur with everything and dominate the network without conveying topical information. Remove them.
\item
  \textbf{Comparing co-occurrence counts across corpora of different sizes.} Larger corpora produce higher co-occurrence counts mechanically. Normalise or use relative measures.
\item
  \textbf{Over-reading small clusters.} A cluster with two or three keywords may reflect a single paper, not a research theme. Report cluster size alongside content.
\end{itemize}

\section{Exercises}\label{coword-exercises}

\begin{enumerate}
\def\labelenumi{\arabic{enumi}.}
\item
  \textbf{Author keywords vs.~concepts.} Build co-occurrence networks from both author keywords (if available) and OpenAlex concepts for the same corpus. Compare the resulting topical clusters.
\item
  \textbf{Temporal evolution.} Split the corpus into yearly slices. For each year, build a co-word network and identify the top 5 keywords by degree. Which keywords are consistently central, and which emerge over time?
\item
  \textbf{Strategic diagram.} Compute density (internal cohesion) and centrality (external connections) for each keyword cluster. Plot them on a strategic diagram (density vs.~centrality quadrants). Which clusters are core themes and which are peripheral?
\item
  \textbf{Normalisation comparison.} Build two networks from the same data: one with raw co-occurrence counts and one with association strength normalisation. How do the most central keywords differ?
\end{enumerate}

\section{Solutions}\label{17-solutions}

Solutions are provided in \ref{solutions}.

\section{Further reading}\label{coword-reading}

\begin{itemize}
\tightlist
\item
  \citet{callon1983} --- The foundational paper on co-word analysis for mapping science.
\item
  \citet{waltman2012} --- Unified framework for bibliometric network construction, including keyword networks.
\item
  \citet{aria2017} --- \texttt{bibliometrix} includes co-word analysis via the \texttt{biblioNetwork()} function.
\item
  \citet{priem2022} --- OpenAlex concepts and topic classification.
\end{itemize}

\section{Session info}\label{coword-sessioninfo}

\begin{Shaded}
\begin{Highlighting}[]
\FunctionTok{sessionInfo}\NormalTok{()}
\end{Highlighting}
\end{Shaded}

\begin{verbatim}
#> R version 4.6.1 (2026-06-24)
#> Platform: x86_64-pc-linux-gnu
#> Running under: Ubuntu 24.04.5 LTS
#> 
#> Matrix products: default
#> BLAS:   /usr/lib/x86_64-linux-gnu/openblas-pthread/libblas.so.3 
#> LAPACK: /usr/lib/x86_64-linux-gnu/openblas-pthread/libopenblasp-r0.3.26.so;  LAPACK version 3.12.0
#> 
#> locale:
#>  [1] LC_CTYPE=C.UTF-8       LC_NUMERIC=C           LC_TIME=C.UTF-8       
#>  [4] LC_COLLATE=C.UTF-8     LC_MONETARY=C.UTF-8    LC_MESSAGES=C.UTF-8   
#>  [7] LC_PAPER=C.UTF-8       LC_NAME=C              LC_ADDRESS=C          
#> [10] LC_TELEPHONE=C         LC_MEASUREMENT=C.UTF-8 LC_IDENTIFICATION=C   
#> 
#> time zone: UTC
#> tzcode source: system (glibc)
#> 
#> attached base packages:
#> [1] stats     graphics  grDevices utils     datasets  methods   base     
#> 
#> other attached packages:
#>  [1] ggraph_2.2.2       tidygraph_1.3.1    igraph_2.3.4       quanteda_4.5.1    
#>  [5] pdftools_3.9.1     arrow_26.0.0       bibliometrix_5.5.0 RefManageR_1.4.0  
#>  [9] bib2df_1.1.2.0     rcrossref_1.2.1    gt_1.4.0           tidytext_0.4.3    
#> [13] glue_1.8.1         openalexR_3.1.2    lubridate_1.9.5    forcats_1.0.1     
#> [17] stringr_1.6.0      dplyr_1.2.1        purrr_1.2.2        readr_2.2.0       
#> [21] tidyr_1.3.2        tibble_3.3.1       ggplot2_4.0.3      tidyverse_2.0.0   
#> 
#> loaded via a namespace (and not attached):
#>   [1] bibtex_0.5.3           RColorBrewer_1.1-3     rstudioapi_0.19.0     
#>   [4] jsonlite_2.0.0         magrittr_2.0.5         farver_2.1.2          
#>   [7] rmarkdown_2.32         fs_2.1.0               vctrs_0.7.3           
#>  [10] memoise_2.0.1          askpass_1.2.1          base64enc_0.1-6       
#>  [13] htmltools_0.5.9        contentanalysis_1.1.1  curl_8.1.0            
#>  [16] janeaustenr_1.0.0      cellranger_1.1.0       htmlwidgets_1.6.4     
#>  [19] tokenizers_0.3.0       plyr_1.8.9             httr2_1.3.0           
#>  [22] cachem_1.1.0           plotly_4.12.1          dimensionsR_0.0.3     
#>  [25] mime_0.13              lifecycle_1.0.5        pkgconfig_2.0.3       
#>  [28] Matrix_1.7-5           R6_2.6.1               fastmap_1.2.0         
#>  [31] shiny_1.14.0           digest_0.6.39          shinycssloaders_1.1.0 
#>  [34] rprojroot_2.1.1        SnowballC_0.7.1        labeling_0.4.3        
#>  [37] urltools_1.7.3.1       timechange_0.4.0       httr_1.4.9            
#>  [40] polyclip_1.10-7        compiler_4.6.1         here_1.0.2            
#>  [43] bit64_4.8.6            withr_3.0.3            S7_0.2.2              
#>  [46] viridis_0.6.5          ggforce_0.5.0          MASS_7.3-65           
#>  [49] bibliometrixData_0.3.0 tools_4.6.1            otel_0.2.0            
#>  [52] stopwords_2.3          zip_3.0.2              rentrez_1.2.4         
#>  [55] httpuv_1.6.17          promises_1.5.0         grid_4.6.1            
#>  [58] stringdist_0.9.17      generics_0.1.4         gtable_0.3.6          
#>  [61] tzdb_0.5.0             rscopus_0.9.0          ca_0.72               
#>  [64] data.table_1.18.6.1    hms_1.1.4              xml2_1.6.0            
#>  [67] utf8_1.2.6             ggrepel_0.9.8          pillar_1.11.1         
#>  [70] later_1.4.8            tweenr_2.0.3           lattice_0.22-9        
#>  [73] bit_4.6.0              tidyselect_1.2.1       miniUI_0.1.2          
#>  [76] knitr_1.52             gridExtra_2.3.1        bookdown_0.48         
#>  [79] crul_1.6.0             xfun_0.62              graphlayouts_1.2.5    
#>  [82] DT_0.34.0              humaniformat_0.6.0     visNetwork_2.1.4      
#>  [85] stringi_1.8.9          qpdf_1.4.1             yaml_2.3.12           
#>  [88] evaluate_1.0.5         httpcode_0.3.0         cli_3.6.6             
#>  [91] xtable_1.8-8           systemfonts_1.3.2      Rcpp_1.1.2            
#>  [94] readxl_1.6.0           triebeard_0.4.1        XML_3.99-0.25         
#>  [97] parallel_4.6.1         assertthat_0.2.1       pubmedR_1.0.2         
#> [100] viridisLite_0.4.3      scales_1.4.0           openxlsx_4.2.9        
#> [103] rlang_1.3.0            fastmatch_1.1-8
\end{verbatim}

\chapter{Community Detection and Backbone Extraction}\label{community-detection-and-backbone-extraction}

\emph{R. Heller · Date: 2026-05-09 · Last changed: 2026-10-31}

\section{Learning objectives}\label{backbone-objectives}

After completing this chapter, you will be able to:

\begin{itemize}
\tightlist
\item
  Explain why dense networks require backbone extraction before community detection
\item
  Apply the disparity filter to extract statistically significant edges
\item
  Tune the resolution parameter for Louvain and Leiden community detection
\item
  Compare community assignments across methods using Normalised Mutual Information
\item
  Interpret hierarchical community structure at multiple resolutions
\end{itemize}

\section{Setup}\label{backbone-setup}

\begin{Shaded}
\begin{Highlighting}[]
\FunctionTok{library}\NormalTok{(tidyverse)}
\FunctionTok{library}\NormalTok{(openalexR)}
\FunctionTok{library}\NormalTok{(igraph)}
\FunctionTok{library}\NormalTok{(tidygraph)}
\FunctionTok{library}\NormalTok{(ggraph)}
\FunctionTok{library}\NormalTok{(glue)}
\FunctionTok{library}\NormalTok{(gt)}

\FunctionTok{set.seed}\NormalTok{(}\DecValTok{20260509}\NormalTok{)}

\FunctionTok{source}\NormalTok{(here}\SpecialCharTok{::}\FunctionTok{here}\NormalTok{(}\StringTok{"R"}\NormalTok{, }\StringTok{"api\_helpers.R"}\NormalTok{))}
\FunctionTok{source}\NormalTok{(here}\SpecialCharTok{::}\FunctionTok{here}\NormalTok{(}\StringTok{"R"}\NormalTok{, }\StringTok{"utils.R"}\NormalTok{))}
\FunctionTok{source}\NormalTok{(here}\SpecialCharTok{::}\FunctionTok{here}\NormalTok{(}\StringTok{"R"}\NormalTok{, }\StringTok{"sci\_palette.R"}\NormalTok{))}
\end{Highlighting}
\end{Shaded}

\section{Conceptual background}\label{backbone-background}

Bibliometric networks are often dense: in a co-citation or bibliographic coupling network, every pair of documents with any shared citation creates an edge. A corpus of 1,000 papers can easily produce a network with hundreds of thousands of edges, most of which represent weak or coincidental relationships. Applying community detection directly to such networks yields poor results --- the algorithms cannot distinguish signal from noise.

\textbf{Backbone extraction} addresses this by pruning edges that are not statistically significant, retaining only the ``skeleton'' of the network. The \textbf{disparity filter} \citep{serrano2009} tests each edge against a null model based on the local weight distribution of its endpoints. An edge is retained if its weight is unexpectedly high given the node's total weight and degree. This preserves multi-scale structure: both high-weight and low-weight nodes can retain their most important connections.

An alternative is simple \textbf{threshold filtering} (remove edges below a fixed weight), but this systematically removes connections between low-activity nodes, biasing the backbone toward already-prominent actors. The disparity filter avoids this bias by using a local significance test.

Once the backbone is extracted, \textbf{community detection} identifies groups of nodes that are more densely connected internally than externally. The Leiden algorithm \citep{traag2019} improves on Louvain \citep{blondel2008} by guaranteeing that all detected communities are internally connected. Both algorithms accept a \textbf{resolution parameter} that controls community granularity: lower values produce fewer, larger communities; higher values produce more, smaller communities. There is no single ``correct'' resolution --- the appropriate value depends on the research question.

\citet{fortunato2010} provides a comprehensive review of community detection methods. In practice, running the algorithm at multiple resolutions and examining the hierarchical structure provides the richest understanding of network organisation.

\section{Worked example}\label{backbone-example}

\subsection{Building a dense network}\label{building-a-dense-network}

We construct a bibliographic coupling network from scientometrics articles.

\begin{Shaded}
\begin{Highlighting}[]
\NormalTok{works }\OtherTok{\textless{}{-}} \FunctionTok{fetch\_openalex}\NormalTok{(}
  \AttributeTok{entity =} \StringTok{"works"}\NormalTok{,}
  \AttributeTok{primary\_location.source.id =} \StringTok{"S148561398"}\NormalTok{,}
  \AttributeTok{from\_publication\_date =} \StringTok{"2020{-}01{-}01"}\NormalTok{,}
  \AttributeTok{to\_publication\_date =} \StringTok{"2023{-}12{-}31"}\NormalTok{,}
  \AttributeTok{options =} \FunctionTok{list}\NormalTok{(}\AttributeTok{sample =} \DecValTok{300}\NormalTok{, }\AttributeTok{seed =} \DecValTok{42}\NormalTok{)}
\NormalTok{)}

\NormalTok{refs }\OtherTok{\textless{}{-}}\NormalTok{ works }\SpecialCharTok{|\textgreater{}}
  \FunctionTok{select}\NormalTok{(}\AttributeTok{citing\_id =}\NormalTok{ id, referenced\_works) }\SpecialCharTok{|\textgreater{}}
  \FunctionTok{unnest}\NormalTok{(referenced\_works) }\SpecialCharTok{|\textgreater{}}
  \FunctionTok{rename}\NormalTok{(}\AttributeTok{cited\_id =}\NormalTok{ referenced\_works)}

\NormalTok{bibcoup }\OtherTok{\textless{}{-}}\NormalTok{ refs }\SpecialCharTok{|\textgreater{}}
  \FunctionTok{inner\_join}\NormalTok{(refs, }\AttributeTok{by =} \StringTok{"cited\_id"}\NormalTok{, }\AttributeTok{suffix =} \FunctionTok{c}\NormalTok{(}\StringTok{"\_a"}\NormalTok{, }\StringTok{"\_b"}\NormalTok{),}
             \AttributeTok{relationship =} \StringTok{"many{-}to{-}many"}\NormalTok{) }\SpecialCharTok{|\textgreater{}}
  \FunctionTok{filter}\NormalTok{(citing\_id\_a }\SpecialCharTok{\textless{}}\NormalTok{ citing\_id\_b) }\SpecialCharTok{|\textgreater{}}
  \FunctionTok{count}\NormalTok{(citing\_id\_a, citing\_id\_b, }\AttributeTok{name =} \StringTok{"shared\_refs"}\NormalTok{)}

\NormalTok{g\_full }\OtherTok{\textless{}{-}} \FunctionTok{graph\_from\_data\_frame}\NormalTok{(}
\NormalTok{  bibcoup }\SpecialCharTok{|\textgreater{}} \FunctionTok{select}\NormalTok{(citing\_id\_a, citing\_id\_b, }\AttributeTok{weight =}\NormalTok{ shared\_refs),}
  \AttributeTok{directed =} \ConstantTok{FALSE}
\NormalTok{) }\SpecialCharTok{|\textgreater{}}
  \FunctionTok{simplify}\NormalTok{(}\AttributeTok{edge.attr.comb =} \FunctionTok{list}\NormalTok{(}\AttributeTok{weight =} \StringTok{"sum"}\NormalTok{))}

\FunctionTok{cat}\NormalTok{(}\FunctionTok{glue}\NormalTok{(}\StringTok{"Full network: \{vcount(g\_full)\} nodes, \{ecount(g\_full)\} edges}\SpecialCharTok{\textbackslash{}n}\StringTok{"}\NormalTok{))}
\end{Highlighting}
\end{Shaded}

\begin{verbatim}
#> Full network: 290 nodes, 4105 edges
\end{verbatim}

\begin{Shaded}
\begin{Highlighting}[]
\FunctionTok{cat}\NormalTok{(}\FunctionTok{glue}\NormalTok{(}\StringTok{"Density: \{round(graph.density(g\_full), 4)\}}\SpecialCharTok{\textbackslash{}n}\StringTok{"}\NormalTok{))}
\end{Highlighting}
\end{Shaded}

\begin{verbatim}
#> Density: 0.098
\end{verbatim}

\subsection{Backbone extraction: disparity filter}\label{backbone-extraction-disparity-filter}

\begin{Shaded}
\begin{Highlighting}[]
\NormalTok{disparity\_filter }\OtherTok{\textless{}{-}} \ControlFlowTok{function}\NormalTok{(g, }\AttributeTok{alpha =} \FloatTok{0.05}\NormalTok{) \{}
\NormalTok{  el }\OtherTok{\textless{}{-}} \FunctionTok{as\_data\_frame}\NormalTok{(g, }\AttributeTok{what =} \StringTok{"edges"}\NormalTok{)}
\NormalTok{  node\_strength }\OtherTok{\textless{}{-}} \FunctionTok{strength}\NormalTok{(g)}

\NormalTok{  keep }\OtherTok{\textless{}{-}} \FunctionTok{map2\_lgl}\NormalTok{(}\FunctionTok{seq\_len}\NormalTok{(}\FunctionTok{nrow}\NormalTok{(el)), el}\SpecialCharTok{$}\NormalTok{weight, }\ControlFlowTok{function}\NormalTok{(i, w) \{}
\NormalTok{    from\_name }\OtherTok{\textless{}{-}}\NormalTok{ el}\SpecialCharTok{$}\NormalTok{from[i]}
\NormalTok{    to\_name }\OtherTok{\textless{}{-}}\NormalTok{ el}\SpecialCharTok{$}\NormalTok{to[i]}
\NormalTok{    k\_from }\OtherTok{\textless{}{-}} \FunctionTok{degree}\NormalTok{(g, from\_name)}
\NormalTok{    k\_to }\OtherTok{\textless{}{-}} \FunctionTok{degree}\NormalTok{(g, to\_name)}
\NormalTok{    s\_from }\OtherTok{\textless{}{-}}\NormalTok{ node\_strength[from\_name]}
\NormalTok{    s\_to }\OtherTok{\textless{}{-}}\NormalTok{ node\_strength[to\_name]}

\NormalTok{    p\_from }\OtherTok{\textless{}{-}}\NormalTok{ (}\DecValTok{1} \SpecialCharTok{{-}}\NormalTok{ w }\SpecialCharTok{/}\NormalTok{ s\_from)}\SpecialCharTok{\^{}}\NormalTok{(k\_from }\SpecialCharTok{{-}} \DecValTok{1}\NormalTok{)}
\NormalTok{    p\_to }\OtherTok{\textless{}{-}}\NormalTok{ (}\DecValTok{1} \SpecialCharTok{{-}}\NormalTok{ w }\SpecialCharTok{/}\NormalTok{ s\_to)}\SpecialCharTok{\^{}}\NormalTok{(k\_to }\SpecialCharTok{{-}} \DecValTok{1}\NormalTok{)}

    \FunctionTok{min}\NormalTok{(p\_from, p\_to) }\SpecialCharTok{\textless{}}\NormalTok{ alpha}
\NormalTok{  \})}

  \FunctionTok{subgraph.edges}\NormalTok{(g, }\FunctionTok{which}\NormalTok{(keep), }\AttributeTok{delete.vertices =} \ConstantTok{TRUE}\NormalTok{)}
\NormalTok{\}}

\NormalTok{g\_backbone }\OtherTok{\textless{}{-}} \FunctionTok{disparity\_filter}\NormalTok{(g\_full, }\AttributeTok{alpha =} \FloatTok{0.05}\NormalTok{)}
\FunctionTok{cat}\NormalTok{(}\FunctionTok{glue}\NormalTok{(}\StringTok{"Backbone: \{vcount(g\_backbone)\} nodes, \{ecount(g\_backbone)\} edges}\SpecialCharTok{\textbackslash{}n}\StringTok{"}\NormalTok{))}
\end{Highlighting}
\end{Shaded}

\begin{verbatim}
#> Backbone: 131 nodes, 132 edges
\end{verbatim}

\begin{Shaded}
\begin{Highlighting}[]
\FunctionTok{cat}\NormalTok{(}\FunctionTok{glue}\NormalTok{(}\StringTok{"Edge retention: \{scales::percent(ecount(g\_backbone) / ecount(g\_full))\}}\SpecialCharTok{\textbackslash{}n}\StringTok{"}\NormalTok{))}
\end{Highlighting}
\end{Shaded}

\begin{verbatim}
#> Edge retention: 3%
\end{verbatim}

\subsection{Threshold filtering for comparison}\label{threshold-filtering-for-comparison}

\begin{Shaded}
\begin{Highlighting}[]
\NormalTok{threshold }\OtherTok{\textless{}{-}} \FunctionTok{quantile}\NormalTok{(}\FunctionTok{E}\NormalTok{(g\_full)}\SpecialCharTok{$}\NormalTok{weight, }\FloatTok{0.75}\NormalTok{)}
\NormalTok{g\_threshold }\OtherTok{\textless{}{-}} \FunctionTok{subgraph.edges}\NormalTok{(}
\NormalTok{  g\_full,}
  \FunctionTok{which}\NormalTok{(}\FunctionTok{E}\NormalTok{(g\_full)}\SpecialCharTok{$}\NormalTok{weight }\SpecialCharTok{\textgreater{}=}\NormalTok{ threshold),}
  \AttributeTok{delete.vertices =} \ConstantTok{TRUE}
\NormalTok{)}

\FunctionTok{cat}\NormalTok{(}\FunctionTok{glue}\NormalTok{(}\StringTok{"Threshold (75th pctile = \{threshold\}): \{vcount(g\_threshold)\} nodes, \{ecount(g\_threshold)\} edges}\SpecialCharTok{\textbackslash{}n}\StringTok{"}\NormalTok{))}
\end{Highlighting}
\end{Shaded}

\begin{verbatim}
#> Threshold (75th pctile = 2): 252 nodes, 1159 edges
\end{verbatim}

\subsection{Community detection at multiple resolutions}\label{community-detection-at-multiple-resolutions}

\begin{Shaded}
\begin{Highlighting}[]
\NormalTok{resolutions }\OtherTok{\textless{}{-}} \FunctionTok{c}\NormalTok{(}\FloatTok{0.5}\NormalTok{, }\FloatTok{0.8}\NormalTok{, }\FloatTok{1.0}\NormalTok{, }\FloatTok{1.5}\NormalTok{, }\FloatTok{2.0}\NormalTok{)}

\NormalTok{sweep\_results }\OtherTok{\textless{}{-}} \FunctionTok{map\_dfr}\NormalTok{(resolutions, }\ControlFlowTok{function}\NormalTok{(res) \{}
\NormalTok{  comm }\OtherTok{\textless{}{-}} \FunctionTok{cluster\_leiden}\NormalTok{(g\_backbone, }\AttributeTok{resolution\_parameter =}\NormalTok{ res,}
                         \AttributeTok{objective\_function =} \StringTok{"modularity"}\NormalTok{)}
  \FunctionTok{tibble}\NormalTok{(}
    \AttributeTok{resolution =}\NormalTok{ res,}
    \AttributeTok{n\_communities =} \FunctionTok{length}\NormalTok{(}\FunctionTok{unique}\NormalTok{(}\FunctionTok{membership}\NormalTok{(comm))),}
    \AttributeTok{modularity =} \FunctionTok{round}\NormalTok{(}\FunctionTok{modularity}\NormalTok{(g\_backbone, }\FunctionTok{membership}\NormalTok{(comm)), }\DecValTok{3}\NormalTok{),}
    \AttributeTok{max\_size =} \FunctionTok{max}\NormalTok{(}\FunctionTok{table}\NormalTok{(}\FunctionTok{membership}\NormalTok{(comm))),}
    \AttributeTok{min\_size =} \FunctionTok{min}\NormalTok{(}\FunctionTok{table}\NormalTok{(}\FunctionTok{membership}\NormalTok{(comm)))}
\NormalTok{  )}
\NormalTok{\})}

\NormalTok{sweep\_results }\SpecialCharTok{|\textgreater{}} \FunctionTok{gt}\NormalTok{()}
\end{Highlighting}
\end{Shaded}

\begingroup
\fontsize{8.0pt}{9.0pt}\selectfont
\begin{longtable}{>{\raggedleft\arraybackslash}p{\dimexpr 0.20\linewidth -2\tabcolsep-1.5\arrayrulewidth}>{\raggedleft\arraybackslash}p{\dimexpr 0.23\linewidth -2\tabcolsep-1.5\arrayrulewidth}>{\raggedleft\arraybackslash}p{\dimexpr 0.20\linewidth -2\tabcolsep-1.5\arrayrulewidth}>{\raggedleft\arraybackslash}p{\dimexpr 0.18\linewidth -2\tabcolsep-1.5\arrayrulewidth}>{\raggedleft\arraybackslash}p{\dimexpr 0.18\linewidth -2\tabcolsep-1.5\arrayrulewidth}}
\toprule
\nolinkurl{resolution} & \nolinkurl{n_communities} & \nolinkurl{modularity} & \nolinkurl{max_size} & \nolinkurl{min_size} \\ 
\midrule\addlinespace[2.5pt]
0.5 & 30 & 0.870 & 14 & 2 \\ 
0.8 & 30 & 0.870 & 12 & 2 \\ 
1.0 & 30 & 0.870 & 12 & 2 \\ 
1.5 & 30 & 0.870 & 13 & 2 \\ 
2.0 & 32 & 0.865 & 12 & 2 \\ 
\bottomrule
\end{longtable}
\endgroup

\begin{Shaded}
\begin{Highlighting}[]
\FunctionTok{ggplot}\NormalTok{(sweep\_results, }\FunctionTok{aes}\NormalTok{(}\AttributeTok{x =}\NormalTok{ resolution)) }\SpecialCharTok{+}
  \FunctionTok{geom\_line}\NormalTok{(}\FunctionTok{aes}\NormalTok{(}\AttributeTok{y =}\NormalTok{ n\_communities), }\AttributeTok{colour =} \FunctionTok{palette\_sci}\NormalTok{(}\DecValTok{2}\NormalTok{)[}\DecValTok{1}\NormalTok{], }\AttributeTok{linewidth =} \DecValTok{1}\NormalTok{) }\SpecialCharTok{+}
  \FunctionTok{geom\_point}\NormalTok{(}\FunctionTok{aes}\NormalTok{(}\AttributeTok{y =}\NormalTok{ n\_communities), }\AttributeTok{colour =} \FunctionTok{palette\_sci}\NormalTok{(}\DecValTok{2}\NormalTok{)[}\DecValTok{1}\NormalTok{], }\AttributeTok{size =} \DecValTok{3}\NormalTok{) }\SpecialCharTok{+}
  \FunctionTok{geom\_line}\NormalTok{(}\FunctionTok{aes}\NormalTok{(}\AttributeTok{y =}\NormalTok{ modularity }\SpecialCharTok{*} \FunctionTok{max}\NormalTok{(n\_communities)),}
            \AttributeTok{colour =} \FunctionTok{palette\_sci}\NormalTok{(}\DecValTok{2}\NormalTok{)[}\DecValTok{2}\NormalTok{], }\AttributeTok{linewidth =} \DecValTok{1}\NormalTok{, }\AttributeTok{linetype =} \StringTok{"dashed"}\NormalTok{) }\SpecialCharTok{+}
  \FunctionTok{scale\_y\_continuous}\NormalTok{(}
    \AttributeTok{name =} \StringTok{"Number of communities"}\NormalTok{,}
    \AttributeTok{sec.axis =} \FunctionTok{sec\_axis}\NormalTok{(}\SpecialCharTok{\textasciitilde{}}\NormalTok{ . }\SpecialCharTok{/} \FunctionTok{max}\NormalTok{(sweep\_results}\SpecialCharTok{$}\NormalTok{n\_communities),}
                        \AttributeTok{name =} \StringTok{"Modularity"}\NormalTok{)}
\NormalTok{  ) }\SpecialCharTok{+}
  \FunctionTok{labs}\NormalTok{(}\AttributeTok{x =} \StringTok{"Resolution parameter"}\NormalTok{) }\SpecialCharTok{+}
  \FunctionTok{theme\_sci}\NormalTok{()}
\end{Highlighting}
\end{Shaded}

\begin{figure}

{\centering \includegraphics[width=0.9\linewidth,alt={Dual-axis line chart showing how the number of detected communities increases and modularity changes as the Leiden resolution parameter increases.}]{scientometrics-in-r_files/figure-latex/18-fig-backbone-resolution-1} 

}

\caption{Number of communities and modularity as a function of resolution parameter.}\label{fig:18-fig-backbone-resolution}
\end{figure}

\subsection{Comparing backbone vs.~full network communities}\label{comparing-backbone-vs.-full-network-communities}

\begin{Shaded}
\begin{Highlighting}[]
\NormalTok{comm\_full }\OtherTok{\textless{}{-}} \FunctionTok{cluster\_leiden}\NormalTok{(g\_full, }\AttributeTok{resolution\_parameter =} \FloatTok{1.0}\NormalTok{,}
                            \AttributeTok{objective\_function =} \StringTok{"modularity"}\NormalTok{)}
\NormalTok{comm\_backbone }\OtherTok{\textless{}{-}} \FunctionTok{cluster\_leiden}\NormalTok{(g\_backbone, }\AttributeTok{resolution\_parameter =} \FloatTok{1.0}\NormalTok{,}
                                \AttributeTok{objective\_function =} \StringTok{"modularity"}\NormalTok{)}

\NormalTok{shared\_nodes }\OtherTok{\textless{}{-}} \FunctionTok{intersect}\NormalTok{(}\FunctionTok{V}\NormalTok{(g\_full)}\SpecialCharTok{$}\NormalTok{name, }\FunctionTok{V}\NormalTok{(g\_backbone)}\SpecialCharTok{$}\NormalTok{name)}
\NormalTok{mem\_full }\OtherTok{\textless{}{-}} \FunctionTok{membership}\NormalTok{(comm\_full)[shared\_nodes]}
\NormalTok{mem\_back }\OtherTok{\textless{}{-}} \FunctionTok{membership}\NormalTok{(comm\_backbone)[shared\_nodes]}

\NormalTok{nmi }\OtherTok{\textless{}{-}} \FunctionTok{compare}\NormalTok{(mem\_full, mem\_back, }\AttributeTok{method =} \StringTok{"nmi"}\NormalTok{)}
\FunctionTok{cat}\NormalTok{(}\FunctionTok{glue}\NormalTok{(}\StringTok{"NMI (full vs backbone, shared nodes): \{round(nmi, 3)\}}\SpecialCharTok{\textbackslash{}n}\StringTok{"}\NormalTok{))}
\end{Highlighting}
\end{Shaded}

\begin{verbatim}
#> NMI (full vs backbone, shared nodes): 0.663
\end{verbatim}

\begin{Shaded}
\begin{Highlighting}[]
\FunctionTok{cat}\NormalTok{(}\FunctionTok{glue}\NormalTok{(}\StringTok{"Communities in full network: \{length(unique(membership(comm\_full)))\}}\SpecialCharTok{\textbackslash{}n}\StringTok{"}\NormalTok{))}
\end{Highlighting}
\end{Shaded}

\begin{verbatim}
#> Communities in full network: 9
\end{verbatim}

\begin{Shaded}
\begin{Highlighting}[]
\FunctionTok{cat}\NormalTok{(}\FunctionTok{glue}\NormalTok{(}\StringTok{"Communities in backbone: \{length(unique(membership(comm\_backbone)))\}}\SpecialCharTok{\textbackslash{}n}\StringTok{"}\NormalTok{))}
\end{Highlighting}
\end{Shaded}

\begin{verbatim}
#> Communities in backbone: 30
\end{verbatim}

\subsection{Visualisation}\label{visualisation-2}

\begin{Shaded}
\begin{Highlighting}[]
\FunctionTok{V}\NormalTok{(g\_backbone)}\SpecialCharTok{$}\NormalTok{community }\OtherTok{\textless{}{-}} \FunctionTok{as.factor}\NormalTok{(}\FunctionTok{membership}\NormalTok{(comm\_backbone))}
\FunctionTok{V}\NormalTok{(g\_backbone)}\SpecialCharTok{$}\NormalTok{degree }\OtherTok{\textless{}{-}} \FunctionTok{degree}\NormalTok{(g\_backbone)}

\FunctionTok{set.seed}\NormalTok{(}\DecValTok{42}\NormalTok{)}
\NormalTok{layout }\OtherTok{\textless{}{-}} \FunctionTok{create\_layout}\NormalTok{(}\FunctionTok{as\_tbl\_graph}\NormalTok{(g\_backbone), }\AttributeTok{layout =} \StringTok{"fr"}\NormalTok{)}

\FunctionTok{ggraph}\NormalTok{(layout) }\SpecialCharTok{+}
  \FunctionTok{geom\_edge\_link}\NormalTok{(}\AttributeTok{alpha =} \FloatTok{0.1}\NormalTok{, }\AttributeTok{colour =} \StringTok{"grey60"}\NormalTok{) }\SpecialCharTok{+}
  \FunctionTok{geom\_node\_point}\NormalTok{(}\FunctionTok{aes}\NormalTok{(}\AttributeTok{size =}\NormalTok{ degree, }\AttributeTok{colour =}\NormalTok{ community), }\AttributeTok{alpha =} \FloatTok{0.8}\NormalTok{) }\SpecialCharTok{+}
  \FunctionTok{scale\_size\_continuous}\NormalTok{(}\AttributeTok{range =} \FunctionTok{c}\NormalTok{(}\DecValTok{1}\NormalTok{, }\DecValTok{5}\NormalTok{), }\AttributeTok{guide =} \StringTok{"none"}\NormalTok{) }\SpecialCharTok{+}
  \FunctionTok{scale\_colour\_manual}\NormalTok{(}\AttributeTok{values =} \FunctionTok{palette\_sci}\NormalTok{(}
    \FunctionTok{n\_distinct}\NormalTok{(}\FunctionTok{V}\NormalTok{(g\_backbone)}\SpecialCharTok{$}\NormalTok{community)}
\NormalTok{  )) }\SpecialCharTok{+}
  \FunctionTok{labs}\NormalTok{(}\AttributeTok{colour =} \StringTok{"Community"}\NormalTok{) }\SpecialCharTok{+}
  \FunctionTok{theme\_void}\NormalTok{(}\AttributeTok{base\_family =} \StringTok{"sans"}\NormalTok{, }\AttributeTok{base\_size =} \DecValTok{11}\NormalTok{) }\SpecialCharTok{+}
  \FunctionTok{theme}\NormalTok{(}\AttributeTok{legend.position =} \StringTok{"bottom"}\NormalTok{)}
\end{Highlighting}
\end{Shaded}

\begin{figure}

{\centering \includegraphics[width=0.9\linewidth,alt={Network graph showing the backbone of the bibliographic coupling network. Nodes coloured by community, sized by degree. Only statistically significant edges are shown.}]{scientometrics-in-r_files/figure-latex/18-fig-backbone-network-1} 

}

\caption{Backbone network with Leiden communities (resolution = 1.0).}\label{fig:18-fig-backbone-network}
\end{figure}

\section{Diagnostics and interpretation}\label{backbone-diagnostics}

\begin{itemize}
\tightlist
\item
  \textbf{Edge retention rate}: The disparity filter typically retains 10--30\% of edges. Retention above 50\% suggests the network may not be dense enough to require backbone extraction.
\item
  \textbf{Isolated nodes}: Backbone extraction removes nodes that lose all their edges. Report how many nodes are lost and whether they represent a biased subset (e.g., low-cited papers).
\item
  \textbf{Resolution plateau}: If modularity remains stable across a range of resolutions, the community structure is robust. Rapid changes suggest sensitivity to the parameter.
\item
  \textbf{Singleton communities}: Communities with one or two nodes are usually noise. Consider merging them into the nearest larger community or excluding them from interpretation.
\end{itemize}

\section{Limitations and responsible use}\label{backbone-limitations}

\section{Limitations and responsible use}\label{limitations-and-responsible-use-17}

\begin{itemize}
\tightlist
\item
  \textbf{The disparity filter has assumptions.} It assumes a uniform null distribution of edge weights across a node's connections. Highly skewed weight distributions may violate this assumption.
\item
  \textbf{Backbone choice affects conclusions.} Different backbone methods (disparity filter, noise-corrected, threshold) retain different edges and can produce different community structures. Report the method and alpha level.
\item
  \textbf{Resolution is a researcher decision.} There is no objectively ``correct'' resolution parameter. The choice determines the granularity of analysis and should be justified by the research question, not by optimising modularity alone \citep{fortunato2010}.
\item
  \textbf{Communities are not ground truth.} Detected communities are algorithmic constructs. They reflect structural patterns in citation data, not necessarily real-world research groups or coherent intellectual traditions \citep{hicks2015}.
\end{itemize}

\section{Common pitfalls}\label{backbone-pitfalls}

\section{Common pitfalls}\label{common-pitfalls-17}

\begin{itemize}
\tightlist
\item
  \textbf{Applying community detection to unfiltered dense networks.} The result is usually a single giant community plus isolated fragments. Always extract the backbone first.
\item
  \textbf{Using a single resolution without justification.} The default resolution of 1.0 is arbitrary. Run a sweep and show the sensitivity analysis.
\item
  \textbf{Comparing modularity across networks of different sizes.} Modularity values are not comparable between different networks. Use NMI to compare partitions.
\item
  \textbf{Interpreting backbone edges as ``strong'' relationships.} The disparity filter retains \emph{locally significant} edges, not necessarily those with the highest absolute weight. A low-weight edge can be retained if it is important relative to its endpoint's other connections.
\end{itemize}

\section{Exercises}\label{backbone-exercises}

\begin{enumerate}
\def\labelenumi{\arabic{enumi}.}
\item
  \textbf{Alpha sensitivity.} Run the disparity filter with alpha values of 0.01, 0.05, 0.10, and 0.20. Plot the number of retained edges and communities against alpha. At what alpha does the network become disconnected?
\item
  \textbf{Threshold vs.~disparity.} Compare threshold-filtered and disparity-filtered communities for the same network. Use NMI to quantify agreement. Which method produces more interpretable communities?
\item
  \textbf{Hierarchical structure.} Run Leiden at resolutions 0.5, 1.0, and 2.0. For each pair of resolutions, check whether the finer communities are nested within the coarser ones (a Sankey or alluvial diagram can help visualise this).
\item
  \textbf{Weighted vs.~unweighted community detection.} Remove edge weights from the backbone and run Leiden again. How much does the community structure change?
\end{enumerate}

\section{Solutions}\label{18-solutions}

Solutions are provided in \ref{solutions}.

\section{Further reading}\label{backbone-reading}

\begin{itemize}
\tightlist
\item
  \citet{serrano2009} --- The disparity filter for extracting multiscale network backbones.
\item
  \citet{traag2019} --- The Leiden algorithm with resolution parameter tuning.
\item
  \citet{blondel2008} --- The Louvain algorithm for fast community detection.
\item
  \citet{fortunato2010} --- Comprehensive review of community detection, including resolution limits.
\item
  \citet{waltman2012} --- Network construction and clustering for bibliometric applications.
\end{itemize}

\section{Session info}\label{backbone-sessioninfo}

\begin{Shaded}
\begin{Highlighting}[]
\FunctionTok{sessionInfo}\NormalTok{()}
\end{Highlighting}
\end{Shaded}

\begin{verbatim}
#> R version 4.6.1 (2026-06-24)
#> Platform: x86_64-pc-linux-gnu
#> Running under: Ubuntu 24.04.5 LTS
#> 
#> Matrix products: default
#> BLAS:   /usr/lib/x86_64-linux-gnu/openblas-pthread/libblas.so.3 
#> LAPACK: /usr/lib/x86_64-linux-gnu/openblas-pthread/libopenblasp-r0.3.26.so;  LAPACK version 3.12.0
#> 
#> locale:
#>  [1] LC_CTYPE=C.UTF-8       LC_NUMERIC=C           LC_TIME=C.UTF-8       
#>  [4] LC_COLLATE=C.UTF-8     LC_MONETARY=C.UTF-8    LC_MESSAGES=C.UTF-8   
#>  [7] LC_PAPER=C.UTF-8       LC_NAME=C              LC_ADDRESS=C          
#> [10] LC_TELEPHONE=C         LC_MEASUREMENT=C.UTF-8 LC_IDENTIFICATION=C   
#> 
#> time zone: UTC
#> tzcode source: system (glibc)
#> 
#> attached base packages:
#> [1] stats     graphics  grDevices utils     datasets  methods   base     
#> 
#> other attached packages:
#>  [1] ggraph_2.2.2       tidygraph_1.3.1    igraph_2.3.4       quanteda_4.5.1    
#>  [5] pdftools_3.9.1     arrow_26.0.0       bibliometrix_5.5.0 RefManageR_1.4.0  
#>  [9] bib2df_1.1.2.0     rcrossref_1.2.1    gt_1.4.0           tidytext_0.4.3    
#> [13] glue_1.8.1         openalexR_3.1.2    lubridate_1.9.5    forcats_1.0.1     
#> [17] stringr_1.6.0      dplyr_1.2.1        purrr_1.2.2        readr_2.2.0       
#> [21] tidyr_1.3.2        tibble_3.3.1       ggplot2_4.0.3      tidyverse_2.0.0   
#> 
#> loaded via a namespace (and not attached):
#>   [1] bibtex_0.5.3           RColorBrewer_1.1-3     rstudioapi_0.19.0     
#>   [4] jsonlite_2.0.0         magrittr_2.0.5         farver_2.1.2          
#>   [7] rmarkdown_2.32         fs_2.1.0               vctrs_0.7.3           
#>  [10] memoise_2.0.1          askpass_1.2.1          base64enc_0.1-6       
#>  [13] htmltools_0.5.9        contentanalysis_1.1.1  curl_8.1.0            
#>  [16] janeaustenr_1.0.0      cellranger_1.1.0       htmlwidgets_1.6.4     
#>  [19] tokenizers_0.3.0       plyr_1.8.9             httr2_1.3.0           
#>  [22] cachem_1.1.0           plotly_4.12.1          dimensionsR_0.0.3     
#>  [25] mime_0.13              lifecycle_1.0.5        pkgconfig_2.0.3       
#>  [28] Matrix_1.7-5           R6_2.6.1               fastmap_1.2.0         
#>  [31] shiny_1.14.0           digest_0.6.39          shinycssloaders_1.1.0 
#>  [34] rprojroot_2.1.1        SnowballC_0.7.1        labeling_0.4.3        
#>  [37] urltools_1.7.3.1       timechange_0.4.0       httr_1.4.9            
#>  [40] polyclip_1.10-7        compiler_4.6.1         here_1.0.2            
#>  [43] bit64_4.8.6            withr_3.0.3            S7_0.2.2              
#>  [46] viridis_0.6.5          ggforce_0.5.0          MASS_7.3-65           
#>  [49] bibliometrixData_0.3.0 tools_4.6.1            otel_0.2.0            
#>  [52] stopwords_2.3          zip_3.0.2              rentrez_1.2.4         
#>  [55] httpuv_1.6.17          promises_1.5.0         grid_4.6.1            
#>  [58] stringdist_0.9.17      generics_0.1.4         gtable_0.3.6          
#>  [61] tzdb_0.5.0             rscopus_0.9.0          ca_0.72               
#>  [64] data.table_1.18.6.1    hms_1.1.4              xml2_1.6.0            
#>  [67] utf8_1.2.6             ggrepel_0.9.8          pillar_1.11.1         
#>  [70] later_1.4.8            tweenr_2.0.3           lattice_0.22-9        
#>  [73] bit_4.6.0              tidyselect_1.2.1       miniUI_0.1.2          
#>  [76] knitr_1.52             gridExtra_2.3.1        bookdown_0.48         
#>  [79] crul_1.6.0             xfun_0.62              graphlayouts_1.2.5    
#>  [82] DT_0.34.0              humaniformat_0.6.0     visNetwork_2.1.4      
#>  [85] stringi_1.8.9          qpdf_1.4.1             yaml_2.3.12           
#>  [88] evaluate_1.0.5         httpcode_0.3.0         cli_3.6.6             
#>  [91] xtable_1.8-8           systemfonts_1.3.2      Rcpp_1.1.2            
#>  [94] readxl_1.6.0           triebeard_0.4.1        XML_3.99-0.25         
#>  [97] parallel_4.6.1         assertthat_0.2.1       pubmedR_1.0.2         
#> [100] viridisLite_0.4.3      scales_1.4.0           openxlsx_4.2.9        
#> [103] rlang_1.3.0            fastmatch_1.1-8
\end{verbatim}

\chapter{Science Mapping and Overlay Maps}\label{science-mapping-and-overlay-maps}

\emph{R. Heller · Date: 2026-05-09 · Last changed: 2026-10-31}

\section{Learning objectives}\label{scimap-objectives}

After completing this chapter, you will be able to:

\begin{itemize}
\tightlist
\item
  Explain the principles of science mapping and its applications
\item
  Construct a journal-level co-citation network from OpenAlex data
\item
  Create an overlay map that positions a research profile on a base map
\item
  Interpret overlay maps to identify research strengths and gaps
\item
  Discuss the limitations of global science maps for local evaluation
\end{itemize}

\section{Setup}\label{scimap-setup}

\begin{Shaded}
\begin{Highlighting}[]
\FunctionTok{library}\NormalTok{(tidyverse)}
\FunctionTok{library}\NormalTok{(openalexR)}
\FunctionTok{library}\NormalTok{(igraph)}
\FunctionTok{library}\NormalTok{(tidygraph)}
\FunctionTok{library}\NormalTok{(ggraph)}
\FunctionTok{library}\NormalTok{(glue)}
\FunctionTok{library}\NormalTok{(gt)}

\FunctionTok{set.seed}\NormalTok{(}\DecValTok{20260509}\NormalTok{)}

\FunctionTok{source}\NormalTok{(here}\SpecialCharTok{::}\FunctionTok{here}\NormalTok{(}\StringTok{"R"}\NormalTok{, }\StringTok{"api\_helpers.R"}\NormalTok{))}
\FunctionTok{source}\NormalTok{(here}\SpecialCharTok{::}\FunctionTok{here}\NormalTok{(}\StringTok{"R"}\NormalTok{, }\StringTok{"utils.R"}\NormalTok{))}
\FunctionTok{source}\NormalTok{(here}\SpecialCharTok{::}\FunctionTok{here}\NormalTok{(}\StringTok{"R"}\NormalTok{, }\StringTok{"sci\_palette.R"}\NormalTok{))}
\end{Highlighting}
\end{Shaded}

\section{Conceptual background}\label{scimap-background}

A \textbf{science map} is a spatial representation of the structure of science, where proximity indicates relatedness. Just as a geographic map shows the layout of countries and continents, a science map shows the layout of disciplines, fields, and research areas. These maps are constructed from bibliometric data --- typically citation links between journals or document sets --- and visualised using network layouts or dimensionality reduction techniques.

The most common base maps use \textbf{journal-level citation networks}. Two journals are linked if papers in one frequently cite papers in the other. Community detection applied to these networks produces clusters corresponding to broad disciplines (biomedical sciences, physical sciences, social sciences, etc.). \citet{waltman2012} developed methods for constructing such maps from large-scale citation data, and the resulting visualisations are widely used in science policy and institutional assessment.

An \textbf{overlay map} \citep{rafols2010} takes a global science map as a fixed background and plots a specific entity's output on top. For example, a university's publications are mapped to the journals in which they appeared; the overlay then shows which areas of the global map the university is active in, and how intensely. This technique is powerful for visualising research profiles, identifying strategic gaps, and comparing institutions or countries without reducing their activity to a single number.

The construction pipeline involves three steps: (1) build or obtain a base map of journals with coordinates and cluster assignments; (2) map the entity's publications to journals in the base map; (3) scale node sizes proportionally and overlay on the base map.

Science maps provide intuitive, holistic overviews of research landscapes. However, they are projections of high-dimensional data into two dimensions, and like any projection, they introduce distortion. Interdisciplinary journals may appear in unexpected locations, and the layout algorithm's random seed affects exact positions \citep{hicks2015}.

\section{Worked example}\label{scimap-example}

\subsection{Building a journal co-citation base map}\label{building-a-journal-co-citation-base-map}

We construct a base map from citation flows between journals in a multidisciplinary sample.

\begin{Shaded}
\begin{Highlighting}[]
\NormalTok{works\_multi }\OtherTok{\textless{}{-}} \FunctionTok{fetch\_openalex}\NormalTok{(}
  \AttributeTok{entity =} \StringTok{"works"}\NormalTok{,}
  \AttributeTok{from\_publication\_date =} \StringTok{"2022{-}01{-}01"}\NormalTok{,}
  \AttributeTok{to\_publication\_date =} \StringTok{"2022{-}12{-}31"}\NormalTok{,}
  \AttributeTok{type =} \StringTok{"article"}\NormalTok{,}
  \AttributeTok{options =} \FunctionTok{list}\NormalTok{(}\AttributeTok{sample =} \DecValTok{500}\NormalTok{, }\AttributeTok{seed =} \DecValTok{42}\NormalTok{)}
\NormalTok{)}

\NormalTok{journal\_data }\OtherTok{\textless{}{-}}\NormalTok{ works\_multi }\SpecialCharTok{|\textgreater{}}
  \FunctionTok{select}\NormalTok{(id, source\_display\_name, referenced\_works) }\SpecialCharTok{|\textgreater{}}
  \FunctionTok{filter}\NormalTok{(}\SpecialCharTok{!}\FunctionTok{is.na}\NormalTok{(source\_display\_name)) }\SpecialCharTok{|\textgreater{}}
  \FunctionTok{unnest}\NormalTok{(referenced\_works) }\SpecialCharTok{|\textgreater{}}
  \FunctionTok{rename}\NormalTok{(}\AttributeTok{cited\_id =}\NormalTok{ referenced\_works)}

\FunctionTok{cat}\NormalTok{(}\FunctionTok{glue}\NormalTok{(}\StringTok{"Citing papers with journal info: \{n\_distinct(journal\_data$id)\}}\SpecialCharTok{\textbackslash{}n}\StringTok{"}\NormalTok{))}
\end{Highlighting}
\end{Shaded}

\begin{verbatim}
#> Citing papers with journal info: 499
\end{verbatim}

\begin{Shaded}
\begin{Highlighting}[]
\FunctionTok{cat}\NormalTok{(}\FunctionTok{glue}\NormalTok{(}\StringTok{"Unique source journals: \{n\_distinct(journal\_data$source\_display\_name)\}}\SpecialCharTok{\textbackslash{}n}\StringTok{"}\NormalTok{))}
\end{Highlighting}
\end{Shaded}

\begin{verbatim}
#> Unique source journals: 470
\end{verbatim}

\begin{Shaded}
\begin{Highlighting}[]
\NormalTok{journal\_cited }\OtherTok{\textless{}{-}}\NormalTok{ journal\_data }\SpecialCharTok{|\textgreater{}}
  \FunctionTok{select}\NormalTok{(}\AttributeTok{citing\_journal =}\NormalTok{ source\_display\_name, cited\_id)}

\NormalTok{journal\_pairs }\OtherTok{\textless{}{-}}\NormalTok{ journal\_cited }\SpecialCharTok{|\textgreater{}}
  \FunctionTok{inner\_join}\NormalTok{(journal\_cited, }\AttributeTok{by =} \StringTok{"cited\_id"}\NormalTok{, }\AttributeTok{suffix =} \FunctionTok{c}\NormalTok{(}\StringTok{"\_a"}\NormalTok{, }\StringTok{"\_b"}\NormalTok{),}
             \AttributeTok{relationship =} \StringTok{"many{-}to{-}many"}\NormalTok{) }\SpecialCharTok{|\textgreater{}}
  \FunctionTok{filter}\NormalTok{(citing\_journal\_a }\SpecialCharTok{\textless{}}\NormalTok{ citing\_journal\_b) }\SpecialCharTok{|\textgreater{}}
  \FunctionTok{count}\NormalTok{(citing\_journal\_a, citing\_journal\_b, }\AttributeTok{name =} \StringTok{"cocitation"}\NormalTok{)}

\NormalTok{journal\_pairs\_top }\OtherTok{\textless{}{-}}\NormalTok{ journal\_pairs }\SpecialCharTok{|\textgreater{}}
  \FunctionTok{filter}\NormalTok{(cocitation }\SpecialCharTok{\textgreater{}=} \DecValTok{3}\NormalTok{)}

\NormalTok{g\_journals }\OtherTok{\textless{}{-}} \FunctionTok{graph\_from\_data\_frame}\NormalTok{(}
\NormalTok{  journal\_pairs\_top }\SpecialCharTok{|\textgreater{}} \FunctionTok{select}\NormalTok{(citing\_journal\_a, citing\_journal\_b,}
                               \AttributeTok{weight =}\NormalTok{ cocitation),}
  \AttributeTok{directed =} \ConstantTok{FALSE}
\NormalTok{) }\SpecialCharTok{|\textgreater{}}
  \FunctionTok{simplify}\NormalTok{(}\AttributeTok{edge.attr.comb =} \FunctionTok{list}\NormalTok{(}\AttributeTok{weight =} \StringTok{"sum"}\NormalTok{))}

\FunctionTok{cat}\NormalTok{(}\FunctionTok{glue}\NormalTok{(}\StringTok{"Journal base map: \{vcount(g\_journals)\} journals, \{ecount(g\_journals)\} edges}\SpecialCharTok{\textbackslash{}n}\StringTok{"}\NormalTok{))}
\end{Highlighting}
\end{Shaded}

\begin{verbatim}
#> Journal base map: 148 journals, 592 edges
\end{verbatim}

\subsection{Clustering journals into disciplines}\label{clustering-journals-into-disciplines}

\begin{Shaded}
\begin{Highlighting}[]
\NormalTok{comm }\OtherTok{\textless{}{-}} \FunctionTok{cluster\_leiden}\NormalTok{(g\_journals, }\AttributeTok{resolution\_parameter =} \FloatTok{0.8}\NormalTok{,}
                       \AttributeTok{objective\_function =} \StringTok{"modularity"}\NormalTok{)}
\FunctionTok{V}\NormalTok{(g\_journals)}\SpecialCharTok{$}\NormalTok{cluster }\OtherTok{\textless{}{-}} \FunctionTok{as.factor}\NormalTok{(}\FunctionTok{membership}\NormalTok{(comm))}
\FunctionTok{V}\NormalTok{(g\_journals)}\SpecialCharTok{$}\NormalTok{degree }\OtherTok{\textless{}{-}} \FunctionTok{degree}\NormalTok{(g\_journals)}

\NormalTok{cluster\_table }\OtherTok{\textless{}{-}} \FunctionTok{tibble}\NormalTok{(}
  \AttributeTok{journal =} \FunctionTok{V}\NormalTok{(g\_journals)}\SpecialCharTok{$}\NormalTok{name,}
  \AttributeTok{cluster =} \FunctionTok{V}\NormalTok{(g\_journals)}\SpecialCharTok{$}\NormalTok{cluster,}
  \AttributeTok{degree =} \FunctionTok{V}\NormalTok{(g\_journals)}\SpecialCharTok{$}\NormalTok{degree}
\NormalTok{) }\SpecialCharTok{|\textgreater{}}
  \FunctionTok{group\_by}\NormalTok{(cluster) }\SpecialCharTok{|\textgreater{}}
  \FunctionTok{summarise}\NormalTok{(}\AttributeTok{n\_journals =} \FunctionTok{n}\NormalTok{(),}
            \AttributeTok{top\_journals =} \FunctionTok{paste}\NormalTok{(}\FunctionTok{head}\NormalTok{(journal[}\FunctionTok{order}\NormalTok{(}\SpecialCharTok{{-}}\NormalTok{degree)], }\DecValTok{3}\NormalTok{),}
                                 \AttributeTok{collapse =} \StringTok{"; "}\NormalTok{),}
            \AttributeTok{.groups =} \StringTok{"drop"}\NormalTok{) }\SpecialCharTok{|\textgreater{}}
  \FunctionTok{arrange}\NormalTok{(}\FunctionTok{desc}\NormalTok{(n\_journals))}

\NormalTok{cluster\_table }\SpecialCharTok{|\textgreater{}} \FunctionTok{head}\NormalTok{(}\DecValTok{8}\NormalTok{) }\SpecialCharTok{|\textgreater{}} \FunctionTok{gt}\NormalTok{()}
\end{Highlighting}
\end{Shaded}

\begingroup
\fontsize{8.0pt}{9.0pt}\selectfont
\begin{longtable}{>{\centering\arraybackslash}p{\dimexpr 0.21\linewidth -2\tabcolsep-1.5\arrayrulewidth}>{\raggedleft\arraybackslash}p{\dimexpr 0.23\linewidth -2\tabcolsep-1.5\arrayrulewidth}>{\raggedright\arraybackslash}p{\dimexpr 0.56\linewidth -2\tabcolsep-1.5\arrayrulewidth}}
\toprule
cluster & \nolinkurl{n_journals} & \nolinkurl{top_journals} \\ 
\midrule\addlinespace[2.5pt]
1 & 143 & Digital Collections of Colorado (Colorado State University); Institutional Repositories DataBase (IRDB); PubMed \\ 
2 & 5 & Bulletin des Médecins Suisses; European Heart Journal; Nature \\ 
\bottomrule
\end{longtable}
\endgroup

\subsection{Visualising the base map}\label{visualising-the-base-map}

\begin{Shaded}
\begin{Highlighting}[]
\FunctionTok{set.seed}\NormalTok{(}\DecValTok{42}\NormalTok{)}
\NormalTok{layout }\OtherTok{\textless{}{-}} \FunctionTok{create\_layout}\NormalTok{(}\FunctionTok{as\_tbl\_graph}\NormalTok{(g\_journals), }\AttributeTok{layout =} \StringTok{"fr"}\NormalTok{)}

\FunctionTok{ggraph}\NormalTok{(layout) }\SpecialCharTok{+}
  \FunctionTok{geom\_edge\_link}\NormalTok{(}\AttributeTok{alpha =} \FloatTok{0.05}\NormalTok{, }\AttributeTok{colour =} \StringTok{"grey70"}\NormalTok{) }\SpecialCharTok{+}
  \FunctionTok{geom\_node\_point}\NormalTok{(}\FunctionTok{aes}\NormalTok{(}\AttributeTok{colour =}\NormalTok{ cluster, }\AttributeTok{size =}\NormalTok{ degree), }\AttributeTok{alpha =} \FloatTok{0.7}\NormalTok{) }\SpecialCharTok{+}
  \FunctionTok{scale\_size\_continuous}\NormalTok{(}\AttributeTok{range =} \FunctionTok{c}\NormalTok{(}\DecValTok{1}\NormalTok{, }\DecValTok{5}\NormalTok{), }\AttributeTok{guide =} \StringTok{"none"}\NormalTok{) }\SpecialCharTok{+}
  \FunctionTok{scale\_colour\_manual}\NormalTok{(}\AttributeTok{values =} \FunctionTok{palette\_sci}\NormalTok{(}
    \FunctionTok{n\_distinct}\NormalTok{(}\FunctionTok{V}\NormalTok{(g\_journals)}\SpecialCharTok{$}\NormalTok{cluster)}
\NormalTok{  )) }\SpecialCharTok{+}
  \FunctionTok{labs}\NormalTok{(}\AttributeTok{colour =} \StringTok{"Cluster"}\NormalTok{) }\SpecialCharTok{+}
  \FunctionTok{theme\_void}\NormalTok{(}\AttributeTok{base\_family =} \StringTok{"sans"}\NormalTok{, }\AttributeTok{base\_size =} \DecValTok{11}\NormalTok{) }\SpecialCharTok{+}
  \FunctionTok{theme}\NormalTok{(}\AttributeTok{legend.position =} \StringTok{"bottom"}\NormalTok{)}
\end{Highlighting}
\end{Shaded}

\begin{figure}

{\centering \includegraphics[width=0.9\linewidth,alt={Network map of journals positioned by co-citation similarity. Colours indicate discipline clusters. Journals that cite each other's papers are positioned nearby.}]{scientometrics-in-r_files/figure-latex/19-fig-scimap-basemap-1} 

}

\caption{Journal co-citation base map coloured by discipline cluster.}\label{fig:19-fig-scimap-basemap}
\end{figure}

\subsection{Creating an overlay map}\label{creating-an-overlay-map}

We overlay the publication profile of an institution.

\begin{Shaded}
\begin{Highlighting}[]
\NormalTok{inst\_works }\OtherTok{\textless{}{-}} \FunctionTok{fetch\_openalex}\NormalTok{(}
  \AttributeTok{entity =} \StringTok{"works"}\NormalTok{,}
  \AttributeTok{authorships.institutions.id =} \StringTok{"I40120149"}\NormalTok{,}
  \AttributeTok{from\_publication\_date =} \StringTok{"2022{-}01{-}01"}\NormalTok{,}
  \AttributeTok{to\_publication\_date =} \StringTok{"2022{-}12{-}31"}\NormalTok{,}
  \AttributeTok{type =} \StringTok{"article"}\NormalTok{,}
  \AttributeTok{options =} \FunctionTok{list}\NormalTok{(}\AttributeTok{sample =} \DecValTok{300}\NormalTok{, }\AttributeTok{seed =} \DecValTok{42}\NormalTok{)}
\NormalTok{)}

\NormalTok{inst\_journals }\OtherTok{\textless{}{-}}\NormalTok{ inst\_works }\SpecialCharTok{|\textgreater{}}
  \FunctionTok{filter}\NormalTok{(}\SpecialCharTok{!}\FunctionTok{is.na}\NormalTok{(source\_display\_name)) }\SpecialCharTok{|\textgreater{}}
  \FunctionTok{count}\NormalTok{(source\_display\_name, }\AttributeTok{name =} \StringTok{"n\_pubs"}\NormalTok{, }\AttributeTok{sort =} \ConstantTok{TRUE}\NormalTok{)}

\NormalTok{basemap\_journals }\OtherTok{\textless{}{-}} \FunctionTok{V}\NormalTok{(g\_journals)}\SpecialCharTok{$}\NormalTok{name}
\NormalTok{inst\_journals\_matched }\OtherTok{\textless{}{-}}\NormalTok{ inst\_journals }\SpecialCharTok{|\textgreater{}}
  \FunctionTok{filter}\NormalTok{(source\_display\_name }\SpecialCharTok{\%in\%}\NormalTok{ basemap\_journals)}

\FunctionTok{cat}\NormalTok{(}\FunctionTok{glue}\NormalTok{(}\StringTok{"Institution journals matched to base map: \{nrow(inst\_journals\_matched)\} / \{nrow(inst\_journals)\}}\SpecialCharTok{\textbackslash{}n}\StringTok{"}\NormalTok{))}
\end{Highlighting}
\end{Shaded}

\begin{verbatim}
#> Institution journals matched to base map: 7 / 259
\end{verbatim}

\begin{Shaded}
\begin{Highlighting}[]
\NormalTok{overlay\_sizes }\OtherTok{\textless{}{-}} \FunctionTok{tibble}\NormalTok{(}\AttributeTok{journal =} \FunctionTok{V}\NormalTok{(g\_journals)}\SpecialCharTok{$}\NormalTok{name) }\SpecialCharTok{|\textgreater{}}
  \FunctionTok{left\_join}\NormalTok{(inst\_journals\_matched, }\AttributeTok{by =} \FunctionTok{c}\NormalTok{(}\StringTok{"journal"} \OtherTok{=} \StringTok{"source\_display\_name"}\NormalTok{)) }\SpecialCharTok{|\textgreater{}}
  \FunctionTok{mutate}\NormalTok{(}\AttributeTok{n\_pubs =} \FunctionTok{replace\_na}\NormalTok{(n\_pubs, }\DecValTok{0}\NormalTok{))}

\FunctionTok{set.seed}\NormalTok{(}\DecValTok{42}\NormalTok{)}
\NormalTok{layout\_overlay }\OtherTok{\textless{}{-}} \FunctionTok{create\_layout}\NormalTok{(}\FunctionTok{as\_tbl\_graph}\NormalTok{(g\_journals), }\AttributeTok{layout =} \StringTok{"fr"}\NormalTok{)}

\FunctionTok{ggraph}\NormalTok{(layout\_overlay) }\SpecialCharTok{+}
  \FunctionTok{geom\_edge\_link}\NormalTok{(}\AttributeTok{alpha =} \FloatTok{0.03}\NormalTok{, }\AttributeTok{colour =} \StringTok{"grey80"}\NormalTok{) }\SpecialCharTok{+}
  \FunctionTok{geom\_node\_point}\NormalTok{(}\AttributeTok{colour =} \StringTok{"grey85"}\NormalTok{, }\AttributeTok{size =} \FloatTok{1.5}\NormalTok{, }\AttributeTok{alpha =} \FloatTok{0.5}\NormalTok{) }\SpecialCharTok{+}
  \FunctionTok{geom\_node\_point}\NormalTok{(}
    \AttributeTok{data =}\NormalTok{ layout\_overlay }\SpecialCharTok{|\textgreater{}} \FunctionTok{left\_join}\NormalTok{(overlay\_sizes, }\AttributeTok{by =} \FunctionTok{c}\NormalTok{(}\StringTok{"name"} \OtherTok{=} \StringTok{"journal"}\NormalTok{)) }\SpecialCharTok{|\textgreater{}}
      \FunctionTok{filter}\NormalTok{(n\_pubs }\SpecialCharTok{\textgreater{}} \DecValTok{0}\NormalTok{),}
    \FunctionTok{aes}\NormalTok{(}\AttributeTok{size =}\NormalTok{ n\_pubs, }\AttributeTok{colour =}\NormalTok{ cluster), }\AttributeTok{alpha =} \FloatTok{0.8}
\NormalTok{  ) }\SpecialCharTok{+}
  \FunctionTok{scale\_size\_continuous}\NormalTok{(}\AttributeTok{range =} \FunctionTok{c}\NormalTok{(}\DecValTok{2}\NormalTok{, }\DecValTok{8}\NormalTok{), }\AttributeTok{name =} \StringTok{"Publications"}\NormalTok{) }\SpecialCharTok{+}
  \FunctionTok{scale\_colour\_manual}\NormalTok{(}\AttributeTok{values =} \FunctionTok{palette\_sci}\NormalTok{(}
    \FunctionTok{n\_distinct}\NormalTok{(}\FunctionTok{V}\NormalTok{(g\_journals)}\SpecialCharTok{$}\NormalTok{cluster)}
\NormalTok{  )) }\SpecialCharTok{+}
  \FunctionTok{labs}\NormalTok{(}\AttributeTok{colour =} \StringTok{"Cluster"}\NormalTok{) }\SpecialCharTok{+}
  \FunctionTok{theme\_void}\NormalTok{(}\AttributeTok{base\_family =} \StringTok{"sans"}\NormalTok{, }\AttributeTok{base\_size =} \DecValTok{11}\NormalTok{) }\SpecialCharTok{+}
  \FunctionTok{theme}\NormalTok{(}\AttributeTok{legend.position =} \StringTok{"bottom"}\NormalTok{)}
\end{Highlighting}
\end{Shaded}

\begin{figure}

{\centering \includegraphics[width=0.9\linewidth,alt={The journal base map with Oxford's publications overlaid as larger nodes. The overlay reveals which areas of science the university is active in.}]{scientometrics-in-r_files/figure-latex/19-fig-scimap-overlay-1} 

}

\caption{Overlay map showing Oxford's 2022 publication profile on the journal base map.}\label{fig:19-fig-scimap-overlay}
\end{figure}

\section{Diagnostics and interpretation}\label{scimap-diagnostics}

\begin{itemize}
\tightlist
\item
  \textbf{Journal matching rate}: The fraction of the entity's publications that match journals in the base map determines overlay completeness. Low matching rates (\textless{} 50\%) indicate the base map is too narrow for the entity's profile.
\item
  \textbf{Layout stability}: Force-directed layouts are sensitive to initial conditions. Always use \texttt{set.seed()} and report the seed. Consider consensus layouts from multiple seeds for important maps.
\item
  \textbf{Cluster labelling}: Algorithmic clusters need human interpretation. Label clusters by inspecting the journals they contain, not by relying on the algorithm.
\item
  \textbf{Overlay concentration}: If the overlay is concentrated in one or two clusters, the institution is highly specialised. Dispersed overlays indicate breadth. Neither is inherently better.
\end{itemize}

\section{Limitations and responsible use}\label{scimap-limitations}

\section{Limitations and responsible use}\label{limitations-and-responsible-use-18}

\begin{itemize}
\tightlist
\item
  \textbf{Projection distortion.} Two-dimensional layouts compress high-dimensional citation relationships. Journals at the boundaries between disciplines may be misplaced. Do not over-interpret exact positions.
\item
  \textbf{Journal-level aggregation.} Mapping publications to journals loses within-journal variation. A paper in \emph{Nature} could be biology, physics, or chemistry --- the overlay cannot distinguish.
\item
  \textbf{Sample-dependent base maps.} Our base map is built from a small sample. Production science maps (e.g., the CWTS journal map) use comprehensive citation data. Treat this example as illustrative.
\item
  \textbf{Not evaluative.} Science maps show \emph{where} research happens, not how good it is. An empty region on a university's overlay map is not a weakness --- it may reflect a deliberate strategic choice \citep{hicks2015}.
\item
  \textbf{Interdisciplinary research is underrepresented.} Journals at disciplinary boundaries may cluster with one discipline, making interdisciplinary work partially invisible \citep{rafols2010}.
\end{itemize}

\section{Common pitfalls}\label{scimap-pitfalls}

\section{Common pitfalls}\label{common-pitfalls-18}

\begin{itemize}
\tightlist
\item
  \textbf{Using too small a base map.} A base map with fewer than 50 journals lacks resolution. Aim for at least 100--200 journals for meaningful discipline separation.
\item
  \textbf{Comparing overlays from different base maps.} Node positions are not transferable between different base maps. To compare institutions, overlay them on the \emph{same} base map.
\item
  \textbf{Ignoring unmatched publications.} Publications in journals not on the base map simply disappear from the overlay. Report the matching rate.
\item
  \textbf{Treating clusters as fixed disciplines.} Cluster assignments depend on the algorithm, resolution parameter, and data sample. They are analytical tools, not natural categories.
\end{itemize}

\section{Exercises}\label{scimap-exercises}

\begin{enumerate}
\def\labelenumi{\arabic{enumi}.}
\item
  \textbf{Multi-institution overlay.} Overlay the profiles of two institutions on the same base map using different colours. Where do their research profiles overlap, and where do they diverge?
\item
  \textbf{Temporal overlay.} Build overlays for the same institution in two different years. Has the research profile shifted toward new areas or remained stable?
\item
  \textbf{Country-level overlay.} Instead of an institution, overlay a country's publication profile. How does a small specialised country (e.g., Singapore) compare to a large diversified one (e.g., the United States)?
\end{enumerate}

\section{Solutions}\label{19-solutions}

Solutions are provided in \ref{solutions}.

\section{Further reading}\label{scimap-reading}

\begin{itemize}
\tightlist
\item
  \citet{rafols2010} --- Overlay maps for visualising research profiles on global science maps.
\item
  \citet{waltman2012} --- Unified approach to constructing and clustering bibliometric maps.
\item
  \citet{hicks2015} --- The Leiden Manifesto: responsible interpretation of quantitative research profiles.
\item
  \citet{fortunato2010} --- Community detection methods relevant to journal clustering.
\end{itemize}

\section{Session info}\label{scimap-sessioninfo}

\begin{Shaded}
\begin{Highlighting}[]
\FunctionTok{sessionInfo}\NormalTok{()}
\end{Highlighting}
\end{Shaded}

\begin{verbatim}
#> R version 4.6.1 (2026-06-24)
#> Platform: x86_64-pc-linux-gnu
#> Running under: Ubuntu 24.04.5 LTS
#> 
#> Matrix products: default
#> BLAS:   /usr/lib/x86_64-linux-gnu/openblas-pthread/libblas.so.3 
#> LAPACK: /usr/lib/x86_64-linux-gnu/openblas-pthread/libopenblasp-r0.3.26.so;  LAPACK version 3.12.0
#> 
#> locale:
#>  [1] LC_CTYPE=C.UTF-8       LC_NUMERIC=C           LC_TIME=C.UTF-8       
#>  [4] LC_COLLATE=C.UTF-8     LC_MONETARY=C.UTF-8    LC_MESSAGES=C.UTF-8   
#>  [7] LC_PAPER=C.UTF-8       LC_NAME=C              LC_ADDRESS=C          
#> [10] LC_TELEPHONE=C         LC_MEASUREMENT=C.UTF-8 LC_IDENTIFICATION=C   
#> 
#> time zone: UTC
#> tzcode source: system (glibc)
#> 
#> attached base packages:
#> [1] stats     graphics  grDevices utils     datasets  methods   base     
#> 
#> other attached packages:
#>  [1] ggraph_2.2.2       tidygraph_1.3.1    igraph_2.3.4       quanteda_4.5.1    
#>  [5] pdftools_3.9.1     arrow_26.0.0       bibliometrix_5.5.0 RefManageR_1.4.0  
#>  [9] bib2df_1.1.2.0     rcrossref_1.2.1    gt_1.4.0           tidytext_0.4.3    
#> [13] glue_1.8.1         openalexR_3.1.2    lubridate_1.9.5    forcats_1.0.1     
#> [17] stringr_1.6.0      dplyr_1.2.1        purrr_1.2.2        readr_2.2.0       
#> [21] tidyr_1.3.2        tibble_3.3.1       ggplot2_4.0.3      tidyverse_2.0.0   
#> 
#> loaded via a namespace (and not attached):
#>   [1] bibtex_0.5.3           RColorBrewer_1.1-3     rstudioapi_0.19.0     
#>   [4] jsonlite_2.0.0         magrittr_2.0.5         farver_2.1.2          
#>   [7] rmarkdown_2.32         fs_2.1.0               vctrs_0.7.3           
#>  [10] memoise_2.0.1          askpass_1.2.1          base64enc_0.1-6       
#>  [13] htmltools_0.5.9        contentanalysis_1.1.1  curl_8.1.0            
#>  [16] janeaustenr_1.0.0      cellranger_1.1.0       htmlwidgets_1.6.4     
#>  [19] tokenizers_0.3.0       plyr_1.8.9             httr2_1.3.0           
#>  [22] cachem_1.1.0           plotly_4.12.1          dimensionsR_0.0.3     
#>  [25] mime_0.13              lifecycle_1.0.5        pkgconfig_2.0.3       
#>  [28] Matrix_1.7-5           R6_2.6.1               fastmap_1.2.0         
#>  [31] shiny_1.14.0           digest_0.6.39          shinycssloaders_1.1.0 
#>  [34] rprojroot_2.1.1        SnowballC_0.7.1        labeling_0.4.3        
#>  [37] urltools_1.7.3.1       timechange_0.4.0       httr_1.4.9            
#>  [40] polyclip_1.10-7        compiler_4.6.1         here_1.0.2            
#>  [43] bit64_4.8.6            withr_3.0.3            S7_0.2.2              
#>  [46] viridis_0.6.5          ggforce_0.5.0          MASS_7.3-65           
#>  [49] bibliometrixData_0.3.0 tools_4.6.1            otel_0.2.0            
#>  [52] stopwords_2.3          zip_3.0.2              rentrez_1.2.4         
#>  [55] httpuv_1.6.17          promises_1.5.0         grid_4.6.1            
#>  [58] stringdist_0.9.17      generics_0.1.4         gtable_0.3.6          
#>  [61] tzdb_0.5.0             rscopus_0.9.0          ca_0.72               
#>  [64] data.table_1.18.6.1    hms_1.1.4              xml2_1.6.0            
#>  [67] utf8_1.2.6             ggrepel_0.9.8          pillar_1.11.1         
#>  [70] later_1.4.8            tweenr_2.0.3           lattice_0.22-9        
#>  [73] bit_4.6.0              tidyselect_1.2.1       miniUI_0.1.2          
#>  [76] knitr_1.52             gridExtra_2.3.1        bookdown_0.48         
#>  [79] crul_1.6.0             xfun_0.62              graphlayouts_1.2.5    
#>  [82] DT_0.34.0              humaniformat_0.6.0     visNetwork_2.1.4      
#>  [85] stringi_1.8.9          qpdf_1.4.1             yaml_2.3.12           
#>  [88] evaluate_1.0.5         httpcode_0.3.0         cli_3.6.6             
#>  [91] xtable_1.8-8           systemfonts_1.3.2      Rcpp_1.1.2            
#>  [94] readxl_1.6.0           triebeard_0.4.1        XML_3.99-0.25         
#>  [97] parallel_4.6.1         assertthat_0.2.1       pubmedR_1.0.2         
#> [100] viridisLite_0.4.3      scales_1.4.0           openxlsx_4.2.9        
#> [103] rlang_1.3.0            fastmatch_1.1-8
\end{verbatim}

\chapter{Network Visualization and Interoperability}\label{network-visualization-and-interoperability}

\emph{R. Heller · Date: 2026-05-09 · Last changed: 2026-10-31}

\section{Learning objectives}\label{viz-objectives}

After completing this chapter, you will be able to:

\begin{itemize}
\tightlist
\item
  Choose appropriate layout algorithms for different network types
\item
  Export igraph networks to Gephi (GraphML), VOSviewer, and Pajek formats
\item
  Create interactive network visualisations with \texttt{visNetwork}
\item
  Apply colour and size encoding that communicates structure clearly
\item
  Produce publication-quality static network figures with \texttt{ggraph}
\end{itemize}

\section{Setup}\label{viz-setup}

\begin{Shaded}
\begin{Highlighting}[]
\FunctionTok{library}\NormalTok{(tidyverse)}
\FunctionTok{library}\NormalTok{(openalexR)}
\FunctionTok{library}\NormalTok{(igraph)}
\FunctionTok{library}\NormalTok{(tidygraph)}
\FunctionTok{library}\NormalTok{(ggraph)}
\FunctionTok{library}\NormalTok{(visNetwork)}
\FunctionTok{library}\NormalTok{(glue)}

\FunctionTok{set.seed}\NormalTok{(}\DecValTok{20260509}\NormalTok{)}

\FunctionTok{source}\NormalTok{(here}\SpecialCharTok{::}\FunctionTok{here}\NormalTok{(}\StringTok{"R"}\NormalTok{, }\StringTok{"api\_helpers.R"}\NormalTok{))}
\FunctionTok{source}\NormalTok{(here}\SpecialCharTok{::}\FunctionTok{here}\NormalTok{(}\StringTok{"R"}\NormalTok{, }\StringTok{"utils.R"}\NormalTok{))}
\FunctionTok{source}\NormalTok{(here}\SpecialCharTok{::}\FunctionTok{here}\NormalTok{(}\StringTok{"R"}\NormalTok{, }\StringTok{"sci\_palette.R"}\NormalTok{))}
\end{Highlighting}
\end{Shaded}

\section{Conceptual background}\label{viz-background}

Network visualisation is both an analytical tool and a communication device. A well-designed network figure can reveal structure --- clusters, bridges, outliers --- at a glance. A poorly designed one obscures these patterns in visual noise.

\textbf{Layout algorithms} determine node positions. Force-directed layouts (Fruchterman-Reingold, Kamada-Kawai) simulate physical forces: connected nodes attract, all nodes repel. They work well for small-to-medium networks (up to \textasciitilde1,000 nodes) and produce aesthetically pleasing results. For larger networks, the DrL (Distributed Recursive Layout) or OpenOrd algorithms scale better. For trees or hierarchies, Reingold-Tilford or Sugiyama layouts are appropriate. The choice of algorithm depends on the network's structure and the visual question being asked.

\textbf{Visual encoding} maps network properties to visual channels: node size (degree, centrality), node colour (community, attribute), edge width (weight), and edge colour or transparency. Effective encoding follows the principle of proportional ink: the visual weight of an element should correspond to its data importance. Avoid redundant encoding (mapping the same variable to both size and colour) and limit the number of distinct colours to what the reader can distinguish (typically 8--12 categories) \citep{fortunato2010}.

\textbf{Interoperability} is essential because no single tool excels at everything. R with igraph and ggraph is excellent for reproducible analysis and publication figures. Gephi \citep{bastian2009} provides a rich GUI for exploration and layout refinement. VOSviewer specialises in bibliometric network visualisation with built-in clustering. Pajek handles very large networks efficiently. A practical workflow often involves constructing and analysing the network in R, then exporting to a specialised tool for specific visualisation tasks.

\textbf{Reproducibility} requires that layouts are deterministic. Force-directed algorithms start from random positions, so \texttt{set.seed()} before every layout computation. For cross-tool reproducibility, export node coordinates alongside the network data.

\section{Worked example}\label{viz-example}

\subsection{Building a sample network}\label{building-a-sample-network}

We reuse a co-authorship network from Scientometrics.

\begin{Shaded}
\begin{Highlighting}[]
\NormalTok{works }\OtherTok{\textless{}{-}} \FunctionTok{fetch\_openalex}\NormalTok{(}
  \AttributeTok{entity =} \StringTok{"works"}\NormalTok{,}
  \AttributeTok{primary\_location.source.id =} \StringTok{"S148561398"}\NormalTok{,}
  \AttributeTok{from\_publication\_date =} \StringTok{"2021{-}01{-}01"}\NormalTok{,}
  \AttributeTok{to\_publication\_date =} \StringTok{"2023{-}12{-}31"}\NormalTok{,}
  \AttributeTok{options =} \FunctionTok{list}\NormalTok{(}\AttributeTok{sample =} \DecValTok{200}\NormalTok{, }\AttributeTok{seed =} \DecValTok{42}\NormalTok{)}
\NormalTok{)}

\NormalTok{author\_data }\OtherTok{\textless{}{-}}\NormalTok{ works }\SpecialCharTok{|\textgreater{}}
  \FunctionTok{select}\NormalTok{(id, authorships) }\SpecialCharTok{|\textgreater{}}
  \FunctionTok{unnest}\NormalTok{(authorships, }\AttributeTok{names\_sep =} \StringTok{"\_"}\NormalTok{) }\SpecialCharTok{|\textgreater{}}
  \FunctionTok{select}\NormalTok{(}\AttributeTok{work\_id =}\NormalTok{ id, }\AttributeTok{author\_id =}\NormalTok{ authorships\_id,}
         \AttributeTok{author\_name =}\NormalTok{ authorships\_display\_name) }\SpecialCharTok{|\textgreater{}}
  \FunctionTok{filter}\NormalTok{(}\SpecialCharTok{!}\FunctionTok{is.na}\NormalTok{(author\_id))}

\NormalTok{edges }\OtherTok{\textless{}{-}}\NormalTok{ author\_data }\SpecialCharTok{|\textgreater{}}
  \FunctionTok{inner\_join}\NormalTok{(author\_data, }\AttributeTok{by =} \StringTok{"work\_id"}\NormalTok{, }\AttributeTok{suffix =} \FunctionTok{c}\NormalTok{(}\StringTok{"\_1"}\NormalTok{, }\StringTok{"\_2"}\NormalTok{),}
             \AttributeTok{relationship =} \StringTok{"many{-}to{-}many"}\NormalTok{) }\SpecialCharTok{|\textgreater{}}
  \FunctionTok{filter}\NormalTok{(author\_id\_1 }\SpecialCharTok{\textless{}}\NormalTok{ author\_id\_2) }\SpecialCharTok{|\textgreater{}}
  \FunctionTok{count}\NormalTok{(author\_id\_1, author\_id\_2, }\AttributeTok{name =} \StringTok{"weight"}\NormalTok{)}

\NormalTok{g }\OtherTok{\textless{}{-}} \FunctionTok{graph\_from\_data\_frame}\NormalTok{(}
\NormalTok{  edges }\SpecialCharTok{|\textgreater{}} \FunctionTok{select}\NormalTok{(author\_id\_1, author\_id\_2, weight),}
  \AttributeTok{directed =} \ConstantTok{FALSE}
\NormalTok{) }\SpecialCharTok{|\textgreater{}}
  \FunctionTok{simplify}\NormalTok{(}\AttributeTok{edge.attr.comb =} \FunctionTok{list}\NormalTok{(}\AttributeTok{weight =} \StringTok{"sum"}\NormalTok{))}

\NormalTok{comp }\OtherTok{\textless{}{-}} \FunctionTok{components}\NormalTok{(g)}
\NormalTok{giant }\OtherTok{\textless{}{-}} \FunctionTok{induced\_subgraph}\NormalTok{(g, }\FunctionTok{which}\NormalTok{(comp}\SpecialCharTok{$}\NormalTok{membership }\SpecialCharTok{==} \FunctionTok{which.max}\NormalTok{(comp}\SpecialCharTok{$}\NormalTok{csize)))}

\NormalTok{comm }\OtherTok{\textless{}{-}} \FunctionTok{cluster\_leiden}\NormalTok{(giant, }\AttributeTok{resolution\_parameter =} \FloatTok{1.0}\NormalTok{,}
                       \AttributeTok{objective\_function =} \StringTok{"modularity"}\NormalTok{)}
\FunctionTok{V}\NormalTok{(giant)}\SpecialCharTok{$}\NormalTok{community }\OtherTok{\textless{}{-}} \FunctionTok{membership}\NormalTok{(comm)}
\FunctionTok{V}\NormalTok{(giant)}\SpecialCharTok{$}\NormalTok{degree }\OtherTok{\textless{}{-}} \FunctionTok{degree}\NormalTok{(giant)}

\NormalTok{author\_lookup }\OtherTok{\textless{}{-}}\NormalTok{ author\_data }\SpecialCharTok{|\textgreater{}} \FunctionTok{distinct}\NormalTok{(author\_id, author\_name)}
\FunctionTok{V}\NormalTok{(giant)}\SpecialCharTok{$}\NormalTok{label }\OtherTok{\textless{}{-}}\NormalTok{ author\_lookup}\SpecialCharTok{$}\NormalTok{author\_name[}
  \FunctionTok{match}\NormalTok{(}\FunctionTok{V}\NormalTok{(giant)}\SpecialCharTok{$}\NormalTok{name, author\_lookup}\SpecialCharTok{$}\NormalTok{author\_id)}
\NormalTok{]}

\FunctionTok{cat}\NormalTok{(}\FunctionTok{glue}\NormalTok{(}\StringTok{"Network: \{vcount(giant)\} nodes, \{ecount(giant)\} edges}\SpecialCharTok{\textbackslash{}n}\StringTok{"}\NormalTok{))}
\end{Highlighting}
\end{Shaded}

\begin{verbatim}
#> Network: 12 nodes, 22 edges
\end{verbatim}

\subsection{Comparing layout algorithms}\label{comparing-layout-algorithms}

\begin{Shaded}
\begin{Highlighting}[]
\NormalTok{tg }\OtherTok{\textless{}{-}} \FunctionTok{as\_tbl\_graph}\NormalTok{(giant) }\SpecialCharTok{|\textgreater{}}
  \FunctionTok{mutate}\NormalTok{(}\AttributeTok{community =} \FunctionTok{as.factor}\NormalTok{(community))}

\NormalTok{layouts }\OtherTok{\textless{}{-}} \FunctionTok{c}\NormalTok{(}\StringTok{"fr"}\NormalTok{, }\StringTok{"kk"}\NormalTok{, }\StringTok{"stress"}\NormalTok{, }\StringTok{"drl"}\NormalTok{)}
\NormalTok{layout\_names }\OtherTok{\textless{}{-}} \FunctionTok{c}\NormalTok{(}\StringTok{"Fruchterman{-}Reingold"}\NormalTok{, }\StringTok{"Kamada{-}Kawai"}\NormalTok{, }\StringTok{"Stress"}\NormalTok{, }\StringTok{"DrL"}\NormalTok{)}

\NormalTok{plots }\OtherTok{\textless{}{-}} \FunctionTok{map2}\NormalTok{(layouts, layout\_names, }\ControlFlowTok{function}\NormalTok{(algo, name) \{}
  \FunctionTok{set.seed}\NormalTok{(}\DecValTok{42}\NormalTok{)}
  \FunctionTok{ggraph}\NormalTok{(tg, }\AttributeTok{layout =}\NormalTok{ algo) }\SpecialCharTok{+}
    \FunctionTok{geom\_edge\_link}\NormalTok{(}\AttributeTok{alpha =} \FloatTok{0.1}\NormalTok{, }\AttributeTok{colour =} \StringTok{"grey60"}\NormalTok{) }\SpecialCharTok{+}
    \FunctionTok{geom\_node\_point}\NormalTok{(}\FunctionTok{aes}\NormalTok{(}\AttributeTok{colour =}\NormalTok{ community, }\AttributeTok{size =}\NormalTok{ degree), }\AttributeTok{alpha =} \FloatTok{0.7}\NormalTok{) }\SpecialCharTok{+}
    \FunctionTok{scale\_size\_continuous}\NormalTok{(}\AttributeTok{range =} \FunctionTok{c}\NormalTok{(}\FloatTok{0.5}\NormalTok{, }\DecValTok{4}\NormalTok{), }\AttributeTok{guide =} \StringTok{"none"}\NormalTok{) }\SpecialCharTok{+}
    \FunctionTok{scale\_colour\_manual}\NormalTok{(}\AttributeTok{values =} \FunctionTok{palette\_sci}\NormalTok{(}
      \FunctionTok{n\_distinct}\NormalTok{(}\FunctionTok{V}\NormalTok{(giant)}\SpecialCharTok{$}\NormalTok{community)}
\NormalTok{    ), }\AttributeTok{guide =} \StringTok{"none"}\NormalTok{) }\SpecialCharTok{+}
    \FunctionTok{labs}\NormalTok{(}\AttributeTok{title =}\NormalTok{ name) }\SpecialCharTok{+}
    \FunctionTok{theme\_void}\NormalTok{(}\AttributeTok{base\_family =} \StringTok{"sans"}\NormalTok{, }\AttributeTok{base\_size =} \DecValTok{11}\NormalTok{)}
\NormalTok{\})}

\NormalTok{patchwork}\SpecialCharTok{::}\FunctionTok{wrap\_plots}\NormalTok{(plots, }\AttributeTok{ncol =} \DecValTok{2}\NormalTok{)}
\end{Highlighting}
\end{Shaded}

\begin{figure}

{\centering \includegraphics[width=0.9\linewidth,alt={A 2x2 panel showing the same co-authorship network plotted with Fruchterman-Reingold, Kamada-Kawai, stress, and DRL layouts. Node colours indicate communities.}]{scientometrics-in-r_files/figure-latex/20-fig-viz-layouts-1} 

}

\caption{The same network rendered with four different layout algorithms.}\label{fig:20-fig-viz-layouts}
\end{figure}

\subsection{Exporting to GraphML (Gephi)}\label{exporting-to-graphml-gephi}

\begin{Shaded}
\begin{Highlighting}[]
\NormalTok{out\_dir }\OtherTok{\textless{}{-}}\NormalTok{ here}\SpecialCharTok{::}\FunctionTok{here}\NormalTok{(}\StringTok{"\_freeze"}\NormalTok{, }\StringTok{"exports"}\NormalTok{)}
\FunctionTok{dir.create}\NormalTok{(out\_dir, }\AttributeTok{recursive =} \ConstantTok{TRUE}\NormalTok{, }\AttributeTok{showWarnings =} \ConstantTok{FALSE}\NormalTok{)}
\NormalTok{graphml\_path }\OtherTok{\textless{}{-}} \FunctionTok{file.path}\NormalTok{(out\_dir, }\StringTok{"coauth\_network.graphml"}\NormalTok{)}
\FunctionTok{write\_graph}\NormalTok{(giant, graphml\_path, }\AttributeTok{format =} \StringTok{"graphml"}\NormalTok{)}
\FunctionTok{cat}\NormalTok{(}\FunctionTok{glue}\NormalTok{(}\StringTok{"Exported GraphML: \{graphml\_path\}}\SpecialCharTok{\textbackslash{}n}\StringTok{"}\NormalTok{))}
\end{Highlighting}
\end{Shaded}

\begin{verbatim}
#> Exported GraphML: /home/runner/work/scientometrics-in-r/scientometrics-in-r/_freeze/exports/coauth_network.graphml
\end{verbatim}

\begin{Shaded}
\begin{Highlighting}[]
\FunctionTok{cat}\NormalTok{(}\FunctionTok{glue}\NormalTok{(}\StringTok{"File size: \{round(file.size(graphml\_path) / 1024, 1)\} KB}\SpecialCharTok{\textbackslash{}n}\StringTok{"}\NormalTok{))}
\end{Highlighting}
\end{Shaded}

\begin{verbatim}
#> File size: 5.1 KB
\end{verbatim}

\subsection{Exporting to Pajek format}\label{exporting-to-pajek-format}

\begin{Shaded}
\begin{Highlighting}[]
\NormalTok{pajek\_path }\OtherTok{\textless{}{-}} \FunctionTok{file.path}\NormalTok{(out\_dir, }\StringTok{"coauth\_network.net"}\NormalTok{)}
\FunctionTok{write\_graph}\NormalTok{(giant, pajek\_path, }\AttributeTok{format =} \StringTok{"pajek"}\NormalTok{)}
\FunctionTok{cat}\NormalTok{(}\FunctionTok{glue}\NormalTok{(}\StringTok{"Exported Pajek: \{pajek\_path\}}\SpecialCharTok{\textbackslash{}n}\StringTok{"}\NormalTok{))}
\end{Highlighting}
\end{Shaded}

\begin{verbatim}
#> Exported Pajek: /home/runner/work/scientometrics-in-r/scientometrics-in-r/_freeze/exports/coauth_network.net
\end{verbatim}

\subsection{Exporting to VOSviewer format}\label{exporting-to-vosviewer-format}

VOSviewer reads tab-separated map and network files.

\begin{Shaded}
\begin{Highlighting}[]
\FunctionTok{set.seed}\NormalTok{(}\DecValTok{42}\NormalTok{)}
\NormalTok{coords }\OtherTok{\textless{}{-}} \FunctionTok{layout\_with\_fr}\NormalTok{(giant)}

\NormalTok{vos\_map }\OtherTok{\textless{}{-}} \FunctionTok{tibble}\NormalTok{(}
  \AttributeTok{id =} \FunctionTok{seq\_len}\NormalTok{(}\FunctionTok{vcount}\NormalTok{(giant)),}
  \AttributeTok{label =} \FunctionTok{V}\NormalTok{(giant)}\SpecialCharTok{$}\NormalTok{label,}
  \AttributeTok{x =}\NormalTok{ coords[, }\DecValTok{1}\NormalTok{],}
  \AttributeTok{y =}\NormalTok{ coords[, }\DecValTok{2}\NormalTok{],}
  \AttributeTok{cluster =} \FunctionTok{V}\NormalTok{(giant)}\SpecialCharTok{$}\NormalTok{community,}
  \AttributeTok{weight =} \FunctionTok{V}\NormalTok{(giant)}\SpecialCharTok{$}\NormalTok{degree}
\NormalTok{)}

\NormalTok{vos\_network }\OtherTok{\textless{}{-}} \FunctionTok{as\_data\_frame}\NormalTok{(giant, }\AttributeTok{what =} \StringTok{"edges"}\NormalTok{) }\SpecialCharTok{|\textgreater{}}
  \FunctionTok{mutate}\NormalTok{(}
    \AttributeTok{from\_idx =} \FunctionTok{match}\NormalTok{(from, }\FunctionTok{V}\NormalTok{(giant)}\SpecialCharTok{$}\NormalTok{name),}
    \AttributeTok{to\_idx =} \FunctionTok{match}\NormalTok{(to, }\FunctionTok{V}\NormalTok{(giant)}\SpecialCharTok{$}\NormalTok{name)}
\NormalTok{  ) }\SpecialCharTok{|\textgreater{}}
  \FunctionTok{select}\NormalTok{(from\_idx, to\_idx, weight)}

\NormalTok{map\_path }\OtherTok{\textless{}{-}} \FunctionTok{file.path}\NormalTok{(out\_dir, }\StringTok{"vosviewer\_map.txt"}\NormalTok{)}
\NormalTok{net\_path }\OtherTok{\textless{}{-}} \FunctionTok{file.path}\NormalTok{(out\_dir, }\StringTok{"vosviewer\_network.txt"}\NormalTok{)}

\FunctionTok{write\_tsv}\NormalTok{(vos\_map, map\_path)}
\FunctionTok{write\_tsv}\NormalTok{(vos\_network, net\_path)}
\FunctionTok{cat}\NormalTok{(}\FunctionTok{glue}\NormalTok{(}\StringTok{"VOSviewer map: \{map\_path\}}\SpecialCharTok{\textbackslash{}n}\StringTok{"}\NormalTok{))}
\end{Highlighting}
\end{Shaded}

\begin{verbatim}
#> VOSviewer map: /home/runner/work/scientometrics-in-r/scientometrics-in-r/_freeze/exports/vosviewer_map.txt
\end{verbatim}

\begin{Shaded}
\begin{Highlighting}[]
\FunctionTok{cat}\NormalTok{(}\FunctionTok{glue}\NormalTok{(}\StringTok{"VOSviewer network: \{net\_path\}}\SpecialCharTok{\textbackslash{}n}\StringTok{"}\NormalTok{))}
\end{Highlighting}
\end{Shaded}

\begin{verbatim}
#> VOSviewer network: /home/runner/work/scientometrics-in-r/scientometrics-in-r/_freeze/exports/vosviewer_network.txt
\end{verbatim}

\subsection{Interactive visualisation with visNetwork}\label{interactive-visualisation-with-visnetwork}

\begin{Shaded}
\begin{Highlighting}[]
\CommentTok{\# HTML widget — only rendered for HTML output (PDF/EPUB skip this chunk).}
\NormalTok{vis\_nodes }\OtherTok{\textless{}{-}} \FunctionTok{tibble}\NormalTok{(}
  \AttributeTok{id =} \FunctionTok{V}\NormalTok{(giant)}\SpecialCharTok{$}\NormalTok{name,}
  \AttributeTok{label =} \FunctionTok{ifelse}\NormalTok{(}\FunctionTok{V}\NormalTok{(giant)}\SpecialCharTok{$}\NormalTok{degree }\SpecialCharTok{\textgreater{}} \FunctionTok{quantile}\NormalTok{(}\FunctionTok{V}\NormalTok{(giant)}\SpecialCharTok{$}\NormalTok{degree, }\FloatTok{0.9}\NormalTok{),}
                 \FunctionTok{V}\NormalTok{(giant)}\SpecialCharTok{$}\NormalTok{label, }\StringTok{""}\NormalTok{),}
  \AttributeTok{title =} \FunctionTok{paste0}\NormalTok{(}\FunctionTok{V}\NormalTok{(giant)}\SpecialCharTok{$}\NormalTok{label, }\StringTok{"\textless{}br\textgreater{}Degree: "}\NormalTok{, }\FunctionTok{V}\NormalTok{(giant)}\SpecialCharTok{$}\NormalTok{degree,}
                 \StringTok{"\textless{}br\textgreater{}Community: "}\NormalTok{, }\FunctionTok{V}\NormalTok{(giant)}\SpecialCharTok{$}\NormalTok{community),}
  \AttributeTok{group =} \FunctionTok{as.character}\NormalTok{(}\FunctionTok{V}\NormalTok{(giant)}\SpecialCharTok{$}\NormalTok{community),}
  \AttributeTok{value =} \FunctionTok{V}\NormalTok{(giant)}\SpecialCharTok{$}\NormalTok{degree}
\NormalTok{)}

\NormalTok{vis\_edges }\OtherTok{\textless{}{-}} \FunctionTok{as\_data\_frame}\NormalTok{(giant, }\AttributeTok{what =} \StringTok{"edges"}\NormalTok{) }\SpecialCharTok{|\textgreater{}}
  \FunctionTok{select}\NormalTok{(from, to, }\AttributeTok{weight =}\NormalTok{ weight)}

\FunctionTok{visNetwork}\NormalTok{(vis\_nodes, vis\_edges, }\AttributeTok{width =} \StringTok{"100\%"}\NormalTok{, }\AttributeTok{height =} \StringTok{"500px"}\NormalTok{) }\SpecialCharTok{|\textgreater{}}
  \FunctionTok{visOptions}\NormalTok{(}\AttributeTok{highlightNearest =} \ConstantTok{TRUE}\NormalTok{, }\AttributeTok{nodesIdSelection =} \ConstantTok{TRUE}\NormalTok{) }\SpecialCharTok{|\textgreater{}}
  \FunctionTok{visPhysics}\NormalTok{(}\AttributeTok{stabilization =} \ConstantTok{FALSE}\NormalTok{) }\SpecialCharTok{|\textgreater{}}
  \FunctionTok{visLayout}\NormalTok{(}\AttributeTok{randomSeed =} \DecValTok{42}\NormalTok{)}
\end{Highlighting}
\end{Shaded}

\subsection{Publication-quality figure}\label{publication-quality-figure}

\begin{Shaded}
\begin{Highlighting}[]
\FunctionTok{set.seed}\NormalTok{(}\DecValTok{42}\NormalTok{)}

\FunctionTok{ggraph}\NormalTok{(tg, }\AttributeTok{layout =} \StringTok{"fr"}\NormalTok{) }\SpecialCharTok{+}
  \FunctionTok{geom\_edge\_link}\NormalTok{(}\AttributeTok{alpha =} \FloatTok{0.08}\NormalTok{, }\AttributeTok{colour =} \StringTok{"grey60"}\NormalTok{) }\SpecialCharTok{+}
  \FunctionTok{geom\_node\_point}\NormalTok{(}\FunctionTok{aes}\NormalTok{(}\AttributeTok{size =}\NormalTok{ degree, }\AttributeTok{colour =}\NormalTok{ community), }\AttributeTok{alpha =} \FloatTok{0.8}\NormalTok{) }\SpecialCharTok{+}
  \FunctionTok{geom\_node\_text}\NormalTok{(}
    \FunctionTok{aes}\NormalTok{(}\AttributeTok{label =} \FunctionTok{ifelse}\NormalTok{(degree }\SpecialCharTok{\textgreater{}} \FunctionTok{quantile}\NormalTok{(degree, }\FloatTok{0.95}\NormalTok{), label, }\ConstantTok{NA}\NormalTok{)),}
    \AttributeTok{repel =} \ConstantTok{TRUE}\NormalTok{, }\AttributeTok{size =} \FloatTok{2.5}\NormalTok{, }\AttributeTok{max.overlaps =} \DecValTok{15}\NormalTok{, }\AttributeTok{na.rm =} \ConstantTok{TRUE}
\NormalTok{  ) }\SpecialCharTok{+}
  \FunctionTok{scale\_size\_continuous}\NormalTok{(}\AttributeTok{range =} \FunctionTok{c}\NormalTok{(}\DecValTok{1}\NormalTok{, }\DecValTok{6}\NormalTok{), }\AttributeTok{guide =} \StringTok{"none"}\NormalTok{) }\SpecialCharTok{+}
  \FunctionTok{scale\_colour\_manual}\NormalTok{(}\AttributeTok{values =} \FunctionTok{palette\_sci}\NormalTok{(}
    \FunctionTok{n\_distinct}\NormalTok{(}\FunctionTok{V}\NormalTok{(giant)}\SpecialCharTok{$}\NormalTok{community)}
\NormalTok{  )) }\SpecialCharTok{+}
  \FunctionTok{labs}\NormalTok{(}\AttributeTok{colour =} \StringTok{"Community"}\NormalTok{) }\SpecialCharTok{+}
  \FunctionTok{theme\_void}\NormalTok{(}\AttributeTok{base\_family =} \StringTok{"sans"}\NormalTok{, }\AttributeTok{base\_size =} \DecValTok{11}\NormalTok{) }\SpecialCharTok{+}
  \FunctionTok{theme}\NormalTok{(}\AttributeTok{legend.position =} \StringTok{"bottom"}\NormalTok{)}
\end{Highlighting}
\end{Shaded}

\begin{figure}

{\centering \includegraphics[width=0.9\linewidth,alt={A clean, publication-quality network visualisation with minimal visual clutter. High-degree nodes are labelled with author names. Community colours are distinct and colour-blind friendly.}]{scientometrics-in-r_files/figure-latex/20-fig-viz-publication-1} 

}

\caption{Publication-ready co-authorship network with labelled high-degree nodes.}\label{fig:20-fig-viz-publication}
\end{figure}

\section{Diagnostics and interpretation}\label{viz-diagnostics}

\begin{itemize}
\tightlist
\item
  \textbf{Layout quality}: A good layout separates communities visually, minimises edge crossings, and distributes nodes evenly. If communities overlap heavily, try a different algorithm or increase the number of layout iterations.
\item
  \textbf{Readability}: If the network is too dense to read, reduce the number of nodes (filter by degree), increase edge transparency, or remove edges below a weight threshold.
\item
  \textbf{Colour distinguishability}: Test figures in greyscale and with colour-blindness simulators. Viridis-based palettes (used by \texttt{palette\_sci()}) are designed for this.
\item
  \textbf{Export verification}: After exporting, open the file in the target tool (Gephi, VOSviewer) to verify that node attributes, edge weights, and labels transferred correctly.
\end{itemize}

\section{Limitations and responsible use}\label{viz-limitations}

\section{Limitations and responsible use}\label{limitations-and-responsible-use-19}

\begin{itemize}
\tightlist
\item
  \textbf{Visualisation is not analysis.} A network figure is a communication tool. The visual impression can mislead if the layout, filtering, or colour encoding obscures structure. Always report the quantitative metrics alongside figures.
\item
  \textbf{Layout is not geography.} Absolute positions in a force-directed layout are meaningless; only relative distances carry information. Never compare positions across different layouts or runs without the same seed.
\item
  \textbf{Large networks are unreadable.} Beyond \textasciitilde500 nodes, static network plots become unintelligible. Use interactive visualisation, backbone extraction (\ref{backbone-example}), or subnetwork extraction for large data.
\item
  \textbf{Export formats lose information.} GraphML preserves most attributes; Pajek format is more limited. Always verify exports and document what was lost \citep{hicks2015}.
\end{itemize}

\section{Common pitfalls}\label{viz-pitfalls}

\section{Common pitfalls}\label{common-pitfalls-19}

\begin{itemize}
\tightlist
\item
  \textbf{Forgetting \texttt{set.seed()}.} Without a fixed seed, force-directed layouts change every render, making figures irreproducible across HTML and PDF.
\item
  \textbf{Encoding too many variables.} Mapping degree to size, community to colour, \emph{and} betweenness to transparency creates visual overload. Choose two visual channels at most.
\item
  \textbf{Using default igraph plots.} Base R \texttt{plot.igraph()} produces low-quality figures. Use \texttt{ggraph} for publication and \texttt{visNetwork} for interactive exploration.
\item
  \textbf{Exporting without coordinates.} If you export a network to Gephi without coordinates, Gephi will recompute the layout, producing a different figure than your R output.
\end{itemize}

\section{Exercises}\label{viz-exercises}

\begin{enumerate}
\def\labelenumi{\arabic{enumi}.}
\item
  \textbf{Layout comparison.} Apply five different layout algorithms to the same network. Which layout best reveals the community structure? Time each layout and report the trade-off between quality and speed.
\item
  \textbf{Interactive filtering.} Create a \texttt{visNetwork} visualisation where users can filter by community membership. Add a dropdown selector for community.
\item
  \textbf{Gephi round-trip.} Export the network to GraphML, import into Gephi, apply a Force Atlas 2 layout, export as PDF. Compare the Gephi figure with the ggraph version.
\item
  \textbf{Colour-blind check.} Render the network figure with 3, 5, and 10 communities. At what point do the viridis colours become hard to distinguish? Test with a colour-blindness simulator.
\end{enumerate}

\section{Solutions}\label{20-solutions}

Solutions are provided in \ref{solutions}.

\section{Further reading}\label{viz-reading}

\begin{itemize}
\tightlist
\item
  \citet{bastian2009} --- Gephi: an open-source platform for network visualisation and exploration.
\item
  \citet{fortunato2010} --- Community detection methods with extensive visualisation examples.
\item
  \citet{waltman2012} --- VOSviewer's approach to bibliometric network visualisation.
\item
  \citet{aria2017} --- \texttt{bibliometrix} includes network visualisation via \texttt{biblioshiny()}.
\end{itemize}

\section{Session info}\label{viz-sessioninfo}

\begin{Shaded}
\begin{Highlighting}[]
\FunctionTok{sessionInfo}\NormalTok{()}
\end{Highlighting}
\end{Shaded}

\begin{verbatim}
#> R version 4.6.1 (2026-06-24)
#> Platform: x86_64-pc-linux-gnu
#> Running under: Ubuntu 24.04.5 LTS
#> 
#> Matrix products: default
#> BLAS:   /usr/lib/x86_64-linux-gnu/openblas-pthread/libblas.so.3 
#> LAPACK: /usr/lib/x86_64-linux-gnu/openblas-pthread/libopenblasp-r0.3.26.so;  LAPACK version 3.12.0
#> 
#> locale:
#>  [1] LC_CTYPE=C.UTF-8       LC_NUMERIC=C           LC_TIME=C.UTF-8       
#>  [4] LC_COLLATE=C.UTF-8     LC_MONETARY=C.UTF-8    LC_MESSAGES=C.UTF-8   
#>  [7] LC_PAPER=C.UTF-8       LC_NAME=C              LC_ADDRESS=C          
#> [10] LC_TELEPHONE=C         LC_MEASUREMENT=C.UTF-8 LC_IDENTIFICATION=C   
#> 
#> time zone: UTC
#> tzcode source: system (glibc)
#> 
#> attached base packages:
#> [1] stats     graphics  grDevices utils     datasets  methods   base     
#> 
#> other attached packages:
#>  [1] visNetwork_2.1.4   ggraph_2.2.2       tidygraph_1.3.1    igraph_2.3.4      
#>  [5] quanteda_4.5.1     pdftools_3.9.1     arrow_26.0.0       bibliometrix_5.5.0
#>  [9] RefManageR_1.4.0   bib2df_1.1.2.0     rcrossref_1.2.1    gt_1.4.0          
#> [13] tidytext_0.4.3     glue_1.8.1         openalexR_3.1.2    lubridate_1.9.5   
#> [17] forcats_1.0.1      stringr_1.6.0      dplyr_1.2.1        purrr_1.2.2       
#> [21] readr_2.2.0        tidyr_1.3.2        tibble_3.3.1       ggplot2_4.0.3     
#> [25] tidyverse_2.0.0   
#> 
#> loaded via a namespace (and not attached):
#>   [1] bibtex_0.5.3           RColorBrewer_1.1-3     rstudioapi_0.19.0     
#>   [4] jsonlite_2.0.0         magrittr_2.0.5         farver_2.1.2          
#>   [7] rmarkdown_2.32         fs_2.1.0               vctrs_0.7.3           
#>  [10] memoise_2.0.1          askpass_1.2.1          base64enc_0.1-6       
#>  [13] htmltools_0.5.9        contentanalysis_1.1.1  curl_8.1.0            
#>  [16] janeaustenr_1.0.0      cellranger_1.1.0       htmlwidgets_1.6.4     
#>  [19] tokenizers_0.3.0       plyr_1.8.9             httr2_1.3.0           
#>  [22] cachem_1.1.0           plotly_4.12.1          dimensionsR_0.0.3     
#>  [25] mime_0.13              lifecycle_1.0.5        pkgconfig_2.0.3       
#>  [28] Matrix_1.7-5           R6_2.6.1               fastmap_1.2.0         
#>  [31] shiny_1.14.0           digest_0.6.39          patchwork_1.3.2       
#>  [34] shinycssloaders_1.1.0  rprojroot_2.1.1        SnowballC_0.7.1       
#>  [37] labeling_0.4.3         urltools_1.7.3.1       timechange_0.4.0      
#>  [40] httr_1.4.9             polyclip_1.10-7        compiler_4.6.1        
#>  [43] here_1.0.2             bit64_4.8.6            withr_3.0.3           
#>  [46] S7_0.2.2               viridis_0.6.5          ggforce_0.5.0         
#>  [49] MASS_7.3-65            bibliometrixData_0.3.0 tools_4.6.1           
#>  [52] otel_0.2.0             stopwords_2.3          zip_3.0.2             
#>  [55] rentrez_1.2.4          httpuv_1.6.17          promises_1.5.0        
#>  [58] grid_4.6.1             stringdist_0.9.17      generics_0.1.4        
#>  [61] gtable_0.3.6           tzdb_0.5.0             rscopus_0.9.0         
#>  [64] ca_0.72                data.table_1.18.6.1    hms_1.1.4             
#>  [67] xml2_1.6.0             utf8_1.2.6             ggrepel_0.9.8         
#>  [70] pillar_1.11.1          vroom_1.7.1            later_1.4.8           
#>  [73] tweenr_2.0.3           lattice_0.22-9         bit_4.6.0             
#>  [76] tidyselect_1.2.1       miniUI_0.1.2           knitr_1.52            
#>  [79] gridExtra_2.3.1        bookdown_0.48          crul_1.6.0            
#>  [82] xfun_0.62              graphlayouts_1.2.5     DT_0.34.0             
#>  [85] humaniformat_0.6.0     stringi_1.8.9          qpdf_1.4.1            
#>  [88] yaml_2.3.12            evaluate_1.0.5         httpcode_0.3.0        
#>  [91] cli_3.6.6              xtable_1.8-8           systemfonts_1.3.2     
#>  [94] Rcpp_1.1.2             readxl_1.6.0           triebeard_0.4.1       
#>  [97] XML_3.99-0.25          parallel_4.6.1         assertthat_0.2.1      
#> [100] pubmedR_1.0.2          viridisLite_0.4.3      scales_1.4.0          
#> [103] crayon_1.5.3           openxlsx_4.2.9         rlang_1.3.0           
#> [106] fastmatch_1.1-8
\end{verbatim}

\part{Text \& Topic Analysis}\label{part-text-topic-analysis}

\chapter{Text Mining Bibliographic Corpora}\label{text-mining-bibliographic-corpora}

\emph{R. Heller · Date: 2026-05-09 · Last changed: 2026-10-31}

\section{Learning objectives}\label{textmine-objectives}

After completing this chapter, you will be able to:

\begin{itemize}
\tightlist
\item
  Build a text corpus from OpenAlex titles and abstracts
\item
  Apply standard preprocessing steps: tokenisation, stopword removal, stemming
\item
  Construct a document-term matrix (DTM) and a document-feature matrix (DFM)
\item
  Compute TF-IDF weights to identify distinctive terms
\item
  Visualise term frequency patterns across time or groups
\end{itemize}

\section{Setup}\label{textmine-setup}

\begin{Shaded}
\begin{Highlighting}[]
\FunctionTok{library}\NormalTok{(tidyverse)}
\FunctionTok{library}\NormalTok{(openalexR)}
\FunctionTok{library}\NormalTok{(quanteda)}
\FunctionTok{library}\NormalTok{(quanteda.textstats)}
\FunctionTok{library}\NormalTok{(glue)}
\FunctionTok{library}\NormalTok{(gt)}

\FunctionTok{set.seed}\NormalTok{(}\DecValTok{20260509}\NormalTok{)}

\FunctionTok{source}\NormalTok{(here}\SpecialCharTok{::}\FunctionTok{here}\NormalTok{(}\StringTok{"R"}\NormalTok{, }\StringTok{"api\_helpers.R"}\NormalTok{))}
\FunctionTok{source}\NormalTok{(here}\SpecialCharTok{::}\FunctionTok{here}\NormalTok{(}\StringTok{"R"}\NormalTok{, }\StringTok{"utils.R"}\NormalTok{))}
\FunctionTok{source}\NormalTok{(here}\SpecialCharTok{::}\FunctionTok{here}\NormalTok{(}\StringTok{"R"}\NormalTok{, }\StringTok{"sci\_palette.R"}\NormalTok{))}
\end{Highlighting}
\end{Shaded}

\section{Conceptual background}\label{textmine-background}

Bibliometric metadata --- titles, abstracts, and keywords --- constitutes a rich but noisy text corpus. Text mining transforms this unstructured data into structured representations suitable for statistical analysis. The pipeline typically follows three stages: preprocessing, representation, and analysis.

\textbf{Preprocessing} converts raw text into a standardised form. Tokenisation splits text into individual words or n-grams. Lowercasing removes case variation. Stopword removal eliminates high-frequency function words (``the'', ``of'', ``and'') that carry little topical information. Stemming or lemmatisation reduces words to their root forms (``computing'', ``computed'', ``computation'' → ``comput'' or ``compute''). Each step involves trade-offs: aggressive stemming can merge distinct concepts, while minimal preprocessing retains noise.

The \textbf{document-term matrix} (DTM) or \textbf{document-feature matrix} (DFM) is the fundamental representation. Rows are documents; columns are terms; cells contain counts or weights. Raw term frequencies overweight common words. \textbf{TF-IDF} (term frequency--inverse document frequency) addresses this by upweighting terms that are frequent in a document but rare across the corpus, highlighting words that distinguish one document from others.

\texttt{quanteda} \citep{aria2017} provides a fast, well-designed toolkit for text analysis in R. Its \texttt{corpus\ →\ tokens\ →\ dfm} pipeline integrates naturally with the tidyverse. For very large corpora, \texttt{quanteda} uses sparse matrix representations that scale to millions of documents.

Text mining complements network-based methods (\ref{coword-background}). Co-word networks reveal term associations; text mining provides the frequency distributions, temporal trends, and discriminative features that characterise a field's vocabulary.

\section{Worked example}\label{textmine-example}

\subsection{Building a text corpus}\label{building-a-text-corpus}

\begin{Shaded}
\begin{Highlighting}[]
\NormalTok{works }\OtherTok{\textless{}{-}} \FunctionTok{fetch\_openalex}\NormalTok{(}
  \AttributeTok{entity =} \StringTok{"works"}\NormalTok{,}
  \AttributeTok{primary\_location.source.id =} \StringTok{"S148561398"}\NormalTok{,}
  \AttributeTok{from\_publication\_date =} \StringTok{"2019{-}01{-}01"}\NormalTok{,}
  \AttributeTok{to\_publication\_date =} \StringTok{"2023{-}12{-}31"}\NormalTok{,}
  \AttributeTok{options =} \FunctionTok{list}\NormalTok{(}\AttributeTok{sample =} \DecValTok{500}\NormalTok{, }\AttributeTok{seed =} \DecValTok{42}\NormalTok{)}
\NormalTok{)}

\NormalTok{text\_df }\OtherTok{\textless{}{-}}\NormalTok{ works }\SpecialCharTok{|\textgreater{}}
  \FunctionTok{filter}\NormalTok{(}\SpecialCharTok{!}\FunctionTok{is.na}\NormalTok{(abstract), }\FunctionTok{nchar}\NormalTok{(abstract) }\SpecialCharTok{\textgreater{}} \DecValTok{50}\NormalTok{) }\SpecialCharTok{|\textgreater{}}
  \FunctionTok{transmute}\NormalTok{(}
    \AttributeTok{doc\_id =}\NormalTok{ id,}
    \AttributeTok{text =} \FunctionTok{paste}\NormalTok{(display\_name, abstract, }\AttributeTok{sep =} \StringTok{". "}\NormalTok{),}
    \AttributeTok{year =} \FunctionTok{year}\NormalTok{(publication\_date)}
\NormalTok{  )}

\FunctionTok{cat}\NormalTok{(}\FunctionTok{glue}\NormalTok{(}\StringTok{"Documents with abstracts: \{nrow(text\_df)\}}\SpecialCharTok{\textbackslash{}n}\StringTok{"}\NormalTok{))}
\end{Highlighting}
\end{Shaded}

\begin{verbatim}
#> Documents with abstracts: 130
\end{verbatim}

\begin{Shaded}
\begin{Highlighting}[]
\NormalTok{corp }\OtherTok{\textless{}{-}} \FunctionTok{corpus}\NormalTok{(text\_df, }\AttributeTok{docid\_field =} \StringTok{"doc\_id"}\NormalTok{, }\AttributeTok{text\_field =} \StringTok{"text"}\NormalTok{)}
\FunctionTok{docvars}\NormalTok{(corp, }\StringTok{"year"}\NormalTok{) }\OtherTok{\textless{}{-}}\NormalTok{ text\_df}\SpecialCharTok{$}\NormalTok{year}

\FunctionTok{cat}\NormalTok{(}\FunctionTok{glue}\NormalTok{(}\StringTok{"Corpus size: \{ndoc(corp)\} documents}\SpecialCharTok{\textbackslash{}n}\StringTok{"}\NormalTok{))}
\end{Highlighting}
\end{Shaded}

\begin{verbatim}
#> Corpus size: 130 documents
\end{verbatim}

\begin{Shaded}
\begin{Highlighting}[]
\FunctionTok{cat}\NormalTok{(}\FunctionTok{glue}\NormalTok{(}\StringTok{"Total tokens: \{sum(ntoken(corp))\}}\SpecialCharTok{\textbackslash{}n}\StringTok{"}\NormalTok{))}
\end{Highlighting}
\end{Shaded}

\begin{verbatim}
#> Total tokens: 32111
\end{verbatim}

\subsection{Tokenisation and preprocessing}\label{tokenisation-and-preprocessing}

\begin{Shaded}
\begin{Highlighting}[]
\NormalTok{toks }\OtherTok{\textless{}{-}} \FunctionTok{tokens}\NormalTok{(corp,}
               \AttributeTok{remove\_punct =} \ConstantTok{TRUE}\NormalTok{,}
               \AttributeTok{remove\_numbers =} \ConstantTok{TRUE}\NormalTok{,}
               \AttributeTok{remove\_symbols =} \ConstantTok{TRUE}\NormalTok{) }\SpecialCharTok{|\textgreater{}}
  \FunctionTok{tokens\_tolower}\NormalTok{() }\SpecialCharTok{|\textgreater{}}
  \FunctionTok{tokens\_remove}\NormalTok{(}\FunctionTok{stopwords}\NormalTok{(}\StringTok{"en"}\NormalTok{)) }\SpecialCharTok{|\textgreater{}}
  \FunctionTok{tokens\_remove}\NormalTok{(}\FunctionTok{c}\NormalTok{(}\StringTok{"also"}\NormalTok{, }\StringTok{"however"}\NormalTok{, }\StringTok{"using"}\NormalTok{, }\StringTok{"based"}\NormalTok{, }\StringTok{"study"}\NormalTok{,}
                   \StringTok{"paper"}\NormalTok{, }\StringTok{"results"}\NormalTok{, }\StringTok{"research"}\NormalTok{, }\StringTok{"analysis"}\NormalTok{)) }\SpecialCharTok{|\textgreater{}}
  \FunctionTok{tokens\_wordstem}\NormalTok{()}

\FunctionTok{cat}\NormalTok{(}\FunctionTok{glue}\NormalTok{(}\StringTok{"Tokens after preprocessing: \{sum(ntoken(toks))\}}\SpecialCharTok{\textbackslash{}n}\StringTok{"}\NormalTok{))}
\end{Highlighting}
\end{Shaded}

\begin{verbatim}
#> Tokens after preprocessing: 16208
\end{verbatim}

\subsection{Document-feature matrix}\label{document-feature-matrix}

\begin{Shaded}
\begin{Highlighting}[]
\NormalTok{dfmat }\OtherTok{\textless{}{-}} \FunctionTok{dfm}\NormalTok{(toks) }\SpecialCharTok{|\textgreater{}}
  \FunctionTok{dfm\_trim}\NormalTok{(}\AttributeTok{min\_termfreq =} \DecValTok{5}\NormalTok{, }\AttributeTok{min\_docfreq =} \DecValTok{3}\NormalTok{)}

\FunctionTok{cat}\NormalTok{(}\FunctionTok{glue}\NormalTok{(}\StringTok{"DFM dimensions: \{nrow(dfmat)\} docs x \{ncol(dfmat)\} features}\SpecialCharTok{\textbackslash{}n}\StringTok{"}\NormalTok{))}
\end{Highlighting}
\end{Shaded}

\begin{verbatim}
#> DFM dimensions: 130 docs x 675 features
\end{verbatim}

\begin{Shaded}
\begin{Highlighting}[]
\FunctionTok{cat}\NormalTok{(}\FunctionTok{glue}\NormalTok{(}\StringTok{"Sparsity: \{scales::percent(sparsity(dfmat))\}}\SpecialCharTok{\textbackslash{}n}\StringTok{"}\NormalTok{))}
\end{Highlighting}
\end{Shaded}

\begin{verbatim}
#> Sparsity: 91%
\end{verbatim}

\begin{Shaded}
\begin{Highlighting}[]
\NormalTok{top\_terms }\OtherTok{\textless{}{-}} \FunctionTok{topfeatures}\NormalTok{(dfmat, }\DecValTok{20}\NormalTok{) }\SpecialCharTok{|\textgreater{}}
  \FunctionTok{enframe}\NormalTok{(}\AttributeTok{name =} \StringTok{"term"}\NormalTok{, }\AttributeTok{value =} \StringTok{"frequency"}\NormalTok{) }\SpecialCharTok{|\textgreater{}}
  \FunctionTok{mutate}\NormalTok{(}\AttributeTok{term =} \FunctionTok{fct\_reorder}\NormalTok{(term, frequency))}

\FunctionTok{ggplot}\NormalTok{(top\_terms, }\FunctionTok{aes}\NormalTok{(}\AttributeTok{x =}\NormalTok{ frequency, }\AttributeTok{y =}\NormalTok{ term)) }\SpecialCharTok{+}
  \FunctionTok{geom\_col}\NormalTok{(}\AttributeTok{fill =} \FunctionTok{palette\_sci}\NormalTok{(}\DecValTok{1}\NormalTok{)) }\SpecialCharTok{+}
  \FunctionTok{labs}\NormalTok{(}\AttributeTok{x =} \StringTok{"Frequency"}\NormalTok{, }\AttributeTok{y =} \ConstantTok{NULL}\NormalTok{) }\SpecialCharTok{+}
  \FunctionTok{theme\_sci}\NormalTok{()}
\end{Highlighting}
\end{Shaded}

\begin{figure}

{\centering \includegraphics[width=0.9\linewidth,alt={Horizontal bar chart showing the 20 most frequent terms after preprocessing, with term frequency on the x-axis.}]{scientometrics-in-r_files/figure-latex/21-fig-textmine-topfeatures-1} 

}

\caption{Top 20 most frequent terms in the Scientometrics corpus.}\label{fig:21-fig-textmine-topfeatures}
\end{figure}

\subsection{TF-IDF weighting}\label{tf-idf-weighting}

\begin{Shaded}
\begin{Highlighting}[]
\NormalTok{dfmat\_tfidf }\OtherTok{\textless{}{-}} \FunctionTok{dfm\_tfidf}\NormalTok{(dfmat)}

\NormalTok{tfidf\_by\_year }\OtherTok{\textless{}{-}}\NormalTok{ text\_df }\SpecialCharTok{|\textgreater{}}
  \FunctionTok{group\_by}\NormalTok{(year) }\SpecialCharTok{|\textgreater{}}
  \FunctionTok{group\_keys}\NormalTok{() }\SpecialCharTok{|\textgreater{}}
  \FunctionTok{pull}\NormalTok{(year) }\SpecialCharTok{|\textgreater{}}
  \FunctionTok{map\_dfr}\NormalTok{(}\ControlFlowTok{function}\NormalTok{(yr) \{}
\NormalTok{    docs }\OtherTok{\textless{}{-}} \FunctionTok{docvars}\NormalTok{(dfmat\_tfidf, }\StringTok{"year"}\NormalTok{) }\SpecialCharTok{==}\NormalTok{ yr}
    \ControlFlowTok{if}\NormalTok{ (}\FunctionTok{sum}\NormalTok{(docs) }\SpecialCharTok{\textless{}} \DecValTok{5}\NormalTok{) }\FunctionTok{return}\NormalTok{(}\FunctionTok{tibble}\NormalTok{())}
\NormalTok{    top }\OtherTok{\textless{}{-}} \FunctionTok{topfeatures}\NormalTok{(dfmat\_tfidf[docs, ], }\DecValTok{10}\NormalTok{)}
    \FunctionTok{tibble}\NormalTok{(}\AttributeTok{year =}\NormalTok{ yr, }\AttributeTok{term =} \FunctionTok{names}\NormalTok{(top), }\AttributeTok{tfidf =} \FunctionTok{unname}\NormalTok{(top))}
\NormalTok{  \})}

\NormalTok{tfidf\_by\_year }\SpecialCharTok{|\textgreater{}}
  \FunctionTok{group\_by}\NormalTok{(year) }\SpecialCharTok{|\textgreater{}}
  \FunctionTok{slice\_max}\NormalTok{(tfidf, }\AttributeTok{n =} \DecValTok{5}\NormalTok{) }\SpecialCharTok{|\textgreater{}}
  \FunctionTok{gt}\NormalTok{() }\SpecialCharTok{|\textgreater{}}
  \FunctionTok{fmt\_number}\NormalTok{(}\AttributeTok{columns =}\NormalTok{ tfidf, }\AttributeTok{decimals =} \DecValTok{1}\NormalTok{)}
\end{Highlighting}
\end{Shaded}

\begingroup
\fontsize{8.0pt}{9.0pt}\selectfont
\begin{longtable}{>{\raggedright\arraybackslash}p{\dimexpr 0.50\linewidth -2\tabcolsep-1.5\arrayrulewidth}>{\raggedleft\arraybackslash}p{\dimexpr 0.50\linewidth -2\tabcolsep-1.5\arrayrulewidth}}
\toprule
term & tfidf \\ 
\midrule\addlinespace[2.5pt]
\multicolumn{2}{l}{2019} \\[2.5pt] 
\midrule\addlinespace[2.5pt]
journal & 20.2 \\ 
econom & 18.6 \\ 
\nolinkurl{manuscript} & 16.6 \\ 
score & 16.6 \\ 
doctor & 16.4 \\ 
\midrule\addlinespace[2.5pt]
\multicolumn{2}{l}{2020} \\[2.5pt] 
\midrule\addlinespace[2.5pt]
lis & 24.6 \\ 
citat & 19.5 \\ 
\nolinkurl{classifi} & 15.7 \\ 
machin & 13.9 \\ 
practic & 13.6 \\ 
\midrule\addlinespace[2.5pt]
\multicolumn{2}{l}{2021} \\[2.5pt] 
\midrule\addlinespace[2.5pt]
retract & 40.8 \\ 
journal & 28.9 \\ 
book & 24.6 \\ 
\nolinkurl{sars-cov-2} & 24.2 \\ 
review & 23.4 \\ 
\midrule\addlinespace[2.5pt]
\multicolumn{2}{l}{2022} \\[2.5pt] 
\midrule\addlinespace[2.5pt]
journal & 29.5 \\ 
p & 25.7 \\ 
public & 25.6 \\ 
citat & 25.5 \\ 
problem & 24.4 \\ 
\midrule\addlinespace[2.5pt]
\multicolumn{2}{l}{2023} \\[2.5pt] 
\midrule\addlinespace[2.5pt]
mention & 31.7 \\ 
polici & 22.0 \\ 
\nolinkurl{document} & 14.5 \\ 
academ & 13.1 \\ 
open & 12.7 \\ 
\bottomrule
\end{longtable}
\endgroup

\subsection{Term frequency over time}\label{term-frequency-over-time}

\begin{Shaded}
\begin{Highlighting}[]
\NormalTok{target\_terms }\OtherTok{\textless{}{-}} \FunctionTok{c}\NormalTok{(}\StringTok{"bibliometr"}\NormalTok{, }\StringTok{"open"}\NormalTok{, }\StringTok{"network"}\NormalTok{, }\StringTok{"impact"}\NormalTok{, }\StringTok{"collabor"}\NormalTok{)}

\NormalTok{freq\_by\_year }\OtherTok{\textless{}{-}} \FunctionTok{map\_dfr}\NormalTok{(}\FunctionTok{unique}\NormalTok{(text\_df}\SpecialCharTok{$}\NormalTok{year), }\ControlFlowTok{function}\NormalTok{(yr) \{}
\NormalTok{  docs }\OtherTok{\textless{}{-}} \FunctionTok{docvars}\NormalTok{(dfmat, }\StringTok{"year"}\NormalTok{) }\SpecialCharTok{==}\NormalTok{ yr}
  \ControlFlowTok{if}\NormalTok{ (}\FunctionTok{sum}\NormalTok{(docs) }\SpecialCharTok{\textless{}} \DecValTok{5}\NormalTok{) }\FunctionTok{return}\NormalTok{(}\FunctionTok{tibble}\NormalTok{())}
\NormalTok{  freq }\OtherTok{\textless{}{-}} \FunctionTok{colSums}\NormalTok{(dfmat[docs, ]) }\SpecialCharTok{/} \FunctionTok{sum}\NormalTok{(dfmat[docs, ])}
  \FunctionTok{tibble}\NormalTok{(}
    \AttributeTok{year =}\NormalTok{ yr,}
    \AttributeTok{term =}\NormalTok{ target\_terms,}
    \AttributeTok{rel\_freq =}\NormalTok{ freq[target\_terms]}
\NormalTok{  )}
\NormalTok{\})}

\FunctionTok{ggplot}\NormalTok{(freq\_by\_year, }\FunctionTok{aes}\NormalTok{(}\AttributeTok{x =}\NormalTok{ year, }\AttributeTok{y =}\NormalTok{ rel\_freq, }\AttributeTok{colour =}\NormalTok{ term)) }\SpecialCharTok{+}
  \FunctionTok{geom\_line}\NormalTok{(}\AttributeTok{linewidth =} \DecValTok{1}\NormalTok{) }\SpecialCharTok{+}
  \FunctionTok{geom\_point}\NormalTok{(}\AttributeTok{size =} \DecValTok{2}\NormalTok{) }\SpecialCharTok{+}
  \FunctionTok{scale\_colour\_manual}\NormalTok{(}\AttributeTok{values =} \FunctionTok{palette\_sci}\NormalTok{(}\FunctionTok{length}\NormalTok{(target\_terms))) }\SpecialCharTok{+}
  \FunctionTok{scale\_y\_continuous}\NormalTok{(}\AttributeTok{labels =}\NormalTok{ scales}\SpecialCharTok{::}\FunctionTok{percent\_format}\NormalTok{(}\AttributeTok{accuracy =} \FloatTok{0.01}\NormalTok{)) }\SpecialCharTok{+}
  \FunctionTok{labs}\NormalTok{(}\AttributeTok{x =} \StringTok{"Year"}\NormalTok{, }\AttributeTok{y =} \StringTok{"Relative frequency"}\NormalTok{, }\AttributeTok{colour =} \StringTok{"Term"}\NormalTok{) }\SpecialCharTok{+}
  \FunctionTok{theme\_sci}\NormalTok{()}
\end{Highlighting}
\end{Shaded}

\begin{figure}

{\centering \includegraphics[width=0.9\linewidth,alt={Line chart showing how the relative frequency of key terms changes across publication years, revealing shifts in research vocabulary.}]{scientometrics-in-r_files/figure-latex/21-fig-textmine-temporal-1} 

}

\caption{Relative frequency of selected terms over time.}\label{fig:21-fig-textmine-temporal}
\end{figure}

\section{Diagnostics and interpretation}\label{textmine-diagnostics}

\begin{itemize}
\tightlist
\item
  \textbf{Vocabulary size vs.~documents}: A typical DFM has far more features than documents. Trim aggressively (minimum document frequency of 3--5) to reduce noise and computation.
\item
  \textbf{Sparsity}: Bibliometric DFMs are typically 95--99\% sparse. This is normal. Sparse matrix representations keep memory usage manageable.
\item
  \textbf{Preprocessing sensitivity}: Results change with preprocessing choices. Always report stopword list, stemming method, and trimming thresholds.
\item
  \textbf{Abstract availability}: Not all OpenAlex records have abstracts. Report the coverage rate and consider whether missing abstracts introduce bias (e.g., older papers, certain publishers).
\end{itemize}

\section{Limitations and responsible use}\label{textmine-limitations}

\section{Limitations and responsible use}\label{limitations-and-responsible-use-20}

\begin{itemize}
\tightlist
\item
  \textbf{Abstracts are summaries.} They capture the main claims but miss nuance, methodology details, and negative results. Full-text analysis (\ref{fulltext-background}) provides richer data but is harder to obtain.
\item
  \textbf{Language bias.} Text mining tools work best for English. Non-English abstracts may be poorly tokenised, incorrectly stemmed, or excluded entirely. This biases results toward Anglophone research \citep{visser2021}.
\item
  \textbf{Stemming conflates concepts.} ``Stem'' and ``stemming'' are merged, but so are ``cell'' (biology) and ``cell'' (spreadsheet). Lemmatisation is more precise but slower.
\item
  \textbf{Bag-of-words ignores context.} TF-IDF and frequency-based methods treat documents as unordered collections of words. Phrase meaning (``machine learning'' vs.~``learning machine'') is lost unless you use n-grams or embeddings (\ref{embed-background}).
\end{itemize}

\section{Common pitfalls}\label{textmine-pitfalls}

\section{Common pitfalls}\label{common-pitfalls-20}

\begin{itemize}
\tightlist
\item
  \textbf{Not removing domain-specific stop words.} Generic stopword lists miss terms like ``study'', ``results'', ``paper'' that dominate academic text without conveying topical information.
\item
  \textbf{Applying TF-IDF before trimming.} Rare terms get extreme TF-IDF scores. Trim the DFM first, then apply TF-IDF.
\item
  \textbf{Comparing raw frequencies across groups of different sizes.} A year with 200 papers will have higher raw counts than a year with 50. Always normalise to relative frequency.
\item
  \textbf{Stemming titles but not abstracts (or vice versa).} Apply identical preprocessing to all text fields.
\end{itemize}

\section{Exercises}\label{textmine-exercises}

\begin{enumerate}
\def\labelenumi{\arabic{enumi}.}
\item
  \textbf{Bigram analysis.} Build a DFM using bigrams (two-word sequences) instead of unigrams. What meaningful phrases emerge that single words miss?
\item
  \textbf{Lexical diversity.} Compute the type-token ratio (TTR) for each publication year. Is the vocabulary growing more diverse or more standardised over time?
\item
  \textbf{Keyness analysis.} Use \texttt{quanteda.textstats::textstat\_keyness()} to identify terms that distinguish one year from all others. What terms characterise 2023 specifically?
\item
  \textbf{Preprocessing sensitivity.} Run the same analysis with and without stemming. How do the top terms differ? Which version is more interpretable?
\end{enumerate}

\section{Solutions}\label{21-solutions}

Solutions are provided in \ref{solutions}.

\section{Further reading}\label{textmine-reading}

\begin{itemize}
\tightlist
\item
  \citet{silge2017} --- \emph{Text Mining with R}; comprehensive guide to tidy text analysis.
\item
  \citet{aria2017} --- \texttt{bibliometrix} text-mining features for bibliometric data.
\item
  \citet{priem2022} --- OpenAlex abstract availability and coverage.
\end{itemize}

\section{Session info}\label{textmine-sessioninfo}

\begin{Shaded}
\begin{Highlighting}[]
\FunctionTok{sessionInfo}\NormalTok{()}
\end{Highlighting}
\end{Shaded}

\begin{verbatim}
#> R version 4.6.1 (2026-06-24)
#> Platform: x86_64-pc-linux-gnu
#> Running under: Ubuntu 24.04.5 LTS
#> 
#> Matrix products: default
#> BLAS:   /usr/lib/x86_64-linux-gnu/openblas-pthread/libblas.so.3 
#> LAPACK: /usr/lib/x86_64-linux-gnu/openblas-pthread/libopenblasp-r0.3.26.so;  LAPACK version 3.12.0
#> 
#> locale:
#>  [1] LC_CTYPE=C.UTF-8       LC_NUMERIC=C           LC_TIME=C.UTF-8       
#>  [4] LC_COLLATE=C.UTF-8     LC_MONETARY=C.UTF-8    LC_MESSAGES=C.UTF-8   
#>  [7] LC_PAPER=C.UTF-8       LC_NAME=C              LC_ADDRESS=C          
#> [10] LC_TELEPHONE=C         LC_MEASUREMENT=C.UTF-8 LC_IDENTIFICATION=C   
#> 
#> time zone: UTC
#> tzcode source: system (glibc)
#> 
#> attached base packages:
#> [1] stats     graphics  grDevices utils     datasets  methods   base     
#> 
#> other attached packages:
#>  [1] quanteda.textstats_0.97.2 visNetwork_2.1.4         
#>  [3] ggraph_2.2.2              tidygraph_1.3.1          
#>  [5] igraph_2.3.4              quanteda_4.5.1           
#>  [7] pdftools_3.9.1            arrow_26.0.0             
#>  [9] bibliometrix_5.5.0        RefManageR_1.4.0         
#> [11] bib2df_1.1.2.0            rcrossref_1.2.1          
#> [13] gt_1.4.0                  tidytext_0.4.3           
#> [15] glue_1.8.1                openalexR_3.1.2          
#> [17] lubridate_1.9.5           forcats_1.0.1            
#> [19] stringr_1.6.0             dplyr_1.2.1              
#> [21] purrr_1.2.2               readr_2.2.0              
#> [23] tidyr_1.3.2               tibble_3.3.1             
#> [25] ggplot2_4.0.3             tidyverse_2.0.0          
#> 
#> loaded via a namespace (and not attached):
#>   [1] bibtex_0.5.3           RColorBrewer_1.1-3     rstudioapi_0.19.0     
#>   [4] jsonlite_2.0.0         magrittr_2.0.5         farver_2.1.2          
#>   [7] rmarkdown_2.32         fs_2.1.0               vctrs_0.7.3           
#>  [10] memoise_2.0.1          askpass_1.2.1          base64enc_0.1-6       
#>  [13] htmltools_0.5.9        contentanalysis_1.1.1  curl_8.1.0            
#>  [16] janeaustenr_1.0.0      cellranger_1.1.0       htmlwidgets_1.6.4     
#>  [19] tokenizers_0.3.0       plyr_1.8.9             httr2_1.3.0           
#>  [22] cachem_1.1.0           plotly_4.12.1          dimensionsR_0.0.3     
#>  [25] mime_0.13              lifecycle_1.0.5        pkgconfig_2.0.3       
#>  [28] Matrix_1.7-5           R6_2.6.1               fastmap_1.2.0         
#>  [31] shiny_1.14.0           digest_0.6.39          patchwork_1.3.2       
#>  [34] shinycssloaders_1.1.0  rprojroot_2.1.1        SnowballC_0.7.1       
#>  [37] labeling_0.4.3         urltools_1.7.3.1       timechange_0.4.0      
#>  [40] httr_1.4.9             polyclip_1.10-7        compiler_4.6.1        
#>  [43] here_1.0.2             bit64_4.8.6            withr_3.0.3           
#>  [46] S7_0.2.2               viridis_0.6.5          ggforce_0.5.0         
#>  [49] MASS_7.3-65            bibliometrixData_0.3.0 tools_4.6.1           
#>  [52] otel_0.2.0             stopwords_2.3          zip_3.0.2             
#>  [55] rentrez_1.2.4          httpuv_1.6.17          promises_1.5.0        
#>  [58] grid_4.6.1             stringdist_0.9.17      generics_0.1.4        
#>  [61] gtable_0.3.6           tzdb_0.5.0             rscopus_0.9.0         
#>  [64] ca_0.72                data.table_1.18.6.1    hms_1.1.4             
#>  [67] xml2_1.6.0             utf8_1.2.6             ggrepel_0.9.8         
#>  [70] pillar_1.11.1          nsyllable_1.0.1        vroom_1.7.1           
#>  [73] later_1.4.8            tweenr_2.0.3           lattice_0.22-9        
#>  [76] bit_4.6.0              tidyselect_1.2.1       miniUI_0.1.2          
#>  [79] knitr_1.52             gridExtra_2.3.1        bookdown_0.48         
#>  [82] crul_1.6.0             xfun_0.62              graphlayouts_1.2.5    
#>  [85] DT_0.34.0              humaniformat_0.6.0     stringi_1.8.9         
#>  [88] qpdf_1.4.1             yaml_2.3.12            evaluate_1.0.5        
#>  [91] httpcode_0.3.0         cli_3.6.6              xtable_1.8-8          
#>  [94] systemfonts_1.3.2      Rcpp_1.1.2             readxl_1.6.0          
#>  [97] triebeard_0.4.1        XML_3.99-0.25          parallel_4.6.1        
#> [100] assertthat_0.2.1       pubmedR_1.0.2          viridisLite_0.4.3     
#> [103] scales_1.4.0           crayon_1.5.3           openxlsx_4.2.9        
#> [106] rlang_1.3.0            fastmatch_1.1-8
\end{verbatim}

\chapter{Topic Modeling}\label{topic-modeling}

\emph{R. Heller · Date: 2026-05-09 · Last changed: 2026-10-31}

\section{Learning objectives}\label{topicmodel-objectives}

After completing this chapter, you will be able to:

\begin{itemize}
\tightlist
\item
  Explain the generative assumptions of LDA and STM
\item
  Prepare a document-feature matrix suitable for topic modelling
\item
  Fit LDA and STM models and select the number of topics
\item
  Evaluate topic quality using coherence and exclusivity metrics
\item
  Interpret and label topics using top words and representative documents
\end{itemize}

\section{Setup}\label{topicmodel-setup}

\begin{Shaded}
\begin{Highlighting}[]
\FunctionTok{library}\NormalTok{(tidyverse)}
\FunctionTok{library}\NormalTok{(openalexR)}
\FunctionTok{library}\NormalTok{(quanteda)}
\FunctionTok{library}\NormalTok{(topicmodels)}
\FunctionTok{library}\NormalTok{(tidytext)}
\FunctionTok{library}\NormalTok{(stm)}
\FunctionTok{library}\NormalTok{(glue)}
\FunctionTok{library}\NormalTok{(gt)}

\FunctionTok{set.seed}\NormalTok{(}\DecValTok{20260509}\NormalTok{)}

\FunctionTok{source}\NormalTok{(here}\SpecialCharTok{::}\FunctionTok{here}\NormalTok{(}\StringTok{"R"}\NormalTok{, }\StringTok{"api\_helpers.R"}\NormalTok{))}
\FunctionTok{source}\NormalTok{(here}\SpecialCharTok{::}\FunctionTok{here}\NormalTok{(}\StringTok{"R"}\NormalTok{, }\StringTok{"utils.R"}\NormalTok{))}
\FunctionTok{source}\NormalTok{(here}\SpecialCharTok{::}\FunctionTok{here}\NormalTok{(}\StringTok{"R"}\NormalTok{, }\StringTok{"sci\_palette.R"}\NormalTok{))}
\end{Highlighting}
\end{Shaded}

\section{Conceptual background}\label{topicmodel-background}

A \textbf{topic model} is an unsupervised statistical model that discovers latent themes (topics) in a collection of documents. Each topic is a probability distribution over words; each document is a mixture of topics. The model ``reverse-engineers'' the writing process: given only the observed word frequencies, it infers which topics exist and how they are distributed across documents.

\textbf{Latent Dirichlet Allocation} (LDA) \citep{blei2003} is the foundational topic model. It assumes a generative process: for each document, draw a topic mixture from a Dirichlet prior; for each word position, draw a topic from that mixture, then draw a word from the topic's word distribution. Inference (fitting the model) recovers the topic-word and document-topic distributions that best explain the observed data.

The \textbf{Structural Topic Model} (STM) \citep{roberts2019} extends LDA by allowing document-level covariates (e.g., publication year, journal) to influence topic prevalence and word usage. This is valuable for bibliometrics because we expect topics to vary systematically with time, discipline, and publication venue. STM also estimates a ``content'' component that lets word usage within a topic vary by covariate.

\textbf{Choosing the number of topics (K)} is the most consequential modelling decision. Too few topics produce overly broad themes; too many fragment coherent themes into noise. Quantitative approaches include held-out log-likelihood, semantic coherence, and exclusivity metrics. In practice, running models at multiple K values and inspecting the resulting topics is essential.

Topic models have been widely used in scientometrics to map research fields, track thematic evolution, identify emerging areas, and compare the topical focus of journals, institutions, or countries.

\section{Worked example}\label{topicmodel-example}

\subsection{Preparing the corpus}\label{preparing-the-corpus}

\begin{Shaded}
\begin{Highlighting}[]
\NormalTok{works }\OtherTok{\textless{}{-}} \FunctionTok{fetch\_openalex}\NormalTok{(}
  \AttributeTok{entity =} \StringTok{"works"}\NormalTok{,}
  \AttributeTok{primary\_location.source.id =} \StringTok{"S148561398"}\NormalTok{,}
  \AttributeTok{from\_publication\_date =} \StringTok{"2018{-}01{-}01"}\NormalTok{,}
  \AttributeTok{to\_publication\_date =} \StringTok{"2023{-}12{-}31"}\NormalTok{,}
  \AttributeTok{options =} \FunctionTok{list}\NormalTok{(}\AttributeTok{sample =} \DecValTok{400}\NormalTok{, }\AttributeTok{seed =} \DecValTok{42}\NormalTok{)}
\NormalTok{)}

\NormalTok{text\_df }\OtherTok{\textless{}{-}}\NormalTok{ works }\SpecialCharTok{|\textgreater{}}
  \FunctionTok{filter}\NormalTok{(}\SpecialCharTok{!}\FunctionTok{is.na}\NormalTok{(abstract), }\FunctionTok{nchar}\NormalTok{(abstract) }\SpecialCharTok{\textgreater{}} \DecValTok{100}\NormalTok{) }\SpecialCharTok{|\textgreater{}}
  \FunctionTok{transmute}\NormalTok{(}
    \AttributeTok{doc\_id =}\NormalTok{ id,}
    \AttributeTok{text =} \FunctionTok{paste}\NormalTok{(display\_name, abstract, }\AttributeTok{sep =} \StringTok{". "}\NormalTok{),}
    \AttributeTok{year =} \FunctionTok{year}\NormalTok{(publication\_date)}
\NormalTok{  )}

\FunctionTok{cat}\NormalTok{(}\FunctionTok{glue}\NormalTok{(}\StringTok{"Documents for topic modelling: \{nrow(text\_df)\}}\SpecialCharTok{\textbackslash{}n}\StringTok{"}\NormalTok{))}
\end{Highlighting}
\end{Shaded}

\begin{verbatim}
#> Documents for topic modelling: 102
\end{verbatim}

\begin{Shaded}
\begin{Highlighting}[]
\NormalTok{corp }\OtherTok{\textless{}{-}} \FunctionTok{corpus}\NormalTok{(text\_df, }\AttributeTok{docid\_field =} \StringTok{"doc\_id"}\NormalTok{, }\AttributeTok{text\_field =} \StringTok{"text"}\NormalTok{)}
\FunctionTok{docvars}\NormalTok{(corp, }\StringTok{"year"}\NormalTok{) }\OtherTok{\textless{}{-}}\NormalTok{ text\_df}\SpecialCharTok{$}\NormalTok{year}

\NormalTok{toks }\OtherTok{\textless{}{-}} \FunctionTok{tokens}\NormalTok{(corp, }\AttributeTok{remove\_punct =} \ConstantTok{TRUE}\NormalTok{, }\AttributeTok{remove\_numbers =} \ConstantTok{TRUE}\NormalTok{) }\SpecialCharTok{|\textgreater{}}
  \FunctionTok{tokens\_tolower}\NormalTok{() }\SpecialCharTok{|\textgreater{}}
  \FunctionTok{tokens\_remove}\NormalTok{(}\FunctionTok{stopwords}\NormalTok{(}\StringTok{"en"}\NormalTok{)) }\SpecialCharTok{|\textgreater{}}
  \FunctionTok{tokens\_remove}\NormalTok{(}\FunctionTok{c}\NormalTok{(}\StringTok{"study"}\NormalTok{, }\StringTok{"paper"}\NormalTok{, }\StringTok{"results"}\NormalTok{, }\StringTok{"research"}\NormalTok{, }\StringTok{"analysis"}\NormalTok{,}
                   \StringTok{"also"}\NormalTok{, }\StringTok{"however"}\NormalTok{, }\StringTok{"using"}\NormalTok{, }\StringTok{"based"}\NormalTok{))}

\NormalTok{dfmat }\OtherTok{\textless{}{-}} \FunctionTok{dfm}\NormalTok{(toks) }\SpecialCharTok{|\textgreater{}}
  \FunctionTok{dfm\_trim}\NormalTok{(}\AttributeTok{min\_termfreq =} \DecValTok{5}\NormalTok{, }\AttributeTok{min\_docfreq =} \DecValTok{3}\NormalTok{)}

\FunctionTok{cat}\NormalTok{(}\FunctionTok{glue}\NormalTok{(}\StringTok{"DFM: \{nrow(dfmat)\} docs x \{ncol(dfmat)\} features}\SpecialCharTok{\textbackslash{}n}\StringTok{"}\NormalTok{))}
\end{Highlighting}
\end{Shaded}

\begin{verbatim}
#> DFM: 102 docs x 576 features
\end{verbatim}

\subsection{Fitting LDA}\label{fitting-lda}

\begin{Shaded}
\begin{Highlighting}[]
\NormalTok{dtm }\OtherTok{\textless{}{-}}\NormalTok{ quanteda}\SpecialCharTok{::}\FunctionTok{convert}\NormalTok{(dfmat, }\AttributeTok{to =} \StringTok{"topicmodels"}\NormalTok{)}

\NormalTok{lda\_model }\OtherTok{\textless{}{-}} \FunctionTok{LDA}\NormalTok{(dtm, }\AttributeTok{k =} \DecValTok{8}\NormalTok{, }\AttributeTok{control =} \FunctionTok{list}\NormalTok{(}\AttributeTok{seed =} \DecValTok{42}\NormalTok{))}

\NormalTok{lda\_topics }\OtherTok{\textless{}{-}} \FunctionTok{tidy}\NormalTok{(lda\_model, }\AttributeTok{matrix =} \StringTok{"beta"}\NormalTok{) }\SpecialCharTok{|\textgreater{}}
  \FunctionTok{rename}\NormalTok{(}\AttributeTok{topic =}\NormalTok{ topic, }\AttributeTok{term =}\NormalTok{ term, }\AttributeTok{beta =}\NormalTok{ beta)}

\FunctionTok{cat}\NormalTok{(}\FunctionTok{glue}\NormalTok{(}\StringTok{"LDA fitted with K = 8 topics}\SpecialCharTok{\textbackslash{}n}\StringTok{"}\NormalTok{))}
\end{Highlighting}
\end{Shaded}

\begin{verbatim}
#> LDA fitted with K = 8 topics
\end{verbatim}

\begin{Shaded}
\begin{Highlighting}[]
\NormalTok{lda\_top\_terms }\OtherTok{\textless{}{-}}\NormalTok{ lda\_topics }\SpecialCharTok{|\textgreater{}}
  \FunctionTok{group\_by}\NormalTok{(topic) }\SpecialCharTok{|\textgreater{}}
  \FunctionTok{slice\_max}\NormalTok{(beta, }\AttributeTok{n =} \DecValTok{10}\NormalTok{) }\SpecialCharTok{|\textgreater{}}
  \FunctionTok{ungroup}\NormalTok{() }\SpecialCharTok{|\textgreater{}}
  \FunctionTok{mutate}\NormalTok{(}\AttributeTok{term =} \FunctionTok{reorder\_within}\NormalTok{(term, beta, topic))}

\FunctionTok{ggplot}\NormalTok{(lda\_top\_terms, }\FunctionTok{aes}\NormalTok{(}\AttributeTok{x =}\NormalTok{ beta, }\AttributeTok{y =}\NormalTok{ term, }\AttributeTok{fill =} \FunctionTok{factor}\NormalTok{(topic))) }\SpecialCharTok{+}
  \FunctionTok{geom\_col}\NormalTok{(}\AttributeTok{show.legend =} \ConstantTok{FALSE}\NormalTok{) }\SpecialCharTok{+}
  \FunctionTok{facet\_wrap}\NormalTok{(}\SpecialCharTok{\textasciitilde{}} \FunctionTok{paste}\NormalTok{(}\StringTok{"Topic"}\NormalTok{, topic), }\AttributeTok{scales =} \StringTok{"free\_y"}\NormalTok{, }\AttributeTok{ncol =} \DecValTok{4}\NormalTok{) }\SpecialCharTok{+}
  \FunctionTok{scale\_y\_reordered}\NormalTok{() }\SpecialCharTok{+}
  \FunctionTok{scale\_fill\_manual}\NormalTok{(}\AttributeTok{values =} \FunctionTok{palette\_sci}\NormalTok{(}\DecValTok{8}\NormalTok{)) }\SpecialCharTok{+}
  \FunctionTok{labs}\NormalTok{(}\AttributeTok{x =} \StringTok{"Word probability (beta)"}\NormalTok{, }\AttributeTok{y =} \ConstantTok{NULL}\NormalTok{) }\SpecialCharTok{+}
  \FunctionTok{theme\_sci}\NormalTok{(}\AttributeTok{base\_size =} \DecValTok{9}\NormalTok{)}
\end{Highlighting}
\end{Shaded}

\begin{figure}

{\centering \includegraphics[width=0.9\linewidth,alt={Faceted bar chart showing the 10 highest-probability terms for each of the 8 LDA topics. Each facet represents one topic.}]{scientometrics-in-r_files/figure-latex/22-fig-topicmodel-lda-1} 

}

\caption{Top 10 terms per LDA topic.}\label{fig:22-fig-topicmodel-lda}
\end{figure}

\subsection{Fitting STM with year covariate}\label{fitting-stm-with-year-covariate}

\begin{Shaded}
\begin{Highlighting}[]
\NormalTok{stm\_dfm }\OtherTok{\textless{}{-}}\NormalTok{ quanteda}\SpecialCharTok{::}\FunctionTok{convert}\NormalTok{(dfmat, }\AttributeTok{to =} \StringTok{"stm"}\NormalTok{)}

\NormalTok{stm\_model }\OtherTok{\textless{}{-}} \FunctionTok{stm}\NormalTok{(}
  \AttributeTok{documents =}\NormalTok{ stm\_dfm}\SpecialCharTok{$}\NormalTok{documents,}
  \AttributeTok{vocab =}\NormalTok{ stm\_dfm}\SpecialCharTok{$}\NormalTok{vocab,}
  \AttributeTok{K =} \DecValTok{8}\NormalTok{,}
  \AttributeTok{prevalence =} \SpecialCharTok{\textasciitilde{}}\NormalTok{ year,}
  \AttributeTok{data =}\NormalTok{ stm\_dfm}\SpecialCharTok{$}\NormalTok{meta,}
  \AttributeTok{seed =} \DecValTok{42}\NormalTok{,}
  \AttributeTok{verbose =} \ConstantTok{FALSE}
\NormalTok{)}

\FunctionTok{cat}\NormalTok{(}\FunctionTok{glue}\NormalTok{(}\StringTok{"STM fitted with K = 8 topics, year as prevalence covariate}\SpecialCharTok{\textbackslash{}n}\StringTok{"}\NormalTok{))}
\end{Highlighting}
\end{Shaded}

\begin{verbatim}
#> STM fitted with K = 8 topics, year as prevalence covariate
\end{verbatim}

\begin{Shaded}
\begin{Highlighting}[]
\FunctionTok{plot}\NormalTok{(stm\_model, }\AttributeTok{type =} \StringTok{"summary"}\NormalTok{, }\AttributeTok{n =} \DecValTok{5}\NormalTok{,}
     \AttributeTok{main =} \StringTok{"STM Topic Proportions and Top FREX Terms"}\NormalTok{)}
\end{Highlighting}
\end{Shaded}

\begin{figure}

{\centering \includegraphics[width=0.9\linewidth,alt={Labelled list showing the top FREX terms for each of 8 STM topics. FREX terms are both frequent within and exclusive to each topic.}]{scientometrics-in-r_files/figure-latex/22-fig-topicmodel-stm-1} 

}

\caption{Top terms per STM topic, showing FREX (frequent and exclusive) words.}\label{fig:22-fig-topicmodel-stm}
\end{figure}

\subsection{Topic prevalence over time}\label{topic-prevalence-over-time}

\begin{Shaded}
\begin{Highlighting}[]
\NormalTok{effect }\OtherTok{\textless{}{-}} \FunctionTok{estimateEffect}\NormalTok{(}\DecValTok{1}\SpecialCharTok{:}\DecValTok{8} \SpecialCharTok{\textasciitilde{}}\NormalTok{ year, stm\_model, }\AttributeTok{meta =}\NormalTok{ stm\_dfm}\SpecialCharTok{$}\NormalTok{meta)}

\NormalTok{effect\_df }\OtherTok{\textless{}{-}} \FunctionTok{map\_dfr}\NormalTok{(}\DecValTok{1}\SpecialCharTok{:}\DecValTok{8}\NormalTok{, }\ControlFlowTok{function}\NormalTok{(k) \{}
\NormalTok{  eff }\OtherTok{\textless{}{-}} \FunctionTok{summary}\NormalTok{(effect, }\AttributeTok{topics =}\NormalTok{ k)}\SpecialCharTok{$}\NormalTok{tables[[}\DecValTok{1}\NormalTok{]]}
  \FunctionTok{tibble}\NormalTok{(}
    \AttributeTok{topic =}\NormalTok{ k,}
    \AttributeTok{term =} \FunctionTok{c}\NormalTok{(}\StringTok{"intercept"}\NormalTok{, }\StringTok{"year"}\NormalTok{),}
    \AttributeTok{estimate =}\NormalTok{ eff[, }\StringTok{"Estimate"}\NormalTok{],}
    \AttributeTok{se =}\NormalTok{ eff[, }\StringTok{"Std. Error"}\NormalTok{]}
\NormalTok{  )}
\NormalTok{\}) }\SpecialCharTok{|\textgreater{}}
  \FunctionTok{filter}\NormalTok{(term }\SpecialCharTok{==} \StringTok{"year"}\NormalTok{)}

\FunctionTok{ggplot}\NormalTok{(effect\_df, }\FunctionTok{aes}\NormalTok{(}\AttributeTok{x =} \FunctionTok{factor}\NormalTok{(topic), }\AttributeTok{y =}\NormalTok{ estimate)) }\SpecialCharTok{+}
  \FunctionTok{geom\_col}\NormalTok{(}\AttributeTok{fill =} \FunctionTok{palette\_sci}\NormalTok{(}\DecValTok{1}\NormalTok{)) }\SpecialCharTok{+}
  \FunctionTok{geom\_errorbar}\NormalTok{(}\FunctionTok{aes}\NormalTok{(}\AttributeTok{ymin =}\NormalTok{ estimate }\SpecialCharTok{{-}} \FloatTok{1.96} \SpecialCharTok{*}\NormalTok{ se, }\AttributeTok{ymax =}\NormalTok{ estimate }\SpecialCharTok{+} \FloatTok{1.96} \SpecialCharTok{*}\NormalTok{ se),}
                \AttributeTok{width =} \FloatTok{0.3}\NormalTok{) }\SpecialCharTok{+}
  \FunctionTok{geom\_hline}\NormalTok{(}\AttributeTok{yintercept =} \DecValTok{0}\NormalTok{, }\AttributeTok{linetype =} \StringTok{"dashed"}\NormalTok{, }\AttributeTok{colour =} \StringTok{"grey50"}\NormalTok{) }\SpecialCharTok{+}
  \FunctionTok{labs}\NormalTok{(}\AttributeTok{x =} \StringTok{"Topic"}\NormalTok{, }\AttributeTok{y =} \StringTok{"Year effect on prevalence"}\NormalTok{) }\SpecialCharTok{+}
  \FunctionTok{theme\_sci}\NormalTok{()}
\end{Highlighting}
\end{Shaded}

\begin{figure}

{\centering \includegraphics[width=0.9\linewidth,alt={Line chart showing how the estimated proportion of each topic changes across publication years. Some topics grow while others decline.}]{scientometrics-in-r_files/figure-latex/22-fig-topicmodel-temporal-1} 

}

\caption{Estimated topic prevalence by publication year (STM).}\label{fig:22-fig-topicmodel-temporal}
\end{figure}

\subsection{Model diagnostics}\label{model-diagnostics}

\begin{Shaded}
\begin{Highlighting}[]
\NormalTok{coherence }\OtherTok{\textless{}{-}} \FunctionTok{semanticCoherence}\NormalTok{(stm\_model, stm\_dfm}\SpecialCharTok{$}\NormalTok{documents)}
\NormalTok{exclusivity }\OtherTok{\textless{}{-}} \FunctionTok{exclusivity}\NormalTok{(stm\_model)}

\NormalTok{diag\_df }\OtherTok{\textless{}{-}} \FunctionTok{tibble}\NormalTok{(}
  \AttributeTok{topic =} \DecValTok{1}\SpecialCharTok{:}\DecValTok{8}\NormalTok{,}
  \AttributeTok{coherence =} \FunctionTok{round}\NormalTok{(coherence, }\DecValTok{2}\NormalTok{),}
  \AttributeTok{exclusivity =} \FunctionTok{round}\NormalTok{(exclusivity, }\DecValTok{2}\NormalTok{)}
\NormalTok{)}

\NormalTok{diag\_df }\SpecialCharTok{|\textgreater{}} \FunctionTok{gt}\NormalTok{()}
\end{Highlighting}
\end{Shaded}

\begingroup
\fontsize{8.0pt}{9.0pt}\selectfont
\begin{longtable}{>{\raggedleft\arraybackslash}p{\dimexpr 0.31\linewidth -2\tabcolsep-1.5\arrayrulewidth}>{\raggedleft\arraybackslash}p{\dimexpr 0.33\linewidth -2\tabcolsep-1.5\arrayrulewidth}>{\raggedleft\arraybackslash}p{\dimexpr 0.36\linewidth -2\tabcolsep-1.5\arrayrulewidth}}
\toprule
topic & \nolinkurl{coherence} & \nolinkurl{exclusivity} \\ 
\midrule\addlinespace[2.5pt]
1 & -95.7 & 9.20 \\ 
2 & -78.3 & 8.88 \\ 
3 & -53.5 & 9.19 \\ 
4 & -95.4 & 9.34 \\ 
5 & -100.5 & 9.64 \\ 
6 & -79.9 & 9.20 \\ 
7 & -88.7 & 8.87 \\ 
8 & -67.2 & 9.10 \\ 
\bottomrule
\end{longtable}
\endgroup

\section{Diagnostics and interpretation}\label{topicmodel-diagnostics}

\begin{itemize}
\tightlist
\item
  \textbf{Coherence}: High coherence means the top words in a topic tend to co-occur in the same documents. Low coherence indicates a ``junk'' topic mixing unrelated terms.
\item
  \textbf{Exclusivity}: High exclusivity means the top words are distinctive to that topic. Topics with low exclusivity share vocabulary with other topics, suggesting they may need to be split.
\item
  \textbf{Coherence-exclusivity trade-off}: Adding more topics generally improves exclusivity but may decrease coherence. Plot both metrics across K values to find a balance.
\item
  \textbf{Human validation}: Always read sample documents assigned to each topic. If a topic cannot be labelled by a domain expert, it may not represent a meaningful theme.
\end{itemize}

\section{Limitations and responsible use}\label{topicmodel-limitations}

\section{Limitations and responsible use}\label{limitations-and-responsible-use-21}

\begin{itemize}
\tightlist
\item
  \textbf{Topics are not ground truth.} Topic models are statistical constructs, not natural categories. Different random seeds, preprocessing choices, or K values produce different topics. Never claim a topic model reveals ``the'' structure of a field.
\item
  \textbf{Abstract-only models miss depth.} Abstracts are compressed summaries. Topics derived from abstracts are coarser than those from full text.
\item
  \textbf{Number of topics is a researcher choice.} There is no objectively correct K. Sensitivity analysis across multiple K values is essential.
\item
  \textbf{Temporal confounds.} In STM, a topic that ``grows'' over time may reflect genuine intellectual shift or simply increased coverage by the database in recent years \citep{priem2022}.
\item
  \textbf{Replication.} Topic models are sensitive to preprocessing and random initialisation. Report all parameters and fix seeds for reproducibility.
\end{itemize}

\section{Common pitfalls}\label{topicmodel-pitfalls}

\section{Common pitfalls}\label{common-pitfalls-21}

\begin{itemize}
\tightlist
\item
  \textbf{Not preprocessing enough.} Rare terms, numbers, and generic academic vocabulary produce noisy topics. Trim the DFM before fitting.
\item
  \textbf{Choosing K by perplexity alone.} Held-out perplexity often favours large K values that produce uninterpretable topics. Use coherence and exclusivity alongside perplexity.
\item
  \textbf{Treating topic proportions as precise measurements.} Document-topic proportions are probabilistic estimates with uncertainty. Do not report them to three decimal places.
\item
  \textbf{Comparing topics across models.} LDA and STM may number and define topics differently. Align topics by content, not by index.
\end{itemize}

\section{Exercises}\label{topicmodel-exercises}

\begin{enumerate}
\def\labelenumi{\arabic{enumi}.}
\item
  \textbf{K selection.} Fit LDA models with K = 5, 8, 12, and 15. Compute coherence for each. Plot coherence against K and select the best value. How sensitive are the results?
\item
  \textbf{STM with journal covariate.} Fetch data from two different journals. Fit an STM with journal as a prevalence covariate. Which topics differ most between journals?
\item
  \textbf{Topic labelling.} For the best LDA model, read the titles of the 5 documents with the highest proportion for each topic. Assign a human-readable label to each topic.
\item
  \textbf{LDA vs.~STM comparison.} Compare the topics from LDA and STM for the same K. Are they similar? What does the STM's year covariate add?
\end{enumerate}

\section{Solutions}\label{22-solutions}

Solutions are provided in \ref{solutions}.

\section{Further reading}\label{topicmodel-reading}

\begin{itemize}
\tightlist
\item
  \citet{blei2003} --- The foundational LDA paper.
\item
  \citet{roberts2019} --- The Structural Topic Model with covariates.
\item
  \citet{silge2017} --- \emph{Text Mining with R}: accessible guide to topic modelling in R.
\item
  \citet{aria2017} --- \texttt{bibliometrix} topic analysis features.
\end{itemize}

\section{Session info}\label{topicmodel-sessioninfo}

\begin{Shaded}
\begin{Highlighting}[]
\FunctionTok{sessionInfo}\NormalTok{()}
\end{Highlighting}
\end{Shaded}

\begin{verbatim}
#> R version 4.6.1 (2026-06-24)
#> Platform: x86_64-pc-linux-gnu
#> Running under: Ubuntu 24.04.5 LTS
#> 
#> Matrix products: default
#> BLAS:   /usr/lib/x86_64-linux-gnu/openblas-pthread/libblas.so.3 
#> LAPACK: /usr/lib/x86_64-linux-gnu/openblas-pthread/libopenblasp-r0.3.26.so;  LAPACK version 3.12.0
#> 
#> locale:
#>  [1] LC_CTYPE=C.UTF-8       LC_NUMERIC=C           LC_TIME=C.UTF-8       
#>  [4] LC_COLLATE=C.UTF-8     LC_MONETARY=C.UTF-8    LC_MESSAGES=C.UTF-8   
#>  [7] LC_PAPER=C.UTF-8       LC_NAME=C              LC_ADDRESS=C          
#> [10] LC_TELEPHONE=C         LC_MEASUREMENT=C.UTF-8 LC_IDENTIFICATION=C   
#> 
#> time zone: UTC
#> tzcode source: system (glibc)
#> 
#> attached base packages:
#> [1] stats     graphics  grDevices utils     datasets  methods   base     
#> 
#> other attached packages:
#>  [1] stm_1.3.8                 topicmodels_0.2-17       
#>  [3] quanteda.textstats_0.97.2 visNetwork_2.1.4         
#>  [5] ggraph_2.2.2              tidygraph_1.3.1          
#>  [7] igraph_2.3.4              quanteda_4.5.1           
#>  [9] pdftools_3.9.1            arrow_26.0.0             
#> [11] bibliometrix_5.5.0        RefManageR_1.4.0         
#> [13] bib2df_1.1.2.0            rcrossref_1.2.1          
#> [15] gt_1.4.0                  tidytext_0.4.3           
#> [17] glue_1.8.1                openalexR_3.1.2          
#> [19] lubridate_1.9.5           forcats_1.0.1            
#> [21] stringr_1.6.0             dplyr_1.2.1              
#> [23] purrr_1.2.2               readr_2.2.0              
#> [25] tidyr_1.3.2               tibble_3.3.1             
#> [27] ggplot2_4.0.3             tidyverse_2.0.0          
#> 
#> loaded via a namespace (and not attached):
#>   [1] bibtex_0.5.3           RColorBrewer_1.1-3     rstudioapi_0.19.0     
#>   [4] jsonlite_2.0.0         magrittr_2.0.5         modeltools_0.2-25     
#>   [7] farver_2.1.2           rmarkdown_2.32         fs_2.1.0              
#>  [10] vctrs_0.7.3            memoise_2.0.1          askpass_1.2.1         
#>  [13] base64enc_0.1-6        htmltools_0.5.9        contentanalysis_1.1.1 
#>  [16] curl_8.1.0             janeaustenr_1.0.0      cellranger_1.1.0      
#>  [19] htmlwidgets_1.6.4      tokenizers_0.3.0       plyr_1.8.9            
#>  [22] httr2_1.3.0            cachem_1.1.0           plotly_4.12.1         
#>  [25] dimensionsR_0.0.3      mime_0.13              lifecycle_1.0.5       
#>  [28] pkgconfig_2.0.3        Matrix_1.7-5           R6_2.6.1              
#>  [31] fastmap_1.2.0          shiny_1.14.0           digest_0.6.39         
#>  [34] patchwork_1.3.2        shinycssloaders_1.1.0  rprojroot_2.1.1       
#>  [37] SnowballC_0.7.1        labeling_0.4.3         urltools_1.7.3.1      
#>  [40] timechange_0.4.0       httr_1.4.9             polyclip_1.10-7       
#>  [43] compiler_4.6.1         here_1.0.2             bit64_4.8.6           
#>  [46] withr_3.0.3            S7_0.2.2               viridis_0.6.5         
#>  [49] ggforce_0.5.0          MASS_7.3-65            bibliometrixData_0.3.0
#>  [52] tools_4.6.1            otel_0.2.0             stopwords_2.3         
#>  [55] zip_3.0.2              rentrez_1.2.4          httpuv_1.6.17         
#>  [58] promises_1.5.0         grid_4.6.1             stringdist_0.9.17     
#>  [61] reshape2_1.4.5         generics_0.1.4         gtable_0.3.6          
#>  [64] tzdb_0.5.0             rscopus_0.9.0          ca_0.72               
#>  [67] data.table_1.18.6.1    hms_1.1.4              xml2_1.6.0            
#>  [70] utf8_1.2.6             ggrepel_0.9.8          pillar_1.11.1         
#>  [73] nsyllable_1.0.1        vroom_1.7.1            later_1.4.8           
#>  [76] tweenr_2.0.3           lattice_0.22-9         bit_4.6.0             
#>  [79] tidyselect_1.2.1       tm_0.7-20              miniUI_0.1.2          
#>  [82] knitr_1.52             gridExtra_2.3.1        NLP_0.3-3             
#>  [85] bookdown_0.48          stats4_4.6.1           crul_1.6.0            
#>  [88] xfun_0.62              graphlayouts_1.2.5     matrixStats_1.5.0     
#>  [91] DT_0.34.0              humaniformat_0.6.0     stringi_1.8.9         
#>  [94] qpdf_1.4.1             yaml_2.3.12            evaluate_1.0.5        
#>  [97] httpcode_0.3.0         cli_3.6.6              xtable_1.8-8          
#> [100] systemfonts_1.3.2      Rcpp_1.1.2             readxl_1.6.0          
#> [103] triebeard_0.4.1        XML_3.99-0.25          parallel_4.6.1        
#> [106] assertthat_0.2.1       pubmedR_1.0.2          slam_0.1-56           
#> [109] viridisLite_0.4.3      scales_1.4.0           crayon_1.5.3          
#> [112] openxlsx_4.2.9         rlang_1.3.0            fastmatch_1.1-8
\end{verbatim}

\chapter{Keyword Extraction}\label{keyword-extraction}

\emph{R. Heller · Date: 2026-05-09 · Last changed: 2026-10-31}

\section{Learning objectives}\label{kwextract-objectives}

After completing this chapter, you will be able to:

\begin{itemize}
\tightlist
\item
  Explain the RAKE and TextRank keyword extraction algorithms
\item
  Apply keyword extraction to bibliometric abstracts using R packages
\item
  Compare automatically extracted keywords with author-assigned keywords
\item
  Evaluate extraction quality using precision, recall, and F1 metrics
\item
  Decide when automatic keyword extraction adds value over existing metadata
\end{itemize}

\section{Setup}\label{kwextract-setup}

\begin{Shaded}
\begin{Highlighting}[]
\FunctionTok{library}\NormalTok{(tidyverse)}
\FunctionTok{library}\NormalTok{(openalexR)}
\FunctionTok{library}\NormalTok{(quanteda)}
\FunctionTok{library}\NormalTok{(quanteda.textstats)}
\FunctionTok{library}\NormalTok{(glue)}
\FunctionTok{library}\NormalTok{(gt)}

\FunctionTok{set.seed}\NormalTok{(}\DecValTok{20260509}\NormalTok{)}

\FunctionTok{source}\NormalTok{(here}\SpecialCharTok{::}\FunctionTok{here}\NormalTok{(}\StringTok{"R"}\NormalTok{, }\StringTok{"api\_helpers.R"}\NormalTok{))}
\FunctionTok{source}\NormalTok{(here}\SpecialCharTok{::}\FunctionTok{here}\NormalTok{(}\StringTok{"R"}\NormalTok{, }\StringTok{"rake\_keywords.R"}\NormalTok{))}
\FunctionTok{source}\NormalTok{(here}\SpecialCharTok{::}\FunctionTok{here}\NormalTok{(}\StringTok{"R"}\NormalTok{, }\StringTok{"utils.R"}\NormalTok{))}
\FunctionTok{source}\NormalTok{(here}\SpecialCharTok{::}\FunctionTok{here}\NormalTok{(}\StringTok{"R"}\NormalTok{, }\StringTok{"sci\_palette.R"}\NormalTok{))}
\end{Highlighting}
\end{Shaded}

\section{Conceptual background}\label{kwextract-background}

Keywords serve as compact descriptors of a document's content. In bibliometrics, they are used for co-word analysis (\ref{coword-background}), topic mapping, and search. Three sources of keywords exist:

\begin{enumerate}
\def\labelenumi{\arabic{enumi}.}
\tightlist
\item
  \textbf{Author keywords}: Selected by the paper's authors. Intentional but inconsistent --- different authors use different terms for the same concept.
\item
  \textbf{Indexed terms}: Assigned by database curators (MeSH in PubMed, descriptors in WoS). Consistent but labour-intensive and available only in some databases.
\item
  \textbf{Extracted keywords}: Automatically derived from titles and abstracts using algorithms. Scalable and consistent but may miss domain nuance.
\end{enumerate}

\textbf{RAKE} (Rapid Automatic Keyword Extraction) identifies candidate keywords by finding sequences of content words delimited by stopwords or punctuation. It scores candidates by the ratio of word weighted degree (including the word itself within candidate phrases) to word frequency, favouring multi-word phrases that appear in specific contexts. RAKE is fast and unsupervised but does not consider document-level or corpus-level statistics.

\textbf{TextRank} adapts the PageRank algorithm to text: words are nodes in a graph, edges connect co-occurring words within a sliding window, and importance is computed iteratively. High-ranking words are selected as keywords. TextRank captures contextual importance better than frequency-based methods but is more computationally expensive.

Both methods extract keywords from individual documents. \textbf{TF-IDF} (\ref{textmine-background}) provides a complementary corpus-level view, identifying terms that distinguish a document from the broader collection.

Evaluating keyword extraction requires ground truth. Author keywords are the most common benchmark, but they are imperfect: authors may omit obvious keywords or include overly specific terms. Agreement depends on the corpus, matching rules, number of keywords, and reference annotation. Define these choices before evaluation and report observed precision and recall; no universal acceptance percentage is assumed here.

\section{Worked example}\label{kwextract-example}

\subsection{Fetching data with keywords}\label{fetching-data-with-keywords}

\begin{Shaded}
\begin{Highlighting}[]
\NormalTok{works }\OtherTok{\textless{}{-}} \FunctionTok{fetch\_openalex}\NormalTok{(}
  \AttributeTok{entity =} \StringTok{"works"}\NormalTok{,}
  \AttributeTok{primary\_location.source.id =} \StringTok{"S148561398"}\NormalTok{,}
  \AttributeTok{from\_publication\_date =} \StringTok{"2021{-}01{-}01"}\NormalTok{,}
  \AttributeTok{to\_publication\_date =} \StringTok{"2023{-}12{-}31"}\NormalTok{,}
  \AttributeTok{options =} \FunctionTok{list}\NormalTok{(}\AttributeTok{sample =} \DecValTok{200}\NormalTok{, }\AttributeTok{seed =} \DecValTok{42}\NormalTok{)}
\NormalTok{)}

\NormalTok{text\_df }\OtherTok{\textless{}{-}}\NormalTok{ works }\SpecialCharTok{|\textgreater{}}
  \FunctionTok{filter}\NormalTok{(}\SpecialCharTok{!}\FunctionTok{is.na}\NormalTok{(abstract), }\FunctionTok{nchar}\NormalTok{(abstract) }\SpecialCharTok{\textgreater{}} \DecValTok{100}\NormalTok{) }\SpecialCharTok{|\textgreater{}}
  \FunctionTok{transmute}\NormalTok{(}
    \AttributeTok{doc\_id =}\NormalTok{ id,}
    \AttributeTok{title =}\NormalTok{ display\_name,}
    \AttributeTok{abstract =}\NormalTok{ abstract,}
    \AttributeTok{text =} \FunctionTok{paste}\NormalTok{(display\_name, abstract, }\AttributeTok{sep =} \StringTok{". "}\NormalTok{)}
\NormalTok{  )}

\FunctionTok{cat}\NormalTok{(}\FunctionTok{glue}\NormalTok{(}\StringTok{"Documents with abstracts: \{nrow(text\_df)\}}\SpecialCharTok{\textbackslash{}n}\StringTok{"}\NormalTok{))}
\end{Highlighting}
\end{Shaded}

\begin{verbatim}
#> Documents with abstracts: 67
\end{verbatim}

\subsection{RAKE keyword extraction}\label{rake-keyword-extraction}

\begin{Shaded}
\begin{Highlighting}[]
\NormalTok{rake\_extract }\OtherTok{\textless{}{-}} \ControlFlowTok{function}\NormalTok{(text, }\AttributeTok{n\_keywords =} \DecValTok{10}\NormalTok{) \{}
  \FunctionTok{rake\_keywords}\NormalTok{(text, }\AttributeTok{stopwords =}\NormalTok{ quanteda}\SpecialCharTok{::}\FunctionTok{stopwords}\NormalTok{(}\StringTok{"en"}\NormalTok{),}
                \AttributeTok{n\_keywords =}\NormalTok{ n\_keywords)}\SpecialCharTok{$}\NormalTok{keyword}
\NormalTok{\}}

\NormalTok{text\_df }\OtherTok{\textless{}{-}}\NormalTok{ text\_df }\SpecialCharTok{|\textgreater{}}
  \FunctionTok{mutate}\NormalTok{(}\AttributeTok{rake\_keywords =} \FunctionTok{map}\NormalTok{(text, \textbackslash{}(t) }\FunctionTok{rake\_extract}\NormalTok{(t, }\AttributeTok{n\_keywords =} \DecValTok{8}\NormalTok{)))}

\NormalTok{text\_df }\SpecialCharTok{|\textgreater{}}
  \FunctionTok{head}\NormalTok{(}\DecValTok{3}\NormalTok{) }\SpecialCharTok{|\textgreater{}}
  \FunctionTok{select}\NormalTok{(title, rake\_keywords) }\SpecialCharTok{|\textgreater{}}
  \FunctionTok{mutate}\NormalTok{(}\AttributeTok{rake\_keywords =} \FunctionTok{map\_chr}\NormalTok{(rake\_keywords, paste, }\AttributeTok{collapse =} \StringTok{"; "}\NormalTok{)) }\SpecialCharTok{|\textgreater{}}
  \FunctionTok{gt}\NormalTok{()}
\end{Highlighting}
\end{Shaded}

\begingroup
\fontsize{8.0pt}{9.0pt}\selectfont
\begin{longtable}{>{\raggedright\arraybackslash}p{\dimexpr 0.50\linewidth -2\tabcolsep-1.5\arrayrulewidth}>{\raggedright\arraybackslash}p{\dimexpr 0.50\linewidth -2\tabcolsep-1.5\arrayrulewidth}}
\toprule
title & \nolinkurl{rake_keywords} \\ 
\midrule\addlinespace[2.5pt]
Market forces in Italian academia today (and yesterday) & better scholars travel longer distances; consistent nesting procedure; greater expected utility; choosing private universities; italian academia today; mapping current scholars; current academic market; sorting effect delineates \\ 
Why do papers from international collaborations get more citations? A bibliometric analysis of Library and Information Science papers & achieve relevant research outputs; countries must therefore consider; science mapping analysis approach; achieve greater scientific impact; abstract scientific activity; become increasingly complex; thus becomes necessary; information science category \\ 
Research on Covid-19: a disruptive phenomenon for bibliometrics & present study also confirms; san francisco declaration; world recent history; present study shows; highest disruptive event; bibliometrics provide methods; journal impact factor; worldwide academic research \\ 
\bottomrule
\end{longtable}
\endgroup

\subsection{TF-IDF based keyword extraction}\label{tf-idf-based-keyword-extraction}

\begin{Shaded}
\begin{Highlighting}[]
\NormalTok{corp }\OtherTok{\textless{}{-}} \FunctionTok{corpus}\NormalTok{(text\_df, }\AttributeTok{docid\_field =} \StringTok{"doc\_id"}\NormalTok{, }\AttributeTok{text\_field =} \StringTok{"text"}\NormalTok{)}

\NormalTok{toks }\OtherTok{\textless{}{-}} \FunctionTok{tokens}\NormalTok{(corp, }\AttributeTok{remove\_punct =} \ConstantTok{TRUE}\NormalTok{, }\AttributeTok{remove\_numbers =} \ConstantTok{TRUE}\NormalTok{) }\SpecialCharTok{|\textgreater{}}
  \FunctionTok{tokens\_tolower}\NormalTok{() }\SpecialCharTok{|\textgreater{}}
  \FunctionTok{tokens\_remove}\NormalTok{(}\FunctionTok{stopwords}\NormalTok{(}\StringTok{"en"}\NormalTok{)) }\SpecialCharTok{|\textgreater{}}
  \FunctionTok{tokens\_remove}\NormalTok{(}\FunctionTok{c}\NormalTok{(}\StringTok{"study"}\NormalTok{, }\StringTok{"paper"}\NormalTok{, }\StringTok{"results"}\NormalTok{, }\StringTok{"research"}\NormalTok{, }\StringTok{"analysis"}\NormalTok{))}

\NormalTok{dfmat }\OtherTok{\textless{}{-}} \FunctionTok{dfm}\NormalTok{(toks) }\SpecialCharTok{|\textgreater{}}
  \FunctionTok{dfm\_trim}\NormalTok{(}\AttributeTok{min\_termfreq =} \DecValTok{2}\NormalTok{)}
\NormalTok{dfmat\_tfidf }\OtherTok{\textless{}{-}} \FunctionTok{dfm\_tfidf}\NormalTok{(dfmat)}

\NormalTok{tfidf\_keywords }\OtherTok{\textless{}{-}} \FunctionTok{map}\NormalTok{(}\FunctionTok{seq\_len}\NormalTok{(}\FunctionTok{nrow}\NormalTok{(dfmat\_tfidf)), }\ControlFlowTok{function}\NormalTok{(i) \{}
\NormalTok{  row }\OtherTok{\textless{}{-}}\NormalTok{ dfmat\_tfidf[i, ]}
\NormalTok{  vals }\OtherTok{\textless{}{-}} \FunctionTok{as.numeric}\NormalTok{(row)}
  \FunctionTok{names}\NormalTok{(vals) }\OtherTok{\textless{}{-}} \FunctionTok{featnames}\NormalTok{(dfmat\_tfidf)}
  \FunctionTok{names}\NormalTok{(}\FunctionTok{sort}\NormalTok{(vals, }\AttributeTok{decreasing =} \ConstantTok{TRUE}\NormalTok{))[}\DecValTok{1}\SpecialCharTok{:}\DecValTok{8}\NormalTok{]}
\NormalTok{\})}

\NormalTok{text\_df}\SpecialCharTok{$}\NormalTok{tfidf\_keywords }\OtherTok{\textless{}{-}}\NormalTok{ tfidf\_keywords}
\end{Highlighting}
\end{Shaded}

\subsection{Corpus-level keyword frequency}\label{corpus-level-keyword-frequency}

\begin{Shaded}
\begin{Highlighting}[]
\NormalTok{all\_kw }\OtherTok{\textless{}{-}}\NormalTok{ text\_df }\SpecialCharTok{|\textgreater{}}
  \FunctionTok{pull}\NormalTok{(rake\_keywords) }\SpecialCharTok{|\textgreater{}}
  \FunctionTok{unlist}\NormalTok{() }\SpecialCharTok{|\textgreater{}}
  \FunctionTok{tibble}\NormalTok{(}\AttributeTok{keyword =}\NormalTok{ \_) }\SpecialCharTok{|\textgreater{}}
  \FunctionTok{count}\NormalTok{(keyword, }\AttributeTok{sort =} \ConstantTok{TRUE}\NormalTok{)}

\NormalTok{all\_kw }\SpecialCharTok{|\textgreater{}}
  \FunctionTok{head}\NormalTok{(}\DecValTok{20}\NormalTok{) }\SpecialCharTok{|\textgreater{}}
  \FunctionTok{mutate}\NormalTok{(}\AttributeTok{keyword =} \FunctionTok{fct\_reorder}\NormalTok{(keyword, n)) }\SpecialCharTok{|\textgreater{}}
  \FunctionTok{ggplot}\NormalTok{(}\FunctionTok{aes}\NormalTok{(}\AttributeTok{x =}\NormalTok{ n, }\AttributeTok{y =}\NormalTok{ keyword)) }\SpecialCharTok{+}
  \FunctionTok{geom\_col}\NormalTok{(}\AttributeTok{fill =} \FunctionTok{palette\_sci}\NormalTok{(}\DecValTok{1}\NormalTok{)) }\SpecialCharTok{+}
  \FunctionTok{labs}\NormalTok{(}\AttributeTok{x =} \StringTok{"Extraction frequency"}\NormalTok{, }\AttributeTok{y =} \ConstantTok{NULL}\NormalTok{) }\SpecialCharTok{+}
  \FunctionTok{theme\_sci}\NormalTok{()}
\end{Highlighting}
\end{Shaded}

\begin{figure}

{\centering \includegraphics[width=0.9\linewidth,alt={Horizontal bar chart showing the 20 most frequently extracted keywords and how often they appear across documents.}]{scientometrics-in-r_files/figure-latex/23-fig-kwextract-frequency-1} 

}

\caption{Top 20 keywords by total extraction frequency across the corpus.}\label{fig:23-fig-kwextract-frequency}
\end{figure}

\subsection{Comparing extraction methods}\label{comparing-extraction-methods}

\begin{Shaded}
\begin{Highlighting}[]
\NormalTok{overlap }\OtherTok{\textless{}{-}} \FunctionTok{map2\_dbl}\NormalTok{(text\_df}\SpecialCharTok{$}\NormalTok{rake\_keywords, text\_df}\SpecialCharTok{$}\NormalTok{tfidf\_keywords,}
\NormalTok{                    \textbackslash{}(r, t) }\FunctionTok{length}\NormalTok{(}\FunctionTok{intersect}\NormalTok{(r, t)) }\SpecialCharTok{/} \FunctionTok{length}\NormalTok{(}\FunctionTok{union}\NormalTok{(r, t)))}

\FunctionTok{cat}\NormalTok{(}\FunctionTok{glue}\NormalTok{(}\StringTok{"Mean Jaccard overlap (RAKE vs TF{-}IDF): \{round(mean(overlap, na.rm = TRUE), 3)\}}\SpecialCharTok{\textbackslash{}n}\StringTok{"}\NormalTok{))}
\end{Highlighting}
\end{Shaded}

\begin{verbatim}
#> Mean Jaccard overlap (RAKE vs TF-IDF): 0.001
\end{verbatim}

\begin{Shaded}
\begin{Highlighting}[]
\FunctionTok{tibble}\NormalTok{(}\AttributeTok{jaccard =}\NormalTok{ overlap) }\SpecialCharTok{|\textgreater{}}
  \FunctionTok{ggplot}\NormalTok{(}\FunctionTok{aes}\NormalTok{(}\AttributeTok{x =}\NormalTok{ jaccard)) }\SpecialCharTok{+}
  \FunctionTok{geom\_histogram}\NormalTok{(}\AttributeTok{binwidth =} \FloatTok{0.05}\NormalTok{, }\AttributeTok{fill =} \FunctionTok{palette\_sci}\NormalTok{(}\DecValTok{1}\NormalTok{), }\AttributeTok{colour =} \StringTok{"white"}\NormalTok{) }\SpecialCharTok{+}
  \FunctionTok{labs}\NormalTok{(}\AttributeTok{x =} \StringTok{"Jaccard similarity"}\NormalTok{, }\AttributeTok{y =} \StringTok{"Documents"}\NormalTok{) }\SpecialCharTok{+}
  \FunctionTok{theme\_sci}\NormalTok{()}
\end{Highlighting}
\end{Shaded}

\begin{figure}

{\centering \includegraphics[width=0.9\linewidth,alt={Histogram showing the distribution of Jaccard similarity scores between the two extraction methods. Most documents show moderate overlap.}]{scientometrics-in-r_files/figure-latex/23-fig-kwextract-overlap-1} 

}

\caption{Distribution of Jaccard similarity between RAKE and TF-IDF keyword sets.}\label{fig:23-fig-kwextract-overlap}
\end{figure}

\section{Diagnostics and interpretation}\label{kwextract-diagnostics}

\begin{itemize}
\tightlist
\item
  \textbf{Extraction quality}: Compare extracted keywords against author keywords for a random sample of 20--30 papers. Specify exact or normalized matching, report precision and recall, and inspect disagreements rather than applying an unsupported universal threshold.
\item
  \textbf{Redundancy}: Check for near-duplicate keywords (singular/plural, abbreviation/full form). Post-processing to merge variants improves quality.
\item
  \textbf{Domain specificity}: Generic terms (``model'', ``method'', ``data'') are often extracted but carry no topical information. Consider a domain-specific stopword list.
\item
  \textbf{Coverage}: Some documents yield few meaningful keywords, especially short abstracts. Flag documents where fewer than 3 keywords were extracted.
\end{itemize}

\section{Limitations and responsible use}\label{kwextract-limitations}

\section{Limitations and responsible use}\label{limitations-and-responsible-use-22}

\begin{itemize}
\tightlist
\item
  \textbf{No method is comprehensive.} Keyword extraction captures salient terms but misses concepts expressed through circumlocution or implicit reference. Human keywords remain valuable.
\item
  \textbf{Language dependence.} RAKE and TextRank work best for English. Stopword lists and tokenisation rules may fail for other languages.
\item
  \textbf{Evaluation is subjective.} There is no single correct keyword set for a document. Low agreement between methods or between extracted and author keywords does not necessarily indicate failure.
\item
  \textbf{Not a replacement for indexing.} Professional indexing (e.g., MeSH) follows controlled vocabularies and domain expertise. Automated extraction is a complement, not a substitute \citep{hicks2015}.
\end{itemize}

\section{Common pitfalls}\label{kwextract-pitfalls}

\section{Common pitfalls}\label{common-pitfalls-22}

\begin{itemize}
\tightlist
\item
  \textbf{Using raw frequencies as keywords.} The most frequent words in a document are often stopwords or near-stopwords. TF-IDF or RAKE scoring is essential.
\item
  \textbf{Not normalising keywords.} ``bibliometric'' and ``bibliometrics'' should be the same keyword. Apply stemming or lemmatisation before evaluation.
\item
  \textbf{Evaluating on titles only.} Titles are too short for reliable keyword extraction. Use abstracts or title + abstract combined.
\item
  \textbf{Ignoring multi-word terms.} Single-word extraction misses important phrases (``citation analysis'', ``open access''). Use n-grams or phrase detection.
\end{itemize}

\section{Exercises}\label{kwextract-exercises}

\begin{enumerate}
\def\labelenumi{\arabic{enumi}.}
\item
  \textbf{Precision and recall.} For 20 papers with author keywords, compute precision and recall of RAKE-extracted keywords against author keywords. Report the mean and range.
\item
  \textbf{Bigram keywords.} Extract bigrams from abstracts and score them by TF-IDF. What multi-word terms emerge that single-word extraction misses?
\item
  \textbf{Keyword enrichment.} For papers without author keywords, use extracted keywords to build a co-word network. Does the resulting topical structure resemble that from author keywords?
\end{enumerate}

\section{Solutions}\label{23-solutions}

Solutions are provided in \ref{solutions}.

\section{Further reading}\label{kwextract-reading}

\begin{itemize}
\tightlist
\item
  \citet{silge2017} --- Keyword and term extraction in R with tidy tools.
\item
  \citet{aria2017} --- \texttt{bibliometrix} keyword analysis capabilities.
\item
  \citet{callon1983} --- Co-word analysis, the downstream use case for extracted keywords.
\item
  \citet{priem2022} --- OpenAlex concept tagging as an alternative to keyword extraction.
\end{itemize}

\section{Session info}\label{kwextract-sessioninfo}

\begin{Shaded}
\begin{Highlighting}[]
\FunctionTok{sessionInfo}\NormalTok{()}
\end{Highlighting}
\end{Shaded}

\begin{verbatim}
#> R version 4.6.1 (2026-06-24)
#> Platform: x86_64-pc-linux-gnu
#> Running under: Ubuntu 24.04.5 LTS
#> 
#> Matrix products: default
#> BLAS:   /usr/lib/x86_64-linux-gnu/openblas-pthread/libblas.so.3 
#> LAPACK: /usr/lib/x86_64-linux-gnu/openblas-pthread/libopenblasp-r0.3.26.so;  LAPACK version 3.12.0
#> 
#> locale:
#>  [1] LC_CTYPE=C.UTF-8       LC_NUMERIC=C           LC_TIME=C.UTF-8       
#>  [4] LC_COLLATE=C.UTF-8     LC_MONETARY=C.UTF-8    LC_MESSAGES=C.UTF-8   
#>  [7] LC_PAPER=C.UTF-8       LC_NAME=C              LC_ADDRESS=C          
#> [10] LC_TELEPHONE=C         LC_MEASUREMENT=C.UTF-8 LC_IDENTIFICATION=C   
#> 
#> time zone: UTC
#> tzcode source: system (glibc)
#> 
#> attached base packages:
#> [1] stats     graphics  grDevices utils     datasets  methods   base     
#> 
#> other attached packages:
#>  [1] stm_1.3.8                 topicmodels_0.2-17       
#>  [3] quanteda.textstats_0.97.2 visNetwork_2.1.4         
#>  [5] ggraph_2.2.2              tidygraph_1.3.1          
#>  [7] igraph_2.3.4              quanteda_4.5.1           
#>  [9] pdftools_3.9.1            arrow_26.0.0             
#> [11] bibliometrix_5.5.0        RefManageR_1.4.0         
#> [13] bib2df_1.1.2.0            rcrossref_1.2.1          
#> [15] gt_1.4.0                  tidytext_0.4.3           
#> [17] glue_1.8.1                openalexR_3.1.2          
#> [19] lubridate_1.9.5           forcats_1.0.1            
#> [21] stringr_1.6.0             dplyr_1.2.1              
#> [23] purrr_1.2.2               readr_2.2.0              
#> [25] tidyr_1.3.2               tibble_3.3.1             
#> [27] ggplot2_4.0.3             tidyverse_2.0.0          
#> 
#> loaded via a namespace (and not attached):
#>   [1] bibtex_0.5.3           RColorBrewer_1.1-3     rstudioapi_0.19.0     
#>   [4] jsonlite_2.0.0         magrittr_2.0.5         modeltools_0.2-25     
#>   [7] farver_2.1.2           rmarkdown_2.32         fs_2.1.0              
#>  [10] vctrs_0.7.3            memoise_2.0.1          askpass_1.2.1         
#>  [13] base64enc_0.1-6        htmltools_0.5.9        contentanalysis_1.1.1 
#>  [16] curl_8.1.0             janeaustenr_1.0.0      cellranger_1.1.0      
#>  [19] htmlwidgets_1.6.4      tokenizers_0.3.0       plyr_1.8.9            
#>  [22] httr2_1.3.0            cachem_1.1.0           plotly_4.12.1         
#>  [25] dimensionsR_0.0.3      mime_0.13              lifecycle_1.0.5       
#>  [28] pkgconfig_2.0.3        Matrix_1.7-5           R6_2.6.1              
#>  [31] fastmap_1.2.0          shiny_1.14.0           digest_0.6.39         
#>  [34] patchwork_1.3.2        shinycssloaders_1.1.0  rprojroot_2.1.1       
#>  [37] SnowballC_0.7.1        labeling_0.4.3         urltools_1.7.3.1      
#>  [40] timechange_0.4.0       httr_1.4.9             polyclip_1.10-7       
#>  [43] compiler_4.6.1         here_1.0.2             bit64_4.8.6           
#>  [46] withr_3.0.3            S7_0.2.2               viridis_0.6.5         
#>  [49] ggforce_0.5.0          MASS_7.3-65            bibliometrixData_0.3.0
#>  [52] tools_4.6.1            otel_0.2.0             stopwords_2.3         
#>  [55] zip_3.0.2              rentrez_1.2.4          httpuv_1.6.17         
#>  [58] promises_1.5.0         grid_4.6.1             stringdist_0.9.17     
#>  [61] reshape2_1.4.5         generics_0.1.4         gtable_0.3.6          
#>  [64] tzdb_0.5.0             rscopus_0.9.0          ca_0.72               
#>  [67] data.table_1.18.6.1    hms_1.1.4              xml2_1.6.0            
#>  [70] utf8_1.2.6             ggrepel_0.9.8          pillar_1.11.1         
#>  [73] nsyllable_1.0.1        vroom_1.7.1            later_1.4.8           
#>  [76] tweenr_2.0.3           lattice_0.22-9         bit_4.6.0             
#>  [79] tidyselect_1.2.1       tm_0.7-20              miniUI_0.1.2          
#>  [82] knitr_1.52             gridExtra_2.3.1        NLP_0.3-3             
#>  [85] bookdown_0.48          stats4_4.6.1           crul_1.6.0            
#>  [88] xfun_0.62              graphlayouts_1.2.5     matrixStats_1.5.0     
#>  [91] DT_0.34.0              humaniformat_0.6.0     stringi_1.8.9         
#>  [94] qpdf_1.4.1             yaml_2.3.12            evaluate_1.0.5        
#>  [97] httpcode_0.3.0         cli_3.6.6              xtable_1.8-8          
#> [100] systemfonts_1.3.2      Rcpp_1.1.2             readxl_1.6.0          
#> [103] triebeard_0.4.1        XML_3.99-0.25          parallel_4.6.1        
#> [106] assertthat_0.2.1       pubmedR_1.0.2          slam_0.1-56           
#> [109] viridisLite_0.4.3      scales_1.4.0           crayon_1.5.3          
#> [112] openxlsx_4.2.9         rlang_1.3.0            fastmatch_1.1-8
\end{verbatim}

\chapter{Embeddings for Scientometrics}\label{embeddings-for-scientometrics}

\emph{R. Heller · Date: 2026-05-09 · Last changed: 2026-10-31}

\section{Learning objectives}\label{embed-objectives}

After completing this chapter, you will be able to:

\begin{itemize}
\tightlist
\item
  Explain how word and document embeddings differ from bag-of-words representations
\item
  Train simple word embeddings on a bibliometric corpus
\item
  Use pre-trained embeddings to compute document similarity
\item
  Apply UMAP or t-SNE to visualise document clusters in embedding space
\item
  Assess when embeddings add value over simpler text representations
\end{itemize}

\section{Setup}\label{embed-setup}

\begin{Shaded}
\begin{Highlighting}[]
\FunctionTok{library}\NormalTok{(tidyverse)}
\FunctionTok{library}\NormalTok{(openalexR)}
\FunctionTok{library}\NormalTok{(quanteda)}
\FunctionTok{library}\NormalTok{(word2vec)}
\FunctionTok{library}\NormalTok{(uwot)}
\FunctionTok{library}\NormalTok{(glue)}
\FunctionTok{library}\NormalTok{(gt)}

\FunctionTok{set.seed}\NormalTok{(}\DecValTok{20260509}\NormalTok{)}

\FunctionTok{source}\NormalTok{(here}\SpecialCharTok{::}\FunctionTok{here}\NormalTok{(}\StringTok{"R"}\NormalTok{, }\StringTok{"api\_helpers.R"}\NormalTok{))}
\FunctionTok{source}\NormalTok{(here}\SpecialCharTok{::}\FunctionTok{here}\NormalTok{(}\StringTok{"R"}\NormalTok{, }\StringTok{"utils.R"}\NormalTok{))}
\FunctionTok{source}\NormalTok{(here}\SpecialCharTok{::}\FunctionTok{here}\NormalTok{(}\StringTok{"R"}\NormalTok{, }\StringTok{"sci\_palette.R"}\NormalTok{))}
\end{Highlighting}
\end{Shaded}

\section{Conceptual background}\label{embed-background}

Bag-of-words representations (\ref{textmine-background}) treat each word as an independent dimension. ``Citation analysis'' and ``bibliometric study'' share no features despite being semantically related. \textbf{Word embeddings} address this by mapping words into a dense, low-dimensional vector space where semantically similar words are close together. The word2vec algorithm learns these representations by predicting words from their context (skip-gram) or context from words (CBOW) in a large training corpus.

\textbf{Document embeddings} extend the idea to whole documents. Simple approaches average the word vectors of a document's words (mean pooling). More sophisticated approaches include Doc2Vec (Paragraph Vector), which learns document-specific vectors alongside word vectors, and transformer-based models like \textbf{SPECTER}, which is pre-trained on scientific paper titles and abstracts using citation-based training signals.

For scientometrics, embeddings enable several applications:

\begin{itemize}
\tightlist
\item
  \textbf{Semantic similarity}: Finding papers with similar meaning, even when they use different terminology.
\item
  \textbf{Clustering}: Grouping papers by content in embedding space, complementing topic models (\ref{topicmodel-background}).
\item
  \textbf{Visualisation}: Projecting embeddings to 2D with UMAP or t-SNE to create interpretable ``maps'' of a corpus.
\item
  \textbf{Anomaly detection}: Identifying papers that are semantically distant from their assigned field or cluster.
\end{itemize}

The trade-off is interpretability: while TF-IDF features are directly readable (the word ``citation'' has a specific weight), embedding dimensions are opaque. Embeddings are powerful for similarity and clustering but resist the kind of term-level interpretation that co-word analysis provides.

\section{Worked example}\label{embed-example}

\subsection{Preparing text data}\label{preparing-text-data}

\begin{Shaded}
\begin{Highlighting}[]
\NormalTok{works }\OtherTok{\textless{}{-}} \FunctionTok{fetch\_openalex}\NormalTok{(}
  \AttributeTok{entity =} \StringTok{"works"}\NormalTok{,}
  \AttributeTok{primary\_location.source.id =} \StringTok{"S148561398"}\NormalTok{,}
  \AttributeTok{from\_publication\_date =} \StringTok{"2019{-}01{-}01"}\NormalTok{,}
  \AttributeTok{to\_publication\_date =} \StringTok{"2023{-}12{-}31"}\NormalTok{,}
  \AttributeTok{options =} \FunctionTok{list}\NormalTok{(}\AttributeTok{sample =} \DecValTok{400}\NormalTok{, }\AttributeTok{seed =} \DecValTok{42}\NormalTok{)}
\NormalTok{)}

\NormalTok{text\_df }\OtherTok{\textless{}{-}}\NormalTok{ works }\SpecialCharTok{|\textgreater{}}
  \FunctionTok{filter}\NormalTok{(}\SpecialCharTok{!}\FunctionTok{is.na}\NormalTok{(abstract), }\FunctionTok{nchar}\NormalTok{(abstract) }\SpecialCharTok{\textgreater{}} \DecValTok{100}\NormalTok{) }\SpecialCharTok{|\textgreater{}}
  \FunctionTok{transmute}\NormalTok{(}
    \AttributeTok{doc\_id =}\NormalTok{ id,}
    \AttributeTok{title =}\NormalTok{ display\_name,}
    \AttributeTok{text =} \FunctionTok{paste}\NormalTok{(display\_name, abstract, }\AttributeTok{sep =} \StringTok{". "}\NormalTok{),}
    \AttributeTok{year =} \FunctionTok{year}\NormalTok{(publication\_date),}
\NormalTok{    cited\_by\_count}
\NormalTok{  )}

\FunctionTok{cat}\NormalTok{(}\FunctionTok{glue}\NormalTok{(}\StringTok{"Documents for embedding: \{nrow(text\_df)\}}\SpecialCharTok{\textbackslash{}n}\StringTok{"}\NormalTok{))}
\end{Highlighting}
\end{Shaded}

\begin{verbatim}
#> Documents for embedding: 108
\end{verbatim}

\subsection{Training word2vec embeddings}\label{training-word2vec-embeddings}

\begin{Shaded}
\begin{Highlighting}[]
\NormalTok{clean\_text }\OtherTok{\textless{}{-}}\NormalTok{ text\_df}\SpecialCharTok{$}\NormalTok{text }\SpecialCharTok{|\textgreater{}}
  \FunctionTok{str\_to\_lower}\NormalTok{() }\SpecialCharTok{|\textgreater{}}
  \FunctionTok{str\_replace\_all}\NormalTok{(}\StringTok{"[\^{}a{-}z ]"}\NormalTok{, }\StringTok{" "}\NormalTok{) }\SpecialCharTok{|\textgreater{}}
  \FunctionTok{str\_squish}\NormalTok{()}

\NormalTok{model }\OtherTok{\textless{}{-}} \FunctionTok{word2vec}\NormalTok{(clean\_text, }\AttributeTok{dim =} \DecValTok{100}\NormalTok{, }\AttributeTok{iter =} \DecValTok{20}\NormalTok{, }\AttributeTok{min\_count =} \DecValTok{3}\NormalTok{,}
                  \AttributeTok{type =} \StringTok{"skip{-}gram"}\NormalTok{, }\AttributeTok{threads =} \DecValTok{1}\NormalTok{)}

\FunctionTok{cat}\NormalTok{(}\FunctionTok{glue}\NormalTok{(}\StringTok{"Vocabulary size: \{nrow(as.matrix(model))\}}\SpecialCharTok{\textbackslash{}n}\StringTok{"}\NormalTok{))}
\end{Highlighting}
\end{Shaded}

\begin{verbatim}
#> Vocabulary size: 1395
\end{verbatim}

\begin{Shaded}
\begin{Highlighting}[]
\ControlFlowTok{if}\NormalTok{ (}\StringTok{"citation"} \SpecialCharTok{\%in\%} \FunctionTok{rownames}\NormalTok{(}\FunctionTok{as.matrix}\NormalTok{(model))) \{}
\NormalTok{  nn }\OtherTok{\textless{}{-}} \FunctionTok{predict}\NormalTok{(model, }\AttributeTok{newdata =} \StringTok{"citation"}\NormalTok{, }\AttributeTok{type =} \StringTok{"nearest"}\NormalTok{, }\AttributeTok{top\_n =} \DecValTok{10}\NormalTok{)}
  \FunctionTok{cat}\NormalTok{(}\StringTok{"Nearest neighbours of \textquotesingle{}citation\textquotesingle{}:}\SpecialCharTok{\textbackslash{}n}\StringTok{"}\NormalTok{)}
  \FunctionTok{print}\NormalTok{(nn)}
\NormalTok{\}}
\end{Highlighting}
\end{Shaded}

\begin{verbatim}
#> Nearest neighbours of 'citation':
#> $citation
#>       term1         term2 similarity rank
#> 1  citation    inequality      0.718    1
#> 2  citation          mean      0.703    2
#> 3  citation     regarding      0.686    3
#> 4  citation        source      0.684    4
#> 5  citation    normalized      0.681    5
#> 6  citation        skewed      0.680    6
#> 7  citation   comparisons      0.675    7
#> 8  citation         links      0.664    8
#> 9  citation    normalised      0.660    9
#> 10 citation heterogeneous      0.654   10
\end{verbatim}

\subsection{Computing document embeddings via mean pooling}\label{computing-document-embeddings-via-mean-pooling}

Apply the same text preprocessing used for training: lowercase before removing
characters outside the vocabulary alphabet. This preserves capitalised words and
acronyms when matching them to the learned word vectors.

\begin{Shaded}
\begin{Highlighting}[]
\NormalTok{word\_matrix }\OtherTok{\textless{}{-}} \FunctionTok{as.matrix}\NormalTok{(model)}

\NormalTok{doc\_embed }\OtherTok{\textless{}{-}} \ControlFlowTok{function}\NormalTok{(text, wm) \{}
\NormalTok{  words }\OtherTok{\textless{}{-}} \FunctionTok{str\_split}\NormalTok{(}\FunctionTok{str\_replace\_all}\NormalTok{(}\FunctionTok{str\_to\_lower}\NormalTok{(text), }\StringTok{"[\^{}a{-}z ]"}\NormalTok{, }\StringTok{" "}\NormalTok{),}
                     \StringTok{"}\SpecialCharTok{\textbackslash{}\textbackslash{}}\StringTok{s+"}\NormalTok{)[[}\DecValTok{1}\NormalTok{]]}
\NormalTok{  matched }\OtherTok{\textless{}{-}}\NormalTok{ words[words }\SpecialCharTok{\%in\%} \FunctionTok{rownames}\NormalTok{(wm)]}
  \ControlFlowTok{if}\NormalTok{ (}\FunctionTok{length}\NormalTok{(matched) }\SpecialCharTok{==} \DecValTok{0}\NormalTok{) }\FunctionTok{return}\NormalTok{(}\FunctionTok{rep}\NormalTok{(}\ConstantTok{NA\_real\_}\NormalTok{, }\FunctionTok{ncol}\NormalTok{(wm)))}
  \FunctionTok{colMeans}\NormalTok{(wm[matched, , }\AttributeTok{drop =} \ConstantTok{FALSE}\NormalTok{])}
\NormalTok{\}}

\NormalTok{doc\_embeddings }\OtherTok{\textless{}{-}} \FunctionTok{map}\NormalTok{(text\_df}\SpecialCharTok{$}\NormalTok{text, \textbackslash{}(t) }\FunctionTok{doc\_embed}\NormalTok{(t, word\_matrix))}
\NormalTok{embed\_mat }\OtherTok{\textless{}{-}} \FunctionTok{do.call}\NormalTok{(rbind, doc\_embeddings)}

\NormalTok{valid }\OtherTok{\textless{}{-}} \FunctionTok{complete.cases}\NormalTok{(embed\_mat)}
\NormalTok{embed\_mat }\OtherTok{\textless{}{-}}\NormalTok{ embed\_mat[valid, ]}
\NormalTok{text\_df\_valid }\OtherTok{\textless{}{-}}\NormalTok{ text\_df[valid, ]}

\FunctionTok{cat}\NormalTok{(}\FunctionTok{glue}\NormalTok{(}\StringTok{"Valid document embeddings: \{nrow(embed\_mat)\}}\SpecialCharTok{\textbackslash{}n}\StringTok{"}\NormalTok{))}
\end{Highlighting}
\end{Shaded}

\begin{verbatim}
#> Valid document embeddings: 108
\end{verbatim}

\subsection{Dimensionality reduction with UMAP}\label{dimensionality-reduction-with-umap}

\begin{Shaded}
\begin{Highlighting}[]
\NormalTok{umap\_result }\OtherTok{\textless{}{-}} \FunctionTok{umap}\NormalTok{(embed\_mat, }\AttributeTok{n\_neighbors =} \DecValTok{15}\NormalTok{, }\AttributeTok{min\_dist =} \FloatTok{0.1}\NormalTok{,}
                    \AttributeTok{n\_components =} \DecValTok{2}\NormalTok{, }\AttributeTok{ret\_model =} \ConstantTok{FALSE}\NormalTok{)}

\NormalTok{text\_df\_valid}\SpecialCharTok{$}\NormalTok{umap\_x }\OtherTok{\textless{}{-}}\NormalTok{ umap\_result[, }\DecValTok{1}\NormalTok{]}
\NormalTok{text\_df\_valid}\SpecialCharTok{$}\NormalTok{umap\_y }\OtherTok{\textless{}{-}}\NormalTok{ umap\_result[, }\DecValTok{2}\NormalTok{]}
\end{Highlighting}
\end{Shaded}

\begin{Shaded}
\begin{Highlighting}[]
\FunctionTok{ggplot}\NormalTok{(text\_df\_valid, }\FunctionTok{aes}\NormalTok{(}\AttributeTok{x =}\NormalTok{ umap\_x, }\AttributeTok{y =}\NormalTok{ umap\_y, }\AttributeTok{colour =} \FunctionTok{factor}\NormalTok{(year))) }\SpecialCharTok{+}
  \FunctionTok{geom\_point}\NormalTok{(}\AttributeTok{alpha =} \FloatTok{0.6}\NormalTok{, }\AttributeTok{size =} \FloatTok{1.5}\NormalTok{) }\SpecialCharTok{+}
  \FunctionTok{scale\_colour\_manual}\NormalTok{(}\AttributeTok{values =} \FunctionTok{palette\_sci}\NormalTok{(}\FunctionTok{n\_distinct}\NormalTok{(text\_df\_valid}\SpecialCharTok{$}\NormalTok{year))) }\SpecialCharTok{+}
  \FunctionTok{labs}\NormalTok{(}\AttributeTok{x =} \StringTok{"UMAP 1"}\NormalTok{, }\AttributeTok{y =} \StringTok{"UMAP 2"}\NormalTok{, }\AttributeTok{colour =} \StringTok{"Year"}\NormalTok{) }\SpecialCharTok{+}
  \FunctionTok{theme\_sci}\NormalTok{() }\SpecialCharTok{+}
  \FunctionTok{theme}\NormalTok{(}\AttributeTok{axis.text =} \FunctionTok{element\_blank}\NormalTok{(), }\AttributeTok{axis.ticks =} \FunctionTok{element\_blank}\NormalTok{())}
\end{Highlighting}
\end{Shaded}

\begin{figure}

{\centering \includegraphics[width=0.9\linewidth,alt={Scatter plot showing documents positioned by semantic similarity in UMAP space. Colour indicates publication year, revealing whether temporal patterns correspond to semantic clusters.}]{scientometrics-in-r_files/figure-latex/24-fig-embed-umap-1} 

}

\caption{UMAP projection of document embeddings, coloured by publication year.}\label{fig:24-fig-embed-umap}
\end{figure}

\subsection{Document similarity search}\label{document-similarity-search}

\begin{Shaded}
\begin{Highlighting}[]
\NormalTok{cosine\_sim }\OtherTok{\textless{}{-}} \ControlFlowTok{function}\NormalTok{(a, b) }\FunctionTok{sum}\NormalTok{(a }\SpecialCharTok{*}\NormalTok{ b) }\SpecialCharTok{/}\NormalTok{ (}\FunctionTok{sqrt}\NormalTok{(}\FunctionTok{sum}\NormalTok{(a}\SpecialCharTok{\^{}}\DecValTok{2}\NormalTok{)) }\SpecialCharTok{*} \FunctionTok{sqrt}\NormalTok{(}\FunctionTok{sum}\NormalTok{(b}\SpecialCharTok{\^{}}\DecValTok{2}\NormalTok{)))}

\NormalTok{query\_idx }\OtherTok{\textless{}{-}} \DecValTok{1}
\NormalTok{sims }\OtherTok{\textless{}{-}} \FunctionTok{map\_dbl}\NormalTok{(}\FunctionTok{seq\_len}\NormalTok{(}\FunctionTok{nrow}\NormalTok{(embed\_mat)),}
\NormalTok{                \textbackslash{}(i) }\FunctionTok{cosine\_sim}\NormalTok{(embed\_mat[query\_idx, ], embed\_mat[i, ]))}

\NormalTok{top\_similar }\OtherTok{\textless{}{-}} \FunctionTok{order}\NormalTok{(sims, }\AttributeTok{decreasing =} \ConstantTok{TRUE}\NormalTok{)[}\DecValTok{2}\SpecialCharTok{:}\DecValTok{6}\NormalTok{]}

\FunctionTok{cat}\NormalTok{(}\FunctionTok{glue}\NormalTok{(}\StringTok{"Query: \{text\_df\_valid$title[query\_idx]\}}\SpecialCharTok{\textbackslash{}n\textbackslash{}n}\StringTok{"}\NormalTok{))}
\end{Highlighting}
\end{Shaded}

\begin{verbatim}
#> Query: Market forces in Italian academia today (and yesterday)
\end{verbatim}

\begin{Shaded}
\begin{Highlighting}[]
\FunctionTok{cat}\NormalTok{(}\StringTok{"Most similar documents:}\SpecialCharTok{\textbackslash{}n}\StringTok{"}\NormalTok{)}
\end{Highlighting}
\end{Shaded}

\begin{verbatim}
#> Most similar documents:
\end{verbatim}

\begin{Shaded}
\begin{Highlighting}[]
\FunctionTok{walk}\NormalTok{(top\_similar, \textbackslash{}(i) \{}
  \FunctionTok{cat}\NormalTok{(}\FunctionTok{glue}\NormalTok{(}\StringTok{"  [\{round(sims[i], 3)\}] \{text\_df\_valid$title[i]\}}\SpecialCharTok{\textbackslash{}n}\StringTok{"}\NormalTok{))}
\NormalTok{\})}
\end{Highlighting}
\end{Shaded}

\begin{verbatim}
#> [0.937] Empirical evidence on the relationship between research and teaching in academia[0.923] Relations between academic reputation and innovation networks[0.909] “Formal and informal networkedness among German Academics”: exploring the role of conferences and co-publications in scientific performance[0.908] Portraying the life cycle of ideas in social psychology through functional (textual) data analysis: a toolkit for digital history[0.907] Measuring academic entities’ impact by content-based citation analysis in a heterogeneous academic network
\end{verbatim}

\subsection{Embedding clusters vs.~citation impact}\label{embedding-clusters-vs.-citation-impact}

\begin{Shaded}
\begin{Highlighting}[]
\FunctionTok{ggplot}\NormalTok{(text\_df\_valid, }\FunctionTok{aes}\NormalTok{(}\AttributeTok{x =}\NormalTok{ umap\_x, }\AttributeTok{y =}\NormalTok{ umap\_y,}
                          \AttributeTok{colour =} \FunctionTok{log1p}\NormalTok{(cited\_by\_count))) }\SpecialCharTok{+}
  \FunctionTok{geom\_point}\NormalTok{(}\AttributeTok{alpha =} \FloatTok{0.6}\NormalTok{, }\AttributeTok{size =} \FloatTok{1.5}\NormalTok{) }\SpecialCharTok{+}
  \FunctionTok{scale\_colour\_viridis\_c}\NormalTok{(}\AttributeTok{option =} \StringTok{"C"}\NormalTok{, }\AttributeTok{name =} \StringTok{"log(cites + 1)"}\NormalTok{) }\SpecialCharTok{+}
  \FunctionTok{labs}\NormalTok{(}\AttributeTok{x =} \StringTok{"UMAP 1"}\NormalTok{, }\AttributeTok{y =} \StringTok{"UMAP 2"}\NormalTok{) }\SpecialCharTok{+}
  \FunctionTok{theme\_sci}\NormalTok{() }\SpecialCharTok{+}
  \FunctionTok{theme}\NormalTok{(}\AttributeTok{axis.text =} \FunctionTok{element\_blank}\NormalTok{(), }\AttributeTok{axis.ticks =} \FunctionTok{element\_blank}\NormalTok{())}
\end{Highlighting}
\end{Shaded}

\begin{figure}

{\centering \includegraphics[width=0.9\linewidth,alt={Scatter plot in UMAP space with documents coloured by log citation count, showing whether highly cited papers cluster together semantically.}]{scientometrics-in-r_files/figure-latex/24-fig-embed-impact-1} 

}

\caption{UMAP embedding coloured by citation impact (log scale).}\label{fig:24-fig-embed-impact}
\end{figure}

\section{Diagnostics and interpretation}\label{embed-diagnostics}

\begin{itemize}
\tightlist
\item
  \textbf{Vocabulary coverage}: Check what fraction of abstract words appear in the embedding vocabulary. Low coverage (\textless{} 80\%) indicates the training corpus is too small or preprocessing is too aggressive.
\item
  \textbf{Nearest neighbours}: Inspect the nearest neighbours of key domain terms. If ``citation'' is closest to ``references'' and ``bibliometric'', the embeddings capture domain structure.
\item
  \textbf{UMAP stability}: UMAP is stochastic. Run with multiple seeds and check whether the global cluster structure is stable. Individual point positions will vary but clusters should persist.
\item
  \textbf{Embedding dimensionality}: 100--300 dimensions is standard. Below 50, embeddings lose semantic resolution; above 300, training becomes slow with diminishing returns.
\end{itemize}

\section{Limitations and responsible use}\label{embed-limitations}

\section{Limitations and responsible use}\label{limitations-and-responsible-use-23}

\begin{itemize}
\tightlist
\item
  \textbf{Embeddings are opaque.} Unlike TF-IDF weights, embedding dimensions have no interpretable meaning. You can measure similarity but not explain \emph{why} two documents are similar without additional analysis.
\item
  \textbf{Training data bias.} Embeddings reflect the biases of their training corpus. If the corpus overrepresents certain topics or perspectives, the embeddings will too.
\item
  \textbf{Small corpora produce poor embeddings.} Word2vec needs tens of thousands of documents for reliable training. For small bibliometric corpora (\textless{} 1,000 papers), use pre-trained embeddings instead.
\item
  \textbf{UMAP/t-SNE distort distances.} These methods preserve local structure but distort global distances. Documents that appear close in 2D may not be the most similar in high-dimensional space \citep{hicks2015}.
\end{itemize}

\section{Common pitfalls}\label{embed-pitfalls}

\section{Common pitfalls}\label{common-pitfalls-23}

\begin{itemize}
\tightlist
\item
  \textbf{Training on too few documents.} Word2vec on 200 abstracts produces unreliable embeddings. Use at least 5,000 documents or switch to pre-trained models.
\item
  \textbf{Averaging without weighting.} Simple mean pooling treats every word equally. TF-IDF-weighted averaging gives more weight to distinctive terms and often produces better document embeddings.
\item
  \textbf{Over-interpreting UMAP clusters.} Visual clusters in UMAP may reflect projection artifacts. Validate clusters with a separate method (e.g., k-means in the full embedding space).
\item
  \textbf{Comparing embeddings from different models.} Embedding spaces are not aligned across different training runs. Cosine similarity is only meaningful within a single embedding space.
\end{itemize}

\section{Exercises}\label{embed-exercises}

\begin{enumerate}
\def\labelenumi{\arabic{enumi}.}
\item
  \textbf{Weighted averaging.} Compute document embeddings using TF-IDF-weighted word vectors instead of simple means. Does the UMAP visualisation change?
\item
  \textbf{K-means clustering.} Apply k-means (k = 5, 8, 12) to the document embeddings. Compare the resulting clusters with LDA topics from \ref{topicmodel-example}. Do they agree?
\item
  \textbf{Pre-trained embeddings.} If available, use SPECTER embeddings (via reticulate and Python) instead of corpus-trained word2vec. How does similarity search quality change?
\item
  \textbf{Analogy test.} Test whether the word2vec model captures analogical relationships (e.g., ``journal'' - ``article'' + ``book'' ≈ ``publisher''). What does the result tell you about the model?
\end{enumerate}

\section{Solutions}\label{24-solutions}

Solutions are provided in \ref{solutions}.

\section{Further reading}\label{embed-reading}

\begin{itemize}
\tightlist
\item
  \citet{silge2017} --- Text embeddings in R, including integration with word2vec.
\item
  \citet{priem2022} --- OpenAlex concept embeddings and their relationship to document content.
\item
  \citet{waltman2012} --- Bibliometric similarity measures as context for embedding-based approaches.
\end{itemize}

\section{Session info}\label{embed-sessioninfo}

\begin{Shaded}
\begin{Highlighting}[]
\FunctionTok{sessionInfo}\NormalTok{()}
\end{Highlighting}
\end{Shaded}

\begin{verbatim}
#> R version 4.6.1 (2026-06-24)
#> Platform: x86_64-pc-linux-gnu
#> Running under: Ubuntu 24.04.5 LTS
#> 
#> Matrix products: default
#> BLAS:   /usr/lib/x86_64-linux-gnu/openblas-pthread/libblas.so.3 
#> LAPACK: /usr/lib/x86_64-linux-gnu/openblas-pthread/libopenblasp-r0.3.26.so;  LAPACK version 3.12.0
#> 
#> locale:
#>  [1] LC_CTYPE=C.UTF-8       LC_NUMERIC=C           LC_TIME=C.UTF-8       
#>  [4] LC_COLLATE=C.UTF-8     LC_MONETARY=C.UTF-8    LC_MESSAGES=C.UTF-8   
#>  [7] LC_PAPER=C.UTF-8       LC_NAME=C              LC_ADDRESS=C          
#> [10] LC_TELEPHONE=C         LC_MEASUREMENT=C.UTF-8 LC_IDENTIFICATION=C   
#> 
#> time zone: UTC
#> tzcode source: system (glibc)
#> 
#> attached base packages:
#> [1] stats     graphics  grDevices utils     datasets  methods   base     
#> 
#> other attached packages:
#>  [1] uwot_0.2.5                Matrix_1.7-5             
#>  [3] word2vec_0.4.1            stm_1.3.8                
#>  [5] topicmodels_0.2-17        quanteda.textstats_0.97.2
#>  [7] visNetwork_2.1.4          ggraph_2.2.2             
#>  [9] tidygraph_1.3.1           igraph_2.3.4             
#> [11] quanteda_4.5.1            pdftools_3.9.1           
#> [13] arrow_26.0.0              bibliometrix_5.5.0       
#> [15] RefManageR_1.4.0          bib2df_1.1.2.0           
#> [17] rcrossref_1.2.1           gt_1.4.0                 
#> [19] tidytext_0.4.3            glue_1.8.1               
#> [21] openalexR_3.1.2           lubridate_1.9.5          
#> [23] forcats_1.0.1             stringr_1.6.0            
#> [25] dplyr_1.2.1               purrr_1.2.2              
#> [27] readr_2.2.0               tidyr_1.3.2              
#> [29] tibble_3.3.1              ggplot2_4.0.3            
#> [31] tidyverse_2.0.0          
#> 
#> loaded via a namespace (and not attached):
#>   [1] bibtex_0.5.3           RColorBrewer_1.1-3     rstudioapi_0.19.0     
#>   [4] jsonlite_2.0.0         magrittr_2.0.5         modeltools_0.2-25     
#>   [7] farver_2.1.2           rmarkdown_2.32         fs_2.1.0              
#>  [10] vctrs_0.7.3            memoise_2.0.1          askpass_1.2.1         
#>  [13] base64enc_0.1-6        htmltools_0.5.9        contentanalysis_1.1.1 
#>  [16] curl_8.1.0             janeaustenr_1.0.0      cellranger_1.1.0      
#>  [19] htmlwidgets_1.6.4      tokenizers_0.3.0       plyr_1.8.9            
#>  [22] httr2_1.3.0            cachem_1.1.0           plotly_4.12.1         
#>  [25] dimensionsR_0.0.3      mime_0.13              lifecycle_1.0.5       
#>  [28] pkgconfig_2.0.3        R6_2.6.1               fastmap_1.2.0         
#>  [31] shiny_1.14.0           digest_0.6.39          patchwork_1.3.2       
#>  [34] shinycssloaders_1.1.0  rprojroot_2.1.1        RSpectra_0.16-2       
#>  [37] SnowballC_0.7.1        labeling_0.4.3         urltools_1.7.3.1      
#>  [40] timechange_0.4.0       httr_1.4.9             polyclip_1.10-7       
#>  [43] compiler_4.6.1         here_1.0.2             bit64_4.8.6           
#>  [46] withr_3.0.3            S7_0.2.2               viridis_0.6.5         
#>  [49] ggforce_0.5.0          MASS_7.3-65            bibliometrixData_0.3.0
#>  [52] tools_4.6.1            otel_0.2.0             stopwords_2.3         
#>  [55] zip_3.0.2              rentrez_1.2.4          httpuv_1.6.17         
#>  [58] promises_1.5.0         grid_4.6.1             stringdist_0.9.17     
#>  [61] reshape2_1.4.5         generics_0.1.4         gtable_0.3.6          
#>  [64] tzdb_0.5.0             rscopus_0.9.0          ca_0.72               
#>  [67] data.table_1.18.6.1    hms_1.1.4              xml2_1.6.0            
#>  [70] utf8_1.2.6             ggrepel_0.9.8          pillar_1.11.1         
#>  [73] nsyllable_1.0.1        vroom_1.7.1            later_1.4.8           
#>  [76] tweenr_2.0.3           lattice_0.22-9         FNN_1.1.4.1           
#>  [79] bit_4.6.0              tidyselect_1.2.1       tm_0.7-20             
#>  [82] miniUI_0.1.2           knitr_1.52             gridExtra_2.3.1       
#>  [85] NLP_0.3-3              bookdown_0.48          stats4_4.6.1          
#>  [88] crul_1.6.0             xfun_0.62              graphlayouts_1.2.5    
#>  [91] matrixStats_1.5.0      DT_0.34.0              humaniformat_0.6.0    
#>  [94] stringi_1.8.9          qpdf_1.4.1             yaml_2.3.12           
#>  [97] evaluate_1.0.5         httpcode_0.3.0         cli_3.6.6             
#> [100] xtable_1.8-8           systemfonts_1.3.2      Rcpp_1.1.2            
#> [103] readxl_1.6.0           triebeard_0.4.1        XML_3.99-0.25         
#> [106] parallel_4.6.1         assertthat_0.2.1       pubmedR_1.0.2         
#> [109] slam_0.1-56            viridisLite_0.4.3      scales_1.4.0          
#> [112] crayon_1.5.3           openxlsx_4.2.9         rlang_1.3.0           
#> [115] fastmatch_1.1-8
\end{verbatim}

\chapter{Detecting Emerging Topics}\label{detecting-emerging-topics}

\emph{R. Heller · Date: 2026-05-09 · Last changed: 2026-10-31}

\section{Learning objectives}\label{emerging-objectives}

After completing this chapter, you will be able to:

\begin{itemize}
\tightlist
\item
  Define what constitutes an ``emerging topic'' in scientometric terms
\item
  Apply burst detection to keyword time series
\item
  Track topic prevalence over time using rolling-window analysis
\item
  Compute novelty and transience indicators for documents
\item
  Distinguish genuine emergence from database coverage artefacts
\end{itemize}

\section{Setup}\label{emerging-setup}

\begin{Shaded}
\begin{Highlighting}[]
\FunctionTok{library}\NormalTok{(tidyverse)}
\FunctionTok{library}\NormalTok{(openalexR)}
\FunctionTok{library}\NormalTok{(quanteda)}
\FunctionTok{library}\NormalTok{(quanteda.textstats)}
\FunctionTok{library}\NormalTok{(tidytext)}
\FunctionTok{library}\NormalTok{(glue)}
\FunctionTok{library}\NormalTok{(gt)}

\FunctionTok{set.seed}\NormalTok{(}\DecValTok{20260509}\NormalTok{)}

\FunctionTok{source}\NormalTok{(here}\SpecialCharTok{::}\FunctionTok{here}\NormalTok{(}\StringTok{"R"}\NormalTok{, }\StringTok{"api\_helpers.R"}\NormalTok{))}
\FunctionTok{source}\NormalTok{(here}\SpecialCharTok{::}\FunctionTok{here}\NormalTok{(}\StringTok{"R"}\NormalTok{, }\StringTok{"utils.R"}\NormalTok{))}
\FunctionTok{source}\NormalTok{(here}\SpecialCharTok{::}\FunctionTok{here}\NormalTok{(}\StringTok{"R"}\NormalTok{, }\StringTok{"sci\_palette.R"}\NormalTok{))}
\end{Highlighting}
\end{Shaded}

\section{Conceptual background}\label{emerging-background}

An \textbf{emerging topic} is a research area experiencing rapid growth in attention --- new publications, new keywords, increasing citation flows. Detecting emergence early has strategic value for researchers, funders, and policymakers. But distinguishing genuine intellectual emergence from database artefacts (a journal newly indexed, a conference proceedings batch-imported) requires careful methodology.

\textbf{Burst detection} formalises the intuition that a keyword or topic has ``burst'' when its frequency increases sharply relative to baseline. \citet{kleinberg2003} introduced a state-based model: a keyword exists in a ``normal'' state (baseline frequency) or a ``burst'' state (elevated frequency). The algorithm identifies time intervals during which each keyword is in its burst state. Cascading bursts at multiple time scales reveal both short-lived spikes and sustained surges.

\textbf{Temporal topic tracking} extends topic models (\ref{topicmodel-background}) by fitting models to time-sliced corpora or using dynamic topic models where topic distributions evolve smoothly over time. A simpler approach uses rolling windows: fit a model to each 3-year window and track how topic proportions change between windows.

\textbf{Novelty and transience indicators} operate at the document level. A paper is \emph{novel} if its content (measured by embedding distance or topic distribution) is distant from prior literature. A paper is \emph{transient} if it is distant from subsequent literature --- it attracted brief attention but did not generate follow-on work. Genuinely emerging topics are characterised by papers that are novel \emph{and} non-transient: they introduce new ideas that the field continues to develop.

\textbf{Growth-rate metrics} are the simplest emergence indicators: the year-over-year growth rate of publications containing a keyword, the acceleration of citation rates, or the increase in the number of distinct institutions working on a topic. These are intuitive but sensitive to baseline effects (a keyword going from 1 to 3 papers is 200\% growth but not meaningful emergence).

\section{Worked example}\label{emerging-example}

\subsection{Preparing temporal keyword data}\label{preparing-temporal-keyword-data}

\begin{Shaded}
\begin{Highlighting}[]
\NormalTok{works }\OtherTok{\textless{}{-}} \FunctionTok{fetch\_openalex}\NormalTok{(}
  \AttributeTok{entity =} \StringTok{"works"}\NormalTok{,}
  \AttributeTok{primary\_location.source.id =} \StringTok{"S148561398"}\NormalTok{,}
  \AttributeTok{from\_publication\_date =} \StringTok{"2015{-}01{-}01"}\NormalTok{,}
  \AttributeTok{to\_publication\_date =} \StringTok{"2023{-}12{-}31"}\NormalTok{,}
  \AttributeTok{options =} \FunctionTok{list}\NormalTok{(}\AttributeTok{sample =} \DecValTok{600}\NormalTok{, }\AttributeTok{seed =} \DecValTok{42}\NormalTok{)}
\NormalTok{)}

\NormalTok{text\_df }\OtherTok{\textless{}{-}}\NormalTok{ works }\SpecialCharTok{|\textgreater{}}
  \FunctionTok{filter}\NormalTok{(}\SpecialCharTok{!}\FunctionTok{is.na}\NormalTok{(abstract), }\FunctionTok{nchar}\NormalTok{(abstract) }\SpecialCharTok{\textgreater{}} \DecValTok{50}\NormalTok{) }\SpecialCharTok{|\textgreater{}}
  \FunctionTok{transmute}\NormalTok{(}
    \AttributeTok{doc\_id =}\NormalTok{ id,}
    \AttributeTok{text =} \FunctionTok{paste}\NormalTok{(display\_name, abstract, }\AttributeTok{sep =} \StringTok{". "}\NormalTok{),}
    \AttributeTok{year =} \FunctionTok{year}\NormalTok{(publication\_date)}
\NormalTok{  )}

\FunctionTok{cat}\NormalTok{(}\FunctionTok{glue}\NormalTok{(}\StringTok{"Documents: \{nrow(text\_df)\}}\SpecialCharTok{\textbackslash{}n}\StringTok{"}\NormalTok{))}
\end{Highlighting}
\end{Shaded}

\begin{verbatim}
#> Documents: 127
\end{verbatim}

\begin{Shaded}
\begin{Highlighting}[]
\FunctionTok{cat}\NormalTok{(}\FunctionTok{glue}\NormalTok{(}\StringTok{"Year range: \{min(text\_df$year)\}{-}{-}\{max(text\_df$year)\}}\SpecialCharTok{\textbackslash{}n}\StringTok{"}\NormalTok{))}
\end{Highlighting}
\end{Shaded}

\begin{verbatim}
#> Year range: 2015--2023
\end{verbatim}

\begin{Shaded}
\begin{Highlighting}[]
\NormalTok{corp }\OtherTok{\textless{}{-}} \FunctionTok{corpus}\NormalTok{(text\_df, }\AttributeTok{docid\_field =} \StringTok{"doc\_id"}\NormalTok{, }\AttributeTok{text\_field =} \StringTok{"text"}\NormalTok{)}
\FunctionTok{docvars}\NormalTok{(corp, }\StringTok{"year"}\NormalTok{) }\OtherTok{\textless{}{-}}\NormalTok{ text\_df}\SpecialCharTok{$}\NormalTok{year}

\NormalTok{toks }\OtherTok{\textless{}{-}} \FunctionTok{tokens}\NormalTok{(corp, }\AttributeTok{remove\_punct =} \ConstantTok{TRUE}\NormalTok{, }\AttributeTok{remove\_numbers =} \ConstantTok{TRUE}\NormalTok{) }\SpecialCharTok{|\textgreater{}}
  \FunctionTok{tokens\_tolower}\NormalTok{() }\SpecialCharTok{|\textgreater{}}
  \FunctionTok{tokens\_remove}\NormalTok{(}\FunctionTok{stopwords}\NormalTok{(}\StringTok{"en"}\NormalTok{)) }\SpecialCharTok{|\textgreater{}}
  \FunctionTok{tokens\_remove}\NormalTok{(}\FunctionTok{c}\NormalTok{(}\StringTok{"study"}\NormalTok{, }\StringTok{"paper"}\NormalTok{, }\StringTok{"results"}\NormalTok{, }\StringTok{"research"}\NormalTok{, }\StringTok{"analysis"}\NormalTok{,}
                   \StringTok{"also"}\NormalTok{, }\StringTok{"however"}\NormalTok{, }\StringTok{"using"}\NormalTok{, }\StringTok{"based"}\NormalTok{))}

\NormalTok{dfmat }\OtherTok{\textless{}{-}} \FunctionTok{dfm}\NormalTok{(toks) }\SpecialCharTok{|\textgreater{}}
  \FunctionTok{dfm\_trim}\NormalTok{(}\AttributeTok{min\_termfreq =} \DecValTok{10}\NormalTok{, }\AttributeTok{min\_docfreq =} \DecValTok{5}\NormalTok{)}

\NormalTok{kw\_by\_year }\OtherTok{\textless{}{-}}\NormalTok{ quanteda}\SpecialCharTok{::}\FunctionTok{convert}\NormalTok{(}\FunctionTok{dfm\_group}\NormalTok{(dfmat, }\AttributeTok{groups =}\NormalTok{ year), }\AttributeTok{to =} \StringTok{"data.frame"}\NormalTok{) }\SpecialCharTok{|\textgreater{}}
  \FunctionTok{pivot\_longer}\NormalTok{(}\SpecialCharTok{{-}}\NormalTok{doc\_id, }\AttributeTok{names\_to =} \StringTok{"term"}\NormalTok{, }\AttributeTok{values\_to =} \StringTok{"count"}\NormalTok{) }\SpecialCharTok{|\textgreater{}}
  \FunctionTok{rename}\NormalTok{(}\AttributeTok{year =}\NormalTok{ doc\_id) }\SpecialCharTok{|\textgreater{}}
  \FunctionTok{mutate}\NormalTok{(}\AttributeTok{year =} \FunctionTok{as.integer}\NormalTok{(year))}

\FunctionTok{cat}\NormalTok{(}\FunctionTok{glue}\NormalTok{(}\StringTok{"Terms tracked: \{n\_distinct(kw\_by\_year$term)\}}\SpecialCharTok{\textbackslash{}n}\StringTok{"}\NormalTok{))}
\end{Highlighting}
\end{Shaded}

\begin{verbatim}
#> Terms tracked: 298
\end{verbatim}

\subsection{Growth rate detection}\label{growth-rate-detection}

\begin{Shaded}
\begin{Highlighting}[]
\NormalTok{growth }\OtherTok{\textless{}{-}}\NormalTok{ kw\_by\_year }\SpecialCharTok{|\textgreater{}}
  \FunctionTok{group\_by}\NormalTok{(term) }\SpecialCharTok{|\textgreater{}}
  \FunctionTok{arrange}\NormalTok{(year) }\SpecialCharTok{|\textgreater{}}
  \FunctionTok{mutate}\NormalTok{(}
    \AttributeTok{growth =}\NormalTok{ (count }\SpecialCharTok{{-}} \FunctionTok{lag}\NormalTok{(count)) }\SpecialCharTok{/} \FunctionTok{pmax}\NormalTok{(}\FunctionTok{lag}\NormalTok{(count), }\DecValTok{1}\NormalTok{),}
    \AttributeTok{total =} \FunctionTok{sum}\NormalTok{(count)}
\NormalTok{  ) }\SpecialCharTok{|\textgreater{}}
  \FunctionTok{ungroup}\NormalTok{() }\SpecialCharTok{|\textgreater{}}
  \FunctionTok{filter}\NormalTok{(}\SpecialCharTok{!}\FunctionTok{is.na}\NormalTok{(growth), total }\SpecialCharTok{\textgreater{}=} \DecValTok{20}\NormalTok{)}

\NormalTok{recent\_growth }\OtherTok{\textless{}{-}}\NormalTok{ growth }\SpecialCharTok{|\textgreater{}}
  \FunctionTok{filter}\NormalTok{(year }\SpecialCharTok{\textgreater{}=} \DecValTok{2021}\NormalTok{) }\SpecialCharTok{|\textgreater{}}
  \FunctionTok{group\_by}\NormalTok{(term) }\SpecialCharTok{|\textgreater{}}
  \FunctionTok{summarise}\NormalTok{(}\AttributeTok{mean\_growth =} \FunctionTok{mean}\NormalTok{(growth), }\AttributeTok{total =} \FunctionTok{first}\NormalTok{(total),}
            \AttributeTok{.groups =} \StringTok{"drop"}\NormalTok{) }\SpecialCharTok{|\textgreater{}}
  \FunctionTok{filter}\NormalTok{(mean\_growth }\SpecialCharTok{\textgreater{}} \DecValTok{0}\NormalTok{) }\SpecialCharTok{|\textgreater{}}
  \FunctionTok{arrange}\NormalTok{(}\FunctionTok{desc}\NormalTok{(mean\_growth))}

\NormalTok{recent\_growth }\SpecialCharTok{|\textgreater{}}
  \FunctionTok{head}\NormalTok{(}\DecValTok{15}\NormalTok{) }\SpecialCharTok{|\textgreater{}}
  \FunctionTok{gt}\NormalTok{() }\SpecialCharTok{|\textgreater{}}
  \FunctionTok{fmt\_percent}\NormalTok{(}\AttributeTok{columns =}\NormalTok{ mean\_growth, }\AttributeTok{decimals =} \DecValTok{0}\NormalTok{)}
\end{Highlighting}
\end{Shaded}

\begingroup
\fontsize{8.0pt}{9.0pt}\selectfont
\begin{longtable}{>{\raggedright\arraybackslash}p{\dimexpr 0.36\linewidth -2\tabcolsep-1.5\arrayrulewidth}>{\raggedleft\arraybackslash}p{\dimexpr 0.35\linewidth -2\tabcolsep-1.5\arrayrulewidth}>{\raggedleft\arraybackslash}p{\dimexpr 0.30\linewidth -2\tabcolsep-1.5\arrayrulewidth}}
\toprule
term & \nolinkurl{mean_growth} & total \\ 
\midrule\addlinespace[2.5pt]
\nolinkurl{countries} & 610\% & 45 \\ 
review & 541\% & 39 \\ 
p & 535\% & 21 \\ 
\nolinkurl{performance} & 438\% & 57 \\ 
peer & 400\% & 27 \\ 
level & 345\% & 30 \\ 
article & 344\% & 46 \\ 
\nolinkurl{universities} & 344\% & 25 \\ 
\nolinkurl{international} & 332\% & 35 \\ 
scores & 300\% & 20 \\ 
\nolinkurl{scientists} & 286\% & 21 \\ 
\nolinkurl{networks} & 277\% & 27 \\ 
\nolinkurl{university} & 277\% & 26 \\ 
top & 241\% & 24 \\ 
\nolinkurl{influence} & 233\% & 24 \\ 
\bottomrule
\end{longtable}
\endgroup

\subsection{Simple burst detection}\label{simple-burst-detection}

\begin{Shaded}
\begin{Highlighting}[]
\NormalTok{detect\_burst }\OtherTok{\textless{}{-}} \ControlFlowTok{function}\NormalTok{(counts, }\AttributeTok{threshold =} \DecValTok{2}\NormalTok{) \{}
\NormalTok{  mean\_count }\OtherTok{\textless{}{-}} \FunctionTok{mean}\NormalTok{(counts)}
\NormalTok{  sd\_count }\OtherTok{\textless{}{-}} \FunctionTok{sd}\NormalTok{(counts)}
  \ControlFlowTok{if}\NormalTok{ (sd\_count }\SpecialCharTok{==} \DecValTok{0}\NormalTok{) }\FunctionTok{return}\NormalTok{(}\FunctionTok{rep}\NormalTok{(}\ConstantTok{FALSE}\NormalTok{, }\FunctionTok{length}\NormalTok{(counts)))}
\NormalTok{  z\_scores }\OtherTok{\textless{}{-}}\NormalTok{ (counts }\SpecialCharTok{{-}}\NormalTok{ mean\_count) }\SpecialCharTok{/}\NormalTok{ sd\_count}
\NormalTok{  z\_scores }\SpecialCharTok{\textgreater{}}\NormalTok{ threshold}
\NormalTok{\}}

\NormalTok{bursts }\OtherTok{\textless{}{-}}\NormalTok{ kw\_by\_year }\SpecialCharTok{|\textgreater{}}
  \FunctionTok{group\_by}\NormalTok{(term) }\SpecialCharTok{|\textgreater{}}
  \FunctionTok{filter}\NormalTok{(}\FunctionTok{n}\NormalTok{() }\SpecialCharTok{\textgreater{}=} \DecValTok{5}\NormalTok{, }\FunctionTok{sum}\NormalTok{(count) }\SpecialCharTok{\textgreater{}=} \DecValTok{20}\NormalTok{) }\SpecialCharTok{|\textgreater{}}
  \FunctionTok{arrange}\NormalTok{(year) }\SpecialCharTok{|\textgreater{}}
  \FunctionTok{mutate}\NormalTok{(}\AttributeTok{is\_burst =} \FunctionTok{detect\_burst}\NormalTok{(count, }\AttributeTok{threshold =} \FloatTok{1.5}\NormalTok{)) }\SpecialCharTok{|\textgreater{}}
  \FunctionTok{ungroup}\NormalTok{()}

\NormalTok{burst\_summary }\OtherTok{\textless{}{-}}\NormalTok{ bursts }\SpecialCharTok{|\textgreater{}}
  \FunctionTok{filter}\NormalTok{(is\_burst) }\SpecialCharTok{|\textgreater{}}
  \FunctionTok{group\_by}\NormalTok{(term) }\SpecialCharTok{|\textgreater{}}
  \FunctionTok{summarise}\NormalTok{(}
    \AttributeTok{burst\_years =} \FunctionTok{paste}\NormalTok{(year, }\AttributeTok{collapse =} \StringTok{", "}\NormalTok{),}
    \AttributeTok{max\_count =} \FunctionTok{max}\NormalTok{(count),}
    \AttributeTok{.groups =} \StringTok{"drop"}
\NormalTok{  ) }\SpecialCharTok{|\textgreater{}}
  \FunctionTok{arrange}\NormalTok{(}\FunctionTok{desc}\NormalTok{(max\_count))}

\NormalTok{burst\_summary }\SpecialCharTok{|\textgreater{}} \FunctionTok{head}\NormalTok{(}\DecValTok{15}\NormalTok{) }\SpecialCharTok{|\textgreater{}} \FunctionTok{gt}\NormalTok{()}
\end{Highlighting}
\end{Shaded}

\begingroup
\fontsize{8.0pt}{9.0pt}\selectfont
\begin{longtable}{>{\raggedright\arraybackslash}p{\dimexpr 0.34\linewidth -2\tabcolsep-1.5\arrayrulewidth}>{\raggedleft\arraybackslash}p{\dimexpr 0.35\linewidth -2\tabcolsep-1.5\arrayrulewidth}>{\raggedleft\arraybackslash}p{\dimexpr 0.31\linewidth -2\tabcolsep-1.5\arrayrulewidth}}
\toprule
term & \nolinkurl{burst_years} & \nolinkurl{max_count} \\ 
\midrule\addlinespace[2.5pt]
\nolinkurl{publications} & 2022 & 37 \\ 
science & 2022 & 36 \\ 
\nolinkurl{scientific} & 2022 & 32 \\ 
\nolinkurl{citation} & 2020 & 29 \\ 
\nolinkurl{citations} & 2022 & 28 \\ 
\nolinkurl{journals} & 2022 & 28 \\ 
authors & 2022 & 27 \\ 
papers & 2017 & 27 \\ 
\nolinkurl{publication} & 2022 & 26 \\ 
\nolinkurl{abstract} & 2022 & 25 \\ 
number & 2022 & 25 \\ 
impact & 2022 & 23 \\ 
journal & 2022 & 22 \\ 
\nolinkurl{different} & 2022 & 21 \\ 
\nolinkurl{countries} & 2022 & 20 \\ 
\bottomrule
\end{longtable}
\endgroup

\subsection{Visualising keyword trajectories}\label{visualising-keyword-trajectories}

\begin{Shaded}
\begin{Highlighting}[]
\NormalTok{top\_burst\_terms }\OtherTok{\textless{}{-}}\NormalTok{ burst\_summary }\SpecialCharTok{|\textgreater{}} \FunctionTok{head}\NormalTok{(}\DecValTok{6}\NormalTok{) }\SpecialCharTok{|\textgreater{}} \FunctionTok{pull}\NormalTok{(term)}

\NormalTok{burst\_plot\_data }\OtherTok{\textless{}{-}}\NormalTok{ bursts }\SpecialCharTok{|\textgreater{}}
  \FunctionTok{filter}\NormalTok{(term }\SpecialCharTok{\%in\%}\NormalTok{ top\_burst\_terms)}

\FunctionTok{ggplot}\NormalTok{(burst\_plot\_data, }\FunctionTok{aes}\NormalTok{(}\AttributeTok{x =}\NormalTok{ year, }\AttributeTok{y =}\NormalTok{ count)) }\SpecialCharTok{+}
  \FunctionTok{geom\_rect}\NormalTok{(}\AttributeTok{data =}\NormalTok{ burst\_plot\_data }\SpecialCharTok{|\textgreater{}} \FunctionTok{filter}\NormalTok{(is\_burst),}
            \FunctionTok{aes}\NormalTok{(}\AttributeTok{xmin =}\NormalTok{ year }\SpecialCharTok{{-}} \FloatTok{0.5}\NormalTok{, }\AttributeTok{xmax =}\NormalTok{ year }\SpecialCharTok{+} \FloatTok{0.5}\NormalTok{,}
                \AttributeTok{ymin =} \SpecialCharTok{{-}}\ConstantTok{Inf}\NormalTok{, }\AttributeTok{ymax =} \ConstantTok{Inf}\NormalTok{),}
            \AttributeTok{fill =} \StringTok{"gold"}\NormalTok{, }\AttributeTok{alpha =} \FloatTok{0.3}\NormalTok{, }\AttributeTok{inherit.aes =} \ConstantTok{FALSE}\NormalTok{) }\SpecialCharTok{+}
  \FunctionTok{geom\_line}\NormalTok{(}\AttributeTok{linewidth =} \FloatTok{0.8}\NormalTok{) }\SpecialCharTok{+}
  \FunctionTok{geom\_point}\NormalTok{(}\AttributeTok{size =} \FloatTok{1.5}\NormalTok{) }\SpecialCharTok{+}
  \FunctionTok{facet\_wrap}\NormalTok{(}\SpecialCharTok{\textasciitilde{}}\NormalTok{ term, }\AttributeTok{scales =} \StringTok{"free\_y"}\NormalTok{, }\AttributeTok{ncol =} \DecValTok{3}\NormalTok{) }\SpecialCharTok{+}
  \FunctionTok{labs}\NormalTok{(}\AttributeTok{x =} \StringTok{"Year"}\NormalTok{, }\AttributeTok{y =} \StringTok{"Frequency"}\NormalTok{) }\SpecialCharTok{+}
  \FunctionTok{theme\_sci}\NormalTok{(}\AttributeTok{base\_size =} \DecValTok{9}\NormalTok{)}
\end{Highlighting}
\end{Shaded}

\begin{figure}

{\centering \includegraphics[width=0.9\linewidth,alt={Line chart showing the frequency of several keywords over time. Burst periods (unusually high frequency) are highlighted with background shading.}]{scientometrics-in-r_files/figure-latex/25-fig-emerging-trajectories-1} 

}

\caption{Temporal trajectories of emerging keywords with burst periods highlighted.}\label{fig:25-fig-emerging-trajectories}
\end{figure}

\subsection{Rolling-window topic tracking}\label{rolling-window-topic-tracking}

\begin{Shaded}
\begin{Highlighting}[]
\NormalTok{window\_size }\OtherTok{\textless{}{-}} \DecValTok{3}
\NormalTok{years }\OtherTok{\textless{}{-}} \FunctionTok{sort}\NormalTok{(}\FunctionTok{unique}\NormalTok{(text\_df}\SpecialCharTok{$}\NormalTok{year))}
\NormalTok{windows }\OtherTok{\textless{}{-}}\NormalTok{ years[}\DecValTok{1}\SpecialCharTok{:}\NormalTok{(}\FunctionTok{length}\NormalTok{(years) }\SpecialCharTok{{-}}\NormalTok{ window\_size }\SpecialCharTok{+} \DecValTok{1}\NormalTok{)]}

\NormalTok{window\_top\_terms }\OtherTok{\textless{}{-}} \FunctionTok{map\_dfr}\NormalTok{(windows, }\ControlFlowTok{function}\NormalTok{(start\_year) \{}
\NormalTok{  end\_year }\OtherTok{\textless{}{-}}\NormalTok{ start\_year }\SpecialCharTok{+}\NormalTok{ window\_size }\SpecialCharTok{{-}} \DecValTok{1}
\NormalTok{  docs }\OtherTok{\textless{}{-}} \FunctionTok{docvars}\NormalTok{(dfmat, }\StringTok{"year"}\NormalTok{) }\SpecialCharTok{\textgreater{}=}\NormalTok{ start\_year }\SpecialCharTok{\&}
          \FunctionTok{docvars}\NormalTok{(dfmat, }\StringTok{"year"}\NormalTok{) }\SpecialCharTok{\textless{}=}\NormalTok{ end\_year}
  \ControlFlowTok{if}\NormalTok{ (}\FunctionTok{sum}\NormalTok{(docs) }\SpecialCharTok{\textless{}} \DecValTok{10}\NormalTok{) }\FunctionTok{return}\NormalTok{(}\FunctionTok{tibble}\NormalTok{())}
\NormalTok{  top }\OtherTok{\textless{}{-}} \FunctionTok{topfeatures}\NormalTok{(dfmat[docs, ], }\DecValTok{10}\NormalTok{)}
  \FunctionTok{tibble}\NormalTok{(}\AttributeTok{window =} \FunctionTok{glue}\NormalTok{(}\StringTok{"\{start\_year\}{-}\{end\_year\}"}\NormalTok{),}
         \AttributeTok{term =} \FunctionTok{names}\NormalTok{(top), }\AttributeTok{frequency =} \FunctionTok{unname}\NormalTok{(top))}
\NormalTok{\})}
\end{Highlighting}
\end{Shaded}

\begin{Shaded}
\begin{Highlighting}[]
\NormalTok{window\_top\_terms }\SpecialCharTok{|\textgreater{}}
  \FunctionTok{group\_by}\NormalTok{(window) }\SpecialCharTok{|\textgreater{}}
  \FunctionTok{slice\_max}\NormalTok{(frequency, }\AttributeTok{n =} \DecValTok{8}\NormalTok{) }\SpecialCharTok{|\textgreater{}}
  \FunctionTok{mutate}\NormalTok{(}\AttributeTok{term =} \FunctionTok{reorder\_within}\NormalTok{(term, frequency, window)) }\SpecialCharTok{|\textgreater{}}
  \FunctionTok{ggplot}\NormalTok{(}\FunctionTok{aes}\NormalTok{(}\AttributeTok{x =}\NormalTok{ frequency, }\AttributeTok{y =}\NormalTok{ term)) }\SpecialCharTok{+}
  \FunctionTok{geom\_col}\NormalTok{(}\AttributeTok{fill =} \FunctionTok{palette\_sci}\NormalTok{(}\DecValTok{1}\NormalTok{)) }\SpecialCharTok{+}
  \FunctionTok{facet\_wrap}\NormalTok{(}\SpecialCharTok{\textasciitilde{}}\NormalTok{ window, }\AttributeTok{scales =} \StringTok{"free\_y"}\NormalTok{, }\AttributeTok{ncol =} \DecValTok{3}\NormalTok{) }\SpecialCharTok{+}
  \FunctionTok{scale\_y\_reordered}\NormalTok{() }\SpecialCharTok{+}
  \FunctionTok{labs}\NormalTok{(}\AttributeTok{x =} \StringTok{"Frequency"}\NormalTok{, }\AttributeTok{y =} \ConstantTok{NULL}\NormalTok{) }\SpecialCharTok{+}
  \FunctionTok{theme\_sci}\NormalTok{(}\AttributeTok{base\_size =} \DecValTok{8}\NormalTok{)}
\end{Highlighting}
\end{Shaded}

\begin{figure}

{\centering \includegraphics[width=0.9\linewidth,alt={Heatmap or faceted bar chart showing the top terms for each rolling time window, revealing which terms gain or lose prominence.}]{scientometrics-in-r_files/figure-latex/25-fig-emerging-windows-1} 

}

\caption{Top terms in each 3-year rolling window, showing vocabulary evolution.}\label{fig:25-fig-emerging-windows}
\end{figure}

\section{Diagnostics and interpretation}\label{emerging-diagnostics}

\begin{itemize}
\tightlist
\item
  \textbf{Baseline effects}: A keyword going from 2 to 6 occurrences shows 200\% growth but may not represent meaningful emergence. Set minimum frequency thresholds before computing growth rates.
\item
  \textbf{Database artefacts}: A newly indexed journal or a batch import of conference proceedings can create spurious ``bursts''. Cross-check with independent evidence.
\item
  \textbf{Burst duration}: Short bursts (1 year) may reflect a single influential paper or a special issue. Sustained bursts (3+ years) more likely indicate genuine emergence.
\item
  \textbf{Window sensitivity}: Rolling-window results change with window width. Report the chosen width and test sensitivity with wider and narrower windows.
\end{itemize}

\section{Limitations and responsible use}\label{emerging-limitations}

\section{Limitations and responsible use}\label{limitations-and-responsible-use-24}

\begin{itemize}
\tightlist
\item
  \textbf{Emergence is retrospective.} By the time bibliometric methods detect emergence, the topic has already gained momentum. These methods identify ``recently emerged'' topics, not future breakthroughs.
\item
  \textbf{Coverage bias.} OpenAlex coverage has expanded rapidly, meaning recent years have more records. Growth rates may reflect database growth rather than intellectual emergence \citep{priem2022}.
\item
  \textbf{Language of emergence.} English-language terms dominate; emerging topics in non-English literatures may be invisible until translated or adopted by English-language publications.
\item
  \textbf{Not predictive.} Burst detection identifies past patterns. It does not predict which topics will continue growing. Many bursts are transient \citep{kleinberg2003}.
\item
  \textbf{Strategic gaming.} If emergence metrics are used for funding decisions, researchers may strategically choose ``hot'' keywords, creating feedback loops \citep{hicks2015}.
\end{itemize}

\section{Common pitfalls}\label{emerging-pitfalls}

\section{Common pitfalls}\label{common-pitfalls-24}

\begin{itemize}
\tightlist
\item
  \textbf{Confusing absolute growth with relative emergence.} A field that publishes 10,000 papers/year growing by 5\% adds 500 papers. A niche area growing from 10 to 30 papers shows 200\% growth. The niche area is ``emerging'' but the large field is adding more absolute volume.
\item
  \textbf{Not controlling for corpus growth.} If the overall number of publications in the corpus grows, all keywords will grow. Use relative frequency (keyword count / total documents) rather than absolute counts.
\item
  \textbf{Overreacting to single-year spikes.} One conference proceedings volume or special issue can create a one-year spike. Require multi-year sustained growth before declaring emergence.
\item
  \textbf{Ignoring synonyms.} ``Deep learning'' and ``neural networks'' may both refer to the same emerging trend. Merge synonyms before analysis.
\end{itemize}

\section{Exercises}\label{emerging-exercises}

\begin{enumerate}
\def\labelenumi{\arabic{enumi}.}
\item
  \textbf{Relative burst detection.} Modify the burst detection to use relative frequency (keyword proportion) instead of absolute counts. Do the same terms still burst?
\item
  \textbf{Emergence by journal.} Compare the most frequent emerging keywords in two different journals. Do they share emerging topics or focus on different frontiers?
\item
  \textbf{Novelty scoring.} For each document, compute its cosine distance to the centroid of all prior-year documents (using TF-IDF vectors). Plot novelty scores over time. Are recent papers becoming more or less novel?
\item
  \textbf{Validation.} Select 5 terms identified as ``emerging'' by burst detection. Search for qualitative evidence (review articles, conference themes) confirming their emergence. What is the false positive rate?
\end{enumerate}

\section{Solutions}\label{25-solutions}

Solutions are provided in \ref{solutions}.

\section{Further reading}\label{emerging-reading}

\begin{itemize}
\tightlist
\item
  \citet{kleinberg2003} --- Bursty patterns in temporal data: the foundational burst detection algorithm.
\item
  \citet{price1963} --- Early observations on the growth of scientific literature and research fronts.
\item
  \citet{priem2022} --- OpenAlex temporal coverage patterns relevant to emergence detection.
\item
  \citet{callon1983} --- Co-word dynamics; an early approach to tracking topical evolution.
\end{itemize}

\section{Session info}\label{emerging-sessioninfo}

\begin{Shaded}
\begin{Highlighting}[]
\FunctionTok{sessionInfo}\NormalTok{()}
\end{Highlighting}
\end{Shaded}

\begin{verbatim}
#> R version 4.6.1 (2026-06-24)
#> Platform: x86_64-pc-linux-gnu
#> Running under: Ubuntu 24.04.5 LTS
#> 
#> Matrix products: default
#> BLAS:   /usr/lib/x86_64-linux-gnu/openblas-pthread/libblas.so.3 
#> LAPACK: /usr/lib/x86_64-linux-gnu/openblas-pthread/libopenblasp-r0.3.26.so;  LAPACK version 3.12.0
#> 
#> locale:
#>  [1] LC_CTYPE=C.UTF-8       LC_NUMERIC=C           LC_TIME=C.UTF-8       
#>  [4] LC_COLLATE=C.UTF-8     LC_MONETARY=C.UTF-8    LC_MESSAGES=C.UTF-8   
#>  [7] LC_PAPER=C.UTF-8       LC_NAME=C              LC_ADDRESS=C          
#> [10] LC_TELEPHONE=C         LC_MEASUREMENT=C.UTF-8 LC_IDENTIFICATION=C   
#> 
#> time zone: UTC
#> tzcode source: system (glibc)
#> 
#> attached base packages:
#> [1] stats     graphics  grDevices utils     datasets  methods   base     
#> 
#> other attached packages:
#>  [1] uwot_0.2.5                Matrix_1.7-5             
#>  [3] word2vec_0.4.1            stm_1.3.8                
#>  [5] topicmodels_0.2-17        quanteda.textstats_0.97.2
#>  [7] visNetwork_2.1.4          ggraph_2.2.2             
#>  [9] tidygraph_1.3.1           igraph_2.3.4             
#> [11] quanteda_4.5.1            pdftools_3.9.1           
#> [13] arrow_26.0.0              bibliometrix_5.5.0       
#> [15] RefManageR_1.4.0          bib2df_1.1.2.0           
#> [17] rcrossref_1.2.1           gt_1.4.0                 
#> [19] tidytext_0.4.3            glue_1.8.1               
#> [21] openalexR_3.1.2           lubridate_1.9.5          
#> [23] forcats_1.0.1             stringr_1.6.0            
#> [25] dplyr_1.2.1               purrr_1.2.2              
#> [27] readr_2.2.0               tidyr_1.3.2              
#> [29] tibble_3.3.1              ggplot2_4.0.3            
#> [31] tidyverse_2.0.0          
#> 
#> loaded via a namespace (and not attached):
#>   [1] bibtex_0.5.3           RColorBrewer_1.1-3     rstudioapi_0.19.0     
#>   [4] jsonlite_2.0.0         magrittr_2.0.5         modeltools_0.2-25     
#>   [7] farver_2.1.2           rmarkdown_2.32         fs_2.1.0              
#>  [10] vctrs_0.7.3            memoise_2.0.1          askpass_1.2.1         
#>  [13] base64enc_0.1-6        htmltools_0.5.9        contentanalysis_1.1.1 
#>  [16] curl_8.1.0             janeaustenr_1.0.0      cellranger_1.1.0      
#>  [19] htmlwidgets_1.6.4      tokenizers_0.3.0       plyr_1.8.9            
#>  [22] httr2_1.3.0            cachem_1.1.0           plotly_4.12.1         
#>  [25] dimensionsR_0.0.3      mime_0.13              lifecycle_1.0.5       
#>  [28] pkgconfig_2.0.3        R6_2.6.1               fastmap_1.2.0         
#>  [31] shiny_1.14.0           digest_0.6.39          patchwork_1.3.2       
#>  [34] shinycssloaders_1.1.0  rprojroot_2.1.1        RSpectra_0.16-2       
#>  [37] SnowballC_0.7.1        labeling_0.4.3         urltools_1.7.3.1      
#>  [40] timechange_0.4.0       httr_1.4.9             polyclip_1.10-7       
#>  [43] compiler_4.6.1         here_1.0.2             bit64_4.8.6           
#>  [46] withr_3.0.3            S7_0.2.2               viridis_0.6.5         
#>  [49] ggforce_0.5.0          MASS_7.3-65            bibliometrixData_0.3.0
#>  [52] tools_4.6.1            otel_0.2.0             stopwords_2.3         
#>  [55] zip_3.0.2              rentrez_1.2.4          httpuv_1.6.17         
#>  [58] promises_1.5.0         grid_4.6.1             stringdist_0.9.17     
#>  [61] reshape2_1.4.5         generics_0.1.4         gtable_0.3.6          
#>  [64] tzdb_0.5.0             rscopus_0.9.0          ca_0.72               
#>  [67] data.table_1.18.6.1    hms_1.1.4              xml2_1.6.0            
#>  [70] utf8_1.2.6             ggrepel_0.9.8          pillar_1.11.1         
#>  [73] nsyllable_1.0.1        vroom_1.7.1            later_1.4.8           
#>  [76] tweenr_2.0.3           lattice_0.22-9         FNN_1.1.4.1           
#>  [79] bit_4.6.0              tidyselect_1.2.1       tm_0.7-20             
#>  [82] miniUI_0.1.2           knitr_1.52             gridExtra_2.3.1       
#>  [85] NLP_0.3-3              bookdown_0.48          stats4_4.6.1          
#>  [88] crul_1.6.0             xfun_0.62              graphlayouts_1.2.5    
#>  [91] matrixStats_1.5.0      DT_0.34.0              humaniformat_0.6.0    
#>  [94] stringi_1.8.9          qpdf_1.4.1             yaml_2.3.12           
#>  [97] evaluate_1.0.5         httpcode_0.3.0         cli_3.6.6             
#> [100] xtable_1.8-8           systemfonts_1.3.2      Rcpp_1.1.2            
#> [103] readxl_1.6.0           triebeard_0.4.1        XML_3.99-0.25         
#> [106] parallel_4.6.1         assertthat_0.2.1       pubmedR_1.0.2         
#> [109] slam_0.1-56            viridisLite_0.4.3      scales_1.4.0          
#> [112] crayon_1.5.3           openxlsx_4.2.9         rlang_1.3.0           
#> [115] fastmatch_1.1-8
\end{verbatim}

\part{Advanced Topics}\label{part-advanced-topics}

\chapter{Altmetrics}\label{altmetrics}

\emph{R. Heller · Date: 2026-05-09 · Last changed: 2026-10-31}

\section{Learning objectives}\label{altmetrics-objectives}

After completing this chapter, you will be able to:

\begin{itemize}
\tightlist
\item
  Define altmetrics and explain how they complement citation-based indicators
\item
  Retrieve altmetric-related data from OpenAlex (open access status, type, counts)
\item
  Compute correlations between citation counts and alternative indicators
\item
  Interpret altmetric signals critically --- distinguishing attention from impact
\item
  Discuss the limitations of altmetrics for research evaluation
\end{itemize}

\section{Setup}\label{altmetrics-setup}

\begin{Shaded}
\begin{Highlighting}[]
\FunctionTok{library}\NormalTok{(tidyverse)}
\FunctionTok{library}\NormalTok{(openalexR)}
\FunctionTok{library}\NormalTok{(glue)}
\FunctionTok{library}\NormalTok{(gt)}

\FunctionTok{set.seed}\NormalTok{(}\DecValTok{20260509}\NormalTok{)}

\FunctionTok{source}\NormalTok{(here}\SpecialCharTok{::}\FunctionTok{here}\NormalTok{(}\StringTok{"R"}\NormalTok{, }\StringTok{"api\_helpers.R"}\NormalTok{))}
\FunctionTok{source}\NormalTok{(here}\SpecialCharTok{::}\FunctionTok{here}\NormalTok{(}\StringTok{"R"}\NormalTok{, }\StringTok{"utils.R"}\NormalTok{))}
\FunctionTok{source}\NormalTok{(here}\SpecialCharTok{::}\FunctionTok{here}\NormalTok{(}\StringTok{"R"}\NormalTok{, }\StringTok{"sci\_palette.R"}\NormalTok{))}
\end{Highlighting}
\end{Shaded}

\section{Conceptual background}\label{altmetrics-background}

Citations measure one dimension of scholarly impact: influence on subsequent research. But scholarly work can also influence policy, education, clinical practice, and public understanding --- none of which are captured by citation counts. \textbf{Altmetrics} (alternative metrics) attempt to fill this gap by tracking online traces of attention: social media mentions, news coverage, blog posts, Wikipedia citations, Mendeley reader counts, and policy document references.

The term was coined by \citet{priem2010}, who argued that the speed and diversity of online scholarly communication demanded new metrics beyond the Journal Impact Factor. Altmetrics are typically faster than citations (social media attention peaks within days of publication, while citations take years) and broader (they capture engagement by non-researchers).

However, altmetrics come with significant caveats. Social media attention correlates weakly with scientific quality. Metrics are easily gamed (bots, self-promotion). Coverage is uneven: biomedical research attracts far more Twitter/X attention than mathematics or engineering. And the platforms providing altmetric data (Altmetric.com, PlumX) are commercial products with proprietary algorithms \citep{hicks2015}.

OpenAlex provides open-access status and citation counts but does not directly provide social media metrics. For altmetric data, external APIs (Altmetric.com, Crossref Event Data) are needed. In this chapter, we demonstrate what can be done with OpenAlex metadata alone and discuss how to extend the analysis with external sources.

\section{Worked example}\label{altmetrics-example}

\subsection{Fetching data with OA indicators}\label{fetching-data-with-oa-indicators}

\begin{Shaded}
\begin{Highlighting}[]
\NormalTok{works }\OtherTok{\textless{}{-}} \FunctionTok{fetch\_openalex}\NormalTok{(}
  \AttributeTok{entity =} \StringTok{"works"}\NormalTok{,}
  \AttributeTok{primary\_location.source.id =} \StringTok{"S148561398"}\NormalTok{,}
  \AttributeTok{from\_publication\_date =} \StringTok{"2020{-}01{-}01"}\NormalTok{,}
  \AttributeTok{to\_publication\_date =} \StringTok{"2023{-}12{-}31"}\NormalTok{,}
  \AttributeTok{options =} \FunctionTok{list}\NormalTok{(}\AttributeTok{sample =} \DecValTok{400}\NormalTok{, }\AttributeTok{seed =} \DecValTok{42}\NormalTok{)}
\NormalTok{)}

\NormalTok{works\_slim }\OtherTok{\textless{}{-}}\NormalTok{ works }\SpecialCharTok{|\textgreater{}}
  \FunctionTok{transmute}\NormalTok{(}
\NormalTok{    id, display\_name,}
    \AttributeTok{year =} \FunctionTok{year}\NormalTok{(publication\_date),}
\NormalTok{    cited\_by\_count,}
    \AttributeTok{oa\_status =}\NormalTok{ oa\_status,}
\NormalTok{    type,}
\NormalTok{    referenced\_works\_count}
\NormalTok{  )}

\FunctionTok{cat}\NormalTok{(}\FunctionTok{glue}\NormalTok{(}\StringTok{"Works retrieved: \{nrow(works\_slim)\}}\SpecialCharTok{\textbackslash{}n}\StringTok{"}\NormalTok{))}
\end{Highlighting}
\end{Shaded}

\begin{verbatim}
#> Works retrieved: 400
\end{verbatim}

\subsection{Open access as an altmetric dimension}\label{open-access-as-an-altmetric-dimension}

\begin{Shaded}
\begin{Highlighting}[]
\NormalTok{oa\_summary }\OtherTok{\textless{}{-}}\NormalTok{ works\_slim }\SpecialCharTok{|\textgreater{}}
  \FunctionTok{count}\NormalTok{(oa\_status, }\AttributeTok{sort =} \ConstantTok{TRUE}\NormalTok{) }\SpecialCharTok{|\textgreater{}}
  \FunctionTok{mutate}\NormalTok{(}\AttributeTok{pct =}\NormalTok{ scales}\SpecialCharTok{::}\FunctionTok{percent}\NormalTok{(n }\SpecialCharTok{/} \FunctionTok{sum}\NormalTok{(n)))}

\NormalTok{oa\_summary }\SpecialCharTok{|\textgreater{}} \FunctionTok{gt}\NormalTok{()}
\end{Highlighting}
\end{Shaded}

\begingroup
\fontsize{8.0pt}{9.0pt}\selectfont
\begin{longtable}{>{\raggedright\arraybackslash}p{\dimexpr 0.35\linewidth -2\tabcolsep-1.5\arrayrulewidth}>{\raggedleft\arraybackslash}p{\dimexpr 0.33\linewidth -2\tabcolsep-1.5\arrayrulewidth}>{\raggedleft\arraybackslash}p{\dimexpr 0.33\linewidth -2\tabcolsep-1.5\arrayrulewidth}}
\toprule
\nolinkurl{oa_status} & n & pct \\ 
\midrule\addlinespace[2.5pt]
closed & 189 & 47\% \\ 
hybrid & 116 & 29\% \\ 
green & 69 & 17\% \\ 
bronze & 26 & 6\% \\ 
\bottomrule
\end{longtable}
\endgroup

\begin{Shaded}
\begin{Highlighting}[]
\NormalTok{works\_slim }\SpecialCharTok{|\textgreater{}}
  \FunctionTok{filter}\NormalTok{(}\SpecialCharTok{!}\FunctionTok{is.na}\NormalTok{(oa\_status)) }\SpecialCharTok{|\textgreater{}}
  \FunctionTok{ggplot}\NormalTok{(}\FunctionTok{aes}\NormalTok{(}\AttributeTok{x =}\NormalTok{ oa\_status, }\AttributeTok{y =}\NormalTok{ cited\_by\_count }\SpecialCharTok{+} \DecValTok{1}\NormalTok{)) }\SpecialCharTok{+}
  \FunctionTok{geom\_boxplot}\NormalTok{(}\AttributeTok{fill =} \FunctionTok{palette\_sci}\NormalTok{(}\DecValTok{1}\NormalTok{), }\AttributeTok{alpha =} \FloatTok{0.7}\NormalTok{) }\SpecialCharTok{+}
  \FunctionTok{scale\_y\_log10}\NormalTok{() }\SpecialCharTok{+}
  \FunctionTok{labs}\NormalTok{(}\AttributeTok{x =} \StringTok{"OA status"}\NormalTok{, }\AttributeTok{y =} \StringTok{"Citations (log scale)"}\NormalTok{) }\SpecialCharTok{+}
  \FunctionTok{theme\_sci}\NormalTok{()}
\end{Highlighting}
\end{Shaded}

\begin{figure}

{\centering \includegraphics[width=0.9\linewidth,alt={Box plot comparing citation counts across different OA categories (gold, green, hybrid, bronze, closed). Shows whether OA articles receive more citations.}]{scientometrics-in-r_files/figure-latex/26-fig-altmetrics-oa-citations-1} 

}

\caption{Citation distributions by open access status.}\label{fig:26-fig-altmetrics-oa-citations}
\end{figure}

\subsection{Citation count distribution}\label{citation-count-distribution}

\begin{Shaded}
\begin{Highlighting}[]
\FunctionTok{ggplot}\NormalTok{(works\_slim, }\FunctionTok{aes}\NormalTok{(}\AttributeTok{x =}\NormalTok{ cited\_by\_count)) }\SpecialCharTok{+}
  \FunctionTok{geom\_histogram}\NormalTok{(}\AttributeTok{binwidth =} \DecValTok{5}\NormalTok{, }\AttributeTok{fill =} \FunctionTok{palette\_sci}\NormalTok{(}\DecValTok{1}\NormalTok{), }\AttributeTok{colour =} \StringTok{"white"}\NormalTok{) }\SpecialCharTok{+}
  \FunctionTok{labs}\NormalTok{(}\AttributeTok{x =} \StringTok{"Citation count"}\NormalTok{, }\AttributeTok{y =} \StringTok{"Number of papers"}\NormalTok{) }\SpecialCharTok{+}
  \FunctionTok{theme\_sci}\NormalTok{()}
\end{Highlighting}
\end{Shaded}

\begin{figure}

{\centering \includegraphics[width=0.9\linewidth,alt={Histogram of citation counts showing a heavy right tail — most papers have few citations, a few have many.}]{scientometrics-in-r_files/figure-latex/26-fig-altmetrics-citedist-1} 

}

\caption{Citation distribution of the sample, showing the characteristic right skew.}\label{fig:26-fig-altmetrics-citedist}
\end{figure}

\subsection{Proxy altmetric: reference list length}\label{proxy-altmetric-reference-list-length}

Reference list length correlates with engagement breadth and can serve as a simple proxy for scholarly thoroughness.

\begin{Shaded}
\begin{Highlighting}[]
\FunctionTok{ggplot}\NormalTok{(works\_slim, }\FunctionTok{aes}\NormalTok{(}\AttributeTok{x =}\NormalTok{ referenced\_works\_count, }\AttributeTok{y =}\NormalTok{ cited\_by\_count)) }\SpecialCharTok{+}
  \FunctionTok{geom\_point}\NormalTok{(}\AttributeTok{alpha =} \FloatTok{0.3}\NormalTok{, }\AttributeTok{colour =} \FunctionTok{palette\_sci}\NormalTok{(}\DecValTok{1}\NormalTok{)) }\SpecialCharTok{+}
  \FunctionTok{geom\_smooth}\NormalTok{(}\AttributeTok{method =} \StringTok{"lm"}\NormalTok{, }\AttributeTok{se =} \ConstantTok{FALSE}\NormalTok{, }\AttributeTok{colour =} \FunctionTok{palette\_sci}\NormalTok{(}\DecValTok{2}\NormalTok{)[}\DecValTok{2}\NormalTok{]) }\SpecialCharTok{+}
  \FunctionTok{labs}\NormalTok{(}\AttributeTok{x =} \StringTok{"Number of references"}\NormalTok{, }\AttributeTok{y =} \StringTok{"Citation count"}\NormalTok{) }\SpecialCharTok{+}
  \FunctionTok{theme\_sci}\NormalTok{()}
\end{Highlighting}
\end{Shaded}

\begin{figure}

{\centering \includegraphics[width=0.9\linewidth,alt={Scatter plot showing a positive but noisy relationship between the number of references a paper cites and the number of citations it receives.}]{scientometrics-in-r_files/figure-latex/26-fig-altmetrics-refcount-1} 

}

\caption{Relationship between reference list length and citation count.}\label{fig:26-fig-altmetrics-refcount}
\end{figure}

\begin{Shaded}
\begin{Highlighting}[]
\NormalTok{cor\_val }\OtherTok{\textless{}{-}} \FunctionTok{cor}\NormalTok{(works\_slim}\SpecialCharTok{$}\NormalTok{referenced\_works\_count, works\_slim}\SpecialCharTok{$}\NormalTok{cited\_by\_count,}
               \AttributeTok{use =} \StringTok{"complete.obs"}\NormalTok{)}
\FunctionTok{cat}\NormalTok{(}\FunctionTok{glue}\NormalTok{(}\StringTok{"Pearson correlation (references vs. citations): \{round(cor\_val, 3)\}}\SpecialCharTok{\textbackslash{}n}\StringTok{"}\NormalTok{))}
\end{Highlighting}
\end{Shaded}

\begin{verbatim}
#> Pearson correlation (references vs. citations): 0.17
\end{verbatim}

\section{Diagnostics and interpretation}\label{altmetrics-diagnostics}

\begin{itemize}
\tightlist
\item
  \textbf{OA advantage}: Open-access articles often receive more citations, but this may reflect self-selection (higher-quality papers are more likely to be made OA) rather than the causal effect of open access.
\item
  \textbf{Platform coverage}: Social media altmetrics are biased toward English-language, biomedical, and controversial topics. Low altmetric scores do not indicate low impact.
\item
  \textbf{Temporal dynamics}: Altmetric attention peaks immediately after publication and decays rapidly. Citations follow a slower, more sustained trajectory.
\item
  \textbf{Gaming risk}: Social media metrics are more easily manipulated than citations. Interpret high altmetric scores with caution.
\end{itemize}

\section{Limitations and responsible use}\label{altmetrics-limitations}

\section{Limitations and responsible use}\label{limitations-and-responsible-use-25}

\begin{itemize}
\tightlist
\item
  \textbf{Attention ≠ impact.} A viral tweet about a paper does not mean the paper is scientifically significant. Altmetrics measure visibility, not quality \citep{hicks2015}.
\item
  \textbf{Disciplinary bias.} Fields with active social media communities (public health, climate science) generate far more altmetric data than fields that do not (pure mathematics, medieval history).
\item
  \textbf{Platform dependence.} Altmetric data depends on commercial platforms whose coverage, algorithms, and existence are not guaranteed.
\item
  \textbf{Never use altmetrics for individual evaluation.} The Leiden Manifesto and DORA both caution against using any single metric --- including altmetrics --- as a proxy for research quality \citep{dora2012}.
\end{itemize}

\section{Common pitfalls}\label{altmetrics-pitfalls}

\section{Common pitfalls}\label{common-pitfalls-25}

\begin{itemize}
\tightlist
\item
  \textbf{Treating Mendeley readers as ``citations.''} Mendeley reader counts reflect who saved a paper, not who cited it. The overlap is partial.
\item
  \textbf{Comparing altmetrics across fields.} A paper in public health with 50 tweets is unremarkable; the same count in algebraic topology is extraordinary.
\item
  \textbf{Ignoring bots and self-promotion.} Some social media mentions are from automated accounts or coordinated campaigns. Commercial altmetric services attempt to filter these, but coverage is imperfect.
\item
  \textbf{Confusing the Altmetric Attention Score with a quality indicator.} The score is a weighted sum of online mentions, not a peer-reviewed quality assessment.
\end{itemize}

\section{Exercises}\label{altmetrics-exercises}

\begin{enumerate}
\def\labelenumi{\arabic{enumi}.}
\item
  \textbf{OA citation advantage.} For two journals with different OA rates, compare mean citations for OA vs.~closed-access papers. Is the ``OA advantage'' consistent across journals?
\item
  \textbf{Temporal attention.} For a set of recent papers, track how citation counts change between 6 months and 2 years after publication. Do papers with high early attention (many references, rapid initial citations) sustain their advantage?
\item
  \textbf{Cross-field comparison.} Fetch samples from a biomedical journal and a mathematics journal. Compare their citation distributions and reference list lengths. What does this suggest about field-specific attention patterns?
\end{enumerate}

\section{Solutions}\label{26-solutions}

Solutions are provided in \ref{solutions}.

\section{Further reading}\label{altmetrics-reading}

\begin{itemize}
\tightlist
\item
  \citet{priem2010} --- The altmetrics manifesto: the original call for alternative impact measures.
\item
  \citet{hicks2015} --- The Leiden Manifesto: responsible metrics, including cautions about altmetrics.
\item
  \citet{dora2012} --- DORA: principles for research assessment beyond the JIF.
\item
  \citet{priem2022} --- OpenAlex open-access classification and metadata.
\end{itemize}

\section{Session info}\label{altmetrics-sessioninfo}

\begin{Shaded}
\begin{Highlighting}[]
\FunctionTok{sessionInfo}\NormalTok{()}
\end{Highlighting}
\end{Shaded}

\begin{verbatim}
#> R version 4.6.1 (2026-06-24)
#> Platform: x86_64-pc-linux-gnu
#> Running under: Ubuntu 24.04.5 LTS
#> 
#> Matrix products: default
#> BLAS:   /usr/lib/x86_64-linux-gnu/openblas-pthread/libblas.so.3 
#> LAPACK: /usr/lib/x86_64-linux-gnu/openblas-pthread/libopenblasp-r0.3.26.so;  LAPACK version 3.12.0
#> 
#> locale:
#>  [1] LC_CTYPE=C.UTF-8       LC_NUMERIC=C           LC_TIME=C.UTF-8       
#>  [4] LC_COLLATE=C.UTF-8     LC_MONETARY=C.UTF-8    LC_MESSAGES=C.UTF-8   
#>  [7] LC_PAPER=C.UTF-8       LC_NAME=C              LC_ADDRESS=C          
#> [10] LC_TELEPHONE=C         LC_MEASUREMENT=C.UTF-8 LC_IDENTIFICATION=C   
#> 
#> time zone: UTC
#> tzcode source: system (glibc)
#> 
#> attached base packages:
#> [1] stats     graphics  grDevices utils     datasets  methods   base     
#> 
#> other attached packages:
#>  [1] uwot_0.2.5                Matrix_1.7-5             
#>  [3] word2vec_0.4.1            stm_1.3.8                
#>  [5] topicmodels_0.2-17        quanteda.textstats_0.97.2
#>  [7] visNetwork_2.1.4          ggraph_2.2.2             
#>  [9] tidygraph_1.3.1           igraph_2.3.4             
#> [11] quanteda_4.5.1            pdftools_3.9.1           
#> [13] arrow_26.0.0              bibliometrix_5.5.0       
#> [15] RefManageR_1.4.0          bib2df_1.1.2.0           
#> [17] rcrossref_1.2.1           gt_1.4.0                 
#> [19] tidytext_0.4.3            glue_1.8.1               
#> [21] openalexR_3.1.2           lubridate_1.9.5          
#> [23] forcats_1.0.1             stringr_1.6.0            
#> [25] dplyr_1.2.1               purrr_1.2.2              
#> [27] readr_2.2.0               tidyr_1.3.2              
#> [29] tibble_3.3.1              ggplot2_4.0.3            
#> [31] tidyverse_2.0.0          
#> 
#> loaded via a namespace (and not attached):
#>   [1] bibtex_0.5.3           RColorBrewer_1.1-3     rstudioapi_0.19.0     
#>   [4] jsonlite_2.0.0         magrittr_2.0.5         modeltools_0.2-25     
#>   [7] farver_2.1.2           rmarkdown_2.32         fs_2.1.0              
#>  [10] vctrs_0.7.3            memoise_2.0.1          askpass_1.2.1         
#>  [13] base64enc_0.1-6        htmltools_0.5.9        contentanalysis_1.1.1 
#>  [16] curl_8.1.0             janeaustenr_1.0.0      cellranger_1.1.0      
#>  [19] htmlwidgets_1.6.4      tokenizers_0.3.0       plyr_1.8.9            
#>  [22] httr2_1.3.0            cachem_1.1.0           plotly_4.12.1         
#>  [25] dimensionsR_0.0.3      mime_0.13              lifecycle_1.0.5       
#>  [28] pkgconfig_2.0.3        R6_2.6.1               fastmap_1.2.0         
#>  [31] shiny_1.14.0           digest_0.6.39          patchwork_1.3.2       
#>  [34] shinycssloaders_1.1.0  rprojroot_2.1.1        RSpectra_0.16-2       
#>  [37] SnowballC_0.7.1        labeling_0.4.3         urltools_1.7.3.1      
#>  [40] timechange_0.4.0       mgcv_1.9-4             httr_1.4.9            
#>  [43] polyclip_1.10-7        compiler_4.6.1         here_1.0.2            
#>  [46] bit64_4.8.6            withr_3.0.3            S7_0.2.2              
#>  [49] viridis_0.6.5          ggforce_0.5.0          MASS_7.3-65           
#>  [52] bibliometrixData_0.3.0 tools_4.6.1            otel_0.2.0            
#>  [55] stopwords_2.3          zip_3.0.2              rentrez_1.2.4         
#>  [58] httpuv_1.6.17          nlme_3.1-169           promises_1.5.0        
#>  [61] grid_4.6.1             stringdist_0.9.17      reshape2_1.4.5        
#>  [64] generics_0.1.4         gtable_0.3.6           tzdb_0.5.0            
#>  [67] rscopus_0.9.0          ca_0.72                data.table_1.18.6.1   
#>  [70] hms_1.1.4              xml2_1.6.0             utf8_1.2.6            
#>  [73] ggrepel_0.9.8          pillar_1.11.1          nsyllable_1.0.1       
#>  [76] vroom_1.7.1            later_1.4.8            splines_4.6.1         
#>  [79] tweenr_2.0.3           lattice_0.22-9         FNN_1.1.4.1           
#>  [82] bit_4.6.0              tidyselect_1.2.1       tm_0.7-20             
#>  [85] miniUI_0.1.2           knitr_1.52             gridExtra_2.3.1       
#>  [88] NLP_0.3-3              bookdown_0.48          stats4_4.6.1          
#>  [91] crul_1.6.0             xfun_0.62              graphlayouts_1.2.5    
#>  [94] matrixStats_1.5.0      DT_0.34.0              humaniformat_0.6.0    
#>  [97] stringi_1.8.9          qpdf_1.4.1             yaml_2.3.12           
#> [100] evaluate_1.0.5         httpcode_0.3.0         cli_3.6.6             
#> [103] xtable_1.8-8           systemfonts_1.3.2      Rcpp_1.1.2            
#> [106] readxl_1.6.0           triebeard_0.4.1        XML_3.99-0.25         
#> [109] parallel_4.6.1         assertthat_0.2.1       pubmedR_1.0.2         
#> [112] slam_0.1-56            viridisLite_0.4.3      scales_1.4.0          
#> [115] crayon_1.5.3           openxlsx_4.2.9         rlang_1.3.0           
#> [118] fastmatch_1.1-8
\end{verbatim}

\chapter{Open Science Indicators}\label{open-science-indicators}

\emph{R. Heller · Date: 2026-05-09 · Last changed: 2026-10-31}

\section{Learning objectives}\label{openscience-objectives}

After completing this chapter, you will be able to:

\begin{itemize}
\tightlist
\item
  Classify open access types (gold, green, hybrid, bronze) using OpenAlex metadata
\item
  Compute OA rates by year, journal, and discipline
\item
  Track the growth of preprint posting and preprint-to-publication pathways
\item
  Assess data availability statement coverage in a corpus
\item
  Discuss the relationship between open science practices and citation impact
\end{itemize}

\section{Setup}\label{openscience-setup}

\begin{Shaded}
\begin{Highlighting}[]
\FunctionTok{library}\NormalTok{(tidyverse)}
\FunctionTok{library}\NormalTok{(openalexR)}
\FunctionTok{library}\NormalTok{(glue)}
\FunctionTok{library}\NormalTok{(gt)}

\FunctionTok{set.seed}\NormalTok{(}\DecValTok{20260509}\NormalTok{)}

\FunctionTok{source}\NormalTok{(here}\SpecialCharTok{::}\FunctionTok{here}\NormalTok{(}\StringTok{"R"}\NormalTok{, }\StringTok{"api\_helpers.R"}\NormalTok{))}
\FunctionTok{source}\NormalTok{(here}\SpecialCharTok{::}\FunctionTok{here}\NormalTok{(}\StringTok{"R"}\NormalTok{, }\StringTok{"utils.R"}\NormalTok{))}
\FunctionTok{source}\NormalTok{(here}\SpecialCharTok{::}\FunctionTok{here}\NormalTok{(}\StringTok{"R"}\NormalTok{, }\StringTok{"sci\_palette.R"}\NormalTok{))}
\end{Highlighting}
\end{Shaded}

\section{Conceptual background}\label{openscience-background}

The open science movement aims to make research processes and outputs freely accessible. Bibliometric analysis can measure the extent to which these ideals are being realised, track trends over time, and identify disparities across disciplines and geographies.

\textbf{Open access (OA)} is the most measurable dimension of open science. OpenAlex classifies each work's OA status using Unpaywall data: \textbf{gold} (published in a fully OA journal), \textbf{green} (available in a repository), \textbf{hybrid} (OA in a subscription journal via APC), \textbf{bronze} (free to read but without an explicit open license), or \textbf{closed} \citep{priem2022}.

\textbf{Preprints} are manuscripts posted publicly before peer review. The growth of preprint servers (bioRxiv, arXiv, medRxiv) has accelerated since 2015, with a dramatic spike during the COVID-19 pandemic. Bibliometric analysis can track preprint posting rates, the time from preprint to journal publication, and whether preprints receive more citations than non-preprint papers.

\textbf{Data availability statements} (DAS) indicate whether the data underlying a study are accessible. Their prevalence is growing as journals adopt mandatory data-sharing policies, but the quality of DAS varies: many state ``data available upon request'' without providing actual access.

Measuring open science is important for policy. Funders increasingly require OA publication and data sharing; institutions report OA rates as indicators of compliance and impact. However, \citet{hicks2015} caution that metrics-driven OA mandates can produce perverse incentives (predatory journals, superficial compliance).

\section{Worked example}\label{openscience-example}

\subsection{OA rates by year}\label{oa-rates-by-year}

\begin{Shaded}
\begin{Highlighting}[]
\NormalTok{works }\OtherTok{\textless{}{-}} \FunctionTok{fetch\_openalex}\NormalTok{(}
  \AttributeTok{entity =} \StringTok{"works"}\NormalTok{,}
  \AttributeTok{primary\_location.source.id =} \StringTok{"S148561398"}\NormalTok{,}
  \AttributeTok{from\_publication\_date =} \StringTok{"2015{-}01{-}01"}\NormalTok{,}
  \AttributeTok{to\_publication\_date =} \StringTok{"2023{-}12{-}31"}\NormalTok{,}
  \AttributeTok{type =} \StringTok{"article"}\NormalTok{,}
  \AttributeTok{options =} \FunctionTok{list}\NormalTok{(}\AttributeTok{sample =} \DecValTok{600}\NormalTok{, }\AttributeTok{seed =} \DecValTok{42}\NormalTok{)}
\NormalTok{)}

\NormalTok{works\_oa }\OtherTok{\textless{}{-}}\NormalTok{ works }\SpecialCharTok{|\textgreater{}}
  \FunctionTok{transmute}\NormalTok{(}
\NormalTok{    id, }\AttributeTok{year =} \FunctionTok{year}\NormalTok{(publication\_date),}
    \AttributeTok{oa\_status =}\NormalTok{ oa\_status,}
\NormalTok{    cited\_by\_count}
\NormalTok{  )}

\FunctionTok{cat}\NormalTok{(}\FunctionTok{glue}\NormalTok{(}\StringTok{"Works: \{nrow(works\_oa)\}}\SpecialCharTok{\textbackslash{}n}\StringTok{"}\NormalTok{))}
\end{Highlighting}
\end{Shaded}

\begin{verbatim}
#> Works: 600
\end{verbatim}

\begin{Shaded}
\begin{Highlighting}[]
\NormalTok{oa\_by\_year }\OtherTok{\textless{}{-}}\NormalTok{ works\_oa }\SpecialCharTok{|\textgreater{}}
  \FunctionTok{count}\NormalTok{(year, oa\_status) }\SpecialCharTok{|\textgreater{}}
  \FunctionTok{group\_by}\NormalTok{(year) }\SpecialCharTok{|\textgreater{}}
  \FunctionTok{mutate}\NormalTok{(}\AttributeTok{pct =}\NormalTok{ n }\SpecialCharTok{/} \FunctionTok{sum}\NormalTok{(n)) }\SpecialCharTok{|\textgreater{}}
  \FunctionTok{ungroup}\NormalTok{()}

\FunctionTok{ggplot}\NormalTok{(oa\_by\_year, }\FunctionTok{aes}\NormalTok{(}\AttributeTok{x =} \FunctionTok{factor}\NormalTok{(year), }\AttributeTok{y =}\NormalTok{ pct, }\AttributeTok{fill =}\NormalTok{ oa\_status)) }\SpecialCharTok{+}
  \FunctionTok{geom\_col}\NormalTok{() }\SpecialCharTok{+}
  \FunctionTok{scale\_y\_continuous}\NormalTok{(}\AttributeTok{labels =}\NormalTok{ scales}\SpecialCharTok{::}\NormalTok{percent) }\SpecialCharTok{+}
  \FunctionTok{scale\_fill\_manual}\NormalTok{(}\AttributeTok{values =} \FunctionTok{palette\_sci}\NormalTok{(}\FunctionTok{n\_distinct}\NormalTok{(oa\_by\_year}\SpecialCharTok{$}\NormalTok{oa\_status))) }\SpecialCharTok{+}
  \FunctionTok{labs}\NormalTok{(}\AttributeTok{x =} \StringTok{"Publication year"}\NormalTok{, }\AttributeTok{y =} \StringTok{"Proportion"}\NormalTok{, }\AttributeTok{fill =} \StringTok{"OA status"}\NormalTok{) }\SpecialCharTok{+}
  \FunctionTok{theme\_sci}\NormalTok{()}
\end{Highlighting}
\end{Shaded}

\begin{figure}

{\centering \includegraphics[width=0.9\linewidth,alt={Stacked bar chart showing the proportion of articles in each OA category (gold, green, hybrid, bronze, closed) for each publication year from 2015 to 2023.}]{scientometrics-in-r_files/figure-latex/27-fig-openscience-oa-trend-1} 

}

\caption{Open access rates by year and OA type.}\label{fig:27-fig-openscience-oa-trend}
\end{figure}

\subsection{OA citation advantage}\label{oa-citation-advantage}

\begin{Shaded}
\begin{Highlighting}[]
\NormalTok{oa\_cites }\OtherTok{\textless{}{-}}\NormalTok{ works\_oa }\SpecialCharTok{|\textgreater{}}
  \FunctionTok{filter}\NormalTok{(}\SpecialCharTok{!}\FunctionTok{is.na}\NormalTok{(oa\_status)) }\SpecialCharTok{|\textgreater{}}
  \FunctionTok{mutate}\NormalTok{(}\AttributeTok{is\_oa =}\NormalTok{ oa\_status }\SpecialCharTok{!=} \StringTok{"closed"}\NormalTok{) }\SpecialCharTok{|\textgreater{}}
  \FunctionTok{group\_by}\NormalTok{(is\_oa) }\SpecialCharTok{|\textgreater{}}
  \FunctionTok{summarise}\NormalTok{(}
    \AttributeTok{n =} \FunctionTok{n}\NormalTok{(),}
    \AttributeTok{mean\_cites =} \FunctionTok{round}\NormalTok{(}\FunctionTok{mean}\NormalTok{(cited\_by\_count), }\DecValTok{1}\NormalTok{),}
    \AttributeTok{median\_cites =} \FunctionTok{median}\NormalTok{(cited\_by\_count),}
    \AttributeTok{.groups =} \StringTok{"drop"}
\NormalTok{  )}

\NormalTok{oa\_cites }\SpecialCharTok{|\textgreater{}} \FunctionTok{gt}\NormalTok{()}
\end{Highlighting}
\end{Shaded}

\begingroup
\fontsize{8.0pt}{9.0pt}\selectfont
\begin{longtable}{>{\centering\arraybackslash}p{\dimexpr 0.23\linewidth -2\tabcolsep-1.5\arrayrulewidth}>{\raggedleft\arraybackslash}p{\dimexpr 0.23\linewidth -2\tabcolsep-1.5\arrayrulewidth}>{\raggedleft\arraybackslash}p{\dimexpr 0.26\linewidth -2\tabcolsep-1.5\arrayrulewidth}>{\raggedleft\arraybackslash}p{\dimexpr 0.28\linewidth -2\tabcolsep-1.5\arrayrulewidth}}
\toprule
is\_oa & n & \nolinkurl{mean_cites} & \nolinkurl{median_cites} \\ 
\midrule\addlinespace[2.5pt]
FALSE & 335 & 31.8 & 19 \\ 
TRUE & 265 & 36.5 & 22 \\ 
\bottomrule
\end{longtable}
\endgroup

\subsection{Document type breakdown}\label{document-type-breakdown}

\begin{Shaded}
\begin{Highlighting}[]
\NormalTok{works }\SpecialCharTok{|\textgreater{}}
  \FunctionTok{count}\NormalTok{(type, oa\_status) }\SpecialCharTok{|\textgreater{}}
  \FunctionTok{group\_by}\NormalTok{(type) }\SpecialCharTok{|\textgreater{}}
  \FunctionTok{mutate}\NormalTok{(}\AttributeTok{pct =}\NormalTok{ n }\SpecialCharTok{/} \FunctionTok{sum}\NormalTok{(n)) }\SpecialCharTok{|\textgreater{}}
  \FunctionTok{ungroup}\NormalTok{() }\SpecialCharTok{|\textgreater{}}
  \FunctionTok{ggplot}\NormalTok{(}\FunctionTok{aes}\NormalTok{(}\AttributeTok{x =}\NormalTok{ pct, }\AttributeTok{y =}\NormalTok{ type, }\AttributeTok{fill =}\NormalTok{ oa\_status)) }\SpecialCharTok{+}
  \FunctionTok{geom\_col}\NormalTok{() }\SpecialCharTok{+}
  \FunctionTok{scale\_x\_continuous}\NormalTok{(}\AttributeTok{labels =}\NormalTok{ scales}\SpecialCharTok{::}\NormalTok{percent) }\SpecialCharTok{+}
  \FunctionTok{scale\_fill\_manual}\NormalTok{(}\AttributeTok{values =} \FunctionTok{palette\_sci}\NormalTok{(}\FunctionTok{n\_distinct}\NormalTok{(works}\SpecialCharTok{$}\NormalTok{oa\_status))) }\SpecialCharTok{+}
  \FunctionTok{labs}\NormalTok{(}\AttributeTok{x =} \StringTok{"Proportion"}\NormalTok{, }\AttributeTok{y =} \ConstantTok{NULL}\NormalTok{, }\AttributeTok{fill =} \StringTok{"OA status"}\NormalTok{) }\SpecialCharTok{+}
  \FunctionTok{theme\_sci}\NormalTok{()}
\end{Highlighting}
\end{Shaded}

\begin{figure}

{\centering \includegraphics[width=0.9\linewidth,alt={Grouped bar chart showing how OA status varies across document types (article, review, etc.).}]{scientometrics-in-r_files/figure-latex/27-fig-openscience-types-1} 

}

\caption{OA status by document type.}\label{fig:27-fig-openscience-types}
\end{figure}

\subsection{Summary table}\label{summary-table-1}

\begin{Shaded}
\begin{Highlighting}[]
\NormalTok{works\_oa }\SpecialCharTok{|\textgreater{}}
  \FunctionTok{count}\NormalTok{(oa\_status, }\AttributeTok{sort =} \ConstantTok{TRUE}\NormalTok{) }\SpecialCharTok{|\textgreater{}}
  \FunctionTok{mutate}\NormalTok{(}\AttributeTok{pct =}\NormalTok{ scales}\SpecialCharTok{::}\FunctionTok{percent}\NormalTok{(n }\SpecialCharTok{/} \FunctionTok{sum}\NormalTok{(n))) }\SpecialCharTok{|\textgreater{}}
  \FunctionTok{gt}\NormalTok{() }\SpecialCharTok{|\textgreater{}}
  \FunctionTok{cols\_label}\NormalTok{(}\AttributeTok{oa\_status =} \StringTok{"OA Status"}\NormalTok{, }\AttributeTok{n =} \StringTok{"Count"}\NormalTok{, }\AttributeTok{pct =} \StringTok{"Percentage"}\NormalTok{)}
\end{Highlighting}
\end{Shaded}

\begingroup
\fontsize{8.0pt}{9.0pt}\selectfont
\begin{longtable}{>{\raggedright\arraybackslash}p{\dimexpr 0.35\linewidth -2\tabcolsep-1.5\arrayrulewidth}>{\raggedleft\arraybackslash}p{\dimexpr 0.33\linewidth -2\tabcolsep-1.5\arrayrulewidth}>{\raggedleft\arraybackslash}p{\dimexpr 0.33\linewidth -2\tabcolsep-1.5\arrayrulewidth}}
\toprule
OA Status & Count & \nolinkurl{Percentage} \\ 
\midrule\addlinespace[2.5pt]
closed & 335 & 55.8\% \\ 
hybrid & 129 & 21.5\% \\ 
green & 119 & 19.8\% \\ 
bronze & 16 & 2.7\% \\ 
diamond & 1 & 0.2\% \\ 
\bottomrule
\end{longtable}
\endgroup

\section{Diagnostics and interpretation}\label{openscience-diagnostics}

\begin{itemize}
\tightlist
\item
  \textbf{Classification accuracy}: OpenAlex OA classification relies on Unpaywall, which may lag behind publisher changes. Spot-check a sample of ``closed'' articles to verify they are truly behind a paywall.
\item
  \textbf{Bronze ambiguity}: Bronze OA means free-to-read without an explicit license. Publishers can revoke access at any time. Do not treat bronze as equivalent to gold or green.
\item
  \textbf{Year effects}: OA rates have risen over time due to funder mandates and policy changes. Compare year-normalised rates rather than pooling across years.
\item
  \textbf{Self-selection}: Higher OA rates among highly cited papers may reflect self-selection (better papers are more likely to be made OA) rather than a causal OA effect.
\end{itemize}

\section{Limitations and responsible use}\label{openscience-limitations}

\section{Limitations and responsible use}\label{limitations-and-responsible-use-26}

\begin{itemize}
\tightlist
\item
  \textbf{OA ≠ quality.} Open access status is a property of the publication venue and access model, not the research quality. Predatory journals are often gold OA.
\item
  \textbf{Coverage gaps.} Green OA via institutional repositories is underreported in some databases. Actual OA rates may be higher than measured.
\item
  \textbf{Data sharing claims vs.~reality.} A data availability statement saying ``data available upon request'' often results in no response when requests are made. Measuring DAS presence overestimates actual data sharing.
\item
  \textbf{Equity implications.} Gold OA via article processing charges shifts costs from readers to authors, potentially disadvantaging researchers without institutional funding \citep{hicks2015}.
\end{itemize}

\section{Common pitfalls}\label{openscience-pitfalls}

\section{Common pitfalls}\label{common-pitfalls-26}

\begin{itemize}
\tightlist
\item
  \textbf{Treating all OA as equivalent.} Gold, green, hybrid, and bronze OA have very different implications for access, sustainability, and cost.
\item
  \textbf{Ignoring embargo periods.} Green OA articles may have 6--24 month embargoes before they become freely available. They are ``closed'' during the embargo.
\item
  \textbf{Double-counting.} An article available in both a gold OA journal and a green repository should be counted once. OpenAlex classifies by the ``best'' OA type.
\item
  \textbf{Using OA rate as a performance indicator.} OA rates reflect policy and funding environment, not research quality. Mandated OA inflates rates without changing the underlying work.
\end{itemize}

\section{Exercises}\label{openscience-exercises}

\begin{enumerate}
\def\labelenumi{\arabic{enumi}.}
\item
  \textbf{Cross-journal OA comparison.} Compare OA rates for three journals in different disciplines. Which discipline has the highest gold OA rate?
\item
  \textbf{Temporal OA growth.} Compute the year-over-year growth rate of gold OA articles in your corpus. Has growth accelerated since 2020?
\item
  \textbf{OA and collaboration.} Is there a relationship between the number of authors on a paper and its OA status? Hypothesise why.
\end{enumerate}

\section{Solutions}\label{27-solutions}

Solutions are provided in \ref{solutions}.

\section{Further reading}\label{openscience-reading}

\begin{itemize}
\tightlist
\item
  \citet{priem2022} --- OpenAlex OA classification methodology.
\item
  \citet{hicks2015} --- The Leiden Manifesto: principles for responsible metrics, including OA considerations.
\item
  \citet{dora2012} --- DORA: recommendations for research assessment beyond journal metrics.
\end{itemize}

\section{Session info}\label{openscience-sessioninfo}

\begin{Shaded}
\begin{Highlighting}[]
\FunctionTok{sessionInfo}\NormalTok{()}
\end{Highlighting}
\end{Shaded}

\begin{verbatim}
#> R version 4.6.1 (2026-06-24)
#> Platform: x86_64-pc-linux-gnu
#> Running under: Ubuntu 24.04.5 LTS
#> 
#> Matrix products: default
#> BLAS:   /usr/lib/x86_64-linux-gnu/openblas-pthread/libblas.so.3 
#> LAPACK: /usr/lib/x86_64-linux-gnu/openblas-pthread/libopenblasp-r0.3.26.so;  LAPACK version 3.12.0
#> 
#> locale:
#>  [1] LC_CTYPE=C.UTF-8       LC_NUMERIC=C           LC_TIME=C.UTF-8       
#>  [4] LC_COLLATE=C.UTF-8     LC_MONETARY=C.UTF-8    LC_MESSAGES=C.UTF-8   
#>  [7] LC_PAPER=C.UTF-8       LC_NAME=C              LC_ADDRESS=C          
#> [10] LC_TELEPHONE=C         LC_MEASUREMENT=C.UTF-8 LC_IDENTIFICATION=C   
#> 
#> time zone: UTC
#> tzcode source: system (glibc)
#> 
#> attached base packages:
#> [1] stats     graphics  grDevices utils     datasets  methods   base     
#> 
#> other attached packages:
#>  [1] uwot_0.2.5                Matrix_1.7-5             
#>  [3] word2vec_0.4.1            stm_1.3.8                
#>  [5] topicmodels_0.2-17        quanteda.textstats_0.97.2
#>  [7] visNetwork_2.1.4          ggraph_2.2.2             
#>  [9] tidygraph_1.3.1           igraph_2.3.4             
#> [11] quanteda_4.5.1            pdftools_3.9.1           
#> [13] arrow_26.0.0              bibliometrix_5.5.0       
#> [15] RefManageR_1.4.0          bib2df_1.1.2.0           
#> [17] rcrossref_1.2.1           gt_1.4.0                 
#> [19] tidytext_0.4.3            glue_1.8.1               
#> [21] openalexR_3.1.2           lubridate_1.9.5          
#> [23] forcats_1.0.1             stringr_1.6.0            
#> [25] dplyr_1.2.1               purrr_1.2.2              
#> [27] readr_2.2.0               tidyr_1.3.2              
#> [29] tibble_3.3.1              ggplot2_4.0.3            
#> [31] tidyverse_2.0.0          
#> 
#> loaded via a namespace (and not attached):
#>   [1] bibtex_0.5.3           RColorBrewer_1.1-3     rstudioapi_0.19.0     
#>   [4] jsonlite_2.0.0         magrittr_2.0.5         modeltools_0.2-25     
#>   [7] farver_2.1.2           rmarkdown_2.32         fs_2.1.0              
#>  [10] vctrs_0.7.3            memoise_2.0.1          askpass_1.2.1         
#>  [13] base64enc_0.1-6        htmltools_0.5.9        contentanalysis_1.1.1 
#>  [16] curl_8.1.0             janeaustenr_1.0.0      cellranger_1.1.0      
#>  [19] htmlwidgets_1.6.4      tokenizers_0.3.0       plyr_1.8.9            
#>  [22] httr2_1.3.0            cachem_1.1.0           plotly_4.12.1         
#>  [25] dimensionsR_0.0.3      mime_0.13              lifecycle_1.0.5       
#>  [28] pkgconfig_2.0.3        R6_2.6.1               fastmap_1.2.0         
#>  [31] shiny_1.14.0           digest_0.6.39          patchwork_1.3.2       
#>  [34] shinycssloaders_1.1.0  rprojroot_2.1.1        RSpectra_0.16-2       
#>  [37] SnowballC_0.7.1        labeling_0.4.3         urltools_1.7.3.1      
#>  [40] timechange_0.4.0       mgcv_1.9-4             httr_1.4.9            
#>  [43] polyclip_1.10-7        compiler_4.6.1         here_1.0.2            
#>  [46] bit64_4.8.6            withr_3.0.3            S7_0.2.2              
#>  [49] viridis_0.6.5          ggforce_0.5.0          MASS_7.3-65           
#>  [52] bibliometrixData_0.3.0 tools_4.6.1            otel_0.2.0            
#>  [55] stopwords_2.3          zip_3.0.2              rentrez_1.2.4         
#>  [58] httpuv_1.6.17          nlme_3.1-169           promises_1.5.0        
#>  [61] grid_4.6.1             stringdist_0.9.17      reshape2_1.4.5        
#>  [64] generics_0.1.4         gtable_0.3.6           tzdb_0.5.0            
#>  [67] rscopus_0.9.0          ca_0.72                data.table_1.18.6.1   
#>  [70] hms_1.1.4              xml2_1.6.0             utf8_1.2.6            
#>  [73] ggrepel_0.9.8          pillar_1.11.1          nsyllable_1.0.1       
#>  [76] vroom_1.7.1            later_1.4.8            splines_4.6.1         
#>  [79] tweenr_2.0.3           lattice_0.22-9         FNN_1.1.4.1           
#>  [82] bit_4.6.0              tidyselect_1.2.1       tm_0.7-20             
#>  [85] miniUI_0.1.2           knitr_1.52             gridExtra_2.3.1       
#>  [88] NLP_0.3-3              bookdown_0.48          stats4_4.6.1          
#>  [91] crul_1.6.0             xfun_0.62              graphlayouts_1.2.5    
#>  [94] matrixStats_1.5.0      DT_0.34.0              humaniformat_0.6.0    
#>  [97] stringi_1.8.9          qpdf_1.4.1             yaml_2.3.12           
#> [100] evaluate_1.0.5         httpcode_0.3.0         cli_3.6.6             
#> [103] xtable_1.8-8           systemfonts_1.3.2      Rcpp_1.1.2            
#> [106] readxl_1.6.0           triebeard_0.4.1        XML_3.99-0.25         
#> [109] parallel_4.6.1         assertthat_0.2.1       pubmedR_1.0.2         
#> [112] slam_0.1-56            viridisLite_0.4.3      scales_1.4.0          
#> [115] crayon_1.5.3           openxlsx_4.2.9         rlang_1.3.0           
#> [118] fastmatch_1.1-8
\end{verbatim}

\chapter{Gender and Equity Analysis}\label{gender-and-equity-analysis}

\emph{R. Heller · Date: 2026-05-09 · Last changed: 2026-10-31}

\section{Learning objectives}\label{gender-objectives}

After completing this chapter, you will be able to:

\begin{itemize}
\tightlist
\item
  Distinguish a name-lookup category from a person's gender identity
\item
  Measure lookup coverage and explain how excluding unmatched records changes a denominator
\item
  Summarise explicitly synthetic authorship records by recorded position
\item
  Identify the data, validation and consent questions required for an actual equity study
\end{itemize}

\section{Setup}\label{gender-setup}

\begin{Shaded}
\begin{Highlighting}[]
\FunctionTok{library}\NormalTok{(tidyverse)}
\FunctionTok{library}\NormalTok{(glue)}
\FunctionTok{library}\NormalTok{(gt)}

\FunctionTok{source}\NormalTok{(here}\SpecialCharTok{::}\FunctionTok{here}\NormalTok{(}\StringTok{"R"}\NormalTok{, }\StringTok{"sci\_palette.R"}\NormalTok{))}
\end{Highlighting}
\end{Shaded}

\section{Conceptual background}\label{gender-background}

Research such as \citet{lariviere2013} examines gender disparities in scientific authorship. Its findings concern particular data, operational definitions and study periods; they are not measurements supplied by this chapter.

Names are not observations of gender identity. A name-classification service produces a model or lookup output, and its coverage and errors depend on its reference data and target population. \href{https://doi.org/10.7717/peerj-cs.156}{Santamaría and Mihaljević (2018)} compare such services and examine their differing performance. A small hand-written dictionary has no demonstrated accuracy for a scholarly population. Increasing its coverage would not establish validity.

Authorship order is also context-dependent. A recorded first or last position does not establish contribution or seniority, and alphabetical conventions differ from contribution-based conventions. Gender-equity analysis therefore requires an explicit population, defensible demographic measurement, appropriate handling of unavailable information, and interpretation of the authorship conventions actually used. Responsible research assessment should not reduce individuals to such categories \citep{hicks2015, dora2012}.

\section{Worked example}\label{gender-example}

\subsection{Explicitly synthetic records}\label{explicitly-synthetic-records}

The following records are invented for a data-handling exercise. They do not describe real authors, papers, journals or gender identities. The lookup categories \textbf{A} and \textbf{B} are arbitrary labels with no demographic meaning. No OpenAlex records or external classification service are used in this example.

\begin{Shaded}
\begin{Highlighting}[]
\NormalTok{synthetic\_authors }\OtherTok{\textless{}{-}} \FunctionTok{tibble}\NormalTok{(}
  \AttributeTok{work\_id =} \FunctionTok{c}\NormalTok{(}\FunctionTok{rep}\NormalTok{(}\StringTok{"example{-}1"}\NormalTok{, }\DecValTok{3}\NormalTok{), }\FunctionTok{rep}\NormalTok{(}\StringTok{"example{-}2"}\NormalTok{, }\DecValTok{2}\NormalTok{), }\StringTok{"example{-}3"}\NormalTok{,}
              \FunctionTok{rep}\NormalTok{(}\StringTok{"example{-}4"}\NormalTok{, }\DecValTok{4}\NormalTok{), }\FunctionTok{rep}\NormalTok{(}\StringTok{"example{-}5"}\NormalTok{, }\DecValTok{2}\NormalTok{)),}
  \AttributeTok{author\_order =} \FunctionTok{c}\NormalTok{(}\DecValTok{1}\SpecialCharTok{:}\DecValTok{3}\NormalTok{, }\DecValTok{1}\SpecialCharTok{:}\DecValTok{2}\NormalTok{, }\DecValTok{1}\NormalTok{, }\DecValTok{1}\SpecialCharTok{:}\DecValTok{4}\NormalTok{, }\DecValTok{1}\SpecialCharTok{:}\DecValTok{2}\NormalTok{),}
  \AttributeTok{first\_name =} \FunctionTok{c}\NormalTok{(}\StringTok{"aria"}\NormalTok{, }\StringTok{"jordan"}\NormalTok{, }\StringTok{"lee"}\NormalTok{, }\StringTok{"noor"}\NormalTok{, }\StringTok{"ari"}\NormalTok{, }\StringTok{"morgan"}\NormalTok{,}
                 \StringTok{"aria"}\NormalTok{, }\StringTok{"nuri"}\NormalTok{, }\StringTok{"sam"}\NormalTok{, }\StringTok{"alex"}\NormalTok{, }\StringTok{"ari"}\NormalTok{, }\StringTok{"noor"}\NormalTok{)}
\NormalTok{)}

\FunctionTok{cat}\NormalTok{(}\FunctionTok{glue}\NormalTok{(}\StringTok{"Synthetic author records: \{nrow(synthetic\_authors)\}}\SpecialCharTok{\textbackslash{}n}\StringTok{"}\NormalTok{))}
\end{Highlighting}
\end{Shaded}

\begin{verbatim}
#> Synthetic author records: 12
\end{verbatim}

\begin{Shaded}
\begin{Highlighting}[]
\FunctionTok{cat}\NormalTok{(}\FunctionTok{glue}\NormalTok{(}\StringTok{"Distinct example names: \{n\_distinct(synthetic\_authors$first\_name)\}}\SpecialCharTok{\textbackslash{}n}\StringTok{"}\NormalTok{))}
\end{Highlighting}
\end{Shaded}

\begin{verbatim}
#> Distinct example names: 9
\end{verbatim}

\subsection{Lookup coverage, not gender classification}\label{lookup-coverage-not-gender-classification}

\begin{Shaded}
\begin{Highlighting}[]
\NormalTok{example\_lookup }\OtherTok{\textless{}{-}} \FunctionTok{tibble}\NormalTok{(}
  \AttributeTok{first\_name =} \FunctionTok{c}\NormalTok{(}\StringTok{"aria"}\NormalTok{, }\StringTok{"ari"}\NormalTok{, }\StringTok{"jordan"}\NormalTok{, }\StringTok{"morgan"}\NormalTok{),}
  \AttributeTok{lookup\_category =} \FunctionTok{c}\NormalTok{(}\StringTok{"A"}\NormalTok{, }\StringTok{"A"}\NormalTok{, }\StringTok{"B"}\NormalTok{, }\StringTok{"B"}\NormalTok{)}
\NormalTok{)}

\NormalTok{lookup\_records }\OtherTok{\textless{}{-}}\NormalTok{ synthetic\_authors }\SpecialCharTok{|\textgreater{}}
  \FunctionTok{left\_join}\NormalTok{(example\_lookup, }\AttributeTok{by =} \StringTok{"first\_name"}\NormalTok{, }\AttributeTok{relationship =} \StringTok{"many{-}to{-}one"}\NormalTok{) }\SpecialCharTok{|\textgreater{}}
  \FunctionTok{mutate}\NormalTok{(}\AttributeTok{lookup\_category =} \FunctionTok{replace\_na}\NormalTok{(lookup\_category, }\StringTok{"Unmatched"}\NormalTok{))}

\NormalTok{coverage }\OtherTok{\textless{}{-}} \FunctionTok{mean}\NormalTok{(lookup\_records}\SpecialCharTok{$}\NormalTok{lookup\_category }\SpecialCharTok{!=} \StringTok{"Unmatched"}\NormalTok{)}
\FunctionTok{cat}\NormalTok{(}\FunctionTok{glue}\NormalTok{(}\StringTok{"Lookup coverage in this invented fixture: \{scales::percent(coverage)\}}\SpecialCharTok{\textbackslash{}n}\StringTok{"}\NormalTok{))}
\end{Highlighting}
\end{Shaded}

\begin{verbatim}
#> Lookup coverage in this invented fixture: 50%
\end{verbatim}

\subsection{Keeping the full denominator}\label{keeping-the-full-denominator}

\begin{Shaded}
\begin{Highlighting}[]
\NormalTok{lookup\_summary }\OtherTok{\textless{}{-}}\NormalTok{ lookup\_records }\SpecialCharTok{|\textgreater{}}
  \FunctionTok{count}\NormalTok{(lookup\_category) }\SpecialCharTok{|\textgreater{}}
  \FunctionTok{mutate}\NormalTok{(}\AttributeTok{pct =}\NormalTok{ n }\SpecialCharTok{/} \FunctionTok{sum}\NormalTok{(n))}

\FunctionTok{ggplot}\NormalTok{(lookup\_summary, }\FunctionTok{aes}\NormalTok{(}\AttributeTok{x =}\NormalTok{ lookup\_category, }\AttributeTok{y =}\NormalTok{ pct, }\AttributeTok{fill =}\NormalTok{ lookup\_category)) }\SpecialCharTok{+}
  \FunctionTok{geom\_col}\NormalTok{(}\AttributeTok{show.legend =} \ConstantTok{FALSE}\NormalTok{) }\SpecialCharTok{+}
  \FunctionTok{scale\_y\_continuous}\NormalTok{(}\AttributeTok{labels =}\NormalTok{ scales}\SpecialCharTok{::}\NormalTok{percent) }\SpecialCharTok{+}
  \FunctionTok{scale\_fill\_manual}\NormalTok{(}\AttributeTok{values =} \FunctionTok{palette\_sci}\NormalTok{(}\DecValTok{3}\NormalTok{)) }\SpecialCharTok{+}
  \FunctionTok{labs}\NormalTok{(}\AttributeTok{x =} \StringTok{"Arbitrary lookup category"}\NormalTok{, }\AttributeTok{y =} \StringTok{"Share of all synthetic records"}\NormalTok{) }\SpecialCharTok{+}
  \FunctionTok{theme\_sci}\NormalTok{()}
\end{Highlighting}
\end{Shaded}

\begin{figure}

{\centering \includegraphics[width=0.9\linewidth,alt={Bar chart with arbitrary categories A, B and Unmatched for a synthetic name-lookup exercise.}]{scientometrics-in-r_files/figure-latex/28-fig-gender-proportion-1} 

}

\caption{Arbitrary lookup categories in invented records, retaining unmatched names in the denominator. These are not gender proportions.}\label{fig:28-fig-gender-proportion}
\end{figure}

\subsection{Recorded authorship positions}\label{recorded-authorship-positions}

The invented order is recorded explicitly. Single-author records are kept separate from first, middle and last positions in multi-author records. These positions are not inferred contribution roles.

\begin{Shaded}
\begin{Highlighting}[]
\NormalTok{position\_data }\OtherTok{\textless{}{-}}\NormalTok{ lookup\_records }\SpecialCharTok{|\textgreater{}}
  \FunctionTok{group\_by}\NormalTok{(work\_id) }\SpecialCharTok{|\textgreater{}}
  \FunctionTok{arrange}\NormalTok{(author\_order, }\AttributeTok{.by\_group =} \ConstantTok{TRUE}\NormalTok{) }\SpecialCharTok{|\textgreater{}}
  \FunctionTok{mutate}\NormalTok{(}
    \AttributeTok{n\_authors =} \FunctionTok{n}\NormalTok{(),}
    \AttributeTok{author\_position =} \FunctionTok{case\_when}\NormalTok{(}
\NormalTok{      n\_authors }\SpecialCharTok{==} \DecValTok{1} \SpecialCharTok{\textasciitilde{}} \StringTok{"single"}\NormalTok{,}
      \FunctionTok{row\_number}\NormalTok{() }\SpecialCharTok{==} \DecValTok{1} \SpecialCharTok{\textasciitilde{}} \StringTok{"first"}\NormalTok{,}
      \FunctionTok{row\_number}\NormalTok{() }\SpecialCharTok{==}\NormalTok{ n\_authors }\SpecialCharTok{\textasciitilde{}} \StringTok{"last"}\NormalTok{,}
      \ConstantTok{TRUE} \SpecialCharTok{\textasciitilde{}} \StringTok{"middle"}
\NormalTok{    )}
\NormalTok{  ) }\SpecialCharTok{|\textgreater{}}
  \FunctionTok{ungroup}\NormalTok{()}

\NormalTok{position\_summary }\OtherTok{\textless{}{-}}\NormalTok{ position\_data }\SpecialCharTok{|\textgreater{}}
  \FunctionTok{count}\NormalTok{(author\_position, lookup\_category) }\SpecialCharTok{|\textgreater{}}
  \FunctionTok{group\_by}\NormalTok{(author\_position) }\SpecialCharTok{|\textgreater{}}
  \FunctionTok{mutate}\NormalTok{(}\AttributeTok{pct =}\NormalTok{ n }\SpecialCharTok{/} \FunctionTok{sum}\NormalTok{(n)) }\SpecialCharTok{|\textgreater{}}
  \FunctionTok{ungroup}\NormalTok{()}

\NormalTok{position\_summary }\SpecialCharTok{|\textgreater{}} \FunctionTok{gt}\NormalTok{()}
\end{Highlighting}
\end{Shaded}

\begingroup
\fontsize{8.0pt}{9.0pt}\selectfont
\begin{longtable}{>{\raggedright\arraybackslash}p{\dimexpr 0.29\linewidth -2\tabcolsep-1.5\arrayrulewidth}>{\raggedright\arraybackslash}p{\dimexpr 0.29\linewidth -2\tabcolsep-1.5\arrayrulewidth}>{\raggedleft\arraybackslash}p{\dimexpr 0.21\linewidth -2\tabcolsep-1.5\arrayrulewidth}>{\raggedleft\arraybackslash}p{\dimexpr 0.21\linewidth -2\tabcolsep-1.5\arrayrulewidth}}
\toprule
\nolinkurl{author_position} & \nolinkurl{lookup_category} & n & pct \\ 
\midrule\addlinespace[2.5pt]
first & A & 3 & 0.750 \\ 
first & \nolinkurl{Unmatched} & 1 & 0.250 \\ 
last & A & 1 & 0.250 \\ 
last & \nolinkurl{Unmatched} & 3 & 0.750 \\ 
middle & B & 1 & 0.333 \\ 
middle & \nolinkurl{Unmatched} & 2 & 0.667 \\ 
single & B & 1 & 1.000 \\ 
\bottomrule
\end{longtable}
\endgroup

\begin{Shaded}
\begin{Highlighting}[]
\FunctionTok{ggplot}\NormalTok{(position\_summary, }\FunctionTok{aes}\NormalTok{(}\AttributeTok{x =}\NormalTok{ author\_position, }\AttributeTok{y =}\NormalTok{ pct, }\AttributeTok{fill =}\NormalTok{ lookup\_category)) }\SpecialCharTok{+}
  \FunctionTok{geom\_col}\NormalTok{(}\AttributeTok{position =} \StringTok{"dodge"}\NormalTok{) }\SpecialCharTok{+}
  \FunctionTok{scale\_y\_continuous}\NormalTok{(}\AttributeTok{labels =}\NormalTok{ scales}\SpecialCharTok{::}\NormalTok{percent) }\SpecialCharTok{+}
  \FunctionTok{scale\_fill\_manual}\NormalTok{(}\AttributeTok{values =} \FunctionTok{palette\_sci}\NormalTok{(}\DecValTok{3}\NormalTok{)) }\SpecialCharTok{+}
  \FunctionTok{labs}\NormalTok{(}\AttributeTok{x =} \StringTok{"Recorded position"}\NormalTok{, }\AttributeTok{y =} \StringTok{"Share of synthetic records"}\NormalTok{, }\AttributeTok{fill =} \StringTok{"Lookup category"}\NormalTok{) }\SpecialCharTok{+}
  \FunctionTok{theme\_sci}\NormalTok{()}
\end{Highlighting}
\end{Shaded}

\begin{figure}

{\centering \includegraphics[width=0.9\linewidth,alt={Grouped bar chart of categories A, B and Unmatched in synthetic first, middle, last and single-author positions.}]{scientometrics-in-r_files/figure-latex/28-fig-gender-position-1} 

}

\caption{Arbitrary lookup categories by recorded position in invented authorship lists. Unmatched records remain in each position's denominator.}\label{fig:28-fig-gender-position}
\end{figure}

\section{Diagnostics and interpretation}\label{gender-diagnostics}

\begin{itemize}
\tightlist
\item
  \textbf{Coverage and validity are different.} Coverage is the proportion assigned a dictionary label. This exercise has no observed gender variable against which accuracy could be estimated.
\item
  \textbf{Denominators matter.} Report unmatched records rather than silently excluding them. In this fixture, half the records are unmatched; the plotted A/B proportions would change if they were discarded.
\item
  \textbf{Units matter.} These are authorship records, not distinct people. Repeated invented names do not establish repeated identity.
\item
  \textbf{Position is descriptive.} Check ordering, single-author records and local authorship conventions before interpreting positions.
\end{itemize}

\section{Limitations and responsible use}\label{gender-limitations}

\section{Limitations and responsible use}\label{limitations-and-responsible-use-27}

This synthetic example establishes no gender distribution, citation gap or equity finding. A real study needs a defensible demographic variable, clearly stated categories and missingness, appropriate consent and data governance, and validation for the population under study. Self-described gender data and the research question may require categories beyond a binary model; a name dictionary cannot supply those observations. Avoid publishing individual demographic assignments inferred from names, and distinguish observed disparity from an explanation of its causes.

\section{Common pitfalls}\label{gender-pitfalls}

\section{Common pitfalls}\label{common-pitfalls-27}

\begin{itemize}
\tightlist
\item
  Treating lookup output as measured gender identity or biological sex.
\item
  Calling high classification coverage evidence of accuracy.
\item
  Removing unmatched names and presenting the remainder as representative.
\item
  Counting authorship records as unique researchers without identity linkage.
\item
  Inferring contribution or seniority from author order without supporting information.
\end{itemize}

\section{Exercises}\label{gender-exercises}

\begin{enumerate}
\def\labelenumi{\arabic{enumi}.}
\item
  \textbf{Denominator comparison.} Recompute the A/B proportions after excluding unmatched records. Explain why those proportions differ from the chart without establishing demographic representation.
\item
  \textbf{Position audit.} Check that the fixture has four first, four last, three middle and one single-author record. Verify that each position's proportions sum to one and still include unmatched records.
\item
  \textbf{Study design.} Specify the population, demographic measurement, consent/missingness plan and authorship conventions needed for a real equity study. Explain which required variables are absent from this fixture.
\end{enumerate}

\section{Solutions}\label{28-solutions}

Solutions are provided in \ref{solutions}.

\section{Further reading}\label{gender-reading}

\begin{itemize}
\tightlist
\item
  \citet{lariviere2013} --- An empirical study of gender disparities in scientific authorship, with its own population and operational definitions.
\item
  \href{https://doi.org/10.7717/peerj-cs.156}{Santamaría and Mihaljević (2018)} --- Comparison and benchmark of name-to-gender inference services.
\item
  \citet{hicks2015} and \citet{dora2012} --- Responsible interpretation of research metrics.
\end{itemize}

\section{Session info}\label{gender-sessioninfo}

\begin{Shaded}
\begin{Highlighting}[]
\FunctionTok{sessionInfo}\NormalTok{()}
\end{Highlighting}
\end{Shaded}

\begin{verbatim}
#> R version 4.6.1 (2026-06-24)
#> Platform: x86_64-pc-linux-gnu
#> Running under: Ubuntu 24.04.5 LTS
#> 
#> Matrix products: default
#> BLAS:   /usr/lib/x86_64-linux-gnu/openblas-pthread/libblas.so.3 
#> LAPACK: /usr/lib/x86_64-linux-gnu/openblas-pthread/libopenblasp-r0.3.26.so;  LAPACK version 3.12.0
#> 
#> locale:
#>  [1] LC_CTYPE=C.UTF-8       LC_NUMERIC=C           LC_TIME=C.UTF-8       
#>  [4] LC_COLLATE=C.UTF-8     LC_MONETARY=C.UTF-8    LC_MESSAGES=C.UTF-8   
#>  [7] LC_PAPER=C.UTF-8       LC_NAME=C              LC_ADDRESS=C          
#> [10] LC_TELEPHONE=C         LC_MEASUREMENT=C.UTF-8 LC_IDENTIFICATION=C   
#> 
#> time zone: UTC
#> tzcode source: system (glibc)
#> 
#> attached base packages:
#> [1] stats     graphics  grDevices utils     datasets  methods   base     
#> 
#> other attached packages:
#>  [1] uwot_0.2.5                Matrix_1.7-5             
#>  [3] word2vec_0.4.1            stm_1.3.8                
#>  [5] topicmodels_0.2-17        quanteda.textstats_0.97.2
#>  [7] visNetwork_2.1.4          ggraph_2.2.2             
#>  [9] tidygraph_1.3.1           igraph_2.3.4             
#> [11] quanteda_4.5.1            pdftools_3.9.1           
#> [13] arrow_26.0.0              bibliometrix_5.5.0       
#> [15] RefManageR_1.4.0          bib2df_1.1.2.0           
#> [17] rcrossref_1.2.1           gt_1.4.0                 
#> [19] tidytext_0.4.3            glue_1.8.1               
#> [21] openalexR_3.1.2           lubridate_1.9.5          
#> [23] forcats_1.0.1             stringr_1.6.0            
#> [25] dplyr_1.2.1               purrr_1.2.2              
#> [27] readr_2.2.0               tidyr_1.3.2              
#> [29] tibble_3.3.1              ggplot2_4.0.3            
#> [31] tidyverse_2.0.0          
#> 
#> loaded via a namespace (and not attached):
#>   [1] bibtex_0.5.3           RColorBrewer_1.1-3     rstudioapi_0.19.0     
#>   [4] jsonlite_2.0.0         magrittr_2.0.5         modeltools_0.2-25     
#>   [7] farver_2.1.2           rmarkdown_2.32         fs_2.1.0              
#>  [10] vctrs_0.7.3            memoise_2.0.1          askpass_1.2.1         
#>  [13] base64enc_0.1-6        htmltools_0.5.9        contentanalysis_1.1.1 
#>  [16] curl_8.1.0             janeaustenr_1.0.0      cellranger_1.1.0      
#>  [19] htmlwidgets_1.6.4      tokenizers_0.3.0       plyr_1.8.9            
#>  [22] httr2_1.3.0            cachem_1.1.0           plotly_4.12.1         
#>  [25] dimensionsR_0.0.3      mime_0.13              lifecycle_1.0.5       
#>  [28] pkgconfig_2.0.3        R6_2.6.1               fastmap_1.2.0         
#>  [31] shiny_1.14.0           digest_0.6.39          patchwork_1.3.2       
#>  [34] shinycssloaders_1.1.0  rprojroot_2.1.1        RSpectra_0.16-2       
#>  [37] SnowballC_0.7.1        labeling_0.4.3         urltools_1.7.3.1      
#>  [40] timechange_0.4.0       mgcv_1.9-4             httr_1.4.9            
#>  [43] polyclip_1.10-7        compiler_4.6.1         here_1.0.2            
#>  [46] bit64_4.8.6            withr_3.0.3            S7_0.2.2              
#>  [49] viridis_0.6.5          ggforce_0.5.0          MASS_7.3-65           
#>  [52] bibliometrixData_0.3.0 tools_4.6.1            otel_0.2.0            
#>  [55] stopwords_2.3          zip_3.0.2              rentrez_1.2.4         
#>  [58] httpuv_1.6.17          nlme_3.1-169           promises_1.5.0        
#>  [61] grid_4.6.1             stringdist_0.9.17      reshape2_1.4.5        
#>  [64] generics_0.1.4         gtable_0.3.6           tzdb_0.5.0            
#>  [67] rscopus_0.9.0          ca_0.72                data.table_1.18.6.1   
#>  [70] hms_1.1.4              xml2_1.6.0             utf8_1.2.6            
#>  [73] ggrepel_0.9.8          pillar_1.11.1          nsyllable_1.0.1       
#>  [76] vroom_1.7.1            later_1.4.8            splines_4.6.1         
#>  [79] tweenr_2.0.3           lattice_0.22-9         FNN_1.1.4.1           
#>  [82] bit_4.6.0              tidyselect_1.2.1       tm_0.7-20             
#>  [85] miniUI_0.1.2           knitr_1.52             gridExtra_2.3.1       
#>  [88] NLP_0.3-3              bookdown_0.48          stats4_4.6.1          
#>  [91] crul_1.6.0             xfun_0.62              graphlayouts_1.2.5    
#>  [94] matrixStats_1.5.0      DT_0.34.0              humaniformat_0.6.0    
#>  [97] stringi_1.8.9          qpdf_1.4.1             yaml_2.3.12           
#> [100] evaluate_1.0.5         httpcode_0.3.0         cli_3.6.6             
#> [103] xtable_1.8-8           systemfonts_1.3.2      Rcpp_1.1.2            
#> [106] readxl_1.6.0           triebeard_0.4.1        XML_3.99-0.25         
#> [109] parallel_4.6.1         assertthat_0.2.1       pubmedR_1.0.2         
#> [112] slam_0.1-56            viridisLite_0.4.3      scales_1.4.0          
#> [115] crayon_1.5.3           openxlsx_4.2.9         rlang_1.3.0           
#> [118] fastmatch_1.1-8
\end{verbatim}

\chapter{Retraction Analysis}\label{retraction-analysis}

\emph{R. Heller · Date: 2026-05-09 · Last changed: 2026-10-31}

\section{Learning objectives}\label{retraction-objectives}

After completing this chapter, you will be able to:

\begin{itemize}
\tightlist
\item
  Identify retracted works in OpenAlex using type and metadata filters
\item
  Analyse retraction rates by year, journal, and field
\item
  Examine citation patterns before and after retraction
\item
  Discuss the limitations of bibliometric approaches to research integrity
\item
  Explain why retracted papers continue to accumulate citations
\end{itemize}

\section{Setup}\label{retraction-setup}

\begin{Shaded}
\begin{Highlighting}[]
\FunctionTok{library}\NormalTok{(tidyverse)}
\FunctionTok{library}\NormalTok{(openalexR)}
\FunctionTok{library}\NormalTok{(glue)}
\FunctionTok{library}\NormalTok{(gt)}

\FunctionTok{set.seed}\NormalTok{(}\DecValTok{20260509}\NormalTok{)}

\FunctionTok{source}\NormalTok{(here}\SpecialCharTok{::}\FunctionTok{here}\NormalTok{(}\StringTok{"R"}\NormalTok{, }\StringTok{"api\_helpers.R"}\NormalTok{))}
\FunctionTok{source}\NormalTok{(here}\SpecialCharTok{::}\FunctionTok{here}\NormalTok{(}\StringTok{"R"}\NormalTok{, }\StringTok{"utils.R"}\NormalTok{))}
\FunctionTok{source}\NormalTok{(here}\SpecialCharTok{::}\FunctionTok{here}\NormalTok{(}\StringTok{"R"}\NormalTok{, }\StringTok{"sci\_palette.R"}\NormalTok{))}
\end{Highlighting}
\end{Shaded}

\section{Conceptual background}\label{retraction-background}

Retraction is the formal mechanism by which the scholarly record corrects itself. A paper is retracted when it is found to contain fabricated data, serious errors, plagiarism, or ethical violations. In theory, retraction should halt the influence of flawed work. In practice, retracted papers continue to be cited --- sometimes for years or decades after retraction.

The growth in retractions has been dramatic: from a few dozen per year in the 1990s to over 10,000 in recent years. This increase reflects both more misconduct detection (plagiarism software, image forensics) and more publications overall. Retraction rates vary by field: biomedical sciences have the highest absolute numbers, partly because of larger publication volumes and greater scrutiny.

\textbf{Citation persistence} is the most studied bibliometric aspect of retractions. Studies consistently find that retracted papers continue to accumulate positive citations (i.e., citations that treat the retracted work as valid) long after retraction. This occurs because: (1) citing authors may not be aware of the retraction; (2) database retraction notices are inconsistently displayed; (3) citations in already-published papers are not retroactively removed.

OpenAlex marks retracted works with a specific type or retraction flag, though coverage depends on publisher reporting. The Retraction Watch database provides the most comprehensive curated list of retractions but is not directly integrated into most bibliometric databases.

Bibliometric analysis of retractions must be approached carefully. The goal is to improve the self-correction mechanism, not to stigmatise individuals or institutions.

\section{Worked example}\label{retraction-example}

\subsection{Finding retracted works}\label{finding-retracted-works}

\begin{Shaded}
\begin{Highlighting}[]
\NormalTok{retracted }\OtherTok{\textless{}{-}} \FunctionTok{fetch\_openalex}\NormalTok{(}
  \AttributeTok{entity =} \StringTok{"works"}\NormalTok{,}
  \AttributeTok{is\_retracted =} \StringTok{"true"}\NormalTok{,}
  \AttributeTok{from\_publication\_date =} \StringTok{"2010{-}01{-}01"}\NormalTok{,}
  \AttributeTok{to\_publication\_date =} \StringTok{"2023{-}12{-}31"}\NormalTok{,}
  \AttributeTok{options =} \FunctionTok{list}\NormalTok{(}\AttributeTok{sample =} \DecValTok{200}\NormalTok{, }\AttributeTok{seed =} \DecValTok{42}\NormalTok{)}
\NormalTok{)}

\FunctionTok{cat}\NormalTok{(}\FunctionTok{glue}\NormalTok{(}\StringTok{"Retracted works retrieved: \{nrow(retracted)\}}\SpecialCharTok{\textbackslash{}n}\StringTok{"}\NormalTok{))}
\end{Highlighting}
\end{Shaded}

\begin{verbatim}
#> Retracted works retrieved: 200
\end{verbatim}

\begin{Shaded}
\begin{Highlighting}[]
\NormalTok{retraction\_years }\OtherTok{\textless{}{-}}\NormalTok{ retracted }\SpecialCharTok{|\textgreater{}}
  \FunctionTok{mutate}\NormalTok{(}\AttributeTok{year =} \FunctionTok{year}\NormalTok{(publication\_date)) }\SpecialCharTok{|\textgreater{}}
  \FunctionTok{count}\NormalTok{(year, }\AttributeTok{sort =} \ConstantTok{FALSE}\NormalTok{)}

\NormalTok{retraction\_years }\SpecialCharTok{|\textgreater{}} \FunctionTok{gt}\NormalTok{()}
\end{Highlighting}
\end{Shaded}

\begingroup
\fontsize{8.0pt}{9.0pt}\selectfont
\begin{longtable}{>{\raggedleft\arraybackslash}p{\dimexpr 0.50\linewidth -2\tabcolsep-1.5\arrayrulewidth}>{\raggedleft\arraybackslash}p{\dimexpr 0.50\linewidth -2\tabcolsep-1.5\arrayrulewidth}}
\toprule
year & n \\ 
\midrule\addlinespace[2.5pt]
2010 & 9 \\ 
2011 & 14 \\ 
2012 & 5 \\ 
2013 & 2 \\ 
2014 & 7 \\ 
2015 & 4 \\ 
2016 & 6 \\ 
2017 & 5 \\ 
2018 & 15 \\ 
2019 & 9 \\ 
2020 & 23 \\ 
2021 & 26 \\ 
2022 & 31 \\ 
2023 & 44 \\ 
\bottomrule
\end{longtable}
\endgroup

\subsection{Citation counts of retracted papers}\label{citation-counts-of-retracted-papers}

\begin{Shaded}
\begin{Highlighting}[]
\FunctionTok{ggplot}\NormalTok{(retracted, }\FunctionTok{aes}\NormalTok{(}\AttributeTok{x =}\NormalTok{ cited\_by\_count)) }\SpecialCharTok{+}
  \FunctionTok{geom\_histogram}\NormalTok{(}\AttributeTok{binwidth =} \DecValTok{5}\NormalTok{, }\AttributeTok{fill =} \FunctionTok{palette\_sci}\NormalTok{(}\DecValTok{1}\NormalTok{), }\AttributeTok{colour =} \StringTok{"white"}\NormalTok{) }\SpecialCharTok{+}
  \FunctionTok{labs}\NormalTok{(}\AttributeTok{x =} \StringTok{"Citations received"}\NormalTok{, }\AttributeTok{y =} \StringTok{"Number of retracted papers"}\NormalTok{) }\SpecialCharTok{+}
  \FunctionTok{theme\_sci}\NormalTok{()}
\end{Highlighting}
\end{Shaded}

\begin{figure}

{\centering \includegraphics[width=0.9\linewidth,alt={Histogram showing how many citations retracted papers have accumulated, revealing that many retracted papers continue to be cited.}]{scientometrics-in-r_files/figure-latex/29-fig-retraction-citations-1} 

}

\caption{Citation count distribution of retracted papers.}\label{fig:29-fig-retraction-citations}
\end{figure}

\begin{Shaded}
\begin{Highlighting}[]
\NormalTok{cite\_stats }\OtherTok{\textless{}{-}}\NormalTok{ retracted }\SpecialCharTok{|\textgreater{}}
  \FunctionTok{summarise}\NormalTok{(}
    \AttributeTok{n =} \FunctionTok{n}\NormalTok{(),}
    \AttributeTok{mean\_cites =} \FunctionTok{round}\NormalTok{(}\FunctionTok{mean}\NormalTok{(cited\_by\_count), }\DecValTok{1}\NormalTok{),}
    \AttributeTok{median\_cites =} \FunctionTok{median}\NormalTok{(cited\_by\_count),}
    \AttributeTok{max\_cites =} \FunctionTok{max}\NormalTok{(cited\_by\_count),}
    \AttributeTok{pct\_uncited =}\NormalTok{ scales}\SpecialCharTok{::}\FunctionTok{percent}\NormalTok{(}\FunctionTok{mean}\NormalTok{(cited\_by\_count }\SpecialCharTok{==} \DecValTok{0}\NormalTok{))}
\NormalTok{  )}

\NormalTok{cite\_stats }\SpecialCharTok{|\textgreater{}} \FunctionTok{gt}\NormalTok{()}
\end{Highlighting}
\end{Shaded}

\begingroup
\fontsize{8.0pt}{9.0pt}\selectfont
\begin{longtable}{>{\raggedleft\arraybackslash}p{\dimexpr 0.18\linewidth -2\tabcolsep-1.5\arrayrulewidth}>{\raggedleft\arraybackslash}p{\dimexpr 0.20\linewidth -2\tabcolsep-1.5\arrayrulewidth}>{\raggedleft\arraybackslash}p{\dimexpr 0.22\linewidth -2\tabcolsep-1.5\arrayrulewidth}>{\raggedleft\arraybackslash}p{\dimexpr 0.19\linewidth -2\tabcolsep-1.5\arrayrulewidth}>{\raggedleft\arraybackslash}p{\dimexpr 0.21\linewidth -2\tabcolsep-1.5\arrayrulewidth}}
\toprule
n & \nolinkurl{mean_cites} & \nolinkurl{median_cites} & \nolinkurl{max_cites} & \nolinkurl{pct_uncited} \\ 
\midrule\addlinespace[2.5pt]
200 & 11.8 & 2 & 576 & 38\% \\ 
\bottomrule
\end{longtable}
\endgroup

\subsection{Retracted papers by source}\label{retracted-papers-by-source}

\begin{Shaded}
\begin{Highlighting}[]
\NormalTok{retracted }\SpecialCharTok{|\textgreater{}}
  \FunctionTok{filter}\NormalTok{(}\SpecialCharTok{!}\FunctionTok{is.na}\NormalTok{(source\_display\_name)) }\SpecialCharTok{|\textgreater{}}
  \FunctionTok{count}\NormalTok{(source\_display\_name, }\AttributeTok{sort =} \ConstantTok{TRUE}\NormalTok{) }\SpecialCharTok{|\textgreater{}}
  \FunctionTok{head}\NormalTok{(}\DecValTok{10}\NormalTok{) }\SpecialCharTok{|\textgreater{}}
  \FunctionTok{mutate}\NormalTok{(}\AttributeTok{source\_display\_name =} \FunctionTok{fct\_reorder}\NormalTok{(source\_display\_name, n)) }\SpecialCharTok{|\textgreater{}}
  \FunctionTok{ggplot}\NormalTok{(}\FunctionTok{aes}\NormalTok{(}\AttributeTok{x =}\NormalTok{ n, }\AttributeTok{y =}\NormalTok{ source\_display\_name)) }\SpecialCharTok{+}
  \FunctionTok{geom\_col}\NormalTok{(}\AttributeTok{fill =} \FunctionTok{palette\_sci}\NormalTok{(}\DecValTok{1}\NormalTok{)) }\SpecialCharTok{+}
  \FunctionTok{labs}\NormalTok{(}\AttributeTok{x =} \StringTok{"Retracted papers"}\NormalTok{, }\AttributeTok{y =} \ConstantTok{NULL}\NormalTok{) }\SpecialCharTok{+}
  \FunctionTok{theme\_sci}\NormalTok{()}
\end{Highlighting}
\end{Shaded}

\begin{figure}

{\centering \includegraphics[width=0.9\linewidth,alt={Horizontal bar chart showing which journals have the most retracted papers in the retrieved sample.}]{scientometrics-in-r_files/figure-latex/29-fig-retraction-sources-1} 

}

\caption{Top journals with the most retracted papers in the sample.}\label{fig:29-fig-retraction-sources}
\end{figure}

\subsection{Citation persistence over time}\label{citation-persistence-over-time}

\begin{Shaded}
\begin{Highlighting}[]
\NormalTok{retracted }\SpecialCharTok{|\textgreater{}}
  \FunctionTok{mutate}\NormalTok{(}\AttributeTok{age =} \DecValTok{2024} \SpecialCharTok{{-}} \FunctionTok{year}\NormalTok{(publication\_date)) }\SpecialCharTok{|\textgreater{}}
  \FunctionTok{filter}\NormalTok{(age }\SpecialCharTok{\textgreater{}=} \DecValTok{1}\NormalTok{, age }\SpecialCharTok{\textless{}=} \DecValTok{12}\NormalTok{) }\SpecialCharTok{|\textgreater{}}
  \FunctionTok{group\_by}\NormalTok{(age) }\SpecialCharTok{|\textgreater{}}
  \FunctionTok{summarise}\NormalTok{(}\AttributeTok{mean\_cites =} \FunctionTok{mean}\NormalTok{(cited\_by\_count), }\AttributeTok{n =} \FunctionTok{n}\NormalTok{(), }\AttributeTok{.groups =} \StringTok{"drop"}\NormalTok{) }\SpecialCharTok{|\textgreater{}}
  \FunctionTok{ggplot}\NormalTok{(}\FunctionTok{aes}\NormalTok{(}\AttributeTok{x =}\NormalTok{ age, }\AttributeTok{y =}\NormalTok{ mean\_cites)) }\SpecialCharTok{+}
  \FunctionTok{geom\_line}\NormalTok{(}\AttributeTok{linewidth =} \DecValTok{1}\NormalTok{, }\AttributeTok{colour =} \FunctionTok{palette\_sci}\NormalTok{(}\DecValTok{1}\NormalTok{)) }\SpecialCharTok{+}
  \FunctionTok{geom\_point}\NormalTok{(}\AttributeTok{size =} \DecValTok{2}\NormalTok{, }\AttributeTok{colour =} \FunctionTok{palette\_sci}\NormalTok{(}\DecValTok{1}\NormalTok{)) }\SpecialCharTok{+}
  \FunctionTok{labs}\NormalTok{(}\AttributeTok{x =} \StringTok{"Years since publication"}\NormalTok{, }\AttributeTok{y =} \StringTok{"Mean citations"}\NormalTok{) }\SpecialCharTok{+}
  \FunctionTok{theme\_sci}\NormalTok{()}
\end{Highlighting}
\end{Shaded}

\begin{figure}

{\centering \includegraphics[width=0.9\linewidth,alt={Line chart showing the mean citation count of retracted papers by years since publication, demonstrating that citations continue to accumulate.}]{scientometrics-in-r_files/figure-latex/29-fig-retraction-persistence-1} 

}

\caption{Mean citations by publication age for retracted papers.}\label{fig:29-fig-retraction-persistence}
\end{figure}

\section{Diagnostics and interpretation}\label{retraction-diagnostics}

\begin{itemize}
\tightlist
\item
  \textbf{Coverage}: OpenAlex retraction coverage depends on publisher metadata. Not all retractions are flagged. Cross-check with Retraction Watch for comprehensive coverage.
\item
  \textbf{Retraction vs.~publication date}: The publication date is when the paper was published, not when it was retracted. The gap between these dates varies from months to decades.
\item
  \textbf{Positive vs.~negative citations}: Citation counts do not distinguish between ``this paper found X'' (positive, problematic) and ``this paper was retracted for Y'' (negative, appropriate). Manual checking of citing contexts is needed.
\item
  \textbf{Denominator}: Retraction \emph{rates} require the total number of publications as denominator. Raw retraction counts favour large journals.
\end{itemize}

\section{Limitations and responsible use}\label{retraction-limitations}

\section{Limitations and responsible use}\label{limitations-and-responsible-use-28}

\begin{itemize}
\tightlist
\item
  \textbf{Retractions are not all equal.} Some retractions reflect honest error, others fraud. The reason for retraction matters for interpretation but is often not coded in bibliometric databases.
\item
  \textbf{Stigmatisation risk.} Listing authors or institutions with high retraction counts can be stigmatising. Report aggregate patterns, not individual-level data, unless specifically studying misconduct.
\item
  \textbf{Database lag.} Retraction notices may take months or years to propagate through databases. Recent retraction data is incomplete.
\item
  \textbf{Self-correction is broader than retraction.} Corrections, expressions of concern, and errata also contribute to self-correction but are harder to track bibliometrically.
\item
  \textbf{Do not equate retractions with research quality.} Fields with more scrutiny detect more problems. A high retraction rate may indicate effective oversight, not worse science \citep{hicks2015}.
\end{itemize}

\section{Common pitfalls}\label{retraction-pitfalls}

\section{Common pitfalls}\label{common-pitfalls-28}

\begin{itemize}
\tightlist
\item
  \textbf{Counting retractions without denominators.} A journal with 10 retractions from 50,000 papers has a lower retraction rate than one with 5 from 1,000.
\item
  \textbf{Assuming all post-retraction citations are problematic.} Some citing papers may specifically discuss the retraction or cite the paper for methodology that remains valid despite data issues.
\item
  \textbf{Ignoring retraction reasons.} Honest errors should be distinguished from misconduct in any analysis of research integrity.
\item
  \textbf{Using retraction analysis for individual evaluation.} Authors of retracted papers may be victims of co-author misconduct. Never use retraction counts as performance indicators.
\end{itemize}

\section{Exercises}\label{retraction-exercises}

\begin{enumerate}
\def\labelenumi{\arabic{enumi}.}
\item
  \textbf{Retraction rate by field.} Fetch retracted works from two different fields. Compare the retraction rate (retractions per 1,000 publications). Account for corpus size.
\item
  \textbf{Citation trajectory.} For a retracted paper with many citations, examine when the citations accumulated. Were most citations received before or after the retraction?
\item
  \textbf{Retraction-citation correlation.} Is there a relationship between a paper's citation count at the time of retraction and the number of post-retraction citations it receives?
\end{enumerate}

\section{Solutions}\label{29-solutions}

Solutions are provided in \ref{solutions}.

\section{Further reading}\label{retraction-reading}

\begin{itemize}
\tightlist
\item
  \citet{fang2012} --- Large-scale analysis of retraction patterns and their causes.
\item
  \citet{hicks2015} --- The Leiden Manifesto: responsible use of metrics in the context of research integrity.
\item
  \citet{garfield1955} --- Citation indexing: the infrastructure that makes retraction tracking possible.
\end{itemize}

\section{Session info}\label{retraction-sessioninfo}

\begin{Shaded}
\begin{Highlighting}[]
\FunctionTok{sessionInfo}\NormalTok{()}
\end{Highlighting}
\end{Shaded}

\begin{verbatim}
#> R version 4.6.1 (2026-06-24)
#> Platform: x86_64-pc-linux-gnu
#> Running under: Ubuntu 24.04.5 LTS
#> 
#> Matrix products: default
#> BLAS:   /usr/lib/x86_64-linux-gnu/openblas-pthread/libblas.so.3 
#> LAPACK: /usr/lib/x86_64-linux-gnu/openblas-pthread/libopenblasp-r0.3.26.so;  LAPACK version 3.12.0
#> 
#> locale:
#>  [1] LC_CTYPE=C.UTF-8       LC_NUMERIC=C           LC_TIME=C.UTF-8       
#>  [4] LC_COLLATE=C.UTF-8     LC_MONETARY=C.UTF-8    LC_MESSAGES=C.UTF-8   
#>  [7] LC_PAPER=C.UTF-8       LC_NAME=C              LC_ADDRESS=C          
#> [10] LC_TELEPHONE=C         LC_MEASUREMENT=C.UTF-8 LC_IDENTIFICATION=C   
#> 
#> time zone: UTC
#> tzcode source: system (glibc)
#> 
#> attached base packages:
#> [1] stats     graphics  grDevices utils     datasets  methods   base     
#> 
#> other attached packages:
#>  [1] uwot_0.2.5                Matrix_1.7-5             
#>  [3] word2vec_0.4.1            stm_1.3.8                
#>  [5] topicmodels_0.2-17        quanteda.textstats_0.97.2
#>  [7] visNetwork_2.1.4          ggraph_2.2.2             
#>  [9] tidygraph_1.3.1           igraph_2.3.4             
#> [11] quanteda_4.5.1            pdftools_3.9.1           
#> [13] arrow_26.0.0              bibliometrix_5.5.0       
#> [15] RefManageR_1.4.0          bib2df_1.1.2.0           
#> [17] rcrossref_1.2.1           gt_1.4.0                 
#> [19] tidytext_0.4.3            glue_1.8.1               
#> [21] openalexR_3.1.2           lubridate_1.9.5          
#> [23] forcats_1.0.1             stringr_1.6.0            
#> [25] dplyr_1.2.1               purrr_1.2.2              
#> [27] readr_2.2.0               tidyr_1.3.2              
#> [29] tibble_3.3.1              ggplot2_4.0.3            
#> [31] tidyverse_2.0.0          
#> 
#> loaded via a namespace (and not attached):
#>   [1] bibtex_0.5.3           RColorBrewer_1.1-3     rstudioapi_0.19.0     
#>   [4] jsonlite_2.0.0         magrittr_2.0.5         modeltools_0.2-25     
#>   [7] farver_2.1.2           rmarkdown_2.32         fs_2.1.0              
#>  [10] vctrs_0.7.3            memoise_2.0.1          askpass_1.2.1         
#>  [13] base64enc_0.1-6        htmltools_0.5.9        contentanalysis_1.1.1 
#>  [16] curl_8.1.0             janeaustenr_1.0.0      cellranger_1.1.0      
#>  [19] htmlwidgets_1.6.4      tokenizers_0.3.0       plyr_1.8.9            
#>  [22] httr2_1.3.0            cachem_1.1.0           plotly_4.12.1         
#>  [25] dimensionsR_0.0.3      mime_0.13              lifecycle_1.0.5       
#>  [28] pkgconfig_2.0.3        R6_2.6.1               fastmap_1.2.0         
#>  [31] shiny_1.14.0           digest_0.6.39          patchwork_1.3.2       
#>  [34] shinycssloaders_1.1.0  rprojroot_2.1.1        RSpectra_0.16-2       
#>  [37] SnowballC_0.7.1        labeling_0.4.3         urltools_1.7.3.1      
#>  [40] timechange_0.4.0       mgcv_1.9-4             httr_1.4.9            
#>  [43] polyclip_1.10-7        compiler_4.6.1         here_1.0.2            
#>  [46] bit64_4.8.6            withr_3.0.3            S7_0.2.2              
#>  [49] viridis_0.6.5          ggforce_0.5.0          MASS_7.3-65           
#>  [52] bibliometrixData_0.3.0 tools_4.6.1            otel_0.2.0            
#>  [55] stopwords_2.3          zip_3.0.2              rentrez_1.2.4         
#>  [58] httpuv_1.6.17          nlme_3.1-169           promises_1.5.0        
#>  [61] grid_4.6.1             stringdist_0.9.17      reshape2_1.4.5        
#>  [64] generics_0.1.4         gtable_0.3.6           tzdb_0.5.0            
#>  [67] rscopus_0.9.0          ca_0.72                data.table_1.18.6.1   
#>  [70] hms_1.1.4              xml2_1.6.0             utf8_1.2.6            
#>  [73] ggrepel_0.9.8          pillar_1.11.1          nsyllable_1.0.1       
#>  [76] vroom_1.7.1            later_1.4.8            splines_4.6.1         
#>  [79] tweenr_2.0.3           lattice_0.22-9         FNN_1.1.4.1           
#>  [82] bit_4.6.0              tidyselect_1.2.1       tm_0.7-20             
#>  [85] miniUI_0.1.2           knitr_1.52             gridExtra_2.3.1       
#>  [88] NLP_0.3-3              bookdown_0.48          stats4_4.6.1          
#>  [91] crul_1.6.0             xfun_0.62              graphlayouts_1.2.5    
#>  [94] matrixStats_1.5.0      DT_0.34.0              humaniformat_0.6.0    
#>  [97] stringi_1.8.9          qpdf_1.4.1             yaml_2.3.12           
#> [100] evaluate_1.0.5         httpcode_0.3.0         cli_3.6.6             
#> [103] xtable_1.8-8           systemfonts_1.3.2      Rcpp_1.1.2            
#> [106] readxl_1.6.0           triebeard_0.4.1        XML_3.99-0.25         
#> [109] parallel_4.6.1         assertthat_0.2.1       pubmedR_1.0.2         
#> [112] slam_0.1-56            viridisLite_0.4.3      scales_1.4.0          
#> [115] crayon_1.5.3           openxlsx_4.2.9         rlang_1.3.0           
#> [118] fastmatch_1.1-8
\end{verbatim}

\chapter{Funding Analysis and Grant-to-Publication Linkage}\label{funding-analysis-and-grant-to-publication-linkage}

\emph{R. Heller · Date: 2026-05-09 · Last changed: 2026-10-31}

\section{Learning objectives}\label{funding-objectives}

After completing this chapter, you will be able to:

\begin{itemize}
\tightlist
\item
  Extract funder information from OpenAlex work records
\item
  Compute funder-level publication counts and citation impact
\item
  Analyse co-funding patterns (papers acknowledging multiple funders)
\item
  Interpret funding data critically, accounting for coverage and attribution challenges
\item
  Discuss why funding-output linkage is inherently imprecise
\end{itemize}

\section{Setup}\label{funding-setup}

\begin{Shaded}
\begin{Highlighting}[]
\FunctionTok{library}\NormalTok{(tidyverse)}
\FunctionTok{library}\NormalTok{(openalexR)}
\FunctionTok{library}\NormalTok{(glue)}
\FunctionTok{library}\NormalTok{(gt)}

\FunctionTok{set.seed}\NormalTok{(}\DecValTok{20260509}\NormalTok{)}

\FunctionTok{source}\NormalTok{(here}\SpecialCharTok{::}\FunctionTok{here}\NormalTok{(}\StringTok{"R"}\NormalTok{, }\StringTok{"api\_helpers.R"}\NormalTok{))}
\FunctionTok{source}\NormalTok{(here}\SpecialCharTok{::}\FunctionTok{here}\NormalTok{(}\StringTok{"R"}\NormalTok{, }\StringTok{"utils.R"}\NormalTok{))}
\FunctionTok{source}\NormalTok{(here}\SpecialCharTok{::}\FunctionTok{here}\NormalTok{(}\StringTok{"R"}\NormalTok{, }\StringTok{"funding\_helpers.R"}\NormalTok{))}
\FunctionTok{source}\NormalTok{(here}\SpecialCharTok{::}\FunctionTok{here}\NormalTok{(}\StringTok{"R"}\NormalTok{, }\StringTok{"sci\_palette.R"}\NormalTok{))}
\end{Highlighting}
\end{Shaded}

\section{Conceptual background}\label{funding-background}

Funding agencies invest billions in research and want to understand what their investment produces. \textbf{Funding acknowledgement analysis} extracts funder names and grant identifiers from publication metadata or full text, linking grants to the papers they supported.

OpenAlex includes funder information for many works, derived from Crossref's funding metadata (which in turn comes from publisher-reported FundRef data). Coverage varies: journals that participate in Crossref's FundRef initiative have good coverage; others may have no funding data at all \citep{priem2022}.

Key analytical questions include: How many publications does a funder support? What is the citation impact of funded vs.~unfunded research? Which funders frequently co-fund research? How does a funder's portfolio map onto the landscape of science?

The fundamental challenge is \textbf{attribution}: a paper acknowledging three funders and seven authors is the result of contributions from all of them. Assigning full credit to each funder double-counts; fractional credit is more principled but harder to implement. Similar issues arise with multi-source funding: a paper may acknowledge an infrastructure grant, a project grant, and a fellowship, each with different levels of contribution.

\section{Worked example}\label{funding-example}

\subsection{Extracting funder data}\label{extracting-funder-data}

\begin{Shaded}
\begin{Highlighting}[]
\NormalTok{works }\OtherTok{\textless{}{-}} \FunctionTok{fetch\_openalex}\NormalTok{(}
  \AttributeTok{entity =} \StringTok{"works"}\NormalTok{,}
  \AttributeTok{primary\_location.source.id =} \StringTok{"S148561398"}\NormalTok{,}
  \AttributeTok{from\_publication\_date =} \StringTok{"2020{-}01{-}01"}\NormalTok{,}
  \AttributeTok{to\_publication\_date =} \StringTok{"2023{-}12{-}31"}\NormalTok{,}
  \AttributeTok{options =} \FunctionTok{list}\NormalTok{(}\AttributeTok{sample =} \DecValTok{400}\NormalTok{, }\AttributeTok{seed =} \DecValTok{42}\NormalTok{)}
\NormalTok{)}

\NormalTok{funding\_works }\OtherTok{\textless{}{-}}\NormalTok{ works }\SpecialCharTok{|\textgreater{}}
  \FunctionTok{mutate}\NormalTok{(}\AttributeTok{funders =} \FunctionTok{map}\NormalTok{(funders, sci\_funder\_rows),}
         \AttributeTok{has\_funder\_metadata =} \FunctionTok{map\_int}\NormalTok{(funders, nrow) }\SpecialCharTok{\textgreater{}} \DecValTok{0}\NormalTok{)}

\NormalTok{funded }\OtherTok{\textless{}{-}}\NormalTok{ funding\_works }\SpecialCharTok{|\textgreater{}}
  \FunctionTok{filter}\NormalTok{(has\_funder\_metadata) }\SpecialCharTok{|\textgreater{}}
  \FunctionTok{select}\NormalTok{(id, display\_name, cited\_by\_count, funders)}

\FunctionTok{cat}\NormalTok{(}\FunctionTok{glue}\NormalTok{(}\StringTok{"Works with observed funder IDs: \{nrow(funded)\} / \{nrow(works)\} (\{scales::percent(nrow(funded)/nrow(works))\})}\SpecialCharTok{\textbackslash{}n}\StringTok{"}\NormalTok{))}
\end{Highlighting}
\end{Shaded}

\begin{verbatim}
#> Works with observed funder IDs: 224 / 400 (56%)
\end{verbatim}

\begin{Shaded}
\begin{Highlighting}[]
\ControlFlowTok{if}\NormalTok{ (}\SpecialCharTok{!}\FunctionTok{nrow}\NormalTok{(funded)) \{}
  \FunctionTok{cat}\NormalTok{(}\StringTok{"No usable funder records occur in this sample. Funder summaries are unavailable; missing metadata does not establish that the papers were unfunded.}\SpecialCharTok{\textbackslash{}n}\StringTok{"}\NormalTok{)}
\NormalTok{\}}
\end{Highlighting}
\end{Shaded}

\begin{Shaded}
\begin{Highlighting}[]
\NormalTok{funder\_data }\OtherTok{\textless{}{-}} \ControlFlowTok{if}\NormalTok{ (}\FunctionTok{nrow}\NormalTok{(funded)) \{}
\NormalTok{  funded }\SpecialCharTok{|\textgreater{}}
    \FunctionTok{unnest}\NormalTok{(funders, }\AttributeTok{names\_sep =} \StringTok{"\_"}\NormalTok{) }\SpecialCharTok{|\textgreater{}}
    \FunctionTok{select}\NormalTok{(}\AttributeTok{work\_id =}\NormalTok{ id, }\AttributeTok{funder\_id =}\NormalTok{ funders\_id,}
           \AttributeTok{funder =}\NormalTok{ funders\_display\_name, cited\_by\_count)}
\NormalTok{\} }\ControlFlowTok{else}\NormalTok{ \{}
  \FunctionTok{tibble}\NormalTok{(}\AttributeTok{work\_id =} \FunctionTok{character}\NormalTok{(), }\AttributeTok{funder\_id =} \FunctionTok{character}\NormalTok{(),}
         \AttributeTok{funder =} \FunctionTok{character}\NormalTok{(), }\AttributeTok{cited\_by\_count =} \FunctionTok{numeric}\NormalTok{())}
\NormalTok{\}}

\FunctionTok{cat}\NormalTok{(}\FunctionTok{glue}\NormalTok{(}\StringTok{"Funder{-}work pairs: \{nrow(funder\_data)\}}\SpecialCharTok{\textbackslash{}n}\StringTok{"}\NormalTok{))}
\end{Highlighting}
\end{Shaded}

\begin{verbatim}
#> Funder-work pairs: 357
\end{verbatim}

\begin{Shaded}
\begin{Highlighting}[]
\FunctionTok{cat}\NormalTok{(}\FunctionTok{glue}\NormalTok{(}\StringTok{"Unique funders: \{n\_distinct(funder\_data$funder\_id)\}}\SpecialCharTok{\textbackslash{}n}\StringTok{"}\NormalTok{))}
\end{Highlighting}
\end{Shaded}

\begin{verbatim}
#> Unique funders: 186
\end{verbatim}

\subsection{Top funders by publication count}\label{top-funders-by-publication-count}

\begin{Shaded}
\begin{Highlighting}[]
\NormalTok{funder\_counts }\OtherTok{\textless{}{-}}\NormalTok{ funder\_data }\SpecialCharTok{|\textgreater{}}
  \FunctionTok{group\_by}\NormalTok{(funder\_id) }\SpecialCharTok{|\textgreater{}}
  \FunctionTok{summarise}\NormalTok{(}\AttributeTok{funder =} \FunctionTok{first}\NormalTok{(funder), }\AttributeTok{n =} \FunctionTok{n\_distinct}\NormalTok{(work\_id), }\AttributeTok{.groups =} \StringTok{"drop"}\NormalTok{) }\SpecialCharTok{|\textgreater{}}
  \FunctionTok{arrange}\NormalTok{(}\FunctionTok{desc}\NormalTok{(n))}

\ControlFlowTok{if}\NormalTok{ (}\FunctionTok{nrow}\NormalTok{(funder\_counts)) \{}
\NormalTok{funder\_counts }\SpecialCharTok{|\textgreater{}}
  \FunctionTok{head}\NormalTok{(}\DecValTok{15}\NormalTok{) }\SpecialCharTok{|\textgreater{}}
  \FunctionTok{mutate}\NormalTok{(}\AttributeTok{funder =} \FunctionTok{fct\_reorder}\NormalTok{(funder, n)) }\SpecialCharTok{|\textgreater{}}
  \FunctionTok{ggplot}\NormalTok{(}\FunctionTok{aes}\NormalTok{(}\AttributeTok{x =}\NormalTok{ n, }\AttributeTok{y =}\NormalTok{ funder)) }\SpecialCharTok{+}
  \FunctionTok{geom\_col}\NormalTok{(}\AttributeTok{fill =} \FunctionTok{palette\_sci}\NormalTok{(}\DecValTok{1}\NormalTok{)) }\SpecialCharTok{+}
  \FunctionTok{labs}\NormalTok{(}\AttributeTok{x =} \StringTok{"Publications supported"}\NormalTok{, }\AttributeTok{y =} \ConstantTok{NULL}\NormalTok{) }\SpecialCharTok{+}
  \FunctionTok{theme\_sci}\NormalTok{()}
\NormalTok{\} }\ControlFlowTok{else}\NormalTok{ \{}
  \FunctionTok{cat}\NormalTok{(}\StringTok{"No funder publication{-}count plot: no observed funder records.}\SpecialCharTok{\textbackslash{}n}\StringTok{"}\NormalTok{)}
\NormalTok{\}}
\end{Highlighting}
\end{Shaded}

\begin{figure}

{\centering \includegraphics[width=0.9\linewidth,alt={Horizontal bar chart showing the 15 funders with the most publications in the sample.}]{scientometrics-in-r_files/figure-latex/30-fig-funding-top-1} 

}

\caption{Up to 15 funders by the number of sampled publications with observed funder metadata.}\label{fig:30-fig-funding-top}
\end{figure}

\subsection{Citation impact by funder}\label{citation-impact-by-funder}

\begin{Shaded}
\begin{Highlighting}[]
\NormalTok{funder\_impact }\OtherTok{\textless{}{-}}\NormalTok{ funder\_data }\SpecialCharTok{|\textgreater{}}
  \FunctionTok{group\_by}\NormalTok{(funder\_id) }\SpecialCharTok{|\textgreater{}}
  \FunctionTok{summarise}\NormalTok{(}
    \AttributeTok{funder =} \FunctionTok{first}\NormalTok{(funder),}
    \AttributeTok{n\_pubs =} \FunctionTok{n}\NormalTok{(),}
    \AttributeTok{mean\_cites =} \FunctionTok{round}\NormalTok{(}\FunctionTok{mean}\NormalTok{(cited\_by\_count), }\DecValTok{1}\NormalTok{),}
    \AttributeTok{median\_cites =} \FunctionTok{median}\NormalTok{(cited\_by\_count),}
    \AttributeTok{.groups =} \StringTok{"drop"}
\NormalTok{  ) }\SpecialCharTok{|\textgreater{}}
  \FunctionTok{filter}\NormalTok{(n\_pubs }\SpecialCharTok{\textgreater{}=} \DecValTok{5}\NormalTok{) }\SpecialCharTok{|\textgreater{}}
  \FunctionTok{arrange}\NormalTok{(}\FunctionTok{desc}\NormalTok{(mean\_cites))}

\ControlFlowTok{if}\NormalTok{ (}\FunctionTok{nrow}\NormalTok{(funder\_impact)) \{}
\NormalTok{  funder\_impact }\SpecialCharTok{|\textgreater{}} \FunctionTok{head}\NormalTok{(}\DecValTok{10}\NormalTok{) }\SpecialCharTok{|\textgreater{}} \FunctionTok{gt}\NormalTok{()}
\NormalTok{\} }\ControlFlowTok{else}\NormalTok{ \{}
  \FunctionTok{cat}\NormalTok{(}\StringTok{"No funder has at least five sampled publications with usable metadata; no funder citation summary is reported.}\SpecialCharTok{\textbackslash{}n}\StringTok{"}\NormalTok{)}
\NormalTok{\}}
\end{Highlighting}
\end{Shaded}

\begingroup
\fontsize{8.0pt}{9.0pt}\selectfont
\begin{longtable}{>{\raggedright\arraybackslash}p{\dimexpr 0.25\linewidth -2\tabcolsep-1.5\arrayrulewidth}>{\raggedright\arraybackslash}p{\dimexpr 0.34\linewidth -2\tabcolsep-1.5\arrayrulewidth}>{\raggedleft\arraybackslash}p{\dimexpr 0.12\linewidth -2\tabcolsep-1.5\arrayrulewidth}>{\raggedleft\arraybackslash}p{\dimexpr 0.14\linewidth -2\tabcolsep-1.5\arrayrulewidth}>{\raggedleft\arraybackslash}p{\dimexpr 0.15\linewidth -2\tabcolsep-1.5\arrayrulewidth}}
\toprule
\nolinkurl{funder_id} & funder & n\_pubs & \nolinkurl{mean_cites} & \nolinkurl{median_cites} \\ 
\midrule\addlinespace[2.5pt]
\nolinkurl{https://openalex.org/F4320335869} & National Social Science Fund of China & 7 & 43.6 & 21 \\ 
\nolinkurl{https://openalex.org/F4320335598} & Agencia Estatal de Investigación & 5 & 35.2 & 22 \\ 
\nolinkurl{https://openalex.org/F4320321001} & National Natural Science Foundation of China & 44 & 27.5 & 25 \\ 
\nolinkurl{https://openalex.org/F4320321091} & Coordenação de Aperfeiçoamento de Pessoal de Nível Superior & 7 & 25.0 & 22 \\ 
\nolinkurl{https://openalex.org/F4320334764} & Japan Society for the Promotion of Science & 5 & 24.0 & 10 \\ 
\nolinkurl{https://openalex.org/F4320321106} & Ministry of Education of the People's Republic of China & 5 & 19.8 & 17 \\ 
\nolinkurl{https://openalex.org/F4320321114} & Bundesministerium für Bildung und Forschung & 14 & 19.1 & 17 \\ 
\nolinkurl{https://openalex.org/F4320335787} & Fundamental Research Funds for the Central Universities & 8 & 18.6 & 16 \\ 
\nolinkurl{https://openalex.org/F4320320300} & European Commission & 9 & 15.9 & 14 \\ 
\nolinkurl{https://openalex.org/F4320321873} & Ministero dell’Istruzione, dell’Università e della Ricerca & 5 & 10.4 & 10 \\ 
\bottomrule
\end{longtable}
\endgroup

\subsection{Comparison by observed funder metadata}\label{comparison-by-observed-funder-metadata}

\begin{Shaded}
\begin{Highlighting}[]
\ControlFlowTok{if}\NormalTok{ (}\FunctionTok{n\_distinct}\NormalTok{(funding\_works}\SpecialCharTok{$}\NormalTok{has\_funder\_metadata) }\SpecialCharTok{==} \DecValTok{2}\DataTypeTok{L}\NormalTok{) \{}
\NormalTok{funding\_works }\SpecialCharTok{|\textgreater{}}
  \FunctionTok{mutate}\NormalTok{(}\AttributeTok{funder\_metadata =} \FunctionTok{if\_else}\NormalTok{(has\_funder\_metadata,}
                                  \StringTok{"Observed funder"}\NormalTok{, }\StringTok{"No observed funder"}\NormalTok{)) }\SpecialCharTok{|\textgreater{}}
  \FunctionTok{ggplot}\NormalTok{(}\FunctionTok{aes}\NormalTok{(}\AttributeTok{x =}\NormalTok{ funder\_metadata, }\AttributeTok{y =}\NormalTok{ cited\_by\_count }\SpecialCharTok{+} \DecValTok{1}\NormalTok{)) }\SpecialCharTok{+}
  \FunctionTok{geom\_boxplot}\NormalTok{(}\AttributeTok{fill =} \FunctionTok{palette\_sci}\NormalTok{(}\DecValTok{1}\NormalTok{), }\AttributeTok{alpha =} \FloatTok{0.7}\NormalTok{) }\SpecialCharTok{+}
  \FunctionTok{scale\_y\_log10}\NormalTok{() }\SpecialCharTok{+}
  \FunctionTok{labs}\NormalTok{(}\AttributeTok{x =} \StringTok{"Funder metadata"}\NormalTok{, }\AttributeTok{y =} \StringTok{"Citations (log scale)"}\NormalTok{) }\SpecialCharTok{+}
  \FunctionTok{theme\_sci}\NormalTok{()}
\NormalTok{\} }\ControlFlowTok{else}\NormalTok{ \{}
  \FunctionTok{cat}\NormalTok{(}\StringTok{"No two{-}group citation comparison: only one funder{-}metadata group is observed.}\SpecialCharTok{\textbackslash{}n}\StringTok{"}\NormalTok{)}
\NormalTok{\}}
\end{Highlighting}
\end{Shaded}

\begin{figure}

{\centering \includegraphics[width=0.9\linewidth,alt={Box plot comparing scholarly citation counts for papers with and without observed funder identifiers in the retrieved metadata.}]{scientometrics-in-r_files/figure-latex/30-fig-funding-comparison-1} 

}

\caption{Citation distribution by presence of observed funder metadata; missing metadata does not establish absence of funding.}\label{fig:30-fig-funding-comparison}
\end{figure}

\section{Diagnostics and interpretation}\label{funding-diagnostics}

\begin{itemize}
\tightlist
\item
  \textbf{Coverage rate}: The fraction of papers with funder data depends on the journal, publisher, and field. Report this rate before interpreting funder statistics.
\item
  \textbf{Funder name normalisation}: The same funder may appear under different names (``NSF'', ``National Science Foundation'', ``US NSF''). Merge variants using funder IDs where available.
\item
  \textbf{Selection and metadata coverage}: Papers with observed funder records may differ systematically from papers without them. Their citation gap can reflect selection and publisher reporting; it does not identify the causal effect of funding or compare confirmed funded and unfunded groups.
\item
  \textbf{Multi-funder attribution}: Papers acknowledging multiple funders inflate the publication count for each. Consider fractional counting.
\end{itemize}

\section{Limitations and responsible use}\label{funding-limitations}

\section{Limitations and responsible use}\label{limitations-and-responsible-use-29}

\begin{itemize}
\tightlist
\item
  \textbf{Coverage is incomplete.} Not all publishers report funding data to Crossref. Missing funder acknowledgements do not mean research was unfunded.
\item
  \textbf{Attribution is ambiguous.} A grant acknowledgement does not specify which part of the work the grant supported. Infrastructure grants, salary support, and project funding all appear identically.
\item
  \textbf{Correlation ≠ causation.} Funded papers receive more citations, but this may reflect that better-resourced labs produce both more funding applications and more cited papers.
\item
  \textbf{Do not rank funders by citation impact.} Funders support different fields, career stages, and risk levels. Comparing mean citations across funders is misleading without field normalisation \citep{hicks2015}.
\end{itemize}

\section{Common pitfalls}\label{funding-pitfalls}

\section{Common pitfalls}\label{common-pitfalls-29}

\begin{itemize}
\tightlist
\item
  \textbf{Ignoring funder coverage rates.} Comparing funder portfolios without accounting for differential coverage produces biased results.
\item
  \textbf{Treating every acknowledgement as equal.} A paper may acknowledge a major project grant and a minor travel grant. Both appear identically in metadata.
\item
  \textbf{Comparing across fields.} A biomedical funder's publications will have higher citation counts than a humanities funder's publications due to field-level citation norms.
\item
  \textbf{Using funder data for individual evaluation.} Whether a researcher has funding reflects many factors beyond research quality: field norms, career stage, institutional support, and luck.
\end{itemize}

\section{Exercises}\label{funding-exercises}

\begin{enumerate}
\def\labelenumi{\arabic{enumi}.}
\item
  \textbf{Funder co-occurrence.} Identify papers with multiple funders. Which funder pairs most frequently co-fund research?
\item
  \textbf{Funder portfolios.} For two major funders, fetch their supported publications and compare topical profiles using keyword analysis.
\item
  \textbf{Temporal trends.} Track the number of papers with funding acknowledgements by year. Is the coverage rate increasing over time?
\end{enumerate}

\section{Solutions}\label{30-solutions}

Solutions are provided in \ref{solutions}.

\section{Further reading}\label{funding-reading}

\begin{itemize}
\tightlist
\item
  \citet{priem2022} --- OpenAlex funder metadata from Crossref FundRef.
\item
  \citet{hicks2015} --- Responsible use of metrics in funding evaluation.
\item
  \citet{waltman2016} --- Citation indicators and their application to funding assessment.
\end{itemize}

\section{Session info}\label{funding-sessioninfo}

\begin{Shaded}
\begin{Highlighting}[]
\FunctionTok{sessionInfo}\NormalTok{()}
\end{Highlighting}
\end{Shaded}

\begin{verbatim}
#> R version 4.6.1 (2026-06-24)
#> Platform: x86_64-pc-linux-gnu
#> Running under: Ubuntu 24.04.5 LTS
#> 
#> Matrix products: default
#> BLAS:   /usr/lib/x86_64-linux-gnu/openblas-pthread/libblas.so.3 
#> LAPACK: /usr/lib/x86_64-linux-gnu/openblas-pthread/libopenblasp-r0.3.26.so;  LAPACK version 3.12.0
#> 
#> locale:
#>  [1] LC_CTYPE=C.UTF-8       LC_NUMERIC=C           LC_TIME=C.UTF-8       
#>  [4] LC_COLLATE=C.UTF-8     LC_MONETARY=C.UTF-8    LC_MESSAGES=C.UTF-8   
#>  [7] LC_PAPER=C.UTF-8       LC_NAME=C              LC_ADDRESS=C          
#> [10] LC_TELEPHONE=C         LC_MEASUREMENT=C.UTF-8 LC_IDENTIFICATION=C   
#> 
#> time zone: UTC
#> tzcode source: system (glibc)
#> 
#> attached base packages:
#> [1] stats     graphics  grDevices utils     datasets  methods   base     
#> 
#> other attached packages:
#>  [1] uwot_0.2.5                Matrix_1.7-5             
#>  [3] word2vec_0.4.1            stm_1.3.8                
#>  [5] topicmodels_0.2-17        quanteda.textstats_0.97.2
#>  [7] visNetwork_2.1.4          ggraph_2.2.2             
#>  [9] tidygraph_1.3.1           igraph_2.3.4             
#> [11] quanteda_4.5.1            pdftools_3.9.1           
#> [13] arrow_26.0.0              bibliometrix_5.5.0       
#> [15] RefManageR_1.4.0          bib2df_1.1.2.0           
#> [17] rcrossref_1.2.1           gt_1.4.0                 
#> [19] tidytext_0.4.3            glue_1.8.1               
#> [21] openalexR_3.1.2           lubridate_1.9.5          
#> [23] forcats_1.0.1             stringr_1.6.0            
#> [25] dplyr_1.2.1               purrr_1.2.2              
#> [27] readr_2.2.0               tidyr_1.3.2              
#> [29] tibble_3.3.1              ggplot2_4.0.3            
#> [31] tidyverse_2.0.0          
#> 
#> loaded via a namespace (and not attached):
#>   [1] bibtex_0.5.3           RColorBrewer_1.1-3     rstudioapi_0.19.0     
#>   [4] jsonlite_2.0.0         magrittr_2.0.5         modeltools_0.2-25     
#>   [7] farver_2.1.2           rmarkdown_2.32         fs_2.1.0              
#>  [10] vctrs_0.7.3            memoise_2.0.1          askpass_1.2.1         
#>  [13] base64enc_0.1-6        htmltools_0.5.9        contentanalysis_1.1.1 
#>  [16] curl_8.1.0             janeaustenr_1.0.0      cellranger_1.1.0      
#>  [19] htmlwidgets_1.6.4      tokenizers_0.3.0       plyr_1.8.9            
#>  [22] httr2_1.3.0            cachem_1.1.0           plotly_4.12.1         
#>  [25] dimensionsR_0.0.3      mime_0.13              lifecycle_1.0.5       
#>  [28] pkgconfig_2.0.3        R6_2.6.1               fastmap_1.2.0         
#>  [31] shiny_1.14.0           digest_0.6.39          patchwork_1.3.2       
#>  [34] shinycssloaders_1.1.0  rprojroot_2.1.1        RSpectra_0.16-2       
#>  [37] SnowballC_0.7.1        labeling_0.4.3         urltools_1.7.3.1      
#>  [40] timechange_0.4.0       mgcv_1.9-4             httr_1.4.9            
#>  [43] polyclip_1.10-7        compiler_4.6.1         here_1.0.2            
#>  [46] bit64_4.8.6            withr_3.0.3            S7_0.2.2              
#>  [49] viridis_0.6.5          ggforce_0.5.0          MASS_7.3-65           
#>  [52] bibliometrixData_0.3.0 tools_4.6.1            otel_0.2.0            
#>  [55] stopwords_2.3          zip_3.0.2              rentrez_1.2.4         
#>  [58] httpuv_1.6.17          nlme_3.1-169           promises_1.5.0        
#>  [61] grid_4.6.1             stringdist_0.9.17      reshape2_1.4.5        
#>  [64] generics_0.1.4         gtable_0.3.6           tzdb_0.5.0            
#>  [67] rscopus_0.9.0          ca_0.72                data.table_1.18.6.1   
#>  [70] hms_1.1.4              xml2_1.6.0             utf8_1.2.6            
#>  [73] ggrepel_0.9.8          pillar_1.11.1          nsyllable_1.0.1       
#>  [76] vroom_1.7.1            later_1.4.8            splines_4.6.1         
#>  [79] tweenr_2.0.3           lattice_0.22-9         FNN_1.1.4.1           
#>  [82] bit_4.6.0              tidyselect_1.2.1       tm_0.7-20             
#>  [85] miniUI_0.1.2           knitr_1.52             gridExtra_2.3.1       
#>  [88] NLP_0.3-3              bookdown_0.48          stats4_4.6.1          
#>  [91] crul_1.6.0             xfun_0.62              graphlayouts_1.2.5    
#>  [94] matrixStats_1.5.0      DT_0.34.0              humaniformat_0.6.0    
#>  [97] stringi_1.8.9          qpdf_1.4.1             yaml_2.3.12           
#> [100] evaluate_1.0.5         httpcode_0.3.0         cli_3.6.6             
#> [103] xtable_1.8-8           systemfonts_1.3.2      Rcpp_1.1.2            
#> [106] readxl_1.6.0           triebeard_0.4.1        XML_3.99-0.25         
#> [109] parallel_4.6.1         assertthat_0.2.1       pubmedR_1.0.2         
#> [112] slam_0.1-56            viridisLite_0.4.3      scales_1.4.0          
#> [115] crayon_1.5.3           openxlsx_4.2.9         rlang_1.3.0           
#> [118] fastmatch_1.1-8
\end{verbatim}

\chapter{Patent--Paper Linkages}\label{patentpaper-linkages}

\emph{R. Heller · Date: 2026-05-09 · Last changed: 2026-10-31}

\section{Learning objectives}\label{patent-objectives}

After completing this chapter, you will be able to:

\begin{itemize}
\tightlist
\item
  Explain the significance of patent--paper citations for understanding research impact
\item
  Specify the paper and patent identifiers needed to establish an observed linkage
\item
  Explore scholarly citation metadata without treating it as patent-citation evidence
\item
  Discuss the limitations of patent citations as indicators of technological impact
\item
  Distinguish between direct citation (paper cited in patent) and indirect influence
\end{itemize}

\section{Setup}\label{patent-setup}

\begin{Shaded}
\begin{Highlighting}[]
\FunctionTok{library}\NormalTok{(tidyverse)}
\FunctionTok{library}\NormalTok{(openalexR)}
\FunctionTok{library}\NormalTok{(glue)}
\FunctionTok{library}\NormalTok{(gt)}

\FunctionTok{set.seed}\NormalTok{(}\DecValTok{20260509}\NormalTok{)}

\FunctionTok{source}\NormalTok{(here}\SpecialCharTok{::}\FunctionTok{here}\NormalTok{(}\StringTok{"R"}\NormalTok{, }\StringTok{"api\_helpers.R"}\NormalTok{))}
\FunctionTok{source}\NormalTok{(here}\SpecialCharTok{::}\FunctionTok{here}\NormalTok{(}\StringTok{"R"}\NormalTok{, }\StringTok{"utils.R"}\NormalTok{))}
\FunctionTok{source}\NormalTok{(here}\SpecialCharTok{::}\FunctionTok{here}\NormalTok{(}\StringTok{"R"}\NormalTok{, }\StringTok{"sci\_palette.R"}\NormalTok{))}
\end{Highlighting}
\end{Shaded}

\section{Conceptual background}\label{patent-background}

When a patent examiner or inventor cites a scientific paper in a patent document, it creates an observable link between knowledge production (science) and knowledge application (technology). These \textbf{patent--paper citations} (also called non-patent literature or NPL citations) are one of the few quantitative measures of the science--technology interface.

Patent citations serve a legal function: they define the prior art against which the patent's novelty is assessed. Not all patent--paper citations reflect genuine intellectual influence --- some are added by examiners for legal completeness, not because the inventor actually used the cited research. Nevertheless, at the aggregate level, patent citation patterns reveal which fields, journals, and institutions produce research that feeds into technological innovation.

The biomedical and chemical sciences dominate patent--paper linkages because pharmaceutical and chemical patents routinely cite scientific literature. Physics and engineering have moderate linkages. Social sciences and humanities are rarely cited in patents, not because they lack impact, but because their impact operates through different channels (policy, education, culture).

The OpenAlex fields retrieved in this example describe scholarly works and scholarly citations. They contain no citing patent identifiers, patent citation dates, or patent-family identifiers. Establishing a patent--paper linkage requires a separate, documented source of patent references and a validated match to the cited paper. The analysis below does not perform that linkage.

\section{Worked example}\label{patent-example}

This executable example explores a 500-article scholarly-citation sample. It is preparation for a possible future linkage analysis, not a substitute for patent data. A high scholarly citation count does not identify a patent-cited paper. No patent-citation prevalence, rate, or technological impact is estimated here.

\subsection{Preparing a scholarly-citation sample}\label{preparing-a-scholarly-citation-sample}

\begin{Shaded}
\begin{Highlighting}[]
\NormalTok{works }\OtherTok{\textless{}{-}} \FunctionTok{fetch\_openalex}\NormalTok{(}
  \AttributeTok{entity =} \StringTok{"works"}\NormalTok{,}
  \AttributeTok{from\_publication\_date =} \StringTok{"2015{-}01{-}01"}\NormalTok{,}
  \AttributeTok{to\_publication\_date =} \StringTok{"2020{-}12{-}31"}\NormalTok{,}
  \AttributeTok{type =} \StringTok{"article"}\NormalTok{,}
  \AttributeTok{options =} \FunctionTok{list}\NormalTok{(}\AttributeTok{sample =} \DecValTok{500}\NormalTok{, }\AttributeTok{seed =} \DecValTok{42}\NormalTok{)}
\NormalTok{)}

\NormalTok{works\_citations }\OtherTok{\textless{}{-}}\NormalTok{ works }\SpecialCharTok{|\textgreater{}}
  \FunctionTok{transmute}\NormalTok{(}
\NormalTok{    id, display\_name, source\_display\_name,}
    \AttributeTok{year =} \FunctionTok{year}\NormalTok{(publication\_date),}
\NormalTok{    cited\_by\_count,}
\NormalTok{    referenced\_works\_count,}
\NormalTok{    type}
\NormalTok{  )}

\FunctionTok{cat}\NormalTok{(}\FunctionTok{glue}\NormalTok{(}\StringTok{"Total works: \{nrow(works\_citations)\}}\SpecialCharTok{\textbackslash{}n}\StringTok{"}\NormalTok{))}
\end{Highlighting}
\end{Shaded}

\begin{verbatim}
#> Total works: 500
\end{verbatim}

\subsection{Scholarly citation distribution}\label{scholarly-citation-distribution}

\begin{Shaded}
\begin{Highlighting}[]
\NormalTok{p90 }\OtherTok{\textless{}{-}} \FunctionTok{quantile}\NormalTok{(works\_citations}\SpecialCharTok{$}\NormalTok{cited\_by\_count, }\FloatTok{0.90}\NormalTok{)}

\FunctionTok{ggplot}\NormalTok{(works\_citations, }\FunctionTok{aes}\NormalTok{(}\AttributeTok{x =}\NormalTok{ cited\_by\_count)) }\SpecialCharTok{+}
  \FunctionTok{geom\_histogram}\NormalTok{(}\AttributeTok{binwidth =} \DecValTok{10}\NormalTok{, }\AttributeTok{fill =} \FunctionTok{palette\_sci}\NormalTok{(}\DecValTok{1}\NormalTok{), }\AttributeTok{colour =} \StringTok{"white"}\NormalTok{) }\SpecialCharTok{+}
  \FunctionTok{geom\_vline}\NormalTok{(}\AttributeTok{xintercept =}\NormalTok{ p90, }\AttributeTok{linetype =} \StringTok{"dashed"}\NormalTok{, }\AttributeTok{colour =} \StringTok{"red"}\NormalTok{) }\SpecialCharTok{+}
  \FunctionTok{annotate}\NormalTok{(}\StringTok{"text"}\NormalTok{, }\AttributeTok{x =}\NormalTok{ p90 }\SpecialCharTok{+} \DecValTok{20}\NormalTok{, }\AttributeTok{y =} \ConstantTok{Inf}\NormalTok{, }\AttributeTok{label =} \StringTok{"90th percentile"}\NormalTok{,}
           \AttributeTok{vjust =} \DecValTok{2}\NormalTok{, }\AttributeTok{colour =} \StringTok{"red"}\NormalTok{, }\AttributeTok{size =} \DecValTok{3}\NormalTok{) }\SpecialCharTok{+}
  \FunctionTok{labs}\NormalTok{(}\AttributeTok{x =} \StringTok{"Scholarly citation count"}\NormalTok{, }\AttributeTok{y =} \StringTok{"Papers"}\NormalTok{) }\SpecialCharTok{+}
  \FunctionTok{theme\_sci}\NormalTok{()}
\end{Highlighting}
\end{Shaded}

\begin{figure}

{\centering \includegraphics[width=0.9\linewidth,alt={Histogram of scholarly citation counts with a vertical line at the sample 90th percentile. The figure contains no patent-citation information.}]{scientometrics-in-r_files/figure-latex/31-fig-patent-citedist-1} 

}

\caption{Scholarly citation counts in the sample, with the sample 90th percentile marked.}\label{fig:31-fig-patent-citedist}
\end{figure}

\subsection{Journals represented among highly cited sampled papers}\label{journals-represented-among-highly-cited-sampled-papers}

\begin{Shaded}
\begin{Highlighting}[]
\NormalTok{works\_citations }\SpecialCharTok{|\textgreater{}}
  \FunctionTok{filter}\NormalTok{(cited\_by\_count }\SpecialCharTok{\textgreater{}=}\NormalTok{ p90, }\SpecialCharTok{!}\FunctionTok{is.na}\NormalTok{(source\_display\_name)) }\SpecialCharTok{|\textgreater{}}
  \FunctionTok{count}\NormalTok{(source\_display\_name, }\AttributeTok{sort =} \ConstantTok{TRUE}\NormalTok{) }\SpecialCharTok{|\textgreater{}}
  \FunctionTok{head}\NormalTok{(}\DecValTok{15}\NormalTok{) }\SpecialCharTok{|\textgreater{}}
  \FunctionTok{mutate}\NormalTok{(}\AttributeTok{source\_display\_name =} \FunctionTok{fct\_reorder}\NormalTok{(source\_display\_name, n)) }\SpecialCharTok{|\textgreater{}}
  \FunctionTok{ggplot}\NormalTok{(}\FunctionTok{aes}\NormalTok{(}\AttributeTok{x =}\NormalTok{ n, }\AttributeTok{y =}\NormalTok{ source\_display\_name)) }\SpecialCharTok{+}
  \FunctionTok{geom\_col}\NormalTok{(}\AttributeTok{fill =} \FunctionTok{palette\_sci}\NormalTok{(}\DecValTok{1}\NormalTok{)) }\SpecialCharTok{+}
  \FunctionTok{labs}\NormalTok{(}\AttributeTok{x =} \StringTok{"Highly cited sampled papers"}\NormalTok{, }\AttributeTok{y =} \ConstantTok{NULL}\NormalTok{) }\SpecialCharTok{+}
  \FunctionTok{theme\_sci}\NormalTok{()}
\end{Highlighting}
\end{Shaded}

\begin{figure}

{\centering \includegraphics[width=0.9\linewidth,alt={Horizontal bar chart counting sampled papers at or above the scholarly-citation threshold, grouped by journal. This is not a ranking of patent-cited research.}]{scientometrics-in-r_files/figure-latex/31-fig-patent-journals-1} 

}

\caption{Up to 15 journals with the most sampled papers at or above the sample scholarly-citation 90th percentile.}\label{fig:31-fig-patent-journals}
\end{figure}

\subsection{Descriptive characteristics by scholarly-citation group}\label{descriptive-characteristics-by-scholarly-citation-group}

\begin{Shaded}
\begin{Highlighting}[]
\NormalTok{works\_citations }\SpecialCharTok{|\textgreater{}}
  \FunctionTok{mutate}\NormalTok{(}\AttributeTok{high\_cite =}\NormalTok{ cited\_by\_count }\SpecialCharTok{\textgreater{}=}\NormalTok{ p90) }\SpecialCharTok{|\textgreater{}}
  \FunctionTok{group\_by}\NormalTok{(high\_cite) }\SpecialCharTok{|\textgreater{}}
  \FunctionTok{summarise}\NormalTok{(}
    \AttributeTok{n =} \FunctionTok{n}\NormalTok{(),}
    \AttributeTok{mean\_refs =} \FunctionTok{round}\NormalTok{(}\FunctionTok{mean}\NormalTok{(referenced\_works\_count, }\AttributeTok{na.rm =} \ConstantTok{TRUE}\NormalTok{), }\DecValTok{1}\NormalTok{),}
    \AttributeTok{mean\_year =} \FunctionTok{round}\NormalTok{(}\FunctionTok{mean}\NormalTok{(year), }\DecValTok{1}\NormalTok{),}
    \AttributeTok{.groups =} \StringTok{"drop"}
\NormalTok{  ) }\SpecialCharTok{|\textgreater{}}
  \FunctionTok{gt}\NormalTok{()}
\end{Highlighting}
\end{Shaded}

\begingroup
\fontsize{8.0pt}{9.0pt}\selectfont
\begin{longtable}{>{\centering\arraybackslash}p{\dimexpr 0.25\linewidth -2\tabcolsep-1.5\arrayrulewidth}>{\raggedleft\arraybackslash}p{\dimexpr 0.24\linewidth -2\tabcolsep-1.5\arrayrulewidth}>{\raggedleft\arraybackslash}p{\dimexpr 0.25\linewidth -2\tabcolsep-1.5\arrayrulewidth}>{\raggedleft\arraybackslash}p{\dimexpr 0.25\linewidth -2\tabcolsep-1.5\arrayrulewidth}}
\toprule
\nolinkurl{high_cite} & n & \nolinkurl{mean_refs} & \nolinkurl{mean_year} \\ 
\midrule\addlinespace[2.5pt]
FALSE & 450 & 10.1 & 2018 \\ 
TRUE & 50 & 54.8 & 2018 \\ 
\bottomrule
\end{longtable}
\endgroup

\section{Diagnostics and interpretation}\label{patent-diagnostics}

\begin{itemize}
\tightlist
\item
  \textbf{Measured quantity}: \texttt{cited\_by\_count} in this analysis counts scholarly citations. Neither the threshold nor the journal plot measures patent citations.
\item
  \textbf{Sample and threshold}: Counts refer only to the retrieved sample, not total journal output. Ties at the 90th percentile can place more than 10\% of papers in the high-citation group.
\item
  \textbf{Citation window}: Current cumulative counts give older publication cohorts more time to receive citations. They are not a fixed-window impact measure.
\item
  \textbf{Field normalisation}: Biomedical papers are far more likely to be cited in patents than social science papers. Compare patent citation rates within fields, not across them.
\item
  \textbf{Time lag}: Estimating a paper-to-patent citation delay requires both dates from actual linked records and a stated observation cutoff; current scholarly citation counts cannot supply it.
\item
  \textbf{Examiner vs.~inventor citations}: Patent examiners add citations that inventors may not have used. This inflates apparent science--technology linkage.
\end{itemize}

\section{Limitations and responsible use}\label{patent-limitations}

\section{Limitations and responsible use}\label{limitations-and-responsible-use-30}

\begin{itemize}
\tightlist
\item
  \textbf{Patent citations are a narrow measure of impact.} Much research influences technology without being cited in patents --- through consulting, training, open-source software, or informal knowledge transfer.
\item
  \textbf{Field coverage differs.} Patent citations capture only some forms of influence. Comparing fields without considering these differences can misrepresent their contributions \citep{hicks2015}.
\item
  \textbf{Geographic bias.} Patent systems vary by country. US and European patents have different citation practices. Analyses based on one patent office may not generalise.
\item
  \textbf{Not evaluative for individuals.} A researcher whose paper is cited in a patent may have had no involvement with the technology. Patent citation reflects relevance of the knowledge, not the researcher's intent to innovate.
\end{itemize}

\section{Common pitfalls}\label{patent-pitfalls}

\section{Common pitfalls}\label{common-pitfalls-30}

\begin{itemize}
\tightlist
\item
  \textbf{Equating patent citations with commercialisation.} A patent citation means the paper is prior art, not that the research led to a product.
\item
  \textbf{Ignoring time lags.} Comparing patent citation rates for papers published in 2023 vs.~2010 is meaningless --- the 2023 papers have not had time to be cited in patents.
\item
  \textbf{Double-counting across patent families.} The same invention may be patented in multiple countries. Count at the patent family level, not the individual patent level.
\item
  \textbf{Substituting scholarly citations for patent linkages.} High scholarly citation counts do not establish that a paper was cited by a patent. Require an observed patent--paper reference before classifying a paper as patent-cited.
\end{itemize}

\section{Exercises}\label{patent-exercises}

\begin{enumerate}
\def\labelenumi{\arabic{enumi}.}
\item
  \textbf{Scholarly-citation comparison.} Compare citation distributions in two journal samples, reporting sample sizes, field differences, publication years, retrieval date and threshold choice. Explain why the result does not estimate patent-citation potential.
\item
  \textbf{Linkage data specification.} Design a table containing a paper identifier, citing patent identifier, patent-family identifier, citation date, source and match status. Describe how you would verify ambiguous references and retain unmatched records. This exercise requires actual patent-reference data before any linkage rates can be calculated.
\item
  \textbf{Institutional patent linkage.} Given an independently validated patent--paper dataset, define the publication denominator, observation window and patent-family deduplication rule for two institutions. Report coverage and uncertainty. If those data are unavailable, identify the missing inputs rather than using scholarly citation percentiles as a proxy.
\end{enumerate}

\section{Solutions}\label{31-solutions}

Solutions are provided in \ref{solutions}.

\section{Further reading}\label{patent-reading}

\begin{itemize}
\tightlist
\item
  \citet{narin1997} --- Foundational analysis of the linkage between science and technology through patent citations.
\item
  \citet{hicks2015} --- Responsible metrics: why patent citations should not be the sole indicator of research impact.
\item
  \citet{priem2022} --- OpenAlex scholarly-work metadata; this example does not establish patent-citation coverage.
\end{itemize}

\section{Session info}\label{patent-sessioninfo}

\begin{Shaded}
\begin{Highlighting}[]
\FunctionTok{sessionInfo}\NormalTok{()}
\end{Highlighting}
\end{Shaded}

\begin{verbatim}
#> R version 4.6.1 (2026-06-24)
#> Platform: x86_64-pc-linux-gnu
#> Running under: Ubuntu 24.04.5 LTS
#> 
#> Matrix products: default
#> BLAS:   /usr/lib/x86_64-linux-gnu/openblas-pthread/libblas.so.3 
#> LAPACK: /usr/lib/x86_64-linux-gnu/openblas-pthread/libopenblasp-r0.3.26.so;  LAPACK version 3.12.0
#> 
#> locale:
#>  [1] LC_CTYPE=C.UTF-8       LC_NUMERIC=C           LC_TIME=C.UTF-8       
#>  [4] LC_COLLATE=C.UTF-8     LC_MONETARY=C.UTF-8    LC_MESSAGES=C.UTF-8   
#>  [7] LC_PAPER=C.UTF-8       LC_NAME=C              LC_ADDRESS=C          
#> [10] LC_TELEPHONE=C         LC_MEASUREMENT=C.UTF-8 LC_IDENTIFICATION=C   
#> 
#> time zone: UTC
#> tzcode source: system (glibc)
#> 
#> attached base packages:
#> [1] stats     graphics  grDevices utils     datasets  methods   base     
#> 
#> other attached packages:
#>  [1] uwot_0.2.5                Matrix_1.7-5             
#>  [3] word2vec_0.4.1            stm_1.3.8                
#>  [5] topicmodels_0.2-17        quanteda.textstats_0.97.2
#>  [7] visNetwork_2.1.4          ggraph_2.2.2             
#>  [9] tidygraph_1.3.1           igraph_2.3.4             
#> [11] quanteda_4.5.1            pdftools_3.9.1           
#> [13] arrow_26.0.0              bibliometrix_5.5.0       
#> [15] RefManageR_1.4.0          bib2df_1.1.2.0           
#> [17] rcrossref_1.2.1           gt_1.4.0                 
#> [19] tidytext_0.4.3            glue_1.8.1               
#> [21] openalexR_3.1.2           lubridate_1.9.5          
#> [23] forcats_1.0.1             stringr_1.6.0            
#> [25] dplyr_1.2.1               purrr_1.2.2              
#> [27] readr_2.2.0               tidyr_1.3.2              
#> [29] tibble_3.3.1              ggplot2_4.0.3            
#> [31] tidyverse_2.0.0          
#> 
#> loaded via a namespace (and not attached):
#>   [1] bibtex_0.5.3           RColorBrewer_1.1-3     rstudioapi_0.19.0     
#>   [4] jsonlite_2.0.0         magrittr_2.0.5         modeltools_0.2-25     
#>   [7] farver_2.1.2           rmarkdown_2.32         fs_2.1.0              
#>  [10] vctrs_0.7.3            memoise_2.0.1          askpass_1.2.1         
#>  [13] base64enc_0.1-6        htmltools_0.5.9        contentanalysis_1.1.1 
#>  [16] curl_8.1.0             janeaustenr_1.0.0      cellranger_1.1.0      
#>  [19] htmlwidgets_1.6.4      tokenizers_0.3.0       plyr_1.8.9            
#>  [22] httr2_1.3.0            cachem_1.1.0           plotly_4.12.1         
#>  [25] dimensionsR_0.0.3      mime_0.13              lifecycle_1.0.5       
#>  [28] pkgconfig_2.0.3        R6_2.6.1               fastmap_1.2.0         
#>  [31] shiny_1.14.0           digest_0.6.39          patchwork_1.3.2       
#>  [34] shinycssloaders_1.1.0  rprojroot_2.1.1        RSpectra_0.16-2       
#>  [37] SnowballC_0.7.1        labeling_0.4.3         urltools_1.7.3.1      
#>  [40] timechange_0.4.0       mgcv_1.9-4             httr_1.4.9            
#>  [43] polyclip_1.10-7        compiler_4.6.1         here_1.0.2            
#>  [46] bit64_4.8.6            withr_3.0.3            S7_0.2.2              
#>  [49] viridis_0.6.5          ggforce_0.5.0          MASS_7.3-65           
#>  [52] bibliometrixData_0.3.0 tools_4.6.1            otel_0.2.0            
#>  [55] stopwords_2.3          zip_3.0.2              rentrez_1.2.4         
#>  [58] httpuv_1.6.17          nlme_3.1-169           promises_1.5.0        
#>  [61] grid_4.6.1             stringdist_0.9.17      reshape2_1.4.5        
#>  [64] generics_0.1.4         gtable_0.3.6           tzdb_0.5.0            
#>  [67] rscopus_0.9.0          ca_0.72                data.table_1.18.6.1   
#>  [70] hms_1.1.4              xml2_1.6.0             utf8_1.2.6            
#>  [73] ggrepel_0.9.8          pillar_1.11.1          nsyllable_1.0.1       
#>  [76] vroom_1.7.1            later_1.4.8            splines_4.6.1         
#>  [79] tweenr_2.0.3           lattice_0.22-9         FNN_1.1.4.1           
#>  [82] bit_4.6.0              tidyselect_1.2.1       tm_0.7-20             
#>  [85] miniUI_0.1.2           knitr_1.52             gridExtra_2.3.1       
#>  [88] NLP_0.3-3              bookdown_0.48          stats4_4.6.1          
#>  [91] crul_1.6.0             xfun_0.62              graphlayouts_1.2.5    
#>  [94] matrixStats_1.5.0      DT_0.34.0              humaniformat_0.6.0    
#>  [97] stringi_1.8.9          qpdf_1.4.1             yaml_2.3.12           
#> [100] evaluate_1.0.5         httpcode_0.3.0         cli_3.6.6             
#> [103] xtable_1.8-8           systemfonts_1.3.2      Rcpp_1.1.2            
#> [106] readxl_1.6.0           triebeard_0.4.1        XML_3.99-0.25         
#> [109] parallel_4.6.1         assertthat_0.2.1       pubmedR_1.0.2         
#> [112] slam_0.1-56            viridisLite_0.4.3      scales_1.4.0          
#> [115] crayon_1.5.3           openxlsx_4.2.9         rlang_1.3.0           
#> [118] fastmatch_1.1-8
\end{verbatim}

\chapter{Causal Inference in Scientometrics}\label{causal-inference-in-scientometrics}

\emph{R. Heller · Date: 2026-05-09 · Last changed: 2026-10-31}

\section{Learning objectives}\label{causal-objectives}

After completing this chapter, you will be able to:

\begin{itemize}
\tightlist
\item
  Distinguish causal from correlational claims in bibliometric research
\item
  Explain the intervention and outcome data required for difference-in-differences
\item
  Compute descriptive OA-by-publication-cohort comparisons without assigning a causal interpretation
\item
  Distinguish a within-year comparison from validated matching or a regression-discontinuity design
\item
  Recognise the assumptions underlying causal inference in observational bibliometric data
\end{itemize}

\section{Setup}\label{causal-setup}

\begin{Shaded}
\begin{Highlighting}[]
\FunctionTok{library}\NormalTok{(tidyverse)}
\FunctionTok{library}\NormalTok{(openalexR)}
\FunctionTok{library}\NormalTok{(glue)}
\FunctionTok{library}\NormalTok{(gt)}

\FunctionTok{set.seed}\NormalTok{(}\DecValTok{20260509}\NormalTok{)}

\FunctionTok{source}\NormalTok{(here}\SpecialCharTok{::}\FunctionTok{here}\NormalTok{(}\StringTok{"R"}\NormalTok{, }\StringTok{"api\_helpers.R"}\NormalTok{))}
\FunctionTok{source}\NormalTok{(here}\SpecialCharTok{::}\FunctionTok{here}\NormalTok{(}\StringTok{"R"}\NormalTok{, }\StringTok{"utils.R"}\NormalTok{))}
\FunctionTok{source}\NormalTok{(here}\SpecialCharTok{::}\FunctionTok{here}\NormalTok{(}\StringTok{"R"}\NormalTok{, }\StringTok{"sci\_palette.R"}\NormalTok{))}
\end{Highlighting}
\end{Shaded}

\section{Conceptual background}\label{causal-background}

Bibliometrics is rich in correlations: OA papers are more cited than closed-access papers; funded research receives more citations; papers in prestigious journals are cited more. But correlation does not imply causation. The OA citation advantage may reflect self-selection (better papers are more likely to be made OA), not the causal effect of open access. Separating cause from confound is essential for evidence-based science policy.

\textbf{Difference-in-differences (DiD)} exploits policy changes as natural experiments. If a funder mandates OA starting in year T, papers from that funder before and after T serve as the treatment group, while papers from a similar funder without the mandate serve as the control. The causal effect is estimated as the difference in trends between treatment and control groups. The key assumption is \emph{parallel trends}: absent the intervention, both groups would have followed the same trajectory.

\textbf{Regression discontinuity design (RDD)} applies when treatment is assigned by a threshold. For example, grants are funded above a score cutoff; papers are accepted above a review score threshold. Comparing outcomes just above and just below the threshold isolates the causal effect of the treatment, under the assumption that units near the threshold are comparable.

\textbf{Matching methods} create comparable groups from observational data. Propensity score matching pairs treated units (e.g., OA papers) with control units (closed-access papers) that have similar observable characteristics (field, journal, year, team size). This reduces confounding but cannot eliminate unmeasured confounders.

All three methods require strong assumptions. In bibliometrics, violations are common: the parallel trends assumption may fail because of concurrent policy changes; the RDD threshold may be manipulable; matching cannot address unmeasured confounders like paper quality. \citet{hicks2015} emphasise that quantitative analysis, however sophisticated, cannot replace expert judgement in research evaluation.

\section{Worked example}\label{causal-example}

\subsection{A descriptive OA-by-publication-cohort comparison}\label{a-descriptive-oa-by-publication-cohort-comparison}

This example samples papers from one journal and compares their current cumulative citation counts. The 2020 split is an arbitrary publication-cohort boundary, not a documented intervention. OA status is observed at retrieval, not recorded treatment assignment. Different publication years contain different papers with different citation exposure. Consequently, neither the interaction regression nor the cohort plot identifies a policy effect or tests a DiD parallel-trends assumption. DiD requires a defensible intervention, comparison design and outcome histories; see \href{https://doi.org/10.1016/j.jeconom.2020.12.001}{Callaway and Sant'Anna (2021)}.

The sample and citation snapshot describe the fetched data only; they are not journal totals or a frozen historical citation panel.

\begin{Shaded}
\begin{Highlighting}[]
\NormalTok{works }\OtherTok{\textless{}{-}} \FunctionTok{fetch\_openalex}\NormalTok{(}
  \AttributeTok{entity =} \StringTok{"works"}\NormalTok{,}
  \AttributeTok{primary\_location.source.id =} \StringTok{"S148561398"}\NormalTok{,}
  \AttributeTok{from\_publication\_date =} \StringTok{"2016{-}01{-}01"}\NormalTok{,}
  \AttributeTok{to\_publication\_date =} \StringTok{"2023{-}12{-}31"}\NormalTok{,}
  \AttributeTok{type =} \StringTok{"article"}\NormalTok{,}
  \AttributeTok{options =} \FunctionTok{list}\NormalTok{(}\AttributeTok{sample =} \DecValTok{500}\NormalTok{, }\AttributeTok{seed =} \DecValTok{42}\NormalTok{)}
\NormalTok{)}

\NormalTok{cohort\_data }\OtherTok{\textless{}{-}}\NormalTok{ works }\SpecialCharTok{|\textgreater{}}
  \FunctionTok{transmute}\NormalTok{(}
\NormalTok{    id, }\AttributeTok{year =} \FunctionTok{year}\NormalTok{(publication\_date),}
\NormalTok{    cited\_by\_count,}
    \AttributeTok{is\_oa =} \FunctionTok{case\_when}\NormalTok{(}
\NormalTok{      oa\_status }\SpecialCharTok{\%in\%} \FunctionTok{c}\NormalTok{(}\StringTok{"gold"}\NormalTok{, }\StringTok{"green"}\NormalTok{, }\StringTok{"hybrid"}\NormalTok{, }\StringTok{"bronze"}\NormalTok{, }\StringTok{"diamond"}\NormalTok{) }\SpecialCharTok{\textasciitilde{}} \ConstantTok{TRUE}\NormalTok{,}
\NormalTok{      oa\_status }\SpecialCharTok{==} \StringTok{"closed"} \SpecialCharTok{\textasciitilde{}} \ConstantTok{FALSE}\NormalTok{,}
      \ConstantTok{TRUE} \SpecialCharTok{\textasciitilde{}} \ConstantTok{NA}
\NormalTok{    ),}
    \AttributeTok{recent\_cohort =}\NormalTok{ year }\SpecialCharTok{\textgreater{}=} \DecValTok{2020}\NormalTok{,}
\NormalTok{    referenced\_works\_count}
\NormalTok{  )}

\NormalTok{excluded\_records }\OtherTok{\textless{}{-}} \FunctionTok{sum}\NormalTok{(}\SpecialCharTok{!}\FunctionTok{complete.cases}\NormalTok{(cohort\_data))}
\NormalTok{cohort\_data }\OtherTok{\textless{}{-}}\NormalTok{ cohort\_data }\SpecialCharTok{|\textgreater{}} \FunctionTok{drop\_na}\NormalTok{()}
\FunctionTok{stopifnot}\NormalTok{(}\SpecialCharTok{!}\FunctionTok{anyDuplicated}\NormalTok{(cohort\_data}\SpecialCharTok{$}\NormalTok{id), }\FunctionTok{nrow}\NormalTok{(cohort\_data) }\SpecialCharTok{\textgreater{}} \DecValTok{0}\NormalTok{)}
\FunctionTok{cat}\NormalTok{(}\FunctionTok{glue}\NormalTok{(}\StringTok{"Excluded incomplete records: \{excluded\_records\}}\SpecialCharTok{\textbackslash{}n}\StringTok{"}\NormalTok{))}
\end{Highlighting}
\end{Shaded}

\begin{verbatim}
#> Excluded incomplete records: 0
\end{verbatim}

\begin{Shaded}
\begin{Highlighting}[]
\FunctionTok{cat}\NormalTok{(}\FunctionTok{glue}\NormalTok{(}\StringTok{"Sampled complete records: \{nrow(cohort\_data)\}}\SpecialCharTok{\textbackslash{}n}\StringTok{"}\NormalTok{))}
\end{Highlighting}
\end{Shaded}

\begin{verbatim}
#> Sampled complete records: 500
\end{verbatim}

\begin{Shaded}
\begin{Highlighting}[]
\FunctionTok{cat}\NormalTok{(}\FunctionTok{glue}\NormalTok{(}\StringTok{"OA: \{sum(cohort\_data$is\_oa)\}, Closed: \{sum(!cohort\_data$is\_oa)\}}\SpecialCharTok{\textbackslash{}n}\StringTok{"}\NormalTok{))}
\end{Highlighting}
\end{Shaded}

\begin{verbatim}
#> OA: 228, Closed: 272
\end{verbatim}

\subsection{Descriptive comparison}\label{descriptive-comparison}

\begin{Shaded}
\begin{Highlighting}[]
\NormalTok{cohort\_summary }\OtherTok{\textless{}{-}}\NormalTok{ cohort\_data }\SpecialCharTok{|\textgreater{}}
  \FunctionTok{group\_by}\NormalTok{(is\_oa, recent\_cohort) }\SpecialCharTok{|\textgreater{}}
  \FunctionTok{summarise}\NormalTok{(}\AttributeTok{mean\_cites =} \FunctionTok{mean}\NormalTok{(cited\_by\_count), }\AttributeTok{n =} \FunctionTok{n}\NormalTok{(), }\AttributeTok{.groups =} \StringTok{"drop"}\NormalTok{) }\SpecialCharTok{|\textgreater{}}
  \FunctionTok{mutate}\NormalTok{(}\AttributeTok{period =} \FunctionTok{ifelse}\NormalTok{(recent\_cohort, }\StringTok{"2020{-}2023"}\NormalTok{, }\StringTok{"2016{-}2019"}\NormalTok{),}
         \AttributeTok{group =} \FunctionTok{ifelse}\NormalTok{(is\_oa, }\StringTok{"Open access"}\NormalTok{, }\StringTok{"Closed"}\NormalTok{))}

\FunctionTok{ggplot}\NormalTok{(cohort\_summary, }\FunctionTok{aes}\NormalTok{(}\AttributeTok{x =}\NormalTok{ period, }\AttributeTok{y =}\NormalTok{ mean\_cites, }\AttributeTok{fill =}\NormalTok{ group)) }\SpecialCharTok{+}
  \FunctionTok{geom\_col}\NormalTok{(}\AttributeTok{position =} \StringTok{"dodge"}\NormalTok{) }\SpecialCharTok{+}
  \FunctionTok{scale\_fill\_manual}\NormalTok{(}\AttributeTok{values =} \FunctionTok{palette\_sci}\NormalTok{(}\DecValTok{2}\NormalTok{)) }\SpecialCharTok{+}
  \FunctionTok{labs}\NormalTok{(}\AttributeTok{x =} \StringTok{"Publication cohort"}\NormalTok{, }\AttributeTok{y =} \StringTok{"Mean cumulative citations at retrieval"}\NormalTok{, }\AttributeTok{fill =} \StringTok{"Access"}\NormalTok{) }\SpecialCharTok{+}
  \FunctionTok{theme\_sci}\NormalTok{()}
\end{Highlighting}
\end{Shaded}

\begin{figure}

{\centering \includegraphics[width=0.9\linewidth,alt={Grouped bar chart of current cumulative citations for OA and closed-access papers published in 2016–2019 or 2020–2023.}]{scientometrics-in-r_files/figure-latex/32-fig-causal-descriptive-1} 

}

\caption{Current cumulative citations in sampled papers, grouped by OA status and publication cohort; not a policy-effect estimate.}\label{fig:32-fig-causal-descriptive}
\end{figure}

\subsection{Descriptive interaction regression}\label{descriptive-interaction-regression}

The interaction describes how the conditional OA/closed citation difference varies between the two sampled publication cohorts. Adding reference count is statistical adjustment, not an identification strategy. The fitted coefficients and conventional linear-model uncertainty summaries must not be reported as causal effects.

\begin{Shaded}
\begin{Highlighting}[]
\NormalTok{interaction\_model }\OtherTok{\textless{}{-}} \FunctionTok{lm}\NormalTok{(cited\_by\_count }\SpecialCharTok{\textasciitilde{}}\NormalTok{ is\_oa }\SpecialCharTok{*}\NormalTok{ recent\_cohort }\SpecialCharTok{+}\NormalTok{ referenced\_works\_count,}
                \AttributeTok{data =}\NormalTok{ cohort\_data)}

\NormalTok{tidy\_interaction }\OtherTok{\textless{}{-}}\NormalTok{ broom}\SpecialCharTok{::}\FunctionTok{tidy}\NormalTok{(interaction\_model)}
\NormalTok{tidy\_interaction }\SpecialCharTok{|\textgreater{}}
  \FunctionTok{mutate}\NormalTok{(}\FunctionTok{across}\NormalTok{(}\FunctionTok{where}\NormalTok{(is.numeric), \textbackslash{}(x) }\FunctionTok{round}\NormalTok{(x, }\DecValTok{3}\NormalTok{))) }\SpecialCharTok{|\textgreater{}}
  \FunctionTok{gt}\NormalTok{()}
\end{Highlighting}
\end{Shaded}

\begingroup
\fontsize{8.0pt}{9.0pt}\selectfont
\begin{longtable}{>{\raggedright\arraybackslash}p{\dimexpr 0.31\linewidth -2\tabcolsep-1.5\arrayrulewidth}>{\raggedleft\arraybackslash}p{\dimexpr 0.17\linewidth -2\tabcolsep-1.5\arrayrulewidth}>{\raggedleft\arraybackslash}p{\dimexpr 0.18\linewidth -2\tabcolsep-1.5\arrayrulewidth}>{\raggedleft\arraybackslash}p{\dimexpr 0.18\linewidth -2\tabcolsep-1.5\arrayrulewidth}>{\raggedleft\arraybackslash}p{\dimexpr 0.17\linewidth -2\tabcolsep-1.5\arrayrulewidth}}
\toprule
term & \nolinkurl{estimate} & \nolinkurl{std.error} & \nolinkurl{statistic} & p.value \\ 
\midrule\addlinespace[2.5pt]
(Intercept) & 24.013 & 5.704 & 4.21 & 0.000 \\ 
\nolinkurl{is_oaTRUE} & 13.196 & 7.324 & 1.80 & 0.072 \\ 
\nolinkurl{recent_cohortTRUE} & -10.102 & 6.869 & -1.47 & 0.142 \\ 
\nolinkurl{referenced_works_count} & 0.205 & 0.074 & 2.75 & 0.006 \\ 
\nolinkurl{is_oaTRUE:recent_cohortTRUE} & -10.422 & 10.185 & -1.02 & 0.307 \\ 
\bottomrule
\end{longtable}
\endgroup

\subsection{Within-year sampling comparison}\label{within-year-sampling-comparison}

For each sampled OA paper, select one closed-access paper from the same publication year. Controls can be reused, and OA papers in years with no closed-access candidate are omitted. This is sampling with replacement within year, not propensity-score matching; it does not establish covariate balance or remove confounding.

\begin{Shaded}
\begin{Highlighting}[]
\NormalTok{oa\_papers }\OtherTok{\textless{}{-}}\NormalTok{ cohort\_data }\SpecialCharTok{|\textgreater{}} \FunctionTok{filter}\NormalTok{(is\_oa)}
\NormalTok{closed\_papers }\OtherTok{\textless{}{-}}\NormalTok{ cohort\_data }\SpecialCharTok{|\textgreater{}} \FunctionTok{filter}\NormalTok{(}\SpecialCharTok{!}\NormalTok{is\_oa)}

\NormalTok{matched }\OtherTok{\textless{}{-}}\NormalTok{ oa\_papers }\SpecialCharTok{|\textgreater{}}
  \FunctionTok{mutate}\NormalTok{(}\AttributeTok{match\_year =}\NormalTok{ year) }\SpecialCharTok{|\textgreater{}}
  \FunctionTok{inner\_join}\NormalTok{(}
\NormalTok{    closed\_papers }\SpecialCharTok{|\textgreater{}} \FunctionTok{select}\NormalTok{(year, }\AttributeTok{control\_id =}\NormalTok{ id, }\AttributeTok{cited\_by\_count\_closed =}\NormalTok{ cited\_by\_count),}
    \AttributeTok{by =} \StringTok{"year"}\NormalTok{,}
    \AttributeTok{relationship =} \StringTok{"many{-}to{-}many"}
\NormalTok{  ) }\SpecialCharTok{|\textgreater{}}
  \FunctionTok{group\_by}\NormalTok{(id) }\SpecialCharTok{|\textgreater{}}
  \FunctionTok{slice\_sample}\NormalTok{(}\AttributeTok{n =} \DecValTok{1}\NormalTok{) }\SpecialCharTok{|\textgreater{}}
  \FunctionTok{ungroup}\NormalTok{()}

\FunctionTok{cat}\NormalTok{(}\FunctionTok{glue}\NormalTok{(}\StringTok{"Within{-}year pairs: \{nrow(matched)\} / \{nrow(oa\_papers)\} OA papers}\SpecialCharTok{\textbackslash{}n}\StringTok{"}\NormalTok{))}
\end{Highlighting}
\end{Shaded}

\begin{verbatim}
#> Within-year pairs: 228 / 228 OA papers
\end{verbatim}

\begin{Shaded}
\begin{Highlighting}[]
\FunctionTok{cat}\NormalTok{(}\FunctionTok{glue}\NormalTok{(}\StringTok{"Distinct closed{-}access controls: \{n\_distinct(matched$control\_id)\}}\SpecialCharTok{\textbackslash{}n}\StringTok{"}\NormalTok{))}
\end{Highlighting}
\end{Shaded}

\begin{verbatim}
#> Distinct closed-access controls: 145
\end{verbatim}

\begin{Shaded}
\begin{Highlighting}[]
\FunctionTok{cat}\NormalTok{(}\FunctionTok{glue}\NormalTok{(}\StringTok{"Mean cites (OA): \{round(mean(matched$cited\_by\_count), 1)\}}\SpecialCharTok{\textbackslash{}n}\StringTok{"}\NormalTok{))}
\end{Highlighting}
\end{Shaded}

\begin{verbatim}
#> Mean cites (OA): 35.2
\end{verbatim}

\begin{Shaded}
\begin{Highlighting}[]
\FunctionTok{cat}\NormalTok{(}\FunctionTok{glue}\NormalTok{(}\StringTok{"Mean cites (closed match): \{round(mean(matched$cited\_by\_count\_closed), 1)\}}\SpecialCharTok{\textbackslash{}n}\StringTok{"}\NormalTok{))}
\end{Highlighting}
\end{Shaded}

\begin{verbatim}
#> Mean cites (closed match): 24.4
\end{verbatim}

\subsection{Publication-cohort profile}\label{publication-cohort-profile}

The horizontal axis is publication year. Each point summarises different papers' cumulative citations at one retrieval time. These are not citations received annually by fixed treated and control groups.

\begin{Shaded}
\begin{Highlighting}[]
\NormalTok{trends }\OtherTok{\textless{}{-}}\NormalTok{ cohort\_data }\SpecialCharTok{|\textgreater{}}
  \FunctionTok{group\_by}\NormalTok{(year, is\_oa) }\SpecialCharTok{|\textgreater{}}
  \FunctionTok{summarise}\NormalTok{(}\AttributeTok{mean\_cites =} \FunctionTok{mean}\NormalTok{(cited\_by\_count), }\AttributeTok{.groups =} \StringTok{"drop"}\NormalTok{) }\SpecialCharTok{|\textgreater{}}
  \FunctionTok{mutate}\NormalTok{(}\AttributeTok{group =} \FunctionTok{ifelse}\NormalTok{(is\_oa, }\StringTok{"OA"}\NormalTok{, }\StringTok{"Closed"}\NormalTok{))}

\FunctionTok{ggplot}\NormalTok{(trends, }\FunctionTok{aes}\NormalTok{(}\AttributeTok{x =}\NormalTok{ year, }\AttributeTok{y =}\NormalTok{ mean\_cites, }\AttributeTok{colour =}\NormalTok{ group)) }\SpecialCharTok{+}
  \FunctionTok{geom\_line}\NormalTok{(}\AttributeTok{linewidth =} \DecValTok{1}\NormalTok{) }\SpecialCharTok{+}
  \FunctionTok{geom\_point}\NormalTok{(}\AttributeTok{size =} \DecValTok{2}\NormalTok{) }\SpecialCharTok{+}
  \FunctionTok{geom\_vline}\NormalTok{(}\AttributeTok{xintercept =} \FloatTok{2019.5}\NormalTok{, }\AttributeTok{linetype =} \StringTok{"dashed"}\NormalTok{, }\AttributeTok{colour =} \StringTok{"grey50"}\NormalTok{) }\SpecialCharTok{+}
  \FunctionTok{scale\_colour\_manual}\NormalTok{(}\AttributeTok{values =} \FunctionTok{palette\_sci}\NormalTok{(}\DecValTok{2}\NormalTok{)) }\SpecialCharTok{+}
  \FunctionTok{labs}\NormalTok{(}\AttributeTok{x =} \StringTok{"Publication year"}\NormalTok{, }\AttributeTok{y =} \StringTok{"Mean cumulative citations at retrieval"}\NormalTok{, }\AttributeTok{colour =} \StringTok{"Access"}\NormalTok{) }\SpecialCharTok{+}
  \FunctionTok{theme\_sci}\NormalTok{()}
\end{Highlighting}
\end{Shaded}

\begin{figure}

{\centering \includegraphics[width=0.9\linewidth,alt={Line chart comparing cumulative citation averages across publication cohorts for OA and closed-access papers; older cohorts have had longer to accrue citations.}]{scientometrics-in-r_files/figure-latex/32-fig-causal-trends-1} 

}

\caption{Current cumulative citations by OA status and publication year in the sample. This cohort profile does not test parallel trends.}\label{fig:32-fig-causal-trends}
\end{figure}

\section{Diagnostics and interpretation}\label{causal-diagnostics}

\begin{itemize}
\tightlist
\item
  \textbf{Selection and exposure}: Report the sampling query, retrieval date, exclusions and group sizes. Older publication cohorts have longer citation exposure; snapshot OA status may differ from status at publication.
\item
  \textbf{Interpret the actual model}: The OA-by-cohort coefficient is a descriptive conditional association. The arbitrary 2020 boundary supplies no intervention and no causal estimand.
\item
  \textbf{Assess the comparison}: Check years without eligible controls, control reuse and covariate differences. One sampled same-year control does not make papers exchangeable.
\item
  \textbf{Causal designs need additional evidence}: Define intervention timing and assignment, comparable outcomes over time, the target effect and identifying assumptions before choosing a DiD, RDD or matching analysis. A non-significant pre-trend test alone does not establish parallel counterfactual trends.
\end{itemize}

\section{Limitations and responsible use}\label{causal-limitations}

\section{Limitations and responsible use}\label{limitations-and-responsible-use-31}

\begin{itemize}
\tightlist
\item
  \textbf{This example is descriptive.} It has no documented intervention or identified causal design. Observational causal designs can support causal interpretation under explicit, defensible assumptions; fitting the shown regression does not supply them.
\item
  \textbf{Unmeasured confounders.} Paper quality, author reputation, and marketing effort are rarely observable in bibliometric data but strongly affect citations.
\item
  \textbf{Treatment definition.} ``OA'' is not a single treatment --- it encompasses gold, green, hybrid, and bronze, each with different mechanisms.
\item
  \textbf{Generalisability.} This sampled single-journal comparison does not establish an OA effect in this journal or elsewhere.
\item
  \textbf{Policy implications require judgement.} Even a well-estimated causal effect is one input into policy decisions, not a sufficient basis for mandates \citep{hicks2015, dora2012}.
\end{itemize}

\section{Common pitfalls}\label{causal-pitfalls}

\section{Common pitfalls}\label{common-pitfalls-31}

\begin{itemize}
\tightlist
\item
  \textbf{Claiming causation from regression coefficients.} A positive OA coefficient in a regression does not prove OA \emph{causes} more citations without addressing confounders.
\item
  \textbf{Calling a cohort profile a pre-trend test.} Current cumulative citations in different publication cohorts do not measure treated/control outcome paths around a policy change.
\item
  \textbf{Treating matching as automatic identification.} Choose covariates from the causal question and timing, evaluate overlap/balance, and avoid conditioning indiscriminately on post-treatment variables.
\item
  \textbf{Over-interpreting p-values.} A p \textless{} 0.05 does not mean the result is practically significant or that the causal assumptions hold.
\end{itemize}

\section{Exercises}\label{causal-exercises}

\begin{enumerate}
\def\labelenumi{\arabic{enumi}.}
\item
  \textbf{Interaction sensitivity.} Compare the OA-by-cohort coefficient with and without adjustment for reference count. Explain why a changed coefficient does not itself diagnose a causal effect or quantify omitted-variable bias.
\item
  \textbf{Boundary sensitivity.} Repeat the descriptive interaction using a 2018 publication-cohort boundary. Explain why this is not a placebo intervention test and why a non-significant coefficient cannot establish parallel trends.
\item
  \textbf{Comparison audit.} Count reused closed-access controls and unmatched OA records in the within-year comparison. Specify the intervention, assignment information and longitudinal outcomes needed for an actual policy-effect study.
\end{enumerate}

\section{Solutions}\label{32-solutions}

Solutions are provided in \ref{solutions}.

\section{Further reading}\label{causal-reading}

\begin{itemize}
\tightlist
\item
  \citet{hicks2015} --- The Leiden Manifesto: caution against mechanistic application of quantitative metrics.
\item
  \citet{dora2012} --- DORA: responsible assessment requires understanding causation.
\item
  \citet{waltman2016} --- Review of citation indicators and their interpretation challenges.
\item
  \href{https://doi.org/10.1016/j.jeconom.2020.12.001}{Callaway and Sant'Anna (2021)} --- Identification, estimation and inference for DiD with multiple periods.
\end{itemize}

\section{Session info}\label{causal-sessioninfo}

\begin{Shaded}
\begin{Highlighting}[]
\FunctionTok{sessionInfo}\NormalTok{()}
\end{Highlighting}
\end{Shaded}

\begin{verbatim}
#> R version 4.6.1 (2026-06-24)
#> Platform: x86_64-pc-linux-gnu
#> Running under: Ubuntu 24.04.5 LTS
#> 
#> Matrix products: default
#> BLAS:   /usr/lib/x86_64-linux-gnu/openblas-pthread/libblas.so.3 
#> LAPACK: /usr/lib/x86_64-linux-gnu/openblas-pthread/libopenblasp-r0.3.26.so;  LAPACK version 3.12.0
#> 
#> locale:
#>  [1] LC_CTYPE=C.UTF-8       LC_NUMERIC=C           LC_TIME=C.UTF-8       
#>  [4] LC_COLLATE=C.UTF-8     LC_MONETARY=C.UTF-8    LC_MESSAGES=C.UTF-8   
#>  [7] LC_PAPER=C.UTF-8       LC_NAME=C              LC_ADDRESS=C          
#> [10] LC_TELEPHONE=C         LC_MEASUREMENT=C.UTF-8 LC_IDENTIFICATION=C   
#> 
#> time zone: UTC
#> tzcode source: system (glibc)
#> 
#> attached base packages:
#> [1] stats     graphics  grDevices utils     datasets  methods   base     
#> 
#> other attached packages:
#>  [1] uwot_0.2.5                Matrix_1.7-5             
#>  [3] word2vec_0.4.1            stm_1.3.8                
#>  [5] topicmodels_0.2-17        quanteda.textstats_0.97.2
#>  [7] visNetwork_2.1.4          ggraph_2.2.2             
#>  [9] tidygraph_1.3.1           igraph_2.3.4             
#> [11] quanteda_4.5.1            pdftools_3.9.1           
#> [13] arrow_26.0.0              bibliometrix_5.5.0       
#> [15] RefManageR_1.4.0          bib2df_1.1.2.0           
#> [17] rcrossref_1.2.1           gt_1.4.0                 
#> [19] tidytext_0.4.3            glue_1.8.1               
#> [21] openalexR_3.1.2           lubridate_1.9.5          
#> [23] forcats_1.0.1             stringr_1.6.0            
#> [25] dplyr_1.2.1               purrr_1.2.2              
#> [27] readr_2.2.0               tidyr_1.3.2              
#> [29] tibble_3.3.1              ggplot2_4.0.3            
#> [31] tidyverse_2.0.0          
#> 
#> loaded via a namespace (and not attached):
#>   [1] bibtex_0.5.3           RColorBrewer_1.1-3     rstudioapi_0.19.0     
#>   [4] jsonlite_2.0.0         magrittr_2.0.5         modeltools_0.2-25     
#>   [7] farver_2.1.2           rmarkdown_2.32         fs_2.1.0              
#>  [10] vctrs_0.7.3            memoise_2.0.1          askpass_1.2.1         
#>  [13] base64enc_0.1-6        htmltools_0.5.9        contentanalysis_1.1.1 
#>  [16] curl_8.1.0             broom_1.0.13           janeaustenr_1.0.0     
#>  [19] cellranger_1.1.0       htmlwidgets_1.6.4      tokenizers_0.3.0      
#>  [22] plyr_1.8.9             httr2_1.3.0            cachem_1.1.0          
#>  [25] plotly_4.12.1          dimensionsR_0.0.3      mime_0.13             
#>  [28] lifecycle_1.0.5        pkgconfig_2.0.3        R6_2.6.1              
#>  [31] fastmap_1.2.0          shiny_1.14.0           digest_0.6.39         
#>  [34] patchwork_1.3.2        shinycssloaders_1.1.0  rprojroot_2.1.1       
#>  [37] RSpectra_0.16-2        SnowballC_0.7.1        labeling_0.4.3        
#>  [40] urltools_1.7.3.1       timechange_0.4.0       mgcv_1.9-4            
#>  [43] httr_1.4.9             polyclip_1.10-7        compiler_4.6.1        
#>  [46] here_1.0.2             bit64_4.8.6            withr_3.0.3           
#>  [49] backports_1.5.1        S7_0.2.2               viridis_0.6.5         
#>  [52] ggforce_0.5.0          MASS_7.3-65            bibliometrixData_0.3.0
#>  [55] tools_4.6.1            otel_0.2.0             stopwords_2.3         
#>  [58] zip_3.0.2              rentrez_1.2.4          httpuv_1.6.17         
#>  [61] nlme_3.1-169           promises_1.5.0         grid_4.6.1            
#>  [64] stringdist_0.9.17      reshape2_1.4.5         generics_0.1.4        
#>  [67] gtable_0.3.6           tzdb_0.5.0             rscopus_0.9.0         
#>  [70] ca_0.72                data.table_1.18.6.1    hms_1.1.4             
#>  [73] xml2_1.6.0             utf8_1.2.6             ggrepel_0.9.8         
#>  [76] pillar_1.11.1          nsyllable_1.0.1        vroom_1.7.1           
#>  [79] later_1.4.8            splines_4.6.1          tweenr_2.0.3          
#>  [82] lattice_0.22-9         FNN_1.1.4.1            bit_4.6.0             
#>  [85] tidyselect_1.2.1       tm_0.7-20              miniUI_0.1.2          
#>  [88] knitr_1.52             gridExtra_2.3.1        NLP_0.3-3             
#>  [91] bookdown_0.48          stats4_4.6.1           crul_1.6.0            
#>  [94] xfun_0.62              graphlayouts_1.2.5     matrixStats_1.5.0     
#>  [97] DT_0.34.0              humaniformat_0.6.0     stringi_1.8.9         
#> [100] qpdf_1.4.1             yaml_2.3.12            evaluate_1.0.5        
#> [103] httpcode_0.3.0         cli_3.6.6              xtable_1.8-8          
#> [106] systemfonts_1.3.2      Rcpp_1.1.2             readxl_1.6.0          
#> [109] triebeard_0.4.1        XML_3.99-0.25          parallel_4.6.1        
#> [112] assertthat_0.2.1       pubmedR_1.0.2          slam_0.1-56           
#> [115] viridisLite_0.4.3      scales_1.4.0           crayon_1.5.3          
#> [118] openxlsx_4.2.9         rlang_1.3.0            fastmatch_1.1-8
\end{verbatim}

\chapter{Reproducibility, RRIDs, and Badges}\label{reproducibility-rrids-and-badges}

\emph{R. Heller · Date: 2026-05-09 · Last changed: 2026-10-31}

\section{Learning objectives}\label{repro-objectives}

After completing this chapter, you will be able to:

\begin{itemize}
\tightlist
\item
  Define reproducibility indicators and explain their role in research evaluation
\item
  Search for data availability statements and open badges in bibliometric metadata
\item
  Analyse the prevalence of reproducibility practices across journals and fields
\item
  Discuss the relationship between transparency practices and citation impact
\item
  Recognise the gap between reproducibility signalling and actual reproducibility
\end{itemize}

\section{Setup}\label{repro-setup}

\begin{Shaded}
\begin{Highlighting}[]
\FunctionTok{library}\NormalTok{(tidyverse)}
\FunctionTok{library}\NormalTok{(openalexR)}
\FunctionTok{library}\NormalTok{(glue)}
\FunctionTok{library}\NormalTok{(gt)}

\FunctionTok{set.seed}\NormalTok{(}\DecValTok{20260509}\NormalTok{)}

\FunctionTok{source}\NormalTok{(here}\SpecialCharTok{::}\FunctionTok{here}\NormalTok{(}\StringTok{"R"}\NormalTok{, }\StringTok{"api\_helpers.R"}\NormalTok{))}
\FunctionTok{source}\NormalTok{(here}\SpecialCharTok{::}\FunctionTok{here}\NormalTok{(}\StringTok{"R"}\NormalTok{, }\StringTok{"utils.R"}\NormalTok{))}
\FunctionTok{source}\NormalTok{(here}\SpecialCharTok{::}\FunctionTok{here}\NormalTok{(}\StringTok{"R"}\NormalTok{, }\StringTok{"sci\_palette.R"}\NormalTok{))}
\end{Highlighting}
\end{Shaded}

\section{Conceptual background}\label{repro-background}

The ``replication crisis'' has prompted widespread adoption of transparency practices: open data, open code, pre-registration, and registered reports. Bibliometrics can measure the uptake of these practices at scale, providing evidence for policy interventions.

\textbf{Open badges} are visual indicators (icons) attached to published papers that certify the availability of data, materials, or pre-registration. Journals like \emph{Psychological Science} pioneered their use. Studies suggest that badges increase data sharing rates, though the quality of shared data varies.

\textbf{Research Resource Identifiers (RRIDs)} are persistent identifiers for research resources: antibodies, cell lines, model organisms, and software tools. Including RRIDs in methods sections improves reproducibility by enabling precise identification of the materials used. The RRID initiative has been adopted by over 1,000 journals.

\textbf{Data availability statements (DAS)} declare whether the data underlying a study are available and where to find them. Mandatory DAS policies have become common, but many statements say ``data available upon request'' without providing actual access. The presence of a DAS is easy to detect bibliometrically; the quality of actual data sharing is much harder to assess.

\textbf{Code availability} is increasingly expected, especially in computational fields. Journals like \emph{Nature} require code availability statements. Platforms like GitHub, Zenodo, and Code Ocean facilitate code sharing, and DOIs for code repositories enable citation and tracking.

These practices are measurable proxies for reproducibility, but they are not reproducibility itself. A paper with an open badge may have unusable data; a paper without one may be perfectly reproducible. \citet{hicks2015} warn against using transparency indicators as mechanical evaluation criteria.

\section{Worked example}\label{repro-example}

\subsection{Analysing access status and metadata coverage}\label{analysing-access-status-and-metadata-coverage}

The example below measures article access status and abstract coverage in
OpenAlex. It does not detect data/code availability, RRIDs, badges, or successful
replication. Those questions require separate full-text/resource extraction
and validation; OA status is not a validated substitute for them.

\begin{Shaded}
\begin{Highlighting}[]
\NormalTok{works }\OtherTok{\textless{}{-}} \FunctionTok{fetch\_openalex}\NormalTok{(}
  \AttributeTok{entity =} \StringTok{"works"}\NormalTok{,}
  \AttributeTok{primary\_location.source.id =} \StringTok{"S148561398"}\NormalTok{,}
  \AttributeTok{from\_publication\_date =} \StringTok{"2018{-}01{-}01"}\NormalTok{,}
  \AttributeTok{to\_publication\_date =} \StringTok{"2023{-}12{-}31"}\NormalTok{,}
  \AttributeTok{type =} \StringTok{"article"}\NormalTok{,}
  \AttributeTok{options =} \FunctionTok{list}\NormalTok{(}\AttributeTok{sample =} \DecValTok{400}\NormalTok{, }\AttributeTok{seed =} \DecValTok{42}\NormalTok{)}
\NormalTok{)}

\NormalTok{works\_repro }\OtherTok{\textless{}{-}}\NormalTok{ works }\SpecialCharTok{|\textgreater{}}
  \FunctionTok{transmute}\NormalTok{(}
\NormalTok{    id, display\_name,}
    \AttributeTok{year =} \FunctionTok{year}\NormalTok{(publication\_date),}
\NormalTok{    cited\_by\_count,}
\NormalTok{    oa\_status,}
\NormalTok{    type,}
    \AttributeTok{has\_abstract =} \SpecialCharTok{!}\FunctionTok{is.na}\NormalTok{(abstract) }\SpecialCharTok{\&} \FunctionTok{nzchar}\NormalTok{(}\FunctionTok{trimws}\NormalTok{(abstract))}
\NormalTok{  )}

\FunctionTok{cat}\NormalTok{(}\FunctionTok{glue}\NormalTok{(}\StringTok{"Works: \{nrow(works\_repro)\}}\SpecialCharTok{\textbackslash{}n}\StringTok{"}\NormalTok{))}
\end{Highlighting}
\end{Shaded}

\begin{verbatim}
#> Works: 400
\end{verbatim}

\subsection{Observed article access status}\label{observed-article-access-status}

\begin{Shaded}
\begin{Highlighting}[]
\NormalTok{works\_repro }\SpecialCharTok{|\textgreater{}}
  \FunctionTok{group\_by}\NormalTok{(year) }\SpecialCharTok{|\textgreater{}}
  \FunctionTok{summarise}\NormalTok{(}\AttributeTok{oa\_rate =} \FunctionTok{mean}\NormalTok{(oa\_status }\SpecialCharTok{!=} \StringTok{"closed"}\NormalTok{, }\AttributeTok{na.rm =} \ConstantTok{TRUE}\NormalTok{),}
            \AttributeTok{.groups =} \StringTok{"drop"}\NormalTok{) }\SpecialCharTok{|\textgreater{}}
  \FunctionTok{ggplot}\NormalTok{(}\FunctionTok{aes}\NormalTok{(}\AttributeTok{x =}\NormalTok{ year, }\AttributeTok{y =}\NormalTok{ oa\_rate)) }\SpecialCharTok{+}
  \FunctionTok{geom\_line}\NormalTok{(}\AttributeTok{linewidth =} \DecValTok{1}\NormalTok{, }\AttributeTok{colour =} \FunctionTok{palette\_sci}\NormalTok{(}\DecValTok{1}\NormalTok{)) }\SpecialCharTok{+}
  \FunctionTok{geom\_point}\NormalTok{(}\AttributeTok{size =} \DecValTok{2}\NormalTok{, }\AttributeTok{colour =} \FunctionTok{palette\_sci}\NormalTok{(}\DecValTok{1}\NormalTok{)) }\SpecialCharTok{+}
  \FunctionTok{scale\_y\_continuous}\NormalTok{(}\AttributeTok{labels =}\NormalTok{ scales}\SpecialCharTok{::}\NormalTok{percent, }\AttributeTok{limits =} \FunctionTok{c}\NormalTok{(}\DecValTok{0}\NormalTok{, }\DecValTok{1}\NormalTok{)) }\SpecialCharTok{+}
  \FunctionTok{labs}\NormalTok{(}\AttributeTok{x =} \StringTok{"Year"}\NormalTok{, }\AttributeTok{y =} \StringTok{"OA rate"}\NormalTok{) }\SpecialCharTok{+}
  \FunctionTok{theme\_sci}\NormalTok{()}
\end{Highlighting}
\end{Shaded}

\begin{figure}

{\centering \includegraphics[width=0.9\linewidth,alt={Line chart of recorded article access status by publication year; it does not measure open data or code.}]{scientometrics-in-r_files/figure-latex/33-fig-repro-oa-trend-1} 

}

\caption{Open-access share among sampled articles with observed access status, by year.}\label{fig:33-fig-repro-oa-trend}
\end{figure}

\subsection{Abstract availability as metadata completeness}\label{abstract-availability-as-metadata-completeness}

\begin{Shaded}
\begin{Highlighting}[]
\NormalTok{works\_repro }\SpecialCharTok{|\textgreater{}}
  \FunctionTok{group\_by}\NormalTok{(year) }\SpecialCharTok{|\textgreater{}}
  \FunctionTok{summarise}\NormalTok{(}\AttributeTok{abstract\_rate =} \FunctionTok{mean}\NormalTok{(has\_abstract), }\AttributeTok{.groups =} \StringTok{"drop"}\NormalTok{) }\SpecialCharTok{|\textgreater{}}
  \FunctionTok{ggplot}\NormalTok{(}\FunctionTok{aes}\NormalTok{(}\AttributeTok{x =}\NormalTok{ year, }\AttributeTok{y =}\NormalTok{ abstract\_rate)) }\SpecialCharTok{+}
  \FunctionTok{geom\_line}\NormalTok{(}\AttributeTok{linewidth =} \DecValTok{1}\NormalTok{, }\AttributeTok{colour =} \FunctionTok{palette\_sci}\NormalTok{(}\DecValTok{1}\NormalTok{)) }\SpecialCharTok{+}
  \FunctionTok{geom\_point}\NormalTok{(}\AttributeTok{size =} \DecValTok{2}\NormalTok{, }\AttributeTok{colour =} \FunctionTok{palette\_sci}\NormalTok{(}\DecValTok{1}\NormalTok{)) }\SpecialCharTok{+}
  \FunctionTok{scale\_y\_continuous}\NormalTok{(}\AttributeTok{labels =}\NormalTok{ scales}\SpecialCharTok{::}\NormalTok{percent, }\AttributeTok{limits =} \FunctionTok{c}\NormalTok{(}\DecValTok{0}\NormalTok{, }\DecValTok{1}\NormalTok{)) }\SpecialCharTok{+}
  \FunctionTok{labs}\NormalTok{(}\AttributeTok{x =} \StringTok{"Year"}\NormalTok{, }\AttributeTok{y =} \StringTok{"Abstract available"}\NormalTok{) }\SpecialCharTok{+}
  \FunctionTok{theme\_sci}\NormalTok{()}
\end{Highlighting}
\end{Shaded}

\begin{figure}

{\centering \includegraphics[width=0.9\linewidth,alt={Line chart showing the proportion of articles with available abstracts over time.}]{scientometrics-in-r_files/figure-latex/33-fig-repro-abstract-1} 

}

\caption{Abstract availability rate by year.}\label{fig:33-fig-repro-abstract}
\end{figure}

\subsection{Access status and cumulative citations}\label{access-status-and-cumulative-citations}

\begin{Shaded}
\begin{Highlighting}[]
\NormalTok{works\_repro }\SpecialCharTok{|\textgreater{}}
  \FunctionTok{mutate}\NormalTok{(}\AttributeTok{is\_oa =} \FunctionTok{ifelse}\NormalTok{(oa\_status }\SpecialCharTok{!=} \StringTok{"closed"}\NormalTok{, }\StringTok{"OA"}\NormalTok{, }\StringTok{"Closed"}\NormalTok{)) }\SpecialCharTok{|\textgreater{}}
  \FunctionTok{ggplot}\NormalTok{(}\FunctionTok{aes}\NormalTok{(}\AttributeTok{x =}\NormalTok{ is\_oa, }\AttributeTok{y =}\NormalTok{ cited\_by\_count }\SpecialCharTok{+} \DecValTok{1}\NormalTok{)) }\SpecialCharTok{+}
  \FunctionTok{geom\_boxplot}\NormalTok{(}\AttributeTok{fill =} \FunctionTok{palette\_sci}\NormalTok{(}\DecValTok{1}\NormalTok{), }\AttributeTok{alpha =} \FloatTok{0.7}\NormalTok{) }\SpecialCharTok{+}
  \FunctionTok{scale\_y\_log10}\NormalTok{() }\SpecialCharTok{+}
  \FunctionTok{labs}\NormalTok{(}\AttributeTok{x =} \StringTok{"Access status"}\NormalTok{, }\AttributeTok{y =} \StringTok{"Citations (log scale)"}\NormalTok{) }\SpecialCharTok{+}
  \FunctionTok{theme\_sci}\NormalTok{()}
\end{Highlighting}
\end{Shaded}

\begin{figure}

{\centering \includegraphics[width=0.9\linewidth,alt={Box plot comparing cumulative citations for OA and closed papers without estimating a causal effect or data-sharing relationship.}]{scientometrics-in-r_files/figure-latex/33-fig-repro-impact-1} 

}

\caption{Current cumulative citation counts by recorded article access status.}\label{fig:33-fig-repro-impact}
\end{figure}

\subsection{Summary statistics}\label{summary-statistics}

\begin{Shaded}
\begin{Highlighting}[]
\NormalTok{works\_repro }\SpecialCharTok{|\textgreater{}}
  \FunctionTok{mutate}\NormalTok{(}\AttributeTok{is\_oa =}\NormalTok{ oa\_status }\SpecialCharTok{!=} \StringTok{"closed"}\NormalTok{) }\SpecialCharTok{|\textgreater{}}
  \FunctionTok{group\_by}\NormalTok{(is\_oa) }\SpecialCharTok{|\textgreater{}}
  \FunctionTok{summarise}\NormalTok{(}
    \AttributeTok{n =} \FunctionTok{n}\NormalTok{(),}
    \AttributeTok{mean\_cites =} \FunctionTok{round}\NormalTok{(}\FunctionTok{mean}\NormalTok{(cited\_by\_count), }\DecValTok{1}\NormalTok{),}
    \AttributeTok{abstract\_rate =}\NormalTok{ scales}\SpecialCharTok{::}\FunctionTok{percent}\NormalTok{(}\FunctionTok{mean}\NormalTok{(has\_abstract)),}
    \AttributeTok{.groups =} \StringTok{"drop"}
\NormalTok{  ) }\SpecialCharTok{|\textgreater{}}
  \FunctionTok{gt}\NormalTok{()}
\end{Highlighting}
\end{Shaded}

\begingroup
\fontsize{8.0pt}{9.0pt}\selectfont
\begin{longtable}{>{\centering\arraybackslash}p{\dimexpr 0.23\linewidth -2\tabcolsep-1.5\arrayrulewidth}>{\raggedleft\arraybackslash}p{\dimexpr 0.23\linewidth -2\tabcolsep-1.5\arrayrulewidth}>{\raggedleft\arraybackslash}p{\dimexpr 0.25\linewidth -2\tabcolsep-1.5\arrayrulewidth}>{\raggedleft\arraybackslash}p{\dimexpr 0.29\linewidth -2\tabcolsep-1.5\arrayrulewidth}}
\toprule
is\_oa & n & \nolinkurl{mean_cites} & \nolinkurl{abstract_rate} \\ 
\midrule\addlinespace[2.5pt]
FALSE & 205 & 24.2 & 0\% \\ 
TRUE & 195 & 32.1 & 54\% \\ 
\bottomrule
\end{longtable}
\endgroup

\section{Diagnostics and interpretation}\label{repro-diagnostics}

\begin{itemize}
\tightlist
\item
  \textbf{Different constructs}: Article access status and abstract coverage do not measure data/code sharing, RRID use, badge adoption or reproducibility. Check full texts and linked resources for those outcomes.
\item
  \textbf{Coverage over time}: Changes in recorded access status or abstract coverage can reflect provider coverage and publication cohorts. The example does not evaluate a DAS or badge policy.
\item
  \textbf{Selection and exposure}: Current cumulative citations mix publication ages, article selection and other confounders. The comparison cannot establish an OA effect or a relationship with open data.
\item
  \textbf{Field differences}: Data sharing norms vary enormously. Genomics shares data routinely; clinical research often cannot due to privacy. Compare within fields.
\end{itemize}

\section{Limitations and responsible use}\label{repro-limitations}

\section{Limitations and responsible use}\label{limitations-and-responsible-use-32}

\begin{itemize}
\tightlist
\item
  \textbf{Signals are not substance.} A data availability statement does not guarantee the data are actually available, usable, or correct. Badges certify process, not outcome.
\item
  \textbf{Mandates inflate compliance without improving practice.} When journals require DAS, compliance rises but ``available upon request'' becomes the default, not genuine sharing.
\item
  \textbf{Reproducibility is multidimensional.} No single indicator captures it. Data availability, code availability, methods transparency, pre-registration, and statistical rigour all contribute independently.
\item
  \textbf{Do not penalise fields where sharing is constrained.} Clinical data (privacy), indigenous data (sovereignty), and sensitive social data cannot be freely shared. Mandating open data policies uniformly across fields is inappropriate \citep{hicks2015}.
\end{itemize}

\section{Common pitfalls}\label{repro-pitfalls}

\section{Common pitfalls}\label{common-pitfalls-32}

\begin{itemize}
\tightlist
\item
  \textbf{Equating open data with reproducibility.} Data can be open but undocumented, incorrectly formatted, or insufficient to reproduce the analysis.
\item
  \textbf{Counting badges without checking quality.} A badge means the journal certified availability at publication time. The link may be broken or the data may have been removed.
\item
  \textbf{Ignoring negative results.} Reproducibility also requires publishing null results and failed replications. These are poorly captured by current transparency indicators.
\item
  \textbf{Using transparency metrics as performance indicators.} Transparency practices should be valued, but mechanically rewarding them creates perverse incentives (superficial compliance).
\end{itemize}

\section{Exercises}\label{repro-exercises}

\begin{enumerate}
\def\labelenumi{\arabic{enumi}.}
\item
  \textbf{DAS detection.} For papers with full abstracts, search for keywords like ``data availability'', ``code availability'', or ``github.com'' in the abstract text. What proportion mention data or code sharing?
\item
  \textbf{Cross-journal comparison.} Compare abstract availability and OA rates for three journals. Do journals with higher OA rates also have better metadata completeness?
\item
  \textbf{Temporal trends.} Track recorded OA status and abstract coverage by publication cohort. Report the usable denominators and explain why these measures do not establish data/code sharing or reproducibility.
\end{enumerate}

\section{Solutions}\label{33-solutions}

Solutions are provided in \ref{solutions}.

\section{Further reading}\label{repro-reading}

\begin{itemize}
\tightlist
\item
  \citet{hicks2015} --- The Leiden Manifesto: responsible use of metrics, including transparency indicators.
\item
  \citet{priem2022} --- OpenAlex metadata as a basis for reproducibility analysis.
\item
  \citet{dora2012} --- DORA: assessment should consider research practices, not just outputs.
\end{itemize}

\section{Session info}\label{repro-sessioninfo}

\begin{Shaded}
\begin{Highlighting}[]
\FunctionTok{sessionInfo}\NormalTok{()}
\end{Highlighting}
\end{Shaded}

\begin{verbatim}
#> R version 4.6.1 (2026-06-24)
#> Platform: x86_64-pc-linux-gnu
#> Running under: Ubuntu 24.04.5 LTS
#> 
#> Matrix products: default
#> BLAS:   /usr/lib/x86_64-linux-gnu/openblas-pthread/libblas.so.3 
#> LAPACK: /usr/lib/x86_64-linux-gnu/openblas-pthread/libopenblasp-r0.3.26.so;  LAPACK version 3.12.0
#> 
#> locale:
#>  [1] LC_CTYPE=C.UTF-8       LC_NUMERIC=C           LC_TIME=C.UTF-8       
#>  [4] LC_COLLATE=C.UTF-8     LC_MONETARY=C.UTF-8    LC_MESSAGES=C.UTF-8   
#>  [7] LC_PAPER=C.UTF-8       LC_NAME=C              LC_ADDRESS=C          
#> [10] LC_TELEPHONE=C         LC_MEASUREMENT=C.UTF-8 LC_IDENTIFICATION=C   
#> 
#> time zone: UTC
#> tzcode source: system (glibc)
#> 
#> attached base packages:
#> [1] stats     graphics  grDevices utils     datasets  methods   base     
#> 
#> other attached packages:
#>  [1] uwot_0.2.5                Matrix_1.7-5             
#>  [3] word2vec_0.4.1            stm_1.3.8                
#>  [5] topicmodels_0.2-17        quanteda.textstats_0.97.2
#>  [7] visNetwork_2.1.4          ggraph_2.2.2             
#>  [9] tidygraph_1.3.1           igraph_2.3.4             
#> [11] quanteda_4.5.1            pdftools_3.9.1           
#> [13] arrow_26.0.0              bibliometrix_5.5.0       
#> [15] RefManageR_1.4.0          bib2df_1.1.2.0           
#> [17] rcrossref_1.2.1           gt_1.4.0                 
#> [19] tidytext_0.4.3            glue_1.8.1               
#> [21] openalexR_3.1.2           lubridate_1.9.5          
#> [23] forcats_1.0.1             stringr_1.6.0            
#> [25] dplyr_1.2.1               purrr_1.2.2              
#> [27] readr_2.2.0               tidyr_1.3.2              
#> [29] tibble_3.3.1              ggplot2_4.0.3            
#> [31] tidyverse_2.0.0          
#> 
#> loaded via a namespace (and not attached):
#>   [1] bibtex_0.5.3           RColorBrewer_1.1-3     rstudioapi_0.19.0     
#>   [4] jsonlite_2.0.0         magrittr_2.0.5         modeltools_0.2-25     
#>   [7] farver_2.1.2           rmarkdown_2.32         fs_2.1.0              
#>  [10] vctrs_0.7.3            memoise_2.0.1          askpass_1.2.1         
#>  [13] base64enc_0.1-6        htmltools_0.5.9        contentanalysis_1.1.1 
#>  [16] curl_8.1.0             broom_1.0.13           janeaustenr_1.0.0     
#>  [19] cellranger_1.1.0       htmlwidgets_1.6.4      tokenizers_0.3.0      
#>  [22] plyr_1.8.9             httr2_1.3.0            cachem_1.1.0          
#>  [25] plotly_4.12.1          dimensionsR_0.0.3      mime_0.13             
#>  [28] lifecycle_1.0.5        pkgconfig_2.0.3        R6_2.6.1              
#>  [31] fastmap_1.2.0          shiny_1.14.0           digest_0.6.39         
#>  [34] patchwork_1.3.2        shinycssloaders_1.1.0  rprojroot_2.1.1       
#>  [37] RSpectra_0.16-2        SnowballC_0.7.1        labeling_0.4.3        
#>  [40] urltools_1.7.3.1       timechange_0.4.0       mgcv_1.9-4            
#>  [43] httr_1.4.9             polyclip_1.10-7        compiler_4.6.1        
#>  [46] here_1.0.2             bit64_4.8.6            withr_3.0.3           
#>  [49] backports_1.5.1        S7_0.2.2               viridis_0.6.5         
#>  [52] ggforce_0.5.0          MASS_7.3-65            bibliometrixData_0.3.0
#>  [55] tools_4.6.1            otel_0.2.0             stopwords_2.3         
#>  [58] zip_3.0.2              rentrez_1.2.4          httpuv_1.6.17         
#>  [61] nlme_3.1-169           promises_1.5.0         grid_4.6.1            
#>  [64] stringdist_0.9.17      reshape2_1.4.5         generics_0.1.4        
#>  [67] gtable_0.3.6           tzdb_0.5.0             rscopus_0.9.0         
#>  [70] ca_0.72                data.table_1.18.6.1    hms_1.1.4             
#>  [73] xml2_1.6.0             utf8_1.2.6             ggrepel_0.9.8         
#>  [76] pillar_1.11.1          nsyllable_1.0.1        vroom_1.7.1           
#>  [79] later_1.4.8            splines_4.6.1          tweenr_2.0.3          
#>  [82] lattice_0.22-9         FNN_1.1.4.1            bit_4.6.0             
#>  [85] tidyselect_1.2.1       tm_0.7-20              miniUI_0.1.2          
#>  [88] knitr_1.52             gridExtra_2.3.1        NLP_0.3-3             
#>  [91] bookdown_0.48          stats4_4.6.1           crul_1.6.0            
#>  [94] xfun_0.62              graphlayouts_1.2.5     matrixStats_1.5.0     
#>  [97] DT_0.34.0              humaniformat_0.6.0     stringi_1.8.9         
#> [100] qpdf_1.4.1             yaml_2.3.12            evaluate_1.0.5        
#> [103] httpcode_0.3.0         cli_3.6.6              xtable_1.8-8          
#> [106] systemfonts_1.3.2      Rcpp_1.1.2             readxl_1.6.0          
#> [109] triebeard_0.4.1        XML_3.99-0.25          parallel_4.6.1        
#> [112] assertthat_0.2.1       pubmedR_1.0.2          slam_0.1-56           
#> [115] viridisLite_0.4.3      scales_1.4.0           crayon_1.5.3          
#> [118] openxlsx_4.2.9         rlang_1.3.0            fastmatch_1.1-8
\end{verbatim}

\part{Production \& Reporting}\label{part-production-reporting}

\chapter{Reproducible Pipelines}\label{reproducible-pipelines}

\emph{R. Heller · Date: 2026-05-09 · Last changed: 2026-10-31}

\section{Learning objectives}\label{pipeline-objectives}

After completing this chapter, you will be able to:

\begin{itemize}
\tightlist
\item
  Explain why ad hoc scripts are insufficient for reproducible bibliometric analysis
\item
  Define a targets pipeline with dependency tracking and selective re-execution
\item
  Use renv to snapshot and restore R package environments
\item
  Understand how Docker containers ensure cross-platform reproducibility
\item
  Structure a bibliometric project for long-term reproducibility
\end{itemize}

\section{Setup}\label{pipeline-setup}

\begin{Shaded}
\begin{Highlighting}[]
\FunctionTok{library}\NormalTok{(tidyverse)}
\FunctionTok{library}\NormalTok{(glue)}

\FunctionTok{set.seed}\NormalTok{(}\DecValTok{20260509}\NormalTok{)}

\FunctionTok{source}\NormalTok{(here}\SpecialCharTok{::}\FunctionTok{here}\NormalTok{(}\StringTok{"R"}\NormalTok{, }\StringTok{"sci\_palette.R"}\NormalTok{))}
\end{Highlighting}
\end{Shaded}

\section{Conceptual background}\label{pipeline-background}

A bibliometric analysis is a multi-step process: fetch data from an API, clean and deduplicate records, compute indicators, build networks, fit models, and generate figures. Each step depends on earlier ones. When the upstream data changes (a new API query, an updated corpus), downstream results must be regenerated --- but only the affected parts.

\textbf{Ad hoc scripts} (running scripts manually in sequence) fail at scale: they do not track dependencies, they re-run everything even when only one input changed, and they break when run on a different machine with different package versions.

\textbf{\texttt{targets}} is an R pipeline tool that solves the dependency problem. You define each analysis step as a ``target'' (a function that takes inputs and produces outputs). \texttt{targets} builds a dependency graph, executes only the targets whose inputs have changed, and caches results. This is both faster (no unnecessary re-computation) and safer (no stale intermediate results).

\textbf{\texttt{renv}} snapshots the exact package versions used in your analysis into a \texttt{renv.lock} file. A collaborator (or your future self) can restore the identical package environment with \texttt{renv::restore()}. This eliminates the ``works on my machine'' problem for R packages.

\textbf{Docker} takes reproducibility further by containerising the entire compute environment: operating system, R version, system libraries, and package versions. A \texttt{Dockerfile} defines the environment declaratively; anyone with Docker can rebuild it identically. Docker is essential for long-term reproducibility, since R versions and system libraries evolve.

These tools complement each other: \texttt{targets} manages the computation, \texttt{renv} manages the packages, and Docker manages the environment. Together, they make a bibliometric analysis fully reproducible years after it was originally run.

\section{Worked example}\label{pipeline-example}

\subsection{Defining a targets pipeline}\label{defining-a-targets-pipeline}

The \texttt{\_targets.R} file defines the pipeline. Here we show the structure for a typical bibliometric analysis.

\begin{Shaded}
\begin{Highlighting}[]
\CommentTok{\# \_targets.R — pipeline definition}
\FunctionTok{library}\NormalTok{(targets)}

\FunctionTok{tar\_option\_set}\NormalTok{(}\AttributeTok{packages =} \FunctionTok{c}\NormalTok{(}\StringTok{"tidyverse"}\NormalTok{, }\StringTok{"openalexR"}\NormalTok{, }\StringTok{"igraph"}\NormalTok{))}

\FunctionTok{list}\NormalTok{(}
  \FunctionTok{tar\_target}\NormalTok{(raw\_works, \{}
    \FunctionTok{fetch\_openalex}\NormalTok{(}
      \AttributeTok{entity =} \StringTok{"works"}\NormalTok{,}
      \AttributeTok{primary\_location.source.id =} \StringTok{"S148561398"}\NormalTok{,}
      \AttributeTok{from\_publication\_date =} \StringTok{"2020{-}01{-}01"}\NormalTok{,}
      \AttributeTok{to\_publication\_date =} \StringTok{"2023{-}12{-}31"}\NormalTok{,}
      \AttributeTok{options =} \FunctionTok{list}\NormalTok{(}\AttributeTok{sample =} \DecValTok{300}\NormalTok{, }\AttributeTok{seed =} \DecValTok{42}\NormalTok{)}
\NormalTok{    )}
\NormalTok{  \}),}

  \FunctionTok{tar\_target}\NormalTok{(clean\_works, \{}
\NormalTok{    raw\_works }\SpecialCharTok{|\textgreater{}}
      \FunctionTok{filter}\NormalTok{(}\SpecialCharTok{!}\FunctionTok{is.na}\NormalTok{(doi)) }\SpecialCharTok{|\textgreater{}}
      \FunctionTok{distinct}\NormalTok{(doi, }\AttributeTok{.keep\_all =} \ConstantTok{TRUE}\NormalTok{) }\SpecialCharTok{|\textgreater{}}
      \FunctionTok{mutate}\NormalTok{(}\AttributeTok{year =} \FunctionTok{year}\NormalTok{(publication\_date))}
\NormalTok{  \}),}

  \FunctionTok{tar\_target}\NormalTok{(citation\_stats, \{}
\NormalTok{    clean\_works }\SpecialCharTok{|\textgreater{}}
      \FunctionTok{group\_by}\NormalTok{(year) }\SpecialCharTok{|\textgreater{}}
      \FunctionTok{summarise}\NormalTok{(}
        \AttributeTok{n =} \FunctionTok{n}\NormalTok{(),}
        \AttributeTok{mean\_cites =} \FunctionTok{mean}\NormalTok{(cited\_by\_count),}
        \AttributeTok{.groups =} \StringTok{"drop"}
\NormalTok{      )}
\NormalTok{  \}),}

  \FunctionTok{tar\_target}\NormalTok{(fig\_trends, \{}
    \FunctionTok{ggplot}\NormalTok{(citation\_stats, }\FunctionTok{aes}\NormalTok{(}\AttributeTok{x =}\NormalTok{ year, }\AttributeTok{y =}\NormalTok{ mean\_cites)) }\SpecialCharTok{+}
      \FunctionTok{geom\_line}\NormalTok{() }\SpecialCharTok{+} \FunctionTok{geom\_point}\NormalTok{() }\SpecialCharTok{+}
      \FunctionTok{labs}\NormalTok{(}\AttributeTok{x =} \StringTok{"Year"}\NormalTok{, }\AttributeTok{y =} \StringTok{"Mean citations"}\NormalTok{)}
\NormalTok{  \})}
\NormalTok{)}
\end{Highlighting}
\end{Shaded}

\subsection{Running and inspecting the pipeline}\label{running-and-inspecting-the-pipeline}

\begin{Shaded}
\begin{Highlighting}[]
\CommentTok{\# Run the pipeline}
\NormalTok{targets}\SpecialCharTok{::}\FunctionTok{tar\_make}\NormalTok{()}

\CommentTok{\# Visualise the dependency graph}
\NormalTok{targets}\SpecialCharTok{::}\FunctionTok{tar\_visnetwork}\NormalTok{()}

\CommentTok{\# Read a specific target}
\NormalTok{targets}\SpecialCharTok{::}\FunctionTok{tar\_read}\NormalTok{(citation\_stats)}
\end{Highlighting}
\end{Shaded}

\subsection{Project structure}\label{project-structure}

\begin{Shaded}
\begin{Highlighting}[]
\CommentTok{\# Recommended project layout:}
\CommentTok{\#}
\CommentTok{\# my{-}bibliometric{-}project/}
\CommentTok{\# ├── \_targets.R          \# Pipeline definition}
\CommentTok{\# ├── R/                  \# Helper functions}
\CommentTok{\# │   ├── fetch.R}
\CommentTok{\# │   ├── clean.R}
\CommentTok{\# │   └── analyse.R}
\CommentTok{\# ├── renv.lock           \# Package snapshot}
\CommentTok{\# ├── Dockerfile          \# Environment definition}
\CommentTok{\# ├── data/               \# Cached data (Parquet)}
\CommentTok{\# ├── output/             \# Figures and tables}
\CommentTok{\# └── report.qmd          \# Quarto report consuming targets}
\end{Highlighting}
\end{Shaded}

\subsection{Using renv}\label{using-renv}

\begin{Shaded}
\begin{Highlighting}[]
\CommentTok{\# Initialise renv in a new project}
\NormalTok{renv}\SpecialCharTok{::}\FunctionTok{init}\NormalTok{()}

\CommentTok{\# Snapshot current package state}
\NormalTok{renv}\SpecialCharTok{::}\FunctionTok{snapshot}\NormalTok{()}

\CommentTok{\# Restore packages on a new machine}
\NormalTok{renv}\SpecialCharTok{::}\FunctionTok{restore}\NormalTok{()}
\end{Highlighting}
\end{Shaded}

\subsection{Dockerfile for bibliometric analysis}\label{dockerfile-for-bibliometric-analysis}

\begin{Shaded}
\begin{Highlighting}[]
\CommentTok{\# Example Dockerfile (not R code — shown for reference)}
\CommentTok{\#}
\CommentTok{\# FROM rocker/verse:4.4.1}
\CommentTok{\# RUN install2.r targets renv openalexR igraph}
\CommentTok{\# COPY renv.lock /project/renv.lock}
\CommentTok{\# WORKDIR /project}
\CommentTok{\# RUN R {-}e "renv::restore()"}
\CommentTok{\# CMD ["R", "{-}e", "targets::tar\_make()"]}
\end{Highlighting}
\end{Shaded}

\section{Diagnostics and interpretation}\label{pipeline-diagnostics}

\begin{itemize}
\tightlist
\item
  \textbf{Target status}: \texttt{tar\_outdated()} lists which targets need re-running. If everything is up to date, \texttt{tar\_make()} completes instantly.
\item
  \textbf{Pipeline visualisation}: \texttt{tar\_visnetwork()} shows the dependency graph. Disconnected subgraphs indicate independent analysis streams.
\item
  \textbf{Cache management}: Targets stores cached results in \texttt{\_targets/}. This directory can grow large. Use \texttt{tar\_prune()} to remove unused targets.
\item
  \textbf{renv consistency}: Run \texttt{renv::status()} to check if your lockfile matches the installed packages.
\end{itemize}

\section{Limitations and responsible use}\label{pipeline-limitations}

\section{Limitations and responsible use}\label{limitations-and-responsible-use-33}

\begin{itemize}
\tightlist
\item
  \textbf{API rate limits}: targets will re-run API calls when inputs change. Use file-based caching (Parquet) to avoid hitting API rate limits on every pipeline run.
\item
  \textbf{Docker learning curve}: Docker adds complexity. For solo projects, renv alone may suffice. Docker becomes essential for team reproducibility and long-term archiving.
\item
  \textbf{Overhead for small projects}: A 50-line analysis script does not need targets. Use pipelines when the analysis has 5+ interdependent steps or will be re-run multiple times.
\item
  \textbf{Reproducibility is not eternal}: Even with Docker, operating system patches and hardware changes can affect results. Archive final outputs alongside the pipeline.
\end{itemize}

\section{Common pitfalls}\label{pipeline-pitfalls}

\section{Common pitfalls}\label{common-pitfalls-33}

\begin{itemize}
\tightlist
\item
  \textbf{Not separating data retrieval from analysis.} API calls in analysis targets re-fetch data on every change. Separate fetching into its own target with file-based caching.
\item
  \textbf{Forgetting to commit renv.lock.} The lockfile is the reproducibility guarantee. Always version-control it.
\item
  **Ignoring the \_targets/ directory in .gitignore.** Cached results should not be committed to git (they can be regenerated). Add \texttt{\_targets/} to \texttt{.gitignore}.
\item
  \textbf{Writing monolithic targets.} Each target should do one thing. Fine-grained targets enable selective re-execution.
\end{itemize}

\section{Exercises}\label{pipeline-exercises}

\begin{enumerate}
\def\labelenumi{\arabic{enumi}.}
\item
  \textbf{Build a pipeline.} Create a \texttt{\_targets.R} file for a simple bibliometric analysis: fetch → clean → summarise → plot. Run it and verify that changing the summary function only re-runs downstream targets.
\item
  \textbf{renv snapshot.} Initialise renv in a project, install a new package, and snapshot. Delete the package and restore from the lockfile.
\item
  \textbf{Branching.} Use \texttt{tar\_target()} with dynamic branching to run the same analysis across multiple journals in parallel.
\end{enumerate}

\section{Solutions}\label{34-solutions}

Solutions are provided in \ref{solutions}.

\section{Further reading}\label{pipeline-reading}

\begin{itemize}
\tightlist
\item
  \citet{hicks2015} --- The Leiden Manifesto: transparency and reproducibility in research evaluation.
\item
  \citet{priem2022} --- OpenAlex as a reproducible data source for bibliometric pipelines.
\end{itemize}

\section{Session info}\label{pipeline-sessioninfo}

\begin{Shaded}
\begin{Highlighting}[]
\FunctionTok{sessionInfo}\NormalTok{()}
\end{Highlighting}
\end{Shaded}

\begin{verbatim}
#> R version 4.6.1 (2026-06-24)
#> Platform: x86_64-pc-linux-gnu
#> Running under: Ubuntu 24.04.5 LTS
#> 
#> Matrix products: default
#> BLAS:   /usr/lib/x86_64-linux-gnu/openblas-pthread/libblas.so.3 
#> LAPACK: /usr/lib/x86_64-linux-gnu/openblas-pthread/libopenblasp-r0.3.26.so;  LAPACK version 3.12.0
#> 
#> locale:
#>  [1] LC_CTYPE=C.UTF-8       LC_NUMERIC=C           LC_TIME=C.UTF-8       
#>  [4] LC_COLLATE=C.UTF-8     LC_MONETARY=C.UTF-8    LC_MESSAGES=C.UTF-8   
#>  [7] LC_PAPER=C.UTF-8       LC_NAME=C              LC_ADDRESS=C          
#> [10] LC_TELEPHONE=C         LC_MEASUREMENT=C.UTF-8 LC_IDENTIFICATION=C   
#> 
#> time zone: UTC
#> tzcode source: system (glibc)
#> 
#> attached base packages:
#> [1] stats     graphics  grDevices utils     datasets  methods   base     
#> 
#> other attached packages:
#>  [1] uwot_0.2.5                Matrix_1.7-5             
#>  [3] word2vec_0.4.1            stm_1.3.8                
#>  [5] topicmodels_0.2-17        quanteda.textstats_0.97.2
#>  [7] visNetwork_2.1.4          ggraph_2.2.2             
#>  [9] tidygraph_1.3.1           igraph_2.3.4             
#> [11] quanteda_4.5.1            pdftools_3.9.1           
#> [13] arrow_26.0.0              bibliometrix_5.5.0       
#> [15] RefManageR_1.4.0          bib2df_1.1.2.0           
#> [17] rcrossref_1.2.1           gt_1.4.0                 
#> [19] tidytext_0.4.3            glue_1.8.1               
#> [21] openalexR_3.1.2           lubridate_1.9.5          
#> [23] forcats_1.0.1             stringr_1.6.0            
#> [25] dplyr_1.2.1               purrr_1.2.2              
#> [27] readr_2.2.0               tidyr_1.3.2              
#> [29] tibble_3.3.1              ggplot2_4.0.3            
#> [31] tidyverse_2.0.0          
#> 
#> loaded via a namespace (and not attached):
#>   [1] bibtex_0.5.3           RColorBrewer_1.1-3     rstudioapi_0.19.0     
#>   [4] jsonlite_2.0.0         magrittr_2.0.5         modeltools_0.2-25     
#>   [7] farver_2.1.2           rmarkdown_2.32         fs_2.1.0              
#>  [10] vctrs_0.7.3            memoise_2.0.1          askpass_1.2.1         
#>  [13] base64enc_0.1-6        htmltools_0.5.9        contentanalysis_1.1.1 
#>  [16] curl_8.1.0             broom_1.0.13           janeaustenr_1.0.0     
#>  [19] cellranger_1.1.0       htmlwidgets_1.6.4      tokenizers_0.3.0      
#>  [22] plyr_1.8.9             httr2_1.3.0            cachem_1.1.0          
#>  [25] plotly_4.12.1          dimensionsR_0.0.3      mime_0.13             
#>  [28] lifecycle_1.0.5        pkgconfig_2.0.3        R6_2.6.1              
#>  [31] fastmap_1.2.0          shiny_1.14.0           digest_0.6.39         
#>  [34] patchwork_1.3.2        shinycssloaders_1.1.0  rprojroot_2.1.1       
#>  [37] RSpectra_0.16-2        SnowballC_0.7.1        labeling_0.4.3        
#>  [40] urltools_1.7.3.1       timechange_0.4.0       mgcv_1.9-4            
#>  [43] httr_1.4.9             polyclip_1.10-7        compiler_4.6.1        
#>  [46] here_1.0.2             bit64_4.8.6            withr_3.0.3           
#>  [49] backports_1.5.1        S7_0.2.2               viridis_0.6.5         
#>  [52] ggforce_0.5.0          MASS_7.3-65            bibliometrixData_0.3.0
#>  [55] tools_4.6.1            otel_0.2.0             stopwords_2.3         
#>  [58] zip_3.0.2              rentrez_1.2.4          httpuv_1.6.17         
#>  [61] nlme_3.1-169           promises_1.5.0         grid_4.6.1            
#>  [64] stringdist_0.9.17      reshape2_1.4.5         generics_0.1.4        
#>  [67] gtable_0.3.6           tzdb_0.5.0             rscopus_0.9.0         
#>  [70] ca_0.72                data.table_1.18.6.1    hms_1.1.4             
#>  [73] xml2_1.6.0             utf8_1.2.6             ggrepel_0.9.8         
#>  [76] pillar_1.11.1          nsyllable_1.0.1        vroom_1.7.1           
#>  [79] later_1.4.8            splines_4.6.1          tweenr_2.0.3          
#>  [82] lattice_0.22-9         FNN_1.1.4.1            bit_4.6.0             
#>  [85] tidyselect_1.2.1       tm_0.7-20              miniUI_0.1.2          
#>  [88] knitr_1.52             gridExtra_2.3.1        NLP_0.3-3             
#>  [91] bookdown_0.48          stats4_4.6.1           crul_1.6.0            
#>  [94] xfun_0.62              graphlayouts_1.2.5     matrixStats_1.5.0     
#>  [97] DT_0.34.0              humaniformat_0.6.0     stringi_1.8.9         
#> [100] qpdf_1.4.1             yaml_2.3.12            evaluate_1.0.5        
#> [103] httpcode_0.3.0         cli_3.6.6              xtable_1.8-8          
#> [106] systemfonts_1.3.2      Rcpp_1.1.2             readxl_1.6.0          
#> [109] triebeard_0.4.1        XML_3.99-0.25          parallel_4.6.1        
#> [112] assertthat_0.2.1       pubmedR_1.0.2          slam_0.1-56           
#> [115] viridisLite_0.4.3      scales_1.4.0           crayon_1.5.3          
#> [118] openxlsx_4.2.9         rlang_1.3.0            fastmatch_1.1-8
\end{verbatim}

\chapter{Interactive Dashboards}\label{interactive-dashboards}

\emph{R. Heller · Date: 2026-05-09 · Last changed: 2026-10-31}

\section{Learning objectives}\label{dashboard-objectives}

After completing this chapter, you will be able to:

\begin{itemize}
\tightlist
\item
  Create a Quarto dashboard layout with multiple panels and pages
\item
  Convert ggplot2 figures to interactive plotly charts
\item
  Build searchable, filterable data tables with the DT package
\item
  Design effective dashboard layouts for bibliometric data exploration
\item
  Deploy dashboards as self-contained HTML files
\end{itemize}

\section{Setup}\label{dashboard-setup}

\begin{Shaded}
\begin{Highlighting}[]
\FunctionTok{library}\NormalTok{(tidyverse)}
\FunctionTok{library}\NormalTok{(openalexR)}
\FunctionTok{library}\NormalTok{(plotly)}
\FunctionTok{library}\NormalTok{(DT)}
\FunctionTok{library}\NormalTok{(glue)}

\FunctionTok{set.seed}\NormalTok{(}\DecValTok{20260509}\NormalTok{)}

\FunctionTok{source}\NormalTok{(here}\SpecialCharTok{::}\FunctionTok{here}\NormalTok{(}\StringTok{"R"}\NormalTok{, }\StringTok{"api\_helpers.R"}\NormalTok{))}
\FunctionTok{source}\NormalTok{(here}\SpecialCharTok{::}\FunctionTok{here}\NormalTok{(}\StringTok{"R"}\NormalTok{, }\StringTok{"utils.R"}\NormalTok{))}
\FunctionTok{source}\NormalTok{(here}\SpecialCharTok{::}\FunctionTok{here}\NormalTok{(}\StringTok{"R"}\NormalTok{, }\StringTok{"sci\_palette.R"}\NormalTok{))}
\end{Highlighting}
\end{Shaded}

\section{Conceptual background}\label{dashboard-background}

Bibliometric analysis often produces more data than can be communicated through static figures and tables. Interactive dashboards allow stakeholders --- research managers, librarians, policy makers --- to explore data on their own terms: filtering by year, hovering for details, sorting tables, and drilling down into specific subsets.

\textbf{Quarto dashboards} use the \texttt{format:\ dashboard} output type to arrange content in cards, rows, and columns. They support static and interactive content side by side, and can be deployed as self-contained HTML files (no server needed) or as server-backed applications.

\textbf{plotly} converts ggplot2 figures into interactive JavaScript visualisations with hover tooltips, zooming, panning, and selection. The \texttt{ggplotly()} function provides a zero-effort upgrade path from static to interactive figures, though fine-tuning often requires direct plotly syntax.

\textbf{DT} (DataTables) renders data frames as interactive HTML tables with search, sort, pagination, and filtering. For bibliometric data, DT is invaluable: users can search for specific authors, sort by citation count, and filter by year --- all without writing code.

The key design principle for dashboards is \textbf{progressive disclosure}: show the most important summary first, then let users drill down into details. Avoid overwhelming users with every metric at once.

\section{Worked example}\label{dashboard-example}

\subsection{Preparing dashboard data}\label{preparing-dashboard-data}

\begin{Shaded}
\begin{Highlighting}[]
\NormalTok{works }\OtherTok{\textless{}{-}} \FunctionTok{fetch\_openalex}\NormalTok{(}
  \AttributeTok{entity =} \StringTok{"works"}\NormalTok{,}
  \AttributeTok{primary\_location.source.id =} \StringTok{"S148561398"}\NormalTok{,}
  \AttributeTok{from\_publication\_date =} \StringTok{"2019{-}01{-}01"}\NormalTok{,}
  \AttributeTok{to\_publication\_date =} \StringTok{"2023{-}12{-}31"}\NormalTok{,}
  \AttributeTok{options =} \FunctionTok{list}\NormalTok{(}\AttributeTok{sample =} \DecValTok{500}\NormalTok{, }\AttributeTok{seed =} \DecValTok{42}\NormalTok{)}
\NormalTok{)}

\NormalTok{dash\_data }\OtherTok{\textless{}{-}}\NormalTok{ works }\SpecialCharTok{|\textgreater{}}
  \FunctionTok{transmute}\NormalTok{(}
\NormalTok{    id, }\AttributeTok{title =}\NormalTok{ display\_name,}
    \AttributeTok{year =} \FunctionTok{year}\NormalTok{(publication\_date),}
\NormalTok{    cited\_by\_count, oa\_status, type, source\_display\_name}
\NormalTok{  )}

\FunctionTok{cat}\NormalTok{(}\FunctionTok{glue}\NormalTok{(}\StringTok{"Dashboard data: \{nrow(dash\_data)\} records}\SpecialCharTok{\textbackslash{}n}\StringTok{"}\NormalTok{))}
\end{Highlighting}
\end{Shaded}

\begin{verbatim}
#> Dashboard data: 500 records
\end{verbatim}

\subsection{Interactive trend chart}\label{interactive-trend-chart}

\begin{Shaded}
\begin{Highlighting}[]
\CommentTok{\# Interactive plotly widget — only rendered for HTML output.}
\NormalTok{trend }\OtherTok{\textless{}{-}}\NormalTok{ dash\_data }\SpecialCharTok{|\textgreater{}}
  \FunctionTok{count}\NormalTok{(year)}

\NormalTok{p }\OtherTok{\textless{}{-}} \FunctionTok{ggplot}\NormalTok{(trend, }\FunctionTok{aes}\NormalTok{(}\AttributeTok{x =}\NormalTok{ year, }\AttributeTok{y =}\NormalTok{ n)) }\SpecialCharTok{+}
  \FunctionTok{geom\_line}\NormalTok{(}\AttributeTok{colour =} \FunctionTok{palette\_sci}\NormalTok{(}\DecValTok{1}\NormalTok{), }\AttributeTok{linewidth =} \DecValTok{1}\NormalTok{) }\SpecialCharTok{+}
  \FunctionTok{geom\_point}\NormalTok{(}\AttributeTok{colour =} \FunctionTok{palette\_sci}\NormalTok{(}\DecValTok{1}\NormalTok{), }\AttributeTok{size =} \DecValTok{3}\NormalTok{) }\SpecialCharTok{+}
  \FunctionTok{labs}\NormalTok{(}\AttributeTok{x =} \StringTok{"Year"}\NormalTok{, }\AttributeTok{y =} \StringTok{"Publications"}\NormalTok{) }\SpecialCharTok{+}
  \FunctionTok{theme\_sci}\NormalTok{()}

\FunctionTok{ggplotly}\NormalTok{(p)}
\end{Highlighting}
\end{Shaded}

\subsection{Interactive citation distribution}\label{interactive-citation-distribution}

\begin{Shaded}
\begin{Highlighting}[]
\CommentTok{\# Interactive plotly widget — only rendered for HTML output.}
\NormalTok{p2 }\OtherTok{\textless{}{-}} \FunctionTok{ggplot}\NormalTok{(dash\_data, }\FunctionTok{aes}\NormalTok{(}\AttributeTok{x =}\NormalTok{ cited\_by\_count)) }\SpecialCharTok{+}
  \FunctionTok{geom\_histogram}\NormalTok{(}\AttributeTok{binwidth =} \DecValTok{5}\NormalTok{, }\AttributeTok{fill =} \FunctionTok{palette\_sci}\NormalTok{(}\DecValTok{1}\NormalTok{), }\AttributeTok{colour =} \StringTok{"white"}\NormalTok{) }\SpecialCharTok{+}
  \FunctionTok{labs}\NormalTok{(}\AttributeTok{x =} \StringTok{"Citation count"}\NormalTok{, }\AttributeTok{y =} \StringTok{"Papers"}\NormalTok{) }\SpecialCharTok{+}
  \FunctionTok{theme\_sci}\NormalTok{()}

\FunctionTok{ggplotly}\NormalTok{(p2)}
\end{Highlighting}
\end{Shaded}

\subsection{Searchable data table}\label{searchable-data-table}

\begin{Shaded}
\begin{Highlighting}[]
\CommentTok{\# Interactive DT widget — only rendered for HTML output.}
\NormalTok{dash\_data }\SpecialCharTok{|\textgreater{}}
  \FunctionTok{select}\NormalTok{(title, year, cited\_by\_count, oa\_status) }\SpecialCharTok{|\textgreater{}}
  \FunctionTok{arrange}\NormalTok{(}\FunctionTok{desc}\NormalTok{(cited\_by\_count)) }\SpecialCharTok{|\textgreater{}}
  \FunctionTok{datatable}\NormalTok{(}
    \AttributeTok{options =} \FunctionTok{list}\NormalTok{(}\AttributeTok{pageLength =} \DecValTok{10}\NormalTok{, }\AttributeTok{searchHighlight =} \ConstantTok{TRUE}\NormalTok{),}
    \AttributeTok{filter =} \StringTok{"top"}\NormalTok{,}
    \AttributeTok{rownames =} \ConstantTok{FALSE}
\NormalTok{  )}
\end{Highlighting}
\end{Shaded}

\subsection{OA status breakdown (interactive)}\label{oa-status-breakdown-interactive}

\begin{Shaded}
\begin{Highlighting}[]
\CommentTok{\# Interactive plotly widget — only rendered for HTML output.}
\NormalTok{oa\_trend }\OtherTok{\textless{}{-}}\NormalTok{ dash\_data }\SpecialCharTok{|\textgreater{}}
  \FunctionTok{count}\NormalTok{(year, oa\_status)}

\NormalTok{p3 }\OtherTok{\textless{}{-}} \FunctionTok{ggplot}\NormalTok{(oa\_trend, }\FunctionTok{aes}\NormalTok{(}\AttributeTok{x =} \FunctionTok{factor}\NormalTok{(year), }\AttributeTok{y =}\NormalTok{ n, }\AttributeTok{fill =}\NormalTok{ oa\_status)) }\SpecialCharTok{+}
  \FunctionTok{geom\_col}\NormalTok{() }\SpecialCharTok{+}
  \FunctionTok{scale\_fill\_manual}\NormalTok{(}\AttributeTok{values =} \FunctionTok{palette\_sci}\NormalTok{(}\FunctionTok{n\_distinct}\NormalTok{(oa\_trend}\SpecialCharTok{$}\NormalTok{oa\_status))) }\SpecialCharTok{+}
  \FunctionTok{labs}\NormalTok{(}\AttributeTok{x =} \StringTok{"Year"}\NormalTok{, }\AttributeTok{y =} \StringTok{"Count"}\NormalTok{, }\AttributeTok{fill =} \StringTok{"OA Status"}\NormalTok{) }\SpecialCharTok{+}
  \FunctionTok{theme\_sci}\NormalTok{()}

\FunctionTok{ggplotly}\NormalTok{(p3)}
\end{Highlighting}
\end{Shaded}

\section{Diagnostics and interpretation}\label{dashboard-diagnostics}

\begin{itemize}
\tightlist
\item
  \textbf{Performance}: Large datasets (\textgreater{} 10,000 rows) can make plotly and DT slow. Aggregate or sample data before rendering interactive widgets.
\item
  \textbf{Accessibility}: Interactive figures are not accessible to screen readers. Always provide alt text and consider static fallbacks.
\item
  \textbf{File size}: Self-contained HTML dashboards embed all data and JavaScript libraries. Files can reach 5--50 MB for complex dashboards.
\item
  \textbf{Browser compatibility}: Test in Chrome, Firefox, and Safari. Some plotly features render differently across browsers.
\end{itemize}

\section{Limitations and responsible use}\label{dashboard-limitations}

\section{Limitations and responsible use}\label{limitations-and-responsible-use-34}

\begin{itemize}
\tightlist
\item
  \textbf{Interactivity is not analysis.} A dashboard lets users explore data, but exploration without statistical rigour can produce misleading impressions. Pair dashboards with proper analysis.
\item
  \textbf{Dashboards can mislead.} Default axis ranges, binning, and colour choices affect interpretation. Design dashboards to prevent misreading, not enable it.
\item
  \textbf{Data privacy.} Interactive tables may expose individual-level data (author names, paper titles). Consider whether the dashboard's audience should see this level of detail.
\item
  \textbf{Maintenance burden.} Dashboards that query live APIs need updating when APIs change. Static dashboards based on cached data are more sustainable \citep{hicks2015}.
\end{itemize}

\section{Common pitfalls}\label{dashboard-pitfalls}

\section{Common pitfalls}\label{common-pitfalls-34}

\begin{itemize}
\tightlist
\item
  \textbf{Too many widgets.} A dashboard with 15 interactive panels is overwhelming. Start with 3--5 and add more only if users need them.
\item
  \textbf{Not filtering large datasets.} Rendering 50,000 rows in DT or 100,000 points in plotly will crash the browser. Pre-aggregate.
\item
  \textbf{Ignoring PDF fallbacks.} If your Quarto book renders to PDF, interactive widgets become static. Ensure the static version is still informative.
\item
  \textbf{Using ggplotly without cleanup.} \texttt{ggplotly()} produces verbose tooltips. Customise hover text with \texttt{tooltip} parameter or direct plotly.
\end{itemize}

\section{Exercises}\label{dashboard-exercises}

\begin{enumerate}
\def\labelenumi{\arabic{enumi}.}
\item
  \textbf{Multi-page dashboard.} Create a Quarto dashboard (format: dashboard) with two pages: ``Overview'' (trend + citation distribution) and ``Details'' (searchable table).
\item
  \textbf{Linked brushing.} Use plotly's selection features to highlight papers in a scatter plot and simultaneously filter a DT table below.
\item
  \textbf{Custom tooltips.} Modify the trend chart to show ``Year: 2023, Papers: 87, Mean cites: 12.3'' on hover instead of the default tooltip.
\end{enumerate}

\section{Solutions}\label{35-solutions}

Solutions are provided in \ref{solutions}.

\section{Further reading}\label{dashboard-reading}

\begin{itemize}
\tightlist
\item
  \citet{aria2017} --- \texttt{bibliometrix} and \texttt{biblioshiny} for interactive bibliometric exploration.
\item
  \citet{priem2022} --- OpenAlex as a data source for dashboard construction.
\end{itemize}

\section{Session info}\label{dashboard-sessioninfo}

\begin{Shaded}
\begin{Highlighting}[]
\FunctionTok{sessionInfo}\NormalTok{()}
\end{Highlighting}
\end{Shaded}

\begin{verbatim}
#> R version 4.6.1 (2026-06-24)
#> Platform: x86_64-pc-linux-gnu
#> Running under: Ubuntu 24.04.5 LTS
#> 
#> Matrix products: default
#> BLAS:   /usr/lib/x86_64-linux-gnu/openblas-pthread/libblas.so.3 
#> LAPACK: /usr/lib/x86_64-linux-gnu/openblas-pthread/libopenblasp-r0.3.26.so;  LAPACK version 3.12.0
#> 
#> locale:
#>  [1] LC_CTYPE=C.UTF-8       LC_NUMERIC=C           LC_TIME=C.UTF-8       
#>  [4] LC_COLLATE=C.UTF-8     LC_MONETARY=C.UTF-8    LC_MESSAGES=C.UTF-8   
#>  [7] LC_PAPER=C.UTF-8       LC_NAME=C              LC_ADDRESS=C          
#> [10] LC_TELEPHONE=C         LC_MEASUREMENT=C.UTF-8 LC_IDENTIFICATION=C   
#> 
#> time zone: UTC
#> tzcode source: system (glibc)
#> 
#> attached base packages:
#> [1] stats     graphics  grDevices utils     datasets  methods   base     
#> 
#> other attached packages:
#>  [1] DT_0.34.0                 plotly_4.12.1            
#>  [3] uwot_0.2.5                Matrix_1.7-5             
#>  [5] word2vec_0.4.1            stm_1.3.8                
#>  [7] topicmodels_0.2-17        quanteda.textstats_0.97.2
#>  [9] visNetwork_2.1.4          ggraph_2.2.2             
#> [11] tidygraph_1.3.1           igraph_2.3.4             
#> [13] quanteda_4.5.1            pdftools_3.9.1           
#> [15] arrow_26.0.0              bibliometrix_5.5.0       
#> [17] RefManageR_1.4.0          bib2df_1.1.2.0           
#> [19] rcrossref_1.2.1           gt_1.4.0                 
#> [21] tidytext_0.4.3            glue_1.8.1               
#> [23] openalexR_3.1.2           lubridate_1.9.5          
#> [25] forcats_1.0.1             stringr_1.6.0            
#> [27] dplyr_1.2.1               purrr_1.2.2              
#> [29] readr_2.2.0               tidyr_1.3.2              
#> [31] tibble_3.3.1              ggplot2_4.0.3            
#> [33] tidyverse_2.0.0          
#> 
#> loaded via a namespace (and not attached):
#>   [1] bibtex_0.5.3           RColorBrewer_1.1-3     rstudioapi_0.19.0     
#>   [4] jsonlite_2.0.0         magrittr_2.0.5         modeltools_0.2-25     
#>   [7] farver_2.1.2           rmarkdown_2.32         fs_2.1.0              
#>  [10] vctrs_0.7.3            memoise_2.0.1          askpass_1.2.1         
#>  [13] base64enc_0.1-6        htmltools_0.5.9        contentanalysis_1.1.1 
#>  [16] curl_8.1.0             broom_1.0.13           janeaustenr_1.0.0     
#>  [19] cellranger_1.1.0       htmlwidgets_1.6.4      tokenizers_0.3.0      
#>  [22] plyr_1.8.9             httr2_1.3.0            cachem_1.1.0          
#>  [25] dimensionsR_0.0.3      mime_0.13              lifecycle_1.0.5       
#>  [28] pkgconfig_2.0.3        R6_2.6.1               fastmap_1.2.0         
#>  [31] shiny_1.14.0           digest_0.6.39          patchwork_1.3.2       
#>  [34] shinycssloaders_1.1.0  rprojroot_2.1.1        RSpectra_0.16-2       
#>  [37] SnowballC_0.7.1        labeling_0.4.3         urltools_1.7.3.1      
#>  [40] timechange_0.4.0       mgcv_1.9-4             httr_1.4.9            
#>  [43] polyclip_1.10-7        compiler_4.6.1         here_1.0.2            
#>  [46] bit64_4.8.6            withr_3.0.3            backports_1.5.1       
#>  [49] S7_0.2.2               viridis_0.6.5          ggforce_0.5.0         
#>  [52] MASS_7.3-65            bibliometrixData_0.3.0 tools_4.6.1           
#>  [55] otel_0.2.0             stopwords_2.3          zip_3.0.2             
#>  [58] rentrez_1.2.4          httpuv_1.6.17          nlme_3.1-169          
#>  [61] promises_1.5.0         grid_4.6.1             stringdist_0.9.17     
#>  [64] reshape2_1.4.5         generics_0.1.4         gtable_0.3.6          
#>  [67] tzdb_0.5.0             rscopus_0.9.0          ca_0.72               
#>  [70] data.table_1.18.6.1    hms_1.1.4              xml2_1.6.0            
#>  [73] utf8_1.2.6             ggrepel_0.9.8          pillar_1.11.1         
#>  [76] nsyllable_1.0.1        vroom_1.7.1            later_1.4.8           
#>  [79] splines_4.6.1          tweenr_2.0.3           lattice_0.22-9        
#>  [82] FNN_1.1.4.1            bit_4.6.0              tidyselect_1.2.1      
#>  [85] tm_0.7-20              miniUI_0.1.2           knitr_1.52            
#>  [88] gridExtra_2.3.1        NLP_0.3-3              bookdown_0.48         
#>  [91] stats4_4.6.1           crul_1.6.0             xfun_0.62             
#>  [94] graphlayouts_1.2.5     matrixStats_1.5.0      humaniformat_0.6.0    
#>  [97] stringi_1.8.9          qpdf_1.4.1             yaml_2.3.12           
#> [100] evaluate_1.0.5         httpcode_0.3.0         cli_3.6.6             
#> [103] xtable_1.8-8           systemfonts_1.3.2      Rcpp_1.1.2            
#> [106] readxl_1.6.0           triebeard_0.4.1        XML_3.99-0.25         
#> [109] parallel_4.6.1         assertthat_0.2.1       pubmedR_1.0.2         
#> [112] slam_0.1-56            viridisLite_0.4.3      scales_1.4.0          
#> [115] crayon_1.5.3           openxlsx_4.2.9         rlang_1.3.0           
#> [118] fastmatch_1.1-8
\end{verbatim}

\chapter{Parameterized Quarto Reports}\label{parameterized-quarto-reports}

\emph{R. Heller · Date: 2026-05-09 · Last changed: 2026-10-31}

\section{Learning objectives}\label{paramreport-objectives}

After completing this chapter, you will be able to:

\begin{itemize}
\tightlist
\item
  Define parameters in a Quarto document's YAML header
\item
  Access parameters in R code chunks to drive data retrieval and analysis
\item
  Render parameterized reports from the command line and from scripts
\item
  Build a template for journal-level bibliometric profiles
\item
  Generate multiple reports in batch using purrr or shell loops
\end{itemize}

\section{Setup}\label{paramreport-setup}

\begin{Shaded}
\begin{Highlighting}[]
\FunctionTok{library}\NormalTok{(tidyverse)}
\FunctionTok{library}\NormalTok{(openalexR)}
\FunctionTok{library}\NormalTok{(glue)}
\FunctionTok{library}\NormalTok{(gt)}

\FunctionTok{set.seed}\NormalTok{(}\DecValTok{20260509}\NormalTok{)}

\FunctionTok{source}\NormalTok{(here}\SpecialCharTok{::}\FunctionTok{here}\NormalTok{(}\StringTok{"R"}\NormalTok{, }\StringTok{"api\_helpers.R"}\NormalTok{))}
\FunctionTok{source}\NormalTok{(here}\SpecialCharTok{::}\FunctionTok{here}\NormalTok{(}\StringTok{"R"}\NormalTok{, }\StringTok{"utils.R"}\NormalTok{))}
\FunctionTok{source}\NormalTok{(here}\SpecialCharTok{::}\FunctionTok{here}\NormalTok{(}\StringTok{"R"}\NormalTok{, }\StringTok{"sci\_palette.R"}\NormalTok{))}
\end{Highlighting}
\end{Shaded}

\section{Conceptual background}\label{paramreport-background}

Many bibliometric tasks are repetitive: generating the same analysis for 10 journals, 5 departments, or 20 funders. Writing separate scripts for each entity is error-prone and hard to maintain. \textbf{Parameterized reports} solve this by defining variables (parameters) that change the report's content without changing its structure.

Quarto supports parameters via the \texttt{params} field in the YAML header. Parameters can be strings, numbers, or lists. Inside the document, \texttt{params\$journal\_id} (or \texttt{params\$institution}) drives data fetching, filtering, and labelling. At render time, parameter values are passed via the command line or a rendering script.

This approach scales naturally: a loop or \texttt{purrr::walk()} call renders dozens of reports with different parameters, producing a complete portfolio of bibliometric profiles with consistent methodology and formatting.

Parameterized reports are especially valuable for institutional reporting (annual bibliometric profiles for each department), journal analytics (quarterly reports for a publisher's journal portfolio), and funding assessment (per-grant output summaries).

\section{Worked example}\label{paramreport-example}

\subsection{Defining parameters}\label{defining-parameters}

A parameterized Quarto document starts with parameters in the YAML header:

\begin{Shaded}
\begin{Highlighting}[]
\CommentTok{\# Example YAML header (not executable R):}
\CommentTok{\#}
\CommentTok{\# {-}{-}{-}}
\CommentTok{\# title: "Journal Bibliometric Profile"}
\CommentTok{\# params:}
\CommentTok{\#   journal\_id: "S148561398"}
\CommentTok{\#   journal\_name: "Scientometrics"}
\CommentTok{\#   start\_year: 2020}
\CommentTok{\#   end\_year: 2023}
\CommentTok{\# {-}{-}{-}}
\end{Highlighting}
\end{Shaded}

\subsection{Using parameters in code}\label{using-parameters-in-code}

\begin{Shaded}
\begin{Highlighting}[]
\NormalTok{journal\_id }\OtherTok{\textless{}{-}} \StringTok{"S148561398"}
\NormalTok{journal\_name }\OtherTok{\textless{}{-}} \StringTok{"Scientometrics"}
\NormalTok{start\_year }\OtherTok{\textless{}{-}} \DecValTok{2020}
\NormalTok{end\_year }\OtherTok{\textless{}{-}} \DecValTok{2023}
\end{Highlighting}
\end{Shaded}

\begin{Shaded}
\begin{Highlighting}[]
\NormalTok{works }\OtherTok{\textless{}{-}} \FunctionTok{fetch\_openalex}\NormalTok{(}
  \AttributeTok{entity =} \StringTok{"works"}\NormalTok{,}
  \AttributeTok{primary\_location.source.id =}\NormalTok{ journal\_id,}
  \AttributeTok{from\_publication\_date =} \FunctionTok{paste0}\NormalTok{(start\_year, }\StringTok{"{-}01{-}01"}\NormalTok{),}
  \AttributeTok{to\_publication\_date =} \FunctionTok{paste0}\NormalTok{(end\_year, }\StringTok{"{-}12{-}31"}\NormalTok{),}
  \AttributeTok{type =} \StringTok{"article"}\NormalTok{,}
  \AttributeTok{options =} \FunctionTok{list}\NormalTok{(}\AttributeTok{sample =} \DecValTok{300}\NormalTok{, }\AttributeTok{seed =} \DecValTok{42}\NormalTok{)}
\NormalTok{)}

\FunctionTok{cat}\NormalTok{(}\FunctionTok{glue}\NormalTok{(}\StringTok{"Journal: \{journal\_name\}}\SpecialCharTok{\textbackslash{}n}\StringTok{"}\NormalTok{))}
\end{Highlighting}
\end{Shaded}

\begin{verbatim}
#> Journal: Scientometrics
\end{verbatim}

\begin{Shaded}
\begin{Highlighting}[]
\FunctionTok{cat}\NormalTok{(}\FunctionTok{glue}\NormalTok{(}\StringTok{"Period: \{start\_year\}{-}{-}\{end\_year\}}\SpecialCharTok{\textbackslash{}n}\StringTok{"}\NormalTok{))}
\end{Highlighting}
\end{Shaded}

\begin{verbatim}
#> Period: 2020--2023
\end{verbatim}

\begin{Shaded}
\begin{Highlighting}[]
\FunctionTok{cat}\NormalTok{(}\FunctionTok{glue}\NormalTok{(}\StringTok{"Articles retrieved: \{nrow(works)\}}\SpecialCharTok{\textbackslash{}n}\StringTok{"}\NormalTok{))}
\end{Highlighting}
\end{Shaded}

\begin{verbatim}
#> Articles retrieved: 300
\end{verbatim}

\subsection{Automated summary statistics}\label{automated-summary-statistics}

\begin{Shaded}
\begin{Highlighting}[]
\NormalTok{summary\_stats }\OtherTok{\textless{}{-}}\NormalTok{ works }\SpecialCharTok{|\textgreater{}}
  \FunctionTok{summarise}\NormalTok{(}
    \AttributeTok{total\_articles =} \FunctionTok{n}\NormalTok{(),}
    \AttributeTok{mean\_citations =} \FunctionTok{round}\NormalTok{(}\FunctionTok{mean}\NormalTok{(cited\_by\_count), }\DecValTok{1}\NormalTok{),}
    \AttributeTok{median\_citations =} \FunctionTok{median}\NormalTok{(cited\_by\_count),}
    \AttributeTok{h\_index =} \FunctionTok{compute\_h\_index}\NormalTok{(cited\_by\_count),}
    \AttributeTok{oa\_rate =} \FunctionTok{round}\NormalTok{(}\FunctionTok{mean}\NormalTok{(oa\_status }\SpecialCharTok{!=} \StringTok{"closed"}\NormalTok{, }\AttributeTok{na.rm =} \ConstantTok{TRUE}\NormalTok{) }\SpecialCharTok{*} \DecValTok{100}\NormalTok{, }\DecValTok{1}\NormalTok{)}
\NormalTok{  ) }\SpecialCharTok{|\textgreater{}}
  \FunctionTok{mutate}\NormalTok{(}\FunctionTok{across}\NormalTok{(}\FunctionTok{everything}\NormalTok{(), as.character)) }\SpecialCharTok{|\textgreater{}}
  \FunctionTok{pivot\_longer}\NormalTok{(}\FunctionTok{everything}\NormalTok{(), }\AttributeTok{names\_to =} \StringTok{"Metric"}\NormalTok{, }\AttributeTok{values\_to =} \StringTok{"Value"}\NormalTok{)}

\NormalTok{summary\_stats }\SpecialCharTok{|\textgreater{}} \FunctionTok{gt}\NormalTok{()}
\end{Highlighting}
\end{Shaded}

\begingroup
\fontsize{8.0pt}{9.0pt}\selectfont
\begin{longtable}{>{\raggedright\arraybackslash}p{\dimexpr 0.59\linewidth -2\tabcolsep-1.5\arrayrulewidth}>{\raggedleft\arraybackslash}p{\dimexpr 0.41\linewidth -2\tabcolsep-1.5\arrayrulewidth}}
\toprule
Metric & Value \\ 
\midrule\addlinespace[2.5pt]
\nolinkurl{total_articles} & 300 \\ 
\nolinkurl{mean_citations} & 25 \\ 
\nolinkurl{median_citations} & 15 \\ 
h\_index & 40 \\ 
oa\_rate & 52 \\ 
\bottomrule
\end{longtable}
\endgroup

\subsection{Parameterized figure}\label{parameterized-figure}

\begin{Shaded}
\begin{Highlighting}[]
\NormalTok{works }\SpecialCharTok{|\textgreater{}}
  \FunctionTok{mutate}\NormalTok{(}\AttributeTok{year =} \FunctionTok{year}\NormalTok{(publication\_date)) }\SpecialCharTok{|\textgreater{}}
  \FunctionTok{count}\NormalTok{(year) }\SpecialCharTok{|\textgreater{}}
  \FunctionTok{ggplot}\NormalTok{(}\FunctionTok{aes}\NormalTok{(}\AttributeTok{x =} \FunctionTok{factor}\NormalTok{(year), }\AttributeTok{y =}\NormalTok{ n)) }\SpecialCharTok{+}
  \FunctionTok{geom\_col}\NormalTok{(}\AttributeTok{fill =} \FunctionTok{palette\_sci}\NormalTok{(}\DecValTok{1}\NormalTok{)) }\SpecialCharTok{+}
  \FunctionTok{labs}\NormalTok{(}\AttributeTok{x =} \StringTok{"Year"}\NormalTok{, }\AttributeTok{y =} \StringTok{"Articles"}\NormalTok{,}
       \AttributeTok{title =} \FunctionTok{glue}\NormalTok{(}\StringTok{"\{journal\_name\} — Annual Output"}\NormalTok{)) }\SpecialCharTok{+}
  \FunctionTok{theme\_sci}\NormalTok{()}
\end{Highlighting}
\end{Shaded}

\begin{figure}

{\centering \includegraphics[width=0.9\linewidth,alt={Bar chart showing the number of articles published per year in the journal specified by the report parameters.}]{scientometrics-in-r_files/figure-latex/36-fig-paramreport-trend-1} 

}

\caption{Annual publication count for the selected journal.}\label{fig:36-fig-paramreport-trend}
\end{figure}

\subsection{Batch rendering}\label{batch-rendering}

\begin{Shaded}
\begin{Highlighting}[]
\CommentTok{\# Render for multiple journals:}
\NormalTok{journals }\OtherTok{\textless{}{-}} \FunctionTok{tibble}\NormalTok{(}
  \AttributeTok{journal\_id =} \FunctionTok{c}\NormalTok{(}\StringTok{"S148561398"}\NormalTok{, }\StringTok{"S202381698"}\NormalTok{, }\StringTok{"S137773608"}\NormalTok{),}
  \AttributeTok{journal\_name =} \FunctionTok{c}\NormalTok{(}\StringTok{"Scientometrics"}\NormalTok{, }\StringTok{"PLOS ONE"}\NormalTok{, }\StringTok{"Nature"}\NormalTok{)}
\NormalTok{)}

\FunctionTok{walk2}\NormalTok{(journals}\SpecialCharTok{$}\NormalTok{journal\_id, journals}\SpecialCharTok{$}\NormalTok{journal\_name, }\ControlFlowTok{function}\NormalTok{(jid, jname) \{}
\NormalTok{  quarto}\SpecialCharTok{::}\FunctionTok{quarto\_render}\NormalTok{(}
    \AttributeTok{input =} \StringTok{"journal\_profile.qmd"}\NormalTok{,}
    \AttributeTok{execute\_params =} \FunctionTok{list}\NormalTok{(}
      \AttributeTok{journal\_id =}\NormalTok{ jid,}
      \AttributeTok{journal\_name =}\NormalTok{ jname,}
      \AttributeTok{start\_year =} \DecValTok{2020}\NormalTok{,}
      \AttributeTok{end\_year =} \DecValTok{2023}
\NormalTok{    ),}
    \AttributeTok{output\_file =} \FunctionTok{glue}\NormalTok{(}\StringTok{"reports/\{jname\}\_profile.html"}\NormalTok{)}
\NormalTok{  )}
\NormalTok{\})}
\end{Highlighting}
\end{Shaded}

\section{Diagnostics and interpretation}\label{paramreport-diagnostics}

\begin{itemize}
\tightlist
\item
  \textbf{Parameter validation}: Check that parameters produce valid API queries. Invalid journal IDs return empty results silently.
\item
  \textbf{Sample consistency}: When using \texttt{options\ =\ list(sample\ =\ N)}, the sample changes if the underlying corpus changes. Use \texttt{seed} for stability.
\item
  \textbf{Output naming}: Automated file naming should avoid special characters. Sanitise journal names before using them in file paths.
\item
  \textbf{Error handling}: Batch rendering should catch and log errors for individual reports without stopping the entire batch.
\end{itemize}

\section{Limitations and responsible use}\label{paramreport-limitations}

\section{Limitations and responsible use}\label{limitations-and-responsible-use-35}

\begin{itemize}
\tightlist
\item
  \textbf{Automated ≠ mindless.} Parameterized reports produce consistent output but still require human review. An automated report with a data error will produce the wrong results consistently across all entities.
\item
  \textbf{One template fits all?} Different entities may need different analyses. A template designed for a large journal may produce misleading results for a small one (e.g., sparse yearly bins).
\item
  \textbf{API costs.} Batch rendering for 100 journals makes 100 API calls. Respect rate limits and cache results.
\item
  \textbf{Report fatigue.} Generating too many reports can overwhelm decision-makers. Produce reports for entities that need them, not for every possible parameter value \citep{hicks2015}.
\end{itemize}

\section{Common pitfalls}\label{paramreport-pitfalls}

\section{Common pitfalls}\label{common-pitfalls-35}

\begin{itemize}
\tightlist
\item
  \textbf{Hardcoding values that should be parameters.} If you find yourself editing a date or ID in the script, it should be a parameter.
\item
  \textbf{Not testing with extreme parameters.} A journal with 5 papers per year needs different visualisations than one with 5,000. Test edge cases.
\item
  \textbf{Rendering without caching.} Each render re-executes all code. Use Quarto's \texttt{cache:\ true} or file-based caching to avoid redundant API calls.
\item
  \textbf{Forgetting output format compatibility.} Interactive widgets work in HTML but not PDF. Ensure your template degrades gracefully.
\end{itemize}

\section{Exercises}\label{paramreport-exercises}

\begin{enumerate}
\def\labelenumi{\arabic{enumi}.}
\item
  \textbf{Institutional template.} Create a parameterized Quarto document for institutional profiles. Parameters: institution ID, name, start year, end year. Include publication count, citation distribution, and top authors.
\item
  \textbf{Batch rendering.} Render the journal profile template for 3 different journals. Compare the resulting reports side by side.
\item
  \textbf{Conditional content.} Add conditional logic: if the sample has fewer than 50 papers, display a warning message instead of the citation distribution figure.
\end{enumerate}

\section{Solutions}\label{36-solutions}

Solutions are provided in \ref{solutions}.

\section{Further reading}\label{paramreport-reading}

\begin{itemize}
\tightlist
\item
  \citet{aria2017} --- \texttt{bibliometrix} reporting features.
\item
  \citet{priem2022} --- OpenAlex API for programmatic data retrieval in reports.
\end{itemize}

\section{Session info}\label{paramreport-sessioninfo}

\begin{Shaded}
\begin{Highlighting}[]
\FunctionTok{sessionInfo}\NormalTok{()}
\end{Highlighting}
\end{Shaded}

\begin{verbatim}
#> R version 4.6.1 (2026-06-24)
#> Platform: x86_64-pc-linux-gnu
#> Running under: Ubuntu 24.04.5 LTS
#> 
#> Matrix products: default
#> BLAS:   /usr/lib/x86_64-linux-gnu/openblas-pthread/libblas.so.3 
#> LAPACK: /usr/lib/x86_64-linux-gnu/openblas-pthread/libopenblasp-r0.3.26.so;  LAPACK version 3.12.0
#> 
#> locale:
#>  [1] LC_CTYPE=C.UTF-8       LC_NUMERIC=C           LC_TIME=C.UTF-8       
#>  [4] LC_COLLATE=C.UTF-8     LC_MONETARY=C.UTF-8    LC_MESSAGES=C.UTF-8   
#>  [7] LC_PAPER=C.UTF-8       LC_NAME=C              LC_ADDRESS=C          
#> [10] LC_TELEPHONE=C         LC_MEASUREMENT=C.UTF-8 LC_IDENTIFICATION=C   
#> 
#> time zone: UTC
#> tzcode source: system (glibc)
#> 
#> attached base packages:
#> [1] stats     graphics  grDevices utils     datasets  methods   base     
#> 
#> other attached packages:
#>  [1] DT_0.34.0                 plotly_4.12.1            
#>  [3] uwot_0.2.5                Matrix_1.7-5             
#>  [5] word2vec_0.4.1            stm_1.3.8                
#>  [7] topicmodels_0.2-17        quanteda.textstats_0.97.2
#>  [9] visNetwork_2.1.4          ggraph_2.2.2             
#> [11] tidygraph_1.3.1           igraph_2.3.4             
#> [13] quanteda_4.5.1            pdftools_3.9.1           
#> [15] arrow_26.0.0              bibliometrix_5.5.0       
#> [17] RefManageR_1.4.0          bib2df_1.1.2.0           
#> [19] rcrossref_1.2.1           gt_1.4.0                 
#> [21] tidytext_0.4.3            glue_1.8.1               
#> [23] openalexR_3.1.2           lubridate_1.9.5          
#> [25] forcats_1.0.1             stringr_1.6.0            
#> [27] dplyr_1.2.1               purrr_1.2.2              
#> [29] readr_2.2.0               tidyr_1.3.2              
#> [31] tibble_3.3.1              ggplot2_4.0.3            
#> [33] tidyverse_2.0.0          
#> 
#> loaded via a namespace (and not attached):
#>   [1] bibtex_0.5.3           RColorBrewer_1.1-3     rstudioapi_0.19.0     
#>   [4] jsonlite_2.0.0         magrittr_2.0.5         modeltools_0.2-25     
#>   [7] farver_2.1.2           rmarkdown_2.32         fs_2.1.0              
#>  [10] vctrs_0.7.3            memoise_2.0.1          askpass_1.2.1         
#>  [13] base64enc_0.1-6        htmltools_0.5.9        contentanalysis_1.1.1 
#>  [16] curl_8.1.0             broom_1.0.13           janeaustenr_1.0.0     
#>  [19] cellranger_1.1.0       htmlwidgets_1.6.4      tokenizers_0.3.0      
#>  [22] plyr_1.8.9             httr2_1.3.0            cachem_1.1.0          
#>  [25] dimensionsR_0.0.3      mime_0.13              lifecycle_1.0.5       
#>  [28] pkgconfig_2.0.3        R6_2.6.1               fastmap_1.2.0         
#>  [31] shiny_1.14.0           digest_0.6.39          patchwork_1.3.2       
#>  [34] shinycssloaders_1.1.0  rprojroot_2.1.1        RSpectra_0.16-2       
#>  [37] SnowballC_0.7.1        labeling_0.4.3         urltools_1.7.3.1      
#>  [40] timechange_0.4.0       mgcv_1.9-4             httr_1.4.9            
#>  [43] polyclip_1.10-7        compiler_4.6.1         here_1.0.2            
#>  [46] bit64_4.8.6            withr_3.0.3            backports_1.5.1       
#>  [49] S7_0.2.2               viridis_0.6.5          ggforce_0.5.0         
#>  [52] MASS_7.3-65            bibliometrixData_0.3.0 tools_4.6.1           
#>  [55] otel_0.2.0             stopwords_2.3          zip_3.0.2             
#>  [58] rentrez_1.2.4          httpuv_1.6.17          nlme_3.1-169          
#>  [61] promises_1.5.0         grid_4.6.1             stringdist_0.9.17     
#>  [64] reshape2_1.4.5         generics_0.1.4         gtable_0.3.6          
#>  [67] tzdb_0.5.0             rscopus_0.9.0          ca_0.72               
#>  [70] data.table_1.18.6.1    hms_1.1.4              xml2_1.6.0            
#>  [73] utf8_1.2.6             ggrepel_0.9.8          pillar_1.11.1         
#>  [76] nsyllable_1.0.1        vroom_1.7.1            later_1.4.8           
#>  [79] splines_4.6.1          tweenr_2.0.3           lattice_0.22-9        
#>  [82] FNN_1.1.4.1            bit_4.6.0              tidyselect_1.2.1      
#>  [85] tm_0.7-20              miniUI_0.1.2           knitr_1.52            
#>  [88] gridExtra_2.3.1        NLP_0.3-3              bookdown_0.48         
#>  [91] stats4_4.6.1           crul_1.6.0             xfun_0.62             
#>  [94] graphlayouts_1.2.5     matrixStats_1.5.0      humaniformat_0.6.0    
#>  [97] stringi_1.8.9          qpdf_1.4.1             yaml_2.3.12           
#> [100] evaluate_1.0.5         httpcode_0.3.0         cli_3.6.6             
#> [103] xtable_1.8-8           systemfonts_1.3.2      Rcpp_1.1.2            
#> [106] readxl_1.6.0           triebeard_0.4.1        XML_3.99-0.25         
#> [109] parallel_4.6.1         assertthat_0.2.1       pubmedR_1.0.2         
#> [112] slam_0.1-56            viridisLite_0.4.3      scales_1.4.0          
#> [115] crayon_1.5.3           openxlsx_4.2.9         rlang_1.3.0           
#> [118] fastmatch_1.1-8
\end{verbatim}

\chapter{Building a Bibliometric Shiny App}\label{building-a-bibliometric-shiny-app}

\emph{R. Heller · Date: 2026-05-09 · Last changed: 2026-10-31}

\section{Learning objectives}\label{shiny-objectives}

After completing this chapter, you will be able to:

\begin{itemize}
\tightlist
\item
  Design a Shiny app architecture for bibliometric analysis
\item
  Implement reactive data fetching from OpenAlex
\item
  Build interactive UI components: dropdowns, sliders, and tabbed panels
\item
  Create server-side logic for citation analysis and network visualisation
\item
  Deploy a Shiny app to shinyapps.io or a self-hosted server
\end{itemize}

\section{Setup}\label{shiny-setup}

\begin{Shaded}
\begin{Highlighting}[]
\FunctionTok{library}\NormalTok{(tidyverse)}
\FunctionTok{library}\NormalTok{(shiny)}
\FunctionTok{library}\NormalTok{(openalexR)}
\FunctionTok{library}\NormalTok{(glue)}

\FunctionTok{set.seed}\NormalTok{(}\DecValTok{20260509}\NormalTok{)}

\FunctionTok{source}\NormalTok{(here}\SpecialCharTok{::}\FunctionTok{here}\NormalTok{(}\StringTok{"R"}\NormalTok{, }\StringTok{"sci\_palette.R"}\NormalTok{))}
\end{Highlighting}
\end{Shaded}

\section{Conceptual background}\label{shiny-background}

Shiny transforms R scripts into web applications. Users interact through widgets (dropdowns, sliders, buttons) and see results update in real time. For bibliometrics, Shiny apps serve two audiences: researchers who want to explore data without writing code, and decision-makers who need accessible analytics dashboards.

A Shiny app has two components: the \textbf{UI} (user interface, defining layout and inputs) and the \textbf{server} (the R logic that reacts to user input and generates output). \textbf{Reactive programming} connects inputs to outputs: when a user changes a dropdown, only the computations that depend on that input are re-executed.

For bibliometric apps, the typical architecture involves:

\begin{enumerate}
\def\labelenumi{\arabic{enumi}.}
\tightlist
\item
  \textbf{Input panel}: Select a journal, institution, or search query; choose date range and sample size.
\item
  \textbf{Data fetching}: Reactively query OpenAlex based on user selections.
\item
  \textbf{Analysis tabs}: Trend charts, citation distributions, author rankings, and optionally network visualisation.
\item
  \textbf{Export}: Download buttons for data (CSV) and figures (PNG).
\end{enumerate}

\textbf{\texttt{biblioshiny()}} from the \texttt{bibliometrix} package \citep{aria2017} provides a ready-made Shiny interface for many standard analyses. Building a custom app gives full control over the analysis pipeline and user experience.

Deployment options include shinyapps.io (hosted by Posit, free tier available), Posit Connect (enterprise), and self-hosted via Shiny Server (open source).

\section{Worked example}\label{shiny-example}

\subsection{App structure}\label{app-structure}

We demonstrate the key components of a bibliometric Shiny app. The code below shows the UI and server logic separately for clarity.

\subsection{UI definition}\label{ui-definition}

\begin{Shaded}
\begin{Highlighting}[]
\NormalTok{ui }\OtherTok{\textless{}{-}} \FunctionTok{fluidPage}\NormalTok{(}
  \FunctionTok{titlePanel}\NormalTok{(}\StringTok{"Bibliometric Explorer"}\NormalTok{),}

  \FunctionTok{sidebarLayout}\NormalTok{(}
    \FunctionTok{sidebarPanel}\NormalTok{(}
      \FunctionTok{textInput}\NormalTok{(}\StringTok{"journal\_id"}\NormalTok{, }\StringTok{"OpenAlex Source ID:"}\NormalTok{,}
                \AttributeTok{value =} \StringTok{"S148561398"}\NormalTok{),}
      \FunctionTok{sliderInput}\NormalTok{(}\StringTok{"year\_range"}\NormalTok{, }\StringTok{"Publication Years:"}\NormalTok{,}
                  \AttributeTok{min =} \DecValTok{2015}\NormalTok{, }\AttributeTok{max =} \DecValTok{2023}\NormalTok{, }\AttributeTok{value =} \FunctionTok{c}\NormalTok{(}\DecValTok{2020}\NormalTok{, }\DecValTok{2023}\NormalTok{)),}
      \FunctionTok{numericInput}\NormalTok{(}\StringTok{"sample\_size"}\NormalTok{, }\StringTok{"Sample Size:"}\NormalTok{, }\AttributeTok{value =} \DecValTok{200}\NormalTok{,}
                   \AttributeTok{min =} \DecValTok{50}\NormalTok{, }\AttributeTok{max =} \DecValTok{1000}\NormalTok{),}
      \FunctionTok{actionButton}\NormalTok{(}\StringTok{"fetch"}\NormalTok{, }\StringTok{"Fetch Data"}\NormalTok{, }\AttributeTok{class =} \StringTok{"btn{-}primary"}\NormalTok{),}
      \FunctionTok{hr}\NormalTok{(),}
      \FunctionTok{downloadButton}\NormalTok{(}\StringTok{"download\_csv"}\NormalTok{, }\StringTok{"Download CSV"}\NormalTok{)}
\NormalTok{    ),}

    \FunctionTok{mainPanel}\NormalTok{(}
      \FunctionTok{tabsetPanel}\NormalTok{(}
        \FunctionTok{tabPanel}\NormalTok{(}\StringTok{"Trends"}\NormalTok{,}
                 \FunctionTok{plotOutput}\NormalTok{(}\StringTok{"trend\_plot"}\NormalTok{),}
                 \FunctionTok{tableOutput}\NormalTok{(}\StringTok{"summary\_table"}\NormalTok{)),}
        \FunctionTok{tabPanel}\NormalTok{(}\StringTok{"Citations"}\NormalTok{,}
                 \FunctionTok{plotOutput}\NormalTok{(}\StringTok{"cite\_dist"}\NormalTok{),}
                 \FunctionTok{plotOutput}\NormalTok{(}\StringTok{"cite\_boxplot"}\NormalTok{)),}
        \FunctionTok{tabPanel}\NormalTok{(}\StringTok{"Data"}\NormalTok{,}
\NormalTok{                 DT}\SpecialCharTok{::}\FunctionTok{dataTableOutput}\NormalTok{(}\StringTok{"data\_table"}\NormalTok{))}
\NormalTok{      )}
\NormalTok{    )}
\NormalTok{  )}
\NormalTok{)}
\end{Highlighting}
\end{Shaded}

\subsection{Server logic}\label{server-logic}

\begin{Shaded}
\begin{Highlighting}[]
\NormalTok{server }\OtherTok{\textless{}{-}} \ControlFlowTok{function}\NormalTok{(input, output, session) \{}

\NormalTok{  works\_data }\OtherTok{\textless{}{-}} \FunctionTok{eventReactive}\NormalTok{(input}\SpecialCharTok{$}\NormalTok{fetch, \{}
    \FunctionTok{withProgress}\NormalTok{(}\AttributeTok{message =} \StringTok{"Fetching from OpenAlex..."}\NormalTok{, \{}
      \FunctionTok{fetch\_openalex}\NormalTok{(}
        \AttributeTok{entity =} \StringTok{"works"}\NormalTok{,}
        \AttributeTok{primary\_location.source.id =}\NormalTok{ input}\SpecialCharTok{$}\NormalTok{journal\_id,}
        \AttributeTok{from\_publication\_date =} \FunctionTok{paste0}\NormalTok{(input}\SpecialCharTok{$}\NormalTok{year\_range[}\DecValTok{1}\NormalTok{], }\StringTok{"{-}01{-}01"}\NormalTok{),}
        \AttributeTok{to\_publication\_date =} \FunctionTok{paste0}\NormalTok{(input}\SpecialCharTok{$}\NormalTok{year\_range[}\DecValTok{2}\NormalTok{], }\StringTok{"{-}12{-}31"}\NormalTok{),}
        \AttributeTok{type =} \StringTok{"article"}\NormalTok{,}
        \AttributeTok{options =} \FunctionTok{list}\NormalTok{(}\AttributeTok{sample =}\NormalTok{ input}\SpecialCharTok{$}\NormalTok{sample\_size, }\AttributeTok{seed =} \DecValTok{42}\NormalTok{)}
\NormalTok{      ) }\SpecialCharTok{|\textgreater{}}
        \FunctionTok{transmute}\NormalTok{(}
\NormalTok{          id, }\AttributeTok{title =}\NormalTok{ display\_name,}
          \AttributeTok{year =} \FunctionTok{year}\NormalTok{(publication\_date),}
\NormalTok{          cited\_by\_count, oa\_status}
\NormalTok{        )}
\NormalTok{    \})}
\NormalTok{  \})}

\NormalTok{  output}\SpecialCharTok{$}\NormalTok{trend\_plot }\OtherTok{\textless{}{-}} \FunctionTok{renderPlot}\NormalTok{(\{}
    \FunctionTok{req}\NormalTok{(}\FunctionTok{works\_data}\NormalTok{())}
    \FunctionTok{works\_data}\NormalTok{() }\SpecialCharTok{|\textgreater{}}
      \FunctionTok{count}\NormalTok{(year) }\SpecialCharTok{|\textgreater{}}
      \FunctionTok{ggplot}\NormalTok{(}\FunctionTok{aes}\NormalTok{(}\AttributeTok{x =} \FunctionTok{factor}\NormalTok{(year), }\AttributeTok{y =}\NormalTok{ n)) }\SpecialCharTok{+}
      \FunctionTok{geom\_col}\NormalTok{(}\AttributeTok{fill =} \FunctionTok{palette\_sci}\NormalTok{(}\DecValTok{1}\NormalTok{)) }\SpecialCharTok{+}
      \FunctionTok{labs}\NormalTok{(}\AttributeTok{x =} \StringTok{"Year"}\NormalTok{, }\AttributeTok{y =} \StringTok{"Publications"}\NormalTok{) }\SpecialCharTok{+}
      \FunctionTok{theme\_sci}\NormalTok{()}
\NormalTok{  \})}

\NormalTok{  output}\SpecialCharTok{$}\NormalTok{cite\_dist }\OtherTok{\textless{}{-}} \FunctionTok{renderPlot}\NormalTok{(\{}
    \FunctionTok{req}\NormalTok{(}\FunctionTok{works\_data}\NormalTok{())}
    \FunctionTok{ggplot}\NormalTok{(}\FunctionTok{works\_data}\NormalTok{(), }\FunctionTok{aes}\NormalTok{(}\AttributeTok{x =}\NormalTok{ cited\_by\_count)) }\SpecialCharTok{+}
      \FunctionTok{geom\_histogram}\NormalTok{(}\AttributeTok{binwidth =} \DecValTok{5}\NormalTok{, }\AttributeTok{fill =} \FunctionTok{palette\_sci}\NormalTok{(}\DecValTok{1}\NormalTok{), }\AttributeTok{colour =} \StringTok{"white"}\NormalTok{) }\SpecialCharTok{+}
      \FunctionTok{labs}\NormalTok{(}\AttributeTok{x =} \StringTok{"Citations"}\NormalTok{, }\AttributeTok{y =} \StringTok{"Count"}\NormalTok{) }\SpecialCharTok{+}
      \FunctionTok{theme\_sci}\NormalTok{()}
\NormalTok{  \})}

\NormalTok{  output}\SpecialCharTok{$}\NormalTok{cite\_boxplot }\OtherTok{\textless{}{-}} \FunctionTok{renderPlot}\NormalTok{(\{}
    \FunctionTok{req}\NormalTok{(}\FunctionTok{works\_data}\NormalTok{())}
    \FunctionTok{works\_data}\NormalTok{() }\SpecialCharTok{|\textgreater{}}
      \FunctionTok{filter}\NormalTok{(}\SpecialCharTok{!}\FunctionTok{is.na}\NormalTok{(oa\_status)) }\SpecialCharTok{|\textgreater{}}
      \FunctionTok{ggplot}\NormalTok{(}\FunctionTok{aes}\NormalTok{(}\AttributeTok{x =}\NormalTok{ oa\_status, }\AttributeTok{y =}\NormalTok{ cited\_by\_count }\SpecialCharTok{+} \DecValTok{1}\NormalTok{)) }\SpecialCharTok{+}
      \FunctionTok{geom\_boxplot}\NormalTok{(}\AttributeTok{fill =} \FunctionTok{palette\_sci}\NormalTok{(}\DecValTok{1}\NormalTok{), }\AttributeTok{alpha =} \FloatTok{0.7}\NormalTok{) }\SpecialCharTok{+}
      \FunctionTok{scale\_y\_log10}\NormalTok{() }\SpecialCharTok{+}
      \FunctionTok{labs}\NormalTok{(}\AttributeTok{x =} \StringTok{"OA Status"}\NormalTok{, }\AttributeTok{y =} \StringTok{"Citations (log)"}\NormalTok{) }\SpecialCharTok{+}
      \FunctionTok{theme\_sci}\NormalTok{()}
\NormalTok{  \})}

\NormalTok{  output}\SpecialCharTok{$}\NormalTok{summary\_table }\OtherTok{\textless{}{-}} \FunctionTok{renderTable}\NormalTok{(\{}
    \FunctionTok{req}\NormalTok{(}\FunctionTok{works\_data}\NormalTok{())}
    \FunctionTok{works\_data}\NormalTok{() }\SpecialCharTok{|\textgreater{}}
      \FunctionTok{summarise}\NormalTok{(}
        \AttributeTok{Articles =} \FunctionTok{n}\NormalTok{(),}
        \StringTok{\textasciigrave{}}\AttributeTok{Mean Citations}\StringTok{\textasciigrave{}} \OtherTok{=} \FunctionTok{round}\NormalTok{(}\FunctionTok{mean}\NormalTok{(cited\_by\_count), }\DecValTok{1}\NormalTok{),}
        \StringTok{\textasciigrave{}}\AttributeTok{Median Citations}\StringTok{\textasciigrave{}} \OtherTok{=} \FunctionTok{median}\NormalTok{(cited\_by\_count),}
        \StringTok{\textasciigrave{}}\AttributeTok{OA Rate}\StringTok{\textasciigrave{}} \OtherTok{=}\NormalTok{ scales}\SpecialCharTok{::}\FunctionTok{percent}\NormalTok{(}\FunctionTok{mean}\NormalTok{(oa\_status }\SpecialCharTok{!=} \StringTok{"closed"}\NormalTok{, }\AttributeTok{na.rm =} \ConstantTok{TRUE}\NormalTok{))}
\NormalTok{      )}
\NormalTok{  \})}

\NormalTok{  output}\SpecialCharTok{$}\NormalTok{data\_table }\OtherTok{\textless{}{-}}\NormalTok{ DT}\SpecialCharTok{::}\FunctionTok{renderDataTable}\NormalTok{(\{}
    \FunctionTok{req}\NormalTok{(}\FunctionTok{works\_data}\NormalTok{())}
    \FunctionTok{works\_data}\NormalTok{() }\SpecialCharTok{|\textgreater{}}
      \FunctionTok{select}\NormalTok{(title, year, cited\_by\_count, oa\_status) }\SpecialCharTok{|\textgreater{}}
      \FunctionTok{arrange}\NormalTok{(}\FunctionTok{desc}\NormalTok{(cited\_by\_count))}
\NormalTok{  \})}

\NormalTok{  output}\SpecialCharTok{$}\NormalTok{download\_csv }\OtherTok{\textless{}{-}} \FunctionTok{downloadHandler}\NormalTok{(}
    \AttributeTok{filename =} \ControlFlowTok{function}\NormalTok{() }\FunctionTok{paste0}\NormalTok{(}\StringTok{"biblio\_data\_"}\NormalTok{, }\FunctionTok{Sys.Date}\NormalTok{(), }\StringTok{".csv"}\NormalTok{),}
    \AttributeTok{content =} \ControlFlowTok{function}\NormalTok{(file) }\FunctionTok{write\_csv}\NormalTok{(}\FunctionTok{works\_data}\NormalTok{(), file)}
\NormalTok{  )}
\NormalTok{\}}
\end{Highlighting}
\end{Shaded}

\subsection{Running the app}\label{running-the-app}

\begin{Shaded}
\begin{Highlighting}[]
\FunctionTok{shinyApp}\NormalTok{(}\AttributeTok{ui =}\NormalTok{ ui, }\AttributeTok{server =}\NormalTok{ server)}
\end{Highlighting}
\end{Shaded}

\subsection{biblioshiny as an alternative}\label{biblioshiny-as-an-alternative}

\begin{Shaded}
\begin{Highlighting}[]
\CommentTok{\# bibliometrix provides a ready{-}made Shiny interface:}
\CommentTok{\# library(bibliometrix)}
\CommentTok{\# biblioshiny()}
\end{Highlighting}
\end{Shaded}

\section{Diagnostics and interpretation}\label{shiny-diagnostics}

\begin{itemize}
\tightlist
\item
  \textbf{Responsiveness}: API calls can take several seconds. Use \texttt{withProgress()} to show loading indicators.
\item
  \textbf{Error handling}: Wrap API calls in \texttt{tryCatch()} to display user-friendly error messages instead of stack traces.
\item
  \textbf{Memory}: Each user session loads data into R memory. For apps with many concurrent users, pre-aggregate data or use database backends.
\item
  \textbf{Testing}: Use \texttt{shinytest2} to automate testing of Shiny apps.
\end{itemize}

\section{Limitations and responsible use}\label{shiny-limitations}

\section{Limitations and responsible use}\label{limitations-and-responsible-use-36}

\begin{itemize}
\tightlist
\item
  \textbf{Shiny apps are not analyses.} They enable exploration but do not substitute for rigorous, documented analytical pipelines. Use Shiny for discovery, not for producing final results.
\item
  \textbf{API rate limits.} Live OpenAlex queries from a Shiny app can hit rate limits if many users fetch data simultaneously. Use caching or pre-fetched data for production apps.
\item
  \textbf{Maintenance.} Shiny apps require ongoing maintenance: API changes, package updates, and server costs. Plan for long-term sustainability before deploying.
\item
  \textbf{Interpretive guardrails.} Non-technical users may misinterpret metrics. Build in tooltips, descriptions, and warnings --- especially around citation-based indicators \citep{hicks2015, dora2012}.
\end{itemize}

\section{Common pitfalls}\label{shiny-pitfalls}

\section{Common pitfalls}\label{common-pitfalls-36}

\begin{itemize}
\tightlist
\item
  \textbf{Fetching data on every input change.} Use \texttt{eventReactive()} with an action button instead of \texttt{reactive()} to prevent re-fetching on every slider adjustment.
\item
  \textbf{Not validating inputs.} Check that the journal ID exists and the date range is valid before making API calls.
\item
  \textbf{Monolithic server functions.} Break the server logic into modules for maintainability as the app grows.
\item
  \textbf{Deploying without testing.} Test with realistic data volumes before deploying. An app that works with 100 records may fail with 10,000.
\end{itemize}

\section{Exercises}\label{shiny-exercises}

\begin{enumerate}
\def\labelenumi{\arabic{enumi}.}
\item
  \textbf{Add network tab.} Extend the app with a tab that builds and displays a co-authorship network using \texttt{visNetwork}.
\item
  \textbf{Caching.} Implement server-side caching so that repeated queries with the same parameters return instantly.
\item
  \textbf{Deployment.} Deploy the app to shinyapps.io. Test with different journal IDs and date ranges.
\item
  \textbf{biblioshiny comparison.} Run \texttt{biblioshiny()} and compare its features with your custom app. What does biblioshiny offer that your app does not?
\end{enumerate}

\section{Solutions}\label{37-solutions}

Solutions are provided in \ref{solutions}.

\section{Further reading}\label{shiny-reading}

\begin{itemize}
\tightlist
\item
  \citet{aria2017} --- \texttt{bibliometrix} and \texttt{biblioshiny()} for turnkey bibliometric exploration.
\item
  \citet{priem2022} --- OpenAlex API for reactive data fetching.
\item
  \citet{hicks2015} --- The Leiden Manifesto: designing interfaces that prevent metric misuse.
\item
  \citet{dora2012} --- DORA: responsible assessment in interactive tools.
\end{itemize}

\section{Session info}\label{shiny-sessioninfo}

\begin{Shaded}
\begin{Highlighting}[]
\FunctionTok{sessionInfo}\NormalTok{()}
\end{Highlighting}
\end{Shaded}

\begin{verbatim}
#> R version 4.6.1 (2026-06-24)
#> Platform: x86_64-pc-linux-gnu
#> Running under: Ubuntu 24.04.5 LTS
#> 
#> Matrix products: default
#> BLAS:   /usr/lib/x86_64-linux-gnu/openblas-pthread/libblas.so.3 
#> LAPACK: /usr/lib/x86_64-linux-gnu/openblas-pthread/libopenblasp-r0.3.26.so;  LAPACK version 3.12.0
#> 
#> locale:
#>  [1] LC_CTYPE=C.UTF-8       LC_NUMERIC=C           LC_TIME=C.UTF-8       
#>  [4] LC_COLLATE=C.UTF-8     LC_MONETARY=C.UTF-8    LC_MESSAGES=C.UTF-8   
#>  [7] LC_PAPER=C.UTF-8       LC_NAME=C              LC_ADDRESS=C          
#> [10] LC_TELEPHONE=C         LC_MEASUREMENT=C.UTF-8 LC_IDENTIFICATION=C   
#> 
#> time zone: UTC
#> tzcode source: system (glibc)
#> 
#> attached base packages:
#> [1] stats     graphics  grDevices utils     datasets  methods   base     
#> 
#> other attached packages:
#>  [1] shiny_1.14.0              DT_0.34.0                
#>  [3] plotly_4.12.1             uwot_0.2.5               
#>  [5] Matrix_1.7-5              word2vec_0.4.1           
#>  [7] stm_1.3.8                 topicmodels_0.2-17       
#>  [9] quanteda.textstats_0.97.2 visNetwork_2.1.4         
#> [11] ggraph_2.2.2              tidygraph_1.3.1          
#> [13] igraph_2.3.4              quanteda_4.5.1           
#> [15] pdftools_3.9.1            arrow_26.0.0             
#> [17] bibliometrix_5.5.0        RefManageR_1.4.0         
#> [19] bib2df_1.1.2.0            rcrossref_1.2.1          
#> [21] gt_1.4.0                  tidytext_0.4.3           
#> [23] glue_1.8.1                openalexR_3.1.2          
#> [25] lubridate_1.9.5           forcats_1.0.1            
#> [27] stringr_1.6.0             dplyr_1.2.1              
#> [29] purrr_1.2.2               readr_2.2.0              
#> [31] tidyr_1.3.2               tibble_3.3.1             
#> [33] ggplot2_4.0.3             tidyverse_2.0.0          
#> 
#> loaded via a namespace (and not attached):
#>   [1] bibtex_0.5.3           RColorBrewer_1.1-3     rstudioapi_0.19.0     
#>   [4] jsonlite_2.0.0         magrittr_2.0.5         modeltools_0.2-25     
#>   [7] farver_2.1.2           rmarkdown_2.32         fs_2.1.0              
#>  [10] vctrs_0.7.3            memoise_2.0.1          askpass_1.2.1         
#>  [13] base64enc_0.1-6        htmltools_0.5.9        contentanalysis_1.1.1 
#>  [16] curl_8.1.0             broom_1.0.13           janeaustenr_1.0.0     
#>  [19] cellranger_1.1.0       htmlwidgets_1.6.4      tokenizers_0.3.0      
#>  [22] plyr_1.8.9             httr2_1.3.0            cachem_1.1.0          
#>  [25] dimensionsR_0.0.3      mime_0.13              lifecycle_1.0.5       
#>  [28] pkgconfig_2.0.3        R6_2.6.1               fastmap_1.2.0         
#>  [31] digest_0.6.39          patchwork_1.3.2        shinycssloaders_1.1.0 
#>  [34] rprojroot_2.1.1        RSpectra_0.16-2        SnowballC_0.7.1       
#>  [37] labeling_0.4.3         urltools_1.7.3.1       timechange_0.4.0      
#>  [40] mgcv_1.9-4             httr_1.4.9             polyclip_1.10-7       
#>  [43] compiler_4.6.1         here_1.0.2             bit64_4.8.6           
#>  [46] withr_3.0.3            backports_1.5.1        S7_0.2.2              
#>  [49] viridis_0.6.5          ggforce_0.5.0          MASS_7.3-65           
#>  [52] bibliometrixData_0.3.0 tools_4.6.1            otel_0.2.0            
#>  [55] stopwords_2.3          zip_3.0.2              rentrez_1.2.4         
#>  [58] httpuv_1.6.17          nlme_3.1-169           promises_1.5.0        
#>  [61] grid_4.6.1             stringdist_0.9.17      reshape2_1.4.5        
#>  [64] generics_0.1.4         gtable_0.3.6           tzdb_0.5.0            
#>  [67] rscopus_0.9.0          ca_0.72                data.table_1.18.6.1   
#>  [70] hms_1.1.4              xml2_1.6.0             utf8_1.2.6            
#>  [73] ggrepel_0.9.8          pillar_1.11.1          nsyllable_1.0.1       
#>  [76] vroom_1.7.1            later_1.4.8            splines_4.6.1         
#>  [79] tweenr_2.0.3           lattice_0.22-9         FNN_1.1.4.1           
#>  [82] bit_4.6.0              tidyselect_1.2.1       tm_0.7-20             
#>  [85] miniUI_0.1.2           knitr_1.52             gridExtra_2.3.1       
#>  [88] NLP_0.3-3              bookdown_0.48          stats4_4.6.1          
#>  [91] crul_1.6.0             xfun_0.62              graphlayouts_1.2.5    
#>  [94] matrixStats_1.5.0      humaniformat_0.6.0     stringi_1.8.9         
#>  [97] qpdf_1.4.1             yaml_2.3.12            evaluate_1.0.5        
#> [100] httpcode_0.3.0         cli_3.6.6              xtable_1.8-8          
#> [103] systemfonts_1.3.2      Rcpp_1.1.2             readxl_1.6.0          
#> [106] triebeard_0.4.1        XML_3.99-0.25          parallel_4.6.1        
#> [109] assertthat_0.2.1       pubmedR_1.0.2          slam_0.1-56           
#> [112] viridisLite_0.4.3      scales_1.4.0           crayon_1.5.3          
#> [115] openxlsx_4.2.9         rlang_1.3.0            fastmatch_1.1-8
\end{verbatim}

\part{Case Studies}\label{part-case-studies}

\chapter{Case Study 1: CRISPR Field Review (2010--2024)}\label{case-study-1-crispr-field-review-20102024}

\emph{R. Heller · Date: 2026-05-09 · Last changed: 2026-10-30}

\section{Objective}\label{objective}

Map the structure, growth, and evolution of CRISPR gene editing research from its emergence to its current state using open bibliometric data.

\section{Setup}\label{setup}

\begin{Shaded}
\begin{Highlighting}[]
\FunctionTok{library}\NormalTok{(tidyverse)}
\FunctionTok{library}\NormalTok{(openalexR)}
\FunctionTok{library}\NormalTok{(igraph)}
\FunctionTok{library}\NormalTok{(tidygraph)}
\FunctionTok{library}\NormalTok{(ggraph)}
\FunctionTok{library}\NormalTok{(quanteda)}
\FunctionTok{library}\NormalTok{(tidytext)}
\FunctionTok{library}\NormalTok{(glue)}
\FunctionTok{library}\NormalTok{(gt)}

\FunctionTok{set.seed}\NormalTok{(}\DecValTok{20260509}\NormalTok{)}

\FunctionTok{source}\NormalTok{(here}\SpecialCharTok{::}\FunctionTok{here}\NormalTok{(}\StringTok{"R"}\NormalTok{, }\StringTok{"api\_helpers.R"}\NormalTok{))}
\FunctionTok{source}\NormalTok{(here}\SpecialCharTok{::}\FunctionTok{here}\NormalTok{(}\StringTok{"R"}\NormalTok{, }\StringTok{"utils.R"}\NormalTok{))}
\FunctionTok{source}\NormalTok{(here}\SpecialCharTok{::}\FunctionTok{here}\NormalTok{(}\StringTok{"R"}\NormalTok{, }\StringTok{"sci\_palette.R"}\NormalTok{))}
\end{Highlighting}
\end{Shaded}

\section{Data acquisition}\label{data-acquisition}

\begin{Shaded}
\begin{Highlighting}[]
\NormalTok{works }\OtherTok{\textless{}{-}} \FunctionTok{fetch\_openalex}\NormalTok{(}
  \AttributeTok{entity =} \StringTok{"works"}\NormalTok{,}
  \AttributeTok{search =} \StringTok{"CRISPR"}\NormalTok{,}
  \AttributeTok{from\_publication\_date =} \StringTok{"2010{-}01{-}01"}\NormalTok{,}
  \AttributeTok{to\_publication\_date =} \StringTok{"2024{-}06{-}30"}\NormalTok{,}
  \AttributeTok{type =} \StringTok{"article"}\NormalTok{,}
  \AttributeTok{options =} \FunctionTok{list}\NormalTok{(}\AttributeTok{sample =} \DecValTok{500}\NormalTok{, }\AttributeTok{seed =} \DecValTok{42}\NormalTok{)}
\NormalTok{)}

\NormalTok{works }\OtherTok{\textless{}{-}}\NormalTok{ works }\SpecialCharTok{|\textgreater{}}
  \FunctionTok{mutate}\NormalTok{(}\AttributeTok{year =} \FunctionTok{year}\NormalTok{(publication\_date))}

\FunctionTok{cat}\NormalTok{(}\FunctionTok{glue}\NormalTok{(}\StringTok{"CRISPR articles retrieved: \{nrow(works)\}}\SpecialCharTok{\textbackslash{}n}\StringTok{"}\NormalTok{))}
\end{Highlighting}
\end{Shaded}

\begin{verbatim}
#> CRISPR articles retrieved: 500
\end{verbatim}

\begin{Shaded}
\begin{Highlighting}[]
\FunctionTok{cat}\NormalTok{(}\FunctionTok{glue}\NormalTok{(}\StringTok{"Year range: \{min(works$year)\}{-}{-}\{max(works$year)\}}\SpecialCharTok{\textbackslash{}n}\StringTok{"}\NormalTok{))}
\end{Highlighting}
\end{Shaded}

\begin{verbatim}
#> Year range: 2012--2024
\end{verbatim}

\section{Publication growth}\label{publication-growth}

\begin{Shaded}
\begin{Highlighting}[]
\NormalTok{works }\SpecialCharTok{|\textgreater{}}
  \FunctionTok{count}\NormalTok{(year) }\SpecialCharTok{|\textgreater{}}
  \FunctionTok{ggplot}\NormalTok{(}\FunctionTok{aes}\NormalTok{(}\AttributeTok{x =}\NormalTok{ year, }\AttributeTok{y =}\NormalTok{ n)) }\SpecialCharTok{+}
  \FunctionTok{geom\_col}\NormalTok{(}\AttributeTok{fill =} \FunctionTok{palette\_sci}\NormalTok{(}\DecValTok{1}\NormalTok{)) }\SpecialCharTok{+}
  \FunctionTok{labs}\NormalTok{(}\AttributeTok{x =} \StringTok{"Year"}\NormalTok{, }\AttributeTok{y =} \StringTok{"Publications"}\NormalTok{) }\SpecialCharTok{+}
  \FunctionTok{theme\_sci}\NormalTok{()}
\end{Highlighting}
\end{Shaded}

\begin{figure}

{\centering \includegraphics[width=0.9\linewidth,alt={Bar chart showing exponential growth in CRISPR publications from 2010 to 2024.}]{scientometrics-in-r_files/figure-latex/cs01-fig-cs01-growth-1} 

}

\caption{Annual publication output in the CRISPR field.}\label{fig:cs01-fig-cs01-growth}
\end{figure}

\section{Citation landscape}\label{citation-landscape}

\begin{Shaded}
\begin{Highlighting}[]
\FunctionTok{ggplot}\NormalTok{(works, }\FunctionTok{aes}\NormalTok{(}\AttributeTok{x =}\NormalTok{ cited\_by\_count)) }\SpecialCharTok{+}
  \FunctionTok{geom\_histogram}\NormalTok{(}\AttributeTok{binwidth =} \DecValTok{10}\NormalTok{, }\AttributeTok{fill =} \FunctionTok{palette\_sci}\NormalTok{(}\DecValTok{1}\NormalTok{), }\AttributeTok{colour =} \StringTok{"white"}\NormalTok{) }\SpecialCharTok{+}
  \FunctionTok{labs}\NormalTok{(}\AttributeTok{x =} \StringTok{"Citations"}\NormalTok{, }\AttributeTok{y =} \StringTok{"Papers"}\NormalTok{) }\SpecialCharTok{+}
  \FunctionTok{theme\_sci}\NormalTok{()}
\end{Highlighting}
\end{Shaded}

\begin{figure}

{\centering \includegraphics[width=0.9\linewidth,alt={Histogram showing the heavily right-skewed citation distribution typical of a rapidly growing field.}]{scientometrics-in-r_files/figure-latex/cs01-fig-cs01-citations-1} 

}

\caption{Citation distribution of CRISPR articles.}\label{fig:cs01-fig-cs01-citations}
\end{figure}

\begin{Shaded}
\begin{Highlighting}[]
\NormalTok{works }\SpecialCharTok{|\textgreater{}}
  \FunctionTok{arrange}\NormalTok{(}\FunctionTok{desc}\NormalTok{(cited\_by\_count)) }\SpecialCharTok{|\textgreater{}}
  \FunctionTok{head}\NormalTok{(}\DecValTok{10}\NormalTok{) }\SpecialCharTok{|\textgreater{}}
  \FunctionTok{select}\NormalTok{(display\_name, year, cited\_by\_count, source\_display\_name) }\SpecialCharTok{|\textgreater{}}
  \FunctionTok{gt}\NormalTok{()}
\end{Highlighting}
\end{Shaded}

\begingroup
\fontsize{8.0pt}{9.0pt}\selectfont
\begin{longtable}{>{\raggedright\arraybackslash}p{\dimexpr 0.41\linewidth -2\tabcolsep-1.5\arrayrulewidth}>{\raggedleft\arraybackslash}p{\dimexpr 0.15\linewidth -2\tabcolsep-1.5\arrayrulewidth}>{\raggedleft\arraybackslash}p{\dimexpr 0.20\linewidth -2\tabcolsep-1.5\arrayrulewidth}>{\raggedright\arraybackslash}p{\dimexpr 0.24\linewidth -2\tabcolsep-1.5\arrayrulewidth}}
\toprule
\nolinkurl{display_name} & year & \nolinkurl{cited_by_count} & \nolinkurl{source_display_name} \\ 
\midrule\addlinespace[2.5pt]
Local regulation of gene expression by lncRNA promoters, transcription and splicing & 2016 & 1267 & Nature \\ 
Parkinson-associated risk variant in distal enhancer of α-synuclein modulates target gene expression & 2016 & 507 & Nature \\ 
The Human Microbiome Project: A Community Resource for the Healthy Human Microbiome & 2012 & 507 & PLoS Biology \\ 
Expanded base editing in rice and wheat using a Cas9-adenosine deaminase fusion & 2018 & 504 & Genome biology \\ 
A reporter for noninvasively monitoring gene expression and plant transformation & 2020 & 461 & Horticulture Research \\ 
Cas9 specifies functional viral targets during CRISPR–Cas adaptation & 2015 & 431 & Nature \\ 
Genome-scale activation screen identifies a lncRNA locus regulating a gene neighbourhood & 2017 & 427 & Nature \\ 
The mechanosensitive Piezo1 channel mediates heart mechano-chemo transduction & 2021 & 304 & Nature Communications \\ 
Structure and Mechanism of the CMR Complex for CRISPR-Mediated Antiviral Immunity & 2012 & 295 & Molecular Cell \\ 
Identification of a Master Regulator of Differentiation in Toxoplasma & 2020 & 264 & Cell \\ 
\bottomrule
\end{longtable}
\endgroup

\section{Co-authorship network}\label{co-authorship-network}

\begin{Shaded}
\begin{Highlighting}[]
\NormalTok{author\_data }\OtherTok{\textless{}{-}}\NormalTok{ works }\SpecialCharTok{|\textgreater{}}
  \FunctionTok{select}\NormalTok{(id, authorships) }\SpecialCharTok{|\textgreater{}}
  \FunctionTok{unnest}\NormalTok{(authorships, }\AttributeTok{names\_sep =} \StringTok{"\_"}\NormalTok{) }\SpecialCharTok{|\textgreater{}}
  \FunctionTok{select}\NormalTok{(}\AttributeTok{work\_id =}\NormalTok{ id, }\AttributeTok{author\_id =}\NormalTok{ authorships\_id,}
         \AttributeTok{author\_name =}\NormalTok{ authorships\_display\_name) }\SpecialCharTok{|\textgreater{}}
  \FunctionTok{filter}\NormalTok{(}\SpecialCharTok{!}\FunctionTok{is.na}\NormalTok{(author\_id))}

\NormalTok{edges }\OtherTok{\textless{}{-}}\NormalTok{ author\_data }\SpecialCharTok{|\textgreater{}}
  \FunctionTok{inner\_join}\NormalTok{(author\_data, }\AttributeTok{by =} \StringTok{"work\_id"}\NormalTok{, }\AttributeTok{suffix =} \FunctionTok{c}\NormalTok{(}\StringTok{"\_1"}\NormalTok{, }\StringTok{"\_2"}\NormalTok{),}
             \AttributeTok{relationship =} \StringTok{"many{-}to{-}many"}\NormalTok{) }\SpecialCharTok{|\textgreater{}}
  \FunctionTok{filter}\NormalTok{(author\_id\_1 }\SpecialCharTok{\textless{}}\NormalTok{ author\_id\_2) }\SpecialCharTok{|\textgreater{}}
  \FunctionTok{count}\NormalTok{(author\_id\_1, author\_id\_2, }\AttributeTok{name =} \StringTok{"weight"}\NormalTok{)}

\NormalTok{g }\OtherTok{\textless{}{-}} \FunctionTok{graph\_from\_data\_frame}\NormalTok{(}
\NormalTok{  edges }\SpecialCharTok{|\textgreater{}} \FunctionTok{select}\NormalTok{(author\_id\_1, author\_id\_2, weight),}
  \AttributeTok{directed =} \ConstantTok{FALSE}
\NormalTok{) }\SpecialCharTok{|\textgreater{}} \FunctionTok{simplify}\NormalTok{(}\AttributeTok{edge.attr.comb =} \FunctionTok{list}\NormalTok{(}\AttributeTok{weight =} \StringTok{"sum"}\NormalTok{))}

\NormalTok{comp }\OtherTok{\textless{}{-}} \FunctionTok{components}\NormalTok{(g)}
\NormalTok{giant }\OtherTok{\textless{}{-}} \FunctionTok{induced\_subgraph}\NormalTok{(g, }\FunctionTok{which}\NormalTok{(comp}\SpecialCharTok{$}\NormalTok{membership }\SpecialCharTok{==} \FunctionTok{which.max}\NormalTok{(comp}\SpecialCharTok{$}\NormalTok{csize)))}
\FunctionTok{V}\NormalTok{(giant)}\SpecialCharTok{$}\NormalTok{community }\OtherTok{\textless{}{-}} \FunctionTok{as.factor}\NormalTok{(}\FunctionTok{membership}\NormalTok{(}
  \FunctionTok{cluster\_leiden}\NormalTok{(giant, }\AttributeTok{resolution\_parameter =} \FloatTok{1.0}\NormalTok{,}
                 \AttributeTok{objective\_function =} \StringTok{"modularity"}\NormalTok{)}
\NormalTok{))}
\FunctionTok{V}\NormalTok{(giant)}\SpecialCharTok{$}\NormalTok{degree }\OtherTok{\textless{}{-}} \FunctionTok{degree}\NormalTok{(giant)}

\FunctionTok{cat}\NormalTok{(}\FunctionTok{glue}\NormalTok{(}\StringTok{"Network: \{vcount(giant)\} nodes, \{ecount(giant)\} edges}\SpecialCharTok{\textbackslash{}n}\StringTok{"}\NormalTok{))}
\end{Highlighting}
\end{Shaded}

\begin{verbatim}
#> Network: 73 nodes, 1061 edges
\end{verbatim}

\begin{Shaded}
\begin{Highlighting}[]
\FunctionTok{set.seed}\NormalTok{(}\DecValTok{42}\NormalTok{)}
\FunctionTok{ggraph}\NormalTok{(}\FunctionTok{as\_tbl\_graph}\NormalTok{(giant), }\AttributeTok{layout =} \StringTok{"fr"}\NormalTok{) }\SpecialCharTok{+}
  \FunctionTok{geom\_edge\_link}\NormalTok{(}\AttributeTok{alpha =} \FloatTok{0.1}\NormalTok{, }\AttributeTok{colour =} \StringTok{"grey60"}\NormalTok{) }\SpecialCharTok{+}
  \FunctionTok{geom\_node\_point}\NormalTok{(}\FunctionTok{aes}\NormalTok{(}\AttributeTok{size =}\NormalTok{ degree, }\AttributeTok{colour =}\NormalTok{ community), }\AttributeTok{alpha =} \FloatTok{0.7}\NormalTok{) }\SpecialCharTok{+}
  \FunctionTok{scale\_size\_continuous}\NormalTok{(}\AttributeTok{range =} \FunctionTok{c}\NormalTok{(}\FloatTok{0.5}\NormalTok{, }\DecValTok{5}\NormalTok{), }\AttributeTok{guide =} \StringTok{"none"}\NormalTok{) }\SpecialCharTok{+}
  \FunctionTok{scale\_colour\_manual}\NormalTok{(}\AttributeTok{values =} \FunctionTok{palette\_sci}\NormalTok{(}\FunctionTok{n\_distinct}\NormalTok{(}\FunctionTok{V}\NormalTok{(giant)}\SpecialCharTok{$}\NormalTok{community))) }\SpecialCharTok{+}
  \FunctionTok{labs}\NormalTok{(}\AttributeTok{colour =} \StringTok{"Community"}\NormalTok{) }\SpecialCharTok{+}
  \FunctionTok{theme\_void}\NormalTok{(}\AttributeTok{base\_family =} \StringTok{"sans"}\NormalTok{, }\AttributeTok{base\_size =} \DecValTok{11}\NormalTok{) }\SpecialCharTok{+} \FunctionTok{theme}\NormalTok{(}\AttributeTok{legend.position =} \StringTok{"bottom"}\NormalTok{)}
\end{Highlighting}
\end{Shaded}

\begin{figure}

{\centering \includegraphics[width=0.9\linewidth,alt={Network graph showing collaboration patterns among CRISPR researchers, with communities coloured by Leiden algorithm.}]{scientometrics-in-r_files/figure-latex/cs01-fig-cs01-network-1} 

}

\caption{Co-authorship network of CRISPR researchers.}\label{fig:cs01-fig-cs01-network}
\end{figure}

\section{Topic evolution}\label{topic-evolution}

\begin{Shaded}
\begin{Highlighting}[]
\NormalTok{text\_df }\OtherTok{\textless{}{-}}\NormalTok{ works }\SpecialCharTok{|\textgreater{}}
  \FunctionTok{filter}\NormalTok{(}\SpecialCharTok{!}\FunctionTok{is.na}\NormalTok{(abstract), }\FunctionTok{nchar}\NormalTok{(abstract) }\SpecialCharTok{\textgreater{}} \DecValTok{50}\NormalTok{) }\SpecialCharTok{|\textgreater{}}
  \FunctionTok{transmute}\NormalTok{(}\AttributeTok{doc\_id =}\NormalTok{ id, }\AttributeTok{text =} \FunctionTok{paste}\NormalTok{(display\_name, abstract, }\AttributeTok{sep =} \StringTok{". "}\NormalTok{), year)}

\NormalTok{corp }\OtherTok{\textless{}{-}} \FunctionTok{corpus}\NormalTok{(text\_df, }\AttributeTok{docid\_field =} \StringTok{"doc\_id"}\NormalTok{, }\AttributeTok{text\_field =} \StringTok{"text"}\NormalTok{)}
\NormalTok{toks }\OtherTok{\textless{}{-}} \FunctionTok{tokens}\NormalTok{(corp, }\AttributeTok{remove\_punct =} \ConstantTok{TRUE}\NormalTok{, }\AttributeTok{remove\_numbers =} \ConstantTok{TRUE}\NormalTok{) }\SpecialCharTok{|\textgreater{}}
  \FunctionTok{tokens\_tolower}\NormalTok{() }\SpecialCharTok{|\textgreater{}}
  \FunctionTok{tokens\_remove}\NormalTok{(}\FunctionTok{stopwords}\NormalTok{(}\StringTok{"en"}\NormalTok{)) }\SpecialCharTok{|\textgreater{}}
  \FunctionTok{tokens\_remove}\NormalTok{(}\FunctionTok{c}\NormalTok{(}\StringTok{"study"}\NormalTok{, }\StringTok{"paper"}\NormalTok{, }\StringTok{"results"}\NormalTok{, }\StringTok{"using"}\NormalTok{, }\StringTok{"based"}\NormalTok{))}

\NormalTok{dfmat }\OtherTok{\textless{}{-}} \FunctionTok{dfm}\NormalTok{(toks) }\SpecialCharTok{|\textgreater{}} \FunctionTok{dfm\_trim}\NormalTok{(}\AttributeTok{min\_termfreq =} \DecValTok{5}\NormalTok{, }\AttributeTok{min\_docfreq =} \DecValTok{3}\NormalTok{)}

\NormalTok{top\_by\_year }\OtherTok{\textless{}{-}} \FunctionTok{map\_dfr}\NormalTok{(}\FunctionTok{unique}\NormalTok{(text\_df}\SpecialCharTok{$}\NormalTok{year), }\ControlFlowTok{function}\NormalTok{(yr) \{}
\NormalTok{  docs }\OtherTok{\textless{}{-}} \FunctionTok{docvars}\NormalTok{(dfmat, }\StringTok{"year"}\NormalTok{) }\SpecialCharTok{==}\NormalTok{ yr}
  \ControlFlowTok{if}\NormalTok{ (}\FunctionTok{sum}\NormalTok{(docs) }\SpecialCharTok{\textless{}} \DecValTok{5}\NormalTok{) }\FunctionTok{return}\NormalTok{(}\FunctionTok{tibble}\NormalTok{())}
\NormalTok{  top }\OtherTok{\textless{}{-}} \FunctionTok{topfeatures}\NormalTok{(dfmat[docs, ], }\DecValTok{5}\NormalTok{)}
  \FunctionTok{tibble}\NormalTok{(}\AttributeTok{year =}\NormalTok{ yr, }\AttributeTok{term =} \FunctionTok{names}\NormalTok{(top), }\AttributeTok{freq =} \FunctionTok{unname}\NormalTok{(top))}
\NormalTok{\})}
\end{Highlighting}
\end{Shaded}

\begin{Shaded}
\begin{Highlighting}[]
\NormalTok{top\_by\_year }\SpecialCharTok{|\textgreater{}}
  \FunctionTok{group\_by}\NormalTok{(year) }\SpecialCharTok{|\textgreater{}}
  \FunctionTok{mutate}\NormalTok{(}\AttributeTok{term =} \FunctionTok{reorder\_within}\NormalTok{(term, freq, year)) }\SpecialCharTok{|\textgreater{}}
  \FunctionTok{ggplot}\NormalTok{(}\FunctionTok{aes}\NormalTok{(}\AttributeTok{x =}\NormalTok{ freq, }\AttributeTok{y =}\NormalTok{ term)) }\SpecialCharTok{+}
  \FunctionTok{geom\_col}\NormalTok{(}\AttributeTok{fill =} \FunctionTok{palette\_sci}\NormalTok{(}\DecValTok{1}\NormalTok{)) }\SpecialCharTok{+}
  \FunctionTok{facet\_wrap}\NormalTok{(}\SpecialCharTok{\textasciitilde{}}\NormalTok{ year, }\AttributeTok{scales =} \StringTok{"free\_y"}\NormalTok{, }\AttributeTok{ncol =} \DecValTok{4}\NormalTok{) }\SpecialCharTok{+}
  \FunctionTok{scale\_y\_reordered}\NormalTok{() }\SpecialCharTok{+}
  \FunctionTok{labs}\NormalTok{(}\AttributeTok{x =} \StringTok{"Frequency"}\NormalTok{, }\AttributeTok{y =} \ConstantTok{NULL}\NormalTok{) }\SpecialCharTok{+}
  \FunctionTok{theme\_sci}\NormalTok{(}\AttributeTok{base\_size =} \DecValTok{8}\NormalTok{)}
\end{Highlighting}
\end{Shaded}

\begin{figure}

{\centering \includegraphics[width=0.9\linewidth,alt={Faceted bar chart showing the top 5 terms in each year of CRISPR research.}]{scientometrics-in-r_files/figure-latex/cs01-fig-cs01-topics-1} 

}

\caption{Top terms by year showing topical evolution.}\label{fig:cs01-fig-cs01-topics}
\end{figure}

\section{Key findings}\label{key-findings}

\begin{enumerate}
\def\labelenumi{\arabic{enumi}.}
\tightlist
\item
  \textbf{Explosive growth}: CRISPR publications grew exponentially from 2012, reflecting the rapid adoption of Cas9-based editing.
\item
  \textbf{Citation concentration}: A small number of foundational papers dominate the citation landscape.
\item
  \textbf{Collaborative structure}: The co-authorship network shows distinct communities, likely corresponding to different application domains (therapeutics, agriculture, basic biology).
\item
  \textbf{Topic evolution}: Early terms focus on methodology; later years shift toward applications and clinical translation.
\end{enumerate}

\section{Lessons learned}\label{lessons-learned}

\begin{itemize}
\tightlist
\item
  OpenAlex sampling provides a representative snapshot but may miss some highly specialised or non-English publications.
\item
  The citation distribution is extreme: median citations are far below the mean, making median-based statistics essential.
\item
  Co-authorship networks in fast-growing fields are fragmented; many research groups work independently.
\end{itemize}

\chapter{Case Study 2: Institutional Benchmarking}\label{case-study-2-institutional-benchmarking}

\emph{R. Heller · Date: 2026-05-09 · Last changed: 2026-10-31}

\section{Objective}\label{objective-1}

Compare citation summaries, fractional credits and open-access metadata in capped publication samples from five universities. These descriptive samples do not measure total productivity or field-normalised impact.

\section{Setup}\label{setup-1}

\begin{Shaded}
\begin{Highlighting}[]
\FunctionTok{library}\NormalTok{(tidyverse)}
\FunctionTok{library}\NormalTok{(openalexR)}
\FunctionTok{library}\NormalTok{(glue)}
\FunctionTok{library}\NormalTok{(gt)}

\FunctionTok{set.seed}\NormalTok{(}\DecValTok{20260509}\NormalTok{)}

\FunctionTok{source}\NormalTok{(here}\SpecialCharTok{::}\FunctionTok{here}\NormalTok{(}\StringTok{"R"}\NormalTok{, }\StringTok{"api\_helpers.R"}\NormalTok{))}
\FunctionTok{source}\NormalTok{(here}\SpecialCharTok{::}\FunctionTok{here}\NormalTok{(}\StringTok{"R"}\NormalTok{, }\StringTok{"utils.R"}\NormalTok{))}
\FunctionTok{source}\NormalTok{(here}\SpecialCharTok{::}\FunctionTok{here}\NormalTok{(}\StringTok{"R"}\NormalTok{, }\StringTok{"sci\_palette.R"}\NormalTok{))}
\FunctionTok{source}\NormalTok{(here}\SpecialCharTok{::}\FunctionTok{here}\NormalTok{(}\StringTok{"R"}\NormalTok{, }\StringTok{"institutional\_samples.R"}\NormalTok{))}
\end{Highlighting}
\end{Shaded}

\section{Data acquisition}\label{data-acquisition-1}

Institution IDs were checked against OpenAlex on 30 October 2026, including the
\href{https://openalex.org/institutions/I17974374}{University of São Paulo record}. Each
request retrieves at most 300 articles; these are samples, not output totals.
A failed or empty request stops the comparison instead of silently dropping an institution.

\begin{Shaded}
\begin{Highlighting}[]
\NormalTok{institutions }\OtherTok{\textless{}{-}} \FunctionTok{tribble}\NormalTok{(}
  \SpecialCharTok{\textasciitilde{}}\NormalTok{short\_name,    }\SpecialCharTok{\textasciitilde{}}\NormalTok{oa\_id,}
  \StringTok{"MIT"}\NormalTok{,          }\StringTok{"I63966007"}\NormalTok{,}
  \StringTok{"Oxford"}\NormalTok{,       }\StringTok{"I40120149"}\NormalTok{,}
  \StringTok{"ETH Zurich"}\NormalTok{,   }\StringTok{"I35440088"}\NormalTok{,}
  \StringTok{"U São Paulo"}\NormalTok{,  }\StringTok{"I17974374"}\NormalTok{,}
  \StringTok{"U Cape Town"}\NormalTok{,  }\StringTok{"I157614274"}
\NormalTok{)}
\end{Highlighting}
\end{Shaded}

\begin{Shaded}
\begin{Highlighting}[]
\NormalTok{fetch\_works }\OtherTok{\textless{}{-}} \ControlFlowTok{function}\NormalTok{(inst\_id) \{}
\NormalTok{  res }\OtherTok{\textless{}{-}} \FunctionTok{fetch\_openalex}\NormalTok{(}
    \AttributeTok{entity =} \StringTok{"works"}\NormalTok{,}
    \AttributeTok{authorships.institutions.id =}\NormalTok{ inst\_id,}
    \AttributeTok{from\_publication\_date =} \StringTok{"2020{-}01{-}01"}\NormalTok{,}
    \AttributeTok{to\_publication\_date =} \StringTok{"2022{-}12{-}31"}\NormalTok{,}
    \AttributeTok{type =} \StringTok{"article"}\NormalTok{,}
    \AttributeTok{options =} \FunctionTok{list}\NormalTok{(}\AttributeTok{sample =} \DecValTok{300}\NormalTok{, }\AttributeTok{seed =} \DecValTok{42}\NormalTok{)}
\NormalTok{  )}
  \ControlFlowTok{if}\NormalTok{ (}\SpecialCharTok{!}\FunctionTok{is.data.frame}\NormalTok{(res) }\SpecialCharTok{||} \SpecialCharTok{!}\FunctionTok{nrow}\NormalTok{(res)) \{}
    \FunctionTok{stop}\NormalTok{(}\FunctionTok{glue}\NormalTok{(}\StringTok{"No usable works returned for \{inst\_id\}; check the request."}\NormalTok{))}
\NormalTok{  \}}
\NormalTok{  res}
\NormalTok{\}}

\NormalTok{inst\_data }\OtherTok{\textless{}{-}}\NormalTok{ institutions }\SpecialCharTok{|\textgreater{}}
  \FunctionTok{mutate}\NormalTok{(}\AttributeTok{works =} \FunctionTok{map}\NormalTok{(oa\_id, fetch\_works))}
\end{Highlighting}
\end{Shaded}

\section{Publication volume}\label{publication-volume}

\begin{Shaded}
\begin{Highlighting}[]
\NormalTok{inst\_data }\SpecialCharTok{|\textgreater{}}
  \FunctionTok{mutate}\NormalTok{(}\AttributeTok{n =} \FunctionTok{map\_int}\NormalTok{(works, nrow)) }\SpecialCharTok{|\textgreater{}}
  \FunctionTok{ggplot}\NormalTok{(}\FunctionTok{aes}\NormalTok{(}\AttributeTok{x =}\NormalTok{ n, }\AttributeTok{y =} \FunctionTok{reorder}\NormalTok{(short\_name, n))) }\SpecialCharTok{+}
  \FunctionTok{geom\_col}\NormalTok{(}\AttributeTok{fill =} \FunctionTok{palette\_sci}\NormalTok{(}\DecValTok{1}\NormalTok{)) }\SpecialCharTok{+}
  \FunctionTok{labs}\NormalTok{(}\AttributeTok{x =} \StringTok{"Sampled articles"}\NormalTok{, }\AttributeTok{y =} \ConstantTok{NULL}\NormalTok{) }\SpecialCharTok{+}
  \FunctionTok{theme\_sci}\NormalTok{()}
\end{Highlighting}
\end{Shaded}

\begin{figure}

{\centering \includegraphics[width=0.9\linewidth,alt={Bar chart comparing sampled publication volumes for five universities.}]{scientometrics-in-r_files/figure-latex/cs02-fig-cs02-volume-1} 

}

\caption{Sampled publication counts by institution.}\label{fig:cs02-fig-cs02-volume}
\end{figure}

\section{Citation summaries}\label{citation-summaries}

As in Chapter 13, the citation ratio uses a common year reference from the union
of the five samples, de-duplicated by work ID. It adjusts for year only, not
field mix or a common citation window. Missing citation counts and undefined
ratios are excluded with their usable denominators shown.

\begin{Shaded}
\begin{Highlighting}[]
\NormalTok{citation\_reference }\OtherTok{\textless{}{-}} \FunctionTok{sci\_pooled\_year\_reference}\NormalTok{(inst\_data}\SpecialCharTok{$}\NormalTok{works)}
\NormalTok{impact }\OtherTok{\textless{}{-}}\NormalTok{ inst\_data }\SpecialCharTok{|\textgreater{}}
  \FunctionTok{mutate}\NormalTok{(}\AttributeTok{summary =} \FunctionTok{map}\NormalTok{(works, sci\_sample\_citation\_summary,}
                       \AttributeTok{reference =}\NormalTok{ citation\_reference)) }\SpecialCharTok{|\textgreater{}}
  \FunctionTok{unnest}\NormalTok{(summary) }\SpecialCharTok{|\textgreater{}}
  \FunctionTok{select}\NormalTok{(short\_name, n\_sampled, n\_citation\_observed, mean\_cites,}
\NormalTok{         median\_cites, n\_ratio\_observed, pooled\_year\_ratio) }\SpecialCharTok{|\textgreater{}}
  \FunctionTok{arrange}\NormalTok{(short\_name)}

\NormalTok{impact }\SpecialCharTok{|\textgreater{}}
  \FunctionTok{gt}\NormalTok{() }\SpecialCharTok{|\textgreater{}}
  \FunctionTok{fmt\_number}\NormalTok{(}\AttributeTok{columns =} \FunctionTok{c}\NormalTok{(mean\_cites, median\_cites, pooled\_year\_ratio), }\AttributeTok{decimals =} \DecValTok{2}\NormalTok{)}
\end{Highlighting}
\end{Shaded}

\begingroup
\fontsize{8.0pt}{9.0pt}\selectfont
\begin{longtable}{>{\raggedright\arraybackslash}p{\dimexpr 0.13\linewidth -2\tabcolsep-1.5\arrayrulewidth}>{\raggedleft\arraybackslash}p{\dimexpr 0.12\linewidth -2\tabcolsep-1.5\arrayrulewidth}>{\raggedleft\arraybackslash}p{\dimexpr 0.17\linewidth -2\tabcolsep-1.5\arrayrulewidth}>{\raggedleft\arraybackslash}p{\dimexpr 0.12\linewidth -2\tabcolsep-1.5\arrayrulewidth}>{\raggedleft\arraybackslash}p{\dimexpr 0.14\linewidth -2\tabcolsep-1.5\arrayrulewidth}>{\raggedleft\arraybackslash}p{\dimexpr 0.16\linewidth -2\tabcolsep-1.5\arrayrulewidth}>{\raggedleft\arraybackslash}p{\dimexpr 0.16\linewidth -2\tabcolsep-1.5\arrayrulewidth}}
\toprule
\nolinkurl{short_name} & \nolinkurl{n_sampled} & \nolinkurl{n_citation_observed} & \nolinkurl{mean_cites} & \nolinkurl{median_cites} & \nolinkurl{n_ratio_observed} & \nolinkurl{pooled_year_ratio} \\ 
\midrule\addlinespace[2.5pt]
ETH Zurich & 300 & 300 & 52.02 & 23.00 & 300 & 1.43 \\ 
MIT & 300 & 300 & 51.59 & 25.00 & 300 & 1.40 \\ 
Oxford & 300 & 300 & 39.69 & 15.00 & 300 & 1.08 \\ 
U Cape Town & 300 & 300 & 29.54 & 12.00 & 300 & 0.77 \\ 
U São Paulo & 300 & 300 & 13.22 & 6.00 & 300 & 0.36 \\ 
\bottomrule
\end{longtable}
\endgroup

\section{Full vs.~fractional counting}\label{full-vs.-fractional-counting}

Count each distinct affiliation once per work and divide its credit equally.
Incomplete affiliations make the fractional total unavailable; the coverage
column shows how many sampled works have usable data. This is bibliometric
credit, not an estimate of each institution's contribution.

\begin{Shaded}
\begin{Highlighting}[]
\NormalTok{counting }\OtherTok{\textless{}{-}}\NormalTok{ inst\_data }\SpecialCharTok{|\textgreater{}}
  \FunctionTok{mutate}\NormalTok{(}\AttributeTok{counts =} \FunctionTok{map2}\NormalTok{(works, oa\_id, sci\_institution\_sample\_counts)) }\SpecialCharTok{|\textgreater{}}
  \FunctionTok{unnest}\NormalTok{(counts) }\SpecialCharTok{|\textgreater{}}
  \FunctionTok{select}\NormalTok{(short\_name, full\_count, affiliation\_covered, frac\_count) }\SpecialCharTok{|\textgreater{}}
  \FunctionTok{mutate}\NormalTok{(}\AttributeTok{ratio =} \FunctionTok{round}\NormalTok{(frac\_count }\SpecialCharTok{/}\NormalTok{ full\_count, }\DecValTok{2}\NormalTok{))}

\NormalTok{counting }\SpecialCharTok{|\textgreater{}} \FunctionTok{gt}\NormalTok{()}
\end{Highlighting}
\end{Shaded}

\begingroup
\fontsize{8.0pt}{9.0pt}\selectfont
\begin{longtable}{>{\raggedright\arraybackslash}p{\dimexpr 0.20\linewidth -2\tabcolsep-1.5\arrayrulewidth}>{\raggedleft\arraybackslash}p{\dimexpr 0.19\linewidth -2\tabcolsep-1.5\arrayrulewidth}>{\raggedleft\arraybackslash}p{\dimexpr 0.26\linewidth -2\tabcolsep-1.5\arrayrulewidth}>{\raggedleft\arraybackslash}p{\dimexpr 0.19\linewidth -2\tabcolsep-1.5\arrayrulewidth}>{\raggedleft\arraybackslash}p{\dimexpr 0.17\linewidth -2\tabcolsep-1.5\arrayrulewidth}}
\toprule
\nolinkurl{short_name} & \nolinkurl{full_count} & \nolinkurl{affiliation_covered} & \nolinkurl{frac_count} & ratio \\ 
\midrule\addlinespace[2.5pt]
MIT & 300 & 251 & NA & NA \\ 
Oxford & 300 & 263 & NA & NA \\ 
ETH Zurich & 300 & 271 & NA & NA \\ 
U São Paulo & 300 & 258 & NA & NA \\ 
U Cape Town & 300 & 255 & NA & NA \\ 
\bottomrule
\end{longtable}
\endgroup

\begin{Shaded}
\begin{Highlighting}[]
\NormalTok{counting }\SpecialCharTok{|\textgreater{}}
  \FunctionTok{pivot\_longer}\NormalTok{(}\FunctionTok{c}\NormalTok{(full\_count, frac\_count), }\AttributeTok{names\_to =} \StringTok{"method"}\NormalTok{, }\AttributeTok{values\_to =} \StringTok{"count"}\NormalTok{) }\SpecialCharTok{|\textgreater{}}
  \FunctionTok{mutate}\NormalTok{(}\AttributeTok{method =} \FunctionTok{recode}\NormalTok{(method, }\AttributeTok{full\_count =} \StringTok{"Full"}\NormalTok{, }\AttributeTok{frac\_count =} \StringTok{"Fractional"}\NormalTok{)) }\SpecialCharTok{|\textgreater{}}
  \FunctionTok{ggplot}\NormalTok{(}\FunctionTok{aes}\NormalTok{(}\AttributeTok{x =}\NormalTok{ count, }\AttributeTok{y =} \FunctionTok{reorder}\NormalTok{(short\_name, count), }\AttributeTok{fill =}\NormalTok{ method)) }\SpecialCharTok{+}
  \FunctionTok{geom\_col}\NormalTok{(}\AttributeTok{position =} \StringTok{"dodge"}\NormalTok{, }\AttributeTok{width =} \FloatTok{0.7}\NormalTok{) }\SpecialCharTok{+}
  \FunctionTok{scale\_fill\_manual}\NormalTok{(}\AttributeTok{values =} \FunctionTok{palette\_sci}\NormalTok{(}\DecValTok{2}\NormalTok{)) }\SpecialCharTok{+}
  \FunctionTok{labs}\NormalTok{(}\AttributeTok{x =} \StringTok{"Count"}\NormalTok{, }\AttributeTok{y =} \ConstantTok{NULL}\NormalTok{, }\AttributeTok{fill =} \StringTok{"Method"}\NormalTok{) }\SpecialCharTok{+}
  \FunctionTok{theme\_sci}\NormalTok{()}
\end{Highlighting}
\end{Shaded}

\begin{figure}

{\centering \includegraphics[width=0.9\linewidth,alt={Grouped bar chart comparing full sample counts with fractional credits where affiliation coverage is complete.}]{scientometrics-in-r_files/figure-latex/cs02-fig-cs02-counting-1} 

}

\caption{Full and fractional credits within the retrieved samples.}\label{fig:cs02-fig-cs02-counting}
\end{figure}

\section{OA rates}\label{oa-rates}

\begin{Shaded}
\begin{Highlighting}[]
\NormalTok{inst\_data }\SpecialCharTok{|\textgreater{}}
  \FunctionTok{mutate}\NormalTok{(}\AttributeTok{oa\_rate =} \FunctionTok{map\_dbl}\NormalTok{(works, \textbackslash{}(w) }\FunctionTok{mean}\NormalTok{(w}\SpecialCharTok{$}\NormalTok{oa\_status }\SpecialCharTok{!=} \StringTok{"closed"}\NormalTok{, }\AttributeTok{na.rm =} \ConstantTok{TRUE}\NormalTok{))) }\SpecialCharTok{|\textgreater{}}
  \FunctionTok{ggplot}\NormalTok{(}\FunctionTok{aes}\NormalTok{(}\AttributeTok{x =}\NormalTok{ oa\_rate, }\AttributeTok{y =} \FunctionTok{reorder}\NormalTok{(short\_name, oa\_rate))) }\SpecialCharTok{+}
  \FunctionTok{geom\_col}\NormalTok{(}\AttributeTok{fill =} \FunctionTok{palette\_sci}\NormalTok{(}\DecValTok{1}\NormalTok{)) }\SpecialCharTok{+}
  \FunctionTok{scale\_x\_continuous}\NormalTok{(}\AttributeTok{labels =}\NormalTok{ scales}\SpecialCharTok{::}\NormalTok{percent) }\SpecialCharTok{+}
  \FunctionTok{labs}\NormalTok{(}\AttributeTok{x =} \StringTok{"OA rate"}\NormalTok{, }\AttributeTok{y =} \ConstantTok{NULL}\NormalTok{) }\SpecialCharTok{+}
  \FunctionTok{theme\_sci}\NormalTok{()}
\end{Highlighting}
\end{Shaded}

\begin{figure}

{\centering \includegraphics[width=0.9\linewidth,alt={Bar chart comparing observed open-access shares in the five university samples.}]{scientometrics-in-r_files/figure-latex/cs02-fig-cs02-oa-1} 

}

\caption{Open-access share among sampled articles with observed OA status.}\label{fig:cs02-fig-cs02-oa}
\end{figure}

\section{Reading the results}\label{reading-the-results}

\begin{enumerate}
\def\labelenumi{\arabic{enumi}.}
\tightlist
\item
  \textbf{Sample size}: Counts are capped by the request and cannot rank institutional productivity.
\item
  \textbf{Citation summaries}: Means and pooled sample/year ratios describe the retrieved records. Neither is field-normalised; an h-index would also depend on sample size and is omitted here.
\item
  \textbf{Fractional credit}: A lower fractional/full ratio means less credit is retained under this equal-institution counting rule, subject to complete affiliation coverage. It does not by itself rank the average affiliation count or collaboration quality.
\item
  \textbf{OA metadata}: Differences in observed OA shares do not establish the effects of national mandates or institutional policies.
\item
  \textbf{No single ranking}: Report coverage, sample design and indicator definitions alongside any comparison \citep{hicks2015}.
\end{enumerate}

\section{Lessons learned}\label{lessons-learned-1}

\begin{itemize}
\tightlist
\item
  Treat these outputs as illustrative sampled comparisons, not population rankings.
\item
  Field mix and differing citation windows confound comparisons.
\item
  Responsible benchmarking requires a justified reference set, transparent counting rules, missing-data coverage and uncertainty analysis.
\end{itemize}

\chapter{Case Study 3: Journal Portfolio Analysis}\label{case-study-3-journal-portfolio-analysis}

\emph{R. Heller · Date: 2026-05-09 · Last changed: 2026-10-31}

\section{Objective}\label{objective-2}

Compare samples from three information-science journals using cumulative citation distributions, publication-cohort summaries, and assigned topics as a starting point for collection review.

\section{Setup}\label{setup-2}

\begin{Shaded}
\begin{Highlighting}[]
\FunctionTok{library}\NormalTok{(tidyverse)}
\FunctionTok{library}\NormalTok{(openalexR)}
\FunctionTok{library}\NormalTok{(tidytext)}
\FunctionTok{library}\NormalTok{(glue)}
\FunctionTok{library}\NormalTok{(gt)}

\FunctionTok{set.seed}\NormalTok{(}\DecValTok{20260509}\NormalTok{)}

\FunctionTok{source}\NormalTok{(here}\SpecialCharTok{::}\FunctionTok{here}\NormalTok{(}\StringTok{"R"}\NormalTok{, }\StringTok{"api\_helpers.R"}\NormalTok{))}
\FunctionTok{source}\NormalTok{(here}\SpecialCharTok{::}\FunctionTok{here}\NormalTok{(}\StringTok{"R"}\NormalTok{, }\StringTok{"utils.R"}\NormalTok{))}
\FunctionTok{source}\NormalTok{(here}\SpecialCharTok{::}\FunctionTok{here}\NormalTok{(}\StringTok{"R"}\NormalTok{, }\StringTok{"sci\_palette.R"}\NormalTok{))}
\end{Highlighting}
\end{Shaded}

\section{Data acquisition}\label{data-acquisition-2}

The source identifiers resolve to \href{https://api.openalex.org/sources/S148561398}{Scientometrics}, \href{https://api.openalex.org/sources/S205292342}{Journal of Informetrics}, and the \href{https://api.openalex.org/sources/S4210197613}{Journal of the Association for Information Science and Technology} (JASIST). These OpenAlex source records were checked on 2026-10-30. Historical titles of JASIST have separate source records and should not be substituted for this recent publication window.

\begin{Shaded}
\begin{Highlighting}[]
\NormalTok{journals }\OtherTok{\textless{}{-}} \FunctionTok{tribble}\NormalTok{(}
  \SpecialCharTok{\textasciitilde{}}\NormalTok{short\_name,        }\SpecialCharTok{\textasciitilde{}}\NormalTok{source\_id,}
  \StringTok{"Scientometrics"}\NormalTok{,   }\StringTok{"S148561398"}\NormalTok{,}
  \StringTok{"J. Informetrics"}\NormalTok{,  }\StringTok{"S205292342"}\NormalTok{,}
  \StringTok{"JASIST"}\NormalTok{,           }\StringTok{"S4210197613"}
\NormalTok{)}
\end{Highlighting}
\end{Shaded}

Request up to 300 randomly sampled articles per journal from 2018-2023, using a fixed seed. These samples do not measure the journals' total output. The seed stabilises sampling within an OpenAlex snapshot; later database updates can change the returned records and citation counts. Preserve the returned works and cache files with any reported analysis.

\begin{Shaded}
\begin{Highlighting}[]
\NormalTok{fetch\_journal }\OtherTok{\textless{}{-}} \ControlFlowTok{function}\NormalTok{(sid) \{}
\NormalTok{  works }\OtherTok{\textless{}{-}} \FunctionTok{fetch\_openalex}\NormalTok{(}
    \AttributeTok{entity =} \StringTok{"works"}\NormalTok{,}
    \AttributeTok{primary\_location.source.id =}\NormalTok{ sid,}
    \AttributeTok{from\_publication\_date =} \StringTok{"2018{-}01{-}01"}\NormalTok{,}
    \AttributeTok{to\_publication\_date =} \StringTok{"2023{-}12{-}31"}\NormalTok{,}
    \AttributeTok{type =} \StringTok{"article"}\NormalTok{,}
    \AttributeTok{options =} \FunctionTok{list}\NormalTok{(}\AttributeTok{sample =} \DecValTok{300}\NormalTok{, }\AttributeTok{seed =} \DecValTok{42}\NormalTok{)}
\NormalTok{  )}
\NormalTok{  required }\OtherTok{\textless{}{-}} \FunctionTok{c}\NormalTok{(}\StringTok{"id"}\NormalTok{, }\StringTok{"publication\_year"}\NormalTok{, }\StringTok{"cited\_by\_count"}\NormalTok{, }\StringTok{"topics"}\NormalTok{)}
  \ControlFlowTok{if}\NormalTok{ (}\SpecialCharTok{!}\FunctionTok{is.data.frame}\NormalTok{(works) }\SpecialCharTok{||} \FunctionTok{nrow}\NormalTok{(works) }\SpecialCharTok{==} \DecValTok{0}\DataTypeTok{L} \SpecialCharTok{||}
      \SpecialCharTok{!}\FunctionTok{all}\NormalTok{(required }\SpecialCharTok{\%in\%} \FunctionTok{names}\NormalTok{(works))) \{}
    \FunctionTok{stop}\NormalTok{(}\StringTok{"No usable journal sample returned for "}\NormalTok{, sid)}
\NormalTok{  \}}
  \ControlFlowTok{if}\NormalTok{ (}\FunctionTok{anyNA}\NormalTok{(works}\SpecialCharTok{$}\NormalTok{id) }\SpecialCharTok{||} \FunctionTok{anyDuplicated}\NormalTok{(works}\SpecialCharTok{$}\NormalTok{id) }\SpecialCharTok{||}
      \FunctionTok{anyNA}\NormalTok{(works}\SpecialCharTok{$}\NormalTok{publication\_year) }\SpecialCharTok{||}
      \FunctionTok{any}\NormalTok{(}\SpecialCharTok{!}\NormalTok{works}\SpecialCharTok{$}\NormalTok{publication\_year }\SpecialCharTok{\%in\%} \DecValTok{2018}\SpecialCharTok{:}\DecValTok{2023}\NormalTok{) }\SpecialCharTok{||}
      \FunctionTok{any}\NormalTok{(}\SpecialCharTok{!}\FunctionTok{is.finite}\NormalTok{(works}\SpecialCharTok{$}\NormalTok{cited\_by\_count)) }\SpecialCharTok{||} \FunctionTok{any}\NormalTok{(works}\SpecialCharTok{$}\NormalTok{cited\_by\_count }\SpecialCharTok{\textless{}} \DecValTok{0}\NormalTok{)) \{}
    \FunctionTok{stop}\NormalTok{(}\StringTok{"Invalid or duplicate records in journal sample for "}\NormalTok{, sid)}
\NormalTok{  \}}
\NormalTok{  works}
\NormalTok{\}}

\NormalTok{journal\_data }\OtherTok{\textless{}{-}}\NormalTok{ journals }\SpecialCharTok{|\textgreater{}}
  \FunctionTok{mutate}\NormalTok{(}\AttributeTok{works =} \FunctionTok{map}\NormalTok{(source\_id, fetch\_journal))}
\end{Highlighting}
\end{Shaded}

\section{Citation distributions in the samples}\label{citation-distributions-in-the-samples}

\texttt{cited\_by\_count} is the cumulative number of citations recorded when the data were retrieved, including citations received after 2023. It does not restrict citing papers to the sampled publication years. The h-index below describes only the returned sample, and all summaries remain unnormalised for field and publication year.

\begin{Shaded}
\begin{Highlighting}[]
\NormalTok{all\_works }\OtherTok{\textless{}{-}}\NormalTok{ journal\_data }\SpecialCharTok{|\textgreater{}}
  \FunctionTok{mutate}\NormalTok{(}\AttributeTok{works =} \FunctionTok{map2}\NormalTok{(works, short\_name, \textbackslash{}(w, n) w }\SpecialCharTok{|\textgreater{}} \FunctionTok{mutate}\NormalTok{(}\AttributeTok{journal =}\NormalTok{ n))) }\SpecialCharTok{|\textgreater{}}
  \FunctionTok{pull}\NormalTok{(works) }\SpecialCharTok{|\textgreater{}}
  \FunctionTok{bind\_rows}\NormalTok{()}

\FunctionTok{ggplot}\NormalTok{(all\_works, }\FunctionTok{aes}\NormalTok{(}\AttributeTok{x =}\NormalTok{ journal, }\AttributeTok{y =}\NormalTok{ cited\_by\_count }\SpecialCharTok{+} \DecValTok{1}\NormalTok{)) }\SpecialCharTok{+}
  \FunctionTok{geom\_boxplot}\NormalTok{(}\AttributeTok{fill =} \FunctionTok{palette\_sci}\NormalTok{(}\DecValTok{1}\NormalTok{), }\AttributeTok{alpha =} \FloatTok{0.7}\NormalTok{) }\SpecialCharTok{+}
  \FunctionTok{scale\_y\_log10}\NormalTok{() }\SpecialCharTok{+}
  \FunctionTok{labs}\NormalTok{(}\AttributeTok{x =} \ConstantTok{NULL}\NormalTok{, }\AttributeTok{y =} \StringTok{"Cumulative citations + 1 (log10 scale)"}\NormalTok{) }\SpecialCharTok{+}
  \FunctionTok{theme\_sci}\NormalTok{()}
\end{Highlighting}
\end{Shaded}

\begin{figure}

{\centering \includegraphics[width=0.9\linewidth,alt={Box plots comparing cumulative citations plus one across the three journal samples.}]{scientometrics-in-r_files/figure-latex/cs03-fig-cs03-impact-1} 

}

\caption{Cumulative citation distributions in the sampled articles, with one added before log scaling.}\label{fig:cs03-fig-cs03-impact}
\end{figure}

\begin{Shaded}
\begin{Highlighting}[]
\NormalTok{all\_works }\SpecialCharTok{|\textgreater{}}
  \FunctionTok{group\_by}\NormalTok{(journal) }\SpecialCharTok{|\textgreater{}}
  \FunctionTok{summarise}\NormalTok{(}
    \AttributeTok{sampled\_articles =} \FunctionTok{n}\NormalTok{(),}
    \AttributeTok{mean\_cites =} \FunctionTok{round}\NormalTok{(}\FunctionTok{mean}\NormalTok{(cited\_by\_count), }\DecValTok{1}\NormalTok{),}
    \AttributeTok{median\_cites =} \FunctionTok{median}\NormalTok{(cited\_by\_count),}
    \AttributeTok{sample\_h\_index =} \FunctionTok{compute\_h\_index}\NormalTok{(cited\_by\_count),}
    \AttributeTok{.groups =} \StringTok{"drop"}
\NormalTok{  ) }\SpecialCharTok{|\textgreater{}}
  \FunctionTok{gt}\NormalTok{() }\SpecialCharTok{|\textgreater{}}
  \FunctionTok{cols\_label}\NormalTok{(}
    \AttributeTok{journal =} \StringTok{"Journal"}\NormalTok{, }\AttributeTok{sampled\_articles =} \StringTok{"Sampled articles"}\NormalTok{,}
    \AttributeTok{mean\_cites =} \StringTok{"Mean cumulative citations"}\NormalTok{, }\AttributeTok{median\_cites =} \StringTok{"Median cumulative citations"}\NormalTok{,}
    \AttributeTok{sample\_h\_index =} \StringTok{"Sample h{-}index"}
\NormalTok{  )}
\end{Highlighting}
\end{Shaded}

\begingroup
\fontsize{8.0pt}{9.0pt}\selectfont
\begin{longtable}{>{\raggedright\arraybackslash}p{\dimexpr 0.21\linewidth -2\tabcolsep-1.5\arrayrulewidth}>{\raggedleft\arraybackslash}p{\dimexpr 0.22\linewidth -2\tabcolsep-1.5\arrayrulewidth}>{\raggedleft\arraybackslash}p{\dimexpr 0.17\linewidth -2\tabcolsep-1.5\arrayrulewidth}>{\raggedleft\arraybackslash}p{\dimexpr 0.19\linewidth -2\tabcolsep-1.5\arrayrulewidth}>{\raggedleft\arraybackslash}p{\dimexpr 0.21\linewidth -2\tabcolsep-1.5\arrayrulewidth}}
\toprule
Journal & Sampled articles & Mean cumulative citations & Median cumulative citations & Sample h-index \\ 
\midrule\addlinespace[2.5pt]
J. Informetrics & 300 & 30.7 & 18.5 & 49 \\ 
JASIST & 300 & 33.6 & 19.5 & 50 \\ 
\nolinkurl{Scientometrics} & 300 & 27.3 & 17.5 & 44 \\ 
\bottomrule
\end{longtable}
\endgroup

\section{Publication-cohort comparison}\label{publication-cohort-comparison}

The following plot compares different publication cohorts at retrieval. Older cohorts have had more time to receive citations, and their composition may also differ. It is not a trajectory following the same articles over time, so it cannot establish citation-aging curves or citation half-lives. That requires time-resolved citation counts and an explicit observation window.

\begin{Shaded}
\begin{Highlighting}[]
\NormalTok{cohorts }\OtherTok{\textless{}{-}}\NormalTok{ all\_works }\SpecialCharTok{|\textgreater{}}
  \FunctionTok{group\_by}\NormalTok{(journal, publication\_year) }\SpecialCharTok{|\textgreater{}}
  \FunctionTok{summarise}\NormalTok{(}\AttributeTok{mean\_cites =} \FunctionTok{mean}\NormalTok{(cited\_by\_count), }\AttributeTok{sampled\_articles =} \FunctionTok{n}\NormalTok{(), }\AttributeTok{.groups =} \StringTok{"drop"}\NormalTok{)}

\FunctionTok{ggplot}\NormalTok{(cohorts, }\FunctionTok{aes}\NormalTok{(}\AttributeTok{x =}\NormalTok{ publication\_year, }\AttributeTok{y =}\NormalTok{ mean\_cites, }\AttributeTok{colour =}\NormalTok{ journal)) }\SpecialCharTok{+}
  \FunctionTok{geom\_line}\NormalTok{(}\AttributeTok{linewidth =} \DecValTok{1}\NormalTok{) }\SpecialCharTok{+}
  \FunctionTok{geom\_point}\NormalTok{(}\FunctionTok{aes}\NormalTok{(}\AttributeTok{size =}\NormalTok{ sampled\_articles)) }\SpecialCharTok{+}
  \FunctionTok{scale\_x\_continuous}\NormalTok{(}\AttributeTok{breaks =} \DecValTok{2018}\SpecialCharTok{:}\DecValTok{2023}\NormalTok{) }\SpecialCharTok{+}
  \FunctionTok{scale\_colour\_manual}\NormalTok{(}\AttributeTok{values =} \FunctionTok{palette\_sci}\NormalTok{(}\DecValTok{3}\NormalTok{)) }\SpecialCharTok{+}
  \FunctionTok{labs}\NormalTok{(}\AttributeTok{x =} \StringTok{"Publication year"}\NormalTok{, }\AttributeTok{y =} \StringTok{"Mean cumulative citations"}\NormalTok{,}
       \AttributeTok{colour =} \StringTok{"Journal"}\NormalTok{, }\AttributeTok{size =} \StringTok{"Sampled articles"}\NormalTok{) }\SpecialCharTok{+}
  \FunctionTok{theme\_sci}\NormalTok{()}
\end{Highlighting}
\end{Shaded}

\begin{figure}

{\centering \includegraphics[width=0.9\linewidth,alt={Lines compare mean cumulative citations for different publication cohorts, with larger points indicating more sampled articles.}]{scientometrics-in-r_files/figure-latex/cs03-fig-cs03-aging-1} 

}

\caption{Mean cumulative citations by publication year in each journal sample; point size shows cohort sample size.}\label{fig:cs03-fig-cs03-aging}
\end{figure}

\section{Topical coverage}\label{topical-coverage}

Count each article once per assigned topic. The flattened \texttt{topics} table also contains the topics' parent subfields, fields, and domains; retain only topic identifiers (\texttt{T...}). An article can contribute to several topics, while articles without topic assignments do not contribute to this chart; topic counts therefore need not sum to the number of sampled articles.

\begin{Shaded}
\begin{Highlighting}[]
\NormalTok{topic\_data }\OtherTok{\textless{}{-}}\NormalTok{ all\_works }\SpecialCharTok{|\textgreater{}}
  \FunctionTok{select}\NormalTok{(id, journal, topics) }\SpecialCharTok{|\textgreater{}}
  \FunctionTok{unnest}\NormalTok{(topics, }\AttributeTok{names\_sep =} \StringTok{"\_"}\NormalTok{) }\SpecialCharTok{|\textgreater{}}
  \FunctionTok{filter}\NormalTok{(}\FunctionTok{str\_detect}\NormalTok{(topics\_id, }\StringTok{"(\^{}|/)T[0{-}9]+$"}\NormalTok{)) }\SpecialCharTok{|\textgreater{}}
  \FunctionTok{select}\NormalTok{(id, journal, }\AttributeTok{topic =}\NormalTok{ topics\_display\_name) }\SpecialCharTok{|\textgreater{}}
  \FunctionTok{mutate}\NormalTok{(}\AttributeTok{topic =} \FunctionTok{str\_to\_lower}\NormalTok{(}\FunctionTok{str\_trim}\NormalTok{(topic))) }\SpecialCharTok{|\textgreater{}}
  \FunctionTok{filter}\NormalTok{(}\SpecialCharTok{!}\FunctionTok{is.na}\NormalTok{(topic), }\FunctionTok{nchar}\NormalTok{(topic) }\SpecialCharTok{\textgreater{}=} \DecValTok{3}\NormalTok{) }\SpecialCharTok{|\textgreater{}}
  \FunctionTok{distinct}\NormalTok{(id, journal, topic)}

\NormalTok{top\_topics }\OtherTok{\textless{}{-}}\NormalTok{ topic\_data }\SpecialCharTok{|\textgreater{}}
  \FunctionTok{count}\NormalTok{(journal, topic, }\AttributeTok{sort =} \ConstantTok{TRUE}\NormalTok{) }\SpecialCharTok{|\textgreater{}}
  \FunctionTok{group\_by}\NormalTok{(journal) }\SpecialCharTok{|\textgreater{}}
  \FunctionTok{slice\_max}\NormalTok{(n, }\AttributeTok{n =} \DecValTok{10}\NormalTok{, }\AttributeTok{with\_ties =} \ConstantTok{FALSE}\NormalTok{) }\SpecialCharTok{|\textgreater{}}
  \FunctionTok{ungroup}\NormalTok{()}
\end{Highlighting}
\end{Shaded}

\begin{Shaded}
\begin{Highlighting}[]
\NormalTok{top\_topics }\SpecialCharTok{|\textgreater{}}
  \FunctionTok{mutate}\NormalTok{(}\AttributeTok{topic =} \FunctionTok{reorder\_within}\NormalTok{(topic, n, journal)) }\SpecialCharTok{|\textgreater{}}
  \FunctionTok{ggplot}\NormalTok{(}\FunctionTok{aes}\NormalTok{(}\AttributeTok{x =}\NormalTok{ n, }\AttributeTok{y =}\NormalTok{ topic)) }\SpecialCharTok{+}
  \FunctionTok{geom\_col}\NormalTok{(}\AttributeTok{fill =} \FunctionTok{palette\_sci}\NormalTok{(}\DecValTok{1}\NormalTok{)) }\SpecialCharTok{+}
  \FunctionTok{facet\_wrap}\NormalTok{(}\SpecialCharTok{\textasciitilde{}}\NormalTok{ journal, }\AttributeTok{scales =} \StringTok{"free\_y"}\NormalTok{, }\AttributeTok{ncol =} \DecValTok{1}\NormalTok{) }\SpecialCharTok{+}
  \FunctionTok{scale\_y\_reordered}\NormalTok{() }\SpecialCharTok{+}
  \FunctionTok{labs}\NormalTok{(}\AttributeTok{x =} \StringTok{"Sampled articles assigned to topic"}\NormalTok{, }\AttributeTok{y =} \ConstantTok{NULL}\NormalTok{) }\SpecialCharTok{+}
  \FunctionTok{theme\_sci}\NormalTok{(}\AttributeTok{base\_size =} \DecValTok{9}\NormalTok{)}
\end{Highlighting}
\end{Shaded}

\begin{figure}

{\centering \includegraphics[width=0.9\linewidth,alt={Faceted bar chart comparing numbers of sampled articles assigned to the most frequent topics in each journal.}]{scientometrics-in-r_files/figure-latex/cs03-fig-cs03-topics-1} 

}

\caption{Up to 10 most frequent assigned topics in each journal sample.}\label{fig:cs03-fig-cs03-topics}
\end{figure}

\section{Reading the results}\label{reading-the-results-1}

\begin{enumerate}
\def\labelenumi{\arabic{enumi}.}
\tightlist
\item
  \textbf{Citation distributions}: Compare the plotted medians with the sample means. A large gap can indicate that a few highly cited articles dominate the mean; journal rankings should be reported only from the observed results.
\item
  \textbf{Publication cohorts}: Compare cohorts alongside their sample sizes. Differences can reflect time available for citations, article composition, or sampling variation; they do not demonstrate a common aging curve.
\item
  \textbf{Assigned topics}: Inspect overlap and differences in the displayed topics, then review representative articles before drawing conclusions about collection coverage.
\end{enumerate}

\section{Lessons learned}\label{lessons-learned-2}

\begin{itemize}
\tightlist
\item
  Journal comparison requires multiple dimensions; no single metric tells the full story.
\item
  Sample-based analysis is illustrative. Dependable journal evaluation requires a defined corpus, justified sampling or complete coverage, appropriate observation windows, and field-normalised indicators where relevant \citep{waltman2016}.
\item
  Collection decisions also depend on local readership, access, cost, and subject needs; citation and topic summaries alone do not measure a journal's value to a library.
\end{itemize}

\chapter{Case Study 4: Gender Gap in Scientific Publishing}\label{case-study-4-gender-gap-in-scientific-publishing}

\emph{R. Heller · Date: 2026-05-09 · Last changed: 2026-10-31}

\section{Objective}\label{objective-3}

Examine why a name lookup cannot establish a gender gap, using an explicitly
synthetic example of two fictional journals. We compare arbitrary lookup
categories, retain unmatched records in every denominator, and track lookup
coverage by recorded authorship position and year. None of the records represents
a real person, paper or journal; no gender identity is observed or inferred.

An actual study of gender equity requires a defensible demographic variable and
its provenance, a defined population, appropriate data governance, and a plan
for validation and missing data. A list of names is not a substitute for those
observations. Chapter \ref{gender-objectives} develops these distinctions.

\section{Setup}\label{setup-3}

\begin{Shaded}
\begin{Highlighting}[]
\FunctionTok{library}\NormalTok{(tidyverse)}
\FunctionTok{library}\NormalTok{(glue)}
\FunctionTok{library}\NormalTok{(gt)}

\FunctionTok{source}\NormalTok{(here}\SpecialCharTok{::}\FunctionTok{here}\NormalTok{(}\StringTok{"R"}\NormalTok{, }\StringTok{"sci\_palette.R"}\NormalTok{))}
\end{Highlighting}
\end{Shaded}

\section{Synthetic teaching records}\label{synthetic-teaching-records}

The six invented papers have explicit publication years and ordered name lists.
The names are placeholders, and repeated names do not establish shared author
identity. Each row created below is an authorship record, not a distinct person.
This example makes no API request.

\begin{Shaded}
\begin{Highlighting}[]
\NormalTok{works\_example }\OtherTok{\textless{}{-}} \FunctionTok{tibble}\NormalTok{(}
  \AttributeTok{work\_id =} \FunctionTok{c}\NormalTok{(}\StringTok{"example{-}a1"}\NormalTok{, }\StringTok{"example{-}a2"}\NormalTok{, }\StringTok{"example{-}a3"}\NormalTok{,}
              \StringTok{"example{-}b1"}\NormalTok{, }\StringTok{"example{-}b2"}\NormalTok{, }\StringTok{"example{-}b3"}\NormalTok{),}
  \AttributeTok{journal =} \FunctionTok{rep}\NormalTok{(}\FunctionTok{c}\NormalTok{(}\StringTok{"Example journal A"}\NormalTok{, }\StringTok{"Example journal B"}\NormalTok{), }\AttributeTok{each =} \DecValTok{3}\NormalTok{),}
  \AttributeTok{year =} \FunctionTok{rep}\NormalTok{(}\DecValTok{2019}\SpecialCharTok{:}\DecValTok{2021}\NormalTok{, }\DecValTok{2}\NormalTok{),}
  \AttributeTok{first\_names =} \FunctionTok{list}\NormalTok{(}\FunctionTok{c}\NormalTok{(}\StringTok{"aria"}\NormalTok{, }\StringTok{"jordan"}\NormalTok{, }\StringTok{"lee"}\NormalTok{), }\FunctionTok{c}\NormalTok{(}\StringTok{"noor"}\NormalTok{, }\StringTok{"ari"}\NormalTok{), }\StringTok{"morgan"}\NormalTok{,}
                     \FunctionTok{c}\NormalTok{(}\StringTok{"aria"}\NormalTok{, }\StringTok{"nuri"}\NormalTok{, }\StringTok{"sam"}\NormalTok{, }\StringTok{"alex"}\NormalTok{), }\FunctionTok{c}\NormalTok{(}\StringTok{"ari"}\NormalTok{, }\StringTok{"noor"}\NormalTok{),}
                     \FunctionTok{c}\NormalTok{(}\StringTok{"lee"}\NormalTok{, }\StringTok{"jordan"}\NormalTok{, }\StringTok{"morgan"}\NormalTok{))}
\NormalTok{)}
\end{Highlighting}
\end{Shaded}

\section{Arbitrary lookup and complete denominators}\label{arbitrary-lookup-and-complete-denominators}

For teaching purposes only, assign \texttt{aria} and \texttt{ari} to A, and \texttt{jordan} and
\texttt{morgan} to B. These arbitrary categories have \textbf{no demographic meaning}.
Every other name is explicitly unmatched. Dropping unmatched records would
change both the denominator and the question answered by a proportion.

\begin{Shaded}
\begin{Highlighting}[]
\NormalTok{authors\_example }\OtherTok{\textless{}{-}}\NormalTok{ works\_example }\SpecialCharTok{|\textgreater{}}
  \FunctionTok{unnest\_longer}\NormalTok{(first\_names, }\AttributeTok{values\_to =} \StringTok{"first\_name"}\NormalTok{, }\AttributeTok{indices\_to =} \StringTok{"author\_order"}\NormalTok{) }\SpecialCharTok{|\textgreater{}}
  \FunctionTok{group\_by}\NormalTok{(work\_id) }\SpecialCharTok{|\textgreater{}}
  \FunctionTok{mutate}\NormalTok{(}
    \AttributeTok{n\_authors =} \FunctionTok{n}\NormalTok{(),}
    \AttributeTok{position =} \FunctionTok{case\_when}\NormalTok{(}
\NormalTok{      n\_authors }\SpecialCharTok{==} \DecValTok{1} \SpecialCharTok{\textasciitilde{}} \StringTok{"single"}\NormalTok{,}
\NormalTok{      author\_order }\SpecialCharTok{==} \DecValTok{1} \SpecialCharTok{\textasciitilde{}} \StringTok{"first"}\NormalTok{,}
\NormalTok{      author\_order }\SpecialCharTok{==}\NormalTok{ n\_authors }\SpecialCharTok{\textasciitilde{}} \StringTok{"last"}\NormalTok{,}
      \ConstantTok{TRUE} \SpecialCharTok{\textasciitilde{}} \StringTok{"middle"}
\NormalTok{    )}
\NormalTok{  ) }\SpecialCharTok{|\textgreater{}}
  \FunctionTok{ungroup}\NormalTok{() }\SpecialCharTok{|\textgreater{}}
  \FunctionTok{mutate}\NormalTok{(}
    \AttributeTok{lookup\_category =} \FunctionTok{case\_when}\NormalTok{(}
\NormalTok{      first\_name }\SpecialCharTok{\%in\%} \FunctionTok{c}\NormalTok{(}\StringTok{"aria"}\NormalTok{, }\StringTok{"ari"}\NormalTok{) }\SpecialCharTok{\textasciitilde{}} \StringTok{"A"}\NormalTok{,}
\NormalTok{      first\_name }\SpecialCharTok{\%in\%} \FunctionTok{c}\NormalTok{(}\StringTok{"jordan"}\NormalTok{, }\StringTok{"morgan"}\NormalTok{) }\SpecialCharTok{\textasciitilde{}} \StringTok{"B"}\NormalTok{,}
      \ConstantTok{TRUE} \SpecialCharTok{\textasciitilde{}} \StringTok{"Unmatched"}
\NormalTok{    ),}
    \AttributeTok{lookup\_category =} \FunctionTok{factor}\NormalTok{(lookup\_category, }\AttributeTok{levels =} \FunctionTok{c}\NormalTok{(}\StringTok{"A"}\NormalTok{, }\StringTok{"B"}\NormalTok{, }\StringTok{"Unmatched"}\NormalTok{)),}
    \AttributeTok{position =} \FunctionTok{factor}\NormalTok{(position, }\AttributeTok{levels =} \FunctionTok{c}\NormalTok{(}\StringTok{"first"}\NormalTok{, }\StringTok{"middle"}\NormalTok{, }\StringTok{"last"}\NormalTok{, }\StringTok{"single"}\NormalTok{))}
\NormalTok{  )}

\NormalTok{coverage\_summary }\OtherTok{\textless{}{-}}\NormalTok{ authors\_example }\SpecialCharTok{|\textgreater{}}
  \FunctionTok{group\_by}\NormalTok{(journal) }\SpecialCharTok{|\textgreater{}}
  \FunctionTok{summarise}\NormalTok{(}\AttributeTok{records =} \FunctionTok{n}\NormalTok{(), }\AttributeTok{matched =} \FunctionTok{sum}\NormalTok{(lookup\_category }\SpecialCharTok{!=} \StringTok{"Unmatched"}\NormalTok{),}
            \AttributeTok{coverage =}\NormalTok{ matched }\SpecialCharTok{/}\NormalTok{ records, }\AttributeTok{.groups =} \StringTok{"drop"}\NormalTok{)}
\NormalTok{coverage\_summary }\SpecialCharTok{|\textgreater{}} \FunctionTok{gt}\NormalTok{() }\SpecialCharTok{|\textgreater{}} \FunctionTok{fmt\_percent}\NormalTok{(}\AttributeTok{columns =}\NormalTok{ coverage, }\AttributeTok{decimals =} \DecValTok{1}\NormalTok{)}
\end{Highlighting}
\end{Shaded}

\begingroup
\fontsize{8.0pt}{9.0pt}\selectfont
\begin{longtable}{>{\raggedright\arraybackslash}p{\dimexpr 0.33\linewidth -2\tabcolsep-1.5\arrayrulewidth}>{\raggedleft\arraybackslash}p{\dimexpr 0.22\linewidth -2\tabcolsep-1.5\arrayrulewidth}>{\raggedleft\arraybackslash}p{\dimexpr 0.22\linewidth -2\tabcolsep-1.5\arrayrulewidth}>{\raggedleft\arraybackslash}p{\dimexpr 0.22\linewidth -2\tabcolsep-1.5\arrayrulewidth}}
\toprule
journal & records & matched & \nolinkurl{coverage} \\ 
\midrule\addlinespace[2.5pt]
Example journal A & 6 & 4 & 66.7\% \\ 
Example journal B & 9 & 4 & 44.4\% \\ 
\bottomrule
\end{longtable}
\endgroup

\section{Lookup categories by fictional journal}\label{lookup-categories-by-fictional-journal}

\begin{Shaded}
\begin{Highlighting}[]
\NormalTok{journal\_summary }\OtherTok{\textless{}{-}}\NormalTok{ authors\_example }\SpecialCharTok{|\textgreater{}}
  \FunctionTok{count}\NormalTok{(journal, lookup\_category, }\AttributeTok{.drop =} \ConstantTok{FALSE}\NormalTok{) }\SpecialCharTok{|\textgreater{}}
  \FunctionTok{group\_by}\NormalTok{(journal) }\SpecialCharTok{|\textgreater{}}
  \FunctionTok{mutate}\NormalTok{(}\AttributeTok{pct =}\NormalTok{ n }\SpecialCharTok{/} \FunctionTok{sum}\NormalTok{(n)) }\SpecialCharTok{|\textgreater{}}
  \FunctionTok{ungroup}\NormalTok{()}

\FunctionTok{ggplot}\NormalTok{(journal\_summary, }\FunctionTok{aes}\NormalTok{(}\AttributeTok{x =}\NormalTok{ journal, }\AttributeTok{y =}\NormalTok{ pct, }\AttributeTok{fill =}\NormalTok{ lookup\_category)) }\SpecialCharTok{+}
  \FunctionTok{geom\_col}\NormalTok{(}\AttributeTok{position =} \StringTok{"dodge"}\NormalTok{) }\SpecialCharTok{+}
  \FunctionTok{scale\_y\_continuous}\NormalTok{(}\AttributeTok{labels =}\NormalTok{ scales}\SpecialCharTok{::}\NormalTok{percent) }\SpecialCharTok{+}
  \FunctionTok{scale\_fill\_manual}\NormalTok{(}\AttributeTok{values =} \FunctionTok{palette\_sci}\NormalTok{(}\DecValTok{3}\NormalTok{), }\AttributeTok{drop =} \ConstantTok{FALSE}\NormalTok{) }\SpecialCharTok{+}
  \FunctionTok{labs}\NormalTok{(}\AttributeTok{x =} \ConstantTok{NULL}\NormalTok{, }\AttributeTok{y =} \StringTok{"Share of all synthetic records"}\NormalTok{, }\AttributeTok{fill =} \StringTok{"Lookup category"}\NormalTok{) }\SpecialCharTok{+}
  \FunctionTok{theme\_sci}\NormalTok{()}
\end{Highlighting}
\end{Shaded}

\begin{figure}

{\centering \includegraphics[width=0.9\linewidth,alt={Grouped bars show A, B and Unmatched shares using every synthetic record in each fictional journal as the denominator.}]{scientometrics-in-r_files/figure-latex/cs04-fig-cs04-journal-1} 

}

\caption{Arbitrary lookup categories among all synthetic authorship records in two fictional journals; these are not gender proportions.}\label{fig:cs04-fig-cs04-journal}
\end{figure}

\section{Recorded authorship position}\label{recorded-authorship-position}

Single-author papers receive their own position label. In a real corpus,
recorded order does not by itself establish seniority or contribution: the
interpretation depends on authorship conventions. The following proportions
retain unmatched records within every observed journal/position group.

\begin{Shaded}
\begin{Highlighting}[]
\NormalTok{position\_summary\_example }\OtherTok{\textless{}{-}}\NormalTok{ authors\_example }\SpecialCharTok{|\textgreater{}}
  \FunctionTok{count}\NormalTok{(journal, position, lookup\_category) }\SpecialCharTok{|\textgreater{}}
  \FunctionTok{group\_by}\NormalTok{(journal, position) }\SpecialCharTok{|\textgreater{}}
  \FunctionTok{mutate}\NormalTok{(}\AttributeTok{pct =}\NormalTok{ n }\SpecialCharTok{/} \FunctionTok{sum}\NormalTok{(n)) }\SpecialCharTok{|\textgreater{}}
  \FunctionTok{ungroup}\NormalTok{()}

\FunctionTok{ggplot}\NormalTok{(position\_summary\_example, }\FunctionTok{aes}\NormalTok{(}\AttributeTok{x =}\NormalTok{ position, }\AttributeTok{y =}\NormalTok{ pct, }\AttributeTok{fill =}\NormalTok{ lookup\_category)) }\SpecialCharTok{+}
  \FunctionTok{geom\_col}\NormalTok{(}\AttributeTok{position =} \StringTok{"dodge"}\NormalTok{) }\SpecialCharTok{+}
  \FunctionTok{facet\_wrap}\NormalTok{(}\SpecialCharTok{\textasciitilde{}}\NormalTok{ journal) }\SpecialCharTok{+}
  \FunctionTok{scale\_x\_discrete}\NormalTok{(}\AttributeTok{drop =} \ConstantTok{FALSE}\NormalTok{) }\SpecialCharTok{+}
  \FunctionTok{scale\_y\_continuous}\NormalTok{(}\AttributeTok{labels =}\NormalTok{ scales}\SpecialCharTok{::}\NormalTok{percent) }\SpecialCharTok{+}
  \FunctionTok{scale\_fill\_manual}\NormalTok{(}\AttributeTok{values =} \FunctionTok{palette\_sci}\NormalTok{(}\DecValTok{3}\NormalTok{), }\AttributeTok{drop =} \ConstantTok{FALSE}\NormalTok{) }\SpecialCharTok{+}
  \FunctionTok{labs}\NormalTok{(}\AttributeTok{x =} \StringTok{"Recorded position"}\NormalTok{, }\AttributeTok{y =} \StringTok{"Share of synthetic records"}\NormalTok{, }\AttributeTok{fill =} \StringTok{"Lookup category"}\NormalTok{) }\SpecialCharTok{+}
  \FunctionTok{theme\_sci}\NormalTok{()}
\end{Highlighting}
\end{Shaded}

\begin{figure}

{\centering \includegraphics[width=0.9\linewidth,alt={Faceted bars compare A, B and Unmatched shares by first, middle, last and single-author positions in two fictional journals.}]{scientometrics-in-r_files/figure-latex/cs04-fig-cs04-position-1} 

}

\caption{Arbitrary lookup categories by recorded position in the synthetic fixture, including unmatched records. Empty positions have no observations.}\label{fig:cs04-fig-cs04-position}
\end{figure}

\section{Coverage by recorded year}\label{coverage-by-recorded-year}

The year comes from the fixture's explicit publication-year column, never from
an identifier. For each fictional journal/year, the denominator includes all
first- or single-author records. Each group has only one such record here, so
coverage is either zero or one. This illustrates denominator accounting; it
supplies no evidence about a demographic trend.

\begin{Shaded}
\begin{Highlighting}[]
\NormalTok{coverage\_by\_year }\OtherTok{\textless{}{-}}\NormalTok{ authors\_example }\SpecialCharTok{|\textgreater{}}
  \FunctionTok{filter}\NormalTok{(position }\SpecialCharTok{\%in\%} \FunctionTok{c}\NormalTok{(}\StringTok{"first"}\NormalTok{, }\StringTok{"single"}\NormalTok{)) }\SpecialCharTok{|\textgreater{}}
  \FunctionTok{group\_by}\NormalTok{(journal, year) }\SpecialCharTok{|\textgreater{}}
  \FunctionTok{summarise}\NormalTok{(}\AttributeTok{records =} \FunctionTok{n}\NormalTok{(), }\AttributeTok{matched =} \FunctionTok{sum}\NormalTok{(lookup\_category }\SpecialCharTok{!=} \StringTok{"Unmatched"}\NormalTok{),}
            \AttributeTok{coverage =}\NormalTok{ matched }\SpecialCharTok{/}\NormalTok{ records, }\AttributeTok{.groups =} \StringTok{"drop"}\NormalTok{)}

\FunctionTok{ggplot}\NormalTok{(coverage\_by\_year, }\FunctionTok{aes}\NormalTok{(}\AttributeTok{x =}\NormalTok{ year, }\AttributeTok{y =}\NormalTok{ coverage, }\AttributeTok{colour =}\NormalTok{ journal)) }\SpecialCharTok{+}
  \FunctionTok{geom\_line}\NormalTok{(}\AttributeTok{linewidth =} \DecValTok{1}\NormalTok{) }\SpecialCharTok{+}
  \FunctionTok{geom\_point}\NormalTok{(}\AttributeTok{size =} \DecValTok{2}\NormalTok{) }\SpecialCharTok{+}
  \FunctionTok{scale\_x\_continuous}\NormalTok{(}\AttributeTok{breaks =} \DecValTok{2019}\SpecialCharTok{:}\DecValTok{2021}\NormalTok{) }\SpecialCharTok{+}
  \FunctionTok{scale\_y\_continuous}\NormalTok{(}\AttributeTok{labels =}\NormalTok{ scales}\SpecialCharTok{::}\NormalTok{percent, }\AttributeTok{limits =} \FunctionTok{c}\NormalTok{(}\DecValTok{0}\NormalTok{, }\DecValTok{1}\NormalTok{)) }\SpecialCharTok{+}
  \FunctionTok{scale\_colour\_manual}\NormalTok{(}\AttributeTok{values =} \FunctionTok{palette\_sci}\NormalTok{(}\DecValTok{2}\NormalTok{)) }\SpecialCharTok{+}
  \FunctionTok{labs}\NormalTok{(}\AttributeTok{x =} \StringTok{"Recorded fixture year"}\NormalTok{, }\AttributeTok{y =} \StringTok{"Lookup coverage"}\NormalTok{, }\AttributeTok{colour =} \StringTok{"Fictional journal"}\NormalTok{) }\SpecialCharTok{+}
  \FunctionTok{theme\_sci}\NormalTok{()}
\end{Highlighting}
\end{Shaded}

\begin{figure}

{\centering \includegraphics[width=0.9\linewidth,alt={Lines show whether the one first-or-single-author record in each fictional journal/year matches an arbitrary lookup category.}]{scientometrics-in-r_files/figure-latex/cs04-fig-cs04-temporal-1} 

}

\caption{Lookup coverage among first- or single-author records by explicit fixture year. Each fictional journal/year has one record; no demographic trend is measured.}\label{fig:cs04-fig-cs04-temporal}
\end{figure}

\section{What the fixture establishes}\label{what-the-fixture-establishes}

\begin{enumerate}
\def\labelenumi{\arabic{enumi}.}
\tightlist
\item
  \textbf{Denominators change the result.} The fixture contains 15 authorship records:
  four A, four B and seven unmatched. Removing unmatched records would make A
  and B each half of the retained records, concealing that almost half of all
  records had no lookup match.
\item
  \textbf{Coverage differs by construction.} Example journal A has four matched
  records out of six; Example journal B has four out of nine. These differences
  follow from the invented name lists, not disciplinary or demographic evidence.
\item
  \textbf{Identity and contribution remain unobserved.} Neither the labels, the
  recorded positions nor their changes across fixture years support a claim
  about a gender gap, seniority or progress toward equity.
\end{enumerate}

\section{Requirements for an actual study}\label{requirements-for-an-actual-study}

A real study needs a scientific question, a defined population and suitable
observations of the demographic variable of interest. Document their provenance,
ethical basis, categories, missingness and validation. Names cannot determine a
person's self-identified gender; a name-classification service, if used for a
carefully justified question, still needs context-specific validation and an
explicit account of what its labels measure. A confidence score alone does not
validate the resulting equity estimate. See Chapter \ref{gender-limitations} and
\href{https://doi.org/10.7717/peerj-cs.156}{Santamaría and Mihaljević (2018)}.

This fixture contains no demographic observations or real journal comparison.
No conclusion about a persistent gender gap, disciplinary differences or
temporal progress can be drawn from it.

\chapter{Case Study 5: Detecting Emerging Topics}\label{case-study-5-detecting-emerging-topics}

\emph{R. Heller · Date: 2026-05-09 · Last changed: 2026-10-30}

\section{Objective}\label{objective-4}

Identify emerging topics in scientometrics research (2015--2023) using keyword burst detection and topic modelling, and validate findings against known developments in the field.

\section{Setup}\label{setup-4}

\begin{Shaded}
\begin{Highlighting}[]
\FunctionTok{library}\NormalTok{(tidyverse)}
\FunctionTok{library}\NormalTok{(openalexR)}
\FunctionTok{library}\NormalTok{(quanteda)}
\FunctionTok{library}\NormalTok{(quanteda.textstats)}
\FunctionTok{library}\NormalTok{(topicmodels)}
\FunctionTok{library}\NormalTok{(tidytext)}
\FunctionTok{library}\NormalTok{(glue)}
\FunctionTok{library}\NormalTok{(gt)}

\FunctionTok{set.seed}\NormalTok{(}\DecValTok{20260509}\NormalTok{)}

\FunctionTok{source}\NormalTok{(here}\SpecialCharTok{::}\FunctionTok{here}\NormalTok{(}\StringTok{"R"}\NormalTok{, }\StringTok{"api\_helpers.R"}\NormalTok{))}
\FunctionTok{source}\NormalTok{(here}\SpecialCharTok{::}\FunctionTok{here}\NormalTok{(}\StringTok{"R"}\NormalTok{, }\StringTok{"utils.R"}\NormalTok{))}
\FunctionTok{source}\NormalTok{(here}\SpecialCharTok{::}\FunctionTok{here}\NormalTok{(}\StringTok{"R"}\NormalTok{, }\StringTok{"sci\_palette.R"}\NormalTok{))}
\end{Highlighting}
\end{Shaded}

\section{Data acquisition}\label{data-acquisition-3}

\begin{Shaded}
\begin{Highlighting}[]
\NormalTok{works }\OtherTok{\textless{}{-}} \FunctionTok{fetch\_openalex}\NormalTok{(}
  \AttributeTok{entity =} \StringTok{"works"}\NormalTok{,}
  \AttributeTok{primary\_location.source.id =} \StringTok{"S148561398"}\NormalTok{,}
  \AttributeTok{from\_publication\_date =} \StringTok{"2015{-}01{-}01"}\NormalTok{,}
  \AttributeTok{to\_publication\_date =} \StringTok{"2023{-}12{-}31"}\NormalTok{,}
  \AttributeTok{type =} \StringTok{"article"}\NormalTok{,}
  \AttributeTok{options =} \FunctionTok{list}\NormalTok{(}\AttributeTok{sample =} \DecValTok{600}\NormalTok{, }\AttributeTok{seed =} \DecValTok{42}\NormalTok{)}
\NormalTok{)}

\NormalTok{text\_df }\OtherTok{\textless{}{-}}\NormalTok{ works }\SpecialCharTok{|\textgreater{}}
  \FunctionTok{filter}\NormalTok{(}\SpecialCharTok{!}\FunctionTok{is.na}\NormalTok{(abstract), }\FunctionTok{nchar}\NormalTok{(abstract) }\SpecialCharTok{\textgreater{}} \DecValTok{50}\NormalTok{) }\SpecialCharTok{|\textgreater{}}
  \FunctionTok{transmute}\NormalTok{(}
    \AttributeTok{doc\_id =}\NormalTok{ id,}
    \AttributeTok{text =} \FunctionTok{paste}\NormalTok{(display\_name, abstract, }\AttributeTok{sep =} \StringTok{". "}\NormalTok{),}
    \AttributeTok{year =} \FunctionTok{year}\NormalTok{(publication\_date)}
\NormalTok{  )}

\FunctionTok{cat}\NormalTok{(}\FunctionTok{glue}\NormalTok{(}\StringTok{"Documents: \{nrow(text\_df)\}}\SpecialCharTok{\textbackslash{}n}\StringTok{"}\NormalTok{))}
\end{Highlighting}
\end{Shaded}

\begin{verbatim}
#> Documents: 131
\end{verbatim}

\section{Keyword burst detection}\label{keyword-burst-detection}

\begin{Shaded}
\begin{Highlighting}[]
\NormalTok{corp }\OtherTok{\textless{}{-}} \FunctionTok{corpus}\NormalTok{(text\_df, }\AttributeTok{docid\_field =} \StringTok{"doc\_id"}\NormalTok{, }\AttributeTok{text\_field =} \StringTok{"text"}\NormalTok{)}
\FunctionTok{docvars}\NormalTok{(corp, }\StringTok{"year"}\NormalTok{) }\OtherTok{\textless{}{-}}\NormalTok{ text\_df}\SpecialCharTok{$}\NormalTok{year}

\NormalTok{toks }\OtherTok{\textless{}{-}} \FunctionTok{tokens}\NormalTok{(corp, }\AttributeTok{remove\_punct =} \ConstantTok{TRUE}\NormalTok{, }\AttributeTok{remove\_numbers =} \ConstantTok{TRUE}\NormalTok{) }\SpecialCharTok{|\textgreater{}}
  \FunctionTok{tokens\_tolower}\NormalTok{() }\SpecialCharTok{|\textgreater{}}
  \FunctionTok{tokens\_remove}\NormalTok{(}\FunctionTok{stopwords}\NormalTok{(}\StringTok{"en"}\NormalTok{)) }\SpecialCharTok{|\textgreater{}}
  \FunctionTok{tokens\_remove}\NormalTok{(}\FunctionTok{c}\NormalTok{(}\StringTok{"study"}\NormalTok{, }\StringTok{"paper"}\NormalTok{, }\StringTok{"results"}\NormalTok{, }\StringTok{"research"}\NormalTok{, }\StringTok{"analysis"}\NormalTok{,}
                   \StringTok{"also"}\NormalTok{, }\StringTok{"however"}\NormalTok{, }\StringTok{"using"}\NormalTok{, }\StringTok{"based"}\NormalTok{))}

\NormalTok{dfmat }\OtherTok{\textless{}{-}} \FunctionTok{dfm}\NormalTok{(toks) }\SpecialCharTok{|\textgreater{}} \FunctionTok{dfm\_trim}\NormalTok{(}\AttributeTok{min\_termfreq =} \DecValTok{10}\NormalTok{, }\AttributeTok{min\_docfreq =} \DecValTok{5}\NormalTok{)}

\NormalTok{kw\_by\_year }\OtherTok{\textless{}{-}}\NormalTok{ quanteda}\SpecialCharTok{::}\FunctionTok{convert}\NormalTok{(}\FunctionTok{dfm\_group}\NormalTok{(dfmat, }\AttributeTok{groups =}\NormalTok{ year), }\AttributeTok{to =} \StringTok{"data.frame"}\NormalTok{) }\SpecialCharTok{|\textgreater{}}
  \FunctionTok{pivot\_longer}\NormalTok{(}\SpecialCharTok{{-}}\NormalTok{doc\_id, }\AttributeTok{names\_to =} \StringTok{"term"}\NormalTok{, }\AttributeTok{values\_to =} \StringTok{"count"}\NormalTok{) }\SpecialCharTok{|\textgreater{}}
  \FunctionTok{rename}\NormalTok{(}\AttributeTok{year =}\NormalTok{ doc\_id) }\SpecialCharTok{|\textgreater{}}
  \FunctionTok{mutate}\NormalTok{(}\AttributeTok{year =} \FunctionTok{as.integer}\NormalTok{(year))}

\NormalTok{growth }\OtherTok{\textless{}{-}}\NormalTok{ kw\_by\_year }\SpecialCharTok{|\textgreater{}}
  \FunctionTok{group\_by}\NormalTok{(term) }\SpecialCharTok{|\textgreater{}}
  \FunctionTok{filter}\NormalTok{(}\FunctionTok{sum}\NormalTok{(count) }\SpecialCharTok{\textgreater{}=} \DecValTok{20}\NormalTok{) }\SpecialCharTok{|\textgreater{}}
  \FunctionTok{arrange}\NormalTok{(year) }\SpecialCharTok{|\textgreater{}}
  \FunctionTok{mutate}\NormalTok{(}\AttributeTok{growth =}\NormalTok{ (count }\SpecialCharTok{{-}} \FunctionTok{lag}\NormalTok{(count)) }\SpecialCharTok{/} \FunctionTok{pmax}\NormalTok{(}\FunctionTok{lag}\NormalTok{(count), }\DecValTok{1}\NormalTok{)) }\SpecialCharTok{|\textgreater{}}
  \FunctionTok{ungroup}\NormalTok{() }\SpecialCharTok{|\textgreater{}}
  \FunctionTok{filter}\NormalTok{(}\SpecialCharTok{!}\FunctionTok{is.na}\NormalTok{(growth))}

\NormalTok{recent\_growth }\OtherTok{\textless{}{-}}\NormalTok{ growth }\SpecialCharTok{|\textgreater{}}
  \FunctionTok{filter}\NormalTok{(year }\SpecialCharTok{\textgreater{}=} \DecValTok{2021}\NormalTok{) }\SpecialCharTok{|\textgreater{}}
  \FunctionTok{group\_by}\NormalTok{(term) }\SpecialCharTok{|\textgreater{}}
  \FunctionTok{summarise}\NormalTok{(}\AttributeTok{mean\_growth =} \FunctionTok{mean}\NormalTok{(growth), }\AttributeTok{total =} \FunctionTok{sum}\NormalTok{(count), }\AttributeTok{.groups =} \StringTok{"drop"}\NormalTok{) }\SpecialCharTok{|\textgreater{}}
  \FunctionTok{filter}\NormalTok{(mean\_growth }\SpecialCharTok{\textgreater{}} \DecValTok{0}\NormalTok{) }\SpecialCharTok{|\textgreater{}}
  \FunctionTok{arrange}\NormalTok{(}\FunctionTok{desc}\NormalTok{(mean\_growth))}
\end{Highlighting}
\end{Shaded}

\begin{Shaded}
\begin{Highlighting}[]
\NormalTok{recent\_growth }\SpecialCharTok{|\textgreater{}}
  \FunctionTok{head}\NormalTok{(}\DecValTok{15}\NormalTok{) }\SpecialCharTok{|\textgreater{}}
  \FunctionTok{mutate}\NormalTok{(}\AttributeTok{term =} \FunctionTok{fct\_reorder}\NormalTok{(term, mean\_growth)) }\SpecialCharTok{|\textgreater{}}
  \FunctionTok{ggplot}\NormalTok{(}\FunctionTok{aes}\NormalTok{(}\AttributeTok{x =}\NormalTok{ mean\_growth, }\AttributeTok{y =}\NormalTok{ term)) }\SpecialCharTok{+}
  \FunctionTok{geom\_col}\NormalTok{(}\AttributeTok{fill =} \FunctionTok{palette\_sci}\NormalTok{(}\DecValTok{1}\NormalTok{)) }\SpecialCharTok{+}
  \FunctionTok{scale\_x\_continuous}\NormalTok{(}\AttributeTok{labels =}\NormalTok{ scales}\SpecialCharTok{::}\NormalTok{percent) }\SpecialCharTok{+}
  \FunctionTok{labs}\NormalTok{(}\AttributeTok{x =} \StringTok{"Mean growth rate (2021–2023)"}\NormalTok{, }\AttributeTok{y =} \ConstantTok{NULL}\NormalTok{) }\SpecialCharTok{+}
  \FunctionTok{theme\_sci}\NormalTok{()}
\end{Highlighting}
\end{Shaded}

\begin{figure}

{\centering \includegraphics[width=0.9\linewidth,alt={Bar chart showing keywords with the highest mean growth rate in recent years.}]{scientometrics-in-r_files/figure-latex/cs05-fig-cs05-growth-1} 

}

\caption{Top 15 fastest-growing keywords (2021–2023).}\label{fig:cs05-fig-cs05-growth}
\end{figure}

\section{Keyword trajectories}\label{keyword-trajectories}

\begin{Shaded}
\begin{Highlighting}[]
\NormalTok{emerging\_terms }\OtherTok{\textless{}{-}}\NormalTok{ recent\_growth }\SpecialCharTok{|\textgreater{}} \FunctionTok{head}\NormalTok{(}\DecValTok{6}\NormalTok{) }\SpecialCharTok{|\textgreater{}} \FunctionTok{pull}\NormalTok{(term)}

\NormalTok{kw\_by\_year }\SpecialCharTok{|\textgreater{}}
  \FunctionTok{filter}\NormalTok{(term }\SpecialCharTok{\%in\%}\NormalTok{ emerging\_terms) }\SpecialCharTok{|\textgreater{}}
  \FunctionTok{ggplot}\NormalTok{(}\FunctionTok{aes}\NormalTok{(}\AttributeTok{x =}\NormalTok{ year, }\AttributeTok{y =}\NormalTok{ count, }\AttributeTok{colour =}\NormalTok{ term)) }\SpecialCharTok{+}
  \FunctionTok{geom\_line}\NormalTok{(}\AttributeTok{linewidth =} \FloatTok{0.8}\NormalTok{) }\SpecialCharTok{+}
  \FunctionTok{geom\_point}\NormalTok{(}\AttributeTok{size =} \FloatTok{1.5}\NormalTok{) }\SpecialCharTok{+}
  \FunctionTok{scale\_colour\_manual}\NormalTok{(}\AttributeTok{values =} \FunctionTok{palette\_sci}\NormalTok{(}\FunctionTok{length}\NormalTok{(emerging\_terms))) }\SpecialCharTok{+}
  \FunctionTok{labs}\NormalTok{(}\AttributeTok{x =} \StringTok{"Year"}\NormalTok{, }\AttributeTok{y =} \StringTok{"Frequency"}\NormalTok{, }\AttributeTok{colour =} \StringTok{"Keyword"}\NormalTok{) }\SpecialCharTok{+}
  \FunctionTok{theme\_sci}\NormalTok{()}
\end{Highlighting}
\end{Shaded}

\begin{figure}

{\centering \includegraphics[width=0.9\linewidth,alt={Line chart showing frequency over time for keywords identified as growing rapidly.}]{scientometrics-in-r_files/figure-latex/cs05-fig-cs05-trajectories-1} 

}

\caption{Temporal trajectories of selected emerging keywords.}\label{fig:cs05-fig-cs05-trajectories}
\end{figure}

\section{Topic modelling}\label{topic-modelling}

\begin{Shaded}
\begin{Highlighting}[]
\NormalTok{dtm }\OtherTok{\textless{}{-}}\NormalTok{ quanteda}\SpecialCharTok{::}\FunctionTok{convert}\NormalTok{(dfmat, }\AttributeTok{to =} \StringTok{"topicmodels"}\NormalTok{)}
\NormalTok{lda }\OtherTok{\textless{}{-}} \FunctionTok{LDA}\NormalTok{(dtm, }\AttributeTok{k =} \DecValTok{8}\NormalTok{, }\AttributeTok{control =} \FunctionTok{list}\NormalTok{(}\AttributeTok{seed =} \DecValTok{42}\NormalTok{))}

\NormalTok{lda\_topics }\OtherTok{\textless{}{-}} \FunctionTok{tidy}\NormalTok{(lda, }\AttributeTok{matrix =} \StringTok{"beta"}\NormalTok{) }\SpecialCharTok{|\textgreater{}}
  \FunctionTok{group\_by}\NormalTok{(topic) }\SpecialCharTok{|\textgreater{}}
  \FunctionTok{slice\_max}\NormalTok{(beta, }\AttributeTok{n =} \DecValTok{8}\NormalTok{) }\SpecialCharTok{|\textgreater{}}
  \FunctionTok{ungroup}\NormalTok{()}
\end{Highlighting}
\end{Shaded}

\begin{Shaded}
\begin{Highlighting}[]
\NormalTok{lda\_topics }\SpecialCharTok{|\textgreater{}}
  \FunctionTok{mutate}\NormalTok{(}\AttributeTok{term =} \FunctionTok{reorder\_within}\NormalTok{(term, beta, topic)) }\SpecialCharTok{|\textgreater{}}
  \FunctionTok{ggplot}\NormalTok{(}\FunctionTok{aes}\NormalTok{(}\AttributeTok{x =}\NormalTok{ beta, }\AttributeTok{y =}\NormalTok{ term, }\AttributeTok{fill =} \FunctionTok{factor}\NormalTok{(topic))) }\SpecialCharTok{+}
  \FunctionTok{geom\_col}\NormalTok{(}\AttributeTok{show.legend =} \ConstantTok{FALSE}\NormalTok{) }\SpecialCharTok{+}
  \FunctionTok{facet\_wrap}\NormalTok{(}\SpecialCharTok{\textasciitilde{}} \FunctionTok{paste}\NormalTok{(}\StringTok{"Topic"}\NormalTok{, topic), }\AttributeTok{scales =} \StringTok{"free\_y"}\NormalTok{, }\AttributeTok{ncol =} \DecValTok{4}\NormalTok{) }\SpecialCharTok{+}
  \FunctionTok{scale\_y\_reordered}\NormalTok{() }\SpecialCharTok{+}
  \FunctionTok{scale\_fill\_manual}\NormalTok{(}\AttributeTok{values =} \FunctionTok{palette\_sci}\NormalTok{(}\DecValTok{8}\NormalTok{)) }\SpecialCharTok{+}
  \FunctionTok{labs}\NormalTok{(}\AttributeTok{x =} \StringTok{"Word probability"}\NormalTok{, }\AttributeTok{y =} \ConstantTok{NULL}\NormalTok{) }\SpecialCharTok{+}
  \FunctionTok{theme\_sci}\NormalTok{(}\AttributeTok{base\_size =} \DecValTok{9}\NormalTok{)}
\end{Highlighting}
\end{Shaded}

\begin{figure}

{\centering \includegraphics[width=0.9\linewidth,alt={Faceted bar chart showing the 8 most probable words in each of 8 LDA topics.}]{scientometrics-in-r_files/figure-latex/cs05-fig-cs05-lda-1} 

}

\caption{Top terms per LDA topic.}\label{fig:cs05-fig-cs05-lda}
\end{figure}

\section{Topic prevalence by year}\label{topic-prevalence-by-year}

\begin{Shaded}
\begin{Highlighting}[]
\NormalTok{doc\_topics }\OtherTok{\textless{}{-}} \FunctionTok{tidy}\NormalTok{(lda, }\AttributeTok{matrix =} \StringTok{"gamma"}\NormalTok{) }\SpecialCharTok{|\textgreater{}}
  \FunctionTok{left\_join}\NormalTok{(text\_df }\SpecialCharTok{|\textgreater{}} \FunctionTok{transmute}\NormalTok{(}\AttributeTok{document =}\NormalTok{ doc\_id, year), }\AttributeTok{by =} \StringTok{"document"}\NormalTok{)}

\NormalTok{topic\_by\_year }\OtherTok{\textless{}{-}}\NormalTok{ doc\_topics }\SpecialCharTok{|\textgreater{}}
  \FunctionTok{group\_by}\NormalTok{(year, topic) }\SpecialCharTok{|\textgreater{}}
  \FunctionTok{summarise}\NormalTok{(}\AttributeTok{mean\_gamma =} \FunctionTok{mean}\NormalTok{(gamma), }\AttributeTok{.groups =} \StringTok{"drop"}\NormalTok{)}

\FunctionTok{ggplot}\NormalTok{(topic\_by\_year, }\FunctionTok{aes}\NormalTok{(}\AttributeTok{x =}\NormalTok{ year, }\AttributeTok{y =}\NormalTok{ mean\_gamma, }\AttributeTok{colour =} \FunctionTok{factor}\NormalTok{(topic))) }\SpecialCharTok{+}
  \FunctionTok{geom\_line}\NormalTok{(}\AttributeTok{linewidth =} \FloatTok{0.7}\NormalTok{) }\SpecialCharTok{+}
  \FunctionTok{scale\_colour\_manual}\NormalTok{(}\AttributeTok{values =} \FunctionTok{palette\_sci}\NormalTok{(}\DecValTok{8}\NormalTok{)) }\SpecialCharTok{+}
  \FunctionTok{labs}\NormalTok{(}\AttributeTok{x =} \StringTok{"Year"}\NormalTok{, }\AttributeTok{y =} \StringTok{"Mean topic proportion"}\NormalTok{, }\AttributeTok{colour =} \StringTok{"Topic"}\NormalTok{) }\SpecialCharTok{+}
  \FunctionTok{theme\_sci}\NormalTok{()}
\end{Highlighting}
\end{Shaded}

\begin{figure}

{\centering \includegraphics[width=0.9\linewidth,alt={Line chart showing how the proportion of each topic changes across publication years.}]{scientometrics-in-r_files/figure-latex/cs05-fig-cs05-prevalence-1} 

}

\caption{Topic prevalence over time.}\label{fig:cs05-fig-cs05-prevalence}
\end{figure}

\section{Key findings}\label{key-findings-1}

\begin{enumerate}
\def\labelenumi{\arabic{enumi}.}
\tightlist
\item
  \textbf{Emerging keywords} align with known developments: open science, AI/machine learning applications, equity and diversity, and preprints.
\item
  \textbf{Topic evolution} shows gradual shifts in emphasis rather than abrupt changes, consistent with a mature field incorporating new tools.
\item
  \textbf{Growth rates} should be interpreted cautiously: terms with low baseline frequencies can show high percentage growth from small increases.
\item
  \textbf{Validation}: the topics identified by LDA correspond to recognisable subfields, lending credibility to the unsupervised approach.
\end{enumerate}

\section{Lessons learned}\label{lessons-learned-3}

\begin{itemize}
\tightlist
\item
  Combining keyword growth rates with topic models provides complementary perspectives: growth rates capture individual terms, while topics capture thematic clusters.
\item
  Database coverage growth can create artefactual ``emergence.'' Always normalise by total corpus size.
\item
  Emerging topics in scientometrics mirror broader trends in science policy (open access mandates, equity initiatives, AI adoption) \citep{kleinberg2003}.
\end{itemize}

\appendix


\chapter{R Setup, renv, and Docker}\label{r-setup-renv-and-docker}

\emph{R. Heller · Date: 2026-05-09 · Last changed: 2026-10-31}

\section{R installation and current first-use path}\label{r-installation-and-current-first-use-path}

Install R from \href{https://cran.r-project.org/}{CRAN}. The current candidate was
locally validated with R 4.6.1. An IDE such as
\href{https://posit.co/download/rstudio-desktop/}{RStudio} is optional. Rendering also
requires Pandoc; PDF output additionally needs XeLaTeX/TinyTeX and R with Cairo support. Compiled packages
may require a C/C++ toolchain and system libraries; the Linux dependencies are
listed in \texttt{.github/workflows/render.yml}.

From the repository root, install and check the book's R dependencies:

\begin{Shaded}
\begin{Highlighting}[]
\VariableTok{RENV\_CONFIG\_AUTOLOADER\_ENABLED}\OperatorTok{=}\NormalTok{false }\ExtensionTok{Rscript} \AttributeTok{{-}{-}vanilla}\NormalTok{ scripts/install{-}book{-}deps.R}
\VariableTok{RENV\_CONFIG\_AUTOLOADER\_ENABLED}\OperatorTok{=}\NormalTok{false }\ExtensionTok{Rscript} \AttributeTok{{-}{-}vanilla}\NormalTok{ scripts/install{-}book{-}deps.R }\AttributeTok{{-}{-}check}
\end{Highlighting}
\end{Shaded}

The installer installs missing packages from current CRAN and checks that their
namespaces load. It does not restore historical versions or certify an immutable
software environment. \texttt{-\/-check} performs no installation or network access;
\texttt{-\/-list} prints the package inventory. Use \texttt{-\/-library=/path/to/library} to select
an R library, or \texttt{-\/-include-optional} for packages used only by unevaluated
examples and exercises. Optional Python, SQL, OCR and deployment prerequisites
must be configured separately.

Internet access is required for installation and uncached provider requests.
Configure an activated OpenAlex key locally as \texttt{OPENALEX\_API\_KEY}, outside Git,
for requests that need authenticated paging. \texttt{SCIENTOMETRICS\_OFFLINE=1} requires
existing exact-query OpenAlex caches in \texttt{\_freeze/openalex\_cache/}; it does not
make other providers offline and cannot supply missing cache entries. See the
repository \texttt{README.md} for the companion package's separate installation path.

\section{Package management with renv}\label{package-management-with-renv}

\textbf{renv} can record and restore a project's R package versions. This repository
retains a historical \texttt{renv.lock} and automatic activation in \texttt{.Rprofile}, but that
lockfile has \textbf{not been qualified for the current book}. The first-use commands
above bypass project startup and use the dependency installer. Do not assume
that restoring the historical lockfile will reproduce the current validated
candidate.

The following is an optional restoration example for a project whose lockfile
you have independently validated, rather than a required setup step for this book:

\begin{Shaded}
\begin{Highlighting}[]
\FunctionTok{install.packages}\NormalTok{(}\StringTok{"renv"}\NormalTok{)}
\NormalTok{renv}\SpecialCharTok{::}\FunctionTok{restore}\NormalTok{()}
\end{Highlighting}
\end{Shaded}

Restoration attempts to install the recorded versions; success also depends on
compatible R, system libraries and available package sources. In a development
environment, add a package and deliberately update the lockfile only after
validating the resulting environment:

\begin{Shaded}
\begin{Highlighting}[]
\NormalTok{renv}\SpecialCharTok{::}\FunctionTok{install}\NormalTok{(}\StringTok{"newpackage"}\NormalTok{) }\CommentTok{\# Replace with the intended package name.}
\NormalTok{renv}\SpecialCharTok{::}\FunctionTok{snapshot}\NormalTok{()}
\end{Highlighting}
\end{Shaded}

A software lockfile does not freeze live API results. Exact reproduction also
requires the same licensed input snapshots and analysis configuration.

\section{Core packages}\label{core-packages}

The table below lists the primary packages used throughout the book.

{\def\LTcaptype{none} % do not increment counter
\begin{longtable}[]{@{}lll@{}}
\toprule\noalign{}
Package & Purpose & Chapter(s) \\
\midrule\noalign{}
\endhead
\bottomrule\noalign{}
\endlastfoot
openalexR & OpenAlex API access & 5--37 \\
tidyverse & Data wrangling and visualisation & All \\
igraph & Network analysis & 15--20 \\
tidygraph & Tidy network manipulation & 15--20 \\
ggraph & Network visualisation & 15--20 \\
quanteda & Text mining and corpus management & 21--25 \\
topicmodels & LDA topic modelling & 22, CS-5 \\
stm & Structural topic models & 22 \\
gt & Publication-quality tables & Throughout \\
targets & Pipeline orchestration & 34 \\
plotly & Interactive charts & 35 \\
DT & Interactive data tables & 35 \\
shiny & Web applications & 37 \\
word2vec & Word embeddings & 24 \\
uwot & UMAP dimensionality reduction & 24 \\
\end{longtable}
}

\section{Docker environment}\label{docker-environment}

The repository's \texttt{Dockerfile} is an \textbf{optional historical recipe}, based on
\texttt{rocker/verse:4.4.1} and \texttt{renv::restore()}. Its restore path has not been qualified
for the current book. It also downloads an amd64 Quarto package, so other
architectures need a compatible recipe. Docker does not remove the need to
validate the lockfile, system dependencies, API access and resulting outputs.

For a deliberate evaluation of that recipe, build it and inspect its R version:

\begin{Shaded}
\begin{Highlighting}[]
\ExtensionTok{docker}\NormalTok{ build }\AttributeTok{{-}t}\NormalTok{ scientometrics{-}in{-}r .}
\ExtensionTok{docker}\NormalTok{ run }\AttributeTok{{-}{-}rm}\NormalTok{ scientometrics{-}in{-}r R }\AttributeTok{{-}{-}version}
\end{Highlighting}
\end{Shaded}

The image's default command calls \texttt{bookdown::render\_book()} for all formats; it
does not launch RStudio Server. Successful construction and execution of this
historical image are not claimed here. Use the current first-use path above
for the locally validated candidate.

\section{Render the book}\label{render-the-book}

This repository is configured for \textbf{bookdown}, with \texttt{\_bookdown.yml} and
\texttt{\_output.yml}. From the repository root, render HTML and refresh its download
manifest:

\begin{Shaded}
\begin{Highlighting}[]
\VariableTok{SCIENTOMETRICS\_REBUILD\_CHUNKS}\OperatorTok{=}\NormalTok{1 }\VariableTok{RENV\_CONFIG\_AUTOLOADER\_ENABLED}\OperatorTok{=}\NormalTok{false }\ExtensionTok{Rscript} \AttributeTok{{-}{-}vanilla} \AttributeTok{{-}e} \StringTok{\textquotesingle{}bookdown::render\_book("index.Rmd", output\_format = "bookdown::bs4\_book")\textquotesingle{}}
\ExtensionTok{Rscript} \AttributeTok{{-}{-}vanilla}\NormalTok{ scripts/build{-}download{-}manifest.R}
\end{Highlighting}
\end{Shaded}

\texttt{SCIENTOMETRICS\_REBUILD\_CHUNKS=1} re-executes analysis chunks while retaining
exact-query provider caches. Missing provider caches still require API access.
The configured output directory is \texttt{docs/}. Serve it locally, for example:

\begin{Shaded}
\begin{Highlighting}[]
\ExtensionTok{python3} \AttributeTok{{-}m}\NormalTok{ http.server 8000 }\AttributeTok{{-}{-}directory}\NormalTok{ docs}
\end{Highlighting}
\end{Shaded}

Then visit \texttt{http://localhost:8000}. PDF additionally needs XeLaTeX/TinyTeX;
\texttt{output\_format\ =\ "all"} requests HTML, PDF and EPUB. Before distribution, run
\texttt{Rscript\ -\/-vanilla\ scripts/repair-epub.R\ docs/scientometrics-in-r.epub} to validate
the ebook's internal links and bundle its citation and licence files, then run
\texttt{Rscript\ -\/-vanilla\ scripts/build-download-manifest.R}. Local outputs do not
update the published site. The README lists the PDF font prerequisites.

Quarto appears in optional examples elsewhere in the book, but \texttt{quarto\ render}
is not the build command for this repository. A Quarto CLI is needed only when
working through those separate examples.

\chapter{SQL Primer for the OpenAlex Snapshot}\label{sql-primer-for-the-openalex-snapshot}

\emph{R. Heller · Date: 2026-05-09 · Last changed: 2026-05-11}

\section{When to use the snapshot}\label{when-to-use-the-snapshot}

The OpenAlex API (Chapter 5) is ideal for moderate queries. For analyses requiring millions of records --- full-database bibliographic coupling, global co-authorship networks, or exhaustive field delineation --- the \href{https://docs.openalex.org/download-all-data/openalex-snapshot}{OpenAlex snapshot} is more practical. The snapshot is a full database dump released monthly in compressed JSON-lines format.

\section{Loading the snapshot}\label{loading-the-snapshot}

OpenAlex provides scripts to load the snapshot into PostgreSQL. After loading:

\begin{Shaded}
\begin{Highlighting}[]
\CommentTok{{-}{-} Count total works}
\KeywordTok{SELECT} \FunctionTok{COUNT}\NormalTok{(}\OperatorTok{*}\NormalTok{) }\KeywordTok{FROM}\NormalTok{ openalex.works;}

\CommentTok{{-}{-} Recent articles with citation counts}
\KeywordTok{SELECT} \KeywordTok{id}\NormalTok{, display\_name, publication\_year, cited\_by\_count}
\KeywordTok{FROM}\NormalTok{ openalex.works}
\KeywordTok{WHERE} \KeywordTok{type} \OperatorTok{=} \StringTok{\textquotesingle{}article\textquotesingle{}}
  \KeywordTok{AND}\NormalTok{ publication\_year }\OperatorTok{\textgreater{}=} \DecValTok{2020}
\KeywordTok{ORDER} \KeywordTok{BY}\NormalTok{ cited\_by\_count }\KeywordTok{DESC}
\KeywordTok{LIMIT} \DecValTok{100}\NormalTok{;}
\end{Highlighting}
\end{Shaded}

\section{Common queries}\label{common-queries}

\subsection{Publication counts by year}\label{publication-counts-by-year}

\begin{Shaded}
\begin{Highlighting}[]
\KeywordTok{SELECT}\NormalTok{ publication\_year, }\FunctionTok{COUNT}\NormalTok{(}\OperatorTok{*}\NormalTok{) }\KeywordTok{AS}\NormalTok{ n\_works}
\KeywordTok{FROM}\NormalTok{ openalex.works}
\KeywordTok{WHERE} \KeywordTok{type} \OperatorTok{=} \StringTok{\textquotesingle{}article\textquotesingle{}}
\KeywordTok{GROUP} \KeywordTok{BY}\NormalTok{ publication\_year}
\KeywordTok{ORDER} \KeywordTok{BY}\NormalTok{ publication\_year;}
\end{Highlighting}
\end{Shaded}

\subsection{Top institutions by output}\label{top-institutions-by-output}

\begin{Shaded}
\begin{Highlighting}[]
\KeywordTok{SELECT}\NormalTok{ i.display\_name, }\FunctionTok{COUNT}\NormalTok{(}\KeywordTok{DISTINCT}\NormalTok{ w.}\KeywordTok{id}\NormalTok{) }\KeywordTok{AS}\NormalTok{ n\_works}
\KeywordTok{FROM}\NormalTok{ openalex.works w}
\KeywordTok{JOIN}\NormalTok{ openalex.works\_authorships wa }\KeywordTok{ON}\NormalTok{ w.}\KeywordTok{id} \OperatorTok{=}\NormalTok{ wa.work\_id}
\KeywordTok{JOIN}\NormalTok{ openalex.works\_authorships\_institutions wai}
  \KeywordTok{ON}\NormalTok{ wa.work\_id }\OperatorTok{=}\NormalTok{ wai.work\_id }\KeywordTok{AND}\NormalTok{ wa.author\_position }\OperatorTok{=}\NormalTok{ wai.author\_position}
\KeywordTok{JOIN}\NormalTok{ openalex.institutions i }\KeywordTok{ON}\NormalTok{ wai.institution\_id }\OperatorTok{=}\NormalTok{ i.}\KeywordTok{id}
\KeywordTok{WHERE}\NormalTok{ w.publication\_year }\OperatorTok{=} \DecValTok{2023}
  \KeywordTok{AND}\NormalTok{ w.}\KeywordTok{type} \OperatorTok{=} \StringTok{\textquotesingle{}article\textquotesingle{}}
\KeywordTok{GROUP} \KeywordTok{BY}\NormalTok{ i.display\_name}
\KeywordTok{ORDER} \KeywordTok{BY}\NormalTok{ n\_works }\KeywordTok{DESC}
\KeywordTok{LIMIT} \DecValTok{20}\NormalTok{;}
\end{Highlighting}
\end{Shaded}

\subsection{Co-authorship pairs}\label{co-authorship-pairs}

\begin{Shaded}
\begin{Highlighting}[]
\KeywordTok{SELECT}\NormalTok{ a1.author\_id }\KeywordTok{AS}\NormalTok{ author\_a,}
\NormalTok{       a2.author\_id }\KeywordTok{AS}\NormalTok{ author\_b,}
       \FunctionTok{COUNT}\NormalTok{(}\KeywordTok{DISTINCT}\NormalTok{ a1.work\_id) }\KeywordTok{AS}\NormalTok{ n\_collabs}
\KeywordTok{FROM}\NormalTok{ openalex.works\_authorships a1}
\KeywordTok{JOIN}\NormalTok{ openalex.works\_authorships a2}
  \KeywordTok{ON}\NormalTok{ a1.work\_id }\OperatorTok{=}\NormalTok{ a2.work\_id }\KeywordTok{AND}\NormalTok{ a1.author\_id }\OperatorTok{\textless{}}\NormalTok{ a2.author\_id}
\KeywordTok{GROUP} \KeywordTok{BY}\NormalTok{ a1.author\_id, a2.author\_id}
\KeywordTok{HAVING} \FunctionTok{COUNT}\NormalTok{(}\KeywordTok{DISTINCT}\NormalTok{ a1.work\_id) }\OperatorTok{\textgreater{}=} \DecValTok{5}
\KeywordTok{ORDER} \KeywordTok{BY}\NormalTok{ n\_collabs }\KeywordTok{DESC}
\KeywordTok{LIMIT} \DecValTok{100}\NormalTok{;}
\end{Highlighting}
\end{Shaded}

\section{Connecting from R}\label{connecting-from-r}

\begin{Shaded}
\begin{Highlighting}[]
\FunctionTok{library}\NormalTok{(DBI)}
\FunctionTok{library}\NormalTok{(RPostgres)}

\NormalTok{con }\OtherTok{\textless{}{-}} \FunctionTok{dbConnect}\NormalTok{(}
  \FunctionTok{Postgres}\NormalTok{(),}
  \AttributeTok{dbname =} \StringTok{"openalex"}\NormalTok{,}
  \AttributeTok{host =} \StringTok{"localhost"}\NormalTok{,}
  \AttributeTok{port =} \DecValTok{5432}\NormalTok{,}
  \AttributeTok{user =} \StringTok{"openalex"}
\NormalTok{)}

\NormalTok{top\_journals }\OtherTok{\textless{}{-}} \FunctionTok{dbGetQuery}\NormalTok{(con, }\StringTok{"}
\StringTok{  SELECT s.display\_name, COUNT(*) AS n\_works}
\StringTok{  FROM openalex.works w}
\StringTok{  JOIN openalex.sources s ON w.primary\_location\_source\_id = s.id}
\StringTok{  WHERE w.publication\_year = 2023}
\StringTok{  GROUP BY s.display\_name}
\StringTok{  ORDER BY n\_works DESC}
\StringTok{  LIMIT 20}
\StringTok{"}\NormalTok{)}

\FunctionTok{dbDisconnect}\NormalTok{(con)}
\end{Highlighting}
\end{Shaded}

\section{Performance tips}\label{performance-tips}

\begin{itemize}
\tightlist
\item
  \textbf{Index key columns}: \texttt{work\_id}, \texttt{author\_id}, \texttt{institution\_id}, \texttt{publication\_year}, and \texttt{cited\_by\_count} should all be indexed.
\item
  \textbf{Partition by year}: For time-series queries, partitioning the works table by \texttt{publication\_year} speeds up range scans substantially.
\item
  \textbf{Materialised views}: Pre-compute expensive joins (e.g., institution--work counts) as materialised views and refresh them after each snapshot update.
\item
  \textbf{Use COPY for bulk loads}: The \texttt{COPY} command is orders of magnitude faster than row-by-row \texttt{INSERT}.
\end{itemize}

\chapter{Indicator Crosswalk Across Databases}\label{indicator-crosswalk-across-databases}

\emph{R. Heller · Date: 2026-05-09 · Last changed: 2026-05-11}

This appendix maps common bibliometric indicators to the data sources and tools that support them. Use it to choose the right source for your analysis.

\section{Publication-level indicators}\label{publication-level-indicators}

{\def\LTcaptype{none} % do not increment counter
\begin{longtable}[]{@{}
  >{\raggedright\arraybackslash}p{(\linewidth - 10\tabcolsep) * \real{0.1667}}
  >{\centering\arraybackslash}p{(\linewidth - 10\tabcolsep) * \real{0.1515}}
  >{\centering\arraybackslash}p{(\linewidth - 10\tabcolsep) * \real{0.2424}}
  >{\centering\arraybackslash}p{(\linewidth - 10\tabcolsep) * \real{0.1212}}
  >{\centering\arraybackslash}p{(\linewidth - 10\tabcolsep) * \real{0.1818}}
  >{\raggedright\arraybackslash}p{(\linewidth - 10\tabcolsep) * \real{0.1364}}@{}}
\toprule\noalign{}
\begin{minipage}[b]{\linewidth}\raggedright
Indicator
\end{minipage} & \begin{minipage}[b]{\linewidth}\centering
OpenAlex
\end{minipage} & \begin{minipage}[b]{\linewidth}\centering
Web of Science
\end{minipage} & \begin{minipage}[b]{\linewidth}\centering
Scopus
\end{minipage} & \begin{minipage}[b]{\linewidth}\centering
Dimensions
\end{minipage} & \begin{minipage}[b]{\linewidth}\raggedright
Chapter
\end{minipage} \\
\midrule\noalign{}
\endhead
\bottomrule\noalign{}
\endlastfoot
Citation count & Y & Y & Y & Y & 9--10 \\
Field-normalised citations (MNCS) & Compute & CNCI (built-in) & FWCI (built-in) & FCR (built-in) & 13 \\
PP(top 10\%) & Compute & Compute & Compute & Compute & 10 \\
Open access status & Y & Y & Y & Y & 27 \\
Funding acknowledgements & Y (partial) & Y & Y & Y & 30 \\
Abstract text & Y & Y & Y & Y & 21--25 \\
References list & Y & Y & Y & Y & 16 \\
Author affiliations & Y & Y & Y & Y & 12--13 \\
\end{longtable}
}

\section{Author-level indicators}\label{author-level-indicators}

{\def\LTcaptype{none} % do not increment counter
\begin{longtable}[]{@{}
  >{\raggedright\arraybackslash}p{(\linewidth - 10\tabcolsep) * \real{0.1667}}
  >{\centering\arraybackslash}p{(\linewidth - 10\tabcolsep) * \real{0.1515}}
  >{\centering\arraybackslash}p{(\linewidth - 10\tabcolsep) * \real{0.2424}}
  >{\centering\arraybackslash}p{(\linewidth - 10\tabcolsep) * \real{0.1212}}
  >{\centering\arraybackslash}p{(\linewidth - 10\tabcolsep) * \real{0.1818}}
  >{\raggedright\arraybackslash}p{(\linewidth - 10\tabcolsep) * \real{0.1364}}@{}}
\toprule\noalign{}
\begin{minipage}[b]{\linewidth}\raggedright
Indicator
\end{minipage} & \begin{minipage}[b]{\linewidth}\centering
OpenAlex
\end{minipage} & \begin{minipage}[b]{\linewidth}\centering
Web of Science
\end{minipage} & \begin{minipage}[b]{\linewidth}\centering
Scopus
\end{minipage} & \begin{minipage}[b]{\linewidth}\centering
Dimensions
\end{minipage} & \begin{minipage}[b]{\linewidth}\raggedright
Chapter
\end{minipage} \\
\midrule\noalign{}
\endhead
\bottomrule\noalign{}
\endlastfoot
h-index & Compute & Y & Y & Compute & 10, 12 \\
Publication count & Y & Y & Y & Y & 12 \\
ORCID integration & Y & Y & Y & Y & 12 \\
Co-author network & Compute & Compute & Compute & Compute & 15 \\
\end{longtable}
}

\section{Journal-level indicators}\label{journal-level-indicators}

{\def\LTcaptype{none} % do not increment counter
\begin{longtable}[]{@{}lcccl@{}}
\toprule\noalign{}
Indicator & OpenAlex & Web of Science & Scopus & Chapter \\
\midrule\noalign{}
\endhead
\bottomrule\noalign{}
\endlastfoot
Journal Impact Factor & --- & Y & --- & 11 \\
CiteScore & --- & --- & Y & 11 \\
SNIP & --- & --- & Y & 11 \\
SJR & --- & --- & Y & 11 \\
Article volume & Y & Y & Y & 11 \\
Citation distribution & Compute & Compute & Compute & 11 \\
\end{longtable}
}

\section{Institution-level indicators}\label{institution-level-indicators}

{\def\LTcaptype{none} % do not increment counter
\begin{longtable}[]{@{}
  >{\raggedright\arraybackslash}p{(\linewidth - 10\tabcolsep) * \real{0.1594}}
  >{\centering\arraybackslash}p{(\linewidth - 10\tabcolsep) * \real{0.1449}}
  >{\centering\arraybackslash}p{(\linewidth - 10\tabcolsep) * \real{0.2174}}
  >{\centering\arraybackslash}p{(\linewidth - 10\tabcolsep) * \real{0.2319}}
  >{\centering\arraybackslash}p{(\linewidth - 10\tabcolsep) * \real{0.1159}}
  >{\raggedright\arraybackslash}p{(\linewidth - 10\tabcolsep) * \real{0.1304}}@{}}
\toprule\noalign{}
\begin{minipage}[b]{\linewidth}\raggedright
Indicator
\end{minipage} & \begin{minipage}[b]{\linewidth}\centering
OpenAlex
\end{minipage} & \begin{minipage}[b]{\linewidth}\centering
Leiden Ranking
\end{minipage} & \begin{minipage}[b]{\linewidth}\centering
Web of Science
\end{minipage} & \begin{minipage}[b]{\linewidth}\centering
Scopus
\end{minipage} & \begin{minipage}[b]{\linewidth}\raggedright
Chapter
\end{minipage} \\
\midrule\noalign{}
\endhead
\bottomrule\noalign{}
\endlastfoot
Publication counts (full/fractional) & Compute & Y & Compute & Compute & 13 \\
MNCS / FWCI & Compute & Y & CNCI & FWCI & 13 \\
PP(top 10\%) & Compute & Y & Compute & Compute & 13 \\
Collaboration rates & Compute & Y & Compute & Compute & 15 \\
Subject profile & Y & Y & Y & Y & 19 \\
\end{longtable}
}

\section{Network indicators}\label{network-indicators}

{\def\LTcaptype{none} % do not increment counter
\begin{longtable}[]{@{}
  >{\raggedright\arraybackslash}p{(\linewidth - 6\tabcolsep) * \real{0.2821}}
  >{\raggedright\arraybackslash}p{(\linewidth - 6\tabcolsep) * \real{0.3077}}
  >{\raggedright\arraybackslash}p{(\linewidth - 6\tabcolsep) * \real{0.1795}}
  >{\raggedright\arraybackslash}p{(\linewidth - 6\tabcolsep) * \real{0.2308}}@{}}
\toprule\noalign{}
\begin{minipage}[b]{\linewidth}\raggedright
Indicator
\end{minipage} & \begin{minipage}[b]{\linewidth}\raggedright
Data needed
\end{minipage} & \begin{minipage}[b]{\linewidth}\raggedright
Source
\end{minipage} & \begin{minipage}[b]{\linewidth}\raggedright
Chapter
\end{minipage} \\
\midrule\noalign{}
\endhead
\bottomrule\noalign{}
\endlastfoot
Co-authorship strength & Author pairs per paper & Any & 15 \\
Bibliographic coupling & Reference lists & OpenAlex, WoS, Scopus & 16 \\
Co-citation & Citing documents & OpenAlex, WoS, Scopus & 16 \\
Keyword co-occurrence & Keywords/topics & Any & 17 \\
Betweenness centrality & Network graph & Computed & 15 \\
Community structure & Network graph & Computed & 18 \\
\end{longtable}
}

\section{Notes}\label{notes}

\begin{itemize}
\tightlist
\item
  \textbf{``Y''} = available directly from the data source.
\item
  \textbf{``Compute''} = the raw data is available but the indicator must be calculated.
\item
  \textbf{``---''} = not available from that source.
\item
  OpenAlex is the only fully free and open source in this table. Web of Science and Scopus require institutional subscriptions. Dimensions offers a free tier with limited functionality.
\item
  Field normalisation always requires reference values. When a source does not provide them natively, you can compute them from the full corpus using the method in Chapter 13.
\end{itemize}

\chapter{Cheatsheets and Quick Reference}\label{cheatsheets-and-quick-reference}

\emph{R. Heller · Date: 2026-05-09 · Last changed: 2026-10-30}

\section{openalexR quick reference}\label{openalexr-quick-reference}

\begin{Shaded}
\begin{Highlighting}[]
\FunctionTok{library}\NormalTok{(openalexR)}

\CommentTok{\# Fetch works by journal (source)}
\NormalTok{works }\OtherTok{\textless{}{-}} \FunctionTok{fetch\_openalex}\NormalTok{(}
  \AttributeTok{entity =} \StringTok{"works"}\NormalTok{,}
  \AttributeTok{primary\_location.source.id =} \StringTok{"S148561398"}\NormalTok{,}
  \AttributeTok{from\_publication\_date =} \StringTok{"2020{-}01{-}01"}\NormalTok{,}
  \AttributeTok{to\_publication\_date =} \StringTok{"2023{-}12{-}31"}\NormalTok{,}
  \AttributeTok{type =} \StringTok{"article"}
\NormalTok{)}

\CommentTok{\# Fetch works by institution}
\NormalTok{works }\OtherTok{\textless{}{-}} \FunctionTok{fetch\_openalex}\NormalTok{(}
  \AttributeTok{entity =} \StringTok{"works"}\NormalTok{,}
  \AttributeTok{authorships.institutions.id =} \StringTok{"I63966007"}\NormalTok{,}
  \AttributeTok{from\_publication\_date =} \StringTok{"2022{-}01{-}01"}\NormalTok{,}
  \AttributeTok{type =} \StringTok{"article"}
\NormalTok{)}

\CommentTok{\# Fetch works by topic}
\NormalTok{works }\OtherTok{\textless{}{-}} \FunctionTok{fetch\_openalex}\NormalTok{(}
  \AttributeTok{entity =} \StringTok{"works"}\NormalTok{,}
  \AttributeTok{topics.id =} \StringTok{"T10102"}\NormalTok{,}
  \AttributeTok{from\_publication\_date =} \StringTok{"2023{-}01{-}01"}
\NormalTok{)}

\CommentTok{\# Fetch author profile}
\NormalTok{author }\OtherTok{\textless{}{-}} \FunctionTok{fetch\_openalex}\NormalTok{(}\AttributeTok{entity =} \StringTok{"authors"}\NormalTok{, }\AttributeTok{id =} \StringTok{"A5023888391"}\NormalTok{)}

\CommentTok{\# Fetch institution profile}
\NormalTok{inst }\OtherTok{\textless{}{-}} \FunctionTok{fetch\_openalex}\NormalTok{(}\AttributeTok{entity =} \StringTok{"institutions"}\NormalTok{, }\AttributeTok{id =} \StringTok{"I63966007"}\NormalTok{)}

\CommentTok{\# Sampling}
\NormalTok{works }\OtherTok{\textless{}{-}} \FunctionTok{fetch\_openalex}\NormalTok{(}
  \AttributeTok{entity =} \StringTok{"works"}\NormalTok{,}
  \AttributeTok{primary\_location.source.id =} \StringTok{"S148561398"}\NormalTok{,}
  \AttributeTok{options =} \FunctionTok{list}\NormalTok{(}\AttributeTok{sample =} \DecValTok{500}\NormalTok{, }\AttributeTok{seed =} \DecValTok{42}\NormalTok{)}
\NormalTok{)}
\end{Highlighting}
\end{Shaded}

\section{quanteda pipeline}\label{quanteda-pipeline}

\begin{Shaded}
\begin{Highlighting}[]
\FunctionTok{library}\NormalTok{(quanteda)}

\CommentTok{\# Full text{-}analysis pipeline}
\NormalTok{corp }\OtherTok{\textless{}{-}} \FunctionTok{corpus}\NormalTok{(df, }\AttributeTok{docid\_field =} \StringTok{"id"}\NormalTok{, }\AttributeTok{text\_field =} \StringTok{"text"}\NormalTok{)}
\NormalTok{toks }\OtherTok{\textless{}{-}} \FunctionTok{tokens}\NormalTok{(corp, }\AttributeTok{remove\_punct =} \ConstantTok{TRUE}\NormalTok{, }\AttributeTok{remove\_numbers =} \ConstantTok{TRUE}\NormalTok{) }\SpecialCharTok{|\textgreater{}}
  \FunctionTok{tokens\_tolower}\NormalTok{() }\SpecialCharTok{|\textgreater{}}
  \FunctionTok{tokens\_remove}\NormalTok{(}\FunctionTok{stopwords}\NormalTok{(}\StringTok{"en"}\NormalTok{))}
\NormalTok{dfmat }\OtherTok{\textless{}{-}} \FunctionTok{dfm}\NormalTok{(toks) }\SpecialCharTok{|\textgreater{}}
  \FunctionTok{dfm\_trim}\NormalTok{(}\AttributeTok{min\_termfreq =} \DecValTok{5}\NormalTok{, }\AttributeTok{min\_docfreq =} \DecValTok{3}\NormalTok{)}

\CommentTok{\# TF{-}IDF weighting}
\NormalTok{dfmat\_tfidf }\OtherTok{\textless{}{-}} \FunctionTok{dfm\_tfidf}\NormalTok{(dfmat)}

\CommentTok{\# Top features}
\FunctionTok{topfeatures}\NormalTok{(dfmat, }\DecValTok{20}\NormalTok{)}

\CommentTok{\# Keyword{-}in{-}context}
\FunctionTok{kwic}\NormalTok{(toks, }\AttributeTok{pattern =} \StringTok{"open access"}\NormalTok{, }\AttributeTok{window =} \DecValTok{5}\NormalTok{)}
\end{Highlighting}
\end{Shaded}

\section{igraph / tidygraph / ggraph}\label{igraph-tidygraph-ggraph}

\begin{Shaded}
\begin{Highlighting}[]
\FunctionTok{library}\NormalTok{(igraph)}
\FunctionTok{library}\NormalTok{(tidygraph)}
\FunctionTok{library}\NormalTok{(ggraph)}

\CommentTok{\# Build graph from edge list}
\NormalTok{g }\OtherTok{\textless{}{-}} \FunctionTok{graph\_from\_data\_frame}\NormalTok{(edges, }\AttributeTok{directed =} \ConstantTok{FALSE}\NormalTok{, }\AttributeTok{vertices =}\NormalTok{ nodes)}

\CommentTok{\# Tidygraph conversion}
\NormalTok{tg }\OtherTok{\textless{}{-}} \FunctionTok{as\_tbl\_graph}\NormalTok{(g) }\SpecialCharTok{|\textgreater{}}
  \FunctionTok{activate}\NormalTok{(nodes) }\SpecialCharTok{|\textgreater{}}
  \FunctionTok{mutate}\NormalTok{(}
    \AttributeTok{degree =} \FunctionTok{centrality\_degree}\NormalTok{(),}
    \AttributeTok{between =} \FunctionTok{centrality\_betweenness}\NormalTok{(),}
    \AttributeTok{community =} \FunctionTok{group\_leiden}\NormalTok{(}\AttributeTok{resolution =} \FloatTok{1.0}\NormalTok{)}
\NormalTok{  )}

\CommentTok{\# ggraph visualisation}
\FunctionTok{ggraph}\NormalTok{(tg, }\AttributeTok{layout =} \StringTok{"fr"}\NormalTok{) }\SpecialCharTok{+}
  \FunctionTok{geom\_edge\_link}\NormalTok{(}\FunctionTok{aes}\NormalTok{(}\AttributeTok{alpha =}\NormalTok{ weight), }\AttributeTok{show.legend =} \ConstantTok{FALSE}\NormalTok{) }\SpecialCharTok{+}
  \FunctionTok{geom\_node\_point}\NormalTok{(}\FunctionTok{aes}\NormalTok{(}\AttributeTok{size =}\NormalTok{ degree, }\AttributeTok{colour =} \FunctionTok{factor}\NormalTok{(community))) }\SpecialCharTok{+}
  \FunctionTok{geom\_node\_text}\NormalTok{(}\FunctionTok{aes}\NormalTok{(}\AttributeTok{label =}\NormalTok{ name), }\AttributeTok{repel =} \ConstantTok{TRUE}\NormalTok{, }\AttributeTok{size =} \DecValTok{3}\NormalTok{) }\SpecialCharTok{+}
  \FunctionTok{theme\_void}\NormalTok{()}
\end{Highlighting}
\end{Shaded}

\section{Common data wrangling patterns}\label{common-data-wrangling-patterns}

\begin{Shaded}
\begin{Highlighting}[]
\FunctionTok{library}\NormalTok{(tidyverse)}

\CommentTok{\# Unnest OpenAlex authorships}
\NormalTok{authors }\OtherTok{\textless{}{-}}\NormalTok{ works }\SpecialCharTok{|\textgreater{}}
  \FunctionTok{select}\NormalTok{(}\AttributeTok{work\_id =}\NormalTok{ id, authorships) }\SpecialCharTok{|\textgreater{}}
  \FunctionTok{unnest}\NormalTok{(authorships, }\AttributeTok{names\_sep =} \StringTok{"\_"}\NormalTok{)}

\CommentTok{\# Unnest OpenAlex topics}
\NormalTok{topics }\OtherTok{\textless{}{-}}\NormalTok{ works }\SpecialCharTok{|\textgreater{}}
  \FunctionTok{select}\NormalTok{(}\AttributeTok{work\_id =}\NormalTok{ id, topics) }\SpecialCharTok{|\textgreater{}}
  \FunctionTok{unnest}\NormalTok{(topics, }\AttributeTok{names\_sep =} \StringTok{"\_"}\NormalTok{)}

\CommentTok{\# Compute h{-}index}
\NormalTok{compute\_h\_index }\OtherTok{\textless{}{-}} \ControlFlowTok{function}\NormalTok{(citations) \{}
\NormalTok{  citations }\OtherTok{\textless{}{-}} \FunctionTok{sort}\NormalTok{(citations, }\AttributeTok{decreasing =} \ConstantTok{TRUE}\NormalTok{)}
  \FunctionTok{sum}\NormalTok{(citations }\SpecialCharTok{\textgreater{}=} \FunctionTok{seq\_along}\NormalTok{(citations))}
\NormalTok{\}}

\CommentTok{\# Field normalisation}
\NormalTok{works }\SpecialCharTok{|\textgreater{}}
  \FunctionTok{group\_by}\NormalTok{(field, year) }\SpecialCharTok{|\textgreater{}}
  \FunctionTok{mutate}\NormalTok{(}
    \AttributeTok{field\_mean =} \FunctionTok{mean}\NormalTok{(cited\_by\_count),}
    \AttributeTok{mncs =}\NormalTok{ cited\_by\_count }\SpecialCharTok{/} \FunctionTok{pmax}\NormalTok{(field\_mean, }\DecValTok{1}\NormalTok{)}
\NormalTok{  )}

\CommentTok{\# Fractional counting}
\NormalTok{works }\SpecialCharTok{|\textgreater{}}
  \FunctionTok{unnest}\NormalTok{(authorships, }\AttributeTok{names\_sep =} \StringTok{"\_"}\NormalTok{) }\SpecialCharTok{|\textgreater{}}
  \FunctionTok{group\_by}\NormalTok{(id) }\SpecialCharTok{|\textgreater{}}
  \FunctionTok{mutate}\NormalTok{(}\AttributeTok{frac =} \DecValTok{1} \SpecialCharTok{/} \FunctionTok{n}\NormalTok{()) }\SpecialCharTok{|\textgreater{}}
  \FunctionTok{ungroup}\NormalTok{()}
\end{Highlighting}
\end{Shaded}

\section{Companion package functions}\label{companion-package-functions}

The \texttt{scientometricsInR} package provides helper functions used throughout the book:

{\def\LTcaptype{none} % do not increment counter
\begin{longtable}[]{@{}
  >{\raggedright\arraybackslash}p{(\linewidth - 4\tabcolsep) * \real{0.3571}}
  >{\raggedright\arraybackslash}p{(\linewidth - 4\tabcolsep) * \real{0.3214}}
  >{\raggedright\arraybackslash}p{(\linewidth - 4\tabcolsep) * \real{0.3214}}@{}}
\toprule\noalign{}
\begin{minipage}[b]{\linewidth}\raggedright
Function
\end{minipage} & \begin{minipage}[b]{\linewidth}\raggedright
Purpose
\end{minipage} & \begin{minipage}[b]{\linewidth}\raggedright
Example
\end{minipage} \\
\midrule\noalign{}
\endhead
\bottomrule\noalign{}
\endlastfoot
\EscVerb[breaklines=true,breakanywhere=true,breakanywheresymbolpre={},breakanywheresymbolpost={}]{fetch\_openalex(...)} & Cached, rate-limited OpenAlex queries & \EscVerb[breaklines=true,breakanywhere=true,breakanywheresymbolpre={},breakanywheresymbolpost={}]{fetch\_openalex(entity\ =\ "works",\ ...)} \\
\EscVerb[breaklines=true,breakanywhere=true,breakanywheresymbolpre={},breakanywheresymbolpost={}]{dedupe\_by\_doi(df)} & Remove duplicate records by DOI & \EscVerb[breaklines=true,breakanywhere=true,breakanywheresymbolpre={},breakanywheresymbolpost={}]{works\ |>\ dedupe\_by\_doi()} \\
\EscVerb[breaklines=true,breakanywhere=true,breakanywheresymbolpre={},breakanywheresymbolpost={}]{compute\_h\_index(citations)} & Compute h-index from citation vector & \EscVerb[breaklines=true,breakanywhere=true,breakanywheresymbolpre={},breakanywheresymbolpost={}]{compute\_h\_index(c(10,\ 8,\ 5,\ 3,\ 1))} → 3 \\
\EscVerb[breaklines=true,breakanywhere=true,breakanywheresymbolpre={},breakanywheresymbolpost={}]{field\_normalize(cites,\ means)} & MNCS normalisation (vector) & \EscVerb[breaklines=true,breakanywhere=true,breakanywheresymbolpre={},breakanywheresymbolpost={}]{field\_normalize(c(20,\ 5),\ c(10,\ 10))} → \EscVerb[breaklines=true,breakanywhere=true,breakanywheresymbolpre={},breakanywheresymbolpost={}]{c(2.0,\ 0.5)} \\
\EscVerb[breaklines=true,breakanywhere=true,breakanywheresymbolpre={},breakanywheresymbolpost={}]{compute\_mncs(df)} & MNCS for a grouped data frame & \EscVerb[breaklines=true,breakanywhere=true,breakanywheresymbolpre={},breakanywheresymbolpost={}]{df\ |>\ compute\_mncs()} \\
\EscVerb[breaklines=true,breakanywhere=true,breakanywheresymbolpre={},breakanywheresymbolpost={}]{build\_coauth\_graph(works)} & Co-authorship igraph from works & \EscVerb[breaklines=true,breakanywhere=true,breakanywheresymbolpre={},breakanywheresymbolpost={}]{build\_coauth\_graph(works)} \\
\EscVerb[breaklines=true,breakanywhere=true,breakanywheresymbolpre={},breakanywheresymbolpost={}]{kleinberg\_bursts(kw,\ dates)} & Keyword burst detection & \EscVerb[breaklines=true,breakanywhere=true,breakanywheresymbolpre={},breakanywheresymbolpost={}]{kleinberg\_bursts(keywords,\ dates)} \\
\EscVerb[breaklines=true,breakanywhere=true,breakanywheresymbolpre={},breakanywheresymbolpost={}]{palette\_sci(n)} & Viridis colour palette & \EscVerb[breaklines=true,breakanywhere=true,breakanywheresymbolpre={},breakanywheresymbolpost={}]{palette\_sci(4)} \\
\EscVerb[breaklines=true,breakanywhere=true,breakanywheresymbolpre={},breakanywheresymbolpost={}]{theme\_sci()} & Minimal ggplot2 theme & \EscVerb[breaklines=true,breakanywhere=true,breakanywheresymbolpre={},breakanywheresymbolpost={}]{+\ theme\_sci()} \\
\end{longtable}
}

\section{Useful keyboard shortcuts}\label{useful-keyboard-shortcuts}

{\def\LTcaptype{none} % do not increment counter
\begin{longtable}[]{@{}lll@{}}
\toprule\noalign{}
Action & RStudio & VS Code \\
\midrule\noalign{}
\endhead
\bottomrule\noalign{}
\endlastfoot
Run current line/selection & Ctrl+Enter & Ctrl+Enter \\
Run current chunk & Ctrl+Shift+Enter & Ctrl+Shift+Enter \\
Insert pipe \EscVerb[breaklines=true,breakanywhere=true,breakanywheresymbolpre={},breakanywheresymbolpost={}]{|>} & Ctrl+Shift+M & --- \\
Insert assignment \EscVerb[breaklines=true,breakanywhere=true,breakanywheresymbolpre={},breakanywheresymbolpost={}]{<-} & Alt+- & --- \\
Render document & Ctrl+Shift+K & --- \\
Go to file & Ctrl+. & Ctrl+P \\
\end{longtable}
}

\chapter{Exercise Solutions}\label{solutions}

\emph{R. Heller · Date: 2026-05-09 · Last changed: 2026-10-31}

This appendix provides suggested solutions for all chapter exercises.
The code sketches below are meant to guide you toward a working answer; they are
not the only valid approach.
Because each solution references data built during the corresponding chapter,
code chunks are \textbf{not evaluated} here --- run them inside the chapter project
after completing the worked example.

\section*{Chapter 1 --- What Is Scientometrics?}\label{chapter-1-what-is-scientometrics}
\addcontentsline{toc}{section}{Chapter 1 --- What Is Scientometrics?}

\textbf{1. Explore a different field.}
Swap the search term and fetch a fresh sample. Comparing the publication trend
and top journals reveals how different fields vary in volume, growth rate, and
journal concentration.

\begin{Shaded}
\begin{Highlighting}[]
\FunctionTok{library}\NormalTok{(openalexR)}
\FunctionTok{library}\NormalTok{(tidyverse)}

\NormalTok{works\_ml }\OtherTok{\textless{}{-}} \FunctionTok{fetch\_openalex}\NormalTok{(}
 \AttributeTok{entity =} \StringTok{"works"}\NormalTok{,}
 \AttributeTok{search =} \StringTok{"machine learning"}\NormalTok{,}
 \AttributeTok{from\_publication\_date =} \StringTok{"2013{-}01{-}01"}\NormalTok{,}
 \AttributeTok{to\_publication\_date =} \StringTok{"2023{-}12{-}31"}\NormalTok{,}
 \AttributeTok{count\_only =} \ConstantTok{FALSE}\NormalTok{,}
 \AttributeTok{per\_page =} \DecValTok{200}
\NormalTok{) }\SpecialCharTok{|\textgreater{}} \FunctionTok{slice\_sample}\NormalTok{(}\AttributeTok{n =} \DecValTok{500}\NormalTok{)}

\NormalTok{works\_ml }\SpecialCharTok{|\textgreater{}}
 \FunctionTok{count}\NormalTok{(publication\_year) }\SpecialCharTok{|\textgreater{}}
 \FunctionTok{ggplot}\NormalTok{(}\FunctionTok{aes}\NormalTok{(publication\_year, n)) }\SpecialCharTok{+}
 \FunctionTok{geom\_col}\NormalTok{() }\SpecialCharTok{+}
 \FunctionTok{labs}\NormalTok{(}\AttributeTok{title =} \StringTok{"Machine{-}learning publications per year"}\NormalTok{,}
      \AttributeTok{x =} \StringTok{"Year"}\NormalTok{, }\AttributeTok{y =} \StringTok{"Works"}\NormalTok{)}
\end{Highlighting}
\end{Shaded}

\emph{Key insight:} High-growth fields like machine learning show steeper annual
increases and a broader spread across journals compared to a niche field like
scientometrics.

\textbf{2. Uncited works.}
Filter to zero-citation records and examine the year distribution.

\begin{Shaded}
\begin{Highlighting}[]
\NormalTok{works }\SpecialCharTok{|\textgreater{}}
 \FunctionTok{filter}\NormalTok{(cited\_by\_count }\SpecialCharTok{==} \DecValTok{0}\NormalTok{) }\SpecialCharTok{|\textgreater{}}
 \FunctionTok{count}\NormalTok{(publication\_year) }\SpecialCharTok{|\textgreater{}}
 \FunctionTok{mutate}\NormalTok{(}\AttributeTok{pct =}\NormalTok{ n }\SpecialCharTok{/} \FunctionTok{sum}\NormalTok{(n)) }\SpecialCharTok{|\textgreater{}}
 \FunctionTok{ggplot}\NormalTok{(}\FunctionTok{aes}\NormalTok{(publication\_year, pct)) }\SpecialCharTok{+}
 \FunctionTok{geom\_col}\NormalTok{() }\SpecialCharTok{+}
 \FunctionTok{scale\_y\_continuous}\NormalTok{(}\AttributeTok{labels =}\NormalTok{ scales}\SpecialCharTok{::}\NormalTok{percent) }\SpecialCharTok{+}
 \FunctionTok{labs}\NormalTok{(}\AttributeTok{title =} \StringTok{"Distribution of uncited works by year"}\NormalTok{,}
      \AttributeTok{x =} \StringTok{"Year"}\NormalTok{, }\AttributeTok{y =} \StringTok{"Share of uncited works"}\NormalTok{)}
\end{Highlighting}
\end{Shaded}

\emph{Key insight:} Recent publication years dominate because papers need time to
accumulate citations --- a window effect.

\textbf{3. Open Access status.}
Compute yearly OA proportions to reveal the trend.

\begin{Shaded}
\begin{Highlighting}[]
\NormalTok{works }\SpecialCharTok{|\textgreater{}}
 \FunctionTok{mutate}\NormalTok{(}\AttributeTok{year =}\NormalTok{ publication\_year) }\SpecialCharTok{|\textgreater{}}
 \FunctionTok{group\_by}\NormalTok{(year) }\SpecialCharTok{|\textgreater{}}
 \FunctionTok{summarise}\NormalTok{(}\AttributeTok{oa\_share =} \FunctionTok{mean}\NormalTok{(is\_oa, }\AttributeTok{na.rm =} \ConstantTok{TRUE}\NormalTok{)) }\SpecialCharTok{|\textgreater{}}
 \FunctionTok{ggplot}\NormalTok{(}\FunctionTok{aes}\NormalTok{(year, oa\_share)) }\SpecialCharTok{+}
 \FunctionTok{geom\_line}\NormalTok{() }\SpecialCharTok{+}
 \FunctionTok{geom\_point}\NormalTok{() }\SpecialCharTok{+}
 \FunctionTok{scale\_y\_continuous}\NormalTok{(}\AttributeTok{labels =}\NormalTok{ scales}\SpecialCharTok{::}\NormalTok{percent) }\SpecialCharTok{+}
 \FunctionTok{labs}\NormalTok{(}\AttributeTok{title =} \StringTok{"Open{-}access share over time"}\NormalTok{,}
      \AttributeTok{x =} \StringTok{"Year"}\NormalTok{, }\AttributeTok{y =} \StringTok{"OA proportion"}\NormalTok{)}
\end{Highlighting}
\end{Shaded}

\emph{Key insight:} OA proportions have risen steadily, driven by funder mandates and
the growth of gold OA journals.

\textbf{4. Growth curve.}
Plot annual counts on a log scale; linearity implies exponential growth.

\begin{Shaded}
\begin{Highlighting}[]
\NormalTok{works }\SpecialCharTok{|\textgreater{}}
 \FunctionTok{filter}\NormalTok{(publication\_year }\SpecialCharTok{\textgreater{}=} \DecValTok{2000}\NormalTok{, publication\_year }\SpecialCharTok{\textless{}=} \DecValTok{2023}\NormalTok{) }\SpecialCharTok{|\textgreater{}}
 \FunctionTok{count}\NormalTok{(publication\_year) }\SpecialCharTok{|\textgreater{}}
 \FunctionTok{ggplot}\NormalTok{(}\FunctionTok{aes}\NormalTok{(publication\_year, n)) }\SpecialCharTok{+}
 \FunctionTok{geom\_point}\NormalTok{() }\SpecialCharTok{+}
 \FunctionTok{geom\_smooth}\NormalTok{(}\AttributeTok{method =} \StringTok{"lm"}\NormalTok{) }\SpecialCharTok{+}
 \FunctionTok{scale\_y\_log10}\NormalTok{() }\SpecialCharTok{+}
 \FunctionTok{labs}\NormalTok{(}\AttributeTok{title =} \StringTok{"Publication growth (log scale)"}\NormalTok{,}
      \AttributeTok{x =} \StringTok{"Year"}\NormalTok{, }\AttributeTok{y =} \StringTok{"Works (log)"}\NormalTok{)}
\end{Highlighting}
\end{Shaded}

\emph{Key insight:} If the points fall roughly on a straight line, growth is
approximately exponential, consistent with Price's law.

\section*{Chapter 2 --- Theoretical Underpinnings}\label{chapter-2-theoretical-underpinnings}
\addcontentsline{toc}{section}{Chapter 2 --- Theoretical Underpinnings}

\textbf{1. Lotka exponent.}
Fit a linear model in log-log space to estimate the exponent.

\begin{Shaded}
\begin{Highlighting}[]
\FunctionTok{library}\NormalTok{(tidyverse)}

\NormalTok{author\_prod }\OtherTok{\textless{}{-}}\NormalTok{ works }\SpecialCharTok{|\textgreater{}}
 \FunctionTok{unnest}\NormalTok{(author) }\SpecialCharTok{|\textgreater{}}
 \FunctionTok{count}\NormalTok{(au\_id, }\AttributeTok{name =} \StringTok{"n\_papers"}\NormalTok{)}

\NormalTok{freq\_tbl }\OtherTok{\textless{}{-}}\NormalTok{ author\_prod }\SpecialCharTok{|\textgreater{}}
 \FunctionTok{count}\NormalTok{(n\_papers, }\AttributeTok{name =} \StringTok{"n\_authors"}\NormalTok{)}

\NormalTok{fit }\OtherTok{\textless{}{-}} \FunctionTok{lm}\NormalTok{(}\FunctionTok{log}\NormalTok{(n\_authors) }\SpecialCharTok{\textasciitilde{}} \FunctionTok{log}\NormalTok{(n\_papers), }\AttributeTok{data =}\NormalTok{ freq\_tbl)}
\FunctionTok{coef}\NormalTok{(fit)  }\CommentTok{\# slope ≈ {-}2 if Lotka holds}

\FunctionTok{ggplot}\NormalTok{(freq\_tbl, }\FunctionTok{aes}\NormalTok{(n\_papers, n\_authors)) }\SpecialCharTok{+}
 \FunctionTok{geom\_point}\NormalTok{() }\SpecialCharTok{+}
 \FunctionTok{scale\_x\_log10}\NormalTok{() }\SpecialCharTok{+}
 \FunctionTok{scale\_y\_log10}\NormalTok{() }\SpecialCharTok{+}
 \FunctionTok{geom\_smooth}\NormalTok{(}\AttributeTok{method =} \StringTok{"lm"}\NormalTok{, }\AttributeTok{se =} \ConstantTok{FALSE}\NormalTok{) }\SpecialCharTok{+}
 \FunctionTok{labs}\NormalTok{(}\AttributeTok{title =} \StringTok{"Lotka\textquotesingle{}s law fit"}\NormalTok{,}
      \AttributeTok{x =} \StringTok{"Papers per author (log)"}\NormalTok{, }\AttributeTok{y =} \StringTok{"Number of authors (log)"}\NormalTok{)}
\end{Highlighting}
\end{Shaded}

\emph{Key insight:} Lotka predicted an exponent near 2; empirical values typically
range from 1.5 to 3.5, depending on the field and time window.

\textbf{2. Bradford zones.}
Divide journals into three zones of roughly equal article counts.

\begin{Shaded}
\begin{Highlighting}[]
\NormalTok{journal\_counts }\OtherTok{\textless{}{-}}\NormalTok{ works }\SpecialCharTok{|\textgreater{}}
 \FunctionTok{count}\NormalTok{(so, }\AttributeTok{sort =} \ConstantTok{TRUE}\NormalTok{) }\SpecialCharTok{|\textgreater{}}
 \FunctionTok{mutate}\NormalTok{(}\AttributeTok{cum =} \FunctionTok{cumsum}\NormalTok{(n),}
        \AttributeTok{total =} \FunctionTok{sum}\NormalTok{(n),}
        \AttributeTok{zone =} \FunctionTok{case\_when}\NormalTok{(}
\NormalTok{         cum }\SpecialCharTok{\textless{}=}\NormalTok{ total }\SpecialCharTok{/} \DecValTok{3} \SpecialCharTok{\textasciitilde{}} \StringTok{"Core"}\NormalTok{,}
\NormalTok{         cum }\SpecialCharTok{\textless{}=} \DecValTok{2} \SpecialCharTok{*}\NormalTok{ total }\SpecialCharTok{/} \DecValTok{3} \SpecialCharTok{\textasciitilde{}} \StringTok{"Mid"}\NormalTok{,}
         \ConstantTok{TRUE} \SpecialCharTok{\textasciitilde{}} \StringTok{"Periphery"}
\NormalTok{        ))}

\NormalTok{journal\_counts }\SpecialCharTok{|\textgreater{}} \FunctionTok{count}\NormalTok{(zone)}
\end{Highlighting}
\end{Shaded}

\emph{Key insight:} The core zone contains very few journals; the periphery contains
many. The ratio between successive zone sizes should approximate Bradford's
multiplier.

\textbf{3. Your field's Zipf distribution.}
Extract title words and plot rank vs.~frequency on log-log axes.

\begin{Shaded}
\begin{Highlighting}[]
\FunctionTok{library}\NormalTok{(tidytext)}

\NormalTok{word\_freq }\OtherTok{\textless{}{-}}\NormalTok{ works }\SpecialCharTok{|\textgreater{}}
 \FunctionTok{select}\NormalTok{(id, display\_name) }\SpecialCharTok{|\textgreater{}}
 \FunctionTok{unnest\_tokens}\NormalTok{(word, display\_name) }\SpecialCharTok{|\textgreater{}}
 \FunctionTok{anti\_join}\NormalTok{(stop\_words, }\AttributeTok{by =} \StringTok{"word"}\NormalTok{) }\SpecialCharTok{|\textgreater{}}
 \FunctionTok{count}\NormalTok{(word, }\AttributeTok{sort =} \ConstantTok{TRUE}\NormalTok{) }\SpecialCharTok{|\textgreater{}}
 \FunctionTok{mutate}\NormalTok{(}\AttributeTok{rank =} \FunctionTok{row\_number}\NormalTok{())}

\FunctionTok{ggplot}\NormalTok{(word\_freq, }\FunctionTok{aes}\NormalTok{(rank, n)) }\SpecialCharTok{+}
 \FunctionTok{geom\_line}\NormalTok{() }\SpecialCharTok{+}
 \FunctionTok{scale\_x\_log10}\NormalTok{() }\SpecialCharTok{+}
 \FunctionTok{scale\_y\_log10}\NormalTok{() }\SpecialCharTok{+}
 \FunctionTok{labs}\NormalTok{(}\AttributeTok{title =} \StringTok{"Zipf\textquotesingle{}s law for title words"}\NormalTok{,}
      \AttributeTok{x =} \StringTok{"Rank (log)"}\NormalTok{, }\AttributeTok{y =} \StringTok{"Frequency (log)"}\NormalTok{)}
\end{Highlighting}
\end{Shaded}

\emph{Key insight:} A roughly linear relationship on log-log axes confirms Zipf's
law. Domain-specific terms dominate the top ranks.

\textbf{4. Matthew Effect simulation.}
Simulate preferential attachment and compare the resulting distribution to
Lotka's law.

\begin{Shaded}
\begin{Highlighting}[]
\FunctionTok{set.seed}\NormalTok{(}\DecValTok{20260509}\NormalTok{)}
\NormalTok{n\_authors }\OtherTok{\textless{}{-}} \DecValTok{100}
\NormalTok{papers }\OtherTok{\textless{}{-}} \FunctionTok{rep}\NormalTok{(}\DecValTok{1}\NormalTok{, n\_authors)}

\ControlFlowTok{for}\NormalTok{ (step }\ControlFlowTok{in} \FunctionTok{seq\_len}\NormalTok{(}\DecValTok{50}\NormalTok{)) \{}
\NormalTok{ probs }\OtherTok{\textless{}{-}}\NormalTok{ papers }\SpecialCharTok{/} \FunctionTok{sum}\NormalTok{(papers)}
\NormalTok{ winner }\OtherTok{\textless{}{-}} \FunctionTok{sample}\NormalTok{(n\_authors, }\DecValTok{1}\NormalTok{, }\AttributeTok{prob =}\NormalTok{ probs)}
\NormalTok{ papers[winner] }\OtherTok{\textless{}{-}}\NormalTok{ papers[winner] }\SpecialCharTok{+} \DecValTok{1}
\NormalTok{\}}

\FunctionTok{tibble}\NormalTok{(}\AttributeTok{papers =}\NormalTok{ papers) }\SpecialCharTok{|\textgreater{}}
 \FunctionTok{count}\NormalTok{(papers, }\AttributeTok{name =} \StringTok{"n\_authors"}\NormalTok{) }\SpecialCharTok{|\textgreater{}}
 \FunctionTok{ggplot}\NormalTok{(}\FunctionTok{aes}\NormalTok{(papers, n\_authors)) }\SpecialCharTok{+}
 \FunctionTok{geom\_point}\NormalTok{() }\SpecialCharTok{+}
 \FunctionTok{scale\_x\_log10}\NormalTok{() }\SpecialCharTok{+}
 \FunctionTok{scale\_y\_log10}\NormalTok{() }\SpecialCharTok{+}
 \FunctionTok{labs}\NormalTok{(}\AttributeTok{title =} \StringTok{"Simulated preferential attachment"}\NormalTok{,}
      \AttributeTok{x =} \StringTok{"Papers (log)"}\NormalTok{, }\AttributeTok{y =} \StringTok{"Authors (log)"}\NormalTok{)}
\end{Highlighting}
\end{Shaded}

\emph{Key insight:} Preferential attachment generates a skewed distribution
resembling Lotka's law, illustrating the Matthew Effect.

\section*{Chapter 3 --- Ethics and Responsible Metrics}\label{chapter-3-ethics-and-responsible-metrics}
\addcontentsline{toc}{section}{Chapter 3 --- Ethics and Responsible Metrics}

\textbf{1. Leiden Manifesto audit.}
This is a qualitative exercise. Evaluate a real hiring policy against the ten
principles.

\begin{Shaded}
\begin{Highlighting}[]
\CommentTok{\# No code required. Suggested structure:}
\CommentTok{\# For each of the 10 Leiden Manifesto principles:}
\CommentTok{\#   1. State the principle}
\CommentTok{\#   2. Quote the relevant policy text (if any)}
\CommentTok{\#   3. Assess: satisfied / partially satisfied / violated}
\CommentTok{\#   4. Explain your reasoning}
\CommentTok{\# Summarise overall compliance in a paragraph.}
\end{Highlighting}
\end{Shaded}

\emph{Key insight:} Most policies rely heavily on the h-index or JIF, violating
principles 6 (field variation) and 7 (qualitative portfolio assessment).

\textbf{2. Self-citation across fields.}
Compare self-citation rates between two different fields.

\begin{Shaded}
\begin{Highlighting}[]
\FunctionTok{library}\NormalTok{(openalexR)}
\FunctionTok{library}\NormalTok{(tidyverse)}

\NormalTok{fields }\OtherTok{\textless{}{-}} \FunctionTok{c}\NormalTok{(}\StringTok{"bibliometrics"}\NormalTok{, }\StringTok{"genomics"}\NormalTok{)}
\NormalTok{results }\OtherTok{\textless{}{-}} \FunctionTok{map}\NormalTok{(fields, \textbackslash{}(term) \{}
\NormalTok{ w }\OtherTok{\textless{}{-}} \FunctionTok{fetch\_openalex}\NormalTok{(}\AttributeTok{entity =} \StringTok{"works"}\NormalTok{, }\AttributeTok{search =}\NormalTok{ term,}
               \AttributeTok{from\_publication\_date =} \StringTok{"2020{-}01{-}01"}\NormalTok{,}
               \AttributeTok{to\_publication\_date =} \StringTok{"2023{-}12{-}31"}\NormalTok{) }\SpecialCharTok{|\textgreater{}}
   \FunctionTok{slice\_sample}\NormalTok{(}\AttributeTok{n =} \DecValTok{200}\NormalTok{)}
\NormalTok{ w }\SpecialCharTok{|\textgreater{}} \FunctionTok{mutate}\NormalTok{(}\AttributeTok{field =}\NormalTok{ term)}
\NormalTok{\}) }\SpecialCharTok{|\textgreater{}} \FunctionTok{list\_rbind}\NormalTok{()}

\CommentTok{\# Self{-}citation proxy: papers citing other papers by the same first author}
\NormalTok{results }\SpecialCharTok{|\textgreater{}}
 \FunctionTok{group\_by}\NormalTok{(field) }\SpecialCharTok{|\textgreater{}}
 \FunctionTok{summarise}\NormalTok{(}\AttributeTok{mean\_cites =} \FunctionTok{mean}\NormalTok{(cited\_by\_count, }\AttributeTok{na.rm =} \ConstantTok{TRUE}\NormalTok{),}
           \AttributeTok{median\_cites =} \FunctionTok{median}\NormalTok{(cited\_by\_count, }\AttributeTok{na.rm =} \ConstantTok{TRUE}\NormalTok{))}
\end{Highlighting}
\end{Shaded}

\emph{Key insight:} Self-citation baselines differ by field. Higher self-citation in
some fields is structural (smaller communities) rather than manipulative.

\textbf{3. DORA compliance check.}
A research exercise using the DORA signatory database.

\begin{Shaded}
\begin{Highlighting}[]
\CommentTok{\# No code required. Visit https://sfdora.org/signers/}
\CommentTok{\# Select three major funders and compare their grant evaluation}
\CommentTok{\# criteria against DORA commitments:}
\CommentTok{\#   {-} Do they avoid using JIF in evaluation?}
\CommentTok{\#   {-} Do they assess research outputs on their own merit?}
\CommentTok{\#   {-} Do they consider a broad range of impact measures?}
\end{Highlighting}
\end{Shaded}

\emph{Key insight:} Many signatories still reference journal prestige in practice,
revealing a gap between stated commitments and implementation.

\textbf{4. Goodhart's Law thought experiment.}
Propose a new indicator and analyse potential gaming strategies.

\begin{Shaded}
\begin{Highlighting}[]
\CommentTok{\# Example indicator: "Replication Impact Score"}
\CommentTok{\#   = fraction of a researcher\textquotesingle{}s papers that have been independently replicated}
\CommentTok{\#}
\CommentTok{\# Potential gaming strategies:}
\CommentTok{\#   {-} Publish trivially replicable results}
\CommentTok{\#   {-} Arrange replication exchanges with collaborators}
\CommentTok{\#   {-} Focus on small, easily replicated studies rather than ambitious work}
\CommentTok{\#}
\CommentTok{\# The numerator (replications) and denominator (total papers) are both gameable.}
\CommentTok{\# This illustrates Goodhart\textquotesingle{}s Law: once a measure becomes a target,}
\CommentTok{\# it ceases to be a good measure.}
\end{Highlighting}
\end{Shaded}

\emph{Key insight:} Any ratio-based metric can be gamed on both numerator and
denominator. Indicator design must anticipate behavioural responses.

\section*{Chapter 4 --- Data Sources Compared}\label{chapter-4-data-sources-compared}
\addcontentsline{toc}{section}{Chapter 4 --- Data Sources Compared}

\textbf{1. Coverage gap analysis.}
Query the same DOIs in OpenAlex and Crossref and compare.

\begin{Shaded}
\begin{Highlighting}[]
\FunctionTok{library}\NormalTok{(openalexR)}
\FunctionTok{library}\NormalTok{(rcrossref)}
\FunctionTok{library}\NormalTok{(tidyverse)}

\NormalTok{dois }\OtherTok{\textless{}{-}}\NormalTok{ works }\SpecialCharTok{|\textgreater{}} \FunctionTok{slice\_sample}\NormalTok{(}\AttributeTok{n =} \DecValTok{20}\NormalTok{) }\SpecialCharTok{|\textgreater{}} \FunctionTok{pull}\NormalTok{(doi) }\SpecialCharTok{|\textgreater{}} \FunctionTok{na.omit}\NormalTok{()}

\NormalTok{oa\_results }\OtherTok{\textless{}{-}} \FunctionTok{map}\NormalTok{(dois, \textbackslash{}(d) \{}
 \FunctionTok{tryCatch}\NormalTok{(}\FunctionTok{fetch\_openalex}\NormalTok{(}\AttributeTok{entity =} \StringTok{"works"}\NormalTok{, }\AttributeTok{doi =}\NormalTok{ d), }\AttributeTok{error =}\NormalTok{ \textbackslash{}(e) }\ConstantTok{NULL}\NormalTok{)}
\NormalTok{\}) }\SpecialCharTok{|\textgreater{}} \FunctionTok{compact}\NormalTok{()}
\NormalTok{oa\_found }\OtherTok{\textless{}{-}} \FunctionTok{length}\NormalTok{(oa\_results)}

\NormalTok{cr\_results }\OtherTok{\textless{}{-}} \FunctionTok{map}\NormalTok{(dois, \textbackslash{}(d) \{}
 \FunctionTok{tryCatch}\NormalTok{(}\FunctionTok{cr\_works}\NormalTok{(d)}\SpecialCharTok{$}\NormalTok{data, }\AttributeTok{error =}\NormalTok{ \textbackslash{}(e) }\ConstantTok{NULL}\NormalTok{)}
\NormalTok{\}) }\SpecialCharTok{|\textgreater{}} \FunctionTok{compact}\NormalTok{()}
\NormalTok{cr\_found }\OtherTok{\textless{}{-}} \FunctionTok{length}\NormalTok{(cr\_results)}

\FunctionTok{tibble}\NormalTok{(}
 \AttributeTok{Source =} \FunctionTok{c}\NormalTok{(}\StringTok{"OpenAlex"}\NormalTok{, }\StringTok{"Crossref"}\NormalTok{),}
 \AttributeTok{Found =} \FunctionTok{c}\NormalTok{(oa\_found, cr\_found),}
 \AttributeTok{Total =} \FunctionTok{length}\NormalTok{(dois),}
 \AttributeTok{Coverage =}\NormalTok{ Found }\SpecialCharTok{/}\NormalTok{ Total}
\NormalTok{)}
\end{Highlighting}
\end{Shaded}

\emph{Key insight:} OpenAlex generally has high DOI coverage because it ingests
Crossref data, but abstract availability may differ.

\textbf{2. Disciplinary bias.}
Compare indexed works across contrasting fields.

\begin{Shaded}
\begin{Highlighting}[]
\NormalTok{fields }\OtherTok{\textless{}{-}} \FunctionTok{list}\NormalTok{(}
 \AttributeTok{biology =} \StringTok{"https://openalex.org/C86803240"}\NormalTok{,}
 \AttributeTok{philosophy =} \StringTok{"https://openalex.org/C138885662"}
\NormalTok{)}

\NormalTok{coverage }\OtherTok{\textless{}{-}} \FunctionTok{imap}\NormalTok{(fields, \textbackslash{}(concept\_id, name) \{}
 \FunctionTok{fetch\_openalex}\NormalTok{(}\AttributeTok{entity =} \StringTok{"works"}\NormalTok{, }\AttributeTok{concept.id =}\NormalTok{ concept\_id,}
          \AttributeTok{count\_only =} \ConstantTok{TRUE}\NormalTok{) }\SpecialCharTok{|\textgreater{}}
   \FunctionTok{mutate}\NormalTok{(}\AttributeTok{field =}\NormalTok{ name)}
\NormalTok{\}) }\SpecialCharTok{|\textgreater{}} \FunctionTok{list\_rbind}\NormalTok{()}

\NormalTok{coverage}
\end{Highlighting}
\end{Shaded}

\emph{Key insight:} STEM fields are substantially better covered than humanities,
reflecting the historical journal-centricity of bibliometric databases.

\textbf{3. Citation count concordance.}
Correlate citation counts from two sources.

\begin{Shaded}
\begin{Highlighting}[]
\FunctionTok{library}\NormalTok{(rcrossref)}

\NormalTok{comparison }\OtherTok{\textless{}{-}}\NormalTok{ works }\SpecialCharTok{|\textgreater{}}
 \FunctionTok{slice\_sample}\NormalTok{(}\AttributeTok{n =} \DecValTok{50}\NormalTok{) }\SpecialCharTok{|\textgreater{}}
 \FunctionTok{filter}\NormalTok{(}\SpecialCharTok{!}\FunctionTok{is.na}\NormalTok{(doi)) }\SpecialCharTok{|\textgreater{}}
 \FunctionTok{rowwise}\NormalTok{() }\SpecialCharTok{|\textgreater{}}
 \FunctionTok{mutate}\NormalTok{(}\AttributeTok{cr\_cites =} \FunctionTok{tryCatch}\NormalTok{(}
   \FunctionTok{cr\_works}\NormalTok{(doi)}\SpecialCharTok{$}\NormalTok{data}\SpecialCharTok{$}\NormalTok{is\_referenced\_by\_count,}
   \AttributeTok{error =}\NormalTok{ \textbackslash{}(e) }\ConstantTok{NA\_integer\_}
\NormalTok{ )) }\SpecialCharTok{|\textgreater{}}
 \FunctionTok{ungroup}\NormalTok{()}

\FunctionTok{cor}\NormalTok{(comparison}\SpecialCharTok{$}\NormalTok{cited\_by\_count, comparison}\SpecialCharTok{$}\NormalTok{cr\_cites,}
    \AttributeTok{use =} \StringTok{"complete.obs"}\NormalTok{, }\AttributeTok{method =} \StringTok{"spearman"}\NormalTok{)}

\FunctionTok{ggplot}\NormalTok{(comparison, }\FunctionTok{aes}\NormalTok{(cited\_by\_count, cr\_cites)) }\SpecialCharTok{+}
 \FunctionTok{geom\_point}\NormalTok{() }\SpecialCharTok{+}
 \FunctionTok{geom\_abline}\NormalTok{(}\AttributeTok{slope =} \DecValTok{1}\NormalTok{, }\AttributeTok{intercept =} \DecValTok{0}\NormalTok{, }\AttributeTok{linetype =} \StringTok{"dashed"}\NormalTok{) }\SpecialCharTok{+}
 \FunctionTok{labs}\NormalTok{(}\AttributeTok{x =} \StringTok{"OpenAlex citations"}\NormalTok{, }\AttributeTok{y =} \StringTok{"Crossref citations"}\NormalTok{,}
      \AttributeTok{title =} \StringTok{"Citation count concordance"}\NormalTok{)}
\end{Highlighting}
\end{Shaded}

\emph{Key insight:} Counts correlate strongly but are not identical; differences
arise from coverage scope and update frequency.

\textbf{4. Preprint tracking.}
Check whether bioRxiv DOIs appear in OpenAlex with published links.

\begin{Shaded}
\begin{Highlighting}[]
\NormalTok{biorxiv\_dois }\OtherTok{\textless{}{-}} \FunctionTok{c}\NormalTok{(}
 \StringTok{"10.1101/2023.01.01.123456"}\NormalTok{,}
 \StringTok{"10.1101/2023.02.15.234567"}
 \CommentTok{\# ... add 8 more real bioRxiv DOIs}
\NormalTok{)}

\NormalTok{preprint\_status }\OtherTok{\textless{}{-}} \FunctionTok{map}\NormalTok{(biorxiv\_dois, \textbackslash{}(d) \{}
 \FunctionTok{tryCatch}\NormalTok{(\{}
\NormalTok{   w }\OtherTok{\textless{}{-}} \FunctionTok{fetch\_openalex}\NormalTok{(}\AttributeTok{entity =} \StringTok{"works"}\NormalTok{, }\AttributeTok{doi =}\NormalTok{ d)}
   \FunctionTok{tibble}\NormalTok{(}\AttributeTok{doi =}\NormalTok{ d, }\AttributeTok{found =} \ConstantTok{TRUE}\NormalTok{,}
          \AttributeTok{has\_published =} \SpecialCharTok{!}\FunctionTok{is.null}\NormalTok{(w}\SpecialCharTok{$}\NormalTok{ids}\SpecialCharTok{$}\NormalTok{doi))}
\NormalTok{ \}, }\AttributeTok{error =}\NormalTok{ \textbackslash{}(e) }\FunctionTok{tibble}\NormalTok{(}\AttributeTok{doi =}\NormalTok{ d, }\AttributeTok{found =} \ConstantTok{FALSE}\NormalTok{, }\AttributeTok{has\_published =} \ConstantTok{FALSE}\NormalTok{))}
\NormalTok{\}) }\SpecialCharTok{|\textgreater{}} \FunctionTok{list\_rbind}\NormalTok{()}

\FunctionTok{mean}\NormalTok{(preprint\_status}\SpecialCharTok{$}\NormalTok{found)}
\FunctionTok{mean}\NormalTok{(preprint\_status}\SpecialCharTok{$}\NormalTok{has\_published)}
\end{Highlighting}
\end{Shaded}

\emph{Key insight:} Most recent bioRxiv preprints appear in OpenAlex; the proportion
linked to a published version depends on recency and field.

\textbf{5. Source selection exercise.}
This is a qualitative exercise requiring no code.

\begin{Shaded}
\begin{Highlighting}[]
\CommentTok{\# Recommended answer structure:}
\CommentTok{\# 1. OpenAlex: free, broad coverage, good for large{-}scale analysis}
\CommentTok{\#    {-} Limitation: may under{-}represent regional journals}
\CommentTok{\# 2. African Journals Online (AJOL): specialised coverage of}
\CommentTok{\#    Sub{-}Saharan African journals}
\CommentTok{\# 3. Combination approach: use OpenAlex for global context,}
\CommentTok{\#    supplement with AJOL for regional coverage}
\CommentTok{\#}
\CommentTok{\# Justification should address:}
\CommentTok{\#   {-} Coverage of nursing literature}
\CommentTok{\#   {-} Representation of Sub{-}Saharan African journals}
\CommentTok{\#   {-} Cost (OpenAlex is free; Scopus/WoS require subscriptions)}
\CommentTok{\#   {-} Reproducibility (OpenAlex is open; others are proprietary)}
\end{Highlighting}
\end{Shaded}

\emph{Key insight:} No single source is sufficient; combining an open global database
with a regional index maximises both coverage and reproducibility.

\section*{Chapter 5 --- APIs and Packages}\label{chapter-5-apis-and-packages}
\addcontentsline{toc}{section}{Chapter 5 --- APIs and Packages}

\textbf{1. Fetch works by institution.}
Find your institution in OpenAlex and retrieve its works.

\begin{Shaded}
\begin{Highlighting}[]
\FunctionTok{library}\NormalTok{(openalexR)}
\FunctionTok{library}\NormalTok{(tidyverse)}

\CommentTok{\# Step 1: find the institution}
\NormalTok{my\_inst }\OtherTok{\textless{}{-}} \FunctionTok{fetch\_openalex}\NormalTok{(}\AttributeTok{entity =} \StringTok{"institutions"}\NormalTok{,}
                    \AttributeTok{search =} \StringTok{"University of Example"}\NormalTok{)}
\NormalTok{my\_inst\_id }\OtherTok{\textless{}{-}}\NormalTok{ my\_inst}\SpecialCharTok{$}\NormalTok{id[}\DecValTok{1}\NormalTok{]}

\CommentTok{\# Step 2: fetch works}
\NormalTok{inst\_works }\OtherTok{\textless{}{-}} \FunctionTok{fetch\_openalex}\NormalTok{(}\AttributeTok{entity =} \StringTok{"works"}\NormalTok{,}
                       \AttributeTok{authorships.institutions.lineage =}\NormalTok{ my\_inst\_id,}
                       \AttributeTok{from\_publication\_date =} \StringTok{"2020{-}01{-}01"}\NormalTok{)}

\FunctionTok{nrow}\NormalTok{(inst\_works)}
\FunctionTok{mean}\NormalTok{(}\SpecialCharTok{!}\FunctionTok{is.na}\NormalTok{(inst\_works}\SpecialCharTok{$}\NormalTok{doi))  }\CommentTok{\# DOI coverage rate}
\end{Highlighting}
\end{Shaded}

\emph{Key insight:} DOI coverage varies by institution; institutions publishing more
in regional or non-English journals tend to have lower DOI rates.

\textbf{2. Compare Crossref metadata.}
Retrieve the same work from both APIs and compare.

\begin{Shaded}
\begin{Highlighting}[]
\FunctionTok{library}\NormalTok{(rcrossref)}

\NormalTok{sample\_doi }\OtherTok{\textless{}{-}}\NormalTok{ inst\_works}\SpecialCharTok{$}\NormalTok{doi[}\DecValTok{1}\NormalTok{]}

\NormalTok{oa\_meta }\OtherTok{\textless{}{-}} \FunctionTok{fetch\_openalex}\NormalTok{(}\AttributeTok{entity =} \StringTok{"works"}\NormalTok{, }\AttributeTok{doi =}\NormalTok{ sample\_doi)}
\NormalTok{cr\_meta }\OtherTok{\textless{}{-}} \FunctionTok{cr\_works}\NormalTok{(sample\_doi)}\SpecialCharTok{$}\NormalTok{data}

\FunctionTok{tibble}\NormalTok{(}
 \AttributeTok{field =} \FunctionTok{c}\NormalTok{(}\StringTok{"citations"}\NormalTok{, }\StringTok{"title\_match"}\NormalTok{),}
 \AttributeTok{OpenAlex =} \FunctionTok{c}\NormalTok{(oa\_meta}\SpecialCharTok{$}\NormalTok{cited\_by\_count,}
\NormalTok{              oa\_meta}\SpecialCharTok{$}\NormalTok{display\_name),}
 \AttributeTok{Crossref =} \FunctionTok{c}\NormalTok{(cr\_meta}\SpecialCharTok{$}\NormalTok{is\_referenced\_by\_count,}
\NormalTok{              cr\_meta}\SpecialCharTok{$}\NormalTok{title)}
\NormalTok{)}
\end{Highlighting}
\end{Shaded}

\emph{Key insight:} Citation counts differ because each database has a different
citation graph. Crossref counts only DOI-based citations; OpenAlex resolves more.

\textbf{3. Build a multi-year corpus.}
Fetch all works from a journal across multiple years.

\begin{Shaded}
\begin{Highlighting}[]
\NormalTok{journal\_works }\OtherTok{\textless{}{-}} \FunctionTok{fetch\_openalex}\NormalTok{(}
 \AttributeTok{entity =} \StringTok{"works"}\NormalTok{,}
 \AttributeTok{primary\_location.source.id =} \StringTok{"S4210174107"}\NormalTok{,  }\CommentTok{\# example journal ID}
 \AttributeTok{from\_publication\_date =} \StringTok{"2015{-}01{-}01"}\NormalTok{,}
 \AttributeTok{to\_publication\_date =} \StringTok{"2023{-}12{-}31"}
\NormalTok{)}

\NormalTok{journal\_works }\SpecialCharTok{|\textgreater{}}
 \FunctionTok{group\_by}\NormalTok{(publication\_year) }\SpecialCharTok{|\textgreater{}}
 \FunctionTok{summarise}\NormalTok{(}\AttributeTok{n =} \FunctionTok{n}\NormalTok{(), }\AttributeTok{median\_cites =} \FunctionTok{median}\NormalTok{(cited\_by\_count)) }\SpecialCharTok{|\textgreater{}}
 \FunctionTok{pivot\_longer}\NormalTok{(}\FunctionTok{c}\NormalTok{(n, median\_cites)) }\SpecialCharTok{|\textgreater{}}
 \FunctionTok{ggplot}\NormalTok{(}\FunctionTok{aes}\NormalTok{(publication\_year, value)) }\SpecialCharTok{+}
 \FunctionTok{geom\_line}\NormalTok{() }\SpecialCharTok{+}
 \FunctionTok{geom\_point}\NormalTok{() }\SpecialCharTok{+}
 \FunctionTok{facet\_wrap}\NormalTok{(}\SpecialCharTok{\textasciitilde{}}\NormalTok{name, }\AttributeTok{scales =} \StringTok{"free\_y"}\NormalTok{) }\SpecialCharTok{+}
 \FunctionTok{labs}\NormalTok{(}\AttributeTok{title =} \StringTok{"Journal trend: publications and citations"}\NormalTok{,}
      \AttributeTok{x =} \StringTok{"Year"}\NormalTok{, }\AttributeTok{y =} \ConstantTok{NULL}\NormalTok{)}
\end{Highlighting}
\end{Shaded}

\emph{Key insight:} Median citation counts typically decrease for recent years due to
the citation time window, while publication counts may rise if the journal is
growing.

\section*{Chapter 6 --- Parsing Native Exports}\label{chapter-6-parsing-native-exports}
\addcontentsline{toc}{section}{Chapter 6 --- Parsing Native Exports}

\textbf{1. Export and parse.}
Parse a BibTeX file from Google Scholar using \texttt{bib2df}.

\begin{Shaded}
\begin{Highlighting}[]
\FunctionTok{library}\NormalTok{(bib2df)}
\FunctionTok{library}\NormalTok{(tidyverse)}

\NormalTok{bib }\OtherTok{\textless{}{-}} \FunctionTok{bib2df}\NormalTok{(}\StringTok{"scholar\_export.bib"}\NormalTok{)}
\FunctionTok{ncol}\NormalTok{(bib)}
\FunctionTok{map\_int}\NormalTok{(bib, \textbackslash{}(x) }\FunctionTok{sum}\NormalTok{(}\SpecialCharTok{!}\FunctionTok{is.na}\NormalTok{(x))) }\SpecialCharTok{|\textgreater{}}
 \FunctionTok{enframe}\NormalTok{(}\StringTok{"field"}\NormalTok{, }\StringTok{"populated"}\NormalTok{) }\SpecialCharTok{|\textgreater{}}
 \FunctionTok{arrange}\NormalTok{(}\FunctionTok{desc}\NormalTok{(populated))}
\end{Highlighting}
\end{Shaded}

\emph{Key insight:} Google Scholar BibTeX exports often lack abstracts and keywords,
making them less useful for text-based analyses.

\textbf{2. Format comparison.}
Parse the same records from CSV and RIS formats and compare fields.

\begin{Shaded}
\begin{Highlighting}[]
\FunctionTok{library}\NormalTok{(bibliometrix)}

\NormalTok{csv\_data }\OtherTok{\textless{}{-}} \FunctionTok{read\_csv}\NormalTok{(}\StringTok{"scopus\_export.csv"}\NormalTok{)}
\NormalTok{ris\_data }\OtherTok{\textless{}{-}} \FunctionTok{convert2df}\NormalTok{(}\StringTok{"scopus\_export.ris"}\NormalTok{, }\AttributeTok{dbsource =} \StringTok{"scopus"}\NormalTok{,}
                       \AttributeTok{format =} \StringTok{"plaintext"}\NormalTok{)}

\FunctionTok{cat}\NormalTok{(}\StringTok{"CSV columns:"}\NormalTok{, }\FunctionTok{ncol}\NormalTok{(csv\_data), }\StringTok{"}\SpecialCharTok{\textbackslash{}n}\StringTok{"}\NormalTok{)}
\FunctionTok{cat}\NormalTok{(}\StringTok{"RIS columns:"}\NormalTok{, }\FunctionTok{ncol}\NormalTok{(ris\_data), }\StringTok{"}\SpecialCharTok{\textbackslash{}n}\StringTok{"}\NormalTok{)}

\CommentTok{\# Compare overlapping fields}
\FunctionTok{intersect}\NormalTok{(}\FunctionTok{names}\NormalTok{(csv\_data), }\FunctionTok{names}\NormalTok{(ris\_data))}
\end{Highlighting}
\end{Shaded}

\emph{Key insight:} RIS files typically include structured reference data and
keywords that CSV exports may omit; CSVs are easier to manipulate in R.

\textbf{3. Standardization function.}
Write a reusable function to normalise columns.

\begin{Shaded}
\begin{Highlighting}[]
\NormalTok{standardise\_biblio }\OtherTok{\textless{}{-}} \ControlFlowTok{function}\NormalTok{(df) \{}
 \FunctionTok{tibble}\NormalTok{(}
   \AttributeTok{doi       =}\NormalTok{ df}\SpecialCharTok{$}\NormalTok{DI }\SpecialCharTok{\%||\%} \ConstantTok{NA\_character\_}\NormalTok{,}
   \AttributeTok{title     =}\NormalTok{ df}\SpecialCharTok{$}\NormalTok{TI }\SpecialCharTok{\%||\%}\NormalTok{ df}\SpecialCharTok{$}\NormalTok{Title }\SpecialCharTok{\%||\%} \ConstantTok{NA\_character\_}\NormalTok{,}
   \AttributeTok{authors   =}\NormalTok{ df}\SpecialCharTok{$}\NormalTok{AU }\SpecialCharTok{\%||\%} \ConstantTok{NA\_character\_}\NormalTok{,}
   \AttributeTok{year      =} \FunctionTok{as.integer}\NormalTok{(df}\SpecialCharTok{$}\NormalTok{PY }\SpecialCharTok{\%||\%}\NormalTok{ df}\SpecialCharTok{$}\NormalTok{Year),}
   \AttributeTok{journal   =}\NormalTok{ df}\SpecialCharTok{$}\NormalTok{SO }\SpecialCharTok{\%||\%}\NormalTok{ df}\SpecialCharTok{$}\StringTok{\textasciigrave{}}\AttributeTok{Source title}\StringTok{\textasciigrave{}} \SpecialCharTok{\%||\%} \ConstantTok{NA\_character\_}\NormalTok{,}
   \AttributeTok{cited\_by  =} \FunctionTok{as.integer}\NormalTok{(df}\SpecialCharTok{$}\NormalTok{TC }\SpecialCharTok{\%||\%}\NormalTok{ df}\SpecialCharTok{$}\StringTok{\textasciigrave{}}\AttributeTok{Cited by}\StringTok{\textasciigrave{}} \SpecialCharTok{\%||\%} \DecValTok{0}\DataTypeTok{L}\NormalTok{)}
\NormalTok{ )}
\NormalTok{\}}

\CommentTok{\# Usage:}
\CommentTok{\# standardise\_biblio(ris\_data)}
\end{Highlighting}
\end{Shaded}

\emph{Key insight:} A standardisation layer decouples downstream analysis from the
input format, improving reproducibility.

\section*{Chapter 7 --- Building Reproducible Corpora}\label{chapter-7-building-reproducible-corpora}
\addcontentsline{toc}{section}{Chapter 7 --- Building Reproducible Corpora}

\textbf{1. DOI coverage by year.}
Compute the DOI coverage rate per publication year.

\begin{Shaded}
\begin{Highlighting}[]
\FunctionTok{library}\NormalTok{(tidyverse)}

\NormalTok{works }\SpecialCharTok{|\textgreater{}}
 \FunctionTok{group\_by}\NormalTok{(publication\_year) }\SpecialCharTok{|\textgreater{}}
 \FunctionTok{summarise}\NormalTok{(}\AttributeTok{doi\_rate =} \FunctionTok{mean}\NormalTok{(}\SpecialCharTok{!}\FunctionTok{is.na}\NormalTok{(doi))) }\SpecialCharTok{|\textgreater{}}
 \FunctionTok{ggplot}\NormalTok{(}\FunctionTok{aes}\NormalTok{(publication\_year, doi\_rate)) }\SpecialCharTok{+}
 \FunctionTok{geom\_line}\NormalTok{() }\SpecialCharTok{+}
 \FunctionTok{geom\_point}\NormalTok{() }\SpecialCharTok{+}
 \FunctionTok{scale\_y\_continuous}\NormalTok{(}\AttributeTok{labels =}\NormalTok{ scales}\SpecialCharTok{::}\NormalTok{percent) }\SpecialCharTok{+}
 \FunctionTok{labs}\NormalTok{(}\AttributeTok{title =} \StringTok{"DOI coverage by year"}\NormalTok{, }\AttributeTok{x =} \StringTok{"Year"}\NormalTok{, }\AttributeTok{y =} \StringTok{"Coverage"}\NormalTok{)}
\end{Highlighting}
\end{Shaded}

\emph{Key insight:} DOI coverage tends to improve for recent years as DOI assignment
became standard practice after the mid-2000s.

\textbf{2. Name variants.}
Detect potential author name duplicates.

\begin{Shaded}
\begin{Highlighting}[]
\FunctionTok{library}\NormalTok{(stringdist)}

\NormalTok{author\_names }\OtherTok{\textless{}{-}}\NormalTok{ works }\SpecialCharTok{|\textgreater{}}
 \FunctionTok{unnest}\NormalTok{(author) }\SpecialCharTok{|\textgreater{}}
 \FunctionTok{distinct}\NormalTok{(au\_id, au\_display\_name) }\SpecialCharTok{|\textgreater{}}
 \FunctionTok{pull}\NormalTok{(au\_display\_name)}

\CommentTok{\# Find pairs with low string distance}
\NormalTok{combos }\OtherTok{\textless{}{-}} \FunctionTok{combn}\NormalTok{(author\_names[}\DecValTok{1}\SpecialCharTok{:}\DecValTok{200}\NormalTok{], }\DecValTok{2}\NormalTok{, }\AttributeTok{simplify =} \ConstantTok{FALSE}\NormalTok{)}
\NormalTok{similar }\OtherTok{\textless{}{-}} \FunctionTok{map}\NormalTok{(combos, \textbackslash{}(pair) \{}
\NormalTok{ d }\OtherTok{\textless{}{-}} \FunctionTok{stringdist}\NormalTok{(pair[}\DecValTok{1}\NormalTok{], pair[}\DecValTok{2}\NormalTok{], }\AttributeTok{method =} \StringTok{"jw"}\NormalTok{)}
 \ControlFlowTok{if}\NormalTok{ (d }\SpecialCharTok{\textless{}} \FloatTok{0.15}\NormalTok{) }\FunctionTok{tibble}\NormalTok{(}\AttributeTok{name1 =}\NormalTok{ pair[}\DecValTok{1}\NormalTok{], }\AttributeTok{name2 =}\NormalTok{ pair[}\DecValTok{2}\NormalTok{], }\AttributeTok{dist =}\NormalTok{ d)}
 \ControlFlowTok{else} \ConstantTok{NULL}
\NormalTok{\}) }\SpecialCharTok{|\textgreater{}} \FunctionTok{compact}\NormalTok{() }\SpecialCharTok{|\textgreater{}} \FunctionTok{list\_rbind}\NormalTok{()}

\NormalTok{similar }\SpecialCharTok{|\textgreater{}} \FunctionTok{arrange}\NormalTok{(dist) }\SpecialCharTok{|\textgreater{}} \FunctionTok{head}\NormalTok{(}\DecValTok{20}\NormalTok{)}
\end{Highlighting}
\end{Shaded}

\emph{Key insight:} Common heuristics include Jaro-Winkler distance, initial
matching, and shared co-author overlap. OpenAlex author IDs help but are not
always correct.

\textbf{3. ROR lookup.}
Use the ROR API to resolve institutions.

\begin{Shaded}
\begin{Highlighting}[]
\FunctionTok{library}\NormalTok{(httr2)}

\NormalTok{top\_insts }\OtherTok{\textless{}{-}}\NormalTok{ works }\SpecialCharTok{|\textgreater{}}
 \FunctionTok{unnest}\NormalTok{(authorships) }\SpecialCharTok{|\textgreater{}}
 \FunctionTok{unnest}\NormalTok{(institutions) }\SpecialCharTok{|\textgreater{}}
 \FunctionTok{count}\NormalTok{(display\_name, }\AttributeTok{sort =} \ConstantTok{TRUE}\NormalTok{) }\SpecialCharTok{|\textgreater{}}
 \FunctionTok{head}\NormalTok{(}\DecValTok{5}\NormalTok{)}

\NormalTok{ror\_lookup }\OtherTok{\textless{}{-}} \ControlFlowTok{function}\NormalTok{(name) \{}
\NormalTok{ resp }\OtherTok{\textless{}{-}} \FunctionTok{request}\NormalTok{(}\StringTok{"https://api.ror.org/organizations"}\NormalTok{) }\SpecialCharTok{|\textgreater{}}
   \FunctionTok{req\_url\_query}\NormalTok{(}\AttributeTok{query =}\NormalTok{ name) }\SpecialCharTok{|\textgreater{}}
   \FunctionTok{req\_perform}\NormalTok{() }\SpecialCharTok{|\textgreater{}}
   \FunctionTok{resp\_body\_json}\NormalTok{()}
 \ControlFlowTok{if}\NormalTok{ (}\FunctionTok{length}\NormalTok{(resp}\SpecialCharTok{$}\NormalTok{items) }\SpecialCharTok{\textgreater{}} \DecValTok{0}\NormalTok{)}
   \FunctionTok{tibble}\NormalTok{(}\AttributeTok{name =}\NormalTok{ name, }\AttributeTok{ror\_id =}\NormalTok{ resp}\SpecialCharTok{$}\NormalTok{items[[}\DecValTok{1}\NormalTok{]]}\SpecialCharTok{$}\NormalTok{id,}
          \AttributeTok{ror\_name =}\NormalTok{ resp}\SpecialCharTok{$}\NormalTok{items[[}\DecValTok{1}\NormalTok{]]}\SpecialCharTok{$}\NormalTok{name)}
 \ControlFlowTok{else}
   \FunctionTok{tibble}\NormalTok{(}\AttributeTok{name =}\NormalTok{ name, }\AttributeTok{ror\_id =} \ConstantTok{NA}\NormalTok{, }\AttributeTok{ror\_name =} \ConstantTok{NA}\NormalTok{)}
\NormalTok{\}}

\FunctionTok{map}\NormalTok{(top\_insts}\SpecialCharTok{$}\NormalTok{display\_name, ror\_lookup) }\SpecialCharTok{|\textgreater{}} \FunctionTok{list\_rbind}\NormalTok{()}
\end{Highlighting}
\end{Shaded}

\emph{Key insight:} ROR provides stable, disambiguated identifiers; matching by name
alone can produce false positives for common institutional names.

\section*{Chapter 8 --- Working with Full Text}\label{chapter-8-working-with-full-text}
\addcontentsline{toc}{section}{Chapter 8 --- Working with Full Text}

\textbf{1. Extract and compare.}
Download a PDF, extract text, and compare the abstract.

\begin{Shaded}
\begin{Highlighting}[]
\FunctionTok{library}\NormalTok{(pdftools)}
\FunctionTok{library}\NormalTok{(tidyverse)}

\NormalTok{pdf\_text\_raw }\OtherTok{\textless{}{-}} \FunctionTok{pdf\_text}\NormalTok{(}\StringTok{"sample\_paper.pdf"}\NormalTok{)}
\NormalTok{extracted\_abstract }\OtherTok{\textless{}{-}} \FunctionTok{str\_extract}\NormalTok{(}
 \FunctionTok{paste}\NormalTok{(pdf\_text\_raw, }\AttributeTok{collapse =} \StringTok{"}\SpecialCharTok{\textbackslash{}n}\StringTok{"}\NormalTok{),}
 \StringTok{"(?i)abstract}\SpecialCharTok{\textbackslash{}\textbackslash{}}\StringTok{s*}\SpecialCharTok{\textbackslash{}\textbackslash{}}\StringTok{.?}\SpecialCharTok{\textbackslash{}\textbackslash{}}\StringTok{s*(.+?)(?=}\SpecialCharTok{\textbackslash{}\textbackslash{}}\StringTok{n}\SpecialCharTok{\textbackslash{}\textbackslash{}}\StringTok{s*(introduction|keywords|1}\SpecialCharTok{\textbackslash{}\textbackslash{}}\StringTok{.))"}\NormalTok{,}
\NormalTok{)}

\CommentTok{\# Compare with OpenAlex abstract}
\NormalTok{oa\_abstract }\OtherTok{\textless{}{-}}\NormalTok{ works }\SpecialCharTok{|\textgreater{}} \FunctionTok{filter}\NormalTok{(doi }\SpecialCharTok{==} \StringTok{"10.xxxx/example"}\NormalTok{) }\SpecialCharTok{|\textgreater{}} \FunctionTok{pull}\NormalTok{(ab)}

\FunctionTok{cat}\NormalTok{(}\StringTok{"Extracted:}\SpecialCharTok{\textbackslash{}n}\StringTok{"}\NormalTok{, }\FunctionTok{str\_trunc}\NormalTok{(extracted\_abstract, }\DecValTok{300}\NormalTok{), }\StringTok{"}\SpecialCharTok{\textbackslash{}n}\StringTok{"}\NormalTok{)}
\FunctionTok{cat}\NormalTok{(}\StringTok{"OpenAlex:}\SpecialCharTok{\textbackslash{}n}\StringTok{"}\NormalTok{, }\FunctionTok{str\_trunc}\NormalTok{(oa\_abstract, }\DecValTok{300}\NormalTok{), }\StringTok{"}\SpecialCharTok{\textbackslash{}n}\StringTok{"}\NormalTok{)}
\end{Highlighting}
\end{Shaded}

\emph{Key insight:} PDF extraction often introduces line-break artefacts and header/
footer noise. The OpenAlex abstract (from publisher metadata) is typically
cleaner.

\textbf{2. Section detection.}
Write a function to split extracted text into sections.

\begin{Shaded}
\begin{Highlighting}[]
\NormalTok{split\_sections }\OtherTok{\textless{}{-}} \ControlFlowTok{function}\NormalTok{(text) \{}
\NormalTok{ lines }\OtherTok{\textless{}{-}} \FunctionTok{str\_split}\NormalTok{(text, }\StringTok{"}\SpecialCharTok{\textbackslash{}n}\StringTok{"}\NormalTok{)[[}\DecValTok{1}\NormalTok{]]}
 \CommentTok{\# Match numbered headings like "1. Introduction" or "2 Methods"}
\NormalTok{ heading\_idx }\OtherTok{\textless{}{-}} \FunctionTok{str\_which}\NormalTok{(lines, }\StringTok{"\^{}}\SpecialCharTok{\textbackslash{}\textbackslash{}}\StringTok{s*}\SpecialCharTok{\textbackslash{}\textbackslash{}}\StringTok{d+}\SpecialCharTok{\textbackslash{}\textbackslash{}}\StringTok{.?}\SpecialCharTok{\textbackslash{}\textbackslash{}}\StringTok{s+[A{-}Z]"}\NormalTok{)}

 \ControlFlowTok{if}\NormalTok{ (}\FunctionTok{length}\NormalTok{(heading\_idx) }\SpecialCharTok{\textless{}} \DecValTok{2}\NormalTok{) }\FunctionTok{return}\NormalTok{(}\FunctionTok{list}\NormalTok{(}\AttributeTok{full\_text =}\NormalTok{ text))}

\NormalTok{ sections }\OtherTok{\textless{}{-}} \FunctionTok{map2}\NormalTok{(heading\_idx, }\FunctionTok{c}\NormalTok{(heading\_idx[}\SpecialCharTok{{-}}\DecValTok{1}\NormalTok{] }\SpecialCharTok{{-}} \DecValTok{1}\NormalTok{, }\FunctionTok{length}\NormalTok{(lines)),}
\NormalTok{                  \textbackslash{}(start, end) \{}
   \FunctionTok{list}\NormalTok{(}
     \AttributeTok{heading =} \FunctionTok{str\_trim}\NormalTok{(lines[start]),}
     \AttributeTok{text =} \FunctionTok{paste}\NormalTok{(lines[(start }\SpecialCharTok{+} \DecValTok{1}\NormalTok{)}\SpecialCharTok{:}\NormalTok{end], }\AttributeTok{collapse =} \StringTok{"}\SpecialCharTok{\textbackslash{}n}\StringTok{"}\NormalTok{)}
\NormalTok{   )}
\NormalTok{ \})}
 \FunctionTok{set\_names}\NormalTok{(}
   \FunctionTok{map\_chr}\NormalTok{(sections, }\StringTok{"text"}\NormalTok{),}
   \FunctionTok{map\_chr}\NormalTok{(sections, }\StringTok{"heading"}\NormalTok{)}
\NormalTok{ )}
\NormalTok{\}}

\NormalTok{sections }\OtherTok{\textless{}{-}} \FunctionTok{split\_sections}\NormalTok{(}\FunctionTok{paste}\NormalTok{(pdf\_text\_raw, }\AttributeTok{collapse =} \StringTok{"}\SpecialCharTok{\textbackslash{}n}\StringTok{"}\NormalTok{))}
\FunctionTok{names}\NormalTok{(sections)}
\end{Highlighting}
\end{Shaded}

\emph{Key insight:} Heading detection is fragile; two-column layouts, supplementary
materials, and non-standard numbering all break simple regex patterns.

\textbf{3. Table extraction.}
Extract a table from a PDF using \texttt{tabulizer}.

\begin{Shaded}
\begin{Highlighting}[]
\CommentTok{\# tabulapdf needs Java (rJava); shown for reference, not executed.}
\CommentTok{\# Install: install.packages("tabulapdf"); requires a working Java runtime.}

\NormalTok{tables }\OtherTok{\textless{}{-}} \FunctionTok{extract\_tables}\NormalTok{(}\StringTok{"sample\_paper.pdf"}\NormalTok{, }\AttributeTok{pages =} \DecValTok{3}\NormalTok{)}
\NormalTok{tbl }\OtherTok{\textless{}{-}} \FunctionTok{as\_tibble}\NormalTok{(tables[[}\DecValTok{1}\NormalTok{]], }\AttributeTok{.name\_repair =} \StringTok{"unique"}\NormalTok{)}

\CommentTok{\# Typical cleaning steps:}
\NormalTok{tbl }\OtherTok{\textless{}{-}}\NormalTok{ tbl }\SpecialCharTok{|\textgreater{}}
 \FunctionTok{row\_to\_names}\NormalTok{(}\AttributeTok{row\_number =} \DecValTok{1}\NormalTok{) }\SpecialCharTok{|\textgreater{}}   \CommentTok{\# janitor::row\_to\_names}
 \FunctionTok{mutate}\NormalTok{(}\FunctionTok{across}\NormalTok{(}\FunctionTok{where}\NormalTok{(is.character), str\_trim)) }\SpecialCharTok{|\textgreater{}}
 \FunctionTok{filter}\NormalTok{(}\FunctionTok{if\_any}\NormalTok{(}\FunctionTok{everything}\NormalTok{(), \textbackslash{}(x) x }\SpecialCharTok{!=} \StringTok{""}\NormalTok{))}
\end{Highlighting}
\end{Shaded}

\emph{Key insight:} Extracted tables almost always need header row promotion,
whitespace trimming, and removal of empty or merged-cell artefact rows.

\section*{Chapter 9 --- Descriptive Bibliometrics}\label{chapter-9-descriptive-bibliometrics}
\addcontentsline{toc}{section}{Chapter 9 --- Descriptive Bibliometrics}

\textbf{1. Lotka's law.}
Fit Lotka's law using \texttt{bibliometrix}.

\begin{Shaded}
\begin{Highlighting}[]
\FunctionTok{library}\NormalTok{(bibliometrix)}
\FunctionTok{library}\NormalTok{(tidyverse)}

\NormalTok{bib\_df }\OtherTok{\textless{}{-}} \FunctionTok{convert2df}\NormalTok{(works)  }\CommentTok{\# or use the already{-}converted object}
\NormalTok{lotka\_result }\OtherTok{\textless{}{-}} \FunctionTok{lotka}\NormalTok{(}\FunctionTok{biblioAnalysis}\NormalTok{(bib\_df))}

\CommentTok{\# Plot observed vs expected}
\FunctionTok{plot}\NormalTok{(lotka\_result)}
\end{Highlighting}
\end{Shaded}

\emph{Key insight:} If the Kolmogorov-Smirnov test p-value is above 0.05, the
observed distribution is consistent with the theoretical inverse-square law.

\textbf{2. Bradford's law.}
Identify core, mid, and peripheral journal zones.

\begin{Shaded}
\begin{Highlighting}[]
\FunctionTok{library}\NormalTok{(openalexR)}
\FunctionTok{library}\NormalTok{(bibliometrix)}

\NormalTok{ml\_works }\OtherTok{\textless{}{-}} \FunctionTok{fetch\_openalex}\NormalTok{(}\AttributeTok{entity =} \StringTok{"works"}\NormalTok{, }\AttributeTok{search =} \StringTok{"machine learning"}\NormalTok{,}
                     \AttributeTok{from\_publication\_date =} \StringTok{"2020{-}01{-}01"}\NormalTok{,}
                     \AttributeTok{to\_publication\_date =} \StringTok{"2023{-}12{-}31"}\NormalTok{) }\SpecialCharTok{|\textgreater{}}
 \FunctionTok{slice\_sample}\NormalTok{(}\AttributeTok{n =} \DecValTok{500}\NormalTok{)}

\NormalTok{bib\_df }\OtherTok{\textless{}{-}} \FunctionTok{convert2df}\NormalTok{(ml\_works)}
\NormalTok{brad }\OtherTok{\textless{}{-}} \FunctionTok{bradford}\NormalTok{(}\FunctionTok{biblioAnalysis}\NormalTok{(bib\_df))}
\NormalTok{brad}\SpecialCharTok{$}\NormalTok{table }\SpecialCharTok{|\textgreater{}} \FunctionTok{head}\NormalTok{(}\DecValTok{20}\NormalTok{)}
\end{Highlighting}
\end{Shaded}

\emph{Key insight:} The core zone typically contains very few journals that account
for a disproportionate share of articles --- Bradford's law in action.

\textbf{3. Three-field exploration.}
Use country instead of author name on the left field.

\begin{Shaded}
\begin{Highlighting}[]
\FunctionTok{library}\NormalTok{(bibliometrix)}

\FunctionTok{threeFieldsPlot}\NormalTok{(bib\_df, }\AttributeTok{fields =} \FunctionTok{c}\NormalTok{(}\StringTok{"AU\_CO"}\NormalTok{, }\StringTok{"SO"}\NormalTok{, }\StringTok{"DE"}\NormalTok{),}
                \AttributeTok{n =} \FunctionTok{c}\NormalTok{(}\DecValTok{10}\NormalTok{, }\DecValTok{10}\NormalTok{, }\DecValTok{10}\NormalTok{))}
\end{Highlighting}
\end{Shaded}

\emph{Key insight:} Country-level Sankey diagrams reveal geographic specialisation
patterns: certain countries cluster around specific journals and keywords.

\textbf{4. Annual growth rate.}
Compute the compound annual growth rate (CAGR).

\begin{Shaded}
\begin{Highlighting}[]
\NormalTok{annual }\OtherTok{\textless{}{-}}\NormalTok{ works }\SpecialCharTok{|\textgreater{}}
 \FunctionTok{count}\NormalTok{(publication\_year) }\SpecialCharTok{|\textgreater{}}
 \FunctionTok{filter}\NormalTok{(publication\_year }\SpecialCharTok{\textgreater{}=} \DecValTok{2010}\NormalTok{, publication\_year }\SpecialCharTok{\textless{}=} \DecValTok{2023}\NormalTok{)}

\NormalTok{first\_year }\OtherTok{\textless{}{-}}\NormalTok{ annual}\SpecialCharTok{$}\NormalTok{n[}\DecValTok{1}\NormalTok{]}
\NormalTok{last\_year  }\OtherTok{\textless{}{-}}\NormalTok{ annual}\SpecialCharTok{$}\NormalTok{n[}\FunctionTok{nrow}\NormalTok{(annual)]}
\NormalTok{n\_years    }\OtherTok{\textless{}{-}} \FunctionTok{nrow}\NormalTok{(annual) }\SpecialCharTok{{-}} \DecValTok{1}

\NormalTok{cagr }\OtherTok{\textless{}{-}}\NormalTok{ (last\_year }\SpecialCharTok{/}\NormalTok{ first\_year)}\SpecialCharTok{\^{}}\NormalTok{(}\DecValTok{1} \SpecialCharTok{/}\NormalTok{ n\_years) }\SpecialCharTok{{-}} \DecValTok{1}
\FunctionTok{cat}\NormalTok{(}\FunctionTok{sprintf}\NormalTok{(}\StringTok{"CAGR: \%.1f\%\%}\SpecialCharTok{\textbackslash{}n}\StringTok{"}\NormalTok{, cagr }\SpecialCharTok{*} \DecValTok{100}\NormalTok{))}
\end{Highlighting}
\end{Shaded}

\emph{Key insight:} A positive CAGR indicates the field is growing; compare against
the global average (\textasciitilde4-5\% for all scientific literature) to assess relative
growth.

\section*{Chapter 10 --- Productivity and Impact Indicators}\label{chapter-10-productivity-and-impact-indicators}
\addcontentsline{toc}{section}{Chapter 10 --- Productivity and Impact Indicators}

\textbf{1. Career h-index.}
Compute h-index, g-index, and m-quotient for a specific author.

\begin{Shaded}
\begin{Highlighting}[]
\FunctionTok{library}\NormalTok{(openalexR)}
\FunctionTok{library}\NormalTok{(tidyverse)}

\NormalTok{author\_works }\OtherTok{\textless{}{-}} \FunctionTok{fetch\_openalex}\NormalTok{(}
 \AttributeTok{entity =} \StringTok{"works"}\NormalTok{,}
 \AttributeTok{author.id =} \StringTok{"A5023888391"}\NormalTok{,  }\CommentTok{\# replace with target author}
 \AttributeTok{per\_page =} \DecValTok{200}
\NormalTok{)}

\NormalTok{cites }\OtherTok{\textless{}{-}} \FunctionTok{sort}\NormalTok{(author\_works}\SpecialCharTok{$}\NormalTok{cited\_by\_count, }\AttributeTok{decreasing =} \ConstantTok{TRUE}\NormalTok{)}

\CommentTok{\# h{-}index}
\NormalTok{h }\OtherTok{\textless{}{-}} \FunctionTok{sum}\NormalTok{(cites }\SpecialCharTok{\textgreater{}=} \FunctionTok{seq\_along}\NormalTok{(cites))}

\CommentTok{\# g{-}index}
\NormalTok{cum\_cites }\OtherTok{\textless{}{-}} \FunctionTok{cumsum}\NormalTok{(cites)}
\NormalTok{g\_candidates }\OtherTok{\textless{}{-}} \FunctionTok{which}\NormalTok{(cum\_cites }\SpecialCharTok{\textgreater{}=} \FunctionTok{seq\_along}\NormalTok{(cites)}\SpecialCharTok{\^{}}\DecValTok{2}\NormalTok{)}
\NormalTok{g }\OtherTok{\textless{}{-}} \ControlFlowTok{if}\NormalTok{ (}\FunctionTok{length}\NormalTok{(g\_candidates)) }\FunctionTok{max}\NormalTok{(g\_candidates) }\ControlFlowTok{else} \DecValTok{0}\DataTypeTok{L}

\CommentTok{\# m{-}quotient}
\NormalTok{career\_years }\OtherTok{\textless{}{-}} \FunctionTok{max}\NormalTok{(lubridate}\SpecialCharTok{::}\FunctionTok{year}\NormalTok{(author\_works}\SpecialCharTok{$}\NormalTok{publication\_date)) }\SpecialCharTok{{-}}
               \FunctionTok{min}\NormalTok{(lubridate}\SpecialCharTok{::}\FunctionTok{year}\NormalTok{(author\_works}\SpecialCharTok{$}\NormalTok{publication\_date)) }\SpecialCharTok{+} \DecValTok{1}
\NormalTok{m }\OtherTok{\textless{}{-}}\NormalTok{ h }\SpecialCharTok{/}\NormalTok{ career\_years}

\FunctionTok{tibble}\NormalTok{(}\AttributeTok{h\_index =}\NormalTok{ h, }\AttributeTok{g\_index =}\NormalTok{ g, }\AttributeTok{m\_quotient =} \FunctionTok{round}\NormalTok{(m, }\DecValTok{2}\NormalTok{))}
\end{Highlighting}
\end{Shaded}

\emph{Key insight:} The h-index rewards consistency, the g-index rewards highly cited
papers, and the m-quotient adjusts for career length.

\textbf{2. Self-citation sensitivity.}
Pseudocode for removing self-citations from the h-index.

\begin{Shaded}
\begin{Highlighting}[]
\CommentTok{\# Pseudocode:}
\CommentTok{\# 1. For each paper by the author, fetch its citing works}
\CommentTok{\# 2. Check if any citing work shares the same author ID}
\CommentTok{\# 3. Count only non{-}self{-}citations}
\CommentTok{\# 4. Recompute h{-}index with adjusted citation counts}
\CommentTok{\#}
\CommentTok{\# Implementation sketch (expensive API calls):}
\CommentTok{\# author\_id \textless{}{-} "A5023888391"}
\CommentTok{\# adjusted\_cites \textless{}{-} map\_int(author\_works$id, \textbackslash{}(work\_id) \{}
\CommentTok{\#   citing \textless{}{-} fetch\_openalex(entity = "works", cites = work\_id)}
\CommentTok{\#   citing |\textgreater{}}
\CommentTok{\#     unnest(author) |\textgreater{}}
\CommentTok{\#     filter(au\_id != author\_id) |\textgreater{}}
\CommentTok{\#     distinct(id) |\textgreater{}}
\CommentTok{\#     nrow()}
\CommentTok{\# \})}
\CommentTok{\# h\_adjusted \textless{}{-} sum(sort(adjusted\_cites, decreasing = TRUE) \textgreater{}=}
\CommentTok{\#                   seq\_along(adjusted\_cites))}
\end{Highlighting}
\end{Shaded}

\emph{Key insight:} Removing citations cannot increase the h-index. Its actual change
requires verified author matches and citation records; no universal change is assumed.

\textbf{3. Empirical resampling of sample ratios.}
Hold the empirical year reference fixed and resample its observed ratios.

\begin{Shaded}
\begin{Highlighting}[]
\FunctionTok{set.seed}\NormalTok{(}\DecValTok{20260509}\NormalTok{)}
\NormalTok{ratios }\OtherTok{\textless{}{-}}\NormalTok{ works\_normalized}\SpecialCharTok{$}\NormalTok{ncs[}\FunctionTok{is.finite}\NormalTok{(works\_normalized}\SpecialCharTok{$}\NormalTok{ncs)]}
\FunctionTok{stopifnot}\NormalTok{(}\FunctionTok{length}\NormalTok{(ratios) }\SpecialCharTok{\textgreater{}} \DecValTok{1}\DataTypeTok{L}\NormalTok{)}
\NormalTok{ratio\_boot }\OtherTok{\textless{}{-}} \FunctionTok{replicate}\NormalTok{(}\DecValTok{1000}\NormalTok{, }\FunctionTok{mean}\NormalTok{(}\FunctionTok{sample}\NormalTok{(ratios, }\AttributeTok{replace =} \ConstantTok{TRUE}\NormalTok{)))}
\FunctionTok{quantile}\NormalTok{(ratio\_boot, }\FunctionTok{c}\NormalTok{(}\FloatTok{0.025}\NormalTok{, }\FloatTok{0.975}\NormalTok{))}
\end{Highlighting}
\end{Shaded}

\emph{Key insight:} This interval describes resampling variation conditional on the
chosen empirical reference. It ignores uncertainty in that reference and does
not establish world-normalized impact or a population confidence interval.

\textbf{4. Normalization and ranks.}
Compare the within-sample mean ratio and percentile selection share by year.

\begin{Shaded}
\begin{Highlighting}[]
\NormalTok{yearly\_indicators }\OtherTok{\textless{}{-}}\NormalTok{ works\_with\_pct }\SpecialCharTok{|\textgreater{}}
 \FunctionTok{group\_by}\NormalTok{(year) }\SpecialCharTok{|\textgreater{}}
 \FunctionTok{summarise}\NormalTok{(}
   \AttributeTok{mean\_sample\_ratio =} \FunctionTok{mean}\NormalTok{(ncs, }\AttributeTok{na.rm =} \ConstantTok{TRUE}\NormalTok{),}
   \AttributeTok{ranked\_works =} \FunctionTok{sum}\NormalTok{(}\SpecialCharTok{!}\FunctionTok{is.na}\NormalTok{(is\_top10)),}
   \AttributeTok{upper\_percentile\_share =} \FunctionTok{mean}\NormalTok{(is\_top10, }\AttributeTok{na.rm =} \ConstantTok{TRUE}\NormalTok{),}
   \AttributeTok{.groups =} \StringTok{"drop"}
\NormalTok{ )}
\NormalTok{yearly\_indicators}
\end{Highlighting}
\end{Shaded}

\emph{Key insight:} The mean ratio is one for each usable year by construction.
Ties and small groups affect the percentile share. These sample summaries do
not rank field-normalized impact; that requires a justified external reference.

\section*{Chapter 11 --- Journal Metrics}\label{chapter-11-journal-metrics}
\addcontentsline{toc}{section}{Chapter 11 --- Journal Metrics}

\textbf{1. Cohort sensitivity.}
Use the verified source table and helper from Chapter 11. This code performs a
new provider request unless that exact query is already cached.

\begin{Shaded}
\begin{Highlighting}[]
\NormalTok{wider\_cohort }\OtherTok{\textless{}{-}} \FunctionTok{fetch\_journal\_works}\NormalTok{(journals}\SpecialCharTok{$}\NormalTok{source\_id[}\DecValTok{1}\NormalTok{], }\DecValTok{2019}\NormalTok{, }\DecValTok{2022}\NormalTok{)}
\FunctionTok{summarise\_journal\_records}\NormalTok{(wider\_cohort)}
\NormalTok{journal\_summary }\SpecialCharTok{|\textgreater{}} \FunctionTok{filter}\NormalTok{(source\_id }\SpecialCharTok{==}\NormalTok{ journals}\SpecialCharTok{$}\NormalTok{source\_id[}\DecValTok{1}\NormalTok{])}
\end{Highlighting}
\end{Shaded}

The summaries concern cumulative citations at retrieval. A wider publication
cohort changes composition and citation exposure; the result need not move in
one direction. CiteScore requires citations received within its specified
four-year interval and eligible document types in both numerator and denominator.
Changing a publication-date filter alone supplies neither requirement.

\textbf{2. Descriptive association.}

\begin{Shaded}
\begin{Highlighting}[]
\NormalTok{complete\_pairs }\OtherTok{\textless{}{-}}\NormalTok{ journal\_summary }\SpecialCharTok{|\textgreater{}}
  \FunctionTok{select}\NormalTok{(mean\_cumulative\_citations, median\_recorded\_references) }\SpecialCharTok{|\textgreater{}}
  \FunctionTok{drop\_na}\NormalTok{()}
\NormalTok{rho }\OtherTok{\textless{}{-}} \ControlFlowTok{if}\NormalTok{ (}\FunctionTok{nrow}\NormalTok{(complete\_pairs) }\SpecialCharTok{\textgreater{}=} \DecValTok{3} \SpecialCharTok{\&\&}
           \FunctionTok{n\_distinct}\NormalTok{(complete\_pairs}\SpecialCharTok{$}\NormalTok{mean\_cumulative\_citations) }\SpecialCharTok{\textgreater{}} \DecValTok{1} \SpecialCharTok{\&\&}
           \FunctionTok{n\_distinct}\NormalTok{(complete\_pairs}\SpecialCharTok{$}\NormalTok{median\_recorded\_references) }\SpecialCharTok{\textgreater{}} \DecValTok{1}\NormalTok{) \{}
  \FunctionTok{cor}\NormalTok{(complete\_pairs}\SpecialCharTok{$}\NormalTok{mean\_cumulative\_citations,}
\NormalTok{      complete\_pairs}\SpecialCharTok{$}\NormalTok{median\_recorded\_references, }\AttributeTok{method =} \StringTok{"spearman"}\NormalTok{)}
\NormalTok{\} }\ControlFlowTok{else} \ConstantTok{NA\_real\_}
\NormalTok{rho}
\end{Highlighting}
\end{Shaded}

The coefficient describes the complete journal pairs only; insufficient or
constant observations leave it unavailable. Its sign and magnitude depend on
the retrieved data. It neither measures citation potential on the citing side
nor demonstrates field normalisation or agreement with SNIP.

\textbf{3. Missing-data denominator.}

\begin{Shaded}
\begin{Highlighting}[]
\NormalTok{invented\_records }\OtherTok{\textless{}{-}} \FunctionTok{tibble}\NormalTok{(}
  \AttributeTok{id =} \FunctionTok{c}\NormalTok{(}\StringTok{"example{-}1"}\NormalTok{, }\StringTok{"example{-}2"}\NormalTok{, }\StringTok{"example{-}3"}\NormalTok{),}
  \AttributeTok{cited\_by\_count =} \FunctionTok{c}\NormalTok{(}\DecValTok{0}\NormalTok{, }\DecValTok{20}\NormalTok{, }\ConstantTok{NA\_real\_}\NormalTok{),}
  \AttributeTok{referenced\_works\_count =} \FunctionTok{c}\NormalTok{(}\DecValTok{10}\NormalTok{, }\ConstantTok{NA\_real\_}\NormalTok{, }\DecValTok{30}\NormalTok{)}
\NormalTok{)}
\FunctionTok{summarise\_journal\_records}\NormalTok{(invented\_records)}
\FunctionTok{summarise\_journal\_records}\NormalTok{(}\FunctionTok{mutate}\NormalTok{(invented\_records, }\AttributeTok{cited\_by\_count =} \ConstantTok{NA\_real\_}\NormalTok{))}
\end{Highlighting}
\end{Shaded}

There are three records but two observed citation counts. Their sum is 20 and
the observed-count mean is 10. Dividing 20 by all three records would give about
6.67 and implicitly count the missing value as zero. When all citation counts
are missing, the observed denominator is zero and the sum, mean and median are
reported as missing. A zero observed count is retained as a valid observation.

\textbf{4. Self-citation design.}

Define the publication cohort and citation-receipt window first. Obtain the
citing--cited work links within that scope and identify each citing work's journal
using validated source IDs. Count each citation edge once. A journal self-citation
share can then use within-journal citation edges in the numerator and all observed
citation edges to the cohort in the denominator, with missing source identities
reported separately. Distinguish a genuinely empty eligible edge set from a failed
request; a zero denominator makes the share undefined.

An average of per-paper self-citation shares weights papers equally and generally
differs from the edge-level share. Neither current cumulative counts nor references
recorded for the cited papers identify the journals that cited them. No fixed
self-citation threshold or empirical result follows from this chapter's data.

\section*{Chapter 12 --- Author-Level Analysis}\label{chapter-12-author-level-analysis}
\addcontentsline{toc}{section}{Chapter 12 --- Author-Level Analysis}

\textbf{1. Profile comparison.}
Compare two researchers on multiple indicators.

\begin{Shaded}
\begin{Highlighting}[]
\FunctionTok{library}\NormalTok{(openalexR)}
\FunctionTok{library}\NormalTok{(tidyverse)}

\NormalTok{ids }\OtherTok{\textless{}{-}} \FunctionTok{c}\NormalTok{(}\StringTok{"A5023888391"}\NormalTok{, }\StringTok{"A5048491430"}\NormalTok{)  }\CommentTok{\# two author IDs}

\NormalTok{profiles }\OtherTok{\textless{}{-}} \FunctionTok{map}\NormalTok{(ids, \textbackslash{}(aid) \{}
\NormalTok{ w }\OtherTok{\textless{}{-}} \FunctionTok{fetch\_openalex}\NormalTok{(}\AttributeTok{entity =} \StringTok{"works"}\NormalTok{, }\AttributeTok{author.id =}\NormalTok{ aid)}
\NormalTok{ cites }\OtherTok{\textless{}{-}} \FunctionTok{sort}\NormalTok{(w}\SpecialCharTok{$}\NormalTok{cited\_by\_count, }\AttributeTok{decreasing =} \ConstantTok{TRUE}\NormalTok{)}
 \FunctionTok{tibble}\NormalTok{(}
   \AttributeTok{author\_id =}\NormalTok{ aid,}
   \AttributeTok{n\_papers  =} \FunctionTok{nrow}\NormalTok{(w),}
   \AttributeTok{h\_index   =} \FunctionTok{sum}\NormalTok{(cites }\SpecialCharTok{\textgreater{}=} \FunctionTok{seq\_along}\NormalTok{(cites)),}
   \AttributeTok{mean\_cites =} \FunctionTok{mean}\NormalTok{(w}\SpecialCharTok{$}\NormalTok{cited\_by\_count)}
\NormalTok{ )}
\NormalTok{\}) }\SpecialCharTok{|\textgreater{}} \FunctionTok{list\_rbind}\NormalTok{()}

\NormalTok{profiles}
\end{Highlighting}
\end{Shaded}

\emph{Key insight:} An author with many papers but a modest h-index publishes
frequently but with less consistent impact; high mean citations suggest a few
highly cited papers.

\textbf{2. ORCID validation.}
Compare works fetched by OpenAlex ID vs.~ORCID filtering.

\begin{Shaded}
\begin{Highlighting}[]
\NormalTok{oa\_by\_id }\OtherTok{\textless{}{-}} \FunctionTok{fetch\_openalex}\NormalTok{(}\AttributeTok{entity =} \StringTok{"works"}\NormalTok{,}
                     \AttributeTok{author.id =} \StringTok{"A5023888391"}\NormalTok{)}

\NormalTok{oa\_by\_orcid }\OtherTok{\textless{}{-}} \FunctionTok{fetch\_openalex}\NormalTok{(}\AttributeTok{entity =} \StringTok{"works"}\NormalTok{,}
                        \AttributeTok{author.orcid =} \StringTok{"0000{-}0001{-}2345{-}6789"}\NormalTok{)}

\FunctionTok{cat}\NormalTok{(}\StringTok{"By ID:"}\NormalTok{, }\FunctionTok{nrow}\NormalTok{(oa\_by\_id), }\StringTok{"works}\SpecialCharTok{\textbackslash{}n}\StringTok{"}\NormalTok{)}
\FunctionTok{cat}\NormalTok{(}\StringTok{"By ORCID:"}\NormalTok{, }\FunctionTok{nrow}\NormalTok{(oa\_by\_orcid), }\StringTok{"works}\SpecialCharTok{\textbackslash{}n}\StringTok{"}\NormalTok{)}

\CommentTok{\# Find discrepancies}
\NormalTok{only\_id }\OtherTok{\textless{}{-}} \FunctionTok{setdiff}\NormalTok{(oa\_by\_id}\SpecialCharTok{$}\NormalTok{id, oa\_by\_orcid}\SpecialCharTok{$}\NormalTok{id)}
\NormalTok{only\_orcid }\OtherTok{\textless{}{-}} \FunctionTok{setdiff}\NormalTok{(oa\_by\_orcid}\SpecialCharTok{$}\NormalTok{id, oa\_by\_id}\SpecialCharTok{$}\NormalTok{id)}
\FunctionTok{cat}\NormalTok{(}\StringTok{"Only in ID set:"}\NormalTok{, }\FunctionTok{length}\NormalTok{(only\_id), }\StringTok{"}\SpecialCharTok{\textbackslash{}n}\StringTok{"}\NormalTok{)}
\FunctionTok{cat}\NormalTok{(}\StringTok{"Only in ORCID set:"}\NormalTok{, }\FunctionTok{length}\NormalTok{(only\_orcid), }\StringTok{"}\SpecialCharTok{\textbackslash{}n}\StringTok{"}\NormalTok{)}
\end{Highlighting}
\end{Shaded}

\emph{Key insight:} Discrepancies arise from incomplete ORCID claiming, name
disambiguation errors, and delayed metadata updates.

\textbf{3. Name variant detection.}
Search for a common name and assess disambiguation.

\begin{Shaded}
\begin{Highlighting}[]
\NormalTok{wang\_wei }\OtherTok{\textless{}{-}} \FunctionTok{fetch\_openalex}\NormalTok{(}\AttributeTok{entity =} \StringTok{"authors"}\NormalTok{,}
                     \AttributeTok{search =} \StringTok{"Wang Wei"}\NormalTok{)}

\FunctionTok{nrow}\NormalTok{(wang\_wei)  }\CommentTok{\# many entities}

\NormalTok{wang\_wei }\SpecialCharTok{|\textgreater{}}
 \FunctionTok{select}\NormalTok{(id, display\_name, works\_count, cited\_by\_count) }\SpecialCharTok{|\textgreater{}}
 \FunctionTok{arrange}\NormalTok{(}\FunctionTok{desc}\NormalTok{(works\_count)) }\SpecialCharTok{|\textgreater{}}
 \FunctionTok{head}\NormalTok{(}\DecValTok{20}\NormalTok{)}

\CommentTok{\# Heuristics for identifying the correct author:}
\CommentTok{\# {-} Match institutional affiliation}
\CommentTok{\# {-} Check ORCID linkage}
\CommentTok{\# {-} Verify co{-}author overlap with known collaborators}
\CommentTok{\# {-} Inspect topic/concept alignment}
\end{Highlighting}
\end{Shaded}

\emph{Key insight:} Common names generate hundreds of author entities; institutional
affiliation and ORCID are the most reliable disambiguation features.

\textbf{4. Fractional productivity.}
Recompute productivity using fractional counting.

\begin{Shaded}
\begin{Highlighting}[]
\NormalTok{author\_works }\OtherTok{\textless{}{-}} \FunctionTok{fetch\_openalex}\NormalTok{(}\AttributeTok{entity =} \StringTok{"works"}\NormalTok{,}
                         \AttributeTok{author.id =} \StringTok{"A5023888391"}\NormalTok{)}

\NormalTok{fractional }\OtherTok{\textless{}{-}}\NormalTok{ author\_works }\SpecialCharTok{|\textgreater{}}
 \FunctionTok{mutate}\NormalTok{(}\AttributeTok{n\_authors =} \FunctionTok{map\_int}\NormalTok{(author, nrow),}
        \AttributeTok{frac\_credit =} \DecValTok{1} \SpecialCharTok{/}\NormalTok{ n\_authors) }\SpecialCharTok{|\textgreater{}}
 \FunctionTok{group\_by}\NormalTok{(publication\_year) }\SpecialCharTok{|\textgreater{}}
 \FunctionTok{summarise}\NormalTok{(}\AttributeTok{full\_count =} \FunctionTok{n}\NormalTok{(),}
           \AttributeTok{frac\_count =} \FunctionTok{sum}\NormalTok{(frac\_credit))}

\FunctionTok{ggplot}\NormalTok{(fractional, }\FunctionTok{aes}\NormalTok{(publication\_year)) }\SpecialCharTok{+}
 \FunctionTok{geom\_line}\NormalTok{(}\FunctionTok{aes}\NormalTok{(}\AttributeTok{y =}\NormalTok{ full\_count, }\AttributeTok{colour =} \StringTok{"Full"}\NormalTok{)) }\SpecialCharTok{+}
 \FunctionTok{geom\_line}\NormalTok{(}\FunctionTok{aes}\NormalTok{(}\AttributeTok{y =}\NormalTok{ frac\_count, }\AttributeTok{colour =} \StringTok{"Fractional"}\NormalTok{)) }\SpecialCharTok{+}
 \FunctionTok{labs}\NormalTok{(}\AttributeTok{title =} \StringTok{"Full vs. fractional productivity"}\NormalTok{,}
      \AttributeTok{x =} \StringTok{"Year"}\NormalTok{, }\AttributeTok{y =} \StringTok{"Papers"}\NormalTok{, }\AttributeTok{colour =} \StringTok{"Counting"}\NormalTok{)}
\end{Highlighting}
\end{Shaded}

\emph{Key insight:} Fractional counting penalises hyper-authored papers. Authors who
collaborate broadly see a larger gap between full and fractional counts.

\textbf{5. Self-citation ratio.}
Estimate the fraction of self-citations for an author.

\begin{Shaded}
\begin{Highlighting}[]
\NormalTok{target\_author }\OtherTok{\textless{}{-}} \StringTok{"A5023888391"}

\CommentTok{\# Sample a subset of the author\textquotesingle{}s works}
\NormalTok{sample\_works }\OtherTok{\textless{}{-}}\NormalTok{ author\_works }\SpecialCharTok{|\textgreater{}} \FunctionTok{slice\_sample}\NormalTok{(}\AttributeTok{n =} \DecValTok{10}\NormalTok{)}

\NormalTok{self\_cite\_frac }\OtherTok{\textless{}{-}} \FunctionTok{map\_dbl}\NormalTok{(sample\_works}\SpecialCharTok{$}\NormalTok{id, \textbackslash{}(wid) \{}
\NormalTok{ citing }\OtherTok{\textless{}{-}} \FunctionTok{tryCatch}\NormalTok{(}
   \FunctionTok{fetch\_openalex}\NormalTok{(}\AttributeTok{entity =} \StringTok{"works"}\NormalTok{, }\AttributeTok{cites =}\NormalTok{ wid),}
   \AttributeTok{error =}\NormalTok{ \textbackslash{}(e) }\FunctionTok{tibble}\NormalTok{()}
\NormalTok{ )}
 \ControlFlowTok{if}\NormalTok{ (}\FunctionTok{nrow}\NormalTok{(citing) }\SpecialCharTok{==} \DecValTok{0}\NormalTok{) }\FunctionTok{return}\NormalTok{(}\ConstantTok{NA\_real\_}\NormalTok{)}
\NormalTok{ citing }\SpecialCharTok{|\textgreater{}}
   \FunctionTok{unnest}\NormalTok{(author) }\SpecialCharTok{|\textgreater{}}
   \FunctionTok{summarise}\NormalTok{(}\AttributeTok{has\_self =} \FunctionTok{any}\NormalTok{(au\_id }\SpecialCharTok{==}\NormalTok{ target\_author)) }\SpecialCharTok{|\textgreater{}}
   \FunctionTok{pull}\NormalTok{(has\_self) }\SpecialCharTok{|\textgreater{}}
   \FunctionTok{mean}\NormalTok{()}
\NormalTok{\})}

\FunctionTok{cat}\NormalTok{(}\FunctionTok{sprintf}\NormalTok{(}\StringTok{"Estimated self{-}citation rate: \%.1f\%\%}\SpecialCharTok{\textbackslash{}n}\StringTok{"}\NormalTok{,}
            \FunctionTok{mean}\NormalTok{(self\_cite\_frac, }\AttributeTok{na.rm =} \ConstantTok{TRUE}\NormalTok{) }\SpecialCharTok{*} \DecValTok{100}\NormalTok{))}
\end{Highlighting}
\end{Shaded}

\emph{Key insight:} Self-citation rates of 5-15\% are typical; rates above 25\%
may signal strategic self-citation behaviour.

\section*{Chapter 13 --- Institutional and Country Analysis}\label{chapter-13-institutional-and-country-analysis}
\addcontentsline{toc}{section}{Chapter 13 --- Institutional and Country Analysis}

\textbf{1. Country-level analysis.}
Aggregate by country using full and fractional counts.

\begin{Shaded}
\begin{Highlighting}[]
\FunctionTok{library}\NormalTok{(tidyverse)}

\NormalTok{country\_counts }\OtherTok{\textless{}{-}}\NormalTok{ works }\SpecialCharTok{|\textgreater{}}
 \FunctionTok{unnest}\NormalTok{(authorships) }\SpecialCharTok{|\textgreater{}}
 \FunctionTok{unnest}\NormalTok{(countries) }\SpecialCharTok{|\textgreater{}}
 \FunctionTok{rename}\NormalTok{(}\AttributeTok{country =}\NormalTok{ countries) }\SpecialCharTok{|\textgreater{}}
 \FunctionTok{mutate}\NormalTok{(}\AttributeTok{n\_countries =} \FunctionTok{n\_distinct}\NormalTok{(country), }\AttributeTok{.by =}\NormalTok{ id) }\SpecialCharTok{|\textgreater{}}
 \FunctionTok{summarise}\NormalTok{(}
   \AttributeTok{full\_count =} \FunctionTok{n\_distinct}\NormalTok{(id),}
   \AttributeTok{frac\_count =} \FunctionTok{sum}\NormalTok{(}\DecValTok{1} \SpecialCharTok{/}\NormalTok{ n\_countries),}
   \AttributeTok{.by =}\NormalTok{ country}
\NormalTok{ ) }\SpecialCharTok{|\textgreater{}}
 \FunctionTok{arrange}\NormalTok{(}\FunctionTok{desc}\NormalTok{(full\_count)) }\SpecialCharTok{|\textgreater{}}
 \FunctionTok{head}\NormalTok{(}\DecValTok{10}\NormalTok{)}

\NormalTok{country\_counts }\SpecialCharTok{|\textgreater{}}
 \FunctionTok{mutate}\NormalTok{(}\AttributeTok{ratio =}\NormalTok{ full\_count }\SpecialCharTok{/}\NormalTok{ frac\_count) }\SpecialCharTok{|\textgreater{}}
 \FunctionTok{ggplot}\NormalTok{(}\FunctionTok{aes}\NormalTok{(}\FunctionTok{reorder}\NormalTok{(country, full\_count), ratio)) }\SpecialCharTok{+}
 \FunctionTok{geom\_col}\NormalTok{() }\SpecialCharTok{+}
 \FunctionTok{coord\_flip}\NormalTok{() }\SpecialCharTok{+}
 \FunctionTok{labs}\NormalTok{(}\AttributeTok{title =} \StringTok{"Full / fractional count ratio by country"}\NormalTok{,}
      \AttributeTok{x =} \ConstantTok{NULL}\NormalTok{, }\AttributeTok{y =} \StringTok{"Ratio"}\NormalTok{)}
\end{Highlighting}
\end{Shaded}

\emph{Key insight:} Countries with high international collaboration rates (e.g.,
Switzerland, Singapore) show a larger ratio, meaning full counting inflates
their apparent output.

\textbf{2. PP(top 10\%) by institution.}
Compute the proportion of papers in the top 10\% most cited.

\begin{Shaded}
\begin{Highlighting}[]
\CommentTok{\# Compute year{-}specific 90th percentile thresholds}
\NormalTok{thresholds }\OtherTok{\textless{}{-}}\NormalTok{ works }\SpecialCharTok{|\textgreater{}}
 \FunctionTok{group\_by}\NormalTok{(publication\_year) }\SpecialCharTok{|\textgreater{}}
 \FunctionTok{summarise}\NormalTok{(}\AttributeTok{p90 =} \FunctionTok{quantile}\NormalTok{(cited\_by\_count, }\FloatTok{0.9}\NormalTok{))}

\NormalTok{inst\_top10 }\OtherTok{\textless{}{-}}\NormalTok{ works }\SpecialCharTok{|\textgreater{}}
 \FunctionTok{unnest}\NormalTok{(authorships) }\SpecialCharTok{|\textgreater{}}
 \FunctionTok{unnest}\NormalTok{(institutions) }\SpecialCharTok{|\textgreater{}}
 \FunctionTok{left\_join}\NormalTok{(thresholds, }\AttributeTok{by =} \StringTok{"publication\_year"}\NormalTok{) }\SpecialCharTok{|\textgreater{}}
 \FunctionTok{mutate}\NormalTok{(}\AttributeTok{is\_top10 =}\NormalTok{ cited\_by\_count }\SpecialCharTok{\textgreater{}=}\NormalTok{ p90) }\SpecialCharTok{|\textgreater{}}
 \FunctionTok{summarise}\NormalTok{(}
   \AttributeTok{pp\_top10 =} \FunctionTok{mean}\NormalTok{(is\_top10),}
   \AttributeTok{mncs =} \FunctionTok{mean}\NormalTok{(cited\_by\_count),}
   \AttributeTok{.by =}\NormalTok{ display\_name}
\NormalTok{ ) }\SpecialCharTok{|\textgreater{}}
 \FunctionTok{filter}\NormalTok{(}\SpecialCharTok{!}\FunctionTok{is.na}\NormalTok{(display\_name)) }\SpecialCharTok{|\textgreater{}}
 \FunctionTok{arrange}\NormalTok{(}\FunctionTok{desc}\NormalTok{(pp\_top10))}

\FunctionTok{ggplot}\NormalTok{(inst\_top10 }\SpecialCharTok{|\textgreater{}} \FunctionTok{head}\NormalTok{(}\DecValTok{20}\NormalTok{), }\FunctionTok{aes}\NormalTok{(mncs, pp\_top10)) }\SpecialCharTok{+}
 \FunctionTok{geom\_point}\NormalTok{() }\SpecialCharTok{+}
\NormalTok{ ggrepel}\SpecialCharTok{::}\FunctionTok{geom\_text\_repel}\NormalTok{(}\FunctionTok{aes}\NormalTok{(}\AttributeTok{label =}\NormalTok{ display\_name), }\AttributeTok{size =} \DecValTok{3}\NormalTok{) }\SpecialCharTok{+}
 \FunctionTok{labs}\NormalTok{(}\AttributeTok{x =} \StringTok{"MNCS"}\NormalTok{, }\AttributeTok{y =} \StringTok{"PP(top 10\%)"}\NormalTok{)}
\end{Highlighting}
\end{Shaded}

\emph{Key insight:} PP(top 10\%) and MNCS generally correlate but can diverge if an
institution has a few extremely highly cited outliers.

\textbf{3. Collaboration intensity.}
Relate the number of institutional affiliations per paper to counting method
effects.

\begin{Shaded}
\begin{Highlighting}[]
\NormalTok{collab\_intensity }\OtherTok{\textless{}{-}}\NormalTok{ works }\SpecialCharTok{|\textgreater{}}
 \FunctionTok{unnest}\NormalTok{(authorships) }\SpecialCharTok{|\textgreater{}}
 \FunctionTok{unnest}\NormalTok{(institutions) }\SpecialCharTok{|\textgreater{}}
 \FunctionTok{summarise}\NormalTok{(}
   \AttributeTok{n\_insts =} \FunctionTok{n\_distinct}\NormalTok{(display\_name),}
   \AttributeTok{.by =}\NormalTok{ id}
\NormalTok{ )}

\NormalTok{inst\_collab }\OtherTok{\textless{}{-}}\NormalTok{ works }\SpecialCharTok{|\textgreater{}}
 \FunctionTok{unnest}\NormalTok{(authorships) }\SpecialCharTok{|\textgreater{}}
 \FunctionTok{unnest}\NormalTok{(institutions) }\SpecialCharTok{|\textgreater{}}
 \FunctionTok{left\_join}\NormalTok{(collab\_intensity, }\AttributeTok{by =} \StringTok{"id"}\NormalTok{) }\SpecialCharTok{|\textgreater{}}
 \FunctionTok{summarise}\NormalTok{(}
   \AttributeTok{mean\_collab =} \FunctionTok{mean}\NormalTok{(n\_insts),}
   \AttributeTok{full\_count =} \FunctionTok{n\_distinct}\NormalTok{(id),}
   \AttributeTok{frac\_count =} \FunctionTok{sum}\NormalTok{(}\DecValTok{1} \SpecialCharTok{/}\NormalTok{ n\_insts),}
   \AttributeTok{ratio =}\NormalTok{ full\_count }\SpecialCharTok{/}\NormalTok{ frac\_count,}
   \AttributeTok{.by =}\NormalTok{ display\_name}
\NormalTok{ ) }\SpecialCharTok{|\textgreater{}}
 \FunctionTok{filter}\NormalTok{(}\SpecialCharTok{!}\FunctionTok{is.na}\NormalTok{(display\_name))}

\FunctionTok{ggplot}\NormalTok{(inst\_collab, }\FunctionTok{aes}\NormalTok{(mean\_collab, ratio)) }\SpecialCharTok{+}
 \FunctionTok{geom\_point}\NormalTok{(}\AttributeTok{alpha =} \FloatTok{0.5}\NormalTok{) }\SpecialCharTok{+}
 \FunctionTok{geom\_smooth}\NormalTok{(}\AttributeTok{method =} \StringTok{"lm"}\NormalTok{) }\SpecialCharTok{+}
 \FunctionTok{labs}\NormalTok{(}\AttributeTok{x =} \StringTok{"Mean affiliations per paper"}\NormalTok{,}
      \AttributeTok{y =} \StringTok{"Full / fractional ratio"}\NormalTok{,}
      \AttributeTok{title =} \StringTok{"Collaboration intensity vs. counting effect"}\NormalTok{)}
\end{Highlighting}
\end{Shaded}

\emph{Key insight:} There is a positive relationship --- institutions with higher
collaboration intensity are more inflated by full counting.

\section*{Chapter 14 --- Temporal Analysis and Citation Aging}\label{chapter-14-temporal-analysis-and-citation-aging}
\addcontentsline{toc}{section}{Chapter 14 --- Temporal Analysis and Citation Aging}

\textbf{1. Half-life comparison.}
Compare citation half-lives between a mathematics and a biomedical journal.

\begin{Shaded}
\begin{Highlighting}[]
\FunctionTok{library}\NormalTok{(openalexR)}
\FunctionTok{library}\NormalTok{(tidyverse)}

\NormalTok{fetch\_half\_life }\OtherTok{\textless{}{-}} \ControlFlowTok{function}\NormalTok{(journal\_id, label) \{}
\NormalTok{ w }\OtherTok{\textless{}{-}} \FunctionTok{fetch\_openalex}\NormalTok{(}\AttributeTok{entity =} \StringTok{"works"}\NormalTok{,}
               \AttributeTok{primary\_location.source.id =}\NormalTok{ journal\_id,}
               \AttributeTok{from\_publication\_date =} \StringTok{"2015{-}01{-}01"}\NormalTok{,}
               \AttributeTok{to\_publication\_date =} \StringTok{"2018{-}12{-}31"}\NormalTok{)}
\NormalTok{ w }\SpecialCharTok{|\textgreater{}}
   \FunctionTok{mutate}\NormalTok{(}\AttributeTok{age =} \DecValTok{2024} \SpecialCharTok{{-}}\NormalTok{ publication\_year) }\SpecialCharTok{|\textgreater{}}
   \FunctionTok{group\_by}\NormalTok{(age) }\SpecialCharTok{|\textgreater{}}
   \FunctionTok{summarise}\NormalTok{(}\AttributeTok{mean\_cites =} \FunctionTok{mean}\NormalTok{(cited\_by\_count)) }\SpecialCharTok{|\textgreater{}}
   \FunctionTok{mutate}\NormalTok{(}\AttributeTok{cum\_share =} \FunctionTok{cumsum}\NormalTok{(mean\_cites) }\SpecialCharTok{/} \FunctionTok{sum}\NormalTok{(mean\_cites),}
          \AttributeTok{journal =}\NormalTok{ label)}
\NormalTok{\}}

\NormalTok{math\_hl   }\OtherTok{\textless{}{-}} \FunctionTok{fetch\_half\_life}\NormalTok{(}\StringTok{"S154589771"}\NormalTok{, }\StringTok{"Mathematics"}\NormalTok{)}
\NormalTok{biomed\_hl }\OtherTok{\textless{}{-}} \FunctionTok{fetch\_half\_life}\NormalTok{(}\StringTok{"S49861241"}\NormalTok{, }\StringTok{"Biomedicine"}\NormalTok{)}

\FunctionTok{bind\_rows}\NormalTok{(math\_hl, biomed\_hl) }\SpecialCharTok{|\textgreater{}}
 \FunctionTok{ggplot}\NormalTok{(}\FunctionTok{aes}\NormalTok{(age, cum\_share, }\AttributeTok{colour =}\NormalTok{ journal)) }\SpecialCharTok{+}
 \FunctionTok{geom\_line}\NormalTok{() }\SpecialCharTok{+}
 \FunctionTok{geom\_hline}\NormalTok{(}\AttributeTok{yintercept =} \FloatTok{0.5}\NormalTok{, }\AttributeTok{linetype =} \StringTok{"dashed"}\NormalTok{) }\SpecialCharTok{+}
 \FunctionTok{labs}\NormalTok{(}\AttributeTok{title =} \StringTok{"Citation half{-}life comparison"}\NormalTok{,}
      \AttributeTok{x =} \StringTok{"Age (years)"}\NormalTok{, }\AttributeTok{y =} \StringTok{"Cumulative citation share"}\NormalTok{)}
\end{Highlighting}
\end{Shaded}

\emph{Key insight:} Biomedical journals typically have shorter half-lives (3-5 years)
than mathematics journals (8-12 years) because of faster publication cycles.

\textbf{2. Price index computation.}
Compute the fraction of references less than 5 years old.

\begin{Shaded}
\begin{Highlighting}[]
\CommentTok{\# For papers with reference data available}
\NormalTok{sample\_papers }\OtherTok{\textless{}{-}}\NormalTok{ works }\SpecialCharTok{|\textgreater{}}
 \FunctionTok{filter}\NormalTok{(publication\_year }\SpecialCharTok{==} \DecValTok{2020}\NormalTok{) }\SpecialCharTok{|\textgreater{}}
 \FunctionTok{slice\_sample}\NormalTok{(}\AttributeTok{n =} \DecValTok{20}\NormalTok{)}

\NormalTok{price\_indices }\OtherTok{\textless{}{-}} \FunctionTok{map\_dbl}\NormalTok{(sample\_papers}\SpecialCharTok{$}\NormalTok{id, \textbackslash{}(wid) \{}
\NormalTok{ refs }\OtherTok{\textless{}{-}} \FunctionTok{tryCatch}\NormalTok{(}
   \FunctionTok{fetch\_openalex}\NormalTok{(}\AttributeTok{entity =} \StringTok{"works"}\NormalTok{, }\AttributeTok{cited\_by =}\NormalTok{ wid),}
   \AttributeTok{error =}\NormalTok{ \textbackslash{}(e) }\FunctionTok{tibble}\NormalTok{()}
\NormalTok{ )}
 \ControlFlowTok{if}\NormalTok{ (}\FunctionTok{nrow}\NormalTok{(refs) }\SpecialCharTok{==} \DecValTok{0}\NormalTok{) }\FunctionTok{return}\NormalTok{(}\ConstantTok{NA\_real\_}\NormalTok{)}
\NormalTok{ pub\_year }\OtherTok{\textless{}{-}}\NormalTok{ sample\_papers}\SpecialCharTok{$}\NormalTok{publication\_year[sample\_papers}\SpecialCharTok{$}\NormalTok{id }\SpecialCharTok{==}\NormalTok{ wid]}
 \FunctionTok{mean}\NormalTok{(refs}\SpecialCharTok{$}\NormalTok{publication\_year }\SpecialCharTok{\textgreater{}=}\NormalTok{ (pub\_year }\SpecialCharTok{{-}} \DecValTok{5}\NormalTok{), }\AttributeTok{na.rm =} \ConstantTok{TRUE}\NormalTok{)}
\NormalTok{\})}

\FunctionTok{cat}\NormalTok{(}\FunctionTok{sprintf}\NormalTok{(}\StringTok{"Mean Price index: \%.2f}\SpecialCharTok{\textbackslash{}n}\StringTok{"}\NormalTok{, }\FunctionTok{mean}\NormalTok{(price\_indices, }\AttributeTok{na.rm =} \ConstantTok{TRUE}\NormalTok{)))}
\end{Highlighting}
\end{Shaded}

\emph{Key insight:} A high Price index indicates the field relies on recent work
(fast-moving); a low Price index suggests reliance on older, established
literature.

\textbf{3. Sleeping beauty detection.}
Identify delayed-recognition papers.

\begin{Shaded}
\begin{Highlighting}[]
\NormalTok{sleeping }\OtherTok{\textless{}{-}}\NormalTok{ works }\SpecialCharTok{|\textgreater{}}
 \FunctionTok{filter}\NormalTok{(publication\_year }\SpecialCharTok{\textless{}=} \DecValTok{2018}\NormalTok{,}
\NormalTok{        cited\_by\_count }\SpecialCharTok{\textgreater{}=} \DecValTok{20}\NormalTok{) }\SpecialCharTok{|\textgreater{}}
 \FunctionTok{mutate}\NormalTok{(}\AttributeTok{age =} \DecValTok{2024} \SpecialCharTok{{-}}\NormalTok{ publication\_year)}

\CommentTok{\# Approximate: papers with low early citation rate but high total}
\CommentTok{\# (exact per{-}year citation data requires counts\_by\_year from OpenAlex)}
\NormalTok{sleeping\_candidates }\OtherTok{\textless{}{-}}\NormalTok{ sleeping }\SpecialCharTok{|\textgreater{}}
 \FunctionTok{filter}\NormalTok{(}
   \FunctionTok{map\_lgl}\NormalTok{(counts\_by\_year, \textbackslash{}(cy) \{}
     \ControlFlowTok{if}\NormalTok{ (}\FunctionTok{is.null}\NormalTok{(cy)) }\FunctionTok{return}\NormalTok{(}\ConstantTok{FALSE}\NormalTok{)}
\NormalTok{     early }\OtherTok{\textless{}{-}}\NormalTok{ cy }\SpecialCharTok{|\textgreater{}} \FunctionTok{filter}\NormalTok{(year }\SpecialCharTok{\textless{}=}\NormalTok{ publication\_year }\SpecialCharTok{+} \DecValTok{3}\NormalTok{)}
     \FunctionTok{sum}\NormalTok{(early}\SpecialCharTok{$}\NormalTok{cited\_by\_count, }\AttributeTok{na.rm =} \ConstantTok{TRUE}\NormalTok{) }\SpecialCharTok{\textless{}} \DecValTok{2}
\NormalTok{   \})}
\NormalTok{ )}

\NormalTok{sleeping\_candidates }\SpecialCharTok{|\textgreater{}}
 \FunctionTok{select}\NormalTok{(display\_name, publication\_year, cited\_by\_count) }\SpecialCharTok{|\textgreater{}}
 \FunctionTok{arrange}\NormalTok{(}\FunctionTok{desc}\NormalTok{(cited\_by\_count))}
\end{Highlighting}
\end{Shaded}

\emph{Key insight:} Sleeping beauties are rare but occur across all fields; they
often represent ahead-of-their-time ideas that gained traction only after a
paradigm shift.

\textbf{4. Median vs.~mean aging.}
Replot aging curves using median citations.

\begin{Shaded}
\begin{Highlighting}[]
\NormalTok{aging\_comparison }\OtherTok{\textless{}{-}}\NormalTok{ works }\SpecialCharTok{|\textgreater{}}
 \FunctionTok{mutate}\NormalTok{(}\AttributeTok{age =} \DecValTok{2024} \SpecialCharTok{{-}}\NormalTok{ publication\_year) }\SpecialCharTok{|\textgreater{}}
 \FunctionTok{group\_by}\NormalTok{(age) }\SpecialCharTok{|\textgreater{}}
 \FunctionTok{summarise}\NormalTok{(}
   \AttributeTok{mean\_cites   =} \FunctionTok{mean}\NormalTok{(cited\_by\_count),}
   \AttributeTok{median\_cites =} \FunctionTok{median}\NormalTok{(cited\_by\_count)}
\NormalTok{ ) }\SpecialCharTok{|\textgreater{}}
 \FunctionTok{pivot\_longer}\NormalTok{(}\FunctionTok{c}\NormalTok{(mean\_cites, median\_cites),}
              \AttributeTok{names\_to =} \StringTok{"stat"}\NormalTok{, }\AttributeTok{values\_to =} \StringTok{"value"}\NormalTok{)}

\FunctionTok{ggplot}\NormalTok{(aging\_comparison, }\FunctionTok{aes}\NormalTok{(age, value, }\AttributeTok{colour =}\NormalTok{ stat)) }\SpecialCharTok{+}
 \FunctionTok{geom\_line}\NormalTok{() }\SpecialCharTok{+}
 \FunctionTok{geom\_point}\NormalTok{() }\SpecialCharTok{+}
 \FunctionTok{labs}\NormalTok{(}\AttributeTok{title =} \StringTok{"Mean vs. median citation aging"}\NormalTok{,}
      \AttributeTok{x =} \StringTok{"Age (years)"}\NormalTok{, }\AttributeTok{y =} \StringTok{"Citations"}\NormalTok{)}
\end{Highlighting}
\end{Shaded}

\emph{Key insight:} The median curve is lower and flatter because the mean is pulled
up by highly cited outliers. Median is more robust for skewed distributions.

\textbf{5. Publication-cohort sensitivity.}
Use the Scientometrics sample fetched in Chapter 14. Its publication dates span
2013--2018; the following cohorts end in 2018 and remain within that retrieved
range. These are current cumulative citation summaries for subsets of a sample.

\begin{Shaded}
\begin{Highlighting}[]
\NormalTok{cohort\_end }\OtherTok{\textless{}{-}} \DecValTok{2018}
\NormalTok{cohort\_lengths }\OtherTok{\textless{}{-}} \FunctionTok{c}\NormalTok{(}\DecValTok{2}\NormalTok{, }\DecValTok{3}\NormalTok{, }\DecValTok{5}\NormalTok{)}
\NormalTok{cohort\_summaries }\OtherTok{\textless{}{-}} \FunctionTok{map\_dfr}\NormalTok{(cohort\_lengths, }\ControlFlowTok{function}\NormalTok{(length\_years) \{}
\NormalTok{  cohort }\OtherTok{\textless{}{-}}\NormalTok{ works\_sciento }\SpecialCharTok{|\textgreater{}}
    \FunctionTok{mutate}\NormalTok{(}\AttributeTok{publication\_year =}\NormalTok{ lubridate}\SpecialCharTok{::}\FunctionTok{year}\NormalTok{(publication\_date)) }\SpecialCharTok{|\textgreater{}}
    \FunctionTok{filter}\NormalTok{(publication\_year }\SpecialCharTok{\textgreater{}=}\NormalTok{ cohort\_end }\SpecialCharTok{{-}}\NormalTok{ length\_years }\SpecialCharTok{+} \DecValTok{1}\NormalTok{,}
\NormalTok{           publication\_year }\SpecialCharTok{\textless{}=}\NormalTok{ cohort\_end)}
\NormalTok{  observed }\OtherTok{\textless{}{-}}\NormalTok{ cohort}\SpecialCharTok{$}\NormalTok{cited\_by\_count[}\SpecialCharTok{!}\FunctionTok{is.na}\NormalTok{(cohort}\SpecialCharTok{$}\NormalTok{cited\_by\_count)]}
  \FunctionTok{tibble}\NormalTok{(}\AttributeTok{cohort\_years =}\NormalTok{ length\_years, }\AttributeTok{n\_records =} \FunctionTok{nrow}\NormalTok{(cohort),}
         \AttributeTok{n\_citations\_observed =} \FunctionTok{length}\NormalTok{(observed),}
         \AttributeTok{mean\_cumulative\_citations =} \ControlFlowTok{if}\NormalTok{ (}\FunctionTok{length}\NormalTok{(observed)) }\FunctionTok{mean}\NormalTok{(observed) }\ControlFlowTok{else} \ConstantTok{NA\_real\_}\NormalTok{)}
\NormalTok{\})}
\NormalTok{cohort\_summaries}

\FunctionTok{ggplot}\NormalTok{(cohort\_summaries, }\FunctionTok{aes}\NormalTok{(}\FunctionTok{factor}\NormalTok{(cohort\_years), mean\_cumulative\_citations)) }\SpecialCharTok{+}
  \FunctionTok{geom\_col}\NormalTok{() }\SpecialCharTok{+}
  \FunctionTok{labs}\NormalTok{(}\AttributeTok{x =} \StringTok{"Publication{-}cohort length (years)"}\NormalTok{,}
       \AttributeTok{y =} \StringTok{"Mean cumulative citations at retrieval"}\NormalTok{,}
       \AttributeTok{title =} \StringTok{"Descriptive sensitivity to publication{-}cohort selection"}\NormalTok{)}
\end{Highlighting}
\end{Shaded}

Changing the publication cohort changes the included papers and their citation
exposure. The mean may rise or fall; it is not guaranteed to increase with cohort
length. A publication-date filter does not restrict citation-receipt dates.
These calculations therefore do not estimate JIF or compare fixed post-publication
citation windows. Missing citation counts remain outside the observed denominator.

\section*{Chapter 15 --- Co-authorship Networks}\label{chapter-15-co-authorship-networks}
\addcontentsline{toc}{section}{Chapter 15 --- Co-authorship Networks}

\textbf{1. Fractional counting.}
Weight edges by 1/(k-1) where k is the number of authors.

\begin{Shaded}
\begin{Highlighting}[]
\FunctionTok{library}\NormalTok{(igraph)}
\FunctionTok{library}\NormalTok{(tidyverse)}

\NormalTok{edges\_frac }\OtherTok{\textless{}{-}}\NormalTok{ works }\SpecialCharTok{|\textgreater{}}
 \FunctionTok{unnest}\NormalTok{(author) }\SpecialCharTok{|\textgreater{}}
 \FunctionTok{group\_by}\NormalTok{(id) }\SpecialCharTok{|\textgreater{}}
 \FunctionTok{mutate}\NormalTok{(}\AttributeTok{k =} \FunctionTok{n}\NormalTok{()) }\SpecialCharTok{|\textgreater{}}
 \FunctionTok{filter}\NormalTok{(k }\SpecialCharTok{\textgreater{}=} \DecValTok{2}\NormalTok{) }\SpecialCharTok{|\textgreater{}}
 \FunctionTok{reframe}\NormalTok{(}
   \AttributeTok{pair =} \FunctionTok{combn}\NormalTok{(au\_id, }\DecValTok{2}\NormalTok{, }\AttributeTok{simplify =} \ConstantTok{FALSE}\NormalTok{),}
   \AttributeTok{weight =} \DecValTok{1} \SpecialCharTok{/}\NormalTok{ (k }\SpecialCharTok{{-}} \DecValTok{1}\NormalTok{)}
\NormalTok{ ) }\SpecialCharTok{|\textgreater{}}
 \FunctionTok{unnest\_wider}\NormalTok{(pair, }\AttributeTok{names\_sep =} \StringTok{"\_"}\NormalTok{) }\SpecialCharTok{|\textgreater{}}
 \FunctionTok{rename}\NormalTok{(}\AttributeTok{from =}\NormalTok{ pair\_1, }\AttributeTok{to =}\NormalTok{ pair\_2)}

\NormalTok{g\_frac }\OtherTok{\textless{}{-}} \FunctionTok{graph\_from\_data\_frame}\NormalTok{(edges\_frac, }\AttributeTok{directed =} \ConstantTok{FALSE}\NormalTok{)}
\FunctionTok{E}\NormalTok{(g\_frac)}\SpecialCharTok{$}\NormalTok{weight }\OtherTok{\textless{}{-}}\NormalTok{ edges\_frac}\SpecialCharTok{$}\NormalTok{weight}

\NormalTok{degree\_frac }\OtherTok{\textless{}{-}} \FunctionTok{strength}\NormalTok{(g\_frac) }\SpecialCharTok{|\textgreater{}}
 \FunctionTok{sort}\NormalTok{(}\AttributeTok{decreasing =} \ConstantTok{TRUE}\NormalTok{) }\SpecialCharTok{|\textgreater{}}
 \FunctionTok{head}\NormalTok{(}\DecValTok{10}\NormalTok{)}
\NormalTok{degree\_frac}
\end{Highlighting}
\end{Shaded}

\emph{Key insight:} Fractional counting reduces the influence of mega-authored
papers. Authors primarily involved in large teams may drop in rank.

\textbf{2. Temporal slicing.}
Build yearly networks and track structural evolution.

\begin{Shaded}
\begin{Highlighting}[]
\NormalTok{yearly\_stats }\OtherTok{\textless{}{-}}\NormalTok{ works }\SpecialCharTok{|\textgreater{}}
 \FunctionTok{group\_by}\NormalTok{(publication\_year) }\SpecialCharTok{|\textgreater{}}
 \FunctionTok{group\_map}\NormalTok{(\textbackslash{}(df, key) \{}
\NormalTok{   edges }\OtherTok{\textless{}{-}}\NormalTok{ df }\SpecialCharTok{|\textgreater{}}
     \FunctionTok{unnest}\NormalTok{(author) }\SpecialCharTok{|\textgreater{}}
     \FunctionTok{group\_by}\NormalTok{(id) }\SpecialCharTok{|\textgreater{}}
     \FunctionTok{filter}\NormalTok{(}\FunctionTok{n}\NormalTok{() }\SpecialCharTok{\textgreater{}=} \DecValTok{2}\NormalTok{) }\SpecialCharTok{|\textgreater{}}
     \FunctionTok{reframe}\NormalTok{(}\AttributeTok{pair =} \FunctionTok{combn}\NormalTok{(au\_id, }\DecValTok{2}\NormalTok{, }\AttributeTok{simplify =} \ConstantTok{FALSE}\NormalTok{)) }\SpecialCharTok{|\textgreater{}}
     \FunctionTok{unnest\_wider}\NormalTok{(pair, }\AttributeTok{names\_sep =} \StringTok{"\_"}\NormalTok{)}
   \ControlFlowTok{if}\NormalTok{ (}\FunctionTok{nrow}\NormalTok{(edges) }\SpecialCharTok{==} \DecValTok{0}\NormalTok{) }\FunctionTok{return}\NormalTok{(}\ConstantTok{NULL}\NormalTok{)}
\NormalTok{   g }\OtherTok{\textless{}{-}} \FunctionTok{graph\_from\_data\_frame}\NormalTok{(edges, }\AttributeTok{directed =} \ConstantTok{FALSE}\NormalTok{)}
   \FunctionTok{tibble}\NormalTok{(}
     \AttributeTok{year =}\NormalTok{ key}\SpecialCharTok{$}\NormalTok{publication\_year,}
     \AttributeTok{n\_nodes =} \FunctionTok{vcount}\NormalTok{(g),}
     \AttributeTok{n\_components =} \FunctionTok{components}\NormalTok{(g)}\SpecialCharTok{$}\NormalTok{no,}
     \AttributeTok{mean\_degree =} \FunctionTok{mean}\NormalTok{(}\FunctionTok{degree}\NormalTok{(g))}
\NormalTok{   )}
\NormalTok{ \}) }\SpecialCharTok{|\textgreater{}} \FunctionTok{list\_rbind}\NormalTok{()}

\FunctionTok{ggplot}\NormalTok{(yearly\_stats, }\FunctionTok{aes}\NormalTok{(year, mean\_degree)) }\SpecialCharTok{+}
 \FunctionTok{geom\_line}\NormalTok{() }\SpecialCharTok{+}
 \FunctionTok{geom\_point}\NormalTok{() }\SpecialCharTok{+}
 \FunctionTok{labs}\NormalTok{(}\AttributeTok{title =} \StringTok{"Mean degree over time"}\NormalTok{, }\AttributeTok{x =} \StringTok{"Year"}\NormalTok{, }\AttributeTok{y =} \StringTok{"Mean degree"}\NormalTok{)}
\end{Highlighting}
\end{Shaded}

\emph{Key insight:} Mean degree typically increases over time as collaboration
becomes more common; the number of components may decrease as the field becomes
more connected.

\textbf{3. Betweenness vs.~degree.}
Compare centrality measures for the giant component.

\begin{Shaded}
\begin{Highlighting}[]
\NormalTok{g }\OtherTok{\textless{}{-}} \FunctionTok{build\_coauth\_graph}\NormalTok{(works)}
\NormalTok{gc }\OtherTok{\textless{}{-}} \FunctionTok{induced\_subgraph}\NormalTok{(g, }\FunctionTok{components}\NormalTok{(g)}\SpecialCharTok{$}\NormalTok{membership }\SpecialCharTok{==}
                         \FunctionTok{which.max}\NormalTok{(}\FunctionTok{components}\NormalTok{(g)}\SpecialCharTok{$}\NormalTok{csize))}

\NormalTok{centrality }\OtherTok{\textless{}{-}} \FunctionTok{tibble}\NormalTok{(}
 \AttributeTok{author =} \FunctionTok{V}\NormalTok{(gc)}\SpecialCharTok{$}\NormalTok{name,}
 \AttributeTok{deg =} \FunctionTok{degree}\NormalTok{(gc),}
 \AttributeTok{btw =} \FunctionTok{betweenness}\NormalTok{(gc, }\AttributeTok{normalized =} \ConstantTok{TRUE}\NormalTok{)}
\NormalTok{)}

\FunctionTok{ggplot}\NormalTok{(centrality, }\FunctionTok{aes}\NormalTok{(deg, btw)) }\SpecialCharTok{+}
 \FunctionTok{geom\_point}\NormalTok{(}\AttributeTok{alpha =} \FloatTok{0.4}\NormalTok{) }\SpecialCharTok{+}
 \FunctionTok{labs}\NormalTok{(}\AttributeTok{x =} \StringTok{"Degree"}\NormalTok{, }\AttributeTok{y =} \StringTok{"Betweenness (normalised)"}\NormalTok{,}
      \AttributeTok{title =} \StringTok{"Betweenness vs. degree in co{-}authorship network"}\NormalTok{)}
\end{Highlighting}
\end{Shaded}

\emph{Key insight:} High-betweenness, low-degree nodes are ``bridge'' authors who
connect otherwise separate groups --- they may be interdisciplinary researchers.

\textbf{4. Resolution parameter.}
Vary the Leiden resolution and compare community counts.

\begin{Shaded}
\begin{Highlighting}[]
\FunctionTok{library}\NormalTok{(igraph)}

\NormalTok{resolutions }\OtherTok{\textless{}{-}} \FunctionTok{c}\NormalTok{(}\FloatTok{0.5}\NormalTok{, }\FloatTok{1.0}\NormalTok{, }\FloatTok{2.0}\NormalTok{)}

\NormalTok{res\_results }\OtherTok{\textless{}{-}} \FunctionTok{map}\NormalTok{(resolutions, \textbackslash{}(r) \{}
\NormalTok{ comm }\OtherTok{\textless{}{-}} \FunctionTok{cluster\_leiden}\NormalTok{(gc, }\AttributeTok{resolution =}\NormalTok{ r,}
                        \AttributeTok{objective\_function =} \StringTok{"modularity"}\NormalTok{)}
 \FunctionTok{tibble}\NormalTok{(}\AttributeTok{resolution =}\NormalTok{ r,}
        \AttributeTok{n\_communities =} \FunctionTok{length}\NormalTok{(comm),}
        \AttributeTok{modularity =} \FunctionTok{modularity}\NormalTok{(comm))}
\NormalTok{\})}

\FunctionTok{list\_rbind}\NormalTok{(res\_results)}
\end{Highlighting}
\end{Shaded}

\emph{Key insight:} Higher resolution produces more communities. The optimal value
depends on the application --- lower resolution reveals broad thematic clusters,
higher resolution reveals tight collaborative teams.

\section*{Chapter 16 --- Co-citation and Bibliographic Coupling}\label{chapter-16-co-citation-and-bibliographic-coupling}
\addcontentsline{toc}{section}{Chapter 16 --- Co-citation and Bibliographic Coupling}

\textbf{1. Cosine normalisation.}
Normalise co-citation strengths by the geometric mean of individual citation
counts.

\begin{Shaded}
\begin{Highlighting}[]
\FunctionTok{library}\NormalTok{(igraph)}
\FunctionTok{library}\NormalTok{(tidyverse)}

\CommentTok{\# Assume cocit\_matrix is the raw co{-}citation matrix}
\NormalTok{diag\_sqrt }\OtherTok{\textless{}{-}} \FunctionTok{sqrt}\NormalTok{(}\FunctionTok{diag}\NormalTok{(cocit\_matrix))}
\NormalTok{cosine\_matrix }\OtherTok{\textless{}{-}}\NormalTok{ cocit\_matrix }\SpecialCharTok{/} \FunctionTok{outer}\NormalTok{(diag\_sqrt, diag\_sqrt)}
\FunctionTok{diag}\NormalTok{(cosine\_matrix) }\OtherTok{\textless{}{-}} \DecValTok{0}

\NormalTok{g\_cosine }\OtherTok{\textless{}{-}} \FunctionTok{graph\_from\_adjacency\_matrix}\NormalTok{(cosine\_matrix,}
                                         \AttributeTok{mode =} \StringTok{"undirected"}\NormalTok{,}
                                         \AttributeTok{weighted =} \ConstantTok{TRUE}\NormalTok{)}

\CommentTok{\# Compare top pairs}
\NormalTok{top\_pairs\_raw }\OtherTok{\textless{}{-}} \FunctionTok{E}\NormalTok{(g\_raw)[}\FunctionTok{order}\NormalTok{(}\SpecialCharTok{{-}}\FunctionTok{E}\NormalTok{(g\_raw)}\SpecialCharTok{$}\NormalTok{weight)][}\DecValTok{1}\SpecialCharTok{:}\DecValTok{10}\NormalTok{]}
\NormalTok{top\_pairs\_cos }\OtherTok{\textless{}{-}} \FunctionTok{E}\NormalTok{(g\_cosine)[}\FunctionTok{order}\NormalTok{(}\SpecialCharTok{{-}}\FunctionTok{E}\NormalTok{(g\_cosine)}\SpecialCharTok{$}\NormalTok{weight)][}\DecValTok{1}\SpecialCharTok{:}\DecValTok{10}\NormalTok{]}
\end{Highlighting}
\end{Shaded}

\emph{Key insight:} Cosine normalisation down-weights pairs that are individually
highly cited, surfacing topically related pairs that raw counts miss.

\textbf{2. Temporal co-citation.}
Split the corpus into two periods and build separate networks.

\begin{Shaded}
\begin{Highlighting}[]
\NormalTok{periods }\OtherTok{\textless{}{-}} \FunctionTok{list}\NormalTok{(}
 \AttributeTok{early =}\NormalTok{ works }\SpecialCharTok{|\textgreater{}} \FunctionTok{filter}\NormalTok{(publication\_year }\SpecialCharTok{\%in\%} \DecValTok{2020}\SpecialCharTok{:}\DecValTok{2021}\NormalTok{),}
 \AttributeTok{late  =}\NormalTok{ works }\SpecialCharTok{|\textgreater{}} \FunctionTok{filter}\NormalTok{(publication\_year }\SpecialCharTok{\%in\%} \DecValTok{2022}\SpecialCharTok{:}\DecValTok{2023}\NormalTok{)}
\NormalTok{)}

\NormalTok{period\_graphs }\OtherTok{\textless{}{-}} \FunctionTok{map}\NormalTok{(periods, \textbackslash{}(df) \{}
 \CommentTok{\# Build co{-}citation network for this period\textquotesingle{}s references}
\NormalTok{ refs }\OtherTok{\textless{}{-}}\NormalTok{ df }\SpecialCharTok{|\textgreater{}} \FunctionTok{unnest}\NormalTok{(referenced\_works)}
\NormalTok{ ref\_pairs }\OtherTok{\textless{}{-}}\NormalTok{ refs }\SpecialCharTok{|\textgreater{}}
   \FunctionTok{group\_by}\NormalTok{(id) }\SpecialCharTok{|\textgreater{}}
   \FunctionTok{filter}\NormalTok{(}\FunctionTok{n}\NormalTok{() }\SpecialCharTok{\textgreater{}=} \DecValTok{2}\NormalTok{) }\SpecialCharTok{|\textgreater{}}
   \FunctionTok{reframe}\NormalTok{(}\AttributeTok{pair =} \FunctionTok{combn}\NormalTok{(referenced\_works, }\DecValTok{2}\NormalTok{, }\AttributeTok{simplify =} \ConstantTok{FALSE}\NormalTok{)) }\SpecialCharTok{|\textgreater{}}
   \FunctionTok{unnest\_wider}\NormalTok{(pair, }\AttributeTok{names\_sep =} \StringTok{"\_"}\NormalTok{)}

\NormalTok{ ref\_pairs }\SpecialCharTok{|\textgreater{}}
   \FunctionTok{count}\NormalTok{(pair\_1, pair\_2, }\AttributeTok{name =} \StringTok{"weight"}\NormalTok{) }\SpecialCharTok{|\textgreater{}}
   \FunctionTok{graph\_from\_data\_frame}\NormalTok{(}\AttributeTok{directed =} \ConstantTok{FALSE}\NormalTok{)}
\NormalTok{\})}

\CommentTok{\# Find stable pairs: edges present in both periods}
\NormalTok{shared\_edges }\OtherTok{\textless{}{-}} \FunctionTok{E}\NormalTok{(period\_graphs}\SpecialCharTok{$}\NormalTok{early) }\SpecialCharTok{\%in\%} \FunctionTok{E}\NormalTok{(period\_graphs}\SpecialCharTok{$}\NormalTok{late)}
\end{Highlighting}
\end{Shaded}

\emph{Key insight:} Stable co-citation pairs represent the intellectual core of the
field; new pairs emerging in the later period may signal research-front shifts.

\textbf{3. Bibliographic coupling by year.}
Track modularity and fragmentation of the intellectual structure.

\begin{Shaded}
\begin{Highlighting}[]
\NormalTok{yearly\_bc }\OtherTok{\textless{}{-}}\NormalTok{ works }\SpecialCharTok{|\textgreater{}}
 \FunctionTok{group\_by}\NormalTok{(publication\_year) }\SpecialCharTok{|\textgreater{}}
 \FunctionTok{group\_map}\NormalTok{(\textbackslash{}(df, key) \{}
\NormalTok{   refs }\OtherTok{\textless{}{-}}\NormalTok{ df }\SpecialCharTok{|\textgreater{}} \FunctionTok{unnest}\NormalTok{(referenced\_works)}
\NormalTok{   pairs }\OtherTok{\textless{}{-}}\NormalTok{ refs }\SpecialCharTok{|\textgreater{}}
     \FunctionTok{group\_by}\NormalTok{(referenced\_works) }\SpecialCharTok{|\textgreater{}}
     \FunctionTok{filter}\NormalTok{(}\FunctionTok{n}\NormalTok{() }\SpecialCharTok{\textgreater{}=} \DecValTok{2}\NormalTok{) }\SpecialCharTok{|\textgreater{}}
     \FunctionTok{reframe}\NormalTok{(}\AttributeTok{pair =} \FunctionTok{combn}\NormalTok{(id, }\DecValTok{2}\NormalTok{, }\AttributeTok{simplify =} \ConstantTok{FALSE}\NormalTok{)) }\SpecialCharTok{|\textgreater{}}
     \FunctionTok{unnest\_wider}\NormalTok{(pair, }\AttributeTok{names\_sep =} \StringTok{"\_"}\NormalTok{) }\SpecialCharTok{|\textgreater{}}
     \FunctionTok{count}\NormalTok{(pair\_1, pair\_2, }\AttributeTok{name =} \StringTok{"weight"}\NormalTok{)}
   \ControlFlowTok{if}\NormalTok{ (}\FunctionTok{nrow}\NormalTok{(pairs) }\SpecialCharTok{==} \DecValTok{0}\NormalTok{) }\FunctionTok{return}\NormalTok{(}\ConstantTok{NULL}\NormalTok{)}
\NormalTok{   g }\OtherTok{\textless{}{-}} \FunctionTok{graph\_from\_data\_frame}\NormalTok{(pairs, }\AttributeTok{directed =} \ConstantTok{FALSE}\NormalTok{)}
\NormalTok{   comm }\OtherTok{\textless{}{-}} \FunctionTok{cluster\_leiden}\NormalTok{(g, }\AttributeTok{resolution =} \FloatTok{1.0}\NormalTok{)}
   \FunctionTok{tibble}\NormalTok{(}\AttributeTok{year =}\NormalTok{ key}\SpecialCharTok{$}\NormalTok{publication\_year,}
          \AttributeTok{n\_components =} \FunctionTok{components}\NormalTok{(g)}\SpecialCharTok{$}\NormalTok{no,}
          \AttributeTok{modularity =} \FunctionTok{modularity}\NormalTok{(comm))}
\NormalTok{ \}) }\SpecialCharTok{|\textgreater{}} \FunctionTok{list\_rbind}\NormalTok{()}

\FunctionTok{ggplot}\NormalTok{(yearly\_bc, }\FunctionTok{aes}\NormalTok{(year, modularity)) }\SpecialCharTok{+}
 \FunctionTok{geom\_line}\NormalTok{() }\SpecialCharTok{+}
 \FunctionTok{labs}\NormalTok{(}\AttributeTok{title =} \StringTok{"Bibliographic coupling modularity over time"}\NormalTok{)}
\end{Highlighting}
\end{Shaded}

\emph{Key insight:} Decreasing modularity suggests convergence across subfields;
increasing modularity suggests specialisation and fragmentation.

\textbf{4. Identifying research fronts.}
Label communities using top title words.

\begin{Shaded}
\begin{Highlighting}[]
\FunctionTok{library}\NormalTok{(tidytext)}

\NormalTok{bc\_comm }\OtherTok{\textless{}{-}} \FunctionTok{cluster\_leiden}\NormalTok{(g\_bc, }\AttributeTok{resolution =} \FloatTok{1.0}\NormalTok{)}
\FunctionTok{V}\NormalTok{(g\_bc)}\SpecialCharTok{$}\NormalTok{community }\OtherTok{\textless{}{-}} \FunctionTok{membership}\NormalTok{(bc\_comm)}

\NormalTok{front\_labels }\OtherTok{\textless{}{-}}\NormalTok{ works }\SpecialCharTok{|\textgreater{}}
 \FunctionTok{filter}\NormalTok{(id }\SpecialCharTok{\%in\%} \FunctionTok{V}\NormalTok{(g\_bc)}\SpecialCharTok{$}\NormalTok{name) }\SpecialCharTok{|\textgreater{}}
 \FunctionTok{left\_join}\NormalTok{(}
   \FunctionTok{tibble}\NormalTok{(}\AttributeTok{id =} \FunctionTok{V}\NormalTok{(g\_bc)}\SpecialCharTok{$}\NormalTok{name, }\AttributeTok{community =} \FunctionTok{V}\NormalTok{(g\_bc)}\SpecialCharTok{$}\NormalTok{community),}
   \AttributeTok{by =} \StringTok{"id"}
\NormalTok{ ) }\SpecialCharTok{|\textgreater{}}
 \FunctionTok{unnest\_tokens}\NormalTok{(word, display\_name) }\SpecialCharTok{|\textgreater{}}
 \FunctionTok{anti\_join}\NormalTok{(stop\_words, }\AttributeTok{by =} \StringTok{"word"}\NormalTok{) }\SpecialCharTok{|\textgreater{}}
 \FunctionTok{count}\NormalTok{(community, word, }\AttributeTok{sort =} \ConstantTok{TRUE}\NormalTok{) }\SpecialCharTok{|\textgreater{}}
 \FunctionTok{group\_by}\NormalTok{(community) }\SpecialCharTok{|\textgreater{}}
 \FunctionTok{slice\_head}\NormalTok{(}\AttributeTok{n =} \DecValTok{3}\NormalTok{) }\SpecialCharTok{|\textgreater{}}
 \FunctionTok{summarise}\NormalTok{(}\AttributeTok{label =} \FunctionTok{paste}\NormalTok{(word, }\AttributeTok{collapse =} \StringTok{", "}\NormalTok{))}

\NormalTok{front\_labels}
\end{Highlighting}
\end{Shaded}

\emph{Key insight:} Community labels based on frequent title words typically
correspond to recognisable subfields or research themes.

\section*{Chapter 17 --- Co-word and Keyword Co-occurrence}\label{chapter-17-co-word-and-keyword-co-occurrence}
\addcontentsline{toc}{section}{Chapter 17 --- Co-word and Keyword Co-occurrence}

\textbf{1. Author keywords vs.~concepts.}
Build and compare two co-occurrence networks.

\begin{Shaded}
\begin{Highlighting}[]
\FunctionTok{library}\NormalTok{(igraph)}
\FunctionTok{library}\NormalTok{(tidyverse)}

\CommentTok{\# Author keywords network}
\NormalTok{kw\_pairs }\OtherTok{\textless{}{-}}\NormalTok{ works }\SpecialCharTok{|\textgreater{}}
 \FunctionTok{unnest}\NormalTok{(keywords) }\SpecialCharTok{|\textgreater{}}
 \FunctionTok{group\_by}\NormalTok{(id) }\SpecialCharTok{|\textgreater{}}
 \FunctionTok{filter}\NormalTok{(}\FunctionTok{n}\NormalTok{() }\SpecialCharTok{\textgreater{}=} \DecValTok{2}\NormalTok{) }\SpecialCharTok{|\textgreater{}}
 \FunctionTok{reframe}\NormalTok{(}\AttributeTok{pair =} \FunctionTok{combn}\NormalTok{(keyword, }\DecValTok{2}\NormalTok{, }\AttributeTok{simplify =} \ConstantTok{FALSE}\NormalTok{)) }\SpecialCharTok{|\textgreater{}}
 \FunctionTok{unnest\_wider}\NormalTok{(pair, }\AttributeTok{names\_sep =} \StringTok{"\_"}\NormalTok{) }\SpecialCharTok{|\textgreater{}}
 \FunctionTok{count}\NormalTok{(pair\_1, pair\_2, }\AttributeTok{name =} \StringTok{"weight"}\NormalTok{, }\AttributeTok{sort =} \ConstantTok{TRUE}\NormalTok{)}

\NormalTok{g\_kw }\OtherTok{\textless{}{-}} \FunctionTok{graph\_from\_data\_frame}\NormalTok{(kw\_pairs, }\AttributeTok{directed =} \ConstantTok{FALSE}\NormalTok{)}

\CommentTok{\# OpenAlex concepts network}
\NormalTok{concept\_pairs }\OtherTok{\textless{}{-}}\NormalTok{ works }\SpecialCharTok{|\textgreater{}}
 \FunctionTok{unnest}\NormalTok{(concepts) }\SpecialCharTok{|\textgreater{}}
 \FunctionTok{group\_by}\NormalTok{(id) }\SpecialCharTok{|\textgreater{}}
 \FunctionTok{filter}\NormalTok{(}\FunctionTok{n}\NormalTok{() }\SpecialCharTok{\textgreater{}=} \DecValTok{2}\NormalTok{) }\SpecialCharTok{|\textgreater{}}
 \FunctionTok{reframe}\NormalTok{(}\AttributeTok{pair =} \FunctionTok{combn}\NormalTok{(display\_name, }\DecValTok{2}\NormalTok{, }\AttributeTok{simplify =} \ConstantTok{FALSE}\NormalTok{)) }\SpecialCharTok{|\textgreater{}}
 \FunctionTok{unnest\_wider}\NormalTok{(pair, }\AttributeTok{names\_sep =} \StringTok{"\_"}\NormalTok{) }\SpecialCharTok{|\textgreater{}}
 \FunctionTok{count}\NormalTok{(pair\_1, pair\_2, }\AttributeTok{name =} \StringTok{"weight"}\NormalTok{, }\AttributeTok{sort =} \ConstantTok{TRUE}\NormalTok{)}

\NormalTok{g\_concept }\OtherTok{\textless{}{-}} \FunctionTok{graph\_from\_data\_frame}\NormalTok{(concept\_pairs, }\AttributeTok{directed =} \ConstantTok{FALSE}\NormalTok{)}

\FunctionTok{cat}\NormalTok{(}\StringTok{"Author KW network:"}\NormalTok{, }\FunctionTok{vcount}\NormalTok{(g\_kw), }\StringTok{"nodes,"}\NormalTok{,}
    \FunctionTok{ecount}\NormalTok{(g\_kw), }\StringTok{"edges}\SpecialCharTok{\textbackslash{}n}\StringTok{"}\NormalTok{)}
\FunctionTok{cat}\NormalTok{(}\StringTok{"Concept network:"}\NormalTok{, }\FunctionTok{vcount}\NormalTok{(g\_concept), }\StringTok{"nodes,"}\NormalTok{,}
    \FunctionTok{ecount}\NormalTok{(g\_concept), }\StringTok{"edges}\SpecialCharTok{\textbackslash{}n}\StringTok{"}\NormalTok{)}
\end{Highlighting}
\end{Shaded}

\emph{Key insight:} Author keywords are more specific but inconsistent; OpenAlex
concepts are more standardised but coarser.

\textbf{2. Temporal evolution.}
Track the top keywords by degree over time.

\begin{Shaded}
\begin{Highlighting}[]
\NormalTok{yearly\_top\_kw }\OtherTok{\textless{}{-}}\NormalTok{ works }\SpecialCharTok{|\textgreater{}}
 \FunctionTok{unnest}\NormalTok{(keywords) }\SpecialCharTok{|\textgreater{}}
 \FunctionTok{group\_by}\NormalTok{(publication\_year) }\SpecialCharTok{|\textgreater{}}
 \FunctionTok{group\_map}\NormalTok{(\textbackslash{}(df, key) \{}
\NormalTok{   pairs }\OtherTok{\textless{}{-}}\NormalTok{ df }\SpecialCharTok{|\textgreater{}}
     \FunctionTok{group\_by}\NormalTok{(id) }\SpecialCharTok{|\textgreater{}}
     \FunctionTok{filter}\NormalTok{(}\FunctionTok{n}\NormalTok{() }\SpecialCharTok{\textgreater{}=} \DecValTok{2}\NormalTok{) }\SpecialCharTok{|\textgreater{}}
     \FunctionTok{reframe}\NormalTok{(}\AttributeTok{pair =} \FunctionTok{combn}\NormalTok{(keyword, }\DecValTok{2}\NormalTok{, }\AttributeTok{simplify =} \ConstantTok{FALSE}\NormalTok{)) }\SpecialCharTok{|\textgreater{}}
     \FunctionTok{unnest\_wider}\NormalTok{(pair, }\AttributeTok{names\_sep =} \StringTok{"\_"}\NormalTok{)}
   \ControlFlowTok{if}\NormalTok{ (}\FunctionTok{nrow}\NormalTok{(pairs) }\SpecialCharTok{==} \DecValTok{0}\NormalTok{) }\FunctionTok{return}\NormalTok{(}\ConstantTok{NULL}\NormalTok{)}
\NormalTok{   g }\OtherTok{\textless{}{-}} \FunctionTok{graph\_from\_data\_frame}\NormalTok{(pairs, }\AttributeTok{directed =} \ConstantTok{FALSE}\NormalTok{)}
   \FunctionTok{tibble}\NormalTok{(}
     \AttributeTok{year =}\NormalTok{ key}\SpecialCharTok{$}\NormalTok{publication\_year,}
     \AttributeTok{keyword =} \FunctionTok{names}\NormalTok{(}\FunctionTok{sort}\NormalTok{(}\FunctionTok{degree}\NormalTok{(g), }\AttributeTok{decreasing =} \ConstantTok{TRUE}\NormalTok{)[}\DecValTok{1}\SpecialCharTok{:}\DecValTok{5}\NormalTok{]),}
     \AttributeTok{degree =} \FunctionTok{sort}\NormalTok{(}\FunctionTok{degree}\NormalTok{(g), }\AttributeTok{decreasing =} \ConstantTok{TRUE}\NormalTok{)[}\DecValTok{1}\SpecialCharTok{:}\DecValTok{5}\NormalTok{]}
\NormalTok{   )}
\NormalTok{ \}) }\SpecialCharTok{|\textgreater{}} \FunctionTok{list\_rbind}\NormalTok{()}

\FunctionTok{ggplot}\NormalTok{(yearly\_top\_kw, }\FunctionTok{aes}\NormalTok{(year, degree, }\AttributeTok{colour =}\NormalTok{ keyword)) }\SpecialCharTok{+}
 \FunctionTok{geom\_line}\NormalTok{() }\SpecialCharTok{+}
 \FunctionTok{labs}\NormalTok{(}\AttributeTok{title =} \StringTok{"Top keywords by degree over time"}\NormalTok{)}
\end{Highlighting}
\end{Shaded}

\emph{Key insight:} Some keywords are perennially central; others emerge or fade,
signalling evolving research interests.

\textbf{3. Strategic diagram.}
Plot clusters on a density-centrality plane.

\begin{Shaded}
\begin{Highlighting}[]
\NormalTok{comm }\OtherTok{\textless{}{-}} \FunctionTok{cluster\_leiden}\NormalTok{(g\_kw, }\AttributeTok{resolution =} \FloatTok{1.0}\NormalTok{)}
\FunctionTok{V}\NormalTok{(g\_kw)}\SpecialCharTok{$}\NormalTok{community }\OtherTok{\textless{}{-}} \FunctionTok{membership}\NormalTok{(comm)}

\NormalTok{strategic }\OtherTok{\textless{}{-}} \FunctionTok{map}\NormalTok{(}\FunctionTok{unique}\NormalTok{(}\FunctionTok{membership}\NormalTok{(comm)), \textbackslash{}(c) \{}
\NormalTok{ nodes }\OtherTok{\textless{}{-}} \FunctionTok{V}\NormalTok{(g\_kw)[}\FunctionTok{V}\NormalTok{(g\_kw)}\SpecialCharTok{$}\NormalTok{community }\SpecialCharTok{==}\NormalTok{ c]}
\NormalTok{ subg }\OtherTok{\textless{}{-}} \FunctionTok{induced\_subgraph}\NormalTok{(g\_kw, nodes)}
 \FunctionTok{tibble}\NormalTok{(}
   \AttributeTok{community =}\NormalTok{ c,}
   \AttributeTok{density =} \FunctionTok{edge\_density}\NormalTok{(subg),}
   \AttributeTok{centrality =} \FunctionTok{mean}\NormalTok{(}\FunctionTok{strength}\NormalTok{(g\_kw, }\AttributeTok{v =}\NormalTok{ nodes))}
\NormalTok{ )}
\NormalTok{\}) }\SpecialCharTok{|\textgreater{}} \FunctionTok{list\_rbind}\NormalTok{()}

\FunctionTok{ggplot}\NormalTok{(strategic, }\FunctionTok{aes}\NormalTok{(centrality, density,}
                      \AttributeTok{label =}\NormalTok{ community)) }\SpecialCharTok{+}
 \FunctionTok{geom\_point}\NormalTok{(}\AttributeTok{size =} \DecValTok{3}\NormalTok{) }\SpecialCharTok{+}
 \FunctionTok{geom\_text}\NormalTok{(}\AttributeTok{nudge\_y =} \FloatTok{0.01}\NormalTok{) }\SpecialCharTok{+}
 \FunctionTok{geom\_hline}\NormalTok{(}\AttributeTok{yintercept =} \FunctionTok{median}\NormalTok{(strategic}\SpecialCharTok{$}\NormalTok{density),}
            \AttributeTok{linetype =} \StringTok{"dashed"}\NormalTok{) }\SpecialCharTok{+}
 \FunctionTok{geom\_vline}\NormalTok{(}\AttributeTok{xintercept =} \FunctionTok{median}\NormalTok{(strategic}\SpecialCharTok{$}\NormalTok{centrality),}
            \AttributeTok{linetype =} \StringTok{"dashed"}\NormalTok{) }\SpecialCharTok{+}
 \FunctionTok{labs}\NormalTok{(}\AttributeTok{title =} \StringTok{"Strategic diagram"}\NormalTok{,}
      \AttributeTok{x =} \StringTok{"Centrality (external)"}\NormalTok{, }\AttributeTok{y =} \StringTok{"Density (internal)"}\NormalTok{)}
\end{Highlighting}
\end{Shaded}

\emph{Key insight:} Upper-right quadrant clusters are core themes (high density and
centrality); lower-left clusters are peripheral or emerging topics.

\textbf{4. Normalisation comparison.}
Compare raw and association-strength-normalised networks.

\begin{Shaded}
\begin{Highlighting}[]
\CommentTok{\# Association strength: w\_ij / (w\_i * w\_j / 2m)}
\NormalTok{total\_weight }\OtherTok{\textless{}{-}} \FunctionTok{sum}\NormalTok{(}\FunctionTok{E}\NormalTok{(g\_kw)}\SpecialCharTok{$}\NormalTok{weight)}
\NormalTok{node\_strength }\OtherTok{\textless{}{-}} \FunctionTok{strength}\NormalTok{(g\_kw)}

\NormalTok{assoc\_edges }\OtherTok{\textless{}{-}} \FunctionTok{as\_data\_frame}\NormalTok{(g\_kw) }\SpecialCharTok{|\textgreater{}}
 \FunctionTok{mutate}\NormalTok{(}
   \AttributeTok{assoc =}\NormalTok{ weight }\SpecialCharTok{/}\NormalTok{ (node\_strength[from] }\SpecialCharTok{*}\NormalTok{ node\_strength[to] }\SpecialCharTok{/}
\NormalTok{                     (}\DecValTok{2} \SpecialCharTok{*}\NormalTok{ total\_weight))}
\NormalTok{ )}

\NormalTok{g\_assoc }\OtherTok{\textless{}{-}} \FunctionTok{graph\_from\_data\_frame}\NormalTok{(assoc\_edges, }\AttributeTok{directed =} \ConstantTok{FALSE}\NormalTok{)}
\FunctionTok{E}\NormalTok{(g\_assoc)}\SpecialCharTok{$}\NormalTok{weight }\OtherTok{\textless{}{-}}\NormalTok{ assoc\_edges}\SpecialCharTok{$}\NormalTok{assoc}

\CommentTok{\# Compare top{-}centrality keywords}
\NormalTok{top\_raw }\OtherTok{\textless{}{-}} \FunctionTok{names}\NormalTok{(}\FunctionTok{sort}\NormalTok{(}\FunctionTok{strength}\NormalTok{(g\_kw), }\AttributeTok{decreasing =} \ConstantTok{TRUE}\NormalTok{))[}\DecValTok{1}\SpecialCharTok{:}\DecValTok{10}\NormalTok{]}
\NormalTok{top\_assoc }\OtherTok{\textless{}{-}} \FunctionTok{names}\NormalTok{(}\FunctionTok{sort}\NormalTok{(}\FunctionTok{strength}\NormalTok{(g\_assoc), }\AttributeTok{decreasing =} \ConstantTok{TRUE}\NormalTok{))[}\DecValTok{1}\SpecialCharTok{:}\DecValTok{10}\NormalTok{]}

\FunctionTok{tibble}\NormalTok{(}\AttributeTok{raw =}\NormalTok{ top\_raw, }\AttributeTok{normalised =}\NormalTok{ top\_assoc)}
\end{Highlighting}
\end{Shaded}

\emph{Key insight:} Normalisation reduces the dominance of high-frequency generic
terms, elevating more specific and meaningful keyword pairs.

\section*{Chapter 18 --- Community Detection and Backbone Extraction}\label{chapter-18-community-detection-and-backbone-extraction}
\addcontentsline{toc}{section}{Chapter 18 --- Community Detection and Backbone Extraction}

\textbf{1. Alpha sensitivity.}
Run the disparity filter across multiple alpha values.

\begin{Shaded}
\begin{Highlighting}[]
\FunctionTok{library}\NormalTok{(igraph)}
\FunctionTok{library}\NormalTok{(tidyverse)}

\NormalTok{alphas }\OtherTok{\textless{}{-}} \FunctionTok{c}\NormalTok{(}\FloatTok{0.01}\NormalTok{, }\FloatTok{0.05}\NormalTok{, }\FloatTok{0.10}\NormalTok{, }\FloatTok{0.20}\NormalTok{)}

\NormalTok{alpha\_results }\OtherTok{\textless{}{-}} \FunctionTok{map}\NormalTok{(alphas, \textbackslash{}(a) \{}
 \CommentTok{\# Disparity filter: keep edge if p\_ij \textless{} alpha}
\NormalTok{ el }\OtherTok{\textless{}{-}} \FunctionTok{as\_data\_frame}\NormalTok{(g) }\SpecialCharTok{|\textgreater{}}
   \FunctionTok{mutate}\NormalTok{(}
     \AttributeTok{s\_from =} \FunctionTok{strength}\NormalTok{(g)[from],}
     \AttributeTok{k\_from =} \FunctionTok{degree}\NormalTok{(g)[from],}
     \AttributeTok{p\_ij =}\NormalTok{ (}\DecValTok{1} \SpecialCharTok{{-}}\NormalTok{ weight }\SpecialCharTok{/}\NormalTok{ s\_from)}\SpecialCharTok{\^{}}\NormalTok{(k\_from }\SpecialCharTok{{-}} \DecValTok{1}\NormalTok{)}
\NormalTok{   ) }\SpecialCharTok{|\textgreater{}}
   \FunctionTok{filter}\NormalTok{(p\_ij }\SpecialCharTok{\textless{}}\NormalTok{ a)}

\NormalTok{ g\_bb }\OtherTok{\textless{}{-}} \FunctionTok{graph\_from\_data\_frame}\NormalTok{(el, }\AttributeTok{directed =} \ConstantTok{FALSE}\NormalTok{)}
\NormalTok{ comm }\OtherTok{\textless{}{-}} \FunctionTok{cluster\_leiden}\NormalTok{(g\_bb, }\AttributeTok{resolution =} \FloatTok{1.0}\NormalTok{)}
 \FunctionTok{tibble}\NormalTok{(}\AttributeTok{alpha =}\NormalTok{ a,}
        \AttributeTok{edges =} \FunctionTok{ecount}\NormalTok{(g\_bb),}
        \AttributeTok{communities =} \FunctionTok{length}\NormalTok{(comm),}
        \AttributeTok{connected =} \FunctionTok{components}\NormalTok{(g\_bb)}\SpecialCharTok{$}\NormalTok{no }\SpecialCharTok{==} \DecValTok{1}\NormalTok{)}
\NormalTok{\}) }\SpecialCharTok{|\textgreater{}} \FunctionTok{list\_rbind}\NormalTok{()}

\NormalTok{alpha\_results}
\end{Highlighting}
\end{Shaded}

\emph{Key insight:} Very low alpha values remove too many edges and disconnect the
network; very high values retain noise. Typical values of 0.05-0.10 balance
signal and connectivity.

\textbf{2. Threshold vs.~disparity.}
Compare two backbone extraction methods using NMI.

\begin{Shaded}
\begin{Highlighting}[]
\FunctionTok{library}\NormalTok{(igraph)}

\CommentTok{\# Threshold filtering: keep edges above median weight}
\NormalTok{threshold }\OtherTok{\textless{}{-}} \FunctionTok{median}\NormalTok{(}\FunctionTok{E}\NormalTok{(g)}\SpecialCharTok{$}\NormalTok{weight)}
\NormalTok{g\_thresh }\OtherTok{\textless{}{-}} \FunctionTok{subgraph.edges}\NormalTok{(g, }\FunctionTok{which}\NormalTok{(}\FunctionTok{E}\NormalTok{(g)}\SpecialCharTok{$}\NormalTok{weight }\SpecialCharTok{\textgreater{}}\NormalTok{ threshold))}
\NormalTok{comm\_thresh }\OtherTok{\textless{}{-}} \FunctionTok{cluster\_leiden}\NormalTok{(g\_thresh, }\AttributeTok{resolution =} \FloatTok{1.0}\NormalTok{)}

\CommentTok{\# Disparity filtering (from exercise 1)}
\NormalTok{comm\_disp }\OtherTok{\textless{}{-}} \FunctionTok{cluster\_leiden}\NormalTok{(g\_bb, }\AttributeTok{resolution =} \FloatTok{1.0}\NormalTok{)}

\CommentTok{\# NMI comparison on shared nodes}
\NormalTok{shared }\OtherTok{\textless{}{-}} \FunctionTok{intersect}\NormalTok{(}\FunctionTok{V}\NormalTok{(g\_thresh)}\SpecialCharTok{$}\NormalTok{name, }\FunctionTok{V}\NormalTok{(g\_bb)}\SpecialCharTok{$}\NormalTok{name)}
\NormalTok{m1 }\OtherTok{\textless{}{-}} \FunctionTok{membership}\NormalTok{(comm\_thresh)[shared]}
\NormalTok{m2 }\OtherTok{\textless{}{-}} \FunctionTok{membership}\NormalTok{(comm\_disp)[shared]}

\FunctionTok{compare}\NormalTok{(m1, m2, }\AttributeTok{method =} \StringTok{"nmi"}\NormalTok{)}
\end{Highlighting}
\end{Shaded}

\emph{Key insight:} The disparity filter preserves locally significant edges
(important for low-degree nodes) while the threshold filter favours globally
strong connections. NMI values below 0.7 indicate meaningfully different
community structures.

\textbf{3. Hierarchical structure.}
Check whether communities at different resolutions are nested.

\begin{Shaded}
\begin{Highlighting}[]
\NormalTok{resolutions }\OtherTok{\textless{}{-}} \FunctionTok{c}\NormalTok{(}\FloatTok{0.5}\NormalTok{, }\FloatTok{1.0}\NormalTok{, }\FloatTok{2.0}\NormalTok{)}

\NormalTok{memberships }\OtherTok{\textless{}{-}} \FunctionTok{map}\NormalTok{(resolutions, \textbackslash{}(r) \{}
\NormalTok{ comm }\OtherTok{\textless{}{-}} \FunctionTok{cluster\_leiden}\NormalTok{(g\_bb, }\AttributeTok{resolution =}\NormalTok{ r)}
 \FunctionTok{tibble}\NormalTok{(}\AttributeTok{node =} \FunctionTok{V}\NormalTok{(g\_bb)}\SpecialCharTok{$}\NormalTok{name, }\AttributeTok{community =} \FunctionTok{membership}\NormalTok{(comm),}
        \AttributeTok{resolution =}\NormalTok{ r)}
\NormalTok{\}) }\SpecialCharTok{|\textgreater{}} \FunctionTok{list\_rbind}\NormalTok{()}

\CommentTok{\# Check nesting: at resolution 2.0, are communities subsets of}
\CommentTok{\# resolution 0.5 communities?}
\NormalTok{nesting }\OtherTok{\textless{}{-}}\NormalTok{ memberships }\SpecialCharTok{|\textgreater{}}
 \FunctionTok{pivot\_wider}\NormalTok{(}\AttributeTok{names\_from =}\NormalTok{ resolution, }\AttributeTok{values\_from =}\NormalTok{ community,}
             \AttributeTok{names\_prefix =} \StringTok{"res\_"}\NormalTok{) }\SpecialCharTok{|\textgreater{}}
 \FunctionTok{count}\NormalTok{(res\_0}\FloatTok{.5}\NormalTok{, res\_2) }\SpecialCharTok{|\textgreater{}}
 \FunctionTok{group\_by}\NormalTok{(res\_2) }\SpecialCharTok{|\textgreater{}}
 \FunctionTok{mutate}\NormalTok{(}\AttributeTok{is\_nested =} \FunctionTok{n\_distinct}\NormalTok{(res\_0}\FloatTok{.5}\NormalTok{) }\SpecialCharTok{==} \DecValTok{1}\NormalTok{)}

\FunctionTok{mean}\NormalTok{(nesting}\SpecialCharTok{$}\NormalTok{is\_nested)  }\CommentTok{\# proportion of fine communities nested}
\end{Highlighting}
\end{Shaded}

\emph{Key insight:} Perfectly nested hierarchies are rare in real networks.
Partial nesting suggests a hierarchical thematic structure with some cross-
cutting topics.

\textbf{4. Weighted vs.~unweighted community detection.}
Remove weights and compare results.

\begin{Shaded}
\begin{Highlighting}[]
\NormalTok{g\_unweighted }\OtherTok{\textless{}{-}}\NormalTok{ g\_bb}
\FunctionTok{E}\NormalTok{(g\_unweighted)}\SpecialCharTok{$}\NormalTok{weight }\OtherTok{\textless{}{-}} \DecValTok{1}

\NormalTok{comm\_weighted   }\OtherTok{\textless{}{-}} \FunctionTok{cluster\_leiden}\NormalTok{(g\_bb, }\AttributeTok{resolution =} \FloatTok{1.0}\NormalTok{)}
\NormalTok{comm\_unweighted }\OtherTok{\textless{}{-}} \FunctionTok{cluster\_leiden}\NormalTok{(g\_unweighted, }\AttributeTok{resolution =} \FloatTok{1.0}\NormalTok{)}

\FunctionTok{compare}\NormalTok{(}\FunctionTok{membership}\NormalTok{(comm\_weighted), }\FunctionTok{membership}\NormalTok{(comm\_unweighted),}
        \AttributeTok{method =} \StringTok{"nmi"}\NormalTok{)}

\FunctionTok{cat}\NormalTok{(}\StringTok{"Weighted communities:"}\NormalTok{, }\FunctionTok{length}\NormalTok{(comm\_weighted), }\StringTok{"}\SpecialCharTok{\textbackslash{}n}\StringTok{"}\NormalTok{)}
\FunctionTok{cat}\NormalTok{(}\StringTok{"Unweighted communities:"}\NormalTok{, }\FunctionTok{length}\NormalTok{(comm\_unweighted), }\StringTok{"}\SpecialCharTok{\textbackslash{}n}\StringTok{"}\NormalTok{)}
\end{Highlighting}
\end{Shaded}

\emph{Key insight:} Removing weights often increases the number of communities because
weak connections that helped merge groups are no longer distinguishable from
strong ones.

\section*{Chapter 19 --- Science Mapping and Overlay Maps}\label{chapter-19-science-mapping-and-overlay-maps}
\addcontentsline{toc}{section}{Chapter 19 --- Science Mapping and Overlay Maps}

\textbf{1. Multi-institution overlay.}
Overlay two institutions on the same base map.

\begin{Shaded}
\begin{Highlighting}[]
\FunctionTok{library}\NormalTok{(tidyverse)}
\FunctionTok{library}\NormalTok{(ggrepel)}

\NormalTok{inst\_ids }\OtherTok{\textless{}{-}} \FunctionTok{c}\NormalTok{(}\StringTok{"I123456789"}\NormalTok{, }\StringTok{"I987654321"}\NormalTok{)}

\NormalTok{profiles }\OtherTok{\textless{}{-}} \FunctionTok{map}\NormalTok{(inst\_ids, \textbackslash{}(iid) \{}
\NormalTok{ w }\OtherTok{\textless{}{-}} \FunctionTok{fetch\_openalex}\NormalTok{(}\AttributeTok{entity =} \StringTok{"works"}\NormalTok{,}
               \AttributeTok{authorships.institutions.lineage =}\NormalTok{ iid,}
               \AttributeTok{from\_publication\_date =} \StringTok{"2020{-}01{-}01"}\NormalTok{)}
\NormalTok{ w }\SpecialCharTok{|\textgreater{}}
   \FunctionTok{unnest}\NormalTok{(concepts) }\SpecialCharTok{|\textgreater{}}
   \FunctionTok{count}\NormalTok{(display\_name, }\AttributeTok{name =} \StringTok{"n\_works"}\NormalTok{) }\SpecialCharTok{|\textgreater{}}
   \FunctionTok{mutate}\NormalTok{(}\AttributeTok{inst =}\NormalTok{ iid)}
\NormalTok{\}) }\SpecialCharTok{|\textgreater{}} \FunctionTok{list\_rbind}\NormalTok{()}

\FunctionTok{ggplot}\NormalTok{(profiles, }\FunctionTok{aes}\NormalTok{(display\_name, n\_works, }\AttributeTok{fill =}\NormalTok{ inst)) }\SpecialCharTok{+}
 \FunctionTok{geom\_col}\NormalTok{(}\AttributeTok{position =} \StringTok{"dodge"}\NormalTok{) }\SpecialCharTok{+}
 \FunctionTok{coord\_flip}\NormalTok{() }\SpecialCharTok{+}
 \FunctionTok{labs}\NormalTok{(}\AttributeTok{title =} \StringTok{"Research profile overlay"}\NormalTok{)}
\end{Highlighting}
\end{Shaded}

\emph{Key insight:} Overlapping areas indicate shared research strengths; divergent
areas reveal institutional specialisation.

\textbf{2. Temporal overlay.}
Compare the same institution's profile across two time periods.

\begin{Shaded}
\begin{Highlighting}[]
\NormalTok{periods }\OtherTok{\textless{}{-}} \FunctionTok{list}\NormalTok{(}
 \AttributeTok{early =} \FunctionTok{c}\NormalTok{(}\StringTok{"2018{-}01{-}01"}\NormalTok{, }\StringTok{"2019{-}12{-}31"}\NormalTok{),}
 \AttributeTok{late  =} \FunctionTok{c}\NormalTok{(}\StringTok{"2022{-}01{-}01"}\NormalTok{, }\StringTok{"2023{-}12{-}31"}\NormalTok{)}
\NormalTok{)}

\NormalTok{temporal\_profiles }\OtherTok{\textless{}{-}} \FunctionTok{imap}\NormalTok{(periods, \textbackslash{}(dates, label) \{}
 \FunctionTok{fetch\_openalex}\NormalTok{(}\AttributeTok{entity =} \StringTok{"works"}\NormalTok{,}
          \AttributeTok{authorships.institutions.lineage =} \StringTok{"I123456789"}\NormalTok{,}
          \AttributeTok{from\_publication\_date =}\NormalTok{ dates[}\DecValTok{1}\NormalTok{],}
          \AttributeTok{to\_publication\_date   =}\NormalTok{ dates[}\DecValTok{2}\NormalTok{]) }\SpecialCharTok{|\textgreater{}}
   \FunctionTok{unnest}\NormalTok{(concepts) }\SpecialCharTok{|\textgreater{}}
   \FunctionTok{count}\NormalTok{(display\_name, }\AttributeTok{name =} \StringTok{"n\_works"}\NormalTok{) }\SpecialCharTok{|\textgreater{}}
   \FunctionTok{mutate}\NormalTok{(}\AttributeTok{period =}\NormalTok{ label)}
\NormalTok{\}) }\SpecialCharTok{|\textgreater{}} \FunctionTok{list\_rbind}\NormalTok{()}

\NormalTok{temporal\_profiles }\SpecialCharTok{|\textgreater{}}
 \FunctionTok{pivot\_wider}\NormalTok{(}\AttributeTok{names\_from =}\NormalTok{ period, }\AttributeTok{values\_from =}\NormalTok{ n\_works,}
             \AttributeTok{values\_fill =} \DecValTok{0}\NormalTok{) }\SpecialCharTok{|\textgreater{}}
 \FunctionTok{mutate}\NormalTok{(}\AttributeTok{shift =}\NormalTok{ late }\SpecialCharTok{{-}}\NormalTok{ early) }\SpecialCharTok{|\textgreater{}}
 \FunctionTok{arrange}\NormalTok{(}\FunctionTok{desc}\NormalTok{(}\FunctionTok{abs}\NormalTok{(shift))) }\SpecialCharTok{|\textgreater{}}
 \FunctionTok{head}\NormalTok{(}\DecValTok{15}\NormalTok{)}
\end{Highlighting}
\end{Shaded}

\emph{Key insight:} Positive shifts indicate growing research areas; negative shifts
suggest declining focus or strategic reorientation.

\textbf{3. Country-level overlay.}
Compare a small and a large country's publication profiles.

\begin{Shaded}
\begin{Highlighting}[]
\NormalTok{countries }\OtherTok{\textless{}{-}} \FunctionTok{c}\NormalTok{(}\StringTok{"SG"}\NormalTok{, }\StringTok{"US"}\NormalTok{)  }\CommentTok{\# Singapore vs. United States}

\NormalTok{country\_profiles }\OtherTok{\textless{}{-}} \FunctionTok{map}\NormalTok{(countries, \textbackslash{}(cc) \{}
 \FunctionTok{fetch\_openalex}\NormalTok{(}\AttributeTok{entity =} \StringTok{"works"}\NormalTok{,}
          \AttributeTok{authorships.countries =}\NormalTok{ cc,}
          \AttributeTok{from\_publication\_date =} \StringTok{"2022{-}01{-}01"}\NormalTok{,}
          \AttributeTok{to\_publication\_date =} \StringTok{"2022{-}12{-}31"}\NormalTok{) }\SpecialCharTok{|\textgreater{}}
   \FunctionTok{slice\_sample}\NormalTok{(}\AttributeTok{n =} \DecValTok{500}\NormalTok{) }\SpecialCharTok{|\textgreater{}}
   \FunctionTok{unnest}\NormalTok{(concepts) }\SpecialCharTok{|\textgreater{}}
   \FunctionTok{count}\NormalTok{(display\_name, }\AttributeTok{name =} \StringTok{"n\_works"}\NormalTok{) }\SpecialCharTok{|\textgreater{}}
   \FunctionTok{mutate}\NormalTok{(}\AttributeTok{share =}\NormalTok{ n\_works }\SpecialCharTok{/} \FunctionTok{sum}\NormalTok{(n\_works),}
          \AttributeTok{country =}\NormalTok{ cc)}
\NormalTok{\}) }\SpecialCharTok{|\textgreater{}} \FunctionTok{list\_rbind}\NormalTok{()}

\NormalTok{country\_profiles }\SpecialCharTok{|\textgreater{}}
 \FunctionTok{pivot\_wider}\NormalTok{(}\AttributeTok{names\_from =}\NormalTok{ country, }\AttributeTok{values\_from =}\NormalTok{ share,}
             \AttributeTok{values\_fill =} \DecValTok{0}\NormalTok{) }\SpecialCharTok{|\textgreater{}}
 \FunctionTok{ggplot}\NormalTok{(}\FunctionTok{aes}\NormalTok{(US, SG)) }\SpecialCharTok{+}
 \FunctionTok{geom\_point}\NormalTok{() }\SpecialCharTok{+}
 \FunctionTok{geom\_abline}\NormalTok{(}\AttributeTok{slope =} \DecValTok{1}\NormalTok{, }\AttributeTok{intercept =} \DecValTok{0}\NormalTok{, }\AttributeTok{linetype =} \StringTok{"dashed"}\NormalTok{) }\SpecialCharTok{+}
\NormalTok{ ggrepel}\SpecialCharTok{::}\FunctionTok{geom\_text\_repel}\NormalTok{(}\FunctionTok{aes}\NormalTok{(}\AttributeTok{label =}\NormalTok{ display\_name), }\AttributeTok{size =} \DecValTok{2}\NormalTok{) }\SpecialCharTok{+}
 \FunctionTok{labs}\NormalTok{(}\AttributeTok{title =} \StringTok{"Country profile: Singapore vs. US"}\NormalTok{)}
\end{Highlighting}
\end{Shaded}

\emph{Key insight:} Small countries tend to show concentrated profiles (above the
diagonal in selected areas); large countries are more evenly distributed.

\section*{Chapter 20 --- Network Visualization and Interoperability}\label{chapter-20-network-visualization-and-interoperability}
\addcontentsline{toc}{section}{Chapter 20 --- Network Visualization and Interoperability}

\textbf{1. Layout comparison.}
Apply five layouts and compare quality and speed.

\begin{Shaded}
\begin{Highlighting}[]
\FunctionTok{library}\NormalTok{(igraph)}
\FunctionTok{library}\NormalTok{(ggraph)}
\FunctionTok{library}\NormalTok{(tidyverse)}

\NormalTok{layouts }\OtherTok{\textless{}{-}} \FunctionTok{c}\NormalTok{(}\StringTok{"fr"}\NormalTok{, }\StringTok{"kk"}\NormalTok{, }\StringTok{"drl"}\NormalTok{, }\StringTok{"graphopt"}\NormalTok{, }\StringTok{"lgl"}\NormalTok{)}

\NormalTok{layout\_results }\OtherTok{\textless{}{-}} \FunctionTok{map}\NormalTok{(layouts, \textbackslash{}(algo) \{}
\NormalTok{ t }\OtherTok{\textless{}{-}} \FunctionTok{system.time}\NormalTok{(\{}
\NormalTok{   p }\OtherTok{\textless{}{-}} \FunctionTok{ggraph}\NormalTok{(gc, }\AttributeTok{layout =}\NormalTok{ algo) }\SpecialCharTok{+}
     \FunctionTok{geom\_edge\_link}\NormalTok{(}\AttributeTok{alpha =} \FloatTok{0.1}\NormalTok{) }\SpecialCharTok{+}
     \FunctionTok{geom\_node\_point}\NormalTok{(}\FunctionTok{aes}\NormalTok{(}\AttributeTok{colour =} \FunctionTok{factor}\NormalTok{(community)), }\AttributeTok{size =} \DecValTok{1}\NormalTok{) }\SpecialCharTok{+}
     \FunctionTok{labs}\NormalTok{(}\AttributeTok{title =}\NormalTok{ algo) }\SpecialCharTok{+}
     \FunctionTok{theme\_void}\NormalTok{()}
\NormalTok{ \})}
 \FunctionTok{list}\NormalTok{(}\AttributeTok{plot =}\NormalTok{ p, }\AttributeTok{time =}\NormalTok{ t[}\StringTok{"elapsed"}\NormalTok{], }\AttributeTok{layout =}\NormalTok{ algo)}
\NormalTok{\})}

\CommentTok{\# Print timing}
\FunctionTok{map}\NormalTok{(layout\_results, \textbackslash{}(x) }\FunctionTok{tibble}\NormalTok{(}\AttributeTok{layout =}\NormalTok{ x}\SpecialCharTok{$}\NormalTok{layout, }\AttributeTok{secs =}\NormalTok{ x}\SpecialCharTok{$}\NormalTok{time)) }\SpecialCharTok{|\textgreater{}}
 \FunctionTok{list\_rbind}\NormalTok{()}
\end{Highlighting}
\end{Shaded}

\emph{Key insight:} Fruchterman-Reingold is a good default; Kamada-Kawai is better
for small networks; DRL scales to large graphs but can produce chaotic layouts.

\textbf{2. Interactive filtering.}
Create a filterable \texttt{visNetwork} visualisation.

\begin{Shaded}
\begin{Highlighting}[]
\FunctionTok{library}\NormalTok{(visNetwork)}

\NormalTok{nodes\_df }\OtherTok{\textless{}{-}} \FunctionTok{tibble}\NormalTok{(}
 \AttributeTok{id =} \FunctionTok{V}\NormalTok{(gc)}\SpecialCharTok{$}\NormalTok{name,}
 \AttributeTok{label =} \FunctionTok{V}\NormalTok{(gc)}\SpecialCharTok{$}\NormalTok{name,}
 \AttributeTok{group =} \FunctionTok{as.character}\NormalTok{(}\FunctionTok{V}\NormalTok{(gc)}\SpecialCharTok{$}\NormalTok{community)}
\NormalTok{)}

\NormalTok{edges\_df }\OtherTok{\textless{}{-}} \FunctionTok{as\_data\_frame}\NormalTok{(gc) }\SpecialCharTok{|\textgreater{}}
 \FunctionTok{select}\NormalTok{(from, to)}

\FunctionTok{visNetwork}\NormalTok{(nodes\_df, edges\_df) }\SpecialCharTok{|\textgreater{}}
 \FunctionTok{visOptions}\NormalTok{(}\AttributeTok{selectedBy =} \StringTok{"group"}\NormalTok{,}
            \AttributeTok{highlightNearest =} \ConstantTok{TRUE}\NormalTok{) }\SpecialCharTok{|\textgreater{}}
 \FunctionTok{visPhysics}\NormalTok{(}\AttributeTok{stabilization =} \ConstantTok{FALSE}\NormalTok{)}
\end{Highlighting}
\end{Shaded}

\emph{Key insight:} Interactive filtering lets users focus on specific communities
and explore inter-community bridges.

\textbf{3. Gephi round-trip.}
Export to GraphML for use in Gephi.

\begin{Shaded}
\begin{Highlighting}[]
\FunctionTok{library}\NormalTok{(igraph)}

\CommentTok{\# Export}
\FunctionTok{write\_graph}\NormalTok{(gc, }\StringTok{"network.graphml"}\NormalTok{, }\AttributeTok{format =} \StringTok{"graphml"}\NormalTok{)}

\CommentTok{\# After Gephi processing, compare:}
\CommentTok{\# {-} Force Atlas 2 layout often reveals community separation better}
\CommentTok{\# {-} R/ggraph offers more reproducible, scriptable figures}
\CommentTok{\# {-} Gephi is better for interactive exploration and publication}
\CommentTok{\#   quality figures via SVG/PDF export}
\end{Highlighting}
\end{Shaded}

\emph{Key insight:} Gephi's Force Atlas 2 often produces cleaner community
separation; R workflows are more reproducible and automatable.

\textbf{4. Colour-blind check.}
Test colour distinguishability with increasing community counts.

\begin{Shaded}
\begin{Highlighting}[]
\FunctionTok{library}\NormalTok{(ggraph)}
\FunctionTok{library}\NormalTok{(viridis)}

\NormalTok{n\_communities }\OtherTok{\textless{}{-}} \FunctionTok{c}\NormalTok{(}\DecValTok{3}\NormalTok{, }\DecValTok{5}\NormalTok{, }\DecValTok{10}\NormalTok{)}

\NormalTok{plots }\OtherTok{\textless{}{-}} \FunctionTok{map}\NormalTok{(n\_communities, \textbackslash{}(k) \{}
\NormalTok{ comm }\OtherTok{\textless{}{-}} \FunctionTok{cluster\_leiden}\NormalTok{(gc, }\AttributeTok{resolution =}\NormalTok{ k }\SpecialCharTok{/} \DecValTok{5}\NormalTok{)}
 \FunctionTok{V}\NormalTok{(gc)}\SpecialCharTok{$}\NormalTok{comm\_k }\OtherTok{\textless{}{-}} \FunctionTok{membership}\NormalTok{(comm)}

 \FunctionTok{ggraph}\NormalTok{(gc, }\AttributeTok{layout =} \StringTok{"fr"}\NormalTok{) }\SpecialCharTok{+}
   \FunctionTok{geom\_edge\_link}\NormalTok{(}\AttributeTok{alpha =} \FloatTok{0.05}\NormalTok{) }\SpecialCharTok{+}
   \FunctionTok{geom\_node\_point}\NormalTok{(}\FunctionTok{aes}\NormalTok{(}\AttributeTok{colour =} \FunctionTok{factor}\NormalTok{(comm\_k)), }\AttributeTok{size =} \DecValTok{2}\NormalTok{) }\SpecialCharTok{+}
   \FunctionTok{scale\_colour\_viridis\_d}\NormalTok{(}\AttributeTok{option =} \StringTok{"D"}\NormalTok{) }\SpecialCharTok{+}
   \FunctionTok{labs}\NormalTok{(}\AttributeTok{title =} \FunctionTok{paste}\NormalTok{(k, }\StringTok{"communities"}\NormalTok{)) }\SpecialCharTok{+}
   \FunctionTok{theme\_void}\NormalTok{() }\SpecialCharTok{+}
   \FunctionTok{theme}\NormalTok{(}\AttributeTok{legend.position =} \StringTok{"none"}\NormalTok{)}
\NormalTok{\})}

\CommentTok{\# Use colorBlindness package or online simulator to test}
\CommentTok{\# library(colorBlindness)}
\CommentTok{\# cvdPlot(plots[[3]])}
\end{Highlighting}
\end{Shaded}

\emph{Key insight:} Viridis palettes remain distinguishable up to about 8 categories;
beyond that, consider shape or faceting as supplementary channels.

\section*{Chapter 21 --- Text Mining Bibliographic Corpora}\label{chapter-21-text-mining-bibliographic-corpora}
\addcontentsline{toc}{section}{Chapter 21 --- Text Mining Bibliographic Corpora}

\textbf{1. Bigram analysis.}
Build a DFM from bigrams.

\begin{Shaded}
\begin{Highlighting}[]
\FunctionTok{library}\NormalTok{(quanteda)}
\FunctionTok{library}\NormalTok{(tidyverse)}

\NormalTok{corp }\OtherTok{\textless{}{-}} \FunctionTok{corpus}\NormalTok{(works, }\AttributeTok{text\_field =} \StringTok{"ab"}\NormalTok{)}
\NormalTok{toks }\OtherTok{\textless{}{-}} \FunctionTok{tokens}\NormalTok{(corp, }\AttributeTok{remove\_punct =} \ConstantTok{TRUE}\NormalTok{, }\AttributeTok{remove\_numbers =} \ConstantTok{TRUE}\NormalTok{) }\SpecialCharTok{|\textgreater{}}
 \FunctionTok{tokens\_tolower}\NormalTok{() }\SpecialCharTok{|\textgreater{}}
 \FunctionTok{tokens\_remove}\NormalTok{(}\FunctionTok{stopwords}\NormalTok{(}\StringTok{"en"}\NormalTok{)) }\SpecialCharTok{|\textgreater{}}
 \FunctionTok{tokens\_ngrams}\NormalTok{(}\AttributeTok{n =} \DecValTok{2}\NormalTok{)}

\NormalTok{dfm\_bi }\OtherTok{\textless{}{-}} \FunctionTok{dfm}\NormalTok{(toks) }\SpecialCharTok{|\textgreater{}} \FunctionTok{dfm\_trim}\NormalTok{(}\AttributeTok{min\_termfreq =} \DecValTok{5}\NormalTok{)}

\FunctionTok{topfeatures}\NormalTok{(dfm\_bi, }\DecValTok{20}\NormalTok{)}
\end{Highlighting}
\end{Shaded}

\emph{Key insight:} Bigrams capture meaningful phrases like ``machine\_learning'' or
``citation\_analysis'' that single words miss.

\textbf{2. Lexical diversity.}
Compute the type-token ratio (TTR) per year.

\begin{Shaded}
\begin{Highlighting}[]
\FunctionTok{library}\NormalTok{(quanteda)}
\FunctionTok{library}\NormalTok{(quanteda.textstats)}

\NormalTok{ttr\_by\_year }\OtherTok{\textless{}{-}}\NormalTok{ works }\SpecialCharTok{|\textgreater{}}
 \FunctionTok{filter}\NormalTok{(}\SpecialCharTok{!}\FunctionTok{is.na}\NormalTok{(ab)) }\SpecialCharTok{|\textgreater{}}
 \FunctionTok{group\_by}\NormalTok{(publication\_year) }\SpecialCharTok{|\textgreater{}}
 \FunctionTok{summarise}\NormalTok{(}
   \AttributeTok{all\_text =} \FunctionTok{paste}\NormalTok{(ab, }\AttributeTok{collapse =} \StringTok{" "}\NormalTok{),}
   \AttributeTok{.groups =} \StringTok{"drop"}
\NormalTok{ ) }\SpecialCharTok{|\textgreater{}}
 \FunctionTok{mutate}\NormalTok{(}
   \AttributeTok{n\_types  =} \FunctionTok{map\_int}\NormalTok{(all\_text, \textbackslash{}(t) \{}
\NormalTok{     toks }\OtherTok{\textless{}{-}} \FunctionTok{tokens}\NormalTok{(t, }\AttributeTok{remove\_punct =} \ConstantTok{TRUE}\NormalTok{) }\SpecialCharTok{|\textgreater{}} \FunctionTok{tokens\_tolower}\NormalTok{()}
     \FunctionTok{ntype}\NormalTok{(}\FunctionTok{dfm}\NormalTok{(toks))}
\NormalTok{   \}),}
   \AttributeTok{n\_tokens =} \FunctionTok{map\_int}\NormalTok{(all\_text, \textbackslash{}(t) \{}
\NormalTok{     toks }\OtherTok{\textless{}{-}} \FunctionTok{tokens}\NormalTok{(t, }\AttributeTok{remove\_punct =} \ConstantTok{TRUE}\NormalTok{)}
     \FunctionTok{ntoken}\NormalTok{(toks)}
\NormalTok{   \}),}
   \AttributeTok{ttr =}\NormalTok{ n\_types }\SpecialCharTok{/}\NormalTok{ n\_tokens}
\NormalTok{ )}

\FunctionTok{ggplot}\NormalTok{(ttr\_by\_year, }\FunctionTok{aes}\NormalTok{(publication\_year, ttr)) }\SpecialCharTok{+}
 \FunctionTok{geom\_line}\NormalTok{() }\SpecialCharTok{+}
 \FunctionTok{labs}\NormalTok{(}\AttributeTok{title =} \StringTok{"Lexical diversity over time"}\NormalTok{,}
      \AttributeTok{x =} \StringTok{"Year"}\NormalTok{, }\AttributeTok{y =} \StringTok{"Type{-}token ratio"}\NormalTok{)}
\end{Highlighting}
\end{Shaded}

\emph{Key insight:} Declining TTR may indicate vocabulary standardisation; increasing
TTR suggests diversifying terminology.

\textbf{3. Keyness analysis.}
Identify terms that distinguish a specific year.

\begin{Shaded}
\begin{Highlighting}[]
\FunctionTok{library}\NormalTok{(quanteda)}
\FunctionTok{library}\NormalTok{(quanteda.textstats)}

\NormalTok{corp }\OtherTok{\textless{}{-}} \FunctionTok{corpus}\NormalTok{(works }\SpecialCharTok{|\textgreater{}} \FunctionTok{filter}\NormalTok{(}\SpecialCharTok{!}\FunctionTok{is.na}\NormalTok{(ab)),}
               \AttributeTok{text\_field =} \StringTok{"ab"}\NormalTok{,}
               \AttributeTok{docid\_field =} \StringTok{"id"}\NormalTok{)}
\FunctionTok{docvars}\NormalTok{(corp, }\StringTok{"year"}\NormalTok{) }\OtherTok{\textless{}{-}}\NormalTok{ works}\SpecialCharTok{$}\NormalTok{publication\_year[}\SpecialCharTok{!}\FunctionTok{is.na}\NormalTok{(works}\SpecialCharTok{$}\NormalTok{ab)]}

\NormalTok{dfm\_all }\OtherTok{\textless{}{-}} \FunctionTok{dfm}\NormalTok{(}\FunctionTok{tokens}\NormalTok{(corp, }\AttributeTok{remove\_punct =} \ConstantTok{TRUE}\NormalTok{) }\SpecialCharTok{|\textgreater{}}
               \FunctionTok{tokens\_tolower}\NormalTok{() }\SpecialCharTok{|\textgreater{}}
               \FunctionTok{tokens\_remove}\NormalTok{(}\FunctionTok{stopwords}\NormalTok{(}\StringTok{"en"}\NormalTok{)))}

\CommentTok{\# Compare 2023 vs. all other years}
\NormalTok{dfm\_grouped }\OtherTok{\textless{}{-}} \FunctionTok{dfm\_group}\NormalTok{(dfm\_all, }\AttributeTok{groups =} \FunctionTok{ifelse}\NormalTok{(}
 \FunctionTok{docvars}\NormalTok{(dfm\_all, }\StringTok{"year"}\NormalTok{) }\SpecialCharTok{==} \DecValTok{2023}\NormalTok{, }\StringTok{"2023"}\NormalTok{, }\StringTok{"other"}
\NormalTok{))}

\NormalTok{keyness }\OtherTok{\textless{}{-}} \FunctionTok{textstat\_keyness}\NormalTok{(dfm\_grouped, }\AttributeTok{target =} \StringTok{"2023"}\NormalTok{)}
\FunctionTok{textplot\_keyness}\NormalTok{(keyness, }\AttributeTok{n =} \DecValTok{20}\NormalTok{)}
\end{Highlighting}
\end{Shaded}

\emph{Key insight:} Keyness analysis surfaces terms unique to a period. For 2023,
terms related to LLMs and AI may dominate.

\textbf{4. Preprocessing sensitivity.}
Compare analyses with and without stemming.

\begin{Shaded}
\begin{Highlighting}[]
\NormalTok{toks\_nostem }\OtherTok{\textless{}{-}} \FunctionTok{tokens}\NormalTok{(corp, }\AttributeTok{remove\_punct =} \ConstantTok{TRUE}\NormalTok{) }\SpecialCharTok{|\textgreater{}}
 \FunctionTok{tokens\_tolower}\NormalTok{() }\SpecialCharTok{|\textgreater{}}
 \FunctionTok{tokens\_remove}\NormalTok{(}\FunctionTok{stopwords}\NormalTok{(}\StringTok{"en"}\NormalTok{))}

\NormalTok{toks\_stem }\OtherTok{\textless{}{-}} \FunctionTok{tokens\_wordstem}\NormalTok{(toks\_nostem)}

\NormalTok{top\_nostem }\OtherTok{\textless{}{-}} \FunctionTok{topfeatures}\NormalTok{(}\FunctionTok{dfm}\NormalTok{(toks\_nostem), }\DecValTok{20}\NormalTok{)}
\NormalTok{top\_stem   }\OtherTok{\textless{}{-}} \FunctionTok{topfeatures}\NormalTok{(}\FunctionTok{dfm}\NormalTok{(toks\_stem), }\DecValTok{20}\NormalTok{)}

\FunctionTok{tibble}\NormalTok{(}
 \AttributeTok{rank =} \DecValTok{1}\SpecialCharTok{:}\DecValTok{20}\NormalTok{,}
 \AttributeTok{unstemmed =} \FunctionTok{names}\NormalTok{(top\_nostem),}
 \AttributeTok{stemmed =} \FunctionTok{names}\NormalTok{(top\_stem)}
\NormalTok{)}
\end{Highlighting}
\end{Shaded}

\emph{Key insight:} Stemming merges variants (e.g., ``citations'' and ``citation'') but
can produce unintelligible stems. For keyword-based analyses, unstemmed terms
are often more interpretable.

\section*{Chapter 22 --- Topic Modeling}\label{chapter-22-topic-modeling}
\addcontentsline{toc}{section}{Chapter 22 --- Topic Modeling}

\textbf{1. K selection.}
Fit LDA models with multiple values of K and compare coherence.

\begin{Shaded}
\begin{Highlighting}[]
\FunctionTok{library}\NormalTok{(topicmodels)}
\FunctionTok{library}\NormalTok{(tidyverse)}
\FunctionTok{set.seed}\NormalTok{(}\DecValTok{20260509}\NormalTok{)}

\NormalTok{ks }\OtherTok{\textless{}{-}} \FunctionTok{c}\NormalTok{(}\DecValTok{5}\NormalTok{, }\DecValTok{8}\NormalTok{, }\DecValTok{12}\NormalTok{, }\DecValTok{15}\NormalTok{)}

\NormalTok{models }\OtherTok{\textless{}{-}} \FunctionTok{map}\NormalTok{(ks, \textbackslash{}(k) \{}
 \FunctionTok{LDA}\NormalTok{(dtm, }\AttributeTok{k =}\NormalTok{ k, }\AttributeTok{method =} \StringTok{"Gibbs"}\NormalTok{,}
     \AttributeTok{control =} \FunctionTok{list}\NormalTok{(}\AttributeTok{seed =} \DecValTok{20260509}\NormalTok{, }\AttributeTok{iter =} \DecValTok{1000}\NormalTok{))}
\NormalTok{\})}

\CommentTok{\# Compute coherence (using top{-}10 terms per topic)}
\NormalTok{coherence\_scores }\OtherTok{\textless{}{-}} \FunctionTok{map\_dbl}\NormalTok{(models, \textbackslash{}(mod) \{}
\NormalTok{ terms }\OtherTok{\textless{}{-}} \FunctionTok{terms}\NormalTok{(mod, }\DecValTok{10}\NormalTok{)}
 \CommentTok{\# Simple PMI{-}based coherence approximation}
 \FunctionTok{mean}\NormalTok{(}\FunctionTok{map\_dbl}\NormalTok{(}\FunctionTok{seq\_len}\NormalTok{(}\FunctionTok{ncol}\NormalTok{(terms)), \textbackslash{}(t) \{}
\NormalTok{   top\_terms }\OtherTok{\textless{}{-}}\NormalTok{ terms[, t]}
\NormalTok{   pairs }\OtherTok{\textless{}{-}} \FunctionTok{combn}\NormalTok{(top\_terms, }\DecValTok{2}\NormalTok{, }\AttributeTok{simplify =} \ConstantTok{FALSE}\NormalTok{)}
   \FunctionTok{mean}\NormalTok{(}\FunctionTok{map\_dbl}\NormalTok{(pairs, \textbackslash{}(p) \{}
     \FunctionTok{log}\NormalTok{((slam}\SpecialCharTok{::}\FunctionTok{col\_sums}\NormalTok{(dtm[, p[}\DecValTok{1}\NormalTok{]] }\SpecialCharTok{\textgreater{}} \DecValTok{0} \SpecialCharTok{\&}\NormalTok{ dtm[, p[}\DecValTok{2}\NormalTok{]] }\SpecialCharTok{\textgreater{}} \DecValTok{0}\NormalTok{) }\SpecialCharTok{+} \DecValTok{1}\NormalTok{) }\SpecialCharTok{/}
\NormalTok{         slam}\SpecialCharTok{::}\FunctionTok{col\_sums}\NormalTok{(dtm[, p[}\DecValTok{1}\NormalTok{]] }\SpecialCharTok{\textgreater{}} \DecValTok{0}\NormalTok{))}
\NormalTok{   \}))}
\NormalTok{ \}))}
\NormalTok{\})}

\FunctionTok{tibble}\NormalTok{(}\AttributeTok{k =}\NormalTok{ ks, }\AttributeTok{coherence =}\NormalTok{ coherence\_scores) }\SpecialCharTok{|\textgreater{}}
 \FunctionTok{ggplot}\NormalTok{(}\FunctionTok{aes}\NormalTok{(k, coherence)) }\SpecialCharTok{+}
 \FunctionTok{geom\_line}\NormalTok{() }\SpecialCharTok{+}
 \FunctionTok{geom\_point}\NormalTok{() }\SpecialCharTok{+}
 \FunctionTok{labs}\NormalTok{(}\AttributeTok{title =} \StringTok{"Topic coherence by K"}\NormalTok{, }\AttributeTok{x =} \StringTok{"K"}\NormalTok{, }\AttributeTok{y =} \StringTok{"Coherence"}\NormalTok{)}
\end{Highlighting}
\end{Shaded}

\emph{Key insight:} The best K balances coherence (higher is better) and
interpretability. There is rarely a single ``correct'' K.

\textbf{2. STM with journal covariate.}
Fit a Structural Topic Model with journal as a prevalence covariate.

\begin{Shaded}
\begin{Highlighting}[]
\FunctionTok{library}\NormalTok{(stm)}
\FunctionTok{set.seed}\NormalTok{(}\DecValTok{20260509}\NormalTok{)}

\CommentTok{\# Prepare data from two journals}
\NormalTok{two\_journals }\OtherTok{\textless{}{-}}\NormalTok{ works }\SpecialCharTok{|\textgreater{}}
 \FunctionTok{filter}\NormalTok{(so }\SpecialCharTok{\%in\%} \FunctionTok{c}\NormalTok{(}\StringTok{"Scientometrics"}\NormalTok{, }\StringTok{"PLOS ONE"}\NormalTok{), }\SpecialCharTok{!}\FunctionTok{is.na}\NormalTok{(ab))}

\NormalTok{processed }\OtherTok{\textless{}{-}} \FunctionTok{textProcessor}\NormalTok{(two\_journals}\SpecialCharTok{$}\NormalTok{ab,}
                           \AttributeTok{metadata =}\NormalTok{ two\_journals)}
\NormalTok{prepped }\OtherTok{\textless{}{-}} \FunctionTok{prepDocuments}\NormalTok{(processed}\SpecialCharTok{$}\NormalTok{documents,}
\NormalTok{                         processed}\SpecialCharTok{$}\NormalTok{vocab,}
\NormalTok{                         processed}\SpecialCharTok{$}\NormalTok{meta)}

\NormalTok{stm\_fit }\OtherTok{\textless{}{-}} \FunctionTok{stm}\NormalTok{(prepped}\SpecialCharTok{$}\NormalTok{documents, prepped}\SpecialCharTok{$}\NormalTok{vocab,}
               \AttributeTok{K =} \DecValTok{10}\NormalTok{,}
               \AttributeTok{prevalence =} \SpecialCharTok{\textasciitilde{}}\NormalTok{ so,}
               \AttributeTok{data =}\NormalTok{ prepped}\SpecialCharTok{$}\NormalTok{meta,}
               \AttributeTok{seed =} \DecValTok{20260509}\NormalTok{)}

\NormalTok{effect }\OtherTok{\textless{}{-}} \FunctionTok{estimateEffect}\NormalTok{(}\DecValTok{1}\SpecialCharTok{:}\DecValTok{10} \SpecialCharTok{\textasciitilde{}}\NormalTok{ so, stm\_fit, }\AttributeTok{metadata =}\NormalTok{ prepped}\SpecialCharTok{$}\NormalTok{meta)}
\FunctionTok{plot}\NormalTok{(effect, }\AttributeTok{covariate =} \StringTok{"so"}\NormalTok{, }\AttributeTok{topics =} \DecValTok{1}\SpecialCharTok{:}\DecValTok{10}\NormalTok{)}
\end{Highlighting}
\end{Shaded}

\emph{Key insight:} Topics with large journal effects represent disciplinary
specialisation; shared topics indicate cross-cutting themes.

\textbf{3. Topic labelling.}
Read top documents per topic to assign labels.

\begin{Shaded}
\begin{Highlighting}[]
\FunctionTok{library}\NormalTok{(topicmodels)}
\FunctionTok{library}\NormalTok{(tidyverse)}

\NormalTok{best\_model }\OtherTok{\textless{}{-}}\NormalTok{ models[[}\FunctionTok{which.max}\NormalTok{(coherence\_scores)]]}
\NormalTok{gamma }\OtherTok{\textless{}{-}} \FunctionTok{posterior}\NormalTok{(best\_model)}\SpecialCharTok{$}\NormalTok{topics}

\CommentTok{\# For each topic, find the top{-}5 documents}
\NormalTok{topic\_labels }\OtherTok{\textless{}{-}} \FunctionTok{map}\NormalTok{(}\FunctionTok{seq\_len}\NormalTok{(}\FunctionTok{ncol}\NormalTok{(gamma)), \textbackslash{}(t) \{}
\NormalTok{ top\_docs }\OtherTok{\textless{}{-}} \FunctionTok{order}\NormalTok{(gamma[, t], }\AttributeTok{decreasing =} \ConstantTok{TRUE}\NormalTok{)[}\DecValTok{1}\SpecialCharTok{:}\DecValTok{5}\NormalTok{]}
\NormalTok{ titles }\OtherTok{\textless{}{-}}\NormalTok{ works}\SpecialCharTok{$}\NormalTok{display\_name[top\_docs]}
 \FunctionTok{tibble}\NormalTok{(}\AttributeTok{topic =}\NormalTok{ t,}
        \AttributeTok{top\_terms =} \FunctionTok{paste}\NormalTok{(}\FunctionTok{terms}\NormalTok{(best\_model, }\DecValTok{5}\NormalTok{)[, t], }\AttributeTok{collapse =} \StringTok{", "}\NormalTok{),}
        \AttributeTok{sample\_titles =} \FunctionTok{paste}\NormalTok{(titles, }\AttributeTok{collapse =} \StringTok{" | "}\NormalTok{))}
\NormalTok{\}) }\SpecialCharTok{|\textgreater{}} \FunctionTok{list\_rbind}\NormalTok{()}

\NormalTok{topic\_labels}
\CommentTok{\# Now assign human{-}readable labels based on inspection}
\end{Highlighting}
\end{Shaded}

\emph{Key insight:} Automated topic labels from top terms are a starting point;
reading representative documents is essential for meaningful labels.

\textbf{4. LDA vs.~STM comparison.}
Compare topics from both models.

\begin{Shaded}
\begin{Highlighting}[]
\CommentTok{\# Extract top terms for each model}
\NormalTok{lda\_terms }\OtherTok{\textless{}{-}} \FunctionTok{terms}\NormalTok{(best\_model, }\DecValTok{10}\NormalTok{)}
\NormalTok{stm\_terms }\OtherTok{\textless{}{-}} \FunctionTok{labelTopics}\NormalTok{(stm\_fit, }\AttributeTok{n =} \DecValTok{10}\NormalTok{)}\SpecialCharTok{$}\NormalTok{prob}

\CommentTok{\# Compute Jaccard similarity between matched topics}
\NormalTok{jaccard }\OtherTok{\textless{}{-}} \ControlFlowTok{function}\NormalTok{(a, b) \{}
 \FunctionTok{length}\NormalTok{(}\FunctionTok{intersect}\NormalTok{(a, b)) }\SpecialCharTok{/} \FunctionTok{length}\NormalTok{(}\FunctionTok{union}\NormalTok{(a, b))}
\NormalTok{\}}

\CommentTok{\# Best{-}match each LDA topic to an STM topic}
\NormalTok{similarities }\OtherTok{\textless{}{-}} \FunctionTok{map\_dbl}\NormalTok{(}\FunctionTok{seq\_len}\NormalTok{(}\FunctionTok{ncol}\NormalTok{(lda\_terms)), \textbackslash{}(i) \{}
 \FunctionTok{max}\NormalTok{(}\FunctionTok{map\_dbl}\NormalTok{(}\FunctionTok{seq\_len}\NormalTok{(}\FunctionTok{nrow}\NormalTok{(stm\_terms)), \textbackslash{}(j) \{}
   \FunctionTok{jaccard}\NormalTok{(lda\_terms[, i], stm\_terms[j, ])}
\NormalTok{ \}))}
\NormalTok{\})}

\FunctionTok{mean}\NormalTok{(similarities)}
\end{Highlighting}
\end{Shaded}

\emph{Key insight:} STM topics may differ because the year covariate allows topic
prevalence to shift over time, producing more temporally coherent topics.

\section*{Chapter 23 --- Keyword Extraction}\label{chapter-23-keyword-extraction}
\addcontentsline{toc}{section}{Chapter 23 --- Keyword Extraction}

\textbf{1. Precision and recall.}
Evaluate RAKE keywords against author keywords.

\begin{Shaded}
\begin{Highlighting}[]
\FunctionTok{library}\NormalTok{(udpipe)}
\FunctionTok{library}\NormalTok{(tidyverse)}

\NormalTok{sample\_papers }\OtherTok{\textless{}{-}}\NormalTok{ works }\SpecialCharTok{|\textgreater{}}
 \FunctionTok{filter}\NormalTok{(}\SpecialCharTok{!}\FunctionTok{is.na}\NormalTok{(ab), }\FunctionTok{map\_lgl}\NormalTok{(keywords, \textbackslash{}(k) }\FunctionTok{length}\NormalTok{(k) }\SpecialCharTok{\textgreater{}} \DecValTok{0}\NormalTok{)) }\SpecialCharTok{|\textgreater{}}
 \FunctionTok{slice\_sample}\NormalTok{(}\AttributeTok{n =} \DecValTok{20}\NormalTok{)}

\NormalTok{eval\_results }\OtherTok{\textless{}{-}} \FunctionTok{map}\NormalTok{(}\FunctionTok{seq\_len}\NormalTok{(}\FunctionTok{nrow}\NormalTok{(sample\_papers)), \textbackslash{}(i) \{}
\NormalTok{ paper }\OtherTok{\textless{}{-}}\NormalTok{ sample\_papers[i, ]}
\NormalTok{ rake\_kw }\OtherTok{\textless{}{-}} \FunctionTok{keywords\_rake}\NormalTok{(}
   \FunctionTok{data.frame}\NormalTok{(}\AttributeTok{doc\_id =}\NormalTok{ paper}\SpecialCharTok{$}\NormalTok{id,}
              \AttributeTok{sentence =}\NormalTok{ paper}\SpecialCharTok{$}\NormalTok{ab,}
              \AttributeTok{stringsAsFactors =} \ConstantTok{FALSE}\NormalTok{),}
   \AttributeTok{term =} \StringTok{"sentence"}\NormalTok{, }\AttributeTok{group =} \StringTok{"doc\_id"}
\NormalTok{ ) }\SpecialCharTok{|\textgreater{}} \FunctionTok{head}\NormalTok{(}\DecValTok{10}\NormalTok{) }\SpecialCharTok{|\textgreater{}} \FunctionTok{pull}\NormalTok{(keyword) }\SpecialCharTok{|\textgreater{}} \FunctionTok{tolower}\NormalTok{()}

\NormalTok{ author\_kw }\OtherTok{\textless{}{-}} \FunctionTok{tolower}\NormalTok{(}\FunctionTok{unlist}\NormalTok{(paper}\SpecialCharTok{$}\NormalTok{keywords))}

 \FunctionTok{tibble}\NormalTok{(}
   \AttributeTok{precision =} \FunctionTok{mean}\NormalTok{(rake\_kw }\SpecialCharTok{\%in\%}\NormalTok{ author\_kw),}
   \AttributeTok{recall =} \FunctionTok{mean}\NormalTok{(author\_kw }\SpecialCharTok{\%in\%}\NormalTok{ rake\_kw)}
\NormalTok{ )}
\NormalTok{\}) }\SpecialCharTok{|\textgreater{}} \FunctionTok{list\_rbind}\NormalTok{()}

\NormalTok{eval\_results }\SpecialCharTok{|\textgreater{}}
 \FunctionTok{summarise}\NormalTok{(}\FunctionTok{across}\NormalTok{(}\FunctionTok{everything}\NormalTok{(), }\FunctionTok{list}\NormalTok{(}\AttributeTok{mean =}\NormalTok{ mean, }\AttributeTok{sd =}\NormalTok{ sd)))}
\end{Highlighting}
\end{Shaded}

\emph{Key insight:} RAKE precision is typically low (0.1-0.3) because algorithms
and authors conceptualise keywords differently. Recall is also limited because
author keywords often include broad terms not in the abstract.

\textbf{2. Bigram keywords.}
Score bigrams by TF-IDF.

\begin{Shaded}
\begin{Highlighting}[]
\FunctionTok{library}\NormalTok{(tidytext)}
\FunctionTok{library}\NormalTok{(tidyverse)}

\NormalTok{bigram\_tfidf }\OtherTok{\textless{}{-}}\NormalTok{ works }\SpecialCharTok{|\textgreater{}}
 \FunctionTok{filter}\NormalTok{(}\SpecialCharTok{!}\FunctionTok{is.na}\NormalTok{(ab)) }\SpecialCharTok{|\textgreater{}}
 \FunctionTok{unnest\_tokens}\NormalTok{(bigram, ab, }\AttributeTok{token =} \StringTok{"ngrams"}\NormalTok{, }\AttributeTok{n =} \DecValTok{2}\NormalTok{) }\SpecialCharTok{|\textgreater{}}
 \FunctionTok{separate}\NormalTok{(bigram, }\FunctionTok{c}\NormalTok{(}\StringTok{"word1"}\NormalTok{, }\StringTok{"word2"}\NormalTok{), }\AttributeTok{sep =} \StringTok{" "}\NormalTok{) }\SpecialCharTok{|\textgreater{}}
 \FunctionTok{filter}\NormalTok{(}\SpecialCharTok{!}\NormalTok{word1 }\SpecialCharTok{\%in\%}\NormalTok{ stop\_words}\SpecialCharTok{$}\NormalTok{word,}
        \SpecialCharTok{!}\NormalTok{word2 }\SpecialCharTok{\%in\%}\NormalTok{ stop\_words}\SpecialCharTok{$}\NormalTok{word) }\SpecialCharTok{|\textgreater{}}
 \FunctionTok{unite}\NormalTok{(bigram, word1, word2, }\AttributeTok{sep =} \StringTok{" "}\NormalTok{) }\SpecialCharTok{|\textgreater{}}
 \FunctionTok{count}\NormalTok{(id, bigram) }\SpecialCharTok{|\textgreater{}}
 \FunctionTok{bind\_tf\_idf}\NormalTok{(bigram, id, n) }\SpecialCharTok{|\textgreater{}}
 \FunctionTok{group\_by}\NormalTok{(bigram) }\SpecialCharTok{|\textgreater{}}
 \FunctionTok{summarise}\NormalTok{(}\AttributeTok{mean\_tfidf =} \FunctionTok{mean}\NormalTok{(tf\_idf)) }\SpecialCharTok{|\textgreater{}}
 \FunctionTok{arrange}\NormalTok{(}\FunctionTok{desc}\NormalTok{(mean\_tfidf))}

\NormalTok{bigram\_tfidf }\SpecialCharTok{|\textgreater{}} \FunctionTok{head}\NormalTok{(}\DecValTok{20}\NormalTok{)}
\end{Highlighting}
\end{Shaded}

\emph{Key insight:} Multi-word terms like ``citation analysis'' and ``research
evaluation'' are invisible to unigram methods but highly informative.

\textbf{3. Keyword enrichment.}
Use extracted keywords to build a co-word network for keyword-poor papers.

\begin{Shaded}
\begin{Highlighting}[]
\FunctionTok{library}\NormalTok{(igraph)}
\FunctionTok{library}\NormalTok{(udpipe)}

\CommentTok{\# Papers without author keywords}
\NormalTok{no\_kw\_papers }\OtherTok{\textless{}{-}}\NormalTok{ works }\SpecialCharTok{|\textgreater{}}
 \FunctionTok{filter}\NormalTok{(}\FunctionTok{map\_lgl}\NormalTok{(keywords, \textbackslash{}(k) }\FunctionTok{length}\NormalTok{(k) }\SpecialCharTok{==} \DecValTok{0}\NormalTok{), }\SpecialCharTok{!}\FunctionTok{is.na}\NormalTok{(ab))}

\CommentTok{\# Extract keywords via RAKE}
\NormalTok{extracted }\OtherTok{\textless{}{-}} \FunctionTok{map}\NormalTok{(}\FunctionTok{seq\_len}\NormalTok{(}\FunctionTok{nrow}\NormalTok{(no\_kw\_papers)), \textbackslash{}(i) \{}
\NormalTok{ kw }\OtherTok{\textless{}{-}} \FunctionTok{keywords\_rake}\NormalTok{(}
   \FunctionTok{data.frame}\NormalTok{(}\AttributeTok{doc\_id =}\NormalTok{ no\_kw\_papers}\SpecialCharTok{$}\NormalTok{id[i],}
              \AttributeTok{sentence =}\NormalTok{ no\_kw\_papers}\SpecialCharTok{$}\NormalTok{ab[i]),}
   \AttributeTok{term =} \StringTok{"sentence"}\NormalTok{, }\AttributeTok{group =} \StringTok{"doc\_id"}
\NormalTok{ ) }\SpecialCharTok{|\textgreater{}} \FunctionTok{head}\NormalTok{(}\DecValTok{5}\NormalTok{) }\SpecialCharTok{|\textgreater{}} \FunctionTok{pull}\NormalTok{(keyword)}
 \FunctionTok{tibble}\NormalTok{(}\AttributeTok{id =}\NormalTok{ no\_kw\_papers}\SpecialCharTok{$}\NormalTok{id[i], }\AttributeTok{keyword =}\NormalTok{ kw)}
\NormalTok{\}) }\SpecialCharTok{|\textgreater{}} \FunctionTok{list\_rbind}\NormalTok{()}

\CommentTok{\# Build co{-}word network}
\NormalTok{pairs }\OtherTok{\textless{}{-}}\NormalTok{ extracted }\SpecialCharTok{|\textgreater{}}
 \FunctionTok{group\_by}\NormalTok{(id) }\SpecialCharTok{|\textgreater{}}
 \FunctionTok{filter}\NormalTok{(}\FunctionTok{n}\NormalTok{() }\SpecialCharTok{\textgreater{}=} \DecValTok{2}\NormalTok{) }\SpecialCharTok{|\textgreater{}}
 \FunctionTok{reframe}\NormalTok{(}\AttributeTok{pair =} \FunctionTok{combn}\NormalTok{(keyword, }\DecValTok{2}\NormalTok{, }\AttributeTok{simplify =} \ConstantTok{FALSE}\NormalTok{)) }\SpecialCharTok{|\textgreater{}}
 \FunctionTok{unnest\_wider}\NormalTok{(pair, }\AttributeTok{names\_sep =} \StringTok{"\_"}\NormalTok{) }\SpecialCharTok{|\textgreater{}}
 \FunctionTok{count}\NormalTok{(pair\_1, pair\_2, }\AttributeTok{sort =} \ConstantTok{TRUE}\NormalTok{)}

\NormalTok{g\_enriched }\OtherTok{\textless{}{-}} \FunctionTok{graph\_from\_data\_frame}\NormalTok{(pairs, }\AttributeTok{directed =} \ConstantTok{FALSE}\NormalTok{)}
\end{Highlighting}
\end{Shaded}

\emph{Key insight:} Extracted keywords produce a topical structure comparable to
author keywords, enabling co-word analysis even for papers without manual
keyword assignments.

\section*{Chapter 24 --- Embeddings for Scientometrics}\label{chapter-24-embeddings-for-scientometrics}
\addcontentsline{toc}{section}{Chapter 24 --- Embeddings for Scientometrics}

\textbf{1. Weighted averaging.}
Use TF-IDF weights for document embedding computation.

\begin{Shaded}
\begin{Highlighting}[]
\FunctionTok{library}\NormalTok{(word2vec)}
\FunctionTok{library}\NormalTok{(tidytext)}
\FunctionTok{library}\NormalTok{(tidyverse)}
\FunctionTok{library}\NormalTok{(umap)}

\CommentTok{\# Compute TF{-}IDF weights}
\NormalTok{tfidf }\OtherTok{\textless{}{-}}\NormalTok{ works }\SpecialCharTok{|\textgreater{}}
 \FunctionTok{filter}\NormalTok{(}\SpecialCharTok{!}\FunctionTok{is.na}\NormalTok{(ab)) }\SpecialCharTok{|\textgreater{}}
 \FunctionTok{unnest\_tokens}\NormalTok{(word, ab) }\SpecialCharTok{|\textgreater{}}
 \FunctionTok{anti\_join}\NormalTok{(stop\_words, }\AttributeTok{by =} \StringTok{"word"}\NormalTok{) }\SpecialCharTok{|\textgreater{}}
 \FunctionTok{count}\NormalTok{(id, word) }\SpecialCharTok{|\textgreater{}}
 \FunctionTok{bind\_tf\_idf}\NormalTok{(word, id, n)}

\CommentTok{\# Get word vectors}
\NormalTok{wv }\OtherTok{\textless{}{-}} \FunctionTok{predict}\NormalTok{(w2v\_model, }\AttributeTok{newdata =} \FunctionTok{unique}\NormalTok{(tfidf}\SpecialCharTok{$}\NormalTok{word), }\AttributeTok{type =} \StringTok{"embedding"}\NormalTok{)}

\CommentTok{\# Weighted average per document}
\NormalTok{doc\_embeddings }\OtherTok{\textless{}{-}}\NormalTok{ tfidf }\SpecialCharTok{|\textgreater{}}
 \FunctionTok{inner\_join}\NormalTok{(}
   \FunctionTok{as\_tibble}\NormalTok{(wv, }\AttributeTok{rownames =} \StringTok{"word"}\NormalTok{),}
   \AttributeTok{by =} \StringTok{"word"}
\NormalTok{ ) }\SpecialCharTok{|\textgreater{}}
 \FunctionTok{group\_by}\NormalTok{(id) }\SpecialCharTok{|\textgreater{}}
 \FunctionTok{summarise}\NormalTok{(}\FunctionTok{across}\NormalTok{(}\FunctionTok{starts\_with}\NormalTok{(}\StringTok{"V"}\NormalTok{), \textbackslash{}(x) }\FunctionTok{weighted.mean}\NormalTok{(x, tf\_idf)))}

\CommentTok{\# UMAP}
\NormalTok{umap\_result }\OtherTok{\textless{}{-}} \FunctionTok{umap}\NormalTok{(}\FunctionTok{as.matrix}\NormalTok{(doc\_embeddings }\SpecialCharTok{|\textgreater{}} \FunctionTok{select}\NormalTok{(}\SpecialCharTok{{-}}\NormalTok{id)))}
\end{Highlighting}
\end{Shaded}

\emph{Key insight:} TF-IDF weighting down-weights common words, often producing
tighter, more meaningful clusters in the UMAP projection.

\textbf{2. K-means clustering.}
Apply k-means to embeddings and compare with LDA topics.

\begin{Shaded}
\begin{Highlighting}[]
\FunctionTok{set.seed}\NormalTok{(}\DecValTok{20260509}\NormalTok{)}

\NormalTok{embedding\_matrix }\OtherTok{\textless{}{-}} \FunctionTok{as.matrix}\NormalTok{(doc\_embeddings }\SpecialCharTok{|\textgreater{}} \FunctionTok{select}\NormalTok{(}\SpecialCharTok{{-}}\NormalTok{id))}

\NormalTok{km\_results }\OtherTok{\textless{}{-}} \FunctionTok{map}\NormalTok{(}\FunctionTok{c}\NormalTok{(}\DecValTok{5}\NormalTok{, }\DecValTok{8}\NormalTok{, }\DecValTok{12}\NormalTok{), \textbackslash{}(k) \{}
\NormalTok{ km }\OtherTok{\textless{}{-}} \FunctionTok{kmeans}\NormalTok{(embedding\_matrix, }\AttributeTok{centers =}\NormalTok{ k, }\AttributeTok{nstart =} \DecValTok{25}\NormalTok{)}
 \FunctionTok{tibble}\NormalTok{(}\AttributeTok{k =}\NormalTok{ k,}
        \AttributeTok{cluster =}\NormalTok{ km}\SpecialCharTok{$}\NormalTok{cluster,}
        \AttributeTok{id =}\NormalTok{ doc\_embeddings}\SpecialCharTok{$}\NormalTok{id)}
\NormalTok{\}) }\SpecialCharTok{|\textgreater{}} \FunctionTok{list\_rbind}\NormalTok{()}

\CommentTok{\# Compare k=8 clusters with LDA topics (K=8)}
\CommentTok{\# Assuming lda\_topics is a vector of topic assignments}
\NormalTok{km8 }\OtherTok{\textless{}{-}}\NormalTok{ km\_results }\SpecialCharTok{|\textgreater{}} \FunctionTok{filter}\NormalTok{(k }\SpecialCharTok{==} \DecValTok{8}\NormalTok{) }\SpecialCharTok{|\textgreater{}} \FunctionTok{pull}\NormalTok{(cluster)}
\NormalTok{compare\_nmi }\OtherTok{\textless{}{-}}\NormalTok{ igraph}\SpecialCharTok{::}\FunctionTok{compare}\NormalTok{(km8, lda\_topics, }\AttributeTok{method =} \StringTok{"nmi"}\NormalTok{)}
\FunctionTok{cat}\NormalTok{(}\FunctionTok{sprintf}\NormalTok{(}\StringTok{"NMI between k{-}means and LDA: \%.3f}\SpecialCharTok{\textbackslash{}n}\StringTok{"}\NormalTok{, compare\_nmi))}
\end{Highlighting}
\end{Shaded}

\emph{Key insight:} Moderate NMI (0.3-0.6) is typical; the two methods capture
overlapping but not identical structure because they model different aspects of
the data.

\textbf{3. Pre-trained embeddings.}
Use SPECTER embeddings via reticulate.

\begin{Shaded}
\begin{Highlighting}[]
\FunctionTok{library}\NormalTok{(reticulate)}

\CommentTok{\# Requires Python with transformers installed}
\NormalTok{transformers }\OtherTok{\textless{}{-}} \FunctionTok{import}\NormalTok{(}\StringTok{"transformers"}\NormalTok{)}
\NormalTok{torch }\OtherTok{\textless{}{-}} \FunctionTok{import}\NormalTok{(}\StringTok{"torch"}\NormalTok{)}

\NormalTok{tokenizer }\OtherTok{\textless{}{-}}\NormalTok{ transformers}\SpecialCharTok{$}\NormalTok{AutoTokenizer}\SpecialCharTok{$}\FunctionTok{from\_pretrained}\NormalTok{(}
 \StringTok{"allenai/specter2"}
\NormalTok{)}
\NormalTok{model }\OtherTok{\textless{}{-}}\NormalTok{ transformers}\SpecialCharTok{$}\NormalTok{AutoModel}\SpecialCharTok{$}\FunctionTok{from\_pretrained}\NormalTok{(}
 \StringTok{"allenai/specter2"}
\NormalTok{)}

\NormalTok{abstracts }\OtherTok{\textless{}{-}}\NormalTok{ works}\SpecialCharTok{$}\NormalTok{ab[}\SpecialCharTok{!}\FunctionTok{is.na}\NormalTok{(works}\SpecialCharTok{$}\NormalTok{ab)][}\DecValTok{1}\SpecialCharTok{:}\DecValTok{100}\NormalTok{]}
\NormalTok{inputs }\OtherTok{\textless{}{-}} \FunctionTok{tokenizer}\NormalTok{(abstracts, }\AttributeTok{padding =} \ConstantTok{TRUE}\NormalTok{, }\AttributeTok{truncation =} \ConstantTok{TRUE}\NormalTok{,}
                    \AttributeTok{return\_tensors =} \StringTok{"pt"}\NormalTok{, }\AttributeTok{max\_length =} \DecValTok{512}\DataTypeTok{L}\NormalTok{)}
\FunctionTok{with}\NormalTok{(torch}\SpecialCharTok{$}\FunctionTok{no\_grad}\NormalTok{(), \{}
\NormalTok{ outputs }\OtherTok{\textless{}{-}} \FunctionTok{model}\NormalTok{(inputs)}
\NormalTok{\})}
\NormalTok{specter\_emb }\OtherTok{\textless{}{-}}\NormalTok{ outputs}\SpecialCharTok{$}\NormalTok{last\_hidden\_state}\SpecialCharTok{$}\FunctionTok{mean}\NormalTok{(}\AttributeTok{dim =} \DecValTok{1}\DataTypeTok{L}\NormalTok{)}\SpecialCharTok{$}\FunctionTok{numpy}\NormalTok{()}
\end{Highlighting}
\end{Shaded}

\emph{Key insight:} Pre-trained scientific embeddings like SPECTER capture domain
semantics better than generic word2vec, especially for similarity search.

\textbf{4. Analogy test.}
Test analogical reasoning in the word2vec model.

\begin{Shaded}
\begin{Highlighting}[]
\FunctionTok{library}\NormalTok{(word2vec)}

\CommentTok{\# Test: "journal" {-} "article" + "book" ≈ ?}
\NormalTok{analogy }\OtherTok{\textless{}{-}} \FunctionTok{predict}\NormalTok{(w2v\_model,}
                   \AttributeTok{newdata =} \FunctionTok{c}\NormalTok{(}\StringTok{"journal"}\NormalTok{, }\StringTok{"article"}\NormalTok{, }\StringTok{"book"}\NormalTok{),}
                   \AttributeTok{type =} \StringTok{"nearest"}\NormalTok{, }\AttributeTok{top\_n =} \DecValTok{5}\NormalTok{)}

\CommentTok{\# Alternative approach using vector arithmetic}
\NormalTok{wv }\OtherTok{\textless{}{-}} \FunctionTok{as.matrix}\NormalTok{(w2v\_model)}
\ControlFlowTok{if}\NormalTok{ (}\FunctionTok{all}\NormalTok{(}\FunctionTok{c}\NormalTok{(}\StringTok{"journal"}\NormalTok{, }\StringTok{"article"}\NormalTok{, }\StringTok{"book"}\NormalTok{) }\SpecialCharTok{\%in\%} \FunctionTok{rownames}\NormalTok{(wv))) \{}
\NormalTok{ target }\OtherTok{\textless{}{-}}\NormalTok{ wv[}\StringTok{"journal"}\NormalTok{, ] }\SpecialCharTok{{-}}\NormalTok{ wv[}\StringTok{"article"}\NormalTok{, ] }\SpecialCharTok{+}\NormalTok{ wv[}\StringTok{"book"}\NormalTok{, ]}
\NormalTok{ sims }\OtherTok{\textless{}{-}} \FunctionTok{word2vec\_similarity}\NormalTok{(}
   \FunctionTok{matrix}\NormalTok{(target, }\AttributeTok{nrow =} \DecValTok{1}\NormalTok{), wv}
\NormalTok{ )}
 \FunctionTok{sort}\NormalTok{(sims[}\DecValTok{1}\NormalTok{, ], }\AttributeTok{decreasing =} \ConstantTok{TRUE}\NormalTok{)[}\DecValTok{1}\SpecialCharTok{:}\DecValTok{5}\NormalTok{]}
\NormalTok{\}}
\end{Highlighting}
\end{Shaded}

\emph{Key insight:} Domain-specific models may capture field-specific analogies (e.g.,
``h-index'' - ``author'' + ``journal'' near ``impact factor'') but fail on generic
analogies, reflecting the training corpus.

\section*{Chapter 25 --- Detecting Emerging Topics}\label{chapter-25-detecting-emerging-topics}
\addcontentsline{toc}{section}{Chapter 25 --- Detecting Emerging Topics}

\textbf{1. Relative burst detection.}
Use proportional frequencies instead of absolute counts.

\begin{Shaded}
\begin{Highlighting}[]
\FunctionTok{library}\NormalTok{(tidyverse)}

\CommentTok{\# Compute keyword proportions by year}
\NormalTok{kw\_yearly }\OtherTok{\textless{}{-}}\NormalTok{ works }\SpecialCharTok{|\textgreater{}}
 \FunctionTok{unnest}\NormalTok{(keywords) }\SpecialCharTok{|\textgreater{}}
 \FunctionTok{count}\NormalTok{(publication\_year, keyword) }\SpecialCharTok{|\textgreater{}}
 \FunctionTok{group\_by}\NormalTok{(publication\_year) }\SpecialCharTok{|\textgreater{}}
 \FunctionTok{mutate}\NormalTok{(}\AttributeTok{proportion =}\NormalTok{ n }\SpecialCharTok{/} \FunctionTok{sum}\NormalTok{(n)) }\SpecialCharTok{|\textgreater{}}
 \FunctionTok{ungroup}\NormalTok{()}

\CommentTok{\# Apply Kleinberg burst detection on proportions}
\NormalTok{bursts\_relative }\OtherTok{\textless{}{-}}\NormalTok{ kw\_yearly }\SpecialCharTok{|\textgreater{}}
 \FunctionTok{group\_by}\NormalTok{(keyword) }\SpecialCharTok{|\textgreater{}}
 \FunctionTok{filter}\NormalTok{(}\FunctionTok{n}\NormalTok{() }\SpecialCharTok{\textgreater{}=} \DecValTok{3}\NormalTok{) }\SpecialCharTok{|\textgreater{}}  \CommentTok{\# need multiple years}
 \FunctionTok{summarise}\NormalTok{(}
   \AttributeTok{burst =} \FunctionTok{list}\NormalTok{(}\FunctionTok{kleinberg\_bursts}\NormalTok{(proportion, publication\_year)),}
   \AttributeTok{.groups =} \StringTok{"drop"}
\NormalTok{ ) }\SpecialCharTok{|\textgreater{}}
 \FunctionTok{unnest}\NormalTok{(burst) }\SpecialCharTok{|\textgreater{}}
 \FunctionTok{arrange}\NormalTok{(}\FunctionTok{desc}\NormalTok{(strength))}

\NormalTok{bursts\_relative }\SpecialCharTok{|\textgreater{}} \FunctionTok{head}\NormalTok{(}\DecValTok{20}\NormalTok{)}
\end{Highlighting}
\end{Shaded}

\emph{Key insight:} Relative frequencies reduce bias toward high-volume years;
some terms that burst in absolute counts may be merely riding overall growth.

\textbf{2. Emergence by journal.}
Compare emerging keywords across two journals.

\begin{Shaded}
\begin{Highlighting}[]
\NormalTok{journals }\OtherTok{\textless{}{-}} \FunctionTok{c}\NormalTok{(}\StringTok{"Scientometrics"}\NormalTok{, }\StringTok{"PLOS ONE"}\NormalTok{)}

\NormalTok{journal\_bursts }\OtherTok{\textless{}{-}} \FunctionTok{map}\NormalTok{(journals, \textbackslash{}(j) \{}
\NormalTok{ kw\_j }\OtherTok{\textless{}{-}}\NormalTok{ works }\SpecialCharTok{|\textgreater{}}
   \FunctionTok{filter}\NormalTok{(so }\SpecialCharTok{==}\NormalTok{ j) }\SpecialCharTok{|\textgreater{}}
   \FunctionTok{unnest}\NormalTok{(keywords) }\SpecialCharTok{|\textgreater{}}
   \FunctionTok{count}\NormalTok{(publication\_year, keyword) }\SpecialCharTok{|\textgreater{}}
   \FunctionTok{group\_by}\NormalTok{(keyword) }\SpecialCharTok{|\textgreater{}}
   \FunctionTok{filter}\NormalTok{(}\FunctionTok{n}\NormalTok{() }\SpecialCharTok{\textgreater{}=} \DecValTok{3}\NormalTok{) }\SpecialCharTok{|\textgreater{}}
   \FunctionTok{summarise}\NormalTok{(}
     \AttributeTok{burst =} \FunctionTok{list}\NormalTok{(}\FunctionTok{kleinberg\_bursts}\NormalTok{(n, publication\_year))}
\NormalTok{   ) }\SpecialCharTok{|\textgreater{}}
   \FunctionTok{unnest}\NormalTok{(burst) }\SpecialCharTok{|\textgreater{}}
   \FunctionTok{mutate}\NormalTok{(}\AttributeTok{journal =}\NormalTok{ j)}
\NormalTok{\}) }\SpecialCharTok{|\textgreater{}} \FunctionTok{list\_rbind}\NormalTok{()}

\CommentTok{\# Compare top bursting terms}
\NormalTok{journal\_bursts }\SpecialCharTok{|\textgreater{}}
 \FunctionTok{group\_by}\NormalTok{(journal) }\SpecialCharTok{|\textgreater{}}
 \FunctionTok{arrange}\NormalTok{(}\FunctionTok{desc}\NormalTok{(strength)) }\SpecialCharTok{|\textgreater{}}
 \FunctionTok{slice\_head}\NormalTok{(}\AttributeTok{n =} \DecValTok{10}\NormalTok{) }\SpecialCharTok{|\textgreater{}}
 \FunctionTok{select}\NormalTok{(journal, keyword, strength)}
\end{Highlighting}
\end{Shaded}

\emph{Key insight:} Different journals serve different research communities; emerging
terms may overlap only for broad methodological advances.

\textbf{3. Novelty scoring.}
Compute document novelty using cosine distance to prior-year centroids.

\begin{Shaded}
\begin{Highlighting}[]
\FunctionTok{library}\NormalTok{(quanteda)}
\FunctionTok{library}\NormalTok{(tidyverse)}

\NormalTok{dfm\_yearly }\OtherTok{\textless{}{-}} \FunctionTok{dfm}\NormalTok{(}\FunctionTok{tokens}\NormalTok{(}\FunctionTok{corpus}\NormalTok{(works, }\AttributeTok{text\_field =} \StringTok{"ab"}\NormalTok{),}
                         \AttributeTok{remove\_punct =} \ConstantTok{TRUE}\NormalTok{) }\SpecialCharTok{|\textgreater{}}
                  \FunctionTok{tokens\_tolower}\NormalTok{() }\SpecialCharTok{|\textgreater{}}
                  \FunctionTok{tokens\_remove}\NormalTok{(}\FunctionTok{stopwords}\NormalTok{(}\StringTok{"en"}\NormalTok{))) }\SpecialCharTok{|\textgreater{}}
 \FunctionTok{dfm\_tfidf}\NormalTok{()}

\NormalTok{years }\OtherTok{\textless{}{-}} \FunctionTok{sort}\NormalTok{(}\FunctionTok{unique}\NormalTok{(works}\SpecialCharTok{$}\NormalTok{publication\_year))}

\NormalTok{novelty }\OtherTok{\textless{}{-}} \FunctionTok{map}\NormalTok{(years[}\SpecialCharTok{{-}}\DecValTok{1}\NormalTok{], \textbackslash{}(yr) \{}
\NormalTok{ prior }\OtherTok{\textless{}{-}} \FunctionTok{dfm\_subset}\NormalTok{(dfm\_yearly, publication\_year }\SpecialCharTok{\textless{}}\NormalTok{ yr)}
\NormalTok{ current }\OtherTok{\textless{}{-}} \FunctionTok{dfm\_subset}\NormalTok{(dfm\_yearly, publication\_year }\SpecialCharTok{==}\NormalTok{ yr)}
 \ControlFlowTok{if}\NormalTok{ (}\FunctionTok{nrow}\NormalTok{(prior) }\SpecialCharTok{==} \DecValTok{0} \SpecialCharTok{||} \FunctionTok{nrow}\NormalTok{(current) }\SpecialCharTok{==} \DecValTok{0}\NormalTok{) }\FunctionTok{return}\NormalTok{(}\ConstantTok{NULL}\NormalTok{)}

\NormalTok{ centroid }\OtherTok{\textless{}{-}} \FunctionTok{colMeans}\NormalTok{(}\FunctionTok{as.matrix}\NormalTok{(prior))}
\NormalTok{ current\_mat }\OtherTok{\textless{}{-}} \FunctionTok{as.matrix}\NormalTok{(current)}

\NormalTok{ cos\_dist }\OtherTok{\textless{}{-}} \DecValTok{1} \SpecialCharTok{{-}}\NormalTok{ (current\_mat }\SpecialCharTok{\%*\%}\NormalTok{ centroid) }\SpecialCharTok{/}
\NormalTok{   (}\FunctionTok{sqrt}\NormalTok{(}\FunctionTok{rowSums}\NormalTok{(current\_mat}\SpecialCharTok{\^{}}\DecValTok{2}\NormalTok{)) }\SpecialCharTok{*} \FunctionTok{sqrt}\NormalTok{(}\FunctionTok{sum}\NormalTok{(centroid}\SpecialCharTok{\^{}}\DecValTok{2}\NormalTok{)))}

 \FunctionTok{tibble}\NormalTok{(}\AttributeTok{year =}\NormalTok{ yr, }\AttributeTok{novelty =} \FunctionTok{as.numeric}\NormalTok{(cos\_dist))}
\NormalTok{\}) }\SpecialCharTok{|\textgreater{}} \FunctionTok{list\_rbind}\NormalTok{()}

\NormalTok{novelty }\SpecialCharTok{|\textgreater{}}
 \FunctionTok{group\_by}\NormalTok{(year) }\SpecialCharTok{|\textgreater{}}
 \FunctionTok{summarise}\NormalTok{(}\AttributeTok{mean\_novelty =} \FunctionTok{mean}\NormalTok{(novelty, }\AttributeTok{na.rm =} \ConstantTok{TRUE}\NormalTok{)) }\SpecialCharTok{|\textgreater{}}
 \FunctionTok{ggplot}\NormalTok{(}\FunctionTok{aes}\NormalTok{(year, mean\_novelty)) }\SpecialCharTok{+}
 \FunctionTok{geom\_line}\NormalTok{() }\SpecialCharTok{+}
 \FunctionTok{labs}\NormalTok{(}\AttributeTok{title =} \StringTok{"Document novelty over time"}\NormalTok{,}
      \AttributeTok{x =} \StringTok{"Year"}\NormalTok{, }\AttributeTok{y =} \StringTok{"Mean cosine distance to prior centroid"}\NormalTok{)}
\end{Highlighting}
\end{Shaded}

\emph{Key insight:} Increasing novelty scores suggest the field is exploring new
territory; stable scores indicate incremental progress within established
paradigms.

\textbf{4. Validation.}
Manually validate burst-detected emerging terms.

\begin{Shaded}
\begin{Highlighting}[]
\CommentTok{\# Select 5 terms identified as "emerging"}
\NormalTok{emerging\_terms }\OtherTok{\textless{}{-}}\NormalTok{ bursts\_relative }\SpecialCharTok{|\textgreater{}}
 \FunctionTok{arrange}\NormalTok{(}\FunctionTok{desc}\NormalTok{(strength)) }\SpecialCharTok{|\textgreater{}}
 \FunctionTok{head}\NormalTok{(}\DecValTok{5}\NormalTok{) }\SpecialCharTok{|\textgreater{}}
 \FunctionTok{pull}\NormalTok{(keyword)}

\CommentTok{\# For each term, search for corroborating evidence:}
\CommentTok{\# {-} Is there a review article about this topic?}
\CommentTok{\# {-} Was it a conference theme in recent years?}
\CommentTok{\# {-} Does Google Trends show a rising interest?}

\CommentTok{\# Create a validation table}
\NormalTok{validation }\OtherTok{\textless{}{-}} \FunctionTok{tibble}\NormalTok{(}
 \AttributeTok{term =}\NormalTok{ emerging\_terms,}
 \AttributeTok{review\_found =} \FunctionTok{c}\NormalTok{(}\ConstantTok{TRUE}\NormalTok{, }\ConstantTok{TRUE}\NormalTok{, }\ConstantTok{FALSE}\NormalTok{, }\ConstantTok{TRUE}\NormalTok{, }\ConstantTok{FALSE}\NormalTok{),}
 \AttributeTok{conference\_theme =} \FunctionTok{c}\NormalTok{(}\ConstantTok{TRUE}\NormalTok{, }\ConstantTok{FALSE}\NormalTok{, }\ConstantTok{FALSE}\NormalTok{, }\ConstantTok{TRUE}\NormalTok{, }\ConstantTok{TRUE}\NormalTok{),}
 \AttributeTok{google\_trend\_rising =} \FunctionTok{c}\NormalTok{(}\ConstantTok{TRUE}\NormalTok{, }\ConstantTok{TRUE}\NormalTok{, }\ConstantTok{TRUE}\NormalTok{, }\ConstantTok{FALSE}\NormalTok{, }\ConstantTok{TRUE}\NormalTok{),}
 \AttributeTok{validated =}\NormalTok{ review\_found }\SpecialCharTok{|}\NormalTok{ conference\_theme}
\NormalTok{)}

\FunctionTok{cat}\NormalTok{(}\FunctionTok{sprintf}\NormalTok{(}\StringTok{"False positive rate: \%.0f\%\%}\SpecialCharTok{\textbackslash{}n}\StringTok{"}\NormalTok{,}
\NormalTok{            (}\DecValTok{1} \SpecialCharTok{{-}} \FunctionTok{mean}\NormalTok{(validation}\SpecialCharTok{$}\NormalTok{validated)) }\SpecialCharTok{*} \DecValTok{100}\NormalTok{))}
\end{Highlighting}
\end{Shaded}

\emph{Key insight:} Burst detection has moderate precision; triangulating with
qualitative sources reduces false positives and increases confidence.

\section*{Chapter 26 --- Altmetrics}\label{chapter-26-altmetrics}
\addcontentsline{toc}{section}{Chapter 26 --- Altmetrics}

\textbf{1. OA citation advantage.}
Compare citations for OA vs.~closed-access papers across journals.

\begin{Shaded}
\begin{Highlighting}[]
\FunctionTok{library}\NormalTok{(openalexR)}
\FunctionTok{library}\NormalTok{(tidyverse)}

\NormalTok{journal\_ids }\OtherTok{\textless{}{-}} \FunctionTok{c}\NormalTok{(}\StringTok{"S4210174107"}\NormalTok{, }\StringTok{"S137773608"}\NormalTok{)}

\NormalTok{oa\_advantage }\OtherTok{\textless{}{-}} \FunctionTok{map}\NormalTok{(journal\_ids, \textbackslash{}(jid) \{}
\NormalTok{ w }\OtherTok{\textless{}{-}} \FunctionTok{fetch\_openalex}\NormalTok{(}\AttributeTok{entity =} \StringTok{"works"}\NormalTok{,}
               \AttributeTok{primary\_location.source.id =}\NormalTok{ jid,}
               \AttributeTok{from\_publication\_date =} \StringTok{"2019{-}01{-}01"}\NormalTok{,}
               \AttributeTok{to\_publication\_date =} \StringTok{"2022{-}12{-}31"}\NormalTok{)}
\NormalTok{ w }\SpecialCharTok{|\textgreater{}}
   \FunctionTok{group\_by}\NormalTok{(is\_oa) }\SpecialCharTok{|\textgreater{}}
   \FunctionTok{summarise}\NormalTok{(}
     \AttributeTok{n =} \FunctionTok{n}\NormalTok{(),}
     \AttributeTok{mean\_cites =} \FunctionTok{mean}\NormalTok{(cited\_by\_count),}
     \AttributeTok{median\_cites =} \FunctionTok{median}\NormalTok{(cited\_by\_count)}
\NormalTok{   ) }\SpecialCharTok{|\textgreater{}}
   \FunctionTok{mutate}\NormalTok{(}\AttributeTok{journal =}\NormalTok{ jid)}
\NormalTok{\}) }\SpecialCharTok{|\textgreater{}} \FunctionTok{list\_rbind}\NormalTok{()}

\NormalTok{oa\_advantage}
\end{Highlighting}
\end{Shaded}

\emph{Key insight:} The OA citation advantage exists in most fields but varies in
magnitude. Self-selection (better papers go OA) may partially explain the
effect.

\textbf{2. Temporal attention.}
Track citation accumulation over time for early-attention papers.

\begin{Shaded}
\begin{Highlighting}[]
\NormalTok{recent\_works }\OtherTok{\textless{}{-}}\NormalTok{ works }\SpecialCharTok{|\textgreater{}}
 \FunctionTok{filter}\NormalTok{(publication\_year }\SpecialCharTok{==} \DecValTok{2022}\NormalTok{, }\SpecialCharTok{!}\FunctionTok{is.null}\NormalTok{(counts\_by\_year))}

\NormalTok{attention\_trajectory }\OtherTok{\textless{}{-}}\NormalTok{ recent\_works }\SpecialCharTok{|\textgreater{}}
 \FunctionTok{unnest}\NormalTok{(counts\_by\_year) }\SpecialCharTok{|\textgreater{}}
 \FunctionTok{group\_by}\NormalTok{(id) }\SpecialCharTok{|\textgreater{}}
 \FunctionTok{arrange}\NormalTok{(year) }\SpecialCharTok{|\textgreater{}}
 \FunctionTok{mutate}\NormalTok{(}\AttributeTok{cum\_cites =} \FunctionTok{cumsum}\NormalTok{(cited\_by\_count)) }\SpecialCharTok{|\textgreater{}}
 \FunctionTok{ungroup}\NormalTok{()}

\CommentTok{\# Classify: high vs. low early attention (first year citations)}
\NormalTok{early\_cites }\OtherTok{\textless{}{-}}\NormalTok{ attention\_trajectory }\SpecialCharTok{|\textgreater{}}
 \FunctionTok{filter}\NormalTok{(year }\SpecialCharTok{==} \DecValTok{2022}\NormalTok{) }\SpecialCharTok{|\textgreater{}}
 \FunctionTok{mutate}\NormalTok{(}\AttributeTok{early\_attention =} \FunctionTok{ifelse}\NormalTok{(cited\_by\_count }\SpecialCharTok{\textgreater{}=} \DecValTok{5}\NormalTok{,}
                                  \StringTok{"High"}\NormalTok{, }\StringTok{"Low"}\NormalTok{))}

\NormalTok{attention\_trajectory }\SpecialCharTok{|\textgreater{}}
 \FunctionTok{left\_join}\NormalTok{(early\_cites }\SpecialCharTok{|\textgreater{}} \FunctionTok{select}\NormalTok{(id, early\_attention), }\AttributeTok{by =} \StringTok{"id"}\NormalTok{) }\SpecialCharTok{|\textgreater{}}
 \FunctionTok{group\_by}\NormalTok{(early\_attention, year) }\SpecialCharTok{|\textgreater{}}
 \FunctionTok{summarise}\NormalTok{(}\AttributeTok{mean\_cum =} \FunctionTok{mean}\NormalTok{(cum\_cites, }\AttributeTok{na.rm =} \ConstantTok{TRUE}\NormalTok{)) }\SpecialCharTok{|\textgreater{}}
 \FunctionTok{ggplot}\NormalTok{(}\FunctionTok{aes}\NormalTok{(year, mean\_cum, }\AttributeTok{colour =}\NormalTok{ early\_attention)) }\SpecialCharTok{+}
 \FunctionTok{geom\_line}\NormalTok{() }\SpecialCharTok{+}
 \FunctionTok{labs}\NormalTok{(}\AttributeTok{title =} \StringTok{"Citation accumulation by early attention"}\NormalTok{)}
\end{Highlighting}
\end{Shaded}

\emph{Key insight:} Papers with high early attention tend to sustain their advantage,
consistent with cumulative advantage (Matthew Effect) in citations.

\textbf{3. Cross-field comparison.}
Compare citation distributions and reference list lengths across fields.

\begin{Shaded}
\begin{Highlighting}[]
\NormalTok{biomed }\OtherTok{\textless{}{-}} \FunctionTok{fetch\_openalex}\NormalTok{(}\AttributeTok{entity =} \StringTok{"works"}\NormalTok{,}
                   \AttributeTok{primary\_location.source.id =} \StringTok{"S49861241"}\NormalTok{,}
                   \AttributeTok{from\_publication\_date =} \StringTok{"2020{-}01{-}01"}\NormalTok{) }\SpecialCharTok{|\textgreater{}}
 \FunctionTok{mutate}\NormalTok{(}\AttributeTok{field =} \StringTok{"Biomedicine"}\NormalTok{) }\SpecialCharTok{|\textgreater{}}
 \FunctionTok{slice\_sample}\NormalTok{(}\AttributeTok{n =} \DecValTok{200}\NormalTok{)}

\NormalTok{math }\OtherTok{\textless{}{-}} \FunctionTok{fetch\_openalex}\NormalTok{(}\AttributeTok{entity =} \StringTok{"works"}\NormalTok{,}
                 \AttributeTok{primary\_location.source.id =} \StringTok{"S154589771"}\NormalTok{,}
                 \AttributeTok{from\_publication\_date =} \StringTok{"2020{-}01{-}01"}\NormalTok{) }\SpecialCharTok{|\textgreater{}}
 \FunctionTok{mutate}\NormalTok{(}\AttributeTok{field =} \StringTok{"Mathematics"}\NormalTok{) }\SpecialCharTok{|\textgreater{}}
 \FunctionTok{slice\_sample}\NormalTok{(}\AttributeTok{n =} \DecValTok{200}\NormalTok{)}

\NormalTok{comparison }\OtherTok{\textless{}{-}} \FunctionTok{bind\_rows}\NormalTok{(biomed, math)}

\NormalTok{comparison }\SpecialCharTok{|\textgreater{}}
 \FunctionTok{group\_by}\NormalTok{(field) }\SpecialCharTok{|\textgreater{}}
 \FunctionTok{summarise}\NormalTok{(}
   \AttributeTok{mean\_cites =} \FunctionTok{mean}\NormalTok{(cited\_by\_count),}
   \AttributeTok{median\_cites =} \FunctionTok{median}\NormalTok{(cited\_by\_count),}
   \AttributeTok{mean\_refs =} \FunctionTok{mean}\NormalTok{(referenced\_works\_count, }\AttributeTok{na.rm =} \ConstantTok{TRUE}\NormalTok{)}
\NormalTok{ )}

\FunctionTok{ggplot}\NormalTok{(comparison, }\FunctionTok{aes}\NormalTok{(cited\_by\_count, }\AttributeTok{fill =}\NormalTok{ field)) }\SpecialCharTok{+}
 \FunctionTok{geom\_histogram}\NormalTok{(}\AttributeTok{position =} \StringTok{"dodge"}\NormalTok{, }\AttributeTok{bins =} \DecValTok{30}\NormalTok{) }\SpecialCharTok{+}
 \FunctionTok{scale\_x\_log10}\NormalTok{() }\SpecialCharTok{+}
 \FunctionTok{labs}\NormalTok{(}\AttributeTok{title =} \StringTok{"Citation distributions by field"}\NormalTok{)}
\end{Highlighting}
\end{Shaded}

\emph{Key insight:} Biomedical papers have higher citation counts and longer reference
lists than mathematics papers. This is why raw citation counts cannot be compared
across fields.

\section*{Chapter 27 --- Open Science Indicators}\label{chapter-27-open-science-indicators}
\addcontentsline{toc}{section}{Chapter 27 --- Open Science Indicators}

\textbf{1. Cross-journal OA comparison.}
Compare OA rates across three journals in different disciplines.

\begin{Shaded}
\begin{Highlighting}[]
\FunctionTok{library}\NormalTok{(openalexR)}
\FunctionTok{library}\NormalTok{(tidyverse)}

\NormalTok{journals }\OtherTok{\textless{}{-}} \FunctionTok{list}\NormalTok{(}
 \AttributeTok{biology =} \StringTok{"S49861241"}\NormalTok{,}
 \AttributeTok{physics =} \StringTok{"S137773608"}\NormalTok{,}
 \AttributeTok{sociology =} \StringTok{"S64468286"}
\NormalTok{)}

\NormalTok{oa\_rates }\OtherTok{\textless{}{-}} \FunctionTok{imap}\NormalTok{(journals, \textbackslash{}(jid, field) \{}
\NormalTok{ w }\OtherTok{\textless{}{-}} \FunctionTok{fetch\_openalex}\NormalTok{(}\AttributeTok{entity =} \StringTok{"works"}\NormalTok{,}
               \AttributeTok{primary\_location.source.id =}\NormalTok{ jid,}
               \AttributeTok{from\_publication\_date =} \StringTok{"2020{-}01{-}01"}\NormalTok{,}
               \AttributeTok{to\_publication\_date =} \StringTok{"2023{-}12{-}31"}\NormalTok{)}
\NormalTok{ w }\SpecialCharTok{|\textgreater{}}
   \FunctionTok{summarise}\NormalTok{(}\AttributeTok{oa\_rate =} \FunctionTok{mean}\NormalTok{(is\_oa, }\AttributeTok{na.rm =} \ConstantTok{TRUE}\NormalTok{),}
             \AttributeTok{gold\_oa =} \FunctionTok{mean}\NormalTok{(oa\_status }\SpecialCharTok{==} \StringTok{"gold"}\NormalTok{, }\AttributeTok{na.rm =} \ConstantTok{TRUE}\NormalTok{)) }\SpecialCharTok{|\textgreater{}}
   \FunctionTok{mutate}\NormalTok{(}\AttributeTok{field =}\NormalTok{ field)}
\NormalTok{\}) }\SpecialCharTok{|\textgreater{}} \FunctionTok{list\_rbind}\NormalTok{()}

\NormalTok{oa\_rates}
\end{Highlighting}
\end{Shaded}

\emph{Key insight:} Biology tends to have the highest gold OA rate (driven by PLOS,
eLife, etc.), while sociology and mathematics lag behind.

\textbf{2. Temporal OA growth.}
Track the year-over-year OA growth rate.

\begin{Shaded}
\begin{Highlighting}[]
\NormalTok{oa\_growth }\OtherTok{\textless{}{-}}\NormalTok{ works }\SpecialCharTok{|\textgreater{}}
 \FunctionTok{filter}\NormalTok{(is\_oa) }\SpecialCharTok{|\textgreater{}}
 \FunctionTok{count}\NormalTok{(publication\_year, }\AttributeTok{name =} \StringTok{"n\_oa"}\NormalTok{) }\SpecialCharTok{|\textgreater{}}
 \FunctionTok{mutate}\NormalTok{(}\AttributeTok{yoy\_growth =}\NormalTok{ (n\_oa }\SpecialCharTok{/} \FunctionTok{lag}\NormalTok{(n\_oa)) }\SpecialCharTok{{-}} \DecValTok{1}\NormalTok{)}

\FunctionTok{ggplot}\NormalTok{(oa\_growth, }\FunctionTok{aes}\NormalTok{(publication\_year, yoy\_growth)) }\SpecialCharTok{+}
 \FunctionTok{geom\_col}\NormalTok{() }\SpecialCharTok{+}
 \FunctionTok{scale\_y\_continuous}\NormalTok{(}\AttributeTok{labels =}\NormalTok{ scales}\SpecialCharTok{::}\NormalTok{percent) }\SpecialCharTok{+}
 \FunctionTok{labs}\NormalTok{(}\AttributeTok{title =} \StringTok{"Year{-}over{-}year OA growth"}\NormalTok{,}
      \AttributeTok{x =} \StringTok{"Year"}\NormalTok{, }\AttributeTok{y =} \StringTok{"Growth rate"}\NormalTok{)}
\end{Highlighting}
\end{Shaded}

\emph{Key insight:} OA growth accelerated around 2020, likely driven by COVID-19
mandates and Plan S implementation.

\textbf{3. OA and collaboration.}
Test whether author count predicts OA status.

\begin{Shaded}
\begin{Highlighting}[]
\NormalTok{works }\SpecialCharTok{|\textgreater{}}
 \FunctionTok{mutate}\NormalTok{(}\AttributeTok{n\_authors =} \FunctionTok{map\_int}\NormalTok{(author, nrow)) }\SpecialCharTok{|\textgreater{}}
 \FunctionTok{group\_by}\NormalTok{(}\AttributeTok{n\_authors =} \FunctionTok{pmin}\NormalTok{(n\_authors, }\DecValTok{10}\NormalTok{)) }\SpecialCharTok{|\textgreater{}}
 \FunctionTok{summarise}\NormalTok{(}\AttributeTok{oa\_rate =} \FunctionTok{mean}\NormalTok{(is\_oa, }\AttributeTok{na.rm =} \ConstantTok{TRUE}\NormalTok{), }\AttributeTok{n =} \FunctionTok{n}\NormalTok{()) }\SpecialCharTok{|\textgreater{}}
 \FunctionTok{ggplot}\NormalTok{(}\FunctionTok{aes}\NormalTok{(n\_authors, oa\_rate)) }\SpecialCharTok{+}
 \FunctionTok{geom\_col}\NormalTok{() }\SpecialCharTok{+}
 \FunctionTok{scale\_y\_continuous}\NormalTok{(}\AttributeTok{labels =}\NormalTok{ scales}\SpecialCharTok{::}\NormalTok{percent) }\SpecialCharTok{+}
 \FunctionTok{labs}\NormalTok{(}\AttributeTok{title =} \StringTok{"OA rate by number of authors"}\NormalTok{,}
      \AttributeTok{x =} \StringTok{"Authors (capped at 10)"}\NormalTok{, }\AttributeTok{y =} \StringTok{"OA rate"}\NormalTok{)}
\end{Highlighting}
\end{Shaded}

\emph{Key insight:} Larger teams may have more funding to cover APCs. International
collaborations also increase OA probability because at least one funder may
mandate it.

\section*{Chapter 28 --- Gender and Equity Analysis}\label{chapter-28-gender-and-equity-analysis}
\addcontentsline{toc}{section}{Chapter 28 --- Gender and Equity Analysis}

\textbf{1. Denominator comparison.}
Use the explicitly synthetic \texttt{lookup\_records} from Chapter 28.

\begin{Shaded}
\begin{Highlighting}[]
\NormalTok{all\_records }\OtherTok{\textless{}{-}}\NormalTok{ lookup\_records }\SpecialCharTok{|\textgreater{}}
  \FunctionTok{count}\NormalTok{(lookup\_category) }\SpecialCharTok{|\textgreater{}}
  \FunctionTok{mutate}\NormalTok{(}\AttributeTok{share =}\NormalTok{ n }\SpecialCharTok{/} \FunctionTok{sum}\NormalTok{(n))}
\NormalTok{matched\_only }\OtherTok{\textless{}{-}}\NormalTok{ lookup\_records }\SpecialCharTok{|\textgreater{}}
  \FunctionTok{filter}\NormalTok{(lookup\_category }\SpecialCharTok{!=} \StringTok{"Unmatched"}\NormalTok{) }\SpecialCharTok{|\textgreater{}}
  \FunctionTok{count}\NormalTok{(lookup\_category) }\SpecialCharTok{|\textgreater{}}
  \FunctionTok{mutate}\NormalTok{(}\AttributeTok{share =}\NormalTok{ n }\SpecialCharTok{/} \FunctionTok{sum}\NormalTok{(n))}
\NormalTok{all\_records}
\NormalTok{matched\_only}
\end{Highlighting}
\end{Shaded}

The fixture has 12 records: four A, two B and six Unmatched. Among all records,
A is one third and B one sixth; among matched records, A is two thirds and B one
third. None of these arbitrary category shares measures gender representation.

\textbf{2. Position audit.}

\begin{Shaded}
\begin{Highlighting}[]
\NormalTok{position\_data }\SpecialCharTok{|\textgreater{}} \FunctionTok{count}\NormalTok{(author\_position)}
\NormalTok{position\_summary }\SpecialCharTok{|\textgreater{}}
  \FunctionTok{group\_by}\NormalTok{(author\_position) }\SpecialCharTok{|\textgreater{}}
  \FunctionTok{summarise}\NormalTok{(}\AttributeTok{proportion\_sum =} \FunctionTok{sum}\NormalTok{(pct), }\AttributeTok{records =} \FunctionTok{sum}\NormalTok{(n), }\AttributeTok{.groups =} \StringTok{"drop"}\NormalTok{)}
\end{Highlighting}
\end{Shaded}

The invented lists produce four first, four last, three middle and one
single-author record. Each position's proportions sum to one because unmatched
records remain included. Recorded position does not establish a contribution role.

\textbf{3. Study design.}

Specify the scientific question and population before defining an analysis.
A real study would need a defensible demographic variable and its provenance,
a consent/data-governance plan, categories appropriate to that question,
missing-data and validation procedures, and evidence about authorship conventions.
The fixture supplies none of those demographic observations, nor real identity
linkage, publication years or citation outcomes. Changing a name dictionary cannot
recover them. No citation-gap or gender-equity conclusion follows from this exercise.

\section*{Chapter 29 --- Retraction Analysis}\label{chapter-29-retraction-analysis}
\addcontentsline{toc}{section}{Chapter 29 --- Retraction Analysis}

\textbf{1. Retraction rate by field.}
Compare retraction rates across fields, normalised by corpus size.

\begin{Shaded}
\begin{Highlighting}[]
\FunctionTok{library}\NormalTok{(openalexR)}
\FunctionTok{library}\NormalTok{(tidyverse)}

\NormalTok{fields }\OtherTok{\textless{}{-}} \FunctionTok{c}\NormalTok{(}\StringTok{"biomedicine"}\NormalTok{, }\StringTok{"computer science"}\NormalTok{)}

\NormalTok{retraction\_rates }\OtherTok{\textless{}{-}} \FunctionTok{map}\NormalTok{(fields, \textbackslash{}(f) \{}
\NormalTok{ all\_works }\OtherTok{\textless{}{-}} \FunctionTok{fetch\_openalex}\NormalTok{(}\AttributeTok{entity =} \StringTok{"works"}\NormalTok{, }\AttributeTok{search =}\NormalTok{ f,}
                       \AttributeTok{from\_publication\_date =} \StringTok{"2015{-}01{-}01"}\NormalTok{,}
                       \AttributeTok{count\_only =} \ConstantTok{TRUE}\NormalTok{)}
\NormalTok{ retracted }\OtherTok{\textless{}{-}} \FunctionTok{fetch\_openalex}\NormalTok{(}\AttributeTok{entity =} \StringTok{"works"}\NormalTok{, }\AttributeTok{search =}\NormalTok{ f,}
                       \AttributeTok{type =} \StringTok{"retraction"}\NormalTok{,}
                       \AttributeTok{from\_publication\_date =} \StringTok{"2015{-}01{-}01"}\NormalTok{,}
                       \AttributeTok{count\_only =} \ConstantTok{TRUE}\NormalTok{)}
 \FunctionTok{tibble}\NormalTok{(}
   \AttributeTok{field =}\NormalTok{ f,}
   \AttributeTok{total =}\NormalTok{ all\_works}\SpecialCharTok{$}\NormalTok{count,}
   \AttributeTok{retracted =}\NormalTok{ retracted}\SpecialCharTok{$}\NormalTok{count,}
   \AttributeTok{rate\_per\_1000 =}\NormalTok{ retracted}\SpecialCharTok{$}\NormalTok{count }\SpecialCharTok{/}\NormalTok{ all\_works}\SpecialCharTok{$}\NormalTok{count }\SpecialCharTok{*} \DecValTok{1000}
\NormalTok{ )}
\NormalTok{\}) }\SpecialCharTok{|\textgreater{}} \FunctionTok{list\_rbind}\NormalTok{()}

\NormalTok{retraction\_rates}
\end{Highlighting}
\end{Shaded}

\emph{Key insight:} Retraction rates are typically higher in biomedical fields (1-2
per 1,000) than in computer science, partly reflecting different scrutiny levels
and reproducibility challenges.

\textbf{2. Citation trajectory.}
Examine when a retracted paper's citations were received.

\begin{Shaded}
\begin{Highlighting}[]
\NormalTok{retracted\_paper }\OtherTok{\textless{}{-}} \FunctionTok{fetch\_openalex}\NormalTok{(}\AttributeTok{entity =} \StringTok{"works"}\NormalTok{,}
                            \AttributeTok{doi =} \StringTok{"10.xxxx/retracted{-}example"}\NormalTok{)}

\ControlFlowTok{if}\NormalTok{ (}\SpecialCharTok{!}\FunctionTok{is.null}\NormalTok{(retracted\_paper}\SpecialCharTok{$}\NormalTok{counts\_by\_year)) \{}
\NormalTok{ retracted\_paper }\SpecialCharTok{|\textgreater{}}
   \FunctionTok{unnest}\NormalTok{(counts\_by\_year) }\SpecialCharTok{|\textgreater{}}
   \FunctionTok{ggplot}\NormalTok{(}\FunctionTok{aes}\NormalTok{(year, cited\_by\_count)) }\SpecialCharTok{+}
   \FunctionTok{geom\_col}\NormalTok{() }\SpecialCharTok{+}
   \FunctionTok{geom\_vline}\NormalTok{(}\AttributeTok{xintercept =} \DecValTok{2020}\NormalTok{, }\AttributeTok{linetype =} \StringTok{"dashed"}\NormalTok{,}
              \AttributeTok{colour =} \StringTok{"red"}\NormalTok{) }\SpecialCharTok{+}  \CommentTok{\# retraction year}
   \FunctionTok{labs}\NormalTok{(}\AttributeTok{title =} \StringTok{"Citation trajectory of a retracted paper"}\NormalTok{,}
        \AttributeTok{x =} \StringTok{"Year"}\NormalTok{, }\AttributeTok{y =} \StringTok{"Citations received"}\NormalTok{)}
\NormalTok{\}}
\end{Highlighting}
\end{Shaded}

\emph{Key insight:} Many retracted papers continue to be cited after retraction,
indicating that the retraction notice does not always reach citing authors.

\textbf{3. Retraction-citation correlation.}
Test whether pre-retraction citation count predicts post-retraction citations.

\begin{Shaded}
\begin{Highlighting}[]
\CommentTok{\# Assuming a set of retracted papers with known retraction years}
\NormalTok{retracted\_set }\OtherTok{\textless{}{-}}\NormalTok{ retracted\_works }\SpecialCharTok{|\textgreater{}}
 \FunctionTok{unnest}\NormalTok{(counts\_by\_year) }\SpecialCharTok{|\textgreater{}}
 \FunctionTok{mutate}\NormalTok{(}\AttributeTok{period =} \FunctionTok{ifelse}\NormalTok{(year }\SpecialCharTok{\textless{}}\NormalTok{ retraction\_year, }\StringTok{"pre"}\NormalTok{, }\StringTok{"post"}\NormalTok{)) }\SpecialCharTok{|\textgreater{}}
 \FunctionTok{group\_by}\NormalTok{(id, period) }\SpecialCharTok{|\textgreater{}}
 \FunctionTok{summarise}\NormalTok{(}\AttributeTok{cites =} \FunctionTok{sum}\NormalTok{(cited\_by\_count), }\AttributeTok{.groups =} \StringTok{"drop"}\NormalTok{) }\SpecialCharTok{|\textgreater{}}
 \FunctionTok{pivot\_wider}\NormalTok{(}\AttributeTok{names\_from =}\NormalTok{ period, }\AttributeTok{values\_from =}\NormalTok{ cites,}
             \AttributeTok{values\_fill =} \DecValTok{0}\NormalTok{)}

\FunctionTok{cor}\NormalTok{(retracted\_set}\SpecialCharTok{$}\NormalTok{pre, retracted\_set}\SpecialCharTok{$}\NormalTok{post, }\AttributeTok{method =} \StringTok{"spearman"}\NormalTok{)}

\FunctionTok{ggplot}\NormalTok{(retracted\_set, }\FunctionTok{aes}\NormalTok{(pre, post)) }\SpecialCharTok{+}
 \FunctionTok{geom\_point}\NormalTok{(}\AttributeTok{alpha =} \FloatTok{0.5}\NormalTok{) }\SpecialCharTok{+}
 \FunctionTok{geom\_smooth}\NormalTok{(}\AttributeTok{method =} \StringTok{"lm"}\NormalTok{) }\SpecialCharTok{+}
 \FunctionTok{labs}\NormalTok{(}\AttributeTok{x =} \StringTok{"Pre{-}retraction citations"}\NormalTok{,}
      \AttributeTok{y =} \StringTok{"Post{-}retraction citations"}\NormalTok{,}
      \AttributeTok{title =} \StringTok{"Pre{-} vs. post{-}retraction citation counts"}\NormalTok{)}
\end{Highlighting}
\end{Shaded}

\emph{Key insight:} A positive correlation suggests that highly visible papers
continue to be cited after retraction --- the ``retraction paradox.''

\section*{Chapter 30 --- Funding Analysis}\label{chapter-30-funding-analysis}
\addcontentsline{toc}{section}{Chapter 30 --- Funding Analysis}

\textbf{1. Funder co-occurrence.}
Identify frequently co-occurring funder pairs.

\begin{Shaded}
\begin{Highlighting}[]
\FunctionTok{library}\NormalTok{(tidyverse)}

\NormalTok{funder\_pairs }\OtherTok{\textless{}{-}} \FunctionTok{inner\_join}\NormalTok{(}
  \FunctionTok{select}\NormalTok{(funder\_data, work\_id, }\AttributeTok{funder\_a =}\NormalTok{ funder\_id),}
  \FunctionTok{select}\NormalTok{(funder\_data, work\_id, }\AttributeTok{funder\_b =}\NormalTok{ funder\_id),}
  \AttributeTok{by =} \StringTok{"work\_id"}\NormalTok{, }\AttributeTok{relationship =} \StringTok{"many{-}to{-}many"}
\NormalTok{) }\SpecialCharTok{|\textgreater{}}
  \FunctionTok{filter}\NormalTok{(funder\_a }\SpecialCharTok{\textless{}}\NormalTok{ funder\_b) }\SpecialCharTok{|\textgreater{}}
  \FunctionTok{distinct}\NormalTok{(work\_id, funder\_a, funder\_b) }\SpecialCharTok{|\textgreater{}}
  \FunctionTok{count}\NormalTok{(funder\_a, funder\_b, }\AttributeTok{sort =} \ConstantTok{TRUE}\NormalTok{)}

\NormalTok{funder\_pairs }\SpecialCharTok{|\textgreater{}} \FunctionTok{head}\NormalTok{(}\DecValTok{10}\NormalTok{)}
\end{Highlighting}
\end{Shaded}

This uses the cleaned, observed funder IDs from Chapter 30 and counts each pair
once per sampled work. An empty result means no co-funding pair was observed in
this sample; it is not evidence that the funders never collaborate.

\textbf{2. Funder portfolios.}
Prepare a bounded comparison for two funders observed in the sample.

\begin{Shaded}
\begin{Highlighting}[]
\FunctionTok{library}\NormalTok{(tidyverse)}

\NormalTok{funder\_ids }\OtherTok{\textless{}{-}} \FunctionTok{head}\NormalTok{(funder\_counts}\SpecialCharTok{$}\NormalTok{funder\_id, }\DecValTok{2}\NormalTok{)}
\ControlFlowTok{if}\NormalTok{ (}\FunctionTok{length}\NormalTok{(funder\_ids) }\SpecialCharTok{\textless{}} \DecValTok{2}\DataTypeTok{L}\NormalTok{) \{}
  \FunctionTok{cat}\NormalTok{(}\StringTok{"A two{-}funder comparison is unavailable in this sample.}\SpecialCharTok{\textbackslash{}n}\StringTok{"}\NormalTok{)}
\NormalTok{\} }\ControlFlowTok{else}\NormalTok{ \{}
\NormalTok{  portfolio\_sample }\OtherTok{\textless{}{-}}\NormalTok{ funder\_data }\SpecialCharTok{|\textgreater{}}
    \FunctionTok{filter}\NormalTok{(funder\_id }\SpecialCharTok{\%in\%}\NormalTok{ funder\_ids) }\SpecialCharTok{|\textgreater{}}
    \FunctionTok{select}\NormalTok{(work\_id, funder\_id) }\SpecialCharTok{|\textgreater{}}
    \FunctionTok{distinct}\NormalTok{() }\SpecialCharTok{|\textgreater{}}
    \FunctionTok{left\_join}\NormalTok{(}\FunctionTok{select}\NormalTok{(works, id, display\_name, abstract),}
              \AttributeTok{by =} \FunctionTok{c}\NormalTok{(}\StringTok{"work\_id"} \OtherTok{=} \StringTok{"id"}\NormalTok{))}
\NormalTok{  portfolio\_sample }\SpecialCharTok{|\textgreater{}} \FunctionTok{count}\NormalTok{(funder\_id, }\AttributeTok{name =} \StringTok{"n\_sampled\_works"}\NormalTok{)}
\NormalTok{\}}
\end{Highlighting}
\end{Shaded}

Apply the text-processing workflow from Chapter 20 to the available titles and
abstracts, reporting text coverage per funder. This corpus contains only the
sampled journal's papers, so it is not either funder's full portfolio. Completing
the exercise for whole portfolios requires a documented retrieval for validated
funder IDs, the same publication window, and verified funding metadata coverage.
Do not substitute invented identifiers or infer strategic priorities from a
small journal sample.

\textbf{3. Temporal trends.}
Track funding acknowledgement coverage over time.

\begin{Shaded}
\begin{Highlighting}[]
\NormalTok{funding\_trend }\OtherTok{\textless{}{-}}\NormalTok{ funding\_works }\SpecialCharTok{|\textgreater{}}
 \FunctionTok{mutate}\NormalTok{(}\AttributeTok{publication\_year =}\NormalTok{ lubridate}\SpecialCharTok{::}\FunctionTok{year}\NormalTok{(publication\_date)) }\SpecialCharTok{|\textgreater{}}
 \FunctionTok{group\_by}\NormalTok{(publication\_year) }\SpecialCharTok{|\textgreater{}}
 \FunctionTok{summarise}\NormalTok{(}\AttributeTok{n\_sampled =} \FunctionTok{n}\NormalTok{(),}
           \AttributeTok{n\_with\_funder =} \FunctionTok{sum}\NormalTok{(has\_funder\_metadata),}
           \AttributeTok{funding\_rate =} \FunctionTok{mean}\NormalTok{(has\_funder\_metadata), }\AttributeTok{.groups =} \StringTok{"drop"}\NormalTok{)}

\FunctionTok{ggplot}\NormalTok{(funding\_trend, }\FunctionTok{aes}\NormalTok{(publication\_year, funding\_rate)) }\SpecialCharTok{+}
 \FunctionTok{geom\_line}\NormalTok{() }\SpecialCharTok{+}
 \FunctionTok{geom\_point}\NormalTok{() }\SpecialCharTok{+}
 \FunctionTok{scale\_y\_continuous}\NormalTok{(}\AttributeTok{labels =}\NormalTok{ scales}\SpecialCharTok{::}\NormalTok{percent) }\SpecialCharTok{+}
 \FunctionTok{labs}\NormalTok{(}\AttributeTok{title =} \StringTok{"Observed funder metadata by publication year in the sample"}\NormalTok{,}
      \AttributeTok{x =} \StringTok{"Publication year"}\NormalTok{, }\AttributeTok{y =} \StringTok{"Share with observed funder IDs"}\NormalTok{)}
\end{Highlighting}
\end{Shaded}

The rate describes metadata observed at retrieval, grouped by publication year.
It does not measure the prevalence of financial support or establish how the
database's coverage changed historically. Keep the annual sample denominators
and report sparse groups and missing publication years.

\section*{Chapter 31 --- Patent-Paper Linkages}\label{chapter-31-patent-paper-linkages}
\addcontentsline{toc}{section}{Chapter 31 --- Patent-Paper Linkages}

\textbf{1. Scholarly-citation comparison.}
Use two journals represented in the Chapter 31 sample. Keep the common sample
threshold when comparing them, and report the publication years and observed
sample sizes. These are descriptive sample summaries, not journal rankings.

\begin{Shaded}
\begin{Highlighting}[]
\FunctionTok{library}\NormalTok{(tidyverse)}

\NormalTok{selected\_journals }\OtherTok{\textless{}{-}}\NormalTok{ works\_citations }\SpecialCharTok{|\textgreater{}}
  \FunctionTok{filter}\NormalTok{(}\SpecialCharTok{!}\FunctionTok{is.na}\NormalTok{(source\_display\_name)) }\SpecialCharTok{|\textgreater{}}
  \FunctionTok{count}\NormalTok{(source\_display\_name, }\AttributeTok{sort =} \ConstantTok{TRUE}\NormalTok{) }\SpecialCharTok{|\textgreater{}}
  \FunctionTok{slice\_head}\NormalTok{(}\AttributeTok{n =} \DecValTok{2}\NormalTok{) }\SpecialCharTok{|\textgreater{}}
  \FunctionTok{pull}\NormalTok{(source\_display\_name)}

\FunctionTok{stopifnot}\NormalTok{(}\FunctionTok{length}\NormalTok{(selected\_journals) }\SpecialCharTok{==} \DecValTok{2}\NormalTok{)}
\NormalTok{common\_p90 }\OtherTok{\textless{}{-}} \FunctionTok{quantile}\NormalTok{(works\_citations}\SpecialCharTok{$}\NormalTok{cited\_by\_count, }\FloatTok{0.9}\NormalTok{, }\AttributeTok{na.rm =} \ConstantTok{TRUE}\NormalTok{)}

\NormalTok{works\_citations }\SpecialCharTok{|\textgreater{}}
  \FunctionTok{filter}\NormalTok{(source\_display\_name }\SpecialCharTok{\%in\%}\NormalTok{ selected\_journals) }\SpecialCharTok{|\textgreater{}}
  \FunctionTok{group\_by}\NormalTok{(source\_display\_name, year) }\SpecialCharTok{|\textgreater{}}
  \FunctionTok{summarise}\NormalTok{(}
    \AttributeTok{n\_sampled =} \FunctionTok{n}\NormalTok{(),}
    \AttributeTok{n\_citation\_observed =} \FunctionTok{sum}\NormalTok{(}\SpecialCharTok{!}\FunctionTok{is.na}\NormalTok{(cited\_by\_count)),}
    \AttributeTok{median\_cites =} \FunctionTok{median}\NormalTok{(cited\_by\_count, }\AttributeTok{na.rm =} \ConstantTok{TRUE}\NormalTok{),}
    \AttributeTok{mean\_cites =} \FunctionTok{mean}\NormalTok{(cited\_by\_count, }\AttributeTok{na.rm =} \ConstantTok{TRUE}\NormalTok{),}
    \AttributeTok{proportion\_at\_or\_above\_common\_p90 =}
      \FunctionTok{mean}\NormalTok{(cited\_by\_count }\SpecialCharTok{\textgreater{}=}\NormalTok{ common\_p90, }\AttributeTok{na.rm =} \ConstantTok{TRUE}\NormalTok{),}
    \AttributeTok{.groups =} \StringTok{"drop"}
\NormalTok{  )}
\end{Highlighting}
\end{Shaded}

Record the original retrieval date from the data provenance; a cached result's
render date is not its retrieval date. Document field classifications separately
before attributing differences to fields: the worked sample does not supply a
field classification for this comparison. The common threshold is the sample's
90th percentile of current cumulative scholarly citations; ties may put more
than 10\% of works at or above it. A separate 90th percentile for each journal
would put approximately 10\% of each journal above its own threshold by
construction. None of these quantities measures patent citations or predicts
technological impact.

\textbf{2. Linkage data specification.}
The linkage table needs actual source records, with at least these fields:

{\def\LTcaptype{none} % do not increment counter
\begin{longtable}[]{@{}
  >{\raggedright\arraybackslash}p{(\linewidth - 2\tabcolsep) * \real{0.5000}}
  >{\raggedright\arraybackslash}p{(\linewidth - 2\tabcolsep) * \real{0.5000}}@{}}
\toprule\noalign{}
\begin{minipage}[b]{\linewidth}\raggedright
Field
\end{minipage} & \begin{minipage}[b]{\linewidth}\raggedright
Purpose
\end{minipage} \\
\midrule\noalign{}
\endhead
\bottomrule\noalign{}
\endlastfoot
\EscVerb[breaklines=true,breakanywhere=true,breakanywheresymbolpre={},breakanywheresymbolpost={}]{paper\_id} & Stable identifier for the matched scholarly work, such as a DOI or OpenAlex work ID. \\
\EscVerb[breaklines=true,breakanywhere=true,breakanywheresymbolpre={},breakanywheresymbolpost={}]{patent\_id} & Identifier for the citing patent document, including the patent office. \\
\EscVerb[breaklines=true,breakanywhere=true,breakanywheresymbolpre={},breakanywheresymbolpost={}]{patent\_family\_id} & Documented family identifier for deduplicating related patent documents. \\
\EscVerb[breaklines=true,breakanywhere=true,breakanywheresymbolpre={},breakanywheresymbolpost={}]{paper\_publication\_date} & Starting date for a clearly defined paper-to-patent interval. \\
\EscVerb[breaklines=true,breakanywhere=true,breakanywheresymbolpre={},breakanywheresymbolpost={}]{citation\_date} & Date associated with the citation in the source, with its meaning documented; if only patent publication date is available, identify it as that proxy. \\
\EscVerb[breaklines=true,breakanywhere=true,breakanywheresymbolpre={},breakanywheresymbolpost={}]{source} and \EscVerb[breaklines=true,breakanywhere=true,breakanywheresymbolpre={},breakanywheresymbolpost={}]{retrieved\_at} & Patent-reference source, version or snapshot, and retrieval date. \\
\EscVerb[breaklines=true,breakanywhere=true,breakanywheresymbolpre={},breakanywheresymbolpost={}]{raw\_reference} & Original reference text or source record identifier for verification. \\
\EscVerb[breaklines=true,breakanywhere=true,breakanywheresymbolpre={},breakanywheresymbolpost={}]{match\_status} & Verified, ambiguous or unmatched, with the matching method and review evidence retained. \\
\end{longtable}
}

Keep ambiguous and unmatched references in the source table. Review candidate
matches against the original bibliographic record, and report matching coverage
before restricting an analysis to verified links. Missing identifiers or dates
must remain missing rather than being inferred from scholarly citation counts.
No patent-reference dataset is supplied by the worked example, so it cannot
yield an observed linkage rate or lag distribution.

\textbf{3. Institutional patent linkage.}
Define an eligible publication set for each institution using the same years,
document types, affiliation rule and data coverage. Deduplicate publications by
paper identifier, and specify how shared publications are credited. Use a common
patent observation cutoff and, where possible, equal follow-up time for each
publication cohort. Count a paper as observed patent-cited only when at least one
verified patent-family link falls in that window; repeated documents within a
family must not inflate linkage counts.

The observed linked-paper proportion is the number of distinct eligible papers
with a verified in-window patent-family link divided by the number of eligible
papers. Report the numerator, denominator, source coverage, unmatched and
ambiguous records, and sensitivity to matching and follow-up choices. A paper
without a link in the observed source is not proof that no patent cites it.

Without the independently validated patent-reference data, the missing inputs
are the citing patent and family identifiers, reference matches, dates and
coverage information. Report the analysis as unavailable. The scholarly-citation
sample and institutional top-decile shares cannot replace those inputs.

\section*{Chapter 32 --- Causal Inference in Scientometrics}\label{chapter-32-causal-inference-in-scientometrics}
\addcontentsline{toc}{section}{Chapter 32 --- Causal Inference in Scientometrics}

\textbf{1. Interaction sensitivity.}
Use the complete sampled records in \texttt{cohort\_data} from Chapter 32.

\begin{Shaded}
\begin{Highlighting}[]
\NormalTok{unadjusted }\OtherTok{\textless{}{-}} \FunctionTok{lm}\NormalTok{(cited\_by\_count }\SpecialCharTok{\textasciitilde{}}\NormalTok{ is\_oa }\SpecialCharTok{*}\NormalTok{ recent\_cohort, }\AttributeTok{data =}\NormalTok{ cohort\_data)}
\NormalTok{adjusted }\OtherTok{\textless{}{-}} \FunctionTok{lm}\NormalTok{(cited\_by\_count }\SpecialCharTok{\textasciitilde{}}\NormalTok{ is\_oa }\SpecialCharTok{*}\NormalTok{ recent\_cohort }\SpecialCharTok{+}\NormalTok{ referenced\_works\_count,}
               \AttributeTok{data =}\NormalTok{ cohort\_data)}
\NormalTok{broom}\SpecialCharTok{::}\FunctionTok{tidy}\NormalTok{(unadjusted)}
\NormalTok{broom}\SpecialCharTok{::}\FunctionTok{tidy}\NormalTok{(adjusted)}
\end{Highlighting}
\end{Shaded}

The interaction coefficients describe associations in two conditional-mean
specifications. A difference between them indicates specification sensitivity;
it does not identify an OA treatment effect or by itself establish the size or
direction of omitted-variable bias.

\textbf{2. Boundary sensitivity.}

\begin{Shaded}
\begin{Highlighting}[]
\NormalTok{alternative }\OtherTok{\textless{}{-}}\NormalTok{ cohort\_data }\SpecialCharTok{|\textgreater{}} \FunctionTok{mutate}\NormalTok{(}\AttributeTok{cohort\_2018 =}\NormalTok{ year }\SpecialCharTok{\textgreater{}=} \DecValTok{2018}\NormalTok{)}
\NormalTok{alternative\_model }\OtherTok{\textless{}{-}} \FunctionTok{lm}\NormalTok{(cited\_by\_count }\SpecialCharTok{\textasciitilde{}}\NormalTok{ is\_oa }\SpecialCharTok{*}\NormalTok{ cohort\_2018 }\SpecialCharTok{+}
\NormalTok{                         referenced\_works\_count, }\AttributeTok{data =}\NormalTok{ alternative)}
\NormalTok{broom}\SpecialCharTok{::}\FunctionTok{tidy}\NormalTok{(alternative\_model)}
\end{Highlighting}
\end{Shaded}

There is no identified intervention at either date. This changes an arbitrary
publication-cohort partition, so it is a descriptive sensitivity check, not a
placebo intervention test. Neither a significant nor a non-significant interaction
establishes parallel counterfactual trends; these data contain different cohorts'
cumulative citations at retrieval, not annual outcomes around a policy change.

\textbf{3. Comparison audit.}

\begin{Shaded}
\begin{Highlighting}[]
\NormalTok{matched }\SpecialCharTok{|\textgreater{}} \FunctionTok{count}\NormalTok{(control\_id, }\AttributeTok{name =} \StringTok{"times\_used"}\NormalTok{) }\SpecialCharTok{|\textgreater{}} \FunctionTok{arrange}\NormalTok{(}\FunctionTok{desc}\NormalTok{(times\_used))}
\FunctionTok{nrow}\NormalTok{(oa\_papers) }\SpecialCharTok{{-}} \FunctionTok{nrow}\NormalTok{(matched)}
\end{Highlighting}
\end{Shaded}

The original within-year draw can reuse a closed-access paper for multiple OA
papers. OA papers from years lacking a closed-access candidate are omitted. Such
selection changes the comparison population and does not establish exchangeability.
An actual policy study must first define treatment assignment and timing, a target
causal effect, outcome measurement over comparable exposure periods, and a defensible
comparison design. It also needs relevant pre-treatment covariates and an explicit
account of identification and inference assumptions.

\section*{Chapter 33 --- Reproducibility, RRIDs, and Badges}\label{chapter-33-reproducibility-rrids-and-badges}
\addcontentsline{toc}{section}{Chapter 33 --- Reproducibility, RRIDs, and Badges}

\textbf{1. DAS detection.}
Screen available abstracts for candidate data/code availability mentions, then
manually validate candidates against full texts and linked resources.

\begin{Shaded}
\begin{Highlighting}[]
\FunctionTok{library}\NormalTok{(tidyverse)}

\NormalTok{das\_keywords }\OtherTok{\textless{}{-}} \FunctionTok{c}\NormalTok{(}\StringTok{"data availability"}\NormalTok{, }\StringTok{"code availability"}\NormalTok{,}
                  \StringTok{"github}\SpecialCharTok{\textbackslash{}\textbackslash{}}\StringTok{.com"}\NormalTok{, }\StringTok{"zenodo"}\NormalTok{, }\StringTok{"figshare"}\NormalTok{,}
                  \StringTok{"data are available"}\NormalTok{, }\StringTok{"code is available"}\NormalTok{)}

\NormalTok{das\_pattern }\OtherTok{\textless{}{-}} \FunctionTok{paste}\NormalTok{(das\_keywords, }\AttributeTok{collapse =} \StringTok{"|"}\NormalTok{)}

\NormalTok{works }\SpecialCharTok{|\textgreater{}}
 \FunctionTok{filter}\NormalTok{(}\SpecialCharTok{!}\FunctionTok{is.na}\NormalTok{(abstract), }\FunctionTok{nzchar}\NormalTok{(}\FunctionTok{trimws}\NormalTok{(abstract))) }\SpecialCharTok{|\textgreater{}}
 \FunctionTok{mutate}\NormalTok{(}\AttributeTok{candidate\_mention =} \FunctionTok{str\_detect}\NormalTok{(}\FunctionTok{tolower}\NormalTok{(abstract), das\_pattern)) }\SpecialCharTok{|\textgreater{}}
 \FunctionTok{summarise}\NormalTok{(}
   \AttributeTok{total =} \FunctionTok{n}\NormalTok{(),}
   \AttributeTok{candidate\_mentions =} \FunctionTok{sum}\NormalTok{(candidate\_mention),}
   \AttributeTok{candidate\_rate =} \FunctionTok{mean}\NormalTok{(candidate\_mention)}
\NormalTok{ )}
\end{Highlighting}
\end{Shaded}

\emph{Key insight:} Abstract screening can miss statements and produce false positives,
so the match rate is neither a validated DAS rate nor a guaranteed lower bound.
Inspect full texts and validate precision/recall on an annotated sample.

\textbf{2. Cross-journal comparison.}
Compare metadata completeness and OA rates across journals.

\begin{Shaded}
\begin{Highlighting}[]
\NormalTok{journal\_ids }\OtherTok{\textless{}{-}} \FunctionTok{c}\NormalTok{(}\StringTok{"S148561398"}\NormalTok{, }\StringTok{"S205292342"}\NormalTok{, }\StringTok{"S4210197613"}\NormalTok{)}

\NormalTok{journal\_quality }\OtherTok{\textless{}{-}} \FunctionTok{map}\NormalTok{(journal\_ids, \textbackslash{}(jid) \{}
\NormalTok{ w }\OtherTok{\textless{}{-}} \FunctionTok{fetch\_openalex}\NormalTok{(}\AttributeTok{entity =} \StringTok{"works"}\NormalTok{,}
               \AttributeTok{primary\_location.source.id =}\NormalTok{ jid,}
               \AttributeTok{from\_publication\_date =} \StringTok{"2020{-}01{-}01"}\NormalTok{,}
               \AttributeTok{options =} \FunctionTok{list}\NormalTok{(}\AttributeTok{sample =} \DecValTok{200}\NormalTok{, }\AttributeTok{seed =} \DecValTok{42}\NormalTok{))}
 \FunctionTok{stopifnot}\NormalTok{(}\FunctionTok{is.data.frame}\NormalTok{(w), }\FunctionTok{nrow}\NormalTok{(w) }\SpecialCharTok{\textgreater{}} \DecValTok{0}\NormalTok{)}
 \FunctionTok{tibble}\NormalTok{(}
   \AttributeTok{journal =}\NormalTok{ jid,}
   \AttributeTok{abstract\_rate =} \FunctionTok{mean}\NormalTok{(}\SpecialCharTok{!}\FunctionTok{is.na}\NormalTok{(w}\SpecialCharTok{$}\NormalTok{abstract) }\SpecialCharTok{\&} \FunctionTok{nzchar}\NormalTok{(}\FunctionTok{trimws}\NormalTok{(w}\SpecialCharTok{$}\NormalTok{abstract))),}
   \AttributeTok{oa\_observed =} \FunctionTok{sum}\NormalTok{(}\SpecialCharTok{!}\FunctionTok{is.na}\NormalTok{(w}\SpecialCharTok{$}\NormalTok{oa\_status)),}
   \AttributeTok{oa\_rate =} \FunctionTok{mean}\NormalTok{(w}\SpecialCharTok{$}\NormalTok{oa\_status }\SpecialCharTok{!=} \StringTok{"closed"}\NormalTok{, }\AttributeTok{na.rm =} \ConstantTok{TRUE}\NormalTok{),}
   \AttributeTok{doi\_rate =} \FunctionTok{mean}\NormalTok{(}\SpecialCharTok{!}\FunctionTok{is.na}\NormalTok{(w}\SpecialCharTok{$}\NormalTok{doi))}
\NormalTok{ )}
\NormalTok{\}) }\SpecialCharTok{|\textgreater{}} \FunctionTok{list\_rbind}\NormalTok{()}

\NormalTok{journal\_quality}
\end{Highlighting}
\end{Shaded}

\emph{Key insight:} Access status and metadata completeness describe different
properties. Inspect actual output and coverage; no association is assumed.

\textbf{3. Temporal trends.}
Track recorded access status and abstract coverage by publication cohort.

\begin{Shaded}
\begin{Highlighting}[]
\NormalTok{metadata\_trends }\OtherTok{\textless{}{-}}\NormalTok{ works\_repro }\SpecialCharTok{|\textgreater{}}
 \FunctionTok{group\_by}\NormalTok{(year) }\SpecialCharTok{|\textgreater{}}
 \FunctionTok{summarise}\NormalTok{(}
   \AttributeTok{abstract\_rate =} \FunctionTok{mean}\NormalTok{(has\_abstract),}
   \AttributeTok{oa\_rate =} \FunctionTok{mean}\NormalTok{(oa\_status }\SpecialCharTok{!=} \StringTok{"closed"}\NormalTok{, }\AttributeTok{na.rm =} \ConstantTok{TRUE}\NormalTok{),}
   \AttributeTok{.groups =} \StringTok{"drop"}
\NormalTok{ ) }\SpecialCharTok{|\textgreater{}}
 \FunctionTok{pivot\_longer}\NormalTok{(}\SpecialCharTok{{-}}\NormalTok{year, }\AttributeTok{names\_to =} \StringTok{"metric"}\NormalTok{,}
              \AttributeTok{values\_to =} \StringTok{"value"}\NormalTok{)}

\FunctionTok{ggplot}\NormalTok{(metadata\_trends, }\FunctionTok{aes}\NormalTok{(year, value,}
                                 \AttributeTok{colour =}\NormalTok{ metric)) }\SpecialCharTok{+}
 \FunctionTok{geom\_line}\NormalTok{() }\SpecialCharTok{+}
 \FunctionTok{scale\_y\_continuous}\NormalTok{(}\AttributeTok{labels =}\NormalTok{ scales}\SpecialCharTok{::}\NormalTok{percent) }\SpecialCharTok{+}
 \FunctionTok{labs}\NormalTok{(}\AttributeTok{title =} \StringTok{"Access status and metadata coverage by cohort"}\NormalTok{,}
      \AttributeTok{x =} \StringTok{"Year"}\NormalTok{, }\AttributeTok{y =} \StringTok{"Proportion"}\NormalTok{)}
\end{Highlighting}
\end{Shaded}

\emph{Key insight:} These are current recorded properties of publication cohorts,
not adoption histories or measures of data/code sharing and reproducibility.

\section*{Chapter 34 --- Reproducible Pipelines}\label{chapter-34-reproducible-pipelines}
\addcontentsline{toc}{section}{Chapter 34 --- Reproducible Pipelines}

\textbf{1. Build a pipeline.}
Create a \texttt{\_targets.R} file for a bibliometric workflow.

\begin{Shaded}
\begin{Highlighting}[]
\FunctionTok{library}\NormalTok{(targets)}

\CommentTok{\# \_targets.R content:}
\FunctionTok{tar\_option\_set}\NormalTok{(}\AttributeTok{packages =} \FunctionTok{c}\NormalTok{(}\StringTok{"openalexR"}\NormalTok{, }\StringTok{"tidyverse"}\NormalTok{))}

\FunctionTok{list}\NormalTok{(}
 \FunctionTok{tar\_target}\NormalTok{(raw\_data, \{}
   \FunctionTok{fetch\_openalex}\NormalTok{(}\AttributeTok{entity =} \StringTok{"works"}\NormalTok{, }\AttributeTok{search =} \StringTok{"bibliometrics"}\NormalTok{,}
            \AttributeTok{from\_publication\_date =} \StringTok{"2020{-}01{-}01"}\NormalTok{)}
\NormalTok{ \}),}

 \FunctionTok{tar\_target}\NormalTok{(clean\_data, \{}
\NormalTok{   raw\_data }\SpecialCharTok{|\textgreater{}}
     \FunctionTok{filter}\NormalTok{(}\SpecialCharTok{!}\FunctionTok{is.na}\NormalTok{(doi), }\SpecialCharTok{!}\FunctionTok{is.na}\NormalTok{(publication\_year)) }\SpecialCharTok{|\textgreater{}}
     \FunctionTok{select}\NormalTok{(id, doi, display\_name, publication\_year,}
\NormalTok{            cited\_by\_count, is\_oa)}
\NormalTok{ \}),}

 \FunctionTok{tar\_target}\NormalTok{(summary\_stats, \{}
\NormalTok{   clean\_data }\SpecialCharTok{|\textgreater{}}
     \FunctionTok{group\_by}\NormalTok{(publication\_year) }\SpecialCharTok{|\textgreater{}}
     \FunctionTok{summarise}\NormalTok{(}\AttributeTok{n =} \FunctionTok{n}\NormalTok{(), }\AttributeTok{mean\_cites =} \FunctionTok{mean}\NormalTok{(cited\_by\_count))}
\NormalTok{ \}),}

 \FunctionTok{tar\_target}\NormalTok{(trend\_plot, \{}
   \FunctionTok{ggplot}\NormalTok{(summary\_stats, }\FunctionTok{aes}\NormalTok{(publication\_year, n)) }\SpecialCharTok{+}
     \FunctionTok{geom\_col}\NormalTok{() }\SpecialCharTok{+}
     \FunctionTok{labs}\NormalTok{(}\AttributeTok{title =} \StringTok{"Publication trend"}\NormalTok{, }\AttributeTok{x =} \StringTok{"Year"}\NormalTok{, }\AttributeTok{y =} \StringTok{"Works"}\NormalTok{)}
\NormalTok{ \})}
\NormalTok{)}

\CommentTok{\# Run with: tar\_make()}
\CommentTok{\# Verify: tar\_visnetwork()}
\end{Highlighting}
\end{Shaded}

\emph{Key insight:} Changing the summary function only re-runs \texttt{summary\_stats} and
\texttt{trend\_plot}, not \texttt{raw\_data} or \texttt{clean\_data} --- this is the power of
dependency-based pipelines.

\textbf{2. renv snapshot.}
Demonstrate dependency management with renv.

\begin{Shaded}
\begin{Highlighting}[]
\CommentTok{\# Step 1: Initialise renv}
\CommentTok{\# renv::init()}

\CommentTok{\# Step 2: Install a package}
\CommentTok{\# install.packages("glue")}
\CommentTok{\# library(glue)}

\CommentTok{\# Step 3: Snapshot}
\CommentTok{\# renv::snapshot()}
\CommentTok{\# This creates/updates renv.lock}

\CommentTok{\# Step 4: Simulate loss {-}{-}{-} remove the package}
\CommentTok{\# remove.packages("glue")}
\CommentTok{\# library(glue)  \# Error: package not found}

\CommentTok{\# Step 5: Restore}
\CommentTok{\# renv::restore()}
\CommentTok{\# library(glue)  \# Works again}
\end{Highlighting}
\end{Shaded}

\emph{Key insight:} renv creates a project-local library and a lockfile, ensuring
exact reproducibility of package versions across machines and time.

\textbf{3. Branching.}
Use dynamic branching to analyse multiple journals in parallel.

\begin{Shaded}
\begin{Highlighting}[]
\FunctionTok{library}\NormalTok{(targets)}

\CommentTok{\# \_targets.R with dynamic branching:}
\FunctionTok{tar\_option\_set}\NormalTok{(}\AttributeTok{packages =} \FunctionTok{c}\NormalTok{(}\StringTok{"openalexR"}\NormalTok{, }\StringTok{"tidyverse"}\NormalTok{))}

\FunctionTok{list}\NormalTok{(}
 \FunctionTok{tar\_target}\NormalTok{(journal\_ids,}
            \FunctionTok{c}\NormalTok{(}\StringTok{"S4210174107"}\NormalTok{, }\StringTok{"S137773608"}\NormalTok{, }\StringTok{"S49861241"}\NormalTok{)),}

 \FunctionTok{tar\_target}\NormalTok{(journal\_data,}
   \FunctionTok{fetch\_openalex}\NormalTok{(}\AttributeTok{entity =} \StringTok{"works"}\NormalTok{,}
            \AttributeTok{primary\_location.source.id =}\NormalTok{ journal\_ids,}
            \AttributeTok{from\_publication\_date =} \StringTok{"2020{-}01{-}01"}\NormalTok{),}
   \AttributeTok{pattern =} \FunctionTok{map}\NormalTok{(journal\_ids)}
\NormalTok{ ),}

 \FunctionTok{tar\_target}\NormalTok{(journal\_summary,}
\NormalTok{   journal\_data }\SpecialCharTok{|\textgreater{}}
     \FunctionTok{summarise}\NormalTok{(}
       \AttributeTok{journal =} \FunctionTok{first}\NormalTok{(so),}
       \AttributeTok{n =} \FunctionTok{n}\NormalTok{(),}
       \AttributeTok{mean\_cites =} \FunctionTok{mean}\NormalTok{(cited\_by\_count)}
\NormalTok{     ),}
   \AttributeTok{pattern =} \FunctionTok{map}\NormalTok{(journal\_data)}
\NormalTok{ )}
\NormalTok{)}
\end{Highlighting}
\end{Shaded}

\emph{Key insight:} Dynamic branching avoids duplicating target definitions; each
journal is processed independently and can run in parallel with
\texttt{tar\_make\_future()}.

\section*{Chapter 35 --- Interactive Dashboards}\label{chapter-35-interactive-dashboards}
\addcontentsline{toc}{section}{Chapter 35 --- Interactive Dashboards}

\textbf{1. Multi-page dashboard.}
Create a Quarto dashboard with two pages.

\begin{Shaded}
\begin{Highlighting}[]
\CommentTok{\# Save as dashboard.qmd:}
\CommentTok{\# {-}{-}{-}}
\CommentTok{\# title: "Bibliometric Dashboard"}
\CommentTok{\# format: dashboard}
\CommentTok{\# {-}{-}{-}}
\CommentTok{\#}
\CommentTok{\# \# Overview}
\CommentTok{\#}
\CommentTok{\# \#\# Row}
\CommentTok{\#}
\CommentTok{\# \textasciigrave{}\textasciigrave{}\textasciigrave{}\{r\}}
\CommentTok{\# library(plotly)}
\CommentTok{\# plot\_ly(summary\_data, x = \textasciitilde{}year, y = \textasciitilde{}n, type = "bar")}
\CommentTok{\# \textasciigrave{}\textasciigrave{}\textasciigrave{}}
\CommentTok{\#}
\CommentTok{\# \#\# Row}
\CommentTok{\#}
\CommentTok{\# \textasciigrave{}\textasciigrave{}\textasciigrave{}\{r\}}
\CommentTok{\# plot\_ly(works, x = \textasciitilde{}cited\_by\_count, type = "histogram")}
\CommentTok{\# \textasciigrave{}\textasciigrave{}\textasciigrave{}}
\CommentTok{\#}
\CommentTok{\# \# Details}
\CommentTok{\#}
\CommentTok{\# \#\# Row}
\CommentTok{\#}
\CommentTok{\# \textasciigrave{}\textasciigrave{}\textasciigrave{}\{r\}}
\CommentTok{\# library(DT)}
\CommentTok{\# datatable(works |\textgreater{} select(display\_name, publication\_year,}
\CommentTok{\#                           cited\_by\_count, is\_oa))}
\CommentTok{\# \textasciigrave{}\textasciigrave{}\textasciigrave{}}
\end{Highlighting}
\end{Shaded}

\emph{Key insight:} Quarto dashboards use page-level headings (\texttt{\#}) for navigation
tabs and row/column-level headings (\texttt{\#\#}) for layout.

\textbf{2. Linked brushing.}
Use plotly selection to filter a DT table.

\begin{Shaded}
\begin{Highlighting}[]
\FunctionTok{library}\NormalTok{(plotly)}
\FunctionTok{library}\NormalTok{(crosstalk)}
\FunctionTok{library}\NormalTok{(DT)}

\NormalTok{shared\_data }\OtherTok{\textless{}{-}}\NormalTok{ SharedData}\SpecialCharTok{$}\FunctionTok{new}\NormalTok{(works)}

\FunctionTok{bscols}\NormalTok{(}
 \FunctionTok{plot\_ly}\NormalTok{(shared\_data, }\AttributeTok{x =} \SpecialCharTok{\textasciitilde{}}\NormalTok{publication\_year,}
         \AttributeTok{y =} \SpecialCharTok{\textasciitilde{}}\NormalTok{cited\_by\_count, }\AttributeTok{type =} \StringTok{"scatter"}\NormalTok{,}
         \AttributeTok{mode =} \StringTok{"markers"}\NormalTok{) }\SpecialCharTok{|\textgreater{}}
   \FunctionTok{layout}\NormalTok{(}\AttributeTok{dragmode =} \StringTok{"select"}\NormalTok{),}

 \FunctionTok{datatable}\NormalTok{(shared\_data,}
           \AttributeTok{options =} \FunctionTok{list}\NormalTok{(}\AttributeTok{pageLength =} \DecValTok{10}\NormalTok{),}
           \AttributeTok{selection =} \StringTok{"multiple"}\NormalTok{)}
\NormalTok{)}
\end{Highlighting}
\end{Shaded}

\emph{Key insight:} \texttt{crosstalk::SharedData} enables linked brushing between any
htmlwidget that supports it, without needing Shiny.

\textbf{3. Custom tooltips.}
Add informative hover text to a plotly chart.

\begin{Shaded}
\begin{Highlighting}[]
\FunctionTok{library}\NormalTok{(plotly)}
\FunctionTok{library}\NormalTok{(tidyverse)}

\NormalTok{tooltip\_data }\OtherTok{\textless{}{-}}\NormalTok{ works }\SpecialCharTok{|\textgreater{}}
 \FunctionTok{group\_by}\NormalTok{(publication\_year) }\SpecialCharTok{|\textgreater{}}
 \FunctionTok{summarise}\NormalTok{(}\AttributeTok{n =} \FunctionTok{n}\NormalTok{(), }\AttributeTok{mean\_cites =} \FunctionTok{round}\NormalTok{(}\FunctionTok{mean}\NormalTok{(cited\_by\_count), }\DecValTok{1}\NormalTok{))}

\FunctionTok{plot\_ly}\NormalTok{(tooltip\_data,}
       \AttributeTok{x =} \SpecialCharTok{\textasciitilde{}}\NormalTok{publication\_year,}
       \AttributeTok{y =} \SpecialCharTok{\textasciitilde{}}\NormalTok{n,}
       \AttributeTok{type =} \StringTok{"bar"}\NormalTok{,}
       \AttributeTok{text =} \SpecialCharTok{\textasciitilde{}}\FunctionTok{paste0}\NormalTok{(}\StringTok{"Year: "}\NormalTok{, publication\_year,}
                      \StringTok{"\textless{}br\textgreater{}Papers: "}\NormalTok{, n,}
                      \StringTok{"\textless{}br\textgreater{}Mean cites: "}\NormalTok{, mean\_cites),}
       \AttributeTok{hoverinfo =} \StringTok{"text"}\NormalTok{)}
\end{Highlighting}
\end{Shaded}

\emph{Key insight:} Custom tooltips using \texttt{text} and \texttt{hoverinfo\ =\ "text"} provide
context-rich information on hover, improving dashboard usability.

\section*{Chapter 36 --- Parameterized Quarto Reports}\label{chapter-36-parameterized-quarto-reports}
\addcontentsline{toc}{section}{Chapter 36 --- Parameterized Quarto Reports}

\textbf{1. Institutional template.}
Create a parameterized report for institutional profiles.

\begin{Shaded}
\begin{Highlighting}[]
\CommentTok{\# Save as inst\_profile.qmd:}
\CommentTok{\# {-}{-}{-}}
\CommentTok{\# title: "Institutional Profile"}
\CommentTok{\# params:}
\CommentTok{\#   inst\_id: "I123456789"}
\CommentTok{\#   inst\_name: "University of Example"}
\CommentTok{\#   start\_year: 2020}
\CommentTok{\#   end\_year: 2023}
\CommentTok{\# {-}{-}{-}}
\CommentTok{\#}
\CommentTok{\# \textasciigrave{}\textasciigrave{}\textasciigrave{}\{r\}}
\CommentTok{\# library(openalexR)}
\CommentTok{\# library(tidyverse)}
\CommentTok{\#}
\CommentTok{\# works \textless{}{-} fetch\_openalex(}
\CommentTok{\#   entity = "works",}
\CommentTok{\#   authorships.institutions.lineage = params$inst\_id,}
\CommentTok{\#   from\_publication\_date = paste0(params$start\_year, "{-}01{-}01"),}
\CommentTok{\#   to\_publication\_date = paste0(params$end\_year, "{-}12{-}31")}
\CommentTok{\# )}
\CommentTok{\#}
\CommentTok{\# cat(sprintf("Total publications: \%d\textbackslash{}n", nrow(works)))}
\CommentTok{\# \textasciigrave{}\textasciigrave{}\textasciigrave{}}
\CommentTok{\#}
\CommentTok{\# \#\# Publication Trend}
\CommentTok{\# \textasciigrave{}\textasciigrave{}\textasciigrave{}\{r\}}
\CommentTok{\# works |\textgreater{}}
\CommentTok{\#   count(publication\_year) |\textgreater{}}
\CommentTok{\#   ggplot(aes(publication\_year, n)) +}
\CommentTok{\#   geom\_col() +}
\CommentTok{\#   labs(title = params$inst\_name)}
\CommentTok{\# \textasciigrave{}\textasciigrave{}\textasciigrave{}}
\CommentTok{\#}
\CommentTok{\# \#\# Citation Distribution}
\CommentTok{\# \textasciigrave{}\textasciigrave{}\textasciigrave{}\{r\}}
\CommentTok{\# ggplot(works, aes(cited\_by\_count)) +}
\CommentTok{\#   geom\_histogram(bins = 30) +}
\CommentTok{\#   scale\_x\_log10()}
\CommentTok{\# \textasciigrave{}\textasciigrave{}\textasciigrave{}}
\CommentTok{\#}
\CommentTok{\# \#\# Top Authors}
\CommentTok{\# \textasciigrave{}\textasciigrave{}\textasciigrave{}\{r\}}
\CommentTok{\# works |\textgreater{} unnest(author) |\textgreater{}}
\CommentTok{\#   count(au\_display\_name, sort = TRUE) |\textgreater{}}
\CommentTok{\#   head(10)}
\CommentTok{\# \textasciigrave{}\textasciigrave{}\textasciigrave{}}
\end{Highlighting}
\end{Shaded}

\emph{Key insight:} Parameters make the same template reusable for any institution;
batch rendering generates a portfolio of reports efficiently.

\textbf{2. Batch rendering.}
Render the template for multiple journals.

\begin{Shaded}
\begin{Highlighting}[]
\FunctionTok{library}\NormalTok{(quarto)}

\NormalTok{journals }\OtherTok{\textless{}{-}} \FunctionTok{list}\NormalTok{(}
 \FunctionTok{list}\NormalTok{(}\AttributeTok{journal\_id =} \StringTok{"S4210174107"}\NormalTok{, }\AttributeTok{journal\_name =} \StringTok{"Scientometrics"}\NormalTok{),}
 \FunctionTok{list}\NormalTok{(}\AttributeTok{journal\_id =} \StringTok{"S137773608"}\NormalTok{, }\AttributeTok{journal\_name =} \StringTok{"Nature Physics"}\NormalTok{),}
 \FunctionTok{list}\NormalTok{(}\AttributeTok{journal\_id =} \StringTok{"S49861241"}\NormalTok{, }\AttributeTok{journal\_name =} \StringTok{"PLOS ONE"}\NormalTok{)}
\NormalTok{)}

\FunctionTok{walk}\NormalTok{(journals, \textbackslash{}(j) \{}
 \FunctionTok{quarto\_render}\NormalTok{(}
   \StringTok{"journal\_profile.qmd"}\NormalTok{,}
   \AttributeTok{execute\_params =} \FunctionTok{list}\NormalTok{(}
     \AttributeTok{journal\_id =}\NormalTok{ j}\SpecialCharTok{$}\NormalTok{journal\_id,}
     \AttributeTok{journal\_name =}\NormalTok{ j}\SpecialCharTok{$}\NormalTok{journal\_name,}
     \AttributeTok{start\_year =} \DecValTok{2020}\NormalTok{,}
     \AttributeTok{end\_year =} \DecValTok{2023}
\NormalTok{   ),}
   \AttributeTok{output\_file =} \FunctionTok{paste0}\NormalTok{(}\StringTok{"report\_"}\NormalTok{, j}\SpecialCharTok{$}\NormalTok{journal\_name, }\StringTok{".html"}\NormalTok{)}
\NormalTok{ )}
\NormalTok{\})}
\end{Highlighting}
\end{Shaded}

\emph{Key insight:} Batch rendering with \texttt{quarto\_render()} and parameterised
templates scales to hundreds of reports with minimal code.

\textbf{3. Conditional content.}
Add a warning when the sample is too small.

\begin{Shaded}
\begin{Highlighting}[]
\CommentTok{\# In the parameterized .qmd file:}
\CommentTok{\#}
\CommentTok{\# \textasciigrave{}\textasciigrave{}\textasciigrave{}\{r\}}
\CommentTok{\# \#| results: asis}
\CommentTok{\# if (nrow(works) \textless{} 50) \{}
\CommentTok{\#   cat("\textless{}div class="callout callout{-}warn"\textgreater{}\textbackslash{}n")}
\CommentTok{\#   cat("This sample contains fewer than 50 papers. ")}
\CommentTok{\#   cat("The citation distribution may not be reliable.\textbackslash{}n")}
\CommentTok{\#   cat(":::\textbackslash{}n")}
\CommentTok{\# \} else \{}
\CommentTok{\#   ggplot(works, aes(cited\_by\_count)) +}
\CommentTok{\#     geom\_histogram(bins = 30) +}
\CommentTok{\#     scale\_x\_log10() +}
\CommentTok{\#     labs(title = "Citation distribution")}
\CommentTok{\# \}}
\CommentTok{\# \textasciigrave{}\textasciigrave{}\textasciigrave{}}
\end{Highlighting}
\end{Shaded}

\emph{Key insight:} Conditional logic with \texttt{results:\ asis} enables dynamic content
switching --- essential for robust automated reports that handle edge cases.

\section*{Chapter 37 --- Building a Bibliometric Shiny App}\label{chapter-37-building-a-bibliometric-shiny-app}
\addcontentsline{toc}{section}{Chapter 37 --- Building a Bibliometric Shiny App}

\textbf{1. Add network tab.}
Extend the app with a co-authorship network visualisation.

\begin{Shaded}
\begin{Highlighting}[]
\FunctionTok{library}\NormalTok{(shiny)}
\FunctionTok{library}\NormalTok{(visNetwork)}
\FunctionTok{library}\NormalTok{(igraph)}
\FunctionTok{library}\NormalTok{(tidyverse)}

\CommentTok{\# Add to UI:}
\CommentTok{\# tabPanel("Network",}
\CommentTok{\#   visNetworkOutput("coauth\_network", height = "600px")}
\CommentTok{\# )}

\CommentTok{\# Add to server:}
\CommentTok{\# output$coauth\_network \textless{}{-} renderVisNetwork(\{}
\CommentTok{\#   req(data())}
\CommentTok{\#   edges \textless{}{-} data() |\textgreater{}}
\CommentTok{\#     unnest(author) |\textgreater{}}
\CommentTok{\#     group\_by(id) |\textgreater{}}
\CommentTok{\#     filter(n() \textgreater{}= 2, n() \textless{}= 10) |\textgreater{}}
\CommentTok{\#     reframe(pair = combn(au\_display\_name, 2, simplify = FALSE)) |\textgreater{}}
\CommentTok{\#     unnest\_wider(pair, names\_sep = "\_") |\textgreater{}}
\CommentTok{\#     count(pair\_1, pair\_2, name = "weight") |\textgreater{}}
\CommentTok{\#     rename(from = pair\_1, to = pair\_2)}
\CommentTok{\#}
\CommentTok{\#   g \textless{}{-} graph\_from\_data\_frame(edges, directed = FALSE)}
\CommentTok{\#   comm \textless{}{-} cluster\_leiden(g, resolution = 1.0)}
\CommentTok{\#}
\CommentTok{\#   nodes \textless{}{-} tibble(}
\CommentTok{\#     id = V(g)$name, label = V(g)$name,}
\CommentTok{\#     group = as.character(membership(comm))}
\CommentTok{\#   )}
\CommentTok{\#}
\CommentTok{\#   visNetwork(nodes, edges) |\textgreater{}}
\CommentTok{\#     visOptions(highlightNearest = TRUE, selectedBy = "group")}
\CommentTok{\# \})}
\end{Highlighting}
\end{Shaded}

\emph{Key insight:} Limiting the network to papers with 2-10 authors avoids
overwhelming the visualisation with hyper-authored papers.

\textbf{2. Caching.}
Implement server-side caching for repeated queries.

\begin{Shaded}
\begin{Highlighting}[]
\FunctionTok{library}\NormalTok{(shiny)}
\FunctionTok{library}\NormalTok{(memoise)}

\CommentTok{\# Create a cached version of the data fetching function}
\NormalTok{fetch\_cached }\OtherTok{\textless{}{-}} \FunctionTok{memoise}\NormalTok{(}\ControlFlowTok{function}\NormalTok{(journal\_id, start\_date, end\_date) \{}
 \FunctionTok{fetch\_openalex}\NormalTok{(}
   \AttributeTok{entity =} \StringTok{"works"}\NormalTok{,}
   \AttributeTok{primary\_location.source.id =}\NormalTok{ journal\_id,}
   \AttributeTok{from\_publication\_date =}\NormalTok{ start\_date,}
   \AttributeTok{to\_publication\_date =}\NormalTok{ end\_date}
\NormalTok{ )}
\NormalTok{\}, }\AttributeTok{cache =}\NormalTok{ cachem}\SpecialCharTok{::}\FunctionTok{cache\_disk}\NormalTok{(}\StringTok{"shiny\_cache"}\NormalTok{))}

\CommentTok{\# In the server:}
\CommentTok{\# data \textless{}{-} reactive(\{}
\CommentTok{\#   fetch\_cached(}
\CommentTok{\#     input$journal\_id,}
\CommentTok{\#     paste0(input$start\_year, "{-}01{-}01"),}
\CommentTok{\#     paste0(input$end\_year, "{-}12{-}31")}
\CommentTok{\#   )}
\CommentTok{\# \})}
\end{Highlighting}
\end{Shaded}

\emph{Key insight:} Disk-backed caching persists across sessions, so the same query
never hits the API twice. Use \texttt{cachem::cache\_disk()} for production apps.

\textbf{3. Deployment.}
Deploy to shinyapps.io.

\begin{Shaded}
\begin{Highlighting}[]
\CommentTok{\# Step 1: Install rsconnect}
\CommentTok{\# install.packages("rsconnect")}

\CommentTok{\# Step 2: Configure account}
\CommentTok{\# rsconnect::setAccountInfo(}
\CommentTok{\#   name = "your{-}account",}
\CommentTok{\#   token = "your{-}token",}
\CommentTok{\#   secret = "your{-}secret"}
\CommentTok{\# )}

\CommentTok{\# Step 3: Deploy}
\CommentTok{\# rsconnect::deployApp(}
\CommentTok{\#   appDir = ".",}
\CommentTok{\#   appName = "bibliometric{-}explorer",}
\CommentTok{\#   appTitle = "Bibliometric Explorer"}
\CommentTok{\# )}

\CommentTok{\# Step 4: Test with different inputs}
\CommentTok{\# {-} Try different journal IDs}
\CommentTok{\# {-} Test with narrow and wide date ranges}
\CommentTok{\# {-} Verify that caching works (second load should be instant)}
\end{Highlighting}
\end{Shaded}

\emph{Key insight:} Test with a variety of inputs before deploying; edge cases (empty
results, very large corpora) should be handled gracefully with validation
messages.

\textbf{4. biblioshiny comparison.}
Run biblioshiny and compare features.

\begin{Shaded}
\begin{Highlighting}[]
\FunctionTok{library}\NormalTok{(bibliometrix)}
\CommentTok{\# biblioshiny()}

\CommentTok{\# Feature comparison table:}
\CommentTok{\# | Feature              | biblioshiny | Custom app |}
\CommentTok{\# |{-}{-}{-}{-}{-}{-}{-}{-}{-}{-}{-}{-}{-}{-}{-}{-}{-}{-}{-}{-}{-}|{-}{-}{-}{-}{-}{-}{-}{-}{-}{-}{-}{-}{-}|{-}{-}{-}{-}{-}{-}{-}{-}{-}{-}{-}{-}|}
\CommentTok{\# | Data sources        | WoS, Scopus | OpenAlex   |}
\CommentTok{\# | Network analysis    | Yes         | Basic      |}
\CommentTok{\# | Topic analysis      | Yes         | No         |}
\CommentTok{\# | Custom indicators   | Limited     | Full       |}
\CommentTok{\# | Parameterization    | No          | Yes        |}
\CommentTok{\# | Deployment          | Local only  | shinyapps  |}
\CommentTok{\# | Reproducibility     | GUI{-}based   | Code{-}based |}
\CommentTok{\#}
\CommentTok{\# biblioshiny is comprehensive but not easily customisable;}
\CommentTok{\# a custom app trades breadth for tailored functionality.}
\end{Highlighting}
\end{Shaded}

\emph{Key insight:} biblioshiny is excellent for exploration but lacks
reproducibility and customisation. A custom app is better for production
workflows with specific analytical requirements.

\chapter*{Glossary}\label{glossary}
\addcontentsline{toc}{chapter}{Glossary}

\emph{R. Heller · Date: 2026-05-09 · Last changed: 2026-05-11}

\textbf{Altmetrics.} Metrics derived from online activity (social media mentions, news coverage, policy documents) rather than citations. See Chapter 26.

\textbf{Bibliographic coupling.} Two documents are bibliographically coupled when they share one or more references. Strength increases with the number of shared references. See Chapter 16.

\textbf{Bibliometrics.} The quantitative study of scholarly publications, their production, dissemination, and impact. A subfield of scientometrics focused on documents.

\textbf{Burst detection.} A method for identifying terms or topics whose frequency increases sharply in a given time window, signalling emerging research fronts. See Chapter 25.

\textbf{Citation aging.} The temporal pattern of how citations accumulate over time after publication. The citation half-life measures how quickly a paper accrues half its total citations. See Chapter 14.

\textbf{Citation window.} The time period during which citations to a publication are counted. Common choices: 3-year, 5-year, or variable windows.

\textbf{Co-authorship network.} A graph where nodes represent authors and edges connect authors who have co-authored at least one publication. See Chapter 15.

\textbf{Co-citation.} Two documents are co-cited when they are both cited by a third document. Frequently co-cited documents are considered intellectually related. See Chapter 16.

\textbf{Co-word analysis.} Network analysis of the co-occurrence of terms (keywords, title words, or descriptors) within documents. See Chapter 17.

\textbf{Community detection.} Algorithmic identification of densely connected subgroups within a network. Common methods: Leiden, Louvain, modularity optimisation. See Chapter 18.

\textbf{Disparity filter.} A network backbone extraction method that retains statistically significant edges based on local weight distributions. See Chapter 18.

\textbf{DORA.} The San Francisco Declaration on Research Assessment. A set of recommendations for improving research evaluation practices, emphasising that journal-based metrics should not be used as surrogates for individual article quality.

\textbf{Field normalisation.} Adjusting citation counts to account for differences in citation practices across research fields and publication years. See Chapter 13.

\textbf{Fractional counting.} An attribution method where credit for a multi-authored publication is divided equally among contributors (1/n per author or institution), in contrast to full counting. See Chapter 13.

\textbf{Full counting.} An attribution method where each author or institution on a publication receives full credit (count of 1). See Chapter 13.

\textbf{h-index.} A researcher has an h-index of \emph{h} if \emph{h} of their papers have each been cited at least \emph{h} times. See Chapter 10.

\textbf{Impact factor (JIF).} The mean number of citations received in year \emph{t} by articles published in a journal in years \emph{t}−1 and \emph{t}−2. See Chapter 11.

\textbf{Leiden Manifesto.} Ten principles for responsible research metrics, emphasising that quantitative evaluation should support qualitative, expert assessment rather than replace it. See Chapter 3.

\textbf{Leiden algorithm.} A community detection method that improves on Louvain by guaranteeing well-connected communities. See Chapter 18.

\textbf{MNCS (Mean Normalised Citation Score).} The average of field-normalised citation scores across a set of publications. An MNCS above 1.0 indicates above-world-average impact. See Chapter 13.

\textbf{Network backbone.} A reduced network retaining only statistically significant or structurally important edges, obtained by filtering methods such as the disparity filter. See Chapter 18.

\textbf{OpenAlex.} A free, open catalogue of the global research system, covering works, authors, institutions, sources, concepts, funders, and publishers. The default data source for this book. See Chapter 5.

\textbf{Overlay map.} A visualisation technique where a subset of data (e.g., one institution's publications) is projected onto a pre-computed base map of science. See Chapter 19.

\textbf{PP(top 10\%).} The proportion of an entity's publications that rank in the top 10\% most cited in their field and year. A percentile-based impact indicator. See Chapter 10.

\textbf{Price index.} The proportion of a paper's references that are no older than 5 years, measuring how current the cited literature is. See Chapter 14.

\textbf{RAKE.} Rapid Automatic Keyword Extraction. An unsupervised method that extracts multi-word keywords based on word co-occurrence within phrases delimited by stopwords. See Chapter 23.

\textbf{Research front.} A cluster of recently published, actively cited papers representing the current edge of a research area.

\textbf{Retraction.} The formal withdrawal of a published paper by its authors or publisher due to errors, misconduct, or irreproducibility. See Chapter 29.

\textbf{Scientometrics.} The quantitative study of science and innovation, encompassing bibliometrics, patent analysis, and the measurement of research systems.

\textbf{STM (Structural Topic Model).} A topic model that allows document-level covariates (e.g., publication year, journal) to affect topic prevalence and content. See Chapter 22.

\textbf{TF-IDF.} Term Frequency--Inverse Document Frequency. A weighting scheme that highlights terms distinctive to individual documents within a corpus. See Chapter 21.

\textbf{Topic modelling.} Unsupervised methods for discovering latent thematic structure in collections of text. Common approaches: LDA, STM, BERTopic. See Chapter 22.

\textbf{UMAP.} Uniform Manifold Approximation and Projection. A dimensionality-reduction technique for visualising high-dimensional data (e.g., document embeddings) in 2D. See Chapter 24.

\chapter*{Colophon}\label{colophon}
\addcontentsline{toc}{chapter}{Colophon}

\emph{R. Heller · Date: 2026-05-11 · Last changed: 2026-06-10}

\begin{verbatim}
#> Built with:
\end{verbatim}

\begin{verbatim}
#>   bookdown:  0.48
\end{verbatim}

\begin{verbatim}
#>   R:  R version 4.6.1 (2026-06-24)
\end{verbatim}

\begin{verbatim}
#>   Build date:  2026-11-10
\end{verbatim}

\begin{verbatim}
#>   Commit:  3096271
\end{verbatim}

\begin{verbatim}
#> 
#> Attached packages:
\end{verbatim}

\begin{verbatim}
#> - Session info ---------------------------------------------------------------
#>  setting  value
#>  version  R version 4.6.1 (2026-06-24)
#>  os       Ubuntu 24.04.5 LTS
#>  system   x86_64, linux-gnu
#>  ui       X11
#>  language (EN)
#>  collate  C.UTF-8
#>  ctype    C.UTF-8
#>  tz       UTC
#>  date     2026-11-10
#>  pandoc   3.8.3 @ /opt/hostedtoolcache/pandoc/3.8.3/x64/ (via rmarkdown)
#>  quarto   NA
#> 
#> - Packages -------------------------------------------------------------------
#>  package            * version date (UTC) lib source
#>  arrow              * 26.0.0  2026-10-24 [1] RSPM (R 4.6.0)
#>  bib2df             * 1.1.2.0 2024-05-16 [1] RSPM (R 4.6.0)
#>  bibliometrix       * 5.5.0   2026-08-29 [1] RSPM (R 4.6.0)
#>  dplyr              * 1.2.1   2026-04-03 [1] RSPM (R 4.6.0)
#>  DT                 * 0.34.0  2025-09-02 [1] RSPM (R 4.6.0)
#>  forcats            * 1.0.1   2025-09-25 [1] RSPM (R 4.6.0)
#>  ggplot2            * 4.0.3   2026-04-22 [1] RSPM (R 4.6.0)
#>  ggraph             * 2.2.2   2025-08-24 [1] RSPM (R 4.6.0)
#>  glue               * 1.8.1   2026-04-17 [1] RSPM (R 4.6.0)
#>  gt                 * 1.4.0   2026-10-15 [1] RSPM (R 4.6.0)
#>  igraph             * 2.3.4   2026-09-30 [1] RSPM (R 4.6.0)
#>  lubridate          * 1.9.5   2026-02-04 [1] RSPM (R 4.6.0)
#>  Matrix             * 1.7-5   2026-03-21 [2] CRAN (R 4.6.1)
#>  openalexR          * 3.1.2   2026-10-14 [1] RSPM (R 4.6.0)
#>  pdftools           * 3.9.1   2026-09-01 [1] RSPM (R 4.6.0)
#>  plotly             * 4.12.1  2026-07-22 [1] RSPM (R 4.6.0)
#>  purrr              * 1.2.2   2026-04-10 [1] RSPM (R 4.6.0)
#>  quanteda           * 4.5.1   2026-10-21 [1] RSPM (R 4.6.0)
#>  quanteda.textstats * 0.97.2  2024-09-03 [1] RSPM (R 4.6.0)
#>  rcrossref          * 1.2.1   2025-08-20 [1] RSPM (R 4.6.0)
#>  readr              * 2.2.0   2026-02-19 [1] RSPM (R 4.6.0)
#>  RefManageR         * 1.4.0   2022-09-30 [1] RSPM (R 4.6.0)
#>  shiny              * 1.14.0  2026-06-21 [1] RSPM (R 4.6.0)
#>  stm                * 1.3.8   2025-09-03 [1] RSPM (R 4.6.0)
#>  stringr            * 1.6.0   2025-11-04 [1] RSPM (R 4.6.0)
#>  tibble             * 3.3.1   2026-01-11 [1] RSPM (R 4.6.0)
#>  tidygraph          * 1.3.1   2024-01-30 [1] RSPM (R 4.6.0)
#>  tidyr              * 1.3.2   2025-12-19 [1] RSPM (R 4.6.0)
#>  tidytext           * 0.4.3   2025-07-25 [1] RSPM (R 4.6.0)
#>  tidyverse          * 2.0.0   2023-02-22 [1] RSPM (R 4.6.0)
#>  topicmodels        * 0.2-17  2024-08-14 [1] RSPM (R 4.6.0)
#>  uwot               * 0.2.5   2026-08-29 [1] RSPM (R 4.6.0)
#>  visNetwork         * 2.1.4   2025-09-04 [1] RSPM (R 4.6.0)
#>  word2vec           * 0.4.1   2025-11-27 [1] RSPM (R 4.6.0)
#> 
#>  [1] /home/runner/work/_temp/Library
#>  [2] /opt/R/4.6.1/lib/R/library
#>  * -- Packages attached to the search path.
#> 
#> ------------------------------------------------------------------------------
\end{verbatim}

LLM assistance (KISSKI's Chat AI service, plus a self-hosted Mistral
Small 24B via Ollama/\texttt{ollamar}) was used for copyediting, formatting,
and scaffolding tasks only --- see Acknowledgments.

\chapter*{Citing this Guide}\label{citing-this-guide-1}
\addcontentsline{toc}{chapter}{Citing this Guide}

\emph{R. Heller · Date: 2026-05-11 · Last changed: 2026-11-09}

The suggested citation is:

\begin{boxempty}
\textbf{Heller, R.} (2026). \emph{Scientometrics in R: Bibliometrics,
Network Science, and Research Evaluation} (Version 0.1.0).
Self-published via GitHub Pages.
\url{https://cttir.github.io/scientometrics-in-r/}.
\end{boxempty}

Download the reference as \href{citation-files/citation.bib}{BibTeX} or \href{citation-files/citation.ris}{.ris}.

\chapter*{References}\label{references}
\addcontentsline{toc}{chapter}{References}

\emph{R. Heller · Date: 2026-05-11 · Last changed: 2026-05-11}

\bibliography{references.bib,packages.bib}

\end{document}
